% 本文件由 LaTeX 排版 Agent 根据有效 translation.json 自动生成。
% author=John Cassian (c. 360-c. 435); work_id=3508; title=话谈

\chapter{话谈}

% structure: p0001 source=h1 target=omitted reason=page-title-already-rendered-by-parent

第1篇 梅瑟院长的第一篇话谈。\newline{}第2篇 梅瑟院长的第二篇话谈。\newline{}第3篇 巴夫努提院长的话谈——论三种弃绝。\newline{}第4篇 达尼尔院长的话谈——论肉身与心神的情欲。\newline{}第5篇 塞拉庇翁院长的话谈——论八种主要恶习。\newline{}第6篇 德奥多罗院长的话谈——论圣人之死。\newline{}第7篇 塞雷努斯院长的第一篇话谈——论心神无常与属神的邪恶。\newline{}第8篇 塞雷努斯院长的第二篇话谈——论率领者。\newline{}第9篇 依撒格院长的第一篇话谈——论祈祷。\newline{}第10篇 依撒格院长的第二篇话谈——论祈祷。\newline{}第11篇 凯雷蒙院长的第一篇话谈——论成全。\newline{}第13篇 凯雷蒙院长的第三篇话谈——论天主的保护。\newline{}第14篇 内斯特罗斯院长的第一篇话谈——论属神的知识。\newline{}第15篇 内斯特罗斯院长的第二篇话谈——论天主的恩赐。\newline{}第16篇 若瑟院长的第一篇话谈——论友谊。\newline{}第17篇 若瑟院长的第二篇话谈——论许诺。\newline{}第18篇 皮亚孟院长的话谈——论三种隐修士。\newline{}第19篇 若望院长的话谈——论修院团体与独修者的目标。\newline{}第20篇 皮努斐院长的话谈——论补赎的终结与赔补的标记。\newline{}第21篇 特奥纳斯院长的第一篇话谈——论五十天内的宽弛。\newline{}第23篇 特奥纳斯院长的第三篇话谈——论无罪。\newline{}第24篇 亚巴郎院长的话谈——论克苦。\par

\SourceRecord{Translated by C.S. Gibson.；Nicene and Post-Nicene Fathers, Second Series；Vol. 11.；Edited by Philip Schaff and Henry Wace.；Buffalo, NY: Christian Literature Publishing Co.,；1894.；<http://www.newadvent.org/fathers/3508.htm>.}

% structure: p0001 source=h1 target=section reason=page-title

\section{第一次会谈}

\Emph{\Person{圣若望·卡西安}}\par

院长\Person{摩西}第一次会谈。\par

在献给真福教宗\Person{卡斯托尔}的前言中，我曾许诺以天主助佑写成十二卷《隐修院规》并论八种主要情欲的治疗；如今，凭我微弱的能力，这义务已尽。我自然希望了解，您和诸位兄弟的判断对这著作是否公允，是否在如此深邃崇高、且从未被论述过的题目上，我们得出了值得你们关注并令所有圣兄弟热切渴望的成果。然而，鉴于上述主教已经离我们而去，归于基督，与此同时，这些最伟大教父——即居住在赛特旷野的独修士们——的十次会谈，他因对圣德无可比拟的热忱，曾命我以同样风格为他撰写（并未顾念他出于深厚情谊，将何等重担搁在难以承受的弱肩上），我认为应将它们特别奉献给您，蒙福的教宗\Person{良提乌斯}和圣兄弟\Person{赫拉迪乌斯}。因为你们中一位曾与我提及之人因弟兄之情、司铎之位，以及更重要的，因神圣研究的热情而与他联合，故对兄弟所欠之债享有继承权；另一位则敢于追随独修士的高尚习俗，不像某些人那样擅自凭借己见，而是受\OriginalTerm{圣神}{Holy Ghost}的感动，抓住正确的教义之路，几乎是未受教而知，并选择不依自己想法，而依他们的传统而学。正因如此，当我如锚泊于静默之港时，一片广阔海洋在我面前展开，使我不得不冒险将一些伟大之人的制度和教导传给后世。因为我的微薄能力之小船将在深海上经历更漫长航程的风险，其程度恰如\OriginalTerm{独修士}{Anchorite}的生活比\OriginalTerm{修院团体}{Cœnobium}更为崇高，而那些无可估量之人所全心投入的\OriginalTerm{默观天主}{contemplation of God}比普通实践生活更为超越。因此，你们当以虔诚祈祷协助我们的努力，惟恐如此神圣的主题——虽以未经尝试却忠实的方式论述——由我们而受损，或我们的淳朴在主题深处迷失。为此，让我们从以前各书中所论述的、关于隐修士外在可见的生活模式，转到内在之人、隐蔽不显的生活；从\OriginalTerm{日课制度}{canonical prayers}，让我们的谈论上升到宗徒所训导的\OriginalTerm{不断祈祷}{unceasing prayer}的持久性，使凡通过阅读我们此前著作，已借驱逐肉性错误而灵性获得\Person{雅各伯}之名者，如今因接受并非出自我们而是教父们的院规，能以纯净的洞察登上以色列的功德和（可以这样说）尊荣，并以同样方式学习在这些完美的高峰上应持守什么。因此，愿你们的祈祷从那位曾使我们堪当看见他们、向他们学习并与他们同住者获得恩赐，使我们完全记起他们的教训，并有口才予以讲述，以便我们能如从他们那里领受时一样美善而精确地解释，并成功将他们本人如同体现在他们自己的院规中、且更重要的，用\OriginalTerm{拉丁语}{Latin tongue}讲话一般，呈现在你们面前。不过，我们尤其希望告知这些会谈以及我们早期作品的读者，若其中有任何事物，因读者自身状况、职业性质，或因习俗及共同生活方式，而似乎不可能或极其困难，则不应以自身能力限度，而应以说话者的价值与完美来衡量；他们先要考察说话者的热忱与目的，因他们真正死于世俗生活，不受肉身亲属情感及世俗事务束缚。其次，当牢记他们所居之地的特性，即他们生活在广袤旷野中，与一切世人隔绝，因而得以启迪心智，默观并说出那些对未入门、未受教者，因其生活方式及习惯平庸而以为不可能的事。但若有人欲对此事给出真正判断，并渴望检验是否能达到如此完美，则让他先努力效法其目的，以同样热忱和同样生活方式，最终他将会发现那些曾以为超出人之能力的事，不仅可能，而且确实令人喜乐。但现在，让我们直接进入他们的会谈和院规。\par

% structure: p0005 source=h2 target=subsection reason=relative-heading-rank; base=h2

\subsection{第一章}

\Emph{关于我们在\Place{赛特}的停留，以及我们向院长\Person{摩西}提出的请求。}\par

当我在斯凯特沙漠时，那里有最卓越的修道之父，一切成全得以兴盛；我与神圣的日耳曼神父在一起（他从我们灵修服务的最初日子起，无论是在会院还是在沙漠中，都是我最亲密的同伴，以至为了显示我们友谊和目标之和谐，人人都会说我们两人身上只有一个心灵和灵魂）。我急于通过他的教导得以扎根，便寻访了院长梅瑟，他在这些灿烂的花朵中不仅因实践之成全，也因默观之成全而卓越；我们一起恳求他为我们发表一篇训诲，以使我们受益；我们并非没有流泪，因为我们深知他决不轻易同意向人开启成全之门，除非那人心怀忠信渴望它，并以全心的痛悔寻求它；因为他恐怕若随意向那些漠不关心或半心半意的人展示它，将那些必要且只应向寻求成全者揭示的事，向不配的人和以轻蔑接受的人揭露，他可能显得自夸，或犯下背叛托付之罪；最后，他被我们的祈祷所感动，便如此开始。\par

% structure: p0008 source=h2 target=subsection reason=relative-heading-rank; base=h2

\subsection{第二章}

\Emph{院长梅瑟的问题：修士的目标和目的是什么。}\par

他说：\Quote{所有技艺和学问都有其目标或标记；各自的目的或终结，勤奋追求每种技艺的人注视着它，从而愉快而平静地忍受各种辛劳、危险和损失。例如农夫，既不躲避一时的炎热骄阳，也不躲避霜冻寒冷，不知疲倦地犁地，一次又一次地用犁铧翻松田里的土块，同时他心中持守着他的目标；即通过勤勉的劳作将土地打碎如细沙，清除所有荆棘，除去所有杂草，因为他相信只有这样他才能达到终极目的，即获得好收成和丰收；他或可无忧地以此为生，或可增加他的财富。又如，当他的谷仓满满实实时，他乐意将其清空，以不停的劳作将种子播入松散的犁沟，为了将来的收成而不计较当前贮存的减少。那些从事商业的人也不惧怕海洋的不确定与风险，不畏惧危险，热切的希望驱使他们奔向获利的目标。此外，那些渴望军旅生活荣誉的人，当他们期待荣誉和权力的目标时，在漂泊中不理会危险和毁灭，也不被当前的损失和战争压垮，因为他们急切地想要获得某种向他们提供的荣誉。我们的修会也有其自身的目标和目的，为此我们甘愿忍受各种辛劳，不仅不厌倦，反而以之悦乐；因此，禁食时的食物匮乏对我们不是考验，守夜的疲倦成为乐趣；阅读并不断默想圣经不会使我们厌倦；此外，不断的劳作、克己、舍弃一切，以及这广大沙漠的恐怖，也吓不倒我们。无疑，正是为此，你们自己藐视亲人之爱，蔑视祖国和世俗的享乐，穿越如此众多的地区，好来到我们这些生活在沙漠中这可怜境况的粗陋之人这里。因此，他说，请回答我：是什么目标和目的激励你们如此愉快地忍受这一切？}\par

% structure: p0011 source=h2 target=subsection reason=relative-heading-rank; base=h2

\subsection{第三章}

\Emph{我们的答复。}\par

当他坚持要我们就此问题发表看法时，我们回答说：\Quote{我们为天国而忍受这一切。}\par

% structure: p0014 source=h2 target=subsection reason=relative-heading-rank; base=h2

\subsection{第四章}

\Emph{院长梅瑟对上述陈述的询问。}\par

他回答说：好，你聪明地讨论了（终极）目标。但我们应不断贴近、藉以获得那目标的具体目标或标靶，你首先应当知道。当我们坦白承认无知时，他接着说：正如我所说，在一切技艺与科学中，首要之事是有一个目标，即为心设定一个标靶，并有一个恒常的心意指向它；因为若人不能以全副勤勉与恒心将此目标摆在眼前，他绝不可能达到他所渴望的终极目的和利益。如我所说，农夫的目标是过无忧而富足的生活，在庄稼生长时，他的直接目的和目标便是：使田地清除一切荆棘与杂草；他不幻想不先通过某种劳作和希望的规划确保他所追求之物，而能获得财富与平安的结局。商人也并未搁置采购货物的欲望，藉此他可以更有利地聚敛财富；因为若他不选择通往财富的道路，他的追求便毫无意义。那些渴望获得世俗荣誉的人，首先决心自己必须投身于何种职责与境况，以便在希望的常规进程中成功获取他们所渴望的荣誉。因此，我们生活方式的终极目的确实是天主的国。但你们必须认真询问那直接目标是什么；因为若我们不以同样方式发现它，我们的奋斗与操劳将归于徒然；因为一个误入歧途的人，只会承受旅途的一切辛苦，却得不到丝毫益处。当我们对此话目瞪口呆时，老人接着说：我们修道的终极目的，如我所说，是天主的国或天国；但直接目标或标靶，是心地纯洁，没有它无人能获得那终极目的。因此，让我们将目光稳固地定在这目标上，如同定在一个明确的标靶上，尽可能笔直地导向它；若我们的思绪稍有偏离，让我们重新注视它，并如同藉着一个可靠的准绳，精确地加以纠正；这准绳将不断使我们一切努力回归到这唯一标靶，并立刻显示我们的心是否偏离了既定的方向。\par

% structure: p0017 source=h2 target=subsection reason=relative-heading-rank; base=h2

\subsection{第五章}

\Emph{试射标靶者的比喻。}\par

如同那些使用兵器的人，当他们要在今世君王面前展示其技艺时，试图将箭或标枪射向绘有奖赏的小靶；因为他们知道，除非藉着瞄准的准线，他们无法获得所期望的终结与奖赏；只有当他们能射中摆在面前的标靶时，才能享受这奖赏；但若靶子从他们视线中消失，尽管他们因技艺不精而徒然偏离正道，却无法察觉已偏离了预定直线方向，因为没有明确的标靶来证明他们瞄准的巧妙，或显示其拙劣。因此，当他们将矢弹白白射向虚空时，他们看不出自己如何失准，或完全失策；因为没有标靶指控他们，显示他们偏离正道多远；也不稳定的目光无法帮助他们纠正并恢复所要求的直线。同样，我们所设定的终极目的，如宗徒所说，是永生；他宣告：\BibleQuote{你们既然有果实归于圣化，结局即是永生。}\BibleRef{罗 6:22} 但直接目标乃是心地纯洁，他并非不公平地称之为\Quote{圣化}；没有它，上述终极目的无法获得。仿佛他用别的话说：你们有直接目标在于心地纯洁，而最终目的则是永生。关于这目标，同一蒙福的宗徒教导我们，并明确使用该词，即\OriginalTerm{目标}{\GreekText{σκοπός}}，如是说：\BibleQuote{忘记背后，努力前面的，向着标竿直跑，要得天主在基督耶稣内从高天召我来的奖赏。}\BibleRef{斐 3:13} 这在希腊文更清楚地表述为\OriginalTerm{向着标竿}{\GreekText{κατὰ} \GreekText{σκοπὸν} \GreekText{διώκω}}，即\BibleQuote{我向着标竿直跑}，仿佛他说：\Quote{带着这个目标，我忘记背后的事，即早年生活的过错，我奋力争取天上的奖赏作为结局。}因此，凡能帮助我们导向这对象，即心地纯洁的，我们当尽全力追随；凡阻碍我们的，当视之为危险有害之物而避免。因为为此我们做一切、忍受一切；为此我们轻视亲属、祖国、荣誉、财富、世俗的享乐及各种欢愉——即为了保持持久的心地纯洁。因此，当这对象摆在我们面前时，我们当始终将我们的行动和思想笔直导向获得它；若它不常定于我们眼前，不仅会使我们的一切劳苦徒劳无益，使之无酬报地白受，还会激起种种互相冲突的思想。因为心若没有固定可回归、并主要系念的点，必然时时刻刻在种种游荡的思想中徘徊，并因外界临到的事物而不断变化为最先呈现给它的状态。\par

% structure: p0020 source=h2 target=subsection reason=relative-heading-rank; base=h2

\subsection{第六章}

\Emph{论那些舍弃世界，却无爱德而追求成全的人。}\par

因为由此就产生了这样的情形：有些人对世间最大的财富嗤之以鼻，不仅是大笔的金银，还有庞大的产业，而我们却看到他们后来竟为了刀、笔、针或笔而心神不宁、激动不已。如果他们能怀着纯洁的心，目光坚定地注视永恒，就绝不会因为琐事而让这种事发生；而为了免于在巨大而珍贵的财富上受苦，他们宁愿完全放弃它们。因为我们也常常看到一些人如此嫉妒地守护他们的书籍，以至于不允许任何人稍微移动或触碰它们，而由此他们遇到了不耐烦和死亡的场合，这些场合警告他们需要获得必要的耐心和爱德；当他们为了基督的爱放弃了全部财富，却仍然在琐事上保持着原来的心态，有时还会因此迅速不安，他们就变得在各方面不结果实、没有果子，如同那些没有宗徒所讲的爱德的人一样：有福的宗徒在精神上预见了这一点，说：\Quote{《虽然》他说，\InnerQuote{我若将全部财产施舍给穷人，又舍身被火焚烧，却没有爱德，于我毫无益处。}}\BibleRef{格前13:3}。由此清楚地看出，成圣并非仅仅通过克己、放弃所有财物和抛弃荣誉就能达到，除非有宗徒所描述的爱德，而这爱德只在于一颗纯洁的心。因为\Quote{不嫉妒，} \Quote{不自夸，不轻易发怒，不作不义的事，不求自己的益处，不为不义而欢喜，不怀恶意}等等，这不就是永远向天主呈上一颗完美而洁净的心，并使之不受任何扰乱吗？\par

% structure: p0023 source=h2 target=subsection reason=relative-heading-rank; base=h2

\subsection{第七章}

\Emph{如何寻求心平气和。}\par

一切事都应为这个目的而做、而寻求。为此我们必须寻求独处，为此我们知道应该接受斋戒、守夜、劳苦、赤身、读经和所有其他德行，好藉它们能预备我们的心，并使之不受一切邪恶情欲的伤害，在这基础上上升到爱德的完美境界。对于这些操练，如果我们偶然从事了一些有益有用的工作，而未能实行我们惯常的纪律，我们就不应被烦恼或愤怒或情欲所胜，而我们要克服这些，正是为了去做我们所忽略的事。因为斋戒的收益不能抵消愤怒的损失，读经的益处也不及轻视弟兄所带来的害处。那些次要的事，如斋戒、守夜、远离世俗、默想圣经，我们应当以我们的主要目的为目标来操练，即心地纯洁，这就是爱德；我们不应为了它们而驱走这个主要的德行，因为只要它在我们内仍然完整无损，我们就不会因那些次要的事被必要地忽略而受伤；因为如果这个我们所说的主要理由——一切事都为之而做——被移走，那么即使做了一切，也毫无用处。因为为了这个缘故，一个人急于为自己取得和准备任何工作的工具，并非只是无目的地拥有它们，也不是好像他把从它们那里期望的收益和益处仅仅看作拥有，而是通过使用它们，可以有效地获得实践知识和那门特定艺术的目的，而它们只是辅助手段。因此，斋戒、守夜、默想圣经、克己和放弃一切财物，并不是完美，而是达至完美的辅助；因为那门科学的终点不在这些事上，而是藉着这些事我们达到终点。那么，如果有人满足于这些，好像它们是至高的善，并将他心中的目的仅仅固定在这些事上，而不努力去达到目的——为这目的这些事才应寻求——那么他实践这些操练就是徒劳的：因为他确实拥有他艺术的手段，却不知道那一切价值所在的目的。因此，任何能扰乱那种心平气和的事——即使它看起来有益而有价值——都应被视为真正有害而躲避，因为凭这个规则，我们将能避免错误和任性的伤害，径直走向所期望的目的并达到它。\par

% structure: p0026 source=h2 target=subsection reason=relative-heading-rank; base=h2

\subsection{第八章}

\Emph{论默观事物的主要努力，及玛尔大和玛利亚的例子。}\par

这应当是我们主要的努力：我们应当不断向往这样的坚定心志，即灵魂永远紧贴天主和天上的事物。凡与此相异者，无论多么重大，都应置于次要地位，甚至视为无关紧要，或可能有害。我们在福音中，关于\Person{玛尔大}和\Person{玛利亚}的事例中，有一个极好的说明这种心态和状况的例证：因为当\Person{玛尔大}从事一项确实神圣的服务时，她是在服侍主和他的门徒，而\Person{玛利亚}只专注于神性教导，紧贴\Person{耶稣}的脚，亲吻并用美好宣认的香膏涂抹它们，主向她表明她选择了更好的一份，那是不应被夺去的：因为当\Person{玛尔大}怀着虔诚的操心忙碌，并为她的服务所困扰时，看到自己独自不足以胜任这样的服务，便向主请求她妹妹的帮助，说：\Quote{祢不介意我的妹妹留下我独自服侍吗？请吩咐她来帮助我}——诚然，她召唤她并非为了不配的工作，而是为了可赞美的服务：然而她从主那里听到了什么？\Quote{\Person{玛尔大}，\Person{玛尔大}，你为许多事操心忙碌：但需要的只有几件，或只有一件。\Person{玛利亚}选择了更好的一份，是不能从她那里夺去的。}你看，主将首要的善置于默想中，即神圣的默观：由此我们看出，其他一切德行都应置于次要地位，即使我们承认它们是必要的、有用的和卓越的，因为它们都是为这一件事而行的。因为当主说：\Quote{你为许多事操心忙碌，但需要的只有几件，或只有一件，}祂将首要的善不放在实践工作中（无论它多么可赞美和富有成果），而放在对祂的默观中，这默观确实是单纯的，且是\Quote{唯一的一件}；祂宣告\Quote{几件事}对于完美的真福是必要的，即那首先通过默想少数圣人而获得的默观：从对他们的默观中，有进步的人藉着天主的助佑上升并达到那被称为\Quote{一件事}的境地，即只考虑天主自己，从而超越那些圣人的行动和服务，只以天主的美善和知识为食粮。\Quote{因此\Person{玛利亚}选择了更好的那份，是不能从她那里夺去的。}这一点必须更仔细地思考。因为当祂说\Person{玛利亚}选择了更好的那份时，虽然祂并没有说什么关于\Person{玛尔大}的话，而且显然没有责备她，但在赞美一人的同时，祂暗示了另一人是次等的。再者，当祂说\Quote{是不能从她那里夺去的}时，祂表明从另一人那里，她的那份是可以被夺去的（因为肉身的服务不能与人永久长存），但教导说，这一人的渴求永无止境。\par

% structure: p0029 source=h2 target=subsection reason=relative-heading-rank; base=h2

\subsection{第九章}

\Emph{一个问题：德行的实践为何不能与人长存。}\par

我们深受感动，回答说：那么怎样？禁食的努力、勤于阅读、慈悲的工作、公义、虔诚和仁慈，这些会从我们这里被夺去，不与行这些事的人长存吗？尤其主自己向这些工作应许天国的赏报，当祂说：\Quote{你们蒙我父降福的，来承受自创世以来为你们预备的国度。因为我饿了，你们给了我吃的；我渴了，你们给了我喝的：}等等。\BibleRef{玛 25:34}那么，这些使行它们的人进入天国的工作，又怎能被夺去呢？\par

% structure: p0032 source=h2 target=subsection reason=relative-heading-rank; base=h2

\subsection{第十章}

\Emph{答复：不是赏报，而是它们的实行将要完结。}\par

\Person{梅瑟}。我并未说过善工的赏报会被夺去，正如主亲自说：\Quote{无论谁，只以门徒的名号，给这些小子中的一个，一杯凉水喝，我实在告诉你们，他决失不了他的赏报。}\BibleRef{玛 10:42}但我坚持说，那些或因身体的需要，或因肉体的冲击，或因这个世界的种种不平等而被迫去做的事，将会被除去。因为勤于阅读和斋戒中的克己，只在这现世对净化心灵和节制肉体有益，只要\Quote{肉体的私欲相反圣神的引导}\BibleRef{迦 5:17}；有时我们甚至看到，在现世中，这些事也会从那些因过度劳苦、身体虚弱或年迈而疲惫不堪的人身上被除去，他们无法再实行这些事。那么，将来岂不更是如此吗？那时\Quote{这可朽坏的将穿上不可朽坏的}\BibleRef{格前 15:53}，如今是\Quote{属生灵的身体}将要复活成为\Quote{属神的身体}\BibleRef{格前 15:44}，并且肉体将开始成为不再与圣神相争的样式。关于这一点，蒙福的宗徒也清楚地说：\Quote{身体的操练益处甚少；惟有虔诚（他显然指爱德）在各方面都有益处，因有今生和来生的应许。}\BibleRef{弟前 4:8}这清楚表明，所谓\Quote{甚少}的事，并非永远实行，也不可能单独赋予那为此劳苦的人以最高的完美境界。因为\Quote{甚少}一词可能指两件事之一：即或指时间的短暂，因为身体操练不可能在今生和来世都与人类共存；或指操练肉体所产生的利益的微小，因为身体的刻苦仅能产生进步的某种开端，而非爱的完美本身——爱拥有今生和来生的应许。因此，我们认为前述工作的实践是必要的，因为没有它们，我们就无法攀登爱德的高峰。因为你们所谓的宗教和慈善工作，在现世中，只要这些不平等和状况的差异仍然存在，就是必要的；但即使在此世，我们也不应期待它们被执行，除非有大量贫穷、困苦和患病的人存在——这是由人的邪恶所致，即那些把万物创造主同样赐予众人的东西，夺取并据为己有（却不使用）之人。因此，只要这种不平等在这个世界上持续，这类工作对实行它的人就是必要且有益的，因为它以善意的目的和虔诚的意愿，带来永恒产业的赏报。但在来世，它将终结，因为那时将充满平等，不再有不平等——正是为了这些不平等，这些事才必须做。那时，众人将从这些多样的实践工作，转向天主的爱德，以及以持续的心灵纯洁默想天上之事。那些急切献身于知识和心灵纯洁的人，已选择竭尽全力投身于这职责，当他们仍在肉体中时，便奔向这职责，他们将在卸下朽坏之后继续从事这职责，并来到主救主的那个应许，他说：\Quote{心里纯洁的人是有福的，因为他们要看见天主。}\BibleRef{玛 5:8}\par

% structure: p0035 source=h2 target=subsection reason=relative-heading-rank; base=h2

\subsection{第十一章}

\Emph{论爱德常存的特质。}\par

你们为何惊奇上述那些职责将会终止呢？圣宗徒告诉我们，连圣神更高的恩赐也会消逝，并指出惟有爱德永存不朽，说：\Quote{先知之恩终必消失；语言之恩终必停止；知识之恩终必消逝，}至于这点他又说：\Quote{爱德永存不朽。}因为所有恩赐都是暂时赐予，按需使用；但当这分施期结束时，它们无疑会立刻消逝。但爱德却永不毁灭。因为它不但在现世对我们有用地运作；在来世，当身体需要的负担卸下后，它将以更强大的力量和更卓越的方式继续，永不会被任何缺陷削弱，反而藉其永恒的不朽性，将更专注、更恳切地依附于天主。\BibleRef{格前 13:8}\par

% structure: p0038 source=h2 target=subsection reason=relative-heading-rank; base=h2

\subsection{第十二章}

\Emph{论在神性默观中持之以恒的问题。}\par

\Person{杰尔马努斯}说：那么，谁在背负我们脆弱的肉体时，能常常如此专注于这默观，以致从不想到有弟兄到来、探访病人、手工劳作，或至少对陌生人及访客表示仁慈呢？最后，谁不被照顾和照料身体所打断呢？或者，心灵如何以及以何种方式能够依附于无形且不可理解的天主？我们愿意了解这一点。\par

% structure: p0041 source=h2 target=subsection reason=relative-heading-rank; base=h2

\subsection{第十三章}

\Emph{关于心灵指向及关乎天主之国和魔鬼之国的答复。}\par


\Person{梅瑟}说：\Quote{人还处在这软弱的肉身中时，要持续不断地依附天主，并如你所说，不可分离地持守对他的默想，是不可能的。但我们应当知道，我们应将心灵的目标固定在哪里，应时常将灵魂的目光转回何方：当心灵能达成这一点时，它可以喜乐；当它偏离此目标时，应忧伤叹息，每当发现自己从注视他中坠落时，应承认自己已从至善中堕落，认为即便片刻离开对基督的注视也是邪淫。当我们的目光稍微偏离他时，让我们将灵魂的眼睛转回他那里，并仿佛沿着一条笔直的路线召回我们心灵的注视。因为一切都取决于内心的状态，当魔鬼从其中被驱逐，罪恶不再在其中作王时，天主的国就建立在我们里面，正如圣史所说：\BibleQuote{天主的国来临时，并非显然可见的，人也不说：看哪，在这里；或说：在那里。因为我实在告诉你们：天主的国就在你们中间。}\BibleRef{路 17:20} 但除了对真理的认识或无知，以及或对罪恶或对美德的爱悦之外，没有什么能\BibleQuote{在你们里面}，通过这些我们在心中为魔鬼或为基督预备一个国度；宗徒描述了这国度的特征，他说：\BibleQuote{因为天主的国不在于吃喝，而在于义德、平安以及在圣神内的喜乐。}\BibleRef{罗 14:17} 所以，如果天主的国在我们内，而天主的国实际就是义德、平安和喜乐，那么住在这三者中的人肯定是在天主的国里；相反，那些生活在不义、纷争以及招致死亡的悲伤中的人，则属于魔鬼的国、地狱和死亡。因为藉着这些标记，天主的国和魔鬼的国得以区分；事实上，如果我们举目仰望高处，思考天上掌权者在天上的生活状态——他们真正在天主的国里——我们应当想象那是什么，除了永恒持久的喜乐？因为真正的幸福所特别专有和相称的，不就是恒常的安宁和永恒的喜乐吗？为了让你确信我们所说的确实如此——不是凭我自己的权威，而是凭主的权威——请听他多么清楚地描述了那世界的特征和状态：\BibleQuote{看，我要创造新天新地，先前的事不再被记忆，也不再被思念；但你们要因我所创造的而永远欢喜快乐。}\BibleRef{依 65:17} 又说：\BibleQuote{其中必找到欢喜和快乐，感谢和赞美的声音；且有月月如此，安息日接连安息日。}\BibleRef{依 66:23} 又说：\BibleQuote{他们必获得欢喜快乐，忧愁和叹息必将逃遁。}\BibleRef{依 35:10} 如果你想要更确切地了解那生命和圣人之城，请听主向天上的耶路撒冷亲自宣告的声音：\BibleQuote{我要使你长官平安，监督正义。你境内不再听见强暴，你疆界内不再有荒凉和毁灭；你的墙壁要称为救恩，你的城门要称为赞美。太阳不再作你白昼的光，月亮的光也不照耀你；上主要作你永远的光，你的天主要作你的光荣。你的太阳不再下落，你的月亮也不退缩；上主要作你永远的光，你悲哀的日子已经结束。}\BibleRef{依 60:17} 因此圣宗徒不是说一般的、无区别的喜乐就是天主的国，而是明确地、强调地指出只有\BibleQuote{在圣神内}\BibleRef{罗 14:17}的喜乐才是。因为他完全知道另一种可憎的喜乐，关于这喜乐我们听到：\BibleQuote{世界要欢乐}\BibleRef{若 16:20}，以及\BibleQuote{你们现今欢笑的人是有祸的，因为你们将要哀哭。}\BibleRef{路 6:25} 实际上，天国必须从三重意义上理解：要么是诸天（即圣人们）将统治其他被征服的事物，正如这句话所说：\BibleQuote{你管五座城，你管十座城}，以及向门徒们说的：\BibleQuote{你们要坐在十二个宝座上，审判以色列十二支派}\BibleRef{玛 19:28}；要么是诸天本身将开始被基督统治，当\BibleQuote{万物都归于他}\BibleRef{格前 15:28}，天主开始成为\BibleQuote{万有中的万有}时；要么是圣人们将在天上与主一同作王。}\par
% structure: p0044 source=h2 target=subsection reason=relative-heading-rank; base=h2

\subsection{第十四章}

\Emph{论灵魂的持续。}\par

因此，凡尚在这肉身内生活的人，就应当已经知道，他必将被托付于那状态和职务，就是他今生使自己成为分享者和追随者的状态和职务；他也不应怀疑，在永世里他将成为那人的伙伴，而在此世他选择了使自己成为那人的仆人和臣仆：依照我主的那句话，他说：\BibleQuote{人若服事我，就当跟随我；我在哪里，我的仆人也必在哪里。} \BibleRef{若 12:26} 因为正如魔鬼的王国是通过同意罪恶而获得的，同样，天主的王国是通过心地纯洁和灵性知识修行德性而达到的。但是，哪里有天主的王国，那里就必定享受永生；哪里有魔鬼的王国，那里无疑就是死亡和坟墓。处在这种状况的人，不能赞美主，正如先知的话告诉我们的：\Quote{死人不能赞美你，上主；凡下入坟墓的（无疑是指罪恶的坟墓）也不能赞美你。但我们}——他说——\Quote{活着的人（实在不是向罪恶或这世界活着，而是向天主活着）将从现今直到永远赞美上主；因为在死亡中无人记念天主；在（罪恶的）坟墓中，谁能承认上主呢？}也就是说，没有人会。因为没有人——即使他自称为基督徒一千次，或自称隐修士——在犯罪时承认天主；没有人容忍主所憎恶的事，记念天主，或真实地自称为他的仆人，而他却以固执的鲁莽藐视他的命令。真福宗徒宣布，贪图享乐的寡妇就陷于这种死亡，他说：\BibleQuote{寡妇若耽于享乐，虽生犹死。} \BibleRef{弟前 5:6} 因此有许多人，虽然还活在这肉身内，却是死的，躺在坟墓里不能赞美天主；反之，有许多人，虽然肉身死了，却在心灵中赞美天主，歌颂他，正如这里所说：\Quote{正义者的众灵和灵魂，请赞美上主}；以及\Quote{每一个灵都要赞美上主。}在\BibleBook{}{默示录}中，被宰杀者的灵魂不仅被说成赞美天主，而且也与他对谈。\BibleRef{默 6:9} 在福音中，主对撒杜塞人说得更清楚：\BibleQuote{你们没有读过天主对你们所说的话吗？他说：我是亚巴郎的天主，依撒格的天主，雅各伯的天主。他不是死人的天主，而是活人的天主：因为所有的人都是为他而活。} \BibleRef{玛 22:31} 关于这些人，宗徒也说：\BibleQuote{因此，天主不以被称为他们的天主为耻：因为他已为他们预备了一座城。} \BibleRef{希 11:16} 因为他们在脱离此肉身之后并非无所事事，也并非没有感知，福音中的比喻向我们证明了这一点，即关于乞丐拉匝禄和身穿紫红袍的富翁的比喻：其中一个获得了福乐的位置，即亚巴郎的怀抱；另一个则被永火的可怕热焰所吞噬。但若你也愿意理解对那强盗所说的话：\BibleQuote{今天你就要与我一同在乐园里。} \BibleRef{路 23:43} 这岂不明确地表明，不但他们先前的智力继续与灵魂同在，而且在他们改变的状态中，他们也享有某种与他们的行为和功过相应的状态吗？因为主若知道他的灵魂在脱离肉身之后会被剥夺知觉或化为虚无，就决不会向他应许这个。因为要与基督一同进入乐园的不是他的肉身，而是他的灵魂。至少我们必须避免，并以最大的恐惧回避异端者那种邪恶的标点法；他们不相信基督在降临地狱的同一天能出现在乐园里，于是这样标点：\Quote{我实在告诉你，今天}，然后停顿，再读\Quote{你就要与我一同在乐园里}，以此方式想象这应许在他离世后并未立即实现，而是在复活后才实现；因为他们不明白，在他复活之前，他曾对那些以为他也像他们一样被人类的困难和肉身的软弱所困扰的犹太人宣告：\BibleQuote{没有人升过天，除了那从天上降下来的人子，他仍在天上。} \BibleRef{若 3:13} 由此他清楚地表明，去世者的灵魂不但没有被剥夺理性，而且甚至并非没有希望、忧愁、喜乐和恐惧等情感；他们已经预先开始品尝在最后审判中为他们所保留的某些东西；他们也不像某些不信之人所认为的那样，在离世后化为虚无，而是过一种更真实的生活，并更加热切地事奉对天主的赞美。事实上，暂且放下圣经的证据，尽我们所能稍微讨论一下灵魂的本性本身，若竟会有——我不说是愚昧，而是所有愚蠢的疯狂——哪怕丝毫怀疑，即人的更高贵的部分（正如真福宗徒所表明的，天主的肖像和模样就在于此）在卸下今世压迫它的肉身的重担后会变得没有知觉，而这部分自身包含理性的一切力量，通过与之结合，使那愚钝无觉的物质肉身有了知觉；尤其是当它——并且理性本身的秩序要求——当心灵脱去如今压着它的粗重肉身后，它将比以往更好地恢复其智力能力，并以比失去它们时更纯净、更精妙的状态接受它们，这难道不是完全超出合理范围吗？但是真福宗徒如此认识到我们所说的真实，以至于他确实愿意离开这肉身；通过与之分离，他将能更热切地联合于主；他说：\Quote{我渴望被解散，与基督同在，这远为更好，因为我们在肉身内时，是远离主的。}因此，\Quote{我们放心大胆，并常怀此心愿，情愿离开肉身，与主同在。为此，我们无论离开或留下，都努力讨他喜悦。}；他确实宣告，灵魂在肉身中的持续存在就是远离主，离开基督，并完全相信，与这肉身的分离和离去意味着与基督同在。同一宗徒又更清楚地论及灵魂的这种非常充满生命的状态：\BibleQuote{但你们却来到了熙雍山，永生天主的城，天上的耶路撒冷，和千万天使的盛会，以及被登记在天上的首生者的教会，和已达到完美的义人的灵魂。} \BibleRef{希 12:22} 他在另一段经文论及这些灵魂说：\Quote{此外，我们曾有过我们肉身的管教者，我们也尊敬他们；难道我们不是更该服从众灵之父而生活吗？}\par

% structure: p0047 source=h2 target=subsection reason=relative-heading-rank; base=h2

\subsection{第十五章}

\Emph{我们该如何默想天主}\par

但是，默观天主可以通过多种方式获得。因为我们不仅通过赞叹祂不可理解的本质（这事仍隐藏在恩许的希望中）来发现天主，而且通过祂创造的伟大、对祂公义的默想以及祂日常眷顾的帮助来看见祂；当我们以纯洁的心灵默想祂在每一世代对圣人所做的一切，当我们以战栗的心钦佩祂统治、指导和管理万物的权能，或祂知识的广大无边，以及祂的那只眼——没有任何心中的秘密能向它隐藏，当我们思考海沙和浪涛的数量（这些由祂测量并为我们所知），当我们在惊异中想到雨滴、岁月的时间、时辰以及过去和未来的一切都呈现于祂的知识；当我们以无限的赞叹凝视祂那不可言喻的仁慈——这仁慈以不知疲倦的忍耐忍受着每时每刻在祂眼前犯下的无数罪过，或祂那召唤——不是因我们任何先前的功绩，而是因祂怜悯的自由恩宠，祂接纳了我们；或者再次，祂赐予那些祂将要收纳的人无数得救的机会——祂使我们如此诞生，以至从我们的摇篮开始，祂的恩宠和祂法律的知识就能赐予我们；祂自己，仅仅为了祂善意的喜悦，在我们内战胜我们的仇敌，以永福和永远的赏报奖赏我们；最后，祂为我们的救恩承担了祂降生成人的救世工程，并将祂圣事的奇妙扩展到万邦。但还有无数其他这类的默想，它们根据我们生活的特质和心灵的纯洁在我们心中升起，通过它们，天主要么被纯洁的眼睛看见，要么被拥抱；这些默想当然没有人能长久保持，如果在他内仍存有任何肉性的情感，因为主说：\Quote{你不能，}主说，\Quote{看见我的面容：因为人看见了我，就不能再活了。} \BibleRef{出 33:20} 即，属于这个世界和属地的情感。\par

% structure: p0050 source=h2 target=subsection reason=relative-heading-rank; base=h2

\subsection{第十六章}

\Emph{一个关于思想易变性质的问题}\par

\Person{热尔马努斯}。那么，为何甚至违背我们的意愿，是的，在我们不知不觉中，闲散的思想如此微妙而隐秘地潜入，以至于不仅驱赶它们，甚至捕捉和抓住它们都极其困难？那么，有时能否找到一颗完全摆脱它们、从未被此类幻觉攻击的心灵？\par

% structure: p0053 source=h2 target=subsection reason=relative-heading-rank; base=h2

\subsection{第十七章}

\Emph{心灵在思想状态中能与不能做什么的答案}\par

\Person{梅瑟}。心灵不可能不被思想接近，但每个认真的人都有权接纳或拒绝它们。既然它们的升起不完全取决于我们自己，那么拒绝或接纳它们则在于我们自己的能力。但是因为我们说过心灵不可能不被思想接近，你们不可将一切都归咎于攻击，或归咎于那些试图将思想灌输给我们的恶灵，否则人内就不会保留任何自由意志，我们改善的努力也不在我们的能力之内。但是，我说，在很大程度上，改善我们思想的特质，并让圣善和属神的思想或属地的思想在我们心中成长，是在我们的能力之内的。因为为此目的，经常阅读和不断默想圣经，以便从中获得属神反省的机会；因此，经常咏唱圣咏，以便从中提供持续的痛悔之情；以及认真的守夜、斋戒和祈祷，以便使心灵谦卑，不思念属地之事，而默想属天之事；因为如果这些被放弃，疏忽潜入我们，心灵因罪的污秽而刚硬，必然倾向肉性的方向而堕落。\par

% structure: p0056 source=h2 target=subsection reason=relative-heading-rank; base=h2

\subsection{第十八章}

\Emph{灵魂与磨石的比喻}\par

而这心灵的动态，用磨轮的比喻来说明并非不恰当：磨轮被湍急的水流冲击而旋转，只要被水的动作推动，就永不停息它的工作。但在管理它的人的能力之内，来决定是用它来磨小麦、大麦还是毒麦。当然，负责那事务的人放入什么，就必须被它碾碎。同样，心灵在今生考验中，也被从四面八方涌来的诱惑激流冲激，不能摆脱思想的流动。但思想的特质——它应当抛弃或接纳——将由它自己努力和勤奋来决定。因为如果我们如我们所说，不断回归对圣经的默想，并提升我们的记忆向属神事物的回想、对成全的渴望和未来永福的希望，属神的思想必然由此升起，并使心灵专注于我们所默想的事物。但是如果我们被懒惰或疏忽胜过后又虚度光阴进行闲谈，或纠缠于今世的挂虑和不必要的焦虑，结果将是某种杂草丛生，给我们的心灵带来有害的占据；正如我们的主和救主所说：\BibleQuote{因为你的宝藏在哪里，你的心也必在哪里。} \BibleRef{玛 6:21}\par

% structure: p0059 source=h2 target=subsection reason=relative-heading-rank; base=h2

\subsection{第十九章}

\Emph{论我们思想的三个来源}\par

首先，我们至少应当知道，我们的思想有三个来源，即来自天主、来自魔鬼和来自我们自己。它们来自天主，当祂愿意以圣神的照耀眷顾我们，提升我们达到更高的进步境界，当我们因懈怠而进步甚少或因行动懒散而被击败时，祂以最有益的痛悔来惩罚我们，或向我们启示天上的奥秘，或使我们的意向和意志转向更好的行为，正如阿哈随鲁王受到主的惩戒后，被促使索要编年史卷，从而记起摩尔德开的善行，擢升他至极高尊荣之位，并立刻撤销了他关于屠杀犹太人最残酷的判决。或如先知所说：\Quote{我要听上主天主在我内所说的话。} 另有一位告诉我们：\BibleQuote{天使对我说，} \BibleRef{匝 1:14} 或当天主子应许祂要与父一同来，并在我们内居住，\BibleRef{若 14:23} 以及\BibleQuote{说话的不是你们，而是你们父的圣神在你们内说话。} \BibleRef{玛 10:20} 还有蒙选之器：\BibleQuote{你们寻求基督在我内说话的明证。} \BibleRef{格后 13:3} 然而，一整系列的思想来自魔鬼，当他试图以罪乐或隐秘攻击毁灭我们时，以狡猾诡计奸诈地将恶显示为善，并为我们将自己扮成光明的天使，\BibleRef{格后 11:4} 正如福音作者告诉我们的：\BibleQuote{晚餐后，魔鬼已使依斯加略人西满的儿子犹达斯决心出卖} \BibleRef{若 13:2} 主；并且再次，他说：\Quote{吃了饼后，}他说，\Quote{撒殚进入了那人的心。} 伯多禄也对阿纳尼雅说：\BibleQuote{为什么撒殚充满了你的心，叫你欺骗圣神？} \BibleRef{宗 5:3} 还有我们在福音中早已读到、由训道篇预言的那句：\BibleQuote{如果统治者的灵魂起来攻击你，不要离开你的位置。} \BibleRef{训 10:4} 以及在列王纪上中对阿哈布所说的邪神之言：\BibleQuote{我要出去，在所有先知的嘴里作一个虚谎的灵。} \BibleRef{列上 22:22} 但我们的思想来自我们自己，当我们自然地回忆正在做、已做过或已听到的事时。蒙福的达味这样说：\Quote{我想起往昔的日子，心中存念远古的岁月；夜间我与心默想，省察自己的神魂。} 又说：\BibleQuote{主知道人的思念，原是虚幻的；} 以及\BibleQuote{义人的思念是公正的。} \BibleRef{箴 12:5} 福音中主也对法利塞人说：\BibleQuote{你们为什么心里思念恶事？} \BibleRef{玛 9:4}\par

% structure: p0062 source=h2 target=subsection reason=relative-heading-rank; base=h2

\subsection{第二十章}

\Emph{关于识别思想的实例，来自一位好货币兑换商的说明。}\par

我们应当仔细留意这一三重秩序，并以明智的辨别力分析心中升起的思虑，首先追寻它们的来源、原因和作者，以便我们能够考虑如何根据那些建议者所应得的，来确定自己该如何顺从它们，从而如同主的命令所吩咐我们的那样，成为精明的兑换银钱的人，他们最高的技能和训练就是检验什么是完全纯金，什么是通常称为\Emph{试过的}金，或者什么是在火中不够精炼的；同时要以无误的技能不被普通的黄铜\OriginalTerm{德纳}{denarius}所欺骗，如果它被涂上明亮的金色而看起来像某种贵重钱币；不仅机敏地认出印有篡位者头像的硬币，还要以更机敏的技艺察觉那些带有真正国王肖像但制作不当的硬币，最后要小心地通过天平的检验来确保它们没有不足重量。所有这一切，那采用此比喻的福音说法指示我们，也应当在精神上遵守；首先，凡进入我们心中的事物，以及我们所接受的一切道理，都应当最仔细地审查，看它是否被圣神的神圣和天上的火所净化，或是否属于犹太人的迷信，或是否来自世俗哲学的骄傲，而只在表面上显示虔诚。如果我们遵守宗徒的劝告，我们就能做到这一点：\BibleQuote{不要信任何神，但要考验这些神是否出于天主。} \BibleRef{若一 4:1}但是，那些在宣誓成为隐修士后，被文采的优雅和哲学家的某些学说所吸引的人，也会被这类事物所欺骗；这些学说乍看之下，由于一些与宗教并不冲突的虔诚含义，像黄金的闪光一样欺骗了那些被表面吸引的听众，却使他们永远成为贫穷可怜的人，如同被假铜钱欺骗的人一样：要么使他们回到这个世界的喧嚣中，要么引诱他们进入异端者的错误和浮夸的妄念中。我们在\WorkTitle{若苏厄书}中读到，这事曾发生在\Person{阿干}身上，\EditorialNote{若苏厄书第七章}当他贪图\Place{培肋舍特人}营中的一块金子，并偷了它时，便受到诅咒，被判处永死。其次，我们应当小心，免得任何错误的解释附在圣经的纯金上，欺骗我们关于金属的价值；魔鬼藉着他的诡计，试图以此欺骗我们的主和救主，好像祂只是一个凡人，当他以恶意的解释曲解了本应普遍适用于所有善人的经文，并试图将其特别应用于祂——那位不需要天使看顾的——身上时，说：\BibleQuote{祂必为你吩咐祂的天使，在你行的一切道路上保护你；他们要用手托着你，免得你的脚碰在石子上。}他巧妙地假设，将圣经宝贵的言论转向并扭曲成危险的意义，与其真意完全相反，以便在可能欺骗我们的金色覆盖下，向我们呈现篡位者的形象和面孔。或者他试图用赝品欺骗我们，例如怂恿我们从事某项虔诚的工作，这工作虽然出自父辈们的真实心意，却在美德的形态下导向恶习；并且，通过无节制或不可能的斋戒、过长的守夜、无度的祈祷或不适当的阅读来欺骗我们，将我们引向坏结局。或者，当他劝我们投身于干涉别人的事务和虔诚的探访，以此驱使我们离开隐修院的精神回廊及其友善安宁的隐秘，并建议我们肩负起贫困虔敬女子的焦虑和挂虑，这样当一位隐修士被这种网罗无可救药地缠住时，他就能以最具伤害性的工作与挂虑来分他的心。或者当他以教化众人和爱慕属灵收益为借口，煽动一个人渴望圣职，从而将我们拉离自己生活的谦逊和严谨。这一切事物，尽管与我们的救恩和修会誓愿相悖，但当它们披上一种怜悯和宗教的幔子时，却轻易欺骗那些缺乏技能和谨慎的人。因为它们模仿真正君王的钱币，起初看似充满虔诚，但并非由有权铸造者——即被认可的天主教父——所铸，也不从接收它们的公共官署发出，而是由魔鬼暗中欺诈制成，并被强塞给不熟练和愚昧的人，造成严重的损害。即使它们起初似乎有用且必要，但如果后来它们开始损害我们修会的健全，并在某种意义上削弱我们整个目标的整体，那么将它们切断并像可能必要却冒犯我们的肢体一样丢弃，是非常好的，这肢体似乎执行右手或脚的职务。因为与其带着全部命令的整块而陷入某种错误，因邪恶的习惯使我们脱离严格的规则和已立志实行的体系，导致如此损失，以至于它永远无法弥补随后的伤害，反而会使我们过去的所有果实和工作的全部身体在地狱之火中被烧毁，不如缺少一条命令的肢体——即其运作或结果——而在其他部分保持安全无恙，并作为软弱者进入天国。\BibleRef{玛 18:8}关于这类幻觉，箴言中说得很好：\BibleQuote{有一条路，人以为正，但它的终局却进入地狱的深处；}又说：\BibleQuote{恶人接近善人时，是有害的，}也就是说，魔鬼在圣善外貌的遮盖下欺骗人：\BibleQuote{但他恨恶守望者的声音，}即来自父辈话语和警告的辨别之力。\par

% structure: p0065 source=h2 target=subsection reason=relative-heading-rank; base=h2

\subsection{第二十一章}

\Emph{论若望院长的幻觉。}\par

我们听说，住在\Place{吕孔}的\Person{若望}院长最近就如此受骗。他的身体因两日禁食未进食物而疲惫衰弱，第三天他正要去用些食物时，魔鬼以肮脏的埃塞俄比亚人形象出现，俯伏在他脚前喊道：\Quote{请宽恕我，因为我将这劳苦指定给了你。}于是这位在明辨方面如此成全的伟大人物明白了：他在不合时宜的禁食伪装下，被魔鬼的诡计所骗，从事了一种禁食，使疲惫的身体遭受不必要的劳苦，实则对灵性有害。他如同被伪币所骗，虽然对其上所印真国王的像表示敬意，却未足够警觉这钱币是否切割、印记正确。而这\Quote{善辨钱币者}的最后职责——如我们之前所提到的，关乎重量的检验——将得以履行：每当我们思想中建议做什么时，我们谨守地再三思考，将它放在心口的秤上，以最精确的天平称量：它是否对所有人充满益处，或因敬畏天主而沉重；或含义完整而健全；或因人的炫耀或某种新奇的自负而轻浮，或愚昧虚荣的骄傲是否减少或减轻了其功绩的重量。于是我们立即将它们放在公共天平上，即通过宗徒和先知的行动与证据来检验，将它们视为完整、完美且足重，或者小心谨慎地拒绝它们，如同不完美、伪劣和重量不足。\par

% structure: p0068 source=h2 target=subsection reason=relative-heading-rank; base=h2

\subsection{第二十二章}

\Emph{论辨别的四种方法。}\par

这种辨别能力将按我们所说的四种方式对我们必要：第一，那材料不让我们忽视它是真金还是镀金；第二，那些虚假许诺宗教行为的思想应如伪造和假冒的钱币被我们拒绝，因为它们印记不正，带有国王的假像；同样，我们应能识别那些在圣经的宝贵金子中，通过虚假和异端的含义，显示的不是真国王而是篡位者形象的思想；第三，那些其重量和价值被虚荣的锈蚀贬低，未能在教父们的天平上通过的钱币，如太轻、虚假和重量不足的，我们应拒绝；这样我们就不会因主的命令而受到警告，丧失一切劳苦的价值与赏报。\BibleQuote{不要在地上为自己积蓄财宝，那里有锈和虫蛀，也有贼挖洞偷窃。}\BibleRef{玛 6:19}因为每当我们为寻求人的荣耀而做某事时，我们知道，如主所说，是在地上为自己积蓄财宝，因而如同埋在地中、藏于土里，必被各种魔鬼毁坏，或被虚荣的蛀锈吞噬，或被骄傲的飞蛾吞吃，对藏它的人毫无用处和益处。因此，我们当不断探察心中所有的内室，以最严密的审察追踪进入其中的一切足迹，免得——不妨这样说——与悟性相关的野兽，或是狮子或是龙，经过时偷偷留下了危险的痕迹，为他人进入内心隐秘处所开了通道，这是由于对思想的疏忽。于是，日复一日、时复一时，我们用福音之犁——即恒常记念主的十字架——翻耕心田，从而得以从心中根除或灭绝毒兽的巢穴和毒蛇的潜伏处。\par

% structure: p0071 source=h2 target=subsection reason=relative-heading-rank; base=h2

\subsection{第二十三章}

\Emph{论教师根据听者的功德而施教。}\par

这时，那位长者见我们惊讶不已，因他讲论的话语而燃起不可满足的渴望，便因我们的赞叹和热诚稍停片刻，随即补充说：\Quote{我儿，既然你们的热情引出了如此长久的讨论，一种好似火焰的东西，按着你们的恳切，给我们的交谈增添了更锐利的兴致，我由此可以清楚看出，你们确实渴望关于成全的教导，我愿再对你们说些关于明辨与恩宠之卓越的事，这明辨与恩宠统治着所有德行并占据其领地；我不仅要通过日常实例证明它的价值与用处，也要借助教父们先前的深思与意见。因为我记得，当人们常常叹息流泪向我请求此类讲论时，我自己也渴望给他们一些教导，却无法办到；不仅我的思想，甚至我的言语能力也衰败了，以致我无法找到如何使他们即使带着些许安慰离开。由这些迹象我们清楚看到，主的恩宠按照听者的功德和热诚赐给讲者话语。又因眼前的夜晚极短，不容我完成讲论，我们不如将之让给身体的休息——如果拒绝合理的部分，整个讲论将不得不耗费在这休息中——并将讲论的完整纲要留给未来某日或某夜不间断地思虑。因为作明辨最佳顾问的人，首当其冲就应在此事上展现其心智的勤勉，并借此证据和耐心证明他们是否拥有且能持有明辨，以免在论述那作为节制之母的德行时，自己反倒陷入其相反的恶习，即过度冗长，以其言行破坏了他们在言语中所推荐的方法和本质的力量。那么，关于这最卓越的明辨——我们仍打算在主赐给我们能力范围内加以探究——首先，当我们辩论它的卓越以及那作为诸德之首存在于其中的节制时，善莫大焉的是不让我们自己超过讲论和时间的适当限度。}\par

于是，真福\Person{梅瑟}以此话止住了我们的交谈，并劝我们——我们虽急切切、侧耳倾听——去小睡片刻；他建议我们就在坐着的席子上躺下，把包裹放在头下代替枕头；因为这些东西是用较厚的纸莎草叶均匀地编成细长束，间隔六英尺，它们有时为弟兄们在礼仪中坐着时提供极低的座位代替脚凳，有时在他们就寝时放在颈下作为头部的支撑，不硬不软，舒适恰好。对于隐修士的这些用途，这些东西被认为特别合适，不仅因为它们稍软、花费金钱与劳力极少——纸莎草沿\Place{尼罗河}两岸到处生长，而且因为它们质地方便、足够轻巧，可随需要搬动或取来。于是，最后在长者的命令下，我们怀着因举行过的交谈而欢乐兴奋的心情，并因所许诺的讨论而满怀希望，极其安静地安歇入睡。\par

\SourceRecord{Translated by C.S. Gibson.；Nicene and Post-Nicene Fathers, Second Series；Vol. 11.；Edited by Philip Schaff and Henry Wace.；Buffalo, NY: Christian Literature Publishing Co.,；1894.；<http://www.newadvent.org/fathers/350801.htm>.}

% structure: p0001 source=h1 target=section reason=page-title

\section{第二篇会谈}

\Signature{圣若望·卡西安}\par

院长摩西的第二篇会谈。\par

% structure: p0004 source=h2 target=subsection reason=relative-heading-rank; base=h2

\subsection{第一章}

\Emph{院长摩西关于分辨之恩的引言。}\par

于是，当我们享受了清晨的睡眠，当黎明的光辉再次愉快地照耀我们，我们开始再次请求他应许的谈话时，有福的摩西就这样开始了：我看见你们燃起如此热切的渴望，以致我相信我本想从我们的灵修会谈中抽出并用于身体休息的那段极短暂的安静，对你们身体的休息毫无益处；当我看清你们的热情时，更大的焦虑也压在我身上。因为我必须更加用心和热忱地偿还我的债务，正与我看到你们越是恳切地要求它相称，如经上所说：\Quote{当你与一位官长同席时，要仔细考虑摆在你面前的是什么，并伸出手，知道你应该准备这样的事。}因此，既然我们要谈论分辨的卓越品质及其德行，这个话题我们昨晚的讨论结束时已经涉及，我们认为最好首先以教父们的意见来确立它的卓越性。当我们的前人对此的想法和说法被展示出来之后，我们就可以提出一些古代和近代的、各种人的沉船和不幸事件，他们因不太注意它而毁灭并绝望地遭难，然后我们尽可能地讨论它的益处和用途：在讨论之后，通过思考其价值和恩宠的重要性，我们将更好地知道如何寻求它并实践它。因为这并非普通的德行，也不是单凭人的努力可以自由获得的，除非有天主祝福的帮助，因为我们读到这也被宗徒列为圣神最高贵的恩赐之一：\BibleQuote{有人由圣神领受智慧的言语，另有人由同一圣神领受知识的言语，另有人由同一圣神领受信德，另有人由同一圣神领受治病的恩赐，}紧接着，\BibleQuote{另有人辨别神恩。}然后在神恩的完整目录之后，他补充说：\BibleQuote{然而这一切都是同一而同一的圣神所运行，随祂的意愿个别分给各人。}\BibleRef{格前 12:8}由此可见，分辨的恩赐并非尘世之物，也非小事，而是天主恩宠最伟大的赏赐。除非一位隐修士以全部的热忱追求它，并获得了以无误的判断辨别在他内升起的各种神类的能力，他必定会走错路，仿佛在黑夜和浓密的黑暗中，不仅会跌入危险的坑穴和悬崖，而且在平坦直接的事情上也经常犯错。\par

% structure: p0007 source=h2 target=subsection reason=relative-heading-rank; base=h2

\subsection{第二章}

\Emph{唯有分辨能给隐修士什么；以及真福安东尼关于此事的论述。}\par

我记得，在我还是一位少年时，在真福安当所居住的\Place{特巴依德地区}，长老们曾到他那里请教成全之道；那次会谈从傍晚一直持续到早晨，而夜间的绝大部分时间都花在这个问题上。因为人们详尽地讨论了哪一种德行或操守能保护隐修士始终不受魔鬼的陷阱和欺骗的伤害，并引领他沿着一条稳妥正直的道路，以坚定的步伐迈向成全的高峰。当每个人都按照自己心意的偏向提出了意见：有些人认为在于斋戒和守夜的热忱，因为灵魂因这些而谦卑，并获得心灵和身体的洁净，就更容易与天主结合；有些人认为在于轻看万物，因为心若完全舍弃万物，就能更自由地来到天主面前，仿佛从此再无任何陷阱能缠住它；另有些人认为，离群索居是必要的，即独居和隐修生活的隐秘，在这样的生活中，人更容易与天主交往，并特别紧贴祂；还有些人主张，应当实践爱德，即仁慈的义务，因为主在福音中特别应许将王国赐给这些人，祂说：\BibleQuote{我父所祝福的，你们来吧！承受自创世以来为你们预备的王国，因为我饿了，你们给了我吃的；我渴了，你们给了我喝的，等等：}\BibleRef{玛 25:35}。当他们以这种方式宣称，通过不同的德行可以获得更确定地接近天主的途径，并且夜间的绝大部分时间都花在这讨论上时，最后，真福安当发言说：你们所提到的这一切，对于渴望天主、想要亲近祂的人，确实是必要且有益的。但是，无数的意外和许多人的经验不容许我们将最重要的恩赐寄托在这些事上。因为，常常有些人尽管在斋戒和守夜上最为严格，又高尚地退隐独处，并且彻底舍弃自己的一切财物，连一天的口粮或一支钱都不留下，又极其虔诚地履行一切仁慈的义务，然而我们却看到他们突然被欺骗，以致不能将他们所开始的工作妥善完成，反而将他们崇高的热忱和可称赞的生活方式引向可怕的结局。因此，如果我们更仔细地追溯他们堕落和受骗的原因，就能清楚地认识到什么才是主要引领人走向天主的。因为当上述那些德行在他们身上满溢时，唯独缺少了\OriginalTerm{分辨}{discretion}，就不允许他们坚持到底。除了因为他们没有从长老那里得到充分的教导，以致未能获得判断和分辨（这分辨超越两端的过度，教导隐修士始终沿着王道行走，既不让他因德行之右而自高，即因过度的热忱而愚蠢地越出适当节制的界限，也不让他因恋慕松懈而转向左边的恶习，即以控制身体为借口，因相反懈怠的精神而变得松弛）之外，再也找不到他们坠落的其他原因了。这就是分辨，它在福音中被称为\Quote{眼睛}，\Quote{身体的光}，按救主的说法：\BibleQuote{你的眼睛就是身体的灯：如果你的眼睛单纯，你全身就都明亮；如果你的眼睛邪恶，你全身就都黑暗。}\BibleRef{玛 6:22}；因为它辨别人的一切思想和行为，看见并审视所有应当做的事。但是，如果在某人身上这眼睛是\Quote{邪恶}的，即没有健全判断和知识的坚固，或因某种错误和自以为是而受骗，就会使我们的全身\Quote{充满黑暗}，即会蒙蔽我们整个的精神视觉和行为，使它们陷入恶习的黑暗和纷扰的阴霾中。因为，祂说：\BibleQuote{如果你内在的光成了黑暗，那黑暗是多么大啊！}\BibleRef{玛 6:22}。因为无人怀疑，当我们内心的判断出错，被无知的黑暗所淹没时，我们的思想和行为（这些是经过思考和分辨的结果）必然被卷入更大罪恶的黑暗中。\par

% structure: p0010 source=h2 target=subsection reason=relative-heading-rank; base=h2

\subsection{第三章}

\Emph{关于\Person{撒乌耳}和\Person{阿哈布}的错误：他们因缺乏分辨而受骗。}\par

最后，在天主判断中第一个配得上祂子民以色列王国的人，因缺乏这分辨的\Quote{眼睛}，仿佛全身充满黑暗，实际上被从王国中推翻了；他被这\Quote{光}的黑暗所欺骗，并错误地以为自己的祭献比服从\Person{撒慕尔}的命令更蒙天主悦纳，就在他本希望借此平息天主尊威的事上遭遇了跌倒的机缘。\EditorialNote{撒慕尔纪上第十五章}。我要说，对分辨的无知又导致以色列王\Person{阿哈布}，在获得天主恩赐的一次大胜和辉煌胜利之后，以为自己仁慈比严格执行天主的命令（在他看来是残忍的规则）更好；受此念头驱使，他本想以仁慈调和一场血腥的胜利，却因这不加分辨的仁慈使全身充满黑暗，并被不可挽回地判处死亡。\EditorialNote{列王纪上第二十章}。\par

% structure: p0013 source=h2 target=subsection reason=relative-heading-rank; base=h2

\subsection{第四章}

\Emph{关于分辨在圣经中的价值所说的。}\par

這就是辨別力，它不僅是\Quote{身體的光}，也被宗徒稱爲太陽，正如所說的：\Quote{不要讓太陽在你們的忿怒上落下。}\BibleRef{弗 4:26}它也被稱爲我們生活的引導：正如所說的：\Quote{沒有引導的人，如落葉般墜落。}它最真實地被稱爲謀略，沒有它，聖經的權威不許我們做任何事，以至於我們甚至不允許在它的調控下取那使人心喜樂的屬神\Quote{酒}：正如所說的：\Quote{一切都要有謀略，喝你的酒要有謀略}，又說：\Quote{像一座城牆被摧毀、沒有圍牆的城，照樣，一個沒有謀略行事的人。}而缺少它對隱修士的損害，引文中的比喻和圖像已表明，將之比作一座被毀、無牆的城。智慧在於此，聰明和理解在於此，沒有它們，我們內在的房屋無法建造，屬神的財富也無法積聚，正如所說的：\Quote{房屋因智慧建造，又因聰明豎立。因理解，倉庫充滿各種寶貴的財物和美物。}我說這就是\Quote{固體食物}，只有那些長大成熟、強壯的人才能領受，正如所說的：\Quote{但固體食物是給長大成熟的人的，他們因習慣而鍛煉了感官，以分辨善惡。}\BibleRef{希 5:14}對我們有益且必要的是，只在於它符合天主聖言及其能力，正如所說的：\Quote{因爲天主聖言是活潑的、有效的，比任何雙刃劍更鋒利，甚至達到靈魂與精神、骨節與骨髓的分割，能辨別心中的思想和意圖。}\BibleRef{希 4:12}由此清楚表明，沒有辨別的恩寵，任何德行都不能完美地獲得或持續。因此，根據真福安當以及所有其他人的判斷，已確定：是辨別力帶領無畏的隱修士逐步走向天主，並使上述德行持續完整無缺；藉由它，人可以較不疲憊地攀登到完美的極高頂峯；沒有它，即使最樂意勞苦的人也無法達到完美的巔峯。因爲辨別力是諸德之母，也是它們的守護者和調節者。\par

% structure: p0016 source=h2 target=subsection reason=relative-heading-rank; base=h2

\subsection{第五章}

論老赫龍的死亡。\par

爲支持真福安當及其他教父古時所作的這項判斷，按我們所承諾的，用一個近代的例子，請回想最近在你們眼前發生的事，即老赫龍如何在幾天前被魔鬼的幻象從高處拋入深淵；我們記得他在這曠野中生活了五十年，以特別的嚴厲恪守嚴格的節制，並以驚人的熱忱追求獨居的隱祕，超過所有住在這裏的人。那麼，他是用什麼詭計或方法被那欺騙者迷惑，歷經如此多的勞苦後，以最慘重的墮落跌倒，使所有住在這曠野的人深陷悲痛？豈不是因爲他缺乏辨別的德行，寧願隨從自己的判斷，而不願聽從弟兄們的勸告、會談和長老的規矩？因他從未間斷地以如此嚴厲的方式修練禁食和齋戒，並如此恆常地堅持獨居和隱修小室，以至於連逾越節的慶祝都不能說服他與弟兄們一同赴宴：按年度的慣例，所有弟兄都留在教堂裏，而他唯獨不加入他們，恐怕因只吃少許豆類而顯出某種程度的鬆懈。因這驕傲受騙，他以極大的尊敬將撒殫的天使當作光明天使接待，盲目奴性地服從其命令，縱身跳入一個深得無法目測的井中，對那保證他德行與勞苦功績如此之大、以致不可能有任何危險的天使之許諾毫無懷疑。爲以自身安全試驗此事最確鑿無疑，他在深夜受騙跳入上述井中，以證明若無傷出來，其德行之偉大。弟兄們費了大力氣把他救出來時，他已幾乎死亡，三天後斷氣；更糟的是，他固執於其執迷，以至於連死亡的經驗也不能說服他，讓他知道被魔鬼的詭計欺騙了。因此，儘管他巨大的勞苦功績和在曠野中度過的歲月，那些以最大同情和仁慈憐憫他結局的人，幾乎無法從院長帕夫努提烏斯那裏爭取到，不將他視爲自殺者，也不認爲他不配獲得爲安息者舉行的紀念和祭獻。\par

% structure: p0019 source=h2 target=subsection reason=relative-heading-rank; base=h2

\subsection{第六章}

論兩位弟兄因缺乏辨別力而毀滅。\par

我还能说什么呢？关于那两位弟兄，他们住在底比斯荒漠那边，就是有福的\Person{安当}曾居住的地方。由于没有充分受到谨慎明辨的影响，当他们穿越那广阔无垠的荒野时，决定除了主自己可能供给他们的食物外，不带任何食物。当他们在荒漠中行走，已经因饥饿而昏厥时，被远处的\Place{玛齐斯人}（一个比几乎所有野蛮部落更加野蛮凶残的民族，因为他们不像其他部落那样因贪图掠夺而流血，而是出于纯粹的残暴心性）发现。他们一反其残暴天性，拿着面包迎接他们。其中一位弟兄因明辨的帮助，欢喜感恩地接受了面包，如同是主所赐，认为这食物是天主所供给的，并且是天主的作为，使那些惯于流血的人向已经昏厥垂死的人提供了生命之粮；但另一位却因食物是由人提供的而拒绝，结果饿死了。虽然这最初源于一种可责怪的信念，但其中一位靠着明辨克服了他鲁莽轻率产生的想法，而另一位坚持愚顽，完全缺乏明辨，给自己带来了死亡，而主本会阻止这死亡，因为他不相信那凶猛的野蛮人忘记其残暴天性、献上面包而非刀剑，乃是出于天主的感动。\par

% structure: p0022 source=h2 target=subsection reason=relative-heading-rank; base=h2

\subsection{第七章}

\Emph{论另一人因缺乏明辨而陷入的幻觉。}\par

为何我也不应提及一人（我们宁愿不提其名，因为他仍在世），他在很长一段时间里，接待了一位伪装成天使光明的魔鬼，并屡次被其无数启示所欺骗，相信他是光明的使者；因为每当这些启示赐下时，魔鬼每晚都在他的小室里提供光亮，无需任何灯盏。最后，魔鬼命令他将与他同住在隐修院的亲生儿子献给天主，以使他的功德借此祭献与圣祖\Person{亚巴郎}相等。他竟被其说服到如此地步，几乎真要犯下谋杀，若非他儿子看见他异常仔细地磨刀，寻找用来捆绑儿子祭献的链子，预感即将发生的罪而惊恐地逃走。\par

% structure: p0025 source=h2 target=subsection reason=relative-heading-rank; base=h2

\subsection{第八章}

\Emph{论美索不达米亚一位隐修士的堕落与欺骗。}\par

讲述那位\Place{美索不达米亚}隐修士受欺骗的经过也很费时。他在该国实行了几乎无人能效仿的克己，多年隐藏在小室里修行，最终却因魔鬼所赐的启示和梦境而被欺骗，在经过如此多的劳苦和善功，超胜了所有同地区的修道人之后，竟悲惨地堕入犹太教和肉身的割礼。因为魔鬼通过使他习惯于异象，意在最终诱使他相信谎言，如同真理的使者，长期向他揭示完全真实的事；最后，它向他显示了基督徒群众，连同我们宗教与信经的首领，即宗徒和殉道者，处在黑暗与污秽中，被一切肮脏所玷污；另一方面，犹太人民同\Person{梅瑟}、圣祖和先知们，充满喜乐地舞蹈，闪耀着夺目的光芒。于是魔鬼说服他，若他想分享他们的赏报与福乐，就必须立即接受割礼。因此，这些人之所以如此悲惨地被欺骗，是因为他们没有努力获得明辨的能力。可见，许多人的不幸和考验表明，缺乏明辨的恩宠是何等危险。\par

% structure: p0028 source=h2 target=subsection reason=relative-heading-rank; base=h2

\subsection{第九章}

\Emph{关于获取真正明辨的问题。}\par

对此，\Person{杰尔玛诺}说：通过最近的实例和古人的决断，已经充分完全地表明了明辨在某种意义上是一切德行的源头和根基。那么，我们想学习如何获得它，或者我们如何分辨它是真正的、出自天主的，还是虚假的、出自魔鬼的；这样（借用你上次讨论的福音比喻中的形象，我们被命令成为良好的兑换银钱者），我们就能看到真君王的像印在钱币上，并识破非真正流通的钱币，并且，如你在昨日谈话中引用的一句普通说法，我们凭借你所充分详尽论述的那种技艺的教导，将其当作伪币拒绝；这技艺理应属于那在福音意义上属灵的、良好的兑换银钱者。因为，若我们不知道如何寻求并获得那德行和恩宠，认识其价值又有何益呢？\par

% structure: p0031 source=h2 target=subsection reason=relative-heading-rank; base=h2

\subsection{第十章}

\Emph{回答如何获得真正的明辨。}\par

于是\Person{梅瑟}说：\Quote{真正的明智只有靠真正的谦卑才能获得。而谦卑的第一个证明，就是将所有的一切（不仅是所做，还有所想）都交由长者审查，绝不信任自己的判断，而是在所有方面顺从他们的决定，并承认什么应凭他们的传承视之为善或恶。这种习惯不仅会教年轻人通过明智的真道行走在正路上，还会使他免受仇敌一切诡计和欺骗的伤害。因为一个人若不凭自己的判断生活，而是依照长者的榜样，就绝不可能受骗；我们那狡猾的仇敌也无法滥用一个人的无知，只要他习惯于不因虚假的谦逊而隐瞒心中升起的所有念头，而是根据长者成熟的判断来抑制或容许它们。因为一个错误的念头在它被发现的那一刻就削弱了；甚至在明智的判断下达之前，那恶蛇就已告明的力量从阴暗的地下洞穴中拖出，可以说被暴露并羞耻地驱逐。因为邪恶的念头只要藏在心中，就会在我们内掌权；为了让你更有效地领会这种判断的力量，我将告诉你\Person{塞拉彼雍}院长做了什么，以及他常为年轻弟兄的教训而向他们讲述的事。}\par

% structure: p0034 source=h2 target=subsection reason=relative-heading-rank; base=h2

\subsection{第11章}

\Emph{\Person{塞拉彼雍}院长关于向他人吐露的念头衰退的言论，以及关于自负的危险。}\par

他说：\Quote{当我还是个少年，与院长\Person{德奥纳斯}同住时，敌人攻击迫使我养成一个习惯：我与那位长者在第九时辰一同用晚餐后，每天偷偷在衣服里藏一块饼，后来背着他悄悄吃掉。尽管我不断因意志的许可而有罪于偷窃，并因根深柢固的渴望而缺乏约束，但当我的非法欲望得到满足后，我便会清醒过来，为所犯的偷窃而折磨自己，其程度超过享用那饼的快乐。当我日复一日、心中不无悲伤地被迫履行这个可说是法老监工所要求的沉重任务（以取代砖块），无法逃脱这残酷的暴政，又羞于向长者透露这秘密偷窃时，由于天主的旨意，恰巧有几个弟兄来到长者的房间，为了接受他的教导，我得以从这自愿的束缚中得到解放。晚餐后，当开始举行神修会议，长者回答他们提出的问题，谈论贪饕之罪和隐秘思想的统治，并展示它们的性质以及只要保持隐秘就具有的可怕力量时，我被那番话的力量所征服，良心刺痛，惊恐万分，因为我认为这些事是他提到的，是因为主向长者揭示了我心中的秘密；首先我暗自叹息，然后我的心痛悔加深，我公然放声哭泣落泪，从我那与之同谋并接收偷窃的衣褶中，拿出那块我因坏习惯而带去偷偷吃掉的饼；我把它放在中间，俯伏在地，恳求宽恕，承认我每天如何偷偷吃掉一块，并痛哭流涕恳求他们祈求主使我摆脱这可怕的奴役。然后长者说：\Quote{孩子，要有信心，}他说，\Quote{不用我多言，你的告解已使你摆脱这奴役。因为今天你战胜了你那胜利的对手，以你的告解将他打倒，这方式足以弥补你先前因沉默而被击倒的损失，正如你从未以自己或别人的回答驳斥他，任由他凌驾于你之上，正如\Person{撒罗满}所说的：\BibleQuote{因为对恶事没有迅速执行判决，世人的心就充满行恶的念头。}因此，在揭露他之后，那恶神再不能烦扰你，那污秽的蛇今后也不能在你内作窝，因为他已被你生命之言的告解从黑暗中拖到光明中。}长者话未说完，看哪！一盏燃烧的灯从我衣褶中出来，充满房间，散发出硫磺气味，其刺鼻的臭味几乎使我们无法待在那里；长者继续他的告诫说：\Quote{看哪！主已向你明显证实了我话语的真实，使你亲眼看见，那造成祂苦难的元凶如何因你生命之言的告解被赶出你的心，并知道那已被揭露的敌人必定不再在你内找到住所，因为他的被逐已显明出来。}因此，正如长者所宣告的，他说，那魔鬼专制统辖对我的统治已被这告解的力量摧毁，永远平息，以致敌人甚至从未试图再向我强加这欲望的回忆，我也从未感觉被那偷偷渴望的激情所抓住。我们看到这个含义在\WorkTitle{训道篇}中用一个比喻巧妙地表达出来。他说：\BibleQuote{如果蛇不嘶嘶声就咬人，对行咒术的人就一无益处。}这表明沉默的蛇咬是危险的，即如果一个来自魔鬼的暗示或思想没有通过告解展示给某个行咒术的人——我是指某个有属神心志的人，知道如何立即用圣经的咒语治愈伤口，并从心中取出蛇的致命毒液——那么就不可能帮助那已经处于危险中且必死的受苦者。因此，通过这种方式，我们将轻易获得真正分辨的知识，以便跟随长老们的足迹，不做任何新奇之事，也不凭自己或自己负责决定任何事，而是在一切事上按照他们的传统和正直生活所教导我们的行事。那被这体系坚固的人，不仅达到完美的分辨方法，而且也将完全安全地避开敌人所有的诡计：因为魔鬼没有别的过失比这更快地拖垮一位隐修士并引他走向死亡，即当他劝诱他轻视长老们的建议，而依靠自己的意见和判断时。因为，如果人巧思所发现的一切技艺和发明，那些仅对今世生活之方便有用的，虽可手触目见，却无人能不通过师长的教训而理解，那么，认为唯独在这不可见、隐秘、只有最纯洁之心才能见的事上不需要导师，是何等愚昧？此错误带来的不是暂时或易于弥补的损失，而是灵魂的毁灭和永死。因为这涉及日夜与不可见、残忍的敌人（而非可见的）作战，是属灵的争战，不仅对抗一两个，而是对抗无数的军队，其失败对所有人来说更危险，因为敌人越凶猛，攻击越隐秘。因此，我们应当始终极其谨慎地跟随长老们的足迹，将心中升起的一切带到他们面前，除去羞愧的面纱。}\par

% structure: p0037 source=h2 target=subsection reason=relative-heading-rank; base=h2

\subsection{第十二章}

\Emph{因谦逊而羞愧于向长老透露思想的告解。}\par

\Person{格尔玛诺}说：这种有害的羞怯，使我们试图隐藏不良思想，其根源尤其在于：我们听说在\Place{叙利亚}地区有一位长老的长上，正如人们所相信的，当一位弟兄向他真诚地宣认、袒露自己的思想后，他后来对此极为愤怒并严厉责备了他。由此导致我们压抑这些思想，羞于向长老们告知，因而无法获得治愈它们的疗法。\par

% structure: p0040 source=h2 target=subsection reason=relative-heading-rank; base=h2

\subsection{第十三章}

\Emph{论践踏羞耻，以及无痛悔者之危险。}\par

梅瑟说：\Quote{并非所有青年人在精神热忱上都相同，在学问和德行上所受的教导也不相等；同样，我们也无法发现所有老年人在完美和卓越上完全一致。因为老年人的真财富不是以白发来衡量，而是以他们青年时的勤奋和过去劳苦的报酬来定。因为经上说：\InnerQuote{你幼年没有收集的，到老年怎能找到呢？}\InnerQuote{可敬的老年不在长寿，也不以年岁来计算；人的智慧才是白发，无瑕的生命才是高寿。}因此，我们不应追随每一位满头白发、仅凭年龄而受尊敬的老年人的步伐，或者采纳他们的传统和劝告，而只应追随那些我们在青年时已发现他们以可嘉许和值得称赞的方式出众，并且不是凭自信而是凭长老们的传统所培育的人。因为有些人（不幸的是他们占多数）在青年时沾染的冷淡和怠惰中度过晚年，因此他们的权威不是来自品格的成熟，而仅仅来自年岁的数量。先知特别针对这些人发出上主的斥责：\InnerQuote{外邦人吞吃了他的力量，他却不知道；白发也已布满在他身上，他却不自知。}\BibleRef{欧 7:9} 我说这些人，不是因他们生活的正直或修道的严谨——这些都是值得赞美和效仿的——而被指给青年作榜样，而仅仅是因为他们的年岁；因此狡猾的仇敌利用他们的白发来欺骗青年，错误地诉诸他们的权威，并狡猾地企图以他们的榜样来动摇和欺骗那些可能因他们或其他人的劝告而被推向完美之路的人；并藉着他们的教导和实践，将他们拖入有害的冷漠或致命的绝望之中。既然我想给你们一个实例，我将告诉你们一个事实，这事实可以供给我们一些有益的教训，但不提当事人的名字，以免我们犯下与那个泄露了向他告明的弟兄罪过的人相似的错误。那时，有一个并非最懒惰的青年，为了灵魂的益处和健康，去拜访一位我们都很熟悉的老人，并坦诚地承认自己受肉欲和行淫之神的困扰，以为会从老人的话语中得到对他努力的安慰，以及对所受创伤的医治。然而，那位老人却以最严厉的斥责攻击他，称他为可怜又可耻的人，不配称为隐修士，因为他竟被这样的罪和欲念所影响；老人非但没有帮助他，反而以斥责伤害他，使他怀着绝望和致命的沮丧离开了他的小室。当这位青年被如此忧伤压垮，陷入沉思，不再思考如何治好他的情欲，而是如何满足他的情欲时，阿颇罗院长——长老中最有经验的一位——遇见了他，从他的神情和忧郁中看出他的烦恼以及他心中暗自翻腾的猛烈攻击，便问他这种不安的原因；当他无法回答老人温和的询问时，老人更加清楚地看出他并非无故想以沉默隐藏如此深切忧伤的原因，于是更急切地问他隐藏悲痛的原因。于是他被逼无奈，承认他正要去一个村庄娶妻，离开修院，回到世俗，因为那位老人告诉他，如果他不能控制肉欲并治好他的情欲，他就不能成为隐修士。于是这位老人以善意的安慰安抚他，告诉他，他自己每天也受着同样的欲望和情欲的刺伤，因此他不应屈服于绝望，也不应对猛烈的攻击感到惊讶；他克胜这攻击，与其说是靠热心的努力，不如说是靠主的仁慈和恩宠；老人恳求他至少推迟一天他的打算，并求他返回他的小室，然后他尽快赶到上述那位老人的修院——走近他时，他伸手流泪祈祷说：\InnerQuote{主啊，惟有你是正义的审判者，是隐秘力量和人性软弱的无形医生，请将这攻击从青年转向老人，让他学会体恤受苦者的软弱，甚至在老年也能同情青年的脆弱。}当他含泪结束祈祷时，他看到一个污秽的埃塞俄比亚人站在那老人的小室对面，向他射出火矢；那老人立刻被射中，走出小室，像疯子或醉汉一样到处乱跑，进进出出，再也不能自制，开始朝着那个青年所去的相同方向匆忙奔去。当阿颇罗院长看到他像被愤怒驱使的疯子一样时，他知道他所看见的魔鬼的火矢已经刺入那老人的心，并以难以忍受的热度在他心中造成这种精神错乱和理解力的混乱；于是他走到他跟前问：\InnerQuote{你急着去哪里？是什么使你忘记了年岁的尊贵，这样幼稚地不安，如此匆忙地奔走？}}\par

他因良心有愧，且被这可耻的兴奋所迷惑，幻想自己心中的情欲已被发现；又因内心的秘密被长老所知，不敢对长老的询问作出任何回答。长老说：\Quote{回到你的静修室去，最终要认清这个事实：直到现在，你一直被魔鬼忽视或鄙视，没有被算在那些每日被激起与魔鬼作战、与他们的努力和热忱搏斗的人之中——你，我不能说不能抵挡，甚至不能推迟一天，他所射向你的那支箭，在你从事这修道职业这么多年之后。}主容许你受此创伤，为的是让你至少在你年老时学会同情你所不熟悉的软弱，并藉你自己的情况和经验，知道如何俯就年轻人的脆弱；然而，当你接待一个受魔鬼攻击困扰的年轻人时，你并没有用任何安慰鼓励他，反而在沮丧和毁灭性的绝望中，把他交到仇敌手中，就你而言，让他被仇敌悲惨地毁灭。但仇敌绝不会以如此猛烈的攻击攻击他——仇敌至今不屑于攻击你——除非是嫉妒他将要取得的进步，它竭力要战胜他本性中的那种德行，用它的火箭摧毁它，因为它毫不怀疑地知道他是更有力量的，既然它认为值得如此猛烈地攻击他。所以，要藉你自己的经验，学会同情那些受苦的人，绝不要用毁灭性的绝望去恐吓那些处于危险中的人，也不要用严厉的言语使他们刚硬，反而要用温柔和蔼的安慰去扶助他们；正如明智的\Person{撒罗满}所说：\Quote{不可不肯拯救那些被拉去处死的人，也不可放弃那些将被杀害的人。}\BibleRef{箴 24:11}并效法我们的救主，不要折断已压伤的芦苇，不要熄灭将残的灯芯，\BibleRef{玛 12:20}并向主祈求那份恩宠，藉着它，你可以在行动和力量上忠实地学会歌唱：\Quote{主赐给我受教的舌头，使我知道怎样用言语扶助疲倦的人。}\BibleRef{依 50:4}因为没有人能忍受仇敌的诡计，或扑灭、压制那些像天然火焰燃烧的肉体情欲，除非天主的恩宠扶助我们的软弱，或保护并支持它。因此，既然这场有益事件的原因已经结束——主藉此要使那年轻人脱离危险的情欲，并教你一些关于他们攻击的猛烈和同情心的事——让我们一同在祈祷中恳求祂，愿祂乐意除去那根鞭子，就是主为你的益处加在你身上的（因为\Quote{祂使人忧愁，也使人痊愈；祂击打，祂的手也医治。祂使人卑微，也使人高升；祂使人死，也使人活；祂使人下降阴府，也使人上来}），并愿祂以圣神丰沛的露水熄灭魔鬼的火箭，这些火箭是祂因我的渴望而容许射伤你的。虽然主因长老的一次祈祷就除去了这诱惑，其速度与祂容许它临到他的速度一样快，但祂藉此清楚证明：一个人的过失被揭露后，不仅不应受责骂，而且受苦者的忧伤也不应被轻忽。因此，绝不要让某一位或几位长老的笨拙或肤浅阻止你，使你离开我们之前提到的生命之道，或离开先贤的传统，即使我们狡猾的仇敌不正当地利用他们的白发来欺骗年轻人；相反，应该毫无羞耻地将一切揭露给长老，并忠实地从他们那里接受伤口的医治，连同生活与言谈的榜样；如果我们完全不凭自己的责任和判断行事，我们就会从其中找到同样的帮助和同样的结果。\par

% structure: p0044 source=h2 target=subsection reason=relative-heading-rank; base=h2

\subsection{第十四章}

\Emph{关于\Person{撒慕尔}的蒙召。}\par

最后，这种看法如此悦乐天主，以至于我们看到这一制度在圣经中并非毫无根据：主不愿亲自以神圣对话的方式教导少年\Person{撒慕尔}（他被拣选担任审判职务），却容许他一两次跑到长老那里；主愿意那位召唤他与之交谈的人，甚至由一位曾得罪天主的人（因他是长老）来教导；主宁愿那蒙祂召唤的人由长老来训练，以考验那蒙召担任神圣职务者的谦卑，并为年轻人树立服从的榜样。\par

% structure: p0047 source=h2 target=subsection reason=relative-heading-rank; base=h2

\subsection{第十五章}

\Emph{关于宗徒\Person{保禄}的蒙召。}\par

当基督亲自呼唤并称呼保禄时，虽然祂本可立即向他揭示成全之道，却宁愿差遣他到阿纳尼雅那里去，并命令他向阿纳尼雅学习真理之道，说：\Quote{起来进城去，在那里必有人告诉你你所当作的事。} \BibleRef{宗 9:6} 祂就这样把保禄送到一位年长者那里，并认为让他受年长者的教导比受祂自己的教导更好，免得在保禄身上可行的事，若各人自以为也该单受天主的训导和教诲，而非受长老们的指导，就会树立自满自足的坏榜样。宗徒本人不仅以书信，也以言行教导我们应当极力避免这种自满自足，正如他所说，他上耶路撒冷唯一的目的就是：私下与同宗徒及在他以前的那些人交流他所传给外邦人的福音——圣神的恩宠以大能神迹奇事伴随着他——如他所说：\Quote{我把我在外邦中所传的福音和他们交流了，免得我白白奔跑或徒然奔跑。} \BibleRef{迦 2:2} 既然这蒙选的器皿也承认自己需要与同宗徒交流，那还有谁如此自满自足和盲目，竟敢信赖自己的判断和智慧呢？由此我们清楚看到，主并不亲自向任何人显示成全之道，若这人有机会学习却轻视长老们的教导和培育，毫不留意那应当最谨慎遵守的话：\Quote{问你的父亲，他必指示你；问你的长老，他们必告诉你。} \BibleRef{申 32:7}\par

% structure: p0050 source=h2 target=subsection reason=relative-heading-rank; base=h2

\subsection{第十六章}

如何寻求辨别力。\par

那么，我们应当尽全力借着谦逊的德能追求辨别之德，因为它能保护我们不受任何极端之害，因有古语说：\OriginalTerm{极端相同}{\GreekText{ἀκρότητες} \GreekText{ἰσότητες}}，即极端相遇。因为过度的斋戒与贪食结果相同，无度地持续守夜对隐修士的损害与沉睡的麻木一样：因为当人因过度节制而衰弱时，他必然会回到那种因疏忽和怠惰而保持的状态，以致我们常见那些未被贪食欺骗的人，却因过度斋戒而毁掉，并因虚弱而陷入他们先前已克服的情欲。不合理的守夜和夜间警醒也毁掉了一些未被睡眠胜过的人。因此，正如宗徒所说：\Quote{藉着左右两手的义德武器，} \BibleRef{格后 6:7} 我们在辨别之德的引导下，以适当的节制前行，行走于两极端之间，以免被诱离为我们划定的节制之途，也不至因不当的疏忽而陷入口腹之乐。\par

% structure: p0053 source=h2 target=subsection reason=relative-heading-rank; base=h2

\subsection{第十七章}

论过度的斋戒与守夜。\par

因为我记得自己曾如此频繁地抗拒对食物的欲望，以至于在禁食两三天后，我的思想甚至不为任何食物的记忆所困扰；而且，因魔鬼的攻击，睡眠也如此远离我的眼睛，以至于连续几天几夜，我常祈求主赐给我双眼一点睡眠。那时我感到，因缺乏食物和睡眠而面临的危险，比与懒惰和贪食的斗争更大。因此，正如我们应当小心，以免因贪图肉身享乐而陷入危险的柔弱，也不在正当时间之前进食或饮食过多；同样，我们应当在适当的时候进食和睡眠，即使我们不喜欢。因为每一次斗争都是由仇敌的诡计引起的；过度的节制比粗心的饱足对我们伤害更大，因为从后者，健康的痛悔能提升我们达到正确的严格尺度，而不能从前者。\par

% structure: p0056 source=h2 target=subsection reason=relative-heading-rank; base=h2

\subsection{第十八章}

关于节制与休息之正当尺度的提问。\par

日耳曼努斯说：那么，保持何种节制的尺度，使我们能以平衡的方式成功地在两极端之间无害地通行呢？\par

% structure: p0059 source=h2 target=subsection reason=relative-heading-rank; base=h2

\subsection{第十九章}

论每日食物的最佳计划。\par

梅瑟说：关于此事，我们知道在我们长老们中间常有讨论。因为当讨论某些隐修士的节制——他们连续以豆子、或仅以蔬菜和水果维持生命时，他们向众人提议仅用面包，其正当尺度定为两块小薄饼，小到几乎不足一磅重。\par

% structure: p0062 source=h2 target=subsection reason=relative-heading-rank; base=h2

\subsection{第二十章}

关于以两块薄饼维持生命的这种节制似乎容易的异议。\par

我们欣然接受这一点，并回答说我们几乎不能认为这种限度是节制，因为我们完全不可能达到它。\par

% structure: p0065 source=h2 target=subsection reason=relative-heading-rank; base=h2

\subsection{第二十一章}

关于经得起考验的节制的价值和尺度的回答。\par

梅瑟说：若你们想检验这条规则的力量，就持续遵守这个限度，即使在主日或安息日，或任何弟兄到来的时候，也从不偏离去吃任何煮熟的食物；因为肉身因这些例外而得到振奋，不仅能在其余日子以更少量的食物维持自己，而且能毫无困难地推迟一切休息，因为它靠超出限度的食物补充而得到支持；而那位总是按上述分量吃饱的人，绝不能做到这一点，也不能将开斋推迟到次日。因为我记得我们的长老们（我也记得我们自己常有同样的经验）认为实践这种节制非常艰难、痛苦，并以如此痛苦和饥饿遵守所定的规则，几乎是不情愿地、含着眼泪和哀叹，才为他们的膳食设定了这个限度。\par

% structure: p0068 source=h2 target=subsection reason=relative-heading-rank; base=h2

\subsection{第二十二章}

通常的节制和进食的限度是什么。\par

但这是通常的节制限度：即每个人根据其体力、体质或年龄的需要，允许自己进食，以维持肉身所需为度，而非为了满足饱胀感。因为若人没有固定规则，时而以寡淡食物和禁食折磨胃腹，时而又暴饮暴食，将遭受极大的损害；正如因缺乏食物而虚弱的心灵在祈祷时失去力量，因肉身过度虚弱而疲惫不堪、被迫打盹一样，同样，当人被暴食压垮时，也无法向天主献上纯洁而自由的祈祷；也无法保持贞洁的纯净不间断，因为在那些似乎以更严厉的禁食来惩罚肉身的日子里，它也用已摄入的食物之燃料点燃了肉欲之火。\par

% structure: p0071 source=h2 target=subsection reason=relative-heading-rank; base=h2

\subsection{第二十三章}

\Emph{如何克制生殖体液之多}\par

因为一旦因食物丰盛而在骨髓中凝结的东西，必然被排出，并被自然法则本身所驱逐；自然法则不容许任何多余体液的过剩像有害和敌对之物那样存留于自身之内，因此我们应以合理且均衡的俭朴来克制我们的身体，以便当我们居住于肉身中而完全无法摆脱这一自然需要时，至少一年之中让我们较少地、不超过三次被这种污秽沾染；而这是由安眠在无任何瘙痒下排出的，而非由隐秘快感的虚假幻象所引发。\par

因此，这就是我们所谈论的适度和均衡的节制量与尺度，它也获得了教父们判断的认可；即每日的饥饿应与每日的餐食相伴，使身体和灵魂保持在同一个状态，不让心灵因禁食之劳累而昏沉，也不被暴食所压迫；最终达到这样一种节俭的饮食，以至有时人在晚上既不注意也不记得他已开斋。\par

% structure: p0075 source=h2 target=subsection reason=relative-heading-rank; base=h2

\subsection{第二十四章}

\Emph{论饮食一致性的困难；以及本雅明兄弟的贪食}\par

这并非毫无困难地做到，以致那些对完美明辨一无所知的人宁愿将禁食延长两天，把今天该吃的留到明天，以便在进食时能随心所欲地享受。而且你知道，最后你的同乡\Person{本雅明}最顽固地坚持这一点：他不愿每天吃两块饼干，也不愿以一贯的自律持续吃他那寡淡的食物，却宁愿总是将禁食延长两天，以便在进食时用双份填满贪婪的胃，通过吃四块饼干享受舒服的饱胀感，并借两天禁食填饱肚子。你无疑记得这人的结局如何——他顽固而执拗地依赖自己的判断而非长老的传统，离弃旷野，回到这世界的虚妄哲学和世俗虚荣中，从而以自己的堕落证实了上述长老的观点，并以他的毁灭教导人：凡信赖自己意见和判断的人，绝不可能攀登完美的顶峰，也必不免被魔鬼危险的诡计所欺骗。\par

% structure: p0078 source=h2 target=subsection reason=relative-heading-rank; base=h2

\subsection{第二十五章}

\Emph{一个问题：如何可能始终遵守同一尺度}\par

\Person{杰尔玛努斯}：那么我们如何能从不违背地遵守这个尺度呢？因为在第九时辰守斋结束后，有时弟兄们来看我们，此时我们必须或为他们的缘故在我们固定的惯常份量上增加一些，或肯定在那我们被告知应向所有人显示的礼貌上有亏欠。\par

% structure: p0081 source=h2 target=subsection reason=relative-heading-rank; base=h2

\subsection{第二十六章}

\Emph{回答：我们应如何不超出适当的食物分量}\par

\Person{梅瑟}：两种义务必须以同样方式、同样谨慎地遵守：因为我们应该极其严格地为我们的节制保留适当的食物份量，同样，出于爱德，应向任何到来的弟兄显示礼貌和鼓励；因为当你向弟兄，甚至向基督自己提供食物时，却不与他一同进食，而让自己与他的餐食隔绝，这绝对荒谬。因此如果我们遵守这个计划，就能避免两方面的过错：即在第九时辰吃两块饼干中的一块——这两块是我们规定的正当份量——，另一块留到晚上，以备类似情况；如果有弟兄来看我们，我们可以与他一起吃那块，这样就不增加我们惯常的份量；通过这一安排，弟兄的到来本应使我们喜悦，却不会给我们带来不便：因为我们将以礼貌所要求的殷勤待他，同时丝毫不放松我们节制的严格。但如果没有人来，我们就可以自由地吃这最后一块饼干，因为它根据我们的会规属于我们；并且由于我们如此节俭（因为第九时辰只吃了一块饼干），晚上胃就不会负担过重——这是那些自以为遵守更严格禁食而将所有餐食推迟到晚上的人常遇到的情况；因为刚吃了东西，会妨碍我们的理智在晚祷和夜祷中保持清醒敏锐；因此在第九时辰，已为食物留出了合适的时间，此时隐修士可以恢复体力，从而不仅在夜间守夜时精神饱满、清醒，而且因为食物已经消化，也能完全准备好进行晚祷。\par

圣\Person{梅瑟}以这有两道菜的筵席，可以说，款待了我们，他不仅藉着当前博学的言谈，向我们展示了分辨的恩宠与能力，而且藉着先前进行的讨论，展示了舍弃的方法以及隐修生活的终向与目的；这样做，使我们以往仅凭精神的热忱与对天主的热情，却闭着眼睛所追求的一切，变得比白昼还要清楚，并使我们感觉到，我们那时已远离了心地纯洁和我们行事的直线有多远，因为所有属于此生的可见技艺的实践，若无对其目的的理解，便不可能存在，也无法在没有明确终向的清晰视界下着手进行。\par

\SourceRecord{Translated by C.S. Gibson.；Nicene and Post-Nicene Fathers, Second Series；Vol. 11.；Edited by Philip Schaff and Henry Wace.；Buffalo, NY: Christian Literature Publishing Co.,；1894.；<http://www.newadvent.org/fathers/350802.htm>.}

% structure: p0001 source=h1 target=section reason=page-title

\section{第三场会议}

\Signature{圣若望•卡西安}\par

\Person{帕夫努提乌斯}院长的会议。论\OriginalTerm{三种弃绝}{three sorts of renunciations}。\par

% structure: p0004 source=h2 target=subsection reason=relative-heading-rank; base=h2

\subsection{第一章}

\Emph{论帕夫努提乌斯院长的生活与品行。}\par

在那如同璀璨星辰照亮此世黑夜的圣人合唱中，我们看见了圣帕夫努提乌斯，他如同伟大的发光体，以知识的辉煌照耀。因为他曾是我们团体的\OriginalTerm{司铎}{presbyter}，我的意思是，那些居住在\Place{斯刻特}旷野中的人之一；他在那里活到极老，从未离开过他的斗室——他年轻时便占据了这斗室，距离教堂五英里，甚至更近的区域；即使年老力衰，距离也未曾阻碍他在星期六或星期日去教堂。但他不愿意空手而归，便扛一桶水在肩上，足够他使用一周，带回他的斗室；即使年过九十，他也不容忍由更年轻者的劳动去取水。他从极幼年时便以如此热忱投身于隐修戒律，以至于只过了很短时间，就已被赋予顺服之德以及一切美善的知识。因为通过谦逊和服从的实践，他抑制了自己的一切欲望，并由此根除了一切过失，获得了隐修体制和古代教父教导所能产生的每一种德行；他因渴望更进一步而燃烧，急切地要深入旷野的深处，以便没有人类同伴打扰他，可以更轻易地与主结合——他渴望即使仍在弟兄团体中生活，就已与他不可分离地结合。而他在那里又一次以过分的热情超越了\OriginalTerm{独修隐士}{Anchorites}的德行，并因对持续神圣默想的渴望而回避与他们会面；他投身于更加荒野和难以接近的孤寂之地，长时间隐藏其中，以至于连独修隐士自己也只是偶尔很难瞥见他，于是人们相信他享受并喜悦于天使的日常陪伴；因这一显著特征，他们给他起了绰号叫\OriginalTerm{水牛}{the Buffalo}。\par

% structure: p0007 source=h2 target=subsection reason=relative-heading-rank; base=h2

\subsection{第二章}

\Emph{论同一位老者的谈话及我们的回答。}\par

我们因此渴望从教导中学习，于是在傍晚时分有些激动地来到他的斗室。短暂沉默后，他开始称赞我们的行动，因为我们为了主的爱离开了家乡，游历了如此多国家，并尽全力忍受旷野的匮乏和考验，效仿他们艰苦的生活——甚至那些生来就处于同样匮乏和贫困状态的人也几乎无法忍受；我们回答说，我们前来是为了他的教导和指教，以便能在某种程度上领受如此伟人的习俗，以及我们已从许多证据得知存在于他身上的那种成全，而不是为了因任何我们无权得到的称赞而受尊崇，或因他的话而心中骄傲（我们有时在自己的斗室里因敌人的暗示而受到这样试探）。因此我们恳请他，宁可向我们呈现能让我们谦卑痛悔的事，而不是奉承和膨胀我们的事。\par

% structure: p0010 source=h2 target=subsection reason=relative-heading-rank; base=h2

\subsection{第三章}

\Emph{\Person{帕夫努提乌斯}院长论\OriginalTerm{三种召叫}{three kinds of vocations}及\OriginalTerm{三种弃绝}{three sorts of renunciations}的陈述。}\par

然后\Person{真福帕夫努提乌斯}说道：有三种召叫。我们也知道有三种弃绝，这对于一个隐修士，无论其召叫为何，都是必要的。我们应当勤勉地首先检查我们为何说有三种召叫的理由，以便当我们确信我们在召叫的第一阶段被召唤到天主的服务中时，我们能够注意我们的生活与我们被召唤的崇高高度相称，因为有一个好的开始而无相应的终结是无用的。但如果我们感到只是在万不得已的情况下才从世俗生活中被拖出，那么，由于我们在宗教方面似乎立足于不太令人满意的开始，我们就必须相应地加倍努力，以属灵的热忱激励自己达到更好的终结。我们也很有理由知道这三重弃绝的方式，因为如果我们不知道它，或者知道它却不尝试在实践中实行，我们就永远无法达到成全。\par

% structure: p0013 source=h2 target=subsection reason=relative-heading-rank; base=h2

\subsection{第四章}

\Emph{对\OriginalTerm{三种召叫}{three callings}的解释。}\par

因此，为阐明这三类\OriginalTerm{召叫}{calling}的主要区别：第一种来自天主，第二种透过人而来，第三种来自\OriginalTerm{强制}{compulsion}。当某种灵感占据了我们的心，甚至在我们睡眠时，也在我们内激起对永生和救恩的渴望，并命令我们跟随天主、以赋予生命的\OriginalTerm{痛悔}{contrition}紧守祂的诫命时，这种召叫便是来自天主的；正如我们在圣经中读到，\Person{亚巴郎}被主的声音从故乡、从所有亲爱的亲友和父家中召叫出来；当时主说：\BibleQuote{离开你的故乡、你的亲族和你父家。}\BibleRef{创 12:1} 我们也听说，真福\Person{安多尼}也是这样被召叫的，他\OriginalTerm{归化}{conversion}的契机唯独来自天主。因为他进入一座圣堂时，在那里从福音中听到主说：\BibleQuote{凡不恼恨父亲、母亲、子女、妻子、田地，甚至自己的性命，就不能做我的门徒}；又说：\BibleQuote{你若愿意成全，去变卖你所有的一切，施舍给穷人，你将有宝藏在天上，然后来跟随我。}他以由衷的痛悔领受了主的这道命令，仿佛是特别针对他的，于是立即舍弃一切跟随了基督，丝毫未受人的劝言和教导所激励。第二种召叫，我们说是透过人而发生的；即当我们被某位圣人的榜样和劝言所激励，从而燃起对救恩的渴望时；借此我们永不会忘记，我们自己也因上述圣人的劝言和善表，蒙主的恩宠被召，将自己交托给这一目标和召叫。同样，我们在圣经中发现，以色列子民正是透过\Person{梅瑟}才脱离埃及的奴役的。但第三种召叫来自强制，即当我们沉迷于此生的财富和享乐，突然有诱惑临到我们，或威胁我们死亡的危险，或打击我们财产的损失和没收，或使我们失去亲爱的亲人，因而最终连我们也不情愿地被驱使转向天主——我们在富裕的日子里曾轻慢祂。这种强制性的召叫，我们在圣经中常见例子：我们读到，因着他们的罪恶，主将以色列子民交在敌人手中；又因敌人的暴虐和残忍，他们回转过来，向主呼求。圣经说：\BibleQuote{主给他们派遣了一位救主，名叫\Person{厄胡得}，是\Person{革辣}的儿子，\Person{本雅明}的儿子，左右手都能使用。}又说：\BibleQuote{他们呼求上主，上主就给他们兴起一位救主，即\Person{加肋布}的幼弟\Person{刻纳次}的儿子\Person{敖特尼耳}，拯救了他们。}圣咏也论及这类事说：\BibleQuote{祂击杀他们的时候，他们就来寻求祂；早晨回到祂那里，记起天主是他们的助佑，至高天主是他们的救赎者。}又说：\BibleQuote{他们在患难中呼求上主，祂就拯救他们脱离困厄。}\par

% structure: p0016 source=h2 target=subsection reason=relative-heading-rank; base=h2

\subsection{第五章}

\Emph{这第一种召叫对懒惰者无益，而最后一种对认真者却无阻碍。}\par

论到这三种召叫，虽然前两种似乎基于更好的原则，但有时我们发现，即使第三种等级——看似最低等、最冷淡的——也使人在心灵上变得成全、最为热忱，与那些在事奉主上有了令人钦佩的开端、并在余下生命中保持最可称赞的心灵热忱的人相似；反之，我们也发现，从较高等级中，许多人冷却下来，常常结局悲惨。正如对前一类人，他们似乎不是出于自由意志，而是出于暴力和强制而悔改，这并未成为阻碍，因为主的慈爱为他们保证了悔改的机会；同样，对后一类人，他们悔改的最初日子如此光明，却毫无益处，因为他们没有小心地将余生带到合宜的结局。例如院长\Person{梅瑟}，他曾住在旷野中一个叫\Place{卡拉穆斯}的地方，因一桩谋杀案，死亡的恐惧迫使他逃进\OriginalTerm{隐修院}{monastery}，但他的功德和圆满福乐一无所缺；因为他如此善用这强制的归化，以至以欣然的热忱将其转化为自愿的归化，并攀登了\OriginalTerm{成全}{perfection}的顶峰。另一方面，对很多人——我不应提他们的名字——他们怀着比此更好的开端事奉主，却毫无益处，因为后来懒惰和心硬悄然侵入，他们陷入危险的\OriginalTerm{麻木}{torpor}状态，以及死亡的无底深渊；我们在宗徒的召叫中清楚看到一个例子。\Person{犹达斯}出于自由意志，以与\Person{伯多禄}及其他宗徒同样的方式接受了\OriginalTerm{宗徒职}{apostolate}的最高品位，但这对他有何益处呢？他竟让他召叫的光辉开端终结于贪得无厌的毁灭结局，甚至像一个残忍的谋杀者，冲向了出卖主的行径？或者，\Person{保禄}突然失明，似乎被违背己意地拖入救恩之路，但这对他有何阻碍呢？他后来以全副心灵的\OriginalTerm{热忱}{fervour}跟随主，开始于强制，却以自由自愿的奉献完成，并以光荣的结局了结了一生，这生命因如此伟大的事迹而显得荣耀。因此一切都取决于结局；一个在起初就被崇高悔改所圣化的人，可能因疏忽而成为失败者；而一个因需要被迫选择隐修生活的人，可能因敬畏天主和诚心而得以成全。\par

% structure: p0019 source=h2 target=subsection reason=relative-heading-rank; base=h2

\subsection{第六章}

\Emph{论三种弃绝。}\par

我们现在必须谈论那些\OriginalTerm{弃绝}{renunciation}，传统和圣经的权威向我们显示了三种弃绝，我们每个人都应以最大的热忱将它们完全实现。第一种是就身体而言，我们轻视世上一切财富与货物；第二种是我们拒绝世俗、邪恶以及灵魂和肉体过去的私欲；第三种是我们使自己的灵魂脱离所有现世和可见的事物，只默思将来之事，并将心专注于不可见之事。我们读到，主命令\Person{亚巴郎}同时做到这三者，当时祂对他说：\BibleQuote{离开你的故乡、你的家族和父家} \BibleRef{创 12:1}。首先祂说\BibleQuote{离开你的故乡}，即离开世上的财富和地上的财宝；其次，\BibleQuote{离开你的家族}，即离开过去的生活、习惯与罪恶，它们从我们出生就依附我们，仿佛因亲和与血缘的纽带而连接；第三，\BibleQuote{离开你的父家}，即离开对世间的所有记忆，那是眼睛所能见到的。因为对于两个父亲，即一位要被舍弃，一位要被寻求，\Person{达味}如此以天主的身份说：\BibleQuote{女儿！请听，请看，请侧耳细听：忘却你的民族、你的父家！}因为说\BibleQuote{女儿！请听}的那位确实是圣父；然而祂作证说，祂敦促应被忘却其家和民族的那位，仍然是祂女儿的父亲。这事发生时，当我们与基督一同死于这个世界的世俗原理，我们不再如宗徒所说的，注目\BibleQuote{那看得见的，而是那看不见的，因为那看不见的是永远的} \BibleRef{格后 4:18}，并且心中离开这个暂时可见的住所，将我们的眼目和心灵转向我们将永远留驻的那一位。当我们行在肉身中，却不再按肉身与主争战，而是以言与行宣告蒙福之宗徒的那句话的真理，\BibleQuote{我们的家乡在天上} \BibleRef{斐 3:20}，那时我们将成功做到这一点。与这三种弃绝相应的，正是\Person{撒罗满}的三卷书。因为\WorkTitle{箴言}对应第一种弃绝，其中抑制了对肉体事物和地上罪恶的欲望；\WorkTitle{训道篇}对应第二种，其中宣称太阳下所行的一切皆是虚妄；\WorkTitle{雅歌}对应第三种，其中灵魂翱翔于所有可见之物之上，借着默观天上之事，实际地与天主的圣言结合。\par

% structure: p0022 source=h2 target=subsection reason=relative-heading-rank; base=h2

\subsection{第七章}

\Emph{我们如何在每一种弃绝中达至成全。}\par

因此，我们虽以极大的虔诚和信德作出第一次舍弃，但若不以此同等的热忱和激情完成第二次舍弃，这对我们将不会有太大益处。因此，当我们在这件事上成功时，我们也将能达到第三次舍弃，即离开我们先前父母的房屋（我们深知，按旧人，从我们出生起，他就是我们的父亲，因为我们\BibleQuote{按本性为义怒之子，和其余的人一样}\BibleRef{弗 2:3}），并将我们全部的心灵目光固定于天上的事物。论及这位父亲，圣经对那些藐视真天主圣父的耶路撒冷说：\BibleQuote{你的父亲是阿摩黎人，你的母亲是赫特人。}\BibleRef{则 16:3}在福音中我们读到：\BibleQuote{你们是出于你们的父亲魔鬼，你们喜欢行你们父亲的欲望。}\BibleRef{若 8:44}当我们离开他，从可见之物转向不可见之物时，我们将能同宗徒一起说：\BibleQuote{因为我们知道：如果我们地上帐幕式的房屋被拆毁，我们必由天主获得一所房舍，一所非人手所造，永远在天上的房舍。}\BibleRef{格后 5:1}以及我们不久前引用的：\BibleQuote{但我们的家乡是在天上，我们也等待主耶稣基督救主从那里降来，他将按他大能的工作，改变我们卑微的身体，使它相似他光荣的身体。}\BibleRef{斐 3:20}以及受祝福的达味的话：\Quote{因为我是世上的旅客}和\Quote{像我所有的祖先一样是外方人。}这样，我们就能按主的话，成为像主在福音中向圣父所提及的那些人一样：\BibleQuote{他们不属于世界，如同我不属于世界一样。}\BibleRef{若 17:16}主又对宗徒们本人说：\BibleQuote{如果你们属于世界，世界必爱属于自己的人；但因为你们不属于世界，故此世界恨你们。}\BibleRef{若 15:19}那么，关于这第三次舍弃，每当我们的灵魂没有被任何肉欲粗鄙的污点所玷污，反而在仔细清除这一切之后，摆脱了一切尘世的性质和欲望，并藉着对神圣事物的不断默想和属神的默观，已如此转向不可见的事物，以致在热切追求天上和属神之事时，不再感到自己被囚禁于这脆弱的肉身和形体之中，反而被提升到这样一种\OriginalTerm{神魂超拔}{ecstasy}的境界：不但耳朵听不到言语，也不忙着注视当前事物的形状，甚至连近在眼前或直对着眼睛的大物体也看不见，那时我们便能在这一舍弃上达到圆满。除了一个人，谁也不能理解这事的真理和力量——这个人曾在经验的教导下经历了上述所说；也就是说，主已将他心灵的眼睛从一切眼前事物上移开，使他不再视它们为即将逝去之物，而视为已经结束之事，并看它们像烟雾一样化为乌有；他像\Person{哈诺客}一样，\Quote{与天主同行}，从人类的生活和方式中被\Quote{提升}，没有\Quote{被寻见}于这人生的虚浮之中。这事在\Person{哈诺客}身上确实以肉身实现了，\BibleBook{创世}{纪}这样告诉我们：\Quote{哈诺客与天主同行，就不见了，因为天主将他提去。}宗徒也说：\Quote{因着信德，哈诺客被提升，使他不见死亡}，即主在福音中所说的死亡：\BibleQuote{活着信从我的人，永远不死。}因此，如果我们渴望获得真正的成全，就应当注意：如同我们外在已藉身体轻视了父母、家园、世上的财富和享乐，我们内心也当同样以心灵弃绝这一切，永不因任何欲望而回到我们所弃绝的事物，如同那些被\Person{梅瑟}带领出来的人，虽然并未实际回去，却被称为在心内回到了埃及；即他们离弃了曾以如此大能奇事带领他们出来的天主，并崇拜他们曾轻视的埃及偶像，正如经上所说：\BibleQuote{他们在心里回到埃及，对亚郎说：请给我们造神，在我们前面引路。}\BibleRef{宗 7:39}因为我们会与那些住在旷野中、在尝过天上玛纳之后仍贪恋罪恶污秽和卑下之食的人一同受到同样的谴责，并且会同他们一起发出同样的怨言：\Quote{我们在埃及过得很好，当时我们坐在肉锅旁，吃葱、蒜、黄瓜和西瓜。}这种说法，虽然原是指那个民族而言，但我们今天看到在我们自己的情况及生活方式中同样应验了：因为每一个舍弃世界后又回到旧欲望、回复先前喜好的人，在内心和行为上都断言了与古人相同的事，说：\Quote{我在埃及过得很好。}我担心这样的人数会与\Person{梅瑟}时代许多倒退者一样多，因为虽然算有六十万三千武装人员从埃及出来，但其中进入应许之地的不过两人。因此，我们应当谨慎，从那些稀少而间隔的人之中选取善行的榜样，因为按照我们在福音中提到的那句话：\BibleQuote{被召的多，但被选的却是少数。}\BibleRef{玛 22:14}因此，如果我们在内心舍弃（这远为高超和宝贵）上未能成功，那么单是身体上的舍弃和仅仅离开埃及改变地点对我们毫无益处。因为我们所谈的单纯身体舍弃，宗徒这样宣告说：\BibleQuote{我若将我所有的财物施舍给人，并舍身被火焚烧，但我若没有爱德，于我毫无益处。}\BibleRef{格前 13:3}受祝福的宗徒绝不会说这话，除非他藉圣神预见到，有些曾施舍全部财物给穷人的人，未能达到福音的成全和\OriginalTerm{爱德}{charity}的高度，因为当骄傲或不耐烦支配他们的心时，他们未能小心地净化自己脱离先前的罪过和无度的习惯，因此永远达不到那永不失败的天主之爱；这些人，既然在第二次舍弃阶段上有亏缺，就更不能达到肯定更高超的第三次舍弃。但你们也当十分谨慎地思考：他并没有简单地说\Quote{我若施舍我的财物}。因为有人可能以为他指的是未曾遵行福音命令、而为自己留下了一些东西的人，如同一些半心半意的人所做的。但他却说\Quote{我若将我所有的财物施舍给穷人}，即，即使我对这些世俗财富的舍弃是完美的。再者，在这舍弃上他还加上了一句更大的：\Quote{我若舍身被火焚烧，但没有爱德，我算不了什么。}好像他用另外的话说：虽然我将我所有的财物施舍给穷人，按照福音的命令，那里告诉我们：\BibleQuote{你若愿意是成全的，去，变卖你所有的，施舍给穷人，你必有宝藏在天上。}\BibleRef{玛 19:21}舍弃它们，使我毫无保留；并且虽然在这种（财物的）分配上，我藉着焚烧肉体加上了殉道，为基督舍身，然而若我不耐烦、易怒、嫉妒、骄傲，或因他人的屈辱而激动，或寻求自己的利益，或沉溺于邪念，或不愿准备和忍耐承受一切加于我身的事，那么在外在的人的舍弃和焚烧上，于我毫无益处，因为内在的人仍被先前的罪过所缠绕；因为在我初皈所怀的热忱中，我轻视了仅仅世俗的财富（这些本身可说不好不坏，是中性之物），却没有同样小心地驱除一个坏心肠的有害力量，或达到主的爱，这爱是忍耐的，是\BibleQuote{仁慈的，不嫉妒，不自大，不轻易发怒，不作不义之事，不求己益，不怀恨}，并且\BibleQuote{凡事包容，凡事忍耐}\BibleRef{格前 13:4}，这爱最后从不允许追随它的人因罪过的欺骗而跌倒。\par

% structure: p0025 source=h2 target=subsection reason=relative-heading-rank; base=h2

\subsection{第八章}

\Emph{论我们真正的财产，灵魂的美或丑从中可见。}\par

我们当竭力使我们的内在人也同样脱去并抛弃它在过去生活中所获得的一切罪的财产；这些罪的财产不断依附于灵魂和肉身，乃是我们的真正财产，除非我们在肉身内就拒绝并砍断它们，它们死后仍将伴随我们。因为就如善德，或产生善德的爱德本身，可在此世获得，并在今生结束后使爱它的人变得美丽而荣耀，同样，我们的过犯将把一颗被污秽颜色染黑弄暗的灵魂传给永恒的纪念。因为灵魂的美丑是其德行或恶习的产物，灵魂从德行或恶习染上的颜色，要么使它如此荣耀，以至它能听到先知说：\BibleQuote{君王要恋慕你的美艳}，要么使它如此黑暗、污秽、丑恶，以至它必然承认自己羞耻的恶臭，说：\BibleQuote{因我的愚昧，我的伤口发臭流脓}，主自己也对她说：\BibleQuote{为何我百姓女儿的伤口不得痊愈？} \BibleRef{耶 8:22} 所以这些才是我们真正的财产，永远留于灵魂，任何君王或仇敌都不能赐予或夺走。这些是我们真正的财产，连死亡本身也不能使灵魂与之分离，但藉着舍弃它们我们能达于成全，藉着紧抓它们我们将承受永死的惩罚。\par

% structure: p0028 source=h2 target=subsection reason=relative-heading-rank; base=h2

\subsection{第九章}

\Emph{论三种财产。}\par

圣经中以三种不同的方式论及财富和财产，即：善的、恶的、中性的。恶的如经上所说：\BibleQuote{富有的，竟忍受饥饿}，以及\BibleQuote{你们富有的人，你们已获得了你们的安慰} \BibleRef{路 6:24}；抛弃这些财富是成全的顶峰；属于那些贫穷者的荣耀，他们在福音中被主的话称赞：\BibleQuote{神贫的人是有福的，因为天国是他们的} \BibleRef{玛 5:3}；在圣咏中：\BibleQuote{这贫穷的人一呼号，上主即俯听}，又说：\BibleQuote{贫乏穷困的人必歌颂你的名}。那些财富是善的，获得它们需要大德和功德，义人拥有它们而受达味称赞说：\BibleQuote{义人的后代必要兴盛；他的家中必有荣耀和财富；他的正义永远存留}，又说：\BibleQuote{人的财富是他性命的赎价} \BibleRef{箴 13:8}。论及这些财富，在\BibleBook{}{默示录}中，对那些没有这些财富而可耻地贫穷赤裸的人说：\BibleQuote{我必要将你从我口中吐出来。因为你说：我是富有的，我已发了财，什么也不缺少；却不知道你是悲惨、可怜、贫穷、瞎眼、赤身的。我劝你向我买经火炼过的金子，好叫你富足；并买白衣穿上，以免你赤身的羞耻显露出来} \BibleRef{默 3:16}。有些财富是中性的，即可善可恶，因为按照使用者的意志和品性，它们可成为善或恶。对此，有福的宗徒说：\BibleQuote{劝告今世的富人，不要心高气傲，也不要寄望于无常的财富，但要寄望于天主（他厚赐百物供我们享受），行善，乐意分施，与人通功，为自己积蓄稳固的基础，以便把握真正的生命} \BibleRef{弟前 6:17}。这些正是福音中那富翁所保留而从未分施给穷人的财富——而乞丐拉匝禄曾躺在他的门口，渴望从他的碎屑中得饱饫；因此他被判受地狱之火无法忍受的火焰和永恒的灼热。\par

% structure: p0031 source=h2 target=subsection reason=relative-heading-rank; base=h2

\subsection{第十章}

\Emph{只靠第一种舍弃程度无人能成为成全的。}\par

既然如此，我们离开现世这些可见的财物，所舍弃的并非自己的财富，而是不属于我们的东西，尽管我们夸口说那是靠自己努力得来的，或是从祖先继承的。因为正如我说过的，除了我们心中所拥有的、紧贴我们灵魂的那份以外，没有什么是真正属于我们的，因此谁也不能从我们身上夺走它。但基督对那些紧抓可见的财富仿佛它们属于自己、不肯与贫乏之人分享的人，用责备的口吻说：\Quote{你们在不义的钱财上不忠信，谁还把属于你们的交给你们呢？}\BibleRef{路 16:12} 显然，不仅日常经验教导我们这些财富不是我们的，而且主的这句话也通过它赋予的称号表明这点。然而，关于可见而无价值的财富，伯多禄对主说：\Quote{看，我们舍弃了一切，跟随了你；那么，我们将得到什么呢？}\BibleRef{玛 19:27} 显然他们舍弃的不过是可怜破旧的渔网。除非这句\Quote{一切}被理解为指那真正伟大而重要的弃绝罪恶，否则我们不会发现宗徒们舍弃了什么有价值的东西，也不会发现主有任何理由赐予他们如此伟大光荣的福分，使他们得以听祂说：\Quote{在复兴的时候，人子坐在祂光荣的宝座上，你们也要坐在十二个宝座上，审判以色列十二支派。}那么，如果那些完全弃绝了地上可见财物的人，尚且不能充分理由达到宗徒的爱德，也不能敏捷有力地攀登那更高、属于少数人的第三级弃绝，那么那些连第一步（其实很容易）也不彻底做到，反而保留着旧的不信、从前肮脏的财富，却幻想可以炫耀隐修士虚名的人，又该怎样看待自己呢？我们所说第一级弃绝舍弃的是不属于自己的东西，因此它本身不足以使弃绝者达于成全，除非他进到第二级，那才是真正属于自己的弃绝。当我们借着除去一切过失而确保这一点后，我们也会登至第三级弃绝的高度，借此不仅超越现世一切作为或特属于人的事物，而且超越我们周围被视为如此光荣的整个宇宙，以心灵之眼俯视它，视之为屈服于虚幻、注定要消逝的；正如宗徒所说：\Quote{我们不是注视那看得见的，而是注视那看不见的；因为那看得见的是暂时的，那看不见的是永远的。}\BibleRef{格后 4:18} 这样我们才堪当听到那对亚巴郎所说的至高话语：\Quote{到我要指示你的地方去。}\BibleRef{创 12:1} 这清楚表明，除非人以全副心神作好上述三种弃绝，否则不能达到这第四种弃绝，它是作为完美的弃绝者的赏报和特恩赐予的，使他堪当进入预许之地，那地不再为他长出罪恶的荆棘与蒺藜；在祛除一切情欲之后，即使在肉身中，也借着心灵的纯洁获得它；任何人的善工或努力都不能赢得它，而是主亲自应许要指示它，说：\Quote{到我要指示你的地方去。}这清楚证明我们得救的开始来自主的召唤，祂说：\Quote{离开你的故乡}；而成全与纯洁的完成同样是祂的恩赐，正如祂说：\Quote{到我要指示你的地方去。}也就是说，不是你自己靠努力能知道或发现的地方，而是我要指示给你的，不仅指示给那不知道的人，甚至指示给那没有寻找的人。由此我们清楚领会：正如我们因主的感动而被催迫走向救恩之路，同样也是在祂的引导和光照下，我们达到至高福乐的成全。\par

% structure: p0034 source=h2 target=subsection reason=relative-heading-rank; base=h2

\subsection{第十一章}

\Emph{关于人的自由意志与天主恩宠的问题。}\par

格玛诺说：那么，哪里还有自由意志的余地？如果天主在我们得救的事上开始并完成一切，我们又如何因自己的努力而值得赞美呢？\par

% structure: p0037 source=h2 target=subsection reason=relative-heading-rank; base=h2

\subsection{第十二章}

\Emph{关于在我们内仍保留自由意志的天主恩宠之经济的答复。}\par

帕夫努提乌斯：如果每项工作和实践之中，开始与结局就是一切，而中间不存在任何东西，那么这话确实会影响我们。然而，我们知道天主以多种方式创造救恩的机会，我们更有力量或更少认真地去利用上天赐予我们的机会。因为正如召叫他的天主的邀请\Quote{离开你的故乡}，同样亚巴郎的服从在于他出去了；正如\Quote{进入那地}的这句话被付诸行动，是服从者的事，同样\Quote{我将指示你的地方}这句话的附加，来自命令或应许它的天主的恩宠。但我们最好确信，尽管我们以不懈的努力实践每种德行，但凭我们所有的努力和热忱，我们永远无法达到成全，单靠人的勤奋和辛劳本身也不足以配得获得福乐的崇高赏报，除非我们藉着主的合作，以及祂将我们的心导向正义而获得它。因此，我们应当每时每刻祈祷，并同达味一起说：\Quote{求你引导我的脚步在你的道路上，使我的脚步不致滑倒。}以及\Quote{祂把我的脚放在磐石上，稳定了我的脚步。}好使那看不见的人类心中主宰，乐意将我们那更倾向于恶——或因缺乏对善的认识，或因情欲的快乐——的意志转向对德行的渴望。我们在先知唱得很清楚的一节经文中读到：\Quote{我被推倒，几乎跌倒，}那里显示了自由意志的软弱。而\Quote{上主扶持了我，}这再次表明主的帮助常与它相连，藉此，当我们跌倒时，祂看见我们，便扶持和支持我们，如同伸手扶持。又说：\Quote{我说：我的脚滑了一下，}即出于意志的滑溜，\Quote{上主，你的仁慈帮助了我。}他再次将天主的帮助与自己的软弱相连，承认不是由于自己的努力，而是依靠天主的仁慈，他信德之脚才没有移动。又说：\Quote{我心中思虑繁多，}这绝对出于我的自由意志，\Quote{你的安慰欢愉了我的灵魂，}即藉着你的灵感进入我心中，并揭示你为那些因你名劳苦的人所预备的未来福乐景象，它们不仅从我心中除去了所有焦虑，而且实际上赐给了它最大的喜悦。又说：\Quote{若不是上主帮助了我，我的灵魂几乎已住在地狱。}他确实表明，由于这自由意志的败坏，他本会住在地狱，若不是藉着主的帮助和保护得救。因为\Quote{人的脚步是由上主引导，}而非自由意志，\Quote{义人纵使跌倒，}至少出于自由意志，\Quote{也不会被抛弃。}原因何在？因为\Quote{上主用手扶持他。}这说得极其清楚：没有义人靠自己足以获得义德，除非每当他们因自由意志的软弱而绊跌时，天主仁慈用手扶持他们，使他们不致完全崩溃而灭亡。\par

% structure: p0040 source=h2 target=subsection reason=relative-heading-rank; base=h2

\subsection{第十三章}

\Emph{我们道路的导向来自天主。}\par

事实上，圣人们从未说过，他们凭自己的努力获得了在走向德行进步和成全的道路上所行之路的导向，而是向主祈求说：\Quote{求你引导我走你的真理之路，}以及\Quote{求你在我面前引导我的道路。}但另一个人不仅藉着信德，也藉着经验，并可以说从事物本性本身，发现了这个事实：\BibleQuote{我知道，上主，人的道路不由自己决定；行路的人也不能定自己的脚步。}\BibleRef{耶 10:23}主自己向以色列说：\BibleQuote{我要引导他像一棵青翠的松树；你的果实从我这里找到。}\BibleRef{欧 14:9}\par

% structure: p0043 source=h2 target=subsection reason=relative-heading-rank; base=h2

\subsection{第十四章}

\Emph{法律的知识是由上主的引导和光照所赐。}\par

法律本身的知识，他们每天努力获得的，不是靠阅读的勤奋，而是藉着天主的引导和光照，正如他们对祂说：\Quote{上主，求你把你的道路指示给我，教给我你的途径；}以及\Quote{求你开启我的眼睛，使我看见你法律中的奇妙；}以及\Quote{求你教我承行你的旨意，因为你是我的天主；}又说：\Quote{是祂教导人知识。}\par

% structure: p0046 source=h2 target=subsection reason=relative-heading-rank; base=h2

\subsection{第十五章}

\Emph{我们可以认识天主诫命的悟性，以及善意的实践，都是主的恩赐。}\par

有福的达味更进一步向主祈求获得那种理解力，藉以认识天主的诫命——他深知这些诫命写在法律书上——他说：\Quote{我是祢的仆人，求祢赐我理解力，使我学习祢的诫命。}他固然拥有本性赋予的理解力，也熟知记载于法律的诫命，但仍向主祈求更透彻地学习，因为他知道本性的所是不能使他满足，除非他的理解力每天都蒙主光照，以属神的方式理解法律，更清楚地认识祂的诫命，正如\Quote{蒙选之器}也清楚表明我们所坚持的这一点。\Quote{因为是天主在你们内工作，使你们愿意，并依照祂的善意行事。}\BibleRef{斐 2:13}还有什么比这更清楚呢？即我们的善意和工作的完成都是主在我们内完全成就的。又说：\Quote{因为为了基督的缘故，赐给你们恩典，不仅为相信祂，而且也为祂受苦。}\BibleRef{斐 1:29}同样，他宣告我们皈依和信德的开端以及忍受苦难，都是主赐予我们的恩赐。达味也知道这事，同样祈求主以慈悲赐予他同一恩赐。他说：\Quote{天主，求祢坚固祢在我们内所成就的工程。}这表明救恩的开始由天主的恩赐和恩宠赐予还不够，除非藉着祂同样的怜悯和不断的帮助持续并完成。因为不是自由意志，而是主\Quote{释放了那些被捆绑的}；不是我们的力量，而是主\Quote{扶起那跌倒的}；不是阅读的勤勉，而是\Quote{主开启瞎子的眼睛}——希腊文作 \OriginalTerm{主使瞎子聪明}{\GreekText{κύριος} \GreekText{σοφοῖ} \GreekText{τυφλούς}}，即\Quote{主使瞎子聪明}；不是我们的照顾，而是\Quote{主保护旅客}；不是我们的勇气，而是\Quote{主扶助（或扶持）所有跌倒的人}。但我们这样说，并非轻视我们的热心、努力和勤勉，好像它们是无用和愚昧的施予，而是要我们知道：没有天主的帮助我们不能奋斗，若未蒙主的助佑和怜悯赐予我们，我们的努力在获得纯洁的大赏报上毫无用处。因为\Quote{马是为作战之日预备的，但援助出自于主}\BibleRef{箴 21:31}，\Quote{因为人不能因力量得胜}\BibleRef{撒上 2:9}。因此，我们应当常与有福的达味同唱：\Quote{我的力量与我的歌颂，不是我的自由意志，而是主，并且祂成了我的救恩。}外邦人的导师也深知此事，他宣布自己能成为新约的仆役，不是凭自己的功德或努力，而是因天主的慈悲。他说：\Quote{并不是我们凭自己能够承担什么事，好像出于自己；而是我们的能力出于天主}——这句可以用不太好的拉丁文但更强烈地表达为\Quote{我们的能力出于天主}——随后接着说：\Quote{并且祂使我们能够做新约的仆役。}\BibleRef{格后 3:5}\par

% structure: p0049 source=h2 target=subsection reason=relative-heading-rank; base=h2

\subsection{第十六章}

\Emph{信德本身也必须由主赐予我们。}\par

宗徒们深知一切关乎救恩的事都由主赐予他们，以致他们甚至请求信德本身由主赐予，说：\Quote{请增加我们的信德}\BibleRef{路 17:5}，因为他们不认为信德能凭自由意志获得，而相信它是天主的白白恩赐。最后，人类救恩的导师教导我们，信德若不藉主的帮助加强，是多么软弱无力、不足，当祂对伯多禄说：\Quote{西满，西满，看，撒殚想要得着你们，好筛你们像筛麦子一样。但我已为祢祈求，使祢的信德不至丧失。}\BibleRef{路 22:31}另有一位发现这事临到自己身上，看见自己的信德被不信的波浪冲向可怕的礁石，就向同一主祈求帮助他的信德，说：\Quote{主，请帮助我的不信。}\BibleRef{谷 9:23}福音中那些宗徒和人们如此清楚地意识到一切美善都是由主的助佑达致成全，而不认为凭自己的力量或自由意志能保守信德无瑕，以致祈求主帮助或赐予信德。若伯多禄需要主的帮助使信德不致丧失，谁敢如此狂妄和盲目，认为自己无需每日从主获得帮助以保守信德呢？尤其主自己在福音中已表明：\Quote{枝条若不留在葡萄树上，凭自己不能结果实；同样，你们若不住在我内，也不能。}\BibleRef{若 15:4}又说：\Quote{因为离了我，你们什么也不能做。}因此，将任何善行归于自己的勤勉而不归于天主的恩宠和助佑，是多么愚昧和邪恶，这由主的话清楚显明，祂规定人若没有祂的感动和合作，不能结出圣神的果实。因为\Quote{一切美好的赠与，一切完善的恩赐，都是从上头来的，从光明之父降下来的。}\BibleRef{雅 1:17}匝加利亚也说：\Quote{因为凡是善的，都是属于祂的，凡是卓越的，都来自祂。}因此，有福的宗徒一贯地说：\Quote{你有什么不是领受的呢？既然领受了，为什么还夸耀，好像没有领受呢？}\BibleRef{格前 4:7}\par

% structure: p0052 source=h2 target=subsection reason=relative-heading-rank; base=h2

\subsection{第十七章}

\Emph{节制和忍受诱惑也必须由主赐予我们。}\par

而蒙福的宗徒如此宣告说，我们所能承受加于我们之试探的一切忍耐，并不那么倚靠我们自己的力量，而是倚靠天主的仁慈和引导：\BibleQuote{你们所受的试探，无非是普通人所能受的。天主是忠信的，他决不许你们受那超过你们能力的试探；而且加给试探，也必开一条出路，叫你们能够承担。}\BibleRef{格前 10:13}同一宗徒又教训说，天主装备并坚强我们的灵魂，为一切善工，并在我们身上运作一切中悦祂的事：\BibleQuote{愿赐平安的天主，就是那凭着永盟之血，领那牧养群羊的大司牧——我们的主耶稣基督从死者中上来的，愿祂在一切善事上成全你们，为承行祂的旨意，在你们身上运行祂所喜悦的事。}\BibleRef{希 13:20}他为得撒洛尼人祈祷同一件事，说：\BibleQuote{愿我们的主耶稣基督自己和那爱我们、并赐给我们永远的安慰和美好的希望、藉着恩宠、天主我们的父，鼓励你们的心，并在一切善言善事上坚固你们。}\BibleRef{得后 2:15}\par

% structure: p0055 source=h2 target=subsection reason=relative-heading-rank; base=h2

\subsection{第十八章}

\Emph{对天主的持续敬畏必须由主赐予我们。}\par

最后，先知耶肋米亚以天主的身份说话，清楚地证明，连我们能持守祂的敬畏，也是由主倾注在我们身上的：他说：\BibleQuote{我要赐给他们一颗心，一条道路，好使他们日日敬畏我，使他们和他们的后代得享福乐。我要与他们订立永久的盟约，不再离开他们，为他们造福；并将我的敬畏放在他们心里，使他们不再离开我。}\BibleRef{耶 32:39}厄则克耳也说：\BibleQuote{我要赐给他们一颗心，在你们五内注入新精神；从他们肉身中除去石心，赐给他们肉心，使他们遵行我的法律，恪守我的诫命，并一一实行，他们将作我的百姓，我将作他们的天主。}\BibleRef{则 11:19}\par

% structure: p0058 source=h2 target=subsection reason=relative-heading-rank; base=h2

\subsection{第十九章}

\Emph{我们善愿的开端和完成都来自天主。}\par

这清楚教导我们，我们善愿的开端是由主自己或借着某人的劝诫或强制，将我们引向救恩之路，藉着祂的感动赐予我们；而我们善功的完成也同样由祂赐予。但以多少热忱来追随天主的鼓励和协助，则在于我们自己的能力；因此，我们或得赏或受罚，都是正义的，因为我们或是粗心，或是谨慎，来配合祂以如此慈爱为我们所定的计划和眷顾安排。这在\BibleBook{}{申命纪}中清楚明白地描述了。他说：\BibleQuote{当上主你的天主领你进入你要去占领的土地，从你面前清除许多民族，即赫特人、基勒加人、阿摩黎人、客纳罕人、培黎齐人、希威人和耶步斯人，七个比你们更大更强的民族，上主你的天主将他们交给你时，你应完全消灭他们，不可与他们立约，也不可与他们联姻。}\BibleRef{申 7:1}可见圣经宣告，是他们被领进了福地，许多民族在他们面前被消灭，比以色列人民更多更强的民族被交在他们手中，这都是天主的白赐。但以色列是否完全消灭他们，或者留他们存活宽恕他们，是否与他们立约，是否与他们联姻，圣经说明这在于他们自己的权力。藉此证言，我们清楚看出，哪些该归因于自由意志，哪些该归因于主的计划和日常协助；即给予我们得救的机会、顺利的开端和胜利，属于天主的恩宠；而以热忱或冷淡来追随天主赐予的福佑，则是我们的事。我们看见同样的事实，在治愈瞎子的例子中也明白教训了。因为耶稣从他们旁边经过，是天主眷顾和屈尊的白赐；但他们呼喊说：\BibleQuote{主，达味之子，可怜我们吧！}\BibleRef{玛 20:31}则是他们自己信德的行为；他们重见光明，是神圣怜悯的恩赐。但在领受任何恩惠之后，天主的恩宠和自由意志的运用仍然存在，十个癞病人的事例向我们显示了这一点。他们全都同样被治愈，但其中一人出于善意的感恩回来感谢时，主寻找那九个，并称赞那一个，表明祂甚至常愿帮助那些忘记祂恩惠的人。因为连感恩者被祂接纳和称赞，忘恩者被祂寻找和责斥，也是祂眷顾的恩赐。\par

% structure: p0061 source=h2 target=subsection reason=relative-heading-rank; base=h2

\subsection{第二十章}

\Emph{没有天主，世上什么也不能成就。}\par

但我们应该坚定不移地相信，世上任何事没有天主便不能成就。因为我们必须承认，一切事或出于祂的旨意，或出于祂的许可；即，凡是善的，都由天主的旨意和祂的帮助而实现；凡是相反的，则出于祂的许可，当我们因自己的罪和内心的顽固，神圣的保护离开我们时，祂容许魔鬼和肉身可耻的情欲来主宰我们。宗徒的话确实教导我们这一点，他说：\BibleQuote{为此，天主任凭他们陷于可耻的情欲中。}又说：\BibleQuote{他们既然不肯承认天主，天主就任凭他们陷于邪恶的心思，去行不正当的事。}主自己藉先知说：\BibleQuote{我的百姓没有听从我的声音，以色列不肯服从我；因此，我任凭他们随从心里的顽梗，让他们照自己的计谋而行。}\par

% structure: p0064 source=h2 target=subsection reason=relative-heading-rank; base=h2

\subsection{第二十一章}

\Emph{一个关于自由意志能力的反驳。}\par

\Person{革尔玛努斯}：这段经文非常清楚地显示自由意志，因为经上记载\Quote{如果我的百姓肯听从我}，别处又称\Quote{但我的百姓不肯听我的声音}。当他说\Quote{如果他们肯听}时，他表明顺从或不顺从的决定掌握在他们自己的手中。那么，如果天主亲自赐给我们力量去听从或不听从，又怎能说我们的救恩不取决于我们自己呢？\par

% structure: p0067 source=h2 target=subsection reason=relative-heading-rank; base=h2

\subsection{第二十二章}

\Emph{回答，即我们的自由意志总需要主的帮助。}\par

\Person{帕夫努提乌斯}：你颇为精明地注意到经上所说的\Quote{如果他们肯听从我}，却没有充分考虑说话的是谁，即对听从或不听从之人说话的那一位；也没有考虑下文：\Quote{我本会迅速制伏他们的敌人，将我的手加在虐待他们的人身上。}因此，任何人都不应试图通过错误的解释来曲解我们所提出的、为证明没有主什么都不能做的经文，也不应利用它来支持自由意志，以致借这段经文妄图夺去天主的恩宠和对人每日的眷顾。这段经文是：\Quote{但我的百姓不肯听我的声音}，以及\Quote{如果我的百姓肯听从我，如果以色列肯行走在我的道路等等}；不过，他应当考虑，正如自由意志的能力因百姓的不服从而显明，同样，那位宣告和劝戒人的天主每日的眷顾也得以显明。因为他说\Quote{如果我的百姓肯听从我}时，清楚地暗示他先前曾向他们说过话。主惯常这样做，不仅通过成文的律法，也通过每日的劝勉，正如依撒意亚所说的：\BibleQuote{我整天向悖逆和反对的百姓伸出我的手。}\BibleRef{依 65:2}因此，这两点都能从这段经文得到支持，因为它说：\Quote{如果我的百姓肯听从，如果以色列行走在我的道路，我本会迅速制伏他们的敌人，将我的手加在虐待他们的人身上。}因为正如自由意志因百姓的不服从而显明，同样，天主的治理和助佑也由这节经文的开端和结尾表明，其中暗示他先前曾向他们说话，而如果他们肯听从，他后来会制伏他们的敌人。因为我们绝无意通过我们所说的来废除人的自由意志，而只是要确立一个事实，即天主的助佑和恩宠对自由意志每日每时都是必要的。当\Person{帕夫努提乌斯}院长用这篇讲论教导我们之后，在午夜前打发我们离开他的斗室，我们心中充满痛悔而非轻松；他坚持认为讲论中的主要教训是：当我们以为藉着完美实行第一次舍弃（我们正竭尽全力去做），就能攀登完美的高峰时，我们发现我们甚至还没有开始梦想隐修士所能达到的高度，因为在我们学习了关于第二次舍弃的一些事情之后，才发现在此之前我们连第三阶段的只言片语都未听过，而这第三阶段包含一切完美，并在许多方面远超过这些较低的阶段。\par

\SourceRecord{Translated by C.S. Gibson.；Nicene and Post-Nicene Fathers, Second Series；Vol. 11.；Edited by Philip Schaff and Henry Wace.；Buffalo, NY: Christian Literature Publishing Co.,；1894.；<http://www.newadvent.org/fathers/350803.htm>.}

% structure: p0001 source=h1 target=section reason=page-title

\section{第四篇会谈}

\Signature{圣若望·卡西安}\par

\Person{院长达尼尔}会谈：论肉情与精神的贪欲。\par

% structure: p0004 source=h2 target=subsection reason=relative-heading-rank; base=h2

\subsection{第一章}

\Emph{\Person{院长达尼尔}的生活。}\par

在基督哲学的其他英雄中，我们也认识\Person{院长达尼尔}，他不仅在斯塞特旷野居住者的各种德行上与众人相等，更以谦逊的恩宠特别突出。此人因他的纯洁和温和，尽管在年龄上是多数人的后辈，却被斯塞特旷野的同一位司铎\Person{真福帕弗努提乌斯}擢升为执事职务：因为\Person{真福帕弗努提乌斯}对他的卓越品质如此欣喜，以至于因知道他在德行和圣洁生活上与自己相等，便急于使他在司铎品级上也与自己平等。既然他不忍看他继续留在较低的职位，又急于在生前为自己预备一个称职的接班人，便提升他到了司铎的尊位。然而，他丝毫没有放弃从前习惯的谦逊，当另一位在场时，他从不因晋升到较高品级而自恃什么，反而在\Person{院长帕弗努提乌斯}奉献属灵祭献时，继续以他以前职务的执事身分行事。尽管如此，\Person{真福帕弗努提乌斯}尽管是一位如此伟大的圣人，在许多事上拥有预知恩宠，但在这次关于继承的希望和自己所作的选择上却失望了，因为他在自己准备作为接班人的那位之后不久，就安息主怀了。\par

% structure: p0007 source=h2 target=subsection reason=relative-heading-rank; base=h2

\subsection{第二章}

\Emph{探究从无法形容的喜乐到极度沮丧的情绪突变之根源。}\par

于是我们请问这位有福的\Person{达尼尔}，为什么当我们坐在小房间里时，有时心中充满极大的喜乐，连同难以言喻的喜悦和最丰盈的至圣情感，以至于我姑且不说\Emph{言语}，就连\Emph{感受}也无法跟上；纯洁的祈祷轻易涌出，心灵充满神果，甚至在睡梦中向\Person{天主}祈祷时，也能感到祈祷轻松有力地达于\Person{天主}；而另一方面，为什么毫无理由地，我们突然充满极大的悲伤，并被无理的沮丧所压垮，以至于我们不仅感到自己也被这样的情感所胜，就连我们的房间也显得可怕，阅读变得乏味，是的，连我们的祈祷也被摇摇晃晃、模糊不清地献上，几乎像沉醉了一样：以致当我们叹息着竭力恢复先前的状态时，我们的心灵却无法做到这一点；它越热切地寻求重新注视\Person{天主}，就越被飘忽不定的思绪剧烈地带走，如此完全丧失了一切神果，以至于甚至不能因对天国的渴望或呈给它的地狱恐惧而从这致命的昏睡中醒来。他对此回答说。\par

% structure: p0010 source=h2 target=subsection reason=relative-heading-rank; base=h2

\subsection{第三章}

\Emph{他对所提问题的回答。}\par

关于你们所提的这种心灵枯干，长老们给出了三重解释。因为它或来自我们方面的疏忽，或来自魔鬼的袭击，或来自\Person{主}的许可和允许。来自我们方面的疏忽，是指当我们因自己的过错而变得冷淡，行事粗心而匆忙，因懒散怠惰而喂养了恶念，从而使我们心田长出荆棘和蒺藜；它们在其中生长，结果使我们变得不结果，并在一切神果和默想上无能为力。来自魔鬼的袭击，是指有时，当我们实际上正专注于善念时，我们的仇敌以诡诈的机巧潜入我们心中，不知不觉、违背我们意愿地，我们从最好的意向中被拖走。\par

% structure: p0013 source=h2 target=subsection reason=relative-heading-rank; base=h2

\subsection{第四章}

\Emph{\Person{天主}的许可和允许如何有两层原因。}\par

但\Person{天主}的许可和允许有两层原因。第一，使我们暂时被\Person{主}离弃，并以全然的谦逊观察我们自己内心的软弱，就不至于因先前因祂眷顾而赐予我们的心纯而自高自大；并且，通过证明当我们被祂离弃时，我们无法靠自己的任何叹息和努力恢复到先前的纯洁状态和喜悦，我们也学到我们先前的内心喜乐不是出于自己的热忱，而是出于祂的恩赐，而且为了当前时刻，必须再次从祂的恩宠和启迪中寻求。而这允许的第二个原因，是为了考验我们的坚忍、心志的坚定和真正的渴望，并在我们身上显示，当圣神离开我们时，我们怀着怎样的心志或祈祷的热诚寻求祂的回归；并且，为使我们一旦发现必须付出怎样的努力寻求那失去的神性喜乐和纯洁的喜乐，我们便学会在获得后更加小心地保存它，并以更紧的把握抓住它。因为人们通常对认为容易重新获得的东西更加疏于保存。\par

% structure: p0016 source=h2 target=subsection reason=relative-heading-rank; base=h2

\subsection{第五章}

\Emph{没有\Person{天主}的帮助，我们的努力和辛劳如何毫无用处。}\par

由此清楚显示，天主的恩宠和仁慈常在我们内成就善事；当它离弃我们时，工人的努力便徒劳无益；无论人多么热切奋力，没有祂的帮助便无法恢复原来的状态；并且这句话常在我们身上应验：\BibleQuote{不在乎那愿意的，也不在乎那奔跑的，只在乎那发慈悲的天主。}\newline{}\BibleRef{罗 9:16} 另一方面，这恩宠有时也不拒绝以你所说的那种神圣灵感，以及丰沛的精神思想，眷顾那些粗心大意和漠不关心的人；却激发不配的人，唤醒沉睡的人，光照被无知蒙蔽的人，并仁慈地责备和惩戒我们，倾注在我们心中，好让我们因祂激起的痛悔而感动，并受推动从懈怠的睡眠中醒来。最后，我们常因祂的突然临在而充满甜美的香气，非人力所能调配——以致灵魂被这些喜乐所沉醉，仿佛被提到神魂超拔的境界，忘却自己仍处于肉身中。\par

% structure: p0019 source=h2 target=subsection reason=relative-heading-rank; base=h2

\subsection{第六章}

\Emph{有时被天主离弃如何对我们有利。}\par

但蒙福的达味认识到，有时我们所谈论的这种离开，以及（仿佛）被天主遗弃，在某种程度上可能对我们有利，以致他不愿祈求不被天主完全遗弃（因他明白这对于他自己以及人类本性趋向完美会是有害的），而是恳求这种遗弃要有分寸和程度，说：\Quote{求祢不要完全抛弃我}，仿佛在说：我知道祢为了圣人们的益处而遗弃他们，为要考验他们，因为若非祢暂时离弃他们，他们绝无可能被魔鬼试探。因此，我并非求祢永不遗弃我，因为不感到自己的软弱并说\Quote{祢使我谦卑，这为我有益}，以及没有争战的机会，对我并不好。如果天主的保护不断持续地遮蔽我，我肯定不会有这种机会。魔鬼在你保护之下不敢攻击我，如同他常对你所爱的人提出这诽谤性的指责和抱怨一样：\BibleQuote{约伯敬畏天主，岂是无缘无故的？祢不是四面围护了他、他的家和一切所有吗？}\newline{}\BibleRef{约 1:9} 但我宁愿恳求祢不要完全抛弃我——希腊文称为\OriginalTerm{直到极度}{\GreekText{ἕως} \GreekText{σφόδρα}}，即太过分了。因为，首先，祢暂时离弃我，以考验我爱情的稳固，这对我有益；同样，如果祢按我的过失和应得而让我过分地被弃，那是危险的，因为人的力量若在试探中长时间缺乏祢的帮助，靠自己的坚定无法忍受，会立即屈服于敌人的势力，除非祢，如同祢深知人的力量，并调节他的争战，\BibleQuote{不让我们受试探超过我们所能承受的，却要在受试探的同时，开一条出路，使我们能担当。}\newline{}\BibleRef{格前 10:13} 我们读到在\WorkTitle{民长纪}中，关于歼灭敌对以色列的精神民族一事，也神秘地设计了类似情况：\BibleQuote{这些是上主留下的民族，为藉他们教导以色列人，使他们学习如何与敌人作战。}\newline{}后来又不久：\BibleQuote{上主留下他们，为要藉他们考验以色列人，看他们是否听从上主藉梅瑟所吩咐他们祖先的诫命。}\newline{}\BibleRef{民 3:1} 天主为以色列保留这场争战，并非嫉妒他们的平安，也不是有意伤害他们，而是因他知道这对他们有益，使他们因常受那些民族的攻击，而不断感到需要上主的助佑，并因此常默想祂并呼求祂的帮助，而不致因懒散安逸而疏忽，失去抵抗的习惯和德性的操练。因为多次，逆境不能征服的人，反因无忧无虑和顺遂而跌倒。\par

% structure: p0022 source=h2 target=subsection reason=relative-heading-rank; base=h2

\subsection{第七章}

\Emph{论宗徒所言肉体与圣神争战之价值。}\par

我们也读到宗徒将这场争战为我们的好处安置在我们肢体中：\BibleQuote{因为肉体的私欲相反圣神，圣神的私欲相反肉体；两者彼此敌对，致使你们不能做你们所愿意的。}\newline{}\BibleRef{迦 5:17} 在此，你也有一个仿佛是上主安排和计划在我们身体内植入的斗争。因为当一件东西存在于所有人中，毫无例外，你还能怎么想，只能认为它属于人类本性受造后的实质，几乎像是自然的；当一件东西被发现与生俱来，并随着成长而增长，我们怎能不相信这是上主的旨意所植入的，不是为损害他们，而是为帮助他们？但这场争战，即肉体与圣神的争战，他告诉我们其原因：使你们不能做你们所愿意的。因此，如果我们实现了天主安排我们不实现的事，即我们不按自己所喜欢的去做，我们怎能不认为这对我们有害？这场争战由造物主的安排植入我们内，在某种方式上对我们有益，呼唤并推动我们达到更高的境界；如果它停止了，那么另一方必定会带来一种危险的平安。\par

% structure: p0025 source=h2 target=subsection reason=relative-heading-rank; base=h2

\subsection{第八章}

\Emph{问题：为何在宗徒那一章中，在讲了肉体和圣神的私欲彼此对立之后，他又增添了第三样东西——人的意志。}\par

杰尔玛诺：虽然我们如今似乎已对那意义略有所悟，但仍未能彻底理解宗徒的意思，故切愿您为我们更清楚地阐明。因为此处似乎指出了三件事：首先是肉身与心神的斗争，其次是心神对肉身的渴望，第三是我们自己的自由意志，它似乎被置于两者之间，且论及它时说道：\Quote{你们不可做你们所愿意的。}关于这主题，虽然如我所说，我们已从您所解释的意义中收集了一些暗示，但既因这次会谈提供了机会，我们仍渴望它被更充分地解释给我们。\par

% structure: p0028 source=h2 target=subsection reason=relative-heading-rank; base=h2

\subsection{第九章}

\Emph{对正确提问者的理解之答复。}\par

达尼尔：辨别问题的区分与旨趣属于理解力；而理解自身的无知乃是知识的主要部分。因此经上记着说：\Quote{愚人若发问，便算为智慧，}\BibleRef{箴 17:28}因为虽则发问者对所提问题的答案无知，但因他明智地发问并学会他所不知的，这事实本身就被算作他的智慧，因为他明智地发现了自己先前所不知的。据此，按照您这区分，看来宗徒在此段中提到了三件事：肉身对心神的贪欲，心神对肉身的贪欲，二者彼此相争，其因由似乎在于：\Quote{使你们不能做你们所愿意的。}如此便剩下第四种情形，即您所忽略的：我们做了我们不愿做的事。现在，我们必须首先发现那两种欲望（即肉身与心神的欲望）的意义，进而学习探讨处于二者之间的我们的自由意志，最后同样地，我们才能看出什么不属于我们的自由意志。\par

% structure: p0031 source=h2 target=subsection reason=relative-heading-rank; base=h2

\subsection{第十章}

\Emph{\Quote{肉身}一词并非只有单一含义。}\par

我们发现圣经中\Quote{肉身}一词有多种不同含义：因为它有时指整个人，即由灵魂与身体所构成者，如\Quote{圣言成了肉身，}\BibleRef{若 1:14}及\Quote{凡有血肉的，都要看见天主的救恩。}\BibleRef{路 3:6}有时指有罪及属血肉的人，如\Quote{我的神不永远住在人里面，因为他们是血肉。}\BibleRef{创 6:3}有时指罪本身，如\Quote{你们不属于血肉，而属于圣神，}\BibleRef{罗 8:9}又\Quote{血肉之体不能承受天主的国：}随后有\Quote{腐朽的不能承受不朽的。}\BibleRef{格前 15:50}有时指血亲及亲属关系，如\Quote{看，我们是你的骨肉，}\BibleRef{撒下 5:1}宗徒说：\Quote{或者我可能激动我同胞的嫉妒，以拯救他们中的一些人。}\BibleRef{罗 11:14}因此我们必须探究，在这四种含义中，我们应取哪一种来理解此处的\Quote{肉身}，因为显然它不可能如同在\Quote{圣言成了肉身}及\Quote{凡有血肉的都要看见天主的救恩}这两处那样。它也不能与\Quote{我的神不永远住在人里面，因为他们是血肉}同义，因为这里的\Quote{肉身}并非如那里仅指有罪的人——当祂说\Quote{肉身贪欲相反圣神，圣神也相反肉身。}\BibleRef{迦 5:17}祂并非谈论物质的东西，而是指在同一个人内，随时间的变换与更迭，或同时或分别相争的现实。\par

% structure: p0034 source=h2 target=subsection reason=relative-heading-rank; base=h2

\subsection{第十一章}

\Emph{宗徒在此处所说的\Quote{肉身}是什么，以及肉身的贪欲是什么。}\par

因此，在此处我们应将\Quote{肉身}理解为不是指人（即其物质成分），而是指血肉的意志和邪恶的欲望，同样\Quote{心神}也不是指任何物质之物，而是指灵魂的良善及属灵的欲望；这一含义蒙福的宗徒刚在开头清楚地指明了，他说：\Quote{但我告诉你们：你们应随从圣神的引导行事，决不要满足本性的私欲，因为本性的私欲相反圣神的引导，圣神的引导相反本性的私欲；二者互相敌对，致使你们不能做你们所愿意的事。}既然这两种（即肉身的欲望和心神的欲望）共存于同一个人内，便在我们内产生了每日进行的内部争战，而肉身的贪欲冒然冲向罪恶，在那些与眼前安逸相关的享乐中放纵。另一方面，心神的欲望与之相反，并愿完全投入属灵的努力中，以致它甚至想要摆脱肉身必要的使用，渴望如此持久地被这些事物所占据，以至于不愿分担对肉身软弱的任何忧虑。肉身喜爱放纵和淫欲；心神甚至不容忍自然的情欲。一个贪睡且饱食；另一个则以守夜和斋戒为滋养，以至于不愿为生活的必要而容许睡眠和食物。一个渴望被万物丰盈所充实，另一个即使缺乏每日微薄食物的供应也满足。一个寻求通过沐浴而容光焕发，每日被众多奉承者围绕；另一个却以污垢和灰尘为乐，喜爱那不可接近的荒漠的独处，并惧怕所有凡人的接近。一个依靠人的尊重和赞美而活，另一个却以所受的伤害和迫害为荣。\par

% structure: p0037 source=h2 target=subsection reason=relative-heading-rank; base=h2

\subsection{第十二章}

\Emph{我们的自由意志是什么，它立于肉身的贪欲与心神之间。}\par

在这两种欲望之间，灵魂的\OriginalTerm{自由意志}{free will}处于一种多少该受责难的中介地位，既不以罪恶的过分为乐，也不安于德行的痛苦。它试图如此克制肉体的情欲，却又不愿承受必要的痛苦；想获得肉身的贞洁，却不克已肉身；想获得心灵的纯洁，却不劳苦守夜；想与肉体的安逸一起充满灵性德行；想达到忍耐的恩宠，却不受矛盾的刺激；想实践基督的谦卑，却不丧失世俗的荣耀；想追求宗教的淳朴，同时怀有世俗的野心；想事奉基督，却不愿失去人的赞美和好感；想宣认真理所要求的严格，却不触怒任何人：总之，它渴望追求将来的福乐，而不失去眼前的福乐。若非从两方面兴起的这些争斗干扰并摧毁了这种\OriginalTerm{冷淡}{lukewarmness}状态，这自由意志绝不会引领我们达到真正的成全，反而会将我们投入最可怜的冷淡境况，使我们如同那些在\BibleBook{}{默示录}中受主谴责的人：\Quote{我知道你的作为：你不冷也不热。我巴不得你或冷或热。但是你现在是温的，我就要立刻从我口中把你吐出来；}\BibleRef{默 3:15}因为当我们顺从这自由意志，想要在这松懈方向上放任自己时，肉体的欲望立刻升起，以它们罪恶的情欲伤害我们，不容我们继续停留在我们所喜爱的纯洁状态中，并吸引我们走上那应当畏惧的冷酷多刺的快乐之路。再者，若我们被灵性的热情所点燃，想要根除肉体的作为，不顾人的软弱全然把自己提升到追求德行的过分努力中，肉体的脆弱便出现，召回我们，约束我们，使我们不过度放纵那对我们有害的灵性；结果，由于这两种欲望在这类争斗中相互矛盾，灵魂的自由意志——它既不愿完全屈服于肉体的欲望，也不愿投身于德行所要求的努力——仿佛被一个公平的天平所调节，而这两者之间的争斗阻碍了那种更危险的灵魂自由意志，并在我们身体的天平上造成某种公平的平衡，极准确地标出血肉与灵性的界限，不容被灵性热情点燃的心灵向右摇摆，也不容许肉体因罪的刺痛向左倾斜。既然这争斗每日在我们内有益地进行，我们就被极有益地驱使我们来到那我们所不喜欢的第四阶段，以便获得心灵的纯洁，不是靠安逸和疏忽，而是靠不断的努力和心灵的痛悔；以长期的禁食、饥渴、醒悟来保持肉身的贞洁；以阅读、守夜、不断祈祷和独处的凄凉来获得心灵的决心；以忍受磨难来保存忍耐；在亵渎和众多的凌辱中事奉我们的造物主；若有必要，在世界的仇恨和敌意中追随真理；同时，由于这样的争斗在我们身体内进行，我们远离了懒惰的疏忽，并被激发去从事那违背本性的努力和对德行的渴望，我们的适当平衡被令人钦佩地保持，一方面我们自由意志的倦怠选择被灵性的热情所调节，另一方面肉体的冰冷被温和的温暖所缓和；而灵性的欲望不容许心灵被拖入放纵的放肆，同样肉体的软弱也不容许灵性被引向对圣德的不合理渴望，以免一方面兴起各种罪恶的诱惑，或另一方面，最早的罪抬起它的头，以更致命的骄傲之箭伤害我们：但从这争斗中会产生适当的平衡，为我们打开一条介于两者之间的安全可靠的德行之路，并教导基督的战士常行在君王的大道上。我们听说主在\BibleBook{}{创世记}中建造那座塔时，类似地安排了此事，那里突然发生语言混乱，制止了人类亵渎和邪恶的企图。因为若没有天主上智的安排，语言的差异在他们中引起扰乱，迫使他们因语言的变异而走向更好的状况，一场愉快而有价值的不和谐便将那些被毁灭性的合一推向灭亡的人召回得救；由于分歧兴起，他们开始体验人的软弱，而此前他们因邪恶的计谋而自高自大，对此一无所知。\par

% structure: p0040 source=h2 target=subsection reason=relative-heading-rank; base=h2

\subsection{第十三章}

\Emph{论因血肉与灵性之争而迟延的益处}\par

但由于这场冲突所造成的不和，产生了一个对我们极为有利的延迟，而这场争斗也为我们带来有益的延缓，以至于当我们在物质身体的抵抗下，无法实施心中邪恶构想之事时，有时我们会因萌生的忏悔或一些更善的念头（这些念头通常在我们执行事情时出现延迟、有反思时间之际降临）而回归更善的心意。最后，那些我们知道不会受到任何肉体阻碍而无法实现其自由意志欲望的存有——我是指魔鬼和属灵的邪恶——它们既然从更高超的天使状态堕落，我们看到其处境比人更糟，相差之大在于：由于每时每刻都有机会满足其欲望，它们不会被延迟而无法最终执行所想象的任何邪恶，因为正如其心思敏捷地构思邪恶，其本质也随时准备好并自由地执行；而它们在实现意愿时既有一条快捷方式，便没有有益的反省来修正其邪恶意图。\par

% structure: p0043 source=h2 target=subsection reason=relative-heading-rank; base=h2

\subsection{第十四章}

\Emph{论属灵邪恶的不可救治的堕落。}\par

因为一个属灵的本质、不受任何物质肉体束缚的，对于内心生出的恶念没有借口，也排除了其罪过的赦免，因为它不像我们那样受外在肉体的刺激而犯罪，而仅仅是被乖僻意志的过错所燃烧。因此它的罪是无赦的，它的软弱是无药可救的。因为它既不受任何世俗之物的引诱而堕落，也就找不到任何宽恕或忏悔的余地。由此我们可以清楚地看出，这在我们内发生的肉体与灵之间的争斗，不仅无害，反而对我们极为有益。\par

% structure: p0046 source=h2 target=subsection reason=relative-heading-rank; base=h2

\subsection{第十五章}

\Emph{论我们身上肉体情欲反对灵的价值。}\par

首先，因为它是对我们懈怠和疏忽的直接谴责，好比一位严厉的导师，从不允许我们偏离严格纪律的路线；倘若我们的疏忽稍有超出应有的严肃界限，它立刻通过欲望的刺激激发我们，责备我们，使我们回到适度的节制。其次，在贞洁和完全纯净方面，当我们因天主的恩宠而看到自己一段时间远离了肉体的玷污时，为了不让我们以为自己再也不会被肉体的运动所困扰，因而心中暗暗自高自大、膨胀起来，仿佛我们不再带着肉体的腐朽，它（肉体的情欲）就谦卑地制止我们，通过它的刺痛提醒我们不过是人。因为我们在其他种类的罪中，甚至在更坏、更有害的罪中，常常不用深思就跌倒，而且不太容易为犯这些罪而感到羞愧，然而在这个特殊的罪上，良心尤其谦卑，借着这种幻觉，它因回想被忽略的情欲而受到刺痛，因为它清楚地看到，自己因自然情绪而变得不洁，而这些情绪在它因属灵的罪而更不洁时，它一无所知；于是它立刻回到先前懒惰的痊愈之道，被警告：既不应信赖过去在纯洁上的成就（因为它看到哪怕稍微偏离主就会失去这成就），也不应认为除了靠天主的恩宠外，还能获得这种纯洁的恩赐，因为实际经验以某种方式教导我们：如果我们渴望达到持久的心之完美，必须不断努力获得谦逊之德。\par

% structure: p0049 source=h2 target=subsection reason=relative-heading-rank; base=h2

\subsection{第十六章}

\Emph{论若无肉体的兴奋使我们谦卑，我们会更严重地跌倒。}\par

因此，我们先前提到的那种权势可以作证：若因这种纯洁而产生的骄傲比一切罪过和邪恶更危险，而我们就此对于任何高度的完美贞洁都得不到赏报，因为这些权势被认为从未经历过这样的肉体情欲，却仅仅因心中的骄傲从高尚的天上地位被永远毁灭。因此，我们将全然绝望地冷淡，因为在我们的身体或心灵中，不会有我们自己懈怠的警告被植入，也不会努力去达到完美的炽热，甚至保持严格的俭朴和禁欲，若非这肉体的兴奋升起，使我们谦卑、挫败我们，促使我们渴望并关注从属灵的罪中洁净自己。\par

% structure: p0052 source=h2 target=subsection reason=relative-heading-rank; base=h2

\subsection{第十七章}

\Emph{论阉人的冷淡。}\par

最后，正因如此，我们在那些阉人身上常常发现这种心灵的冷淡，因为他们可以说免除了肉体的需要，就自以为也不需要身体上的禁欲之苦或心灵的痛悔；由于这种无忧无虑而变得懈怠，他们既不真正努力寻求或获得心灵的完美，甚至也不追求从属灵的缺点中洁净。这种由他们肉体状态而来的状况变成了天性，这完全是更糟的状态。因为从冷淡状态进入不冷不热状态的人，被主的话语定为更可憎的。\par

% structure: p0055 source=h2 target=subsection reason=relative-heading-rank; base=h2

\subsection{第十八章}

\Emph{问题：属血肉的人与属自然的人有何区别。}\par

\Person{Germanus}：您似乎已经非常清楚地阐明了肉体与灵之间争斗的价值，以至于我们相信它可以被我们以某种方式把握；因此我们希望您也同样向我们解释这一点：即属血肉的人和属自然的人有什么区别，或者说为什么属自然的人比属血肉的人更坏。\par

% structure: p0058 source=h2 target=subsection reason=relative-heading-rank; base=h2

\subsection{第十九章}

\Emph{论灵魂的三重状态之回答。}\par

\Person{达尼尔}：按照圣经的陈述，有三种灵魂：第一种是肉性的，第二种是属本性的，第三种是属神的：我们发现宗徒这样描述它们。因为关于肉性的，他说：\Quote{我给你们喝奶，不是吃肉，因为你们那时还不能。其实你们现在也不能，因为你们仍是肉性的。} 又：\Quote{因为在你们中间有嫉妒和纷争，你们岂不是肉性的吗？} \BibleRef{格前 3:2} 关于属本性的，他也这样说：\Quote{然而属本性的人不能领受天主圣神的事，因为在他看那是愚妄。} 但关于属神的：\Quote{但属神的人判断万事，却没有人能判断他。} \BibleRef{格前 2:14} 又：\Quote{你们属神的人，要以温柔的心神教导这样的人。} \BibleRef{迦 6:1} 因此，虽然我们在弃绝时已不再属肉性，即我们开始与世俗中的人隔离，与肉身的公开玷污无关，我们仍须小心，竭尽全力立即达到属神的境界，以免我们自欺，因为我们在外表上似乎已弃绝了这个世界，摆脱了肉性淫乱的污秽，以为由此已达到完美的顶峰；于是变得粗心大意，不再关心洁除其他情欲，因而停留在两者之间，无法达到属神境界的进步；或者因为我们以为，若在外表上与世界及其享乐隔离，就算充分完美，或者因为我们脱离了败坏和肉性的交往，从而发现自己处于最糟糕的温吞状态，发现我们被主从口中吐出来，正如祂所说的：\Quote{我愿你或冷或热。但你是温吞的，我必将你从我口中吐出。} \BibleRef{默 3:15} 主宣称，那些祂曾慈爱地怀揣在腹中、却可耻地变温吞的人，将被吐出并弃绝于祂的怀抱，这不是没有理由的：因为他们本可给祂提供有益健康的养料，却宁愿被从祂心中撕去：因此他们变得远不如那些从未成为主口中食物的人，正如我们厌恶因恶心而反胃的东西。因为冷的东西入口后会被温暖，且令人满意地接纳并产生好结果。但曾经因其可悲的温吞而被拒绝的，我们不仅不敢用嘴唇触碰，甚至远远看到都会极度厌恶。所以，说这样的人更糟是恰当的，因为一个肉性的人，即世俗人和外邦人，更容易被引向得救的皈依和完美的顶峰，胜过那些发了隐修士愿，却未按规则掌握完美之道，从而永远从属神热忱之火退缩的人。因为前者最终被肉身的罪恶击垮，承认自己的污秽，在痛悔中从肉性的污秽奔向真正洁净的源泉和完美的顶峰，在惊惧于自己所处的无信仰之冰冷状态时，被圣神之火点燃，更容易飞向完美。因为一个起初温吞，并开始滥用隐修士之名，未曾以应有的谦卑和热忱踏上这条修道路的人，一旦被这可悲的瘟疫感染，仿佛因此而松弛，便无法再自行辨别何为完美，也无法听从他人的劝告。因为他在心里说主告诉我们的：\Quote{因为我是富足的，已经发了财，一样都不缺；} 于是立刻对他应用以下的话：\Quote{但你是那困苦、可怜、贫穷、瞎眼、赤身的。} \BibleRef{默 3:17} 他比世俗人更糟之处在于：他不知道自己困苦、瞎眼、赤身或需要洁净，也不需要任何人的指导和教诲；因此他听不进中肯的劝告，因为他没有意识到自己因隐修士之名而受重负，在众人眼中被贬低，而实际上每个人都认为他是圣人，视他为天主的仆人，他将来必受更严格的审判和惩罚。最后，我们何必再详述那些我们已经充分发现并经验证实的事呢？我们常见那些冷淡和肉性的人，即世俗人和外邦人，获得属神的温暖；但温吞和\Quote{属本性}的人却从未如此。我们在先知书中也读到，这些人为主所憎恶，以致命令属神和有学问的人停止警告和教导他们，不要将生命之言的种子撒在贫瘠、不结果、被毒荆棘缠绕的土地上；而应鄙视这种土地，反而耕种休耕地，即将他们所有的关怀、教导和对生命之言的热忱转移到外邦人和世俗人身上：正如我们所读：\Quote{主对犹大人和耶路撒冷的居民这样说：你们要开垦荒地，不要撒种在荆棘中。} \BibleRef{耶 4:3}\par

% structure: p0061 source=h2 target=subsection reason=relative-heading-rank; base=h2

\subsection{第二十章}

\Emph{论那些弃绝世界不善的人。}\par

最后，我羞惭地说，我们发现许多人所作的弃绝是这样的：他们并未改变以往的罪过和习惯，只是改变了生活状态和世俗的服装。因为他们热切地积聚他们从未有过的财富，或者不放弃他们已经拥有的，更可悲的是，他们甚至以这样的借口努力增加财富：他们断言，他们应当用这些财富永久地支持他们的亲戚或弟兄，或者他们以建立修会的名义囤积财富，幻想自己能以院长身份主持这些修会。但如果他们真诚地寻求成全之道，他们反而会竭尽全力达到这一点：即他们不仅脱去财富，也脱去所有以往的爱好和事务，毫无保留地完全将自己置于长老们的指导下，从而不仅不为他人操心，甚至不为自身担忧。然而相反，我们发现他们热切地想凌驾于弟兄之上，却从不服从自己的长老，以骄傲为起点，他们很乐意教导别人，却不费心去学习或实践他们应当教导的内容：因此其结果必然如救主所说，成了\Quote{瞎子领瞎子}，结果\Quote{两人都掉在坑里}。\BibleRef{玛 15:14}这种骄傲虽然只有一种，却采取两种形式。一种形式不断装出严肃庄重的样子，另一种则放荡不羁地发出傻笑和大笑。前者喜欢不说话；后者觉得被约束在静默中很难受，毫不顾忌地谈论不合宜和愚昧的事，却耻于被人认为不如别人或不如别人有见识。前者因骄傲寻求圣职，后者鄙视圣职，因为它认为自己不适合或低于自己以前的尊荣、生活和出身应得的功绩。这两种中哪一种更坏，每个人必须自己考虑和判断。无论如何，不顺服的性质是一样的：一个人违背长老的命令，无论是出于工作的热忱还是贪图安逸；为了睡觉而破坏修院的规矩与为了警醒而破坏它同样有害；为了读书而违抗院长的命令与为了睡觉而轻慢它没有区别；若因你的斋戒或你的早餐而忽视一个弟兄，在骄傲的动机上没有差别：只是那些似乎披着德性外衣、以属灵形式显现的过错，比那些公开出于肉欲享乐的过错更严重、更难治愈。因为后者如同完全明显可见的疾病，被对付并治愈；而前者因被德性的外衣遮盖，未被治愈，使受害者陷入更危险和致命的病患中。\par

% structure: p0064 source=h2 target=subsection reason=relative-heading-rank; base=h2

\subsection{第二十一章}

\Emph{那些轻视大事而忙于琐事的人。}\par

因为我们怎能说明，我们看到一些人，在初期的弃绝热情中，他们舍弃了田产、巨额财富和世俗服务，投奔隐修院，却仍然热切地投身于那些不能完全割舍、而且在此生活方式中无法缺少的事情，即使它们是微小琐碎的事？以致对他们而言，对这些琐事的焦虑超过了他们对所有财产的爱。他们确实不会因忽视了更大的财富和产业而获益，如果他们只是将对它们的感情（正是为了轻视它们）转移到微小琐碎的事上。因为贪恋和贪婪之罪，他们在真正贵重的事上可能不会犯，却在更普通的事上保留着，从而表明他们并未除去以前的贪婪，只是改变了对象。如果他们过分关心席子、篮子、毯子、书籍以及诸如此类的琐事，同样的情欲仍然束缚着他们。他们竟如此嫉妒地维护自己对这些东西的权利，以致因它们与弟兄们生气，更糟的是，他们不羞耻于为此争吵。他们仍被从前贪婪的恶果所困扰，不满足于拥有因身体需要和要求而迫使隐修士拥有的东西，按照共同的数量和尺度，而是也在此表现出他们内心的贪婪，试图拥有那些他们必须使用的东西，比别人的更好；或者，超越一切应有的界限，将本应属于所有弟兄共同拥有的东西，作为自己特殊和私有的财产，防范他人触碰。仿佛金属的差异，而不是贪婪的情欲才是重要的；仿佛因大事而生气是错的，而因琐事生气则无罪；仿佛我们舍弃所有贵重物品，正是为了更容易学会对琐事毫不在意！因为，一个人是在巨大华丽的事物上还是在最微小的琐事上向贪恋让步，有什么区别呢？只不过，如果一个人轻视了大事却为小事所奴役，我们应当认为他更糟。因此，这种弃绝世俗并未达到内心的成全，因为它虽被视为贫穷，却仍保持财富的心。\par

\SourceRecord{Translated by C.S. Gibson.；Nicene and Post-Nicene Fathers, Second Series；Vol. 11.；Edited by Philip Schaff and Henry Wace.；Buffalo, NY: Christian Literature Publishing Co.,；1894.；<http://www.newadvent.org/fathers/350804.htm>.}

% structure: p0001 source=h1 target=section reason=page-title

\section{第五次会谈}

\Signature{圣若望·卡西安}\par

塞拉皮雍院长的会谈：论八种主要过犯。\par

% structure: p0004 source=h2 target=subsection reason=relative-heading-rank; base=h2

\subsection{第一章}

\Emph{我们抵达塞拉皮雍院长的静室，并就各类过犯及克服之道进行询问。}\par

在那长老及长者的集会中，有一人名叫塞拉皮雍，特别被赋以明辨的恩宠，我认为将其会谈记录下来是值得的。因为当我们恳请他讲论克服我们过犯的方法，以便其起源和原因能更清楚地展现给我们时，他便这样开始了。\par

% structure: p0007 source=h2 target=subsection reason=relative-heading-rank; base=h2

\subsection{第二章}

\Emph{塞拉皮雍院长列举八种主要过犯。}\par

攻击人类的过犯共有八种，即：第一，\OriginalTerm{贪饕}{gastrimargia}，意为暴食；第二，邪淫；第三，\OriginalTerm{贪财}{philargyria}，即吝啬或贪爱钱财；第四，愤怒；第五，忧伤；第六，\OriginalTerm{倦怠}{acedia}，即无精打采或消沉；第七，\OriginalTerm{虚荣}{cenodoxia}，即自夸或虚妄的光荣；第八，骄傲。\par

% structure: p0010 source=h2 target=subsection reason=relative-heading-rank; base=h2

\subsection{第三章}

\Emph{论过犯的两类及其作用于我们的四种方式。}\par

在这些过犯中，有两类。因为有些是自然而来的，如贪饕；有些则是在本性之外的，如贪婪。但它们作用于我们的方式有四种。因为有些过犯若没有肉体的行动便不能完成，如贪饕和邪淫；而有些则无需任何身体行动便可完成，如骄傲和虚荣。有些过犯激动的缘由在我们之外，如贪婪和愤怒；有些则由内心感受所激起，如倦怠和忧伤。\par

% structure: p0013 source=h2 target=subsection reason=relative-heading-rank; base=h2

\subsection{第四章}

\Emph{重新审视贪饕与邪淫的情欲及其疗方。}\par

为更清楚地证明这一点，不仅凭我力所能及的简短论述，更藉圣经的见证：贪饕和邪淫虽然在我们的本性之中（因为它们有时不经心灵的挑动，仅凭肉体的冲动和诱惑便会产生），但若要完成，则必须找到外在对象，从而只通过身体行为才能生效。因为\BibleQuote{每个人受自己的私欲所诱惑。私欲怀孕后，便生出罪；罪长成后，便生出死亡。}\BibleRef{雅 1:14}第一个亚当若不是手边有物质食物并误用，就不会成为贪饕的牺牲品；第二个亚当若没有某物的引诱，也不会受到试探，当时对祂说：\BibleQuote{祢若是天主子，就命这些石头变成饼吧。}\BibleRef{玛 4:3}而且，人人都清楚，邪淫也只是通过身体行为完成的，正如天主对蒙福的约伯谈到这个灵时说：\BibleQuote{他的力量在他的腰间，他的能力在他的肚脐。}\BibleRef{约 40:16}因此，这两种过犯尤其需要借助肉体才能实现，故不仅需要灵魂的灵性照料，也需要肉体的克制；因为仅仅心灵的决心不足以抵抗它们的攻击（有时愤怒、忧郁或其他情欲可单凭心灵的效劳而不需任何肉体克制就能克服）；但也必须运用肉体的惩罚，通过斋戒、守夜和痛悔行为来实现；除此之外，还需变换环境，因为既然这些罪过是心灵和肉体双方缺陷的结果，就只能由双方共同努力来克服。虽然蒙福的宗徒普遍称一切过犯为肉体的，因为他把仇恨、愤怒、异端等列为肉体的行为\BibleRef{迦 5:19}，但为了医治它们并更精确地发现其性质，我们将它们分为两类：一类称为肉体的，一类称为灵性的。我们称那些特别与纵容肉体欲望有关，并使肉体如此着迷和满足的为肉体的，以至于有时它扰乱安息的心灵，甚至违背心灵意愿将其拖入同意其欲望。蒙福的宗徒论及这些时说：\BibleQuote{我们从前也都在它们中生活，放纵肉体的私欲，随从肉体和思念的意愿，生来就是义怒之子，和别人一样。}\BibleRef{弗 2:3}但这些仅由心灵冲动而生，不仅对肉体毫无益处，反而带来有害的软弱，并只以最悲惨的快乐为食物来喂养患病心灵的，我们称为灵性的。因此，这些只需要内心的单一药物；但那些肉体的，如我们所说，只能通过双重治疗来治愈。因此，对于那些渴望纯洁的人，极有益处的是先从自身撤去滋养这些肉体情欲的物质，因为通过这些物质，仍受邪恶影响的灵魂可能会产生这些欲望的机会或回忆。因为复杂的疾病需要复杂的治疗。必须从身体撤去会诱惑它的对象和物质，以免私欲企图爆发为行动；在心灵方面，我们也应同样谨慎地放置对圣经的勤勉默想、警惕的焦虑和退入独处，以免连思想中也产生欲望。但其他过犯，与同伴交往并非障碍，反而对真正希望摆脱它们的人有极大的用处，因为在与别人混合时，他们更经常受到责备，并且由于更频繁地被激怒，过犯的存在变得明显，从而以迅速的治疗得以痊愈。\par

% structure: p0016 source=h2 target=subsection reason=relative-heading-rank; base=h2

\subsection{第五章}

\Emph{我们的主如何独自受试探而没有犯罪。}\par

因此，我们的主耶稣基督，虽被宗徒之言宣告在各方面都曾像我们一样受过试探，却被说成是\Quote{没有罪}，\BibleRef{希 4:15}即没有这种私欲的感染，因为祂对肉性情欲的煽动一无所知，而我们即使违背自己的意愿、在不知不觉中也确实会为此困扰；因为总领天神如此描述祂受孕的方式：\Quote{圣神要临于你，至高者的能力要庇荫你，因此那要诞生的圣者，将称为天主子。}\BibleRef{路 1:35}\par

% structure: p0019 source=h2 target=subsection reason=relative-heading-rank; base=h2

\subsection{第六章}

\Emph{论主被魔鬼攻击时受试探的方式。}\par

因为那拥有天主完美肖像者，理应亲自受那些情欲的试探，即亚当在他还保持天主肖像未损时也曾受过的试探，即贪饕、虚荣、骄傲；而不是那些在违命之后、当天主肖像受损时，他因自己的过失而陷入并纠缠其中的情欲。因为正是贪饕使他取了禁树的果子，虚荣使他听见了\Quote{你们的眼将要开了}，骄傲使他听见了\Quote{你们将如同天主，知道善恶}。\BibleRef{创 3:5}我们读到，主我们的救主也曾受这三种罪的试探：贪饕，当魔鬼对祂说：\Quote{命这些石头变成饼}；虚荣：\Quote{你若是天主子，就跳下去}；骄傲，当他将世上万国及其荣华指给祂看并说：\Quote{你若俯伏朝拜我，这一切我都要给你。}这是为藉祂的榜样教我们，当我们受类似试探攻击时，应如何战胜那诱惑者。因此前后两者都被称为亚当：前者是毁灭与死亡之首，后者是复活与生命之首。藉前者全人类被定罪，藉后者全人类得释放。前者由原始未成形的泥土塑造，后者由童贞玛利亚所生。祂虽理当受试探，却不必在试探中跌倒。那战胜贪饕者，不可能受奸淫的诱惑，因为奸淫源于过剩和贪饕之根，连第一个亚当也不至于被它毁灭，除非在它产生之前就被魔鬼的诡计所欺骗而成为情欲的牺牲品。因此，天主子并非绝对地说成了\Quote{有罪的肉身}，而是\Quote{有罪肉身的形状}，因为祂虽有真正的肉身，也吃、喝、睡觉，并真正接受了钉痕，但祂内并无从堕落继承来的真正罪，只有类似的东西。因为祂未曾经历肉欲的火箭——这些火箭在我们身上甚至违背我们的意愿、从本性的构造中升起——但祂因分享我们的本性，而承担了类似的东西。因为祂既真实地履行了属于我们的一切功能，承受了所有人类的软弱，故也被认为曾有此感受，为使祂藉这些软弱似乎在祂自己的肉身上也承担了这种过犯和罪的状态。最后，魔鬼只用他曾欺骗第一个亚当的那些罪试探祂，推断祂若在那些曾推翻其前驱者的试探中失败，也会在其他事上同样受骗。但因他在第一次交锋中被击败，便不能将那从第一次过犯之根生长出来的第二次软弱加于祂。因为他看到祂甚至连产生这软弱的任何东西都未接受，期望从祂得到罪的果实是徒劳的，因为他看到祂丝毫未在自己内接纳其种子或根。然而，根据路加——他把那句\Quote{你若是天主子，就跳下去}\BibleRef{路 4:9}的试探放在最后——我们可以将之理解为骄傲之感；而较早的那个（玛窦放在第三位），即路加福音所述魔鬼在一瞬间将世上万国指给祂看并许给祂的试探，则可视为贪财之感；因为战胜贪饕后，他不敢再以奸淫试探祂，转而以贪财（他知道这是万恶之根，\BibleRef{弟前 6:10}）试探；在此再次被击败后，他不敢再以任何后续的罪攻击祂，因为他很清楚这些罪都是从这一根源而生；于是他转向最后的情欲，即骄傲，他知道即那些完美并克服了所有其他罪的人也可能受其影响，他记得自己身为路济弗尔以及许多其他天使正是因此从天庭坠落，而未受任何先前情欲的试探。因此，在我们提到的这一顺序（路加福音所给的顺序）中，那最狡猾的敌人攻击第一个和第二个亚当时所使用的诱惑和形式完全一致。因为对前者他说：\Quote{你们的眼将要开了}；对后者他\Quote{将世上万国及其荣华指给祂看}。对前者他说：\Quote{你们将如同天主}；对后者他说：\Quote{你若是天主子。}\par

% structure: p0022 source=h2 target=subsection reason=relative-heading-rank; base=h2

\subsection{第七章}

\Emph{论虚荣与骄傲如何无需肉身的协助便可成就。}\par

为循我们所拟定的顺序，继续论述其他情欲的活动方式（此分析因论述贪饕及主的试探而中断），虚荣与骄傲甚至无需身体的丝毫协助，便能臻于圆满。因为，这两种情欲需要肉体的什么行动呢？它们仅凭灵魂的同意与欲求从人获得赞美与荣耀，便足以给被掳的灵魂带来彻底的毁灭。至于前述路济弗尔\OriginalTerm{路济弗尔}{Lucifer}的那古老骄傲，又有什么肉体的行动呢？正如先知告诉我们，他只是在心中思想：\BibleQuote{你心里曾说：我要升到天上，我要高举我的宝座在天主的众星以上；我要升到高云之上，我要与至高者同等。}\BibleRef{依 14:13}正如无人煽动他骄傲，他的思想独自便成就了这完整的罪过与永恒的堕落；尤其因为他所图谋的统治权并未随后实现。\par

% structure: p0025 source=h2 target=subsection reason=relative-heading-rank; base=h2

\subsection{第八章}

\Emph{论贪婪——这非本性的情欲，以及它与我们本性的诸罪之间的区别。}\par

贪婪与愤怒，虽非同一性质（因为前者是我们本性之外的事物，而后者似乎在我们内里有其种子），却以同样方式萌发，因为大多数情况下，它们是在我们之外的事物中找到被激起的理由。因那些尚软弱的人常抱怨说，他们因他人的刺激与挑唆而陷入这些罪，并被别人的挑衅一头栽进愤怒与贪婪的情欲中。但贪婪确是我们本性之外的事物，我们可以清楚地看到这一点：即它被证明其最初起点不在我们内里，也不源于维持身魂合一及生命存在所需之物。因为显而易见，除每日饮食外，没有什么是我们共同生活实际的需求与必需品；而其他一切，无论我们以何等热忱与细心保存，生命本身的需要都表明它们是与人需要不同的事物。因此，这诱惑作为本性之外的东西，只攻击那些冷淡且根基不牢的隐修士；而那些本性之罪却连最优秀的隐修士和独修士也不停地困扰。这显然是真的，我们发现有些民族完全不受此贪欲的侵扰，因为他们从未因习惯与习俗而接受这一缺失与弱点。我们相信，洪水之前的古老世界，长久以来并不知晓这种欲望的疯狂。而对于我们中每一个彻底弃绝世界的人，我们知道这欲望无需任何困难便能根除，只要人放弃他全部财产，并以不给自己留下一文钱的方式寻求修道戒律。我们有成千上万的见证者，他们在瞬间放弃了所有财产，并如此彻底地根除了这一情欲，以至于此后丝毫不再受其困扰，尽管他们一生都要与贪饕斗争，若非以最警醒的心与身体的节制努力，便无法安全免于贪饕。\par

% structure: p0028 source=h2 target=subsection reason=relative-heading-rank; base=h2

\subsection{第九章}

\Emph{论沮丧与怠惰\OriginalTerm{怠惰}{accidie}如何像其他罪过一样，通常毫无外在刺激而出现。}\par

沮丧与怠惰通常毫无外在刺激而出现，如同我们刚才谈论的那些情欲一样：因为我们深知，它们常以最令人痛苦的方式困扰独修士和那些定居在旷野、与其他人无来往的人。任何生活在旷野并经验过内心争战的人，都能轻易凭经验证明此事实。\par

% structure: p0031 source=h2 target=subsection reason=relative-heading-rank; base=h2

\subsection{第十章}

\Emph{论六种罪过如何相互关联，而另外两种不同的罪如何彼此相似。}\par

在這八種毛病中，雖然它們的起源和對我們的影響方式各不相同，但前六種，即貪饕、邪淫、悭吝、忿怒、憂傷、懈怠，彼此之間有某種聯繫，可謂環環相扣，以至於任何一種的過度都為下一種提供了起點。因為從貪饕的過度必然生出邪淫，從邪淫生出悭吝，從悭吝生出忿怒，從忿怒生出憂傷，從憂傷生出懈怠。因此，我們必須以同樣的方式和方法與它們作戰；攻克一種之後，應始終迎戰下一種。因為一棵高大茂密的有害樹木，若其賴以生長的根部先被暴露或切斷，就更容易枯萎；一處危險的水潭，若其源頭和流動的管道被仔細堵塞，就會立刻乾涸。因此，為了克服懈怠，你必須先戰勝憂傷；為了除去憂傷，必須先驅逐忿怒；為了壓制忿怒，必須先踐踏悭吝；為了根除悭吝，必須先抑制邪淫；為了消滅邪淫，必須先懲戒貪饕之罪。但剩下的兩種毛病，即虛榮和驕傲，以與我們所論的其他毛病類似的方式相連，以致一者的增長為另一者提供起點（因為虛榮的過度激起驕傲）；但它們與前六種毛病完全不同，不屬於同一類別，因為它們不僅沒有從這些毛病中產生的機會，反而以完全不同的方式和樣式被激發。因為當其他毛病被根除時，後者反而更加茁壯；從其他毛病的死亡中，它們更強有力地萌發和生長。因此，我們受到這兩種毛病的攻擊，方式截然不同。因為我們陷入前六種毛病中的每一種，總是在被前面的毛病擊敗之時；但我們陷入後兩種，卻是在我們取得勝利，尤其是在輝煌的勝利之後。因此，所有毛病的規律是：正如它們從前面毛病的增長中產生，也因除去前面的毛病而被根除。這樣，為了驅逐驕傲，必須抑制虛榮；所以，如果我們總是戰勝前面的毛病，後面的毛病就會被抑制；通過消滅那些帶頭的，其餘的情慾將輕易平息。雖然我們所論的這八種毛病以我們所展示的方式相連和結合，但它們可以更精確地分為四組和子類。因為貪饕與邪淫有特殊聯繫；悭吝與忿怒，憂傷與懈怠，虛榮與驕傲則緊密相連。\par

% structure: p0034 source=h2 target=subsection reason=relative-heading-rank; base=h2

\subsection{第十一章}

\Emph{關於這些毛病的起源和特性。}\par

现在，逐一分述各类过失：贪饕有三类：(1) 促使隐修士在规定和适当时间之前进食的；(2) 只求填满肚腹，狼吞虎咽各种食物的；(3) 寻求珍馐美味的。这三类给隐修士带来不小的困扰，除非他同样小心谨慎地摆脱这一切。因为正如人绝不应在适当时间之前打破斋戒，我们也必须完全避免进食时的任何贪婪，以及饮食的选择和精致烹制；因为从这三个原因产生了不同却极其危险的灵魂状态。从第一种生出对隐修院的厌恶，进而对那里的生活产生不满和无法容忍，这必然很快导致退却和迅速离开。第二种点燃了奢侈和淫欲的火箭。第三种也为它的囚犯的网罗编织了贪婪的纠缠之网，并始终阻碍隐修士跟随基督完美的自我舍弃。当我们心中有这种情欲的痕迹时，我们可以通过以下方式识别它：如果一位弟兄留我们吃饭，我们不满足于吃他预备并提供给我们的美味食物，反而不可原谅地擅自请求添加别的调料或额外的东西，这是绝不应该做的，原因有三：(1) 因为隐修士的思想应当时常习惯于实践忍耐和节制，并像宗徒一样，学会在任何境况中都知足。\BibleRef{斐 4:11} 因为一个因尝到一次难以下咽的食物就心烦意乱，甚至不能短暂克服自己口味细腻的人，永远无法成功约束身体中隐秘而更重要的欲望；(2) 因为有时主人正好没有我们要求的那个东西，我们暴露了他所宁愿只有天主知道的贫困，使他为自己餐桌的匮乏和简陋感到羞愧；(3) 因为有时其他人并不在意我们要求的调味品，于是我们在满足自己味觉欲望时，却惹恼了大多数人。因此，我们务必避免这种自由。关于邪淫有三类：(1) 通过性交完成的；(2) 不接触妇女发生的，我们读到祖犹大的儿子敖难因此被主击杀；圣经称之为污秽；宗徒论及此事说：\Quote{但我对未婚者和寡妇说，他们若能像我一样，为他们更好。但他们若不能自制，就让他们结婚；因为结婚比燃烧更好。}\BibleRef{格前 7:8} (3) 在心里和意念中构思的，主在福音中说：\Quote{凡注视妇女而贪恋她的，已在心中与她犯奸淫了。}\BibleRef{玛 5:28} 有福的宗徒告诉我们，这三类必须以同一种方式灭绝。他说：\Quote{致死，}他说，\Quote{你们在地上的肢体，即淫乱、污秽、邪情等。}\BibleRef{哥 3:5} 论及其中两类，他又对厄弗所人说：\Quote{至于淫乱和污秽，在你们中连提也不可；}又说：\Quote{但你们要知道，凡淫乱者、污秽者、贪得者，即拜偶像者，在基督和天主的国里没有产业。}\BibleRef{弗 5:3} 正如我们必须同样谨慎地避免这三类，它们也都同样将我们排斥在基督的国之外。关于贪财有三类：(1) 阻止弃绝者让自己完全舍弃财物和产业的；(2) 引诱我们过分急切地重新占有那些我们已经施舍给穷人的东西的；(3) 使人贪得并获取他以前从未拥有之物的。关于愤怒有三类：一种在内心沸腾，希腊文称为\OriginalTerm{怒气}{\GreekText{θυμός}}；另一种在言语、行为、行动中爆发出来，他们称之为\OriginalTerm{愤怒}{\GreekText{ὀργή}}；宗徒论及后者说：\Quote{但现在你们要戒绝一切愤怒和怨恨；}\BibleRef{哥 3:8} 第三种不像前两种那样一时沸腾、随即消散，而是持续数日甚至长期，称之为\OriginalTerm{积怨}{\GreekText{μῆνις}}。所有这三类我们都必须同样恐惧地谴责。关于忧闷有两类：一类在愤怒平息后产生，或是因遭受损失，或目标受挫受阻所致；另一类出于无理的焦虑或绝望。关于倦怠有两类：一类使受影响的人入睡；另一类使他们放弃静修小室而逃离。关于虚荣，虽然它呈现多种形式和形态，分为不同类别，但主要有两类：(1) 我们因肉体之事和可见之事而骄傲；(2) 我们因属灵的和不可见的事而燃起虚浮赞誉的欲望。\par

% structure: p0037 source=h2 target=subsection reason=relative-heading-rank; base=h2

\subsection{第十二章}

\Emph{虚荣如何能有益于我们}\par

但在某件事上，虚荣对初学者却有用处。我指的是那些仍受肉身之罪困扰的人，例如，当他们受到邪淫之神的困扰时，他们便设想司铎之尊贵，或在众人中的声誉，好让人以为他们是圣人、无瑕的；于是，藉着这些思虑，他们驱散邪淫的污秽暗示，视之为卑贱，至少与他们的地位和名声不相称；如此，他们以较小的恶战胜了较大的恶。因为一个人受虚荣之罪的困扰，总比陷入邪淫的欲望要好——一旦陷入，要么完全无法恢复，要么在跌倒后极难恢复。一位先知以天主的口吻精妙地表达了这思想，说：\BibleQuote{为了我名的缘故，我远远地止息我的愤怒；我以我的赞颂箝制你，免得你灭亡。}\BibleRef{依 48:9} 也就是说，当你被虚荣的赞颂所束缚时，你就不可能冲入地狱的深渊，或无可挽回地陷入大罪之中。我们不必惊讶，这种情欲竟能阻止人陷入邪淫之罪；因为许多例子一再证明，一旦人被它的毒害和瘟疫所感染，它便使人不知疲倦，以致连两三天禁食都几乎感觉不到。我们常知道一些住在这旷野的人承认，当他们还在叙利亚的隐修院时，可以毫不困难地五天不进食；而如今，他们甚至刚到第三时辰就被饥饿压倒，几乎无法将每日禁食坚持到第九时辰。关于此事，院长玛加略对一位问他为何在旷野第三时辰就感到饥饿，而在隐修院时却常常整周不吃也不饿的人，作了一个极妙的回答。他说：\Quote{因为在这里，没有人看见你禁食，并用对你的称赞来喂养和支持你；但在那里，你因他人的注意和虚荣的食物而肥胖。}关于我们所说的虚荣的侵袭如何阻止邪淫之罪，\BibleBook{列王}{纪}中有一个精彩而富有寓意的形象：当以色列子民被埃及王乃苛俘虏后，亚述王拿步高上来，将他们从埃及边界带回本国，他并非要恢复他们先前的自由和故乡，而是要带他们到自己的地方，并将他们迁到比他们被囚的埃及更远的国家。这个比喻恰好适用于我们眼前的情况。因为虽然屈服于虚荣之罪比屈服于邪淫危害较小，但摆脱虚荣之统治却更为困难。因为被带得更远的俘虏，回返故乡和父辈自由自然更加困难，先知责备的话也恰当地指向他：\BibleQuote{你为什么在异乡衰老？}\BibleRef{巴 3:11} 因为一个人若没有清理自己过错的田地，就可说他确实在异乡衰老了。骄傲有两种：一是肉身的骄傲，二是精神的骄傲，后者更糟。因为它尤其攻击那些在某些善行上似乎已有进步的人。\par

% structure: p0040 source=h2 target=subsection reason=relative-heading-rank; base=h2

\subsection{第十三章}

\Emph{论这些过错以不同方式攻击我们。}\par

虽然这八种过错困扰各种人，但它们并非以同样的方式攻击所有人。因为在一个人身上，邪淫之灵占据首位；在另一个人身上，愤怒横行；在另一个人身上，虚荣称王；在另一个人身上，骄傲当道。虽然我们显然都受到所有这一切的攻击，但困难以非常不同的方式和形态临到我们各人。\par

% structure: p0043 source=h2 target=subsection reason=relative-heading-rank; base=h2

\subsection{第十四章}

\Emph{论我们必须在过错攻击时与之争战。}\par

因此，我们必须以这样的方式对抗这些毛病：每个人都应发现自己的缠扰之罪，将主要攻击指向它，用全部心思和警觉守护自己以防其侵袭，每日用斋戒的武器对抗它，时刻从心中发出叹息和呻吟的箭矢攻击它，用守夜的努力和内心的默想对付它，此外不断向天主倾洒眼泪和祈祷，持续而明确地祈求从它的攻击中获得解脱。因为一个人除非首先清楚认识到，尽管他必须日夜不懈地付出最大谨慎和警惕，以便使自己变得纯洁，但他不可能凭自己的力量和努力在这场斗争中取得胜利，否则他无法战胜任何一种情欲。即使他感到已经摆脱了这个毛病，他仍应以同样的目的探查内心深处，找出仍存在的最坏毛病，并集中圣神的所有力量特别攻击它，这样通过总是战胜较强的情欲，他将迅速而轻易地战胜其余的情欲，因为灵魂通过一系列的胜利变得更加强健，而且下次对抗较弱的情欲将确保他在斗争中更迅速地成功——就像那些为了赏报而习惯在世俗君王面前面对各种野兽的人一样，这种表演通常被称为\Quote{\OriginalTerm{混合演出}{pancarpus}}。我说，这些人首先攻击他们所见最强大最凶猛的野兽，当这些被解决后，他们就能更容易地打倒其余不太可怕和强大的野兽。同样，通过总是战胜较强的情欲，当较弱的情欲取而代之，我们将毫无风险地获得完全的胜利。我们不应认为，如果一个人特别对抗\Emph{一个}毛病，而似乎过于大意，不防备其他毛病的袭击，他就会轻易地被突然袭击所伤，因为这是不可能发生的。因为一个人若渴望洁净自己的心，并坚定心的意志对抗任何毛病的袭击，他不可能不普遍地惧怕所有其他毛病，并对它们同样小心。因为如果一个人因污染自己而玷污了其他毛病，使自己不配得到纯洁的赏报，他怎么可能成功战胜那个他渴望摆脱的情欲呢？但是，当我们心的主要意图选定一个情欲作为特别攻击的对象时，我们将更迫切地为它祈祷，以特别的焦虑和热忱恳求我们能够特别警惕它，从而获得迅速胜利。因为律法的颁布者亲自教导我们，在斗争中应遵循这一计划，不依赖自己的力量；正如他所说：\BibleQuote{你不要害怕他们，因为上主你的天主在你们中间，是一位大而可畏的天主；他要在你面前渐渐消灭这些民族，一小部分一小部分地；你不可能一下子消灭他们，免得野兽增多危害你。但上主你的天主必将他们交在你面前，击杀他们，直到他们完全灭亡。}\BibleRef{申 7:21}\par

% structure: p0046 source=h2 target=subsection reason=relative-heading-rank; base=h2

\subsection{第十五章}

\Emph{没有天主的助佑，我们无法对抗自己的毛病，并且我们不应因战胜它们而骄傲。}\par

我们也不可因战胜它们而自高自大，祂同样告诫我们，说：\BibleQuote{你们吃饱了，建造了华美的房屋居住，有了牛群羊群，金银财物充盈的时候，心高气傲，忘记了那领你们出埃及地、脱离奴役之家的上主你们的天主；祂曾带领你们走过那大而可怕的旷野。}\BibleRef{申 8:12} 撒罗满在\BibleBook{}{箴言}中也说：\BibleQuote{你的仇敌跌倒时，你不必欢喜；他倾覆时，你心不可骄傲，免得上主看见而不悦，于是转离他的怒气。}即：免得祂见到你心骄气傲，就停止攻击他，而你开始被祂抛弃，以致再次被那因天主的恩宠你先前已克服的情欲所困扰。因为先知不会这样祈祷说：\BibleQuote{上主，求你不要将向你认罪的灵魂交给野兽}，除非他知道有些人因心骄气傲被重新交给那些他们已克服的过犯，为使他们谦卑。因此，对我们来说，既有实际经验证实，又有无数圣经章节教导，我们不可能靠自己的力量战胜如此强大的敌人，除非单靠天主的援助；并且我们应当每日将整个胜利归于天主自己，正如主亲自通过梅瑟在这点上教导我们：\BibleQuote{当上主你的天主在你面前消灭他们时，你心里不要说：我因自己的义德，上主领我来占据这地；其实这些民族是因他们的邪恶而被消灭。你进去占据他们的土地，不是因为你的义德，也不是因为你心中的正直，而是因为他们行恶，所以你进去时他们被消灭。}\BibleRef{申 9:4} 请问，还有什么比这更清晰地反驳我们那种不敬的观念和狂妄，即想把我们的一切作为归功于自己的自由意志和努力？祂告诉我们：\BibleQuote{当上主你的天主在你面前消灭他们时，你心里不要说：我因自己的义德，上主领我来占据这地。}对那些睁开眼睛、侧耳倾听的人，这不明明地说：当你与肉性过犯的争战顺利，看到自己摆脱了它们的污秽和世俗的风气时，不要因争战和胜利的成功而自高自大，将之归功于自己的能力和智慧，也不要以为你靠自己的努力、精力和自由意志战胜了属灵的邪恶和肉性的罪？因为毫无疑问，若非得到主的帮助的支持和保护，你根本不可能成功。\par

% structure: p0049 source=h2 target=subsection reason=relative-heading-rank; base=h2

\subsection{第十六章}

\Emph{以色列占据其土地的七个民族的意义，以及为何有时称为\Quote{七个}，有时称为\Quote{许多}}\par

这就是主应许要给以色列子民的那七国土地，当他们出埃及时。并且，正如宗徒所说，发生他们在身上的这一切都是\BibleQuote{预像}\BibleRef{格前 10:6}，我们应当看作是为教训我们而写的。因为如此记载：\BibleQuote{上主你的天主领你进入你要去占领的土地，在你面前消灭许多民族，即赫特人、基尔加斯人、阿摩黎人、客纳罕人、培黎齐人、希威人和耶步斯人，七个比你们更多、更强的民族；上主你的天主把他们交给你，你要彻底消灭他们。}\BibleRef{申 7:1}而他们说这些民族更多，是因为过犯在数量上远远超过德行，因此在列举时，这些民族算作七个，但谈到攻击时，并不给出数目，因为这样记载：\BibleQuote{会在你面前消灭许多民族。}因为从这七重刺激和罪恶根源生出的肉性情欲的种类，比以色列的后裔更多。因为由此生出凶杀、争斗、异端、偷窃、假见证、亵渎、暴饮暴食、酗酒、诽谤、戏谑、污秽言论、谎言、伪誓、愚妄谈话、粗鄙、不安、贪吝、苦毒、喧嚷、恼怒、轻蔑、抱怨、试探、绝望，以及其他许多过犯，不胜枚举。若我们倾向于认为这些是小事，就让我们听听宗徒如何看待：\BibleQuote{你们不可抱怨，}他说，\BibleQuote{如同他们中有些人抱怨，就被毁灭者所灭；}以及论试探：\BibleQuote{我们也不可试探基督，如同他们中有些人试探了，就被蛇所灭。}\BibleRef{格前 10:9}关于诽谤：\BibleQuote{不要爱诽谤，免得你被连根拔除。}关于绝望：\BibleQuote{他们既已绝望，便放纵情欲，行各种不洁的污秽。}\BibleRef{弗 4:19}而喧嚷以及忿怒、暴怒和亵渎都是被谴责的，同一宗徒如此告诫我们：\BibleQuote{一切苦毒、忿怒、暴怒、喧嚷、亵渎，连同一切恶毒，都该从你们中除掉。}\BibleRef{弗 4:31}以及更多类似的事。尽管这些远多于德行，但若那八种主要罪恶——我们知道这些是从它们而生——先被克服，那么所有这些便立即沉沦，并与它们一同被永远消灭。因为从\OriginalTerm{贪饕}{gluttony}生出暴饮暴食和酗酒。从\OriginalTerm{邪淫}{fornication}生出污秽言论、粗鄙、戏谑和愚妄谈话。从\OriginalTerm{悭吝}{covetousness}生出谎言、诡诈、偷窃、伪誓、贪求不义之财、假见证、强暴、不人道、贪婪。从\OriginalTerm{忿怒}{anger}生出凶杀、喧嚷和暴怒。从\OriginalTerm{忧愁}{dejection}生出怨恨、懦弱、苦毒、绝望。从\OriginalTerm{懈怠}{acedia}生出懒惰、嗜睡、粗鲁、不安、游荡、身心不定、喋喋不休、好管闲事。从\OriginalTerm{虚荣}{vainglory}生出争论、异端、自夸、标新立异的自信。从\OriginalTerm{骄傲}{pride}生出轻蔑、嫉妒、不顺服、亵渎、抱怨、诽谤。至于所有这些祸患比我们更强，我们可以从它们攻击我们的方式非常清楚地知道。因为对肉性情欲的乐趣在我们肢体中的战斗力量，比德行的渴望更强大，后者只有在内心和身体最深的忏悔中才能获得。但若你只用精神之眼注视我们敌人的无数军队，宗徒在说到\BibleQuote{因为我们战斗不是对抗血和肉，而是对抗率领者，对抗掌权者，对抗这黑暗世界的霸主，对抗天界里邪恶的鬼神}时列举的，以及我们在\BibleBook{}{第十九篇圣咏}中关于义人所说的：\BibleQuote{虽有千人仆倒在你左边，万人仆倒在你右边}，那么你将清楚看见他们远比我们——我们不过是血肉和属土的受造物——更多、更强，而他们被赋予了一种属灵和非物质的实体。\par

% structure: p0052 source=h2 target=subsection reason=relative-heading-rank; base=h2

\subsection{第十七章}

\Emph{关于将列国与八种过犯相比的一个问题。}\par

\Person{革尔玛努斯}：那么，怎么会有\Emph{八种}过犯攻击我们，而梅瑟将对抗以色列百姓的列国算作\Emph{七国}，我们如何能恰当地夺取我们过犯的领土呢？\par

% structure: p0055 source=h2 target=subsection reason=relative-heading-rank; base=h2

\subsection{第十八章}

\Person{塞拉皮雍}：每个人都完全同意，有八种主要过犯影响隐修士。而且它们并非全部包含在列国的预像中，原因如下：因为在\WorkTitle{申命纪}中，梅瑟——更正确地说是主通过他——是对那些已经离开埃及、从最强大的民族——我指的是埃及人——中获得自由的人说话。我们发现这个预像在我们身上也成立，因为当我们摆脱了世界的陷阱时，我们发现我们已从贪饕，即肚腹和口舌的罪中得释放；像他们一样，我们要与这剩下的七国争战，完全不计已经克服的那一国。那国的土地并未赐给以色列为产业，但主的命令规定他们应立即离开并出去。因此，我们的斋戒应当适度，以免因过度禁欲——这源于肉身的软弱和虚弱——我们有必要再次回到埃及地，即回到我们先前在弃绝世界时所舍弃的贪食和肉欲。这在预像中发生在自那班人身上，他们进入德行的旷野后，又贪恋他们曾在埃及时坐过的肉锅。\par

% structure: p0057 source=h2 target=subsection reason=relative-heading-rank; base=h2

\subsection{第十九章}

\Emph{为何一国应被撇弃，而七国被命消灭之理由。}\par

但为什么那生育以色列子民的民族不被命令完全消灭，而只被命令离开其土地，而那七个民族却被命令完全消灭呢？原因在此：无论我们因之进入诸德之荒野的灵性热忱多么巨大，我们仍然不可能完全摆脱贪饕的近傍或它的服事，并且可以这样说，摆脱与它的日常来往。因为对美味和精致食物的喜好将继续存在于我们内，如同某种自然而固有之物，即使我们努力割除所有多余的食欲和欲望，既然它们不能被完全消灭，就应当回避和躲避。关于这些，我们读到：\BibleQuote{不要挂虑肉身，而满足它的欲望。}\BibleRef{罗 13:14} 既然我们仍保留着对这种挂虑的感觉，而它被命令不全然割除，只在没有欲望的情况下保留，那么显然，我们并没有消灭埃及民族，而是以某种方式与之分离，不思想任何关于奢侈和精致盛宴的事，而是如宗徒所说：\BibleQuote{有衣有食，就当知足。}\BibleRef{弟前 6:8} 这在法律中以预象方式这样命令：\BibleQuote{不可憎恶埃及人，因为你在他的地上作过旅客。}\BibleRef{申 23:7} 因为必要的食物不被拒绝给身体，否则会对身体有危险，对灵魂有罪。但关于那七个麻烦的毛病，我们必须尽一切方式从我们灵魂的最深处根除它们的情感。因为关于它们，我们读到：\BibleQuote{一切毒辣、忿怒、愤恨、喧嚷、亵渎，连同一切恶毒，都应当从你们中除去}；又：\BibleQuote{至于邪淫、一切不洁和贪婪，连提都不该在你们中间提起，或淫猥、轻浮之言、戏谑之语。} 我们能够割除这些从外面嫁接到我们本性的毛病的根，但却不可能割除贪饕的机会。因为无论我们进步到何种地步，我们都难免是我们所生的人。这不仅是像我们这样微末之人的生活，而且是一切达到成全者的生活和习俗都能证明的：即使他们摆脱了一切其他情欲的诱因，以完美的灵性热忱和肉身克己退隐到旷野，却仍然不能不挂虑每日的餐食和岁岁年年的食物准备。\par

% structure: p0060 source=h2 target=subsection reason=relative-heading-rank; base=h2

\subsection{第二十章}

\Emph{论贪饕的本性，可用鹰的比喻来说明。}\par

关于这情欲，无论隐修士多么有灵性和出众，都必将受其困扰，有一个极好的说明见于鹰的比喻。因为这只鸟在高空翱翔，飞越最高的云层，远离所有凡人的眼目和整个地面，却仍因肚腹的需要被迫降下，落于地面，以腐肉和死尸为食。这清楚说明，贪饕的精神不能像其他一切毛病那样被完全根除，也不能像它们一样被彻底摧毁，而只能以心神的力量抑制和约束它的一切诱因和一切多余的欲望。\par

% structure: p0063 source=h2 target=subsection reason=relative-heading-rank; base=h2

\subsection{第二十一章}

\Emph{论贪饕的持久性，如向某些哲学家所描述的那样。}\par

因为这毛病的本性被一位长老在与某些哲学家的讨论中以下述谜题巧妙地表达出来，这些哲学家以为可以因其基督徒的纯朴而像嘲弄乡下佬一样嘲笑他。\Quote{神父，}他说，\Quote{我父亲留我在众多债主的掌握中。所有其他债主我都还清了，摆脱了他们所有的追逼；但有一个债主，即使我每日偿还，也无法满足。}当他们不明白这谜题的意思，急切恳求他解释时，\Quote{我生来，}他说，\Quote{处于许多毛病的包围中。但当天主激发我渴望自由时，我弃绝了这个世界，同时放弃了我从父亲那里继承的一切财产，如此我像对待催逼的债主一样满足了它们全部，完全摆脱了它们。但我从未能完全摆脱贪饕的诱惑。因为即使我把所取食物的量减到最小，还是无法避免它每日的催促，必须不断被它\InnerQuote{讨债}，就像通过不断满足它而作无休止的偿还，并在它的索求下缴纳不绝的租税。}于是他们宣称，这个他们至今视为愚钝和乡下佬的人，竟彻底掌握了哲学的首要原则，即伦理训练，而且他们惊奇他凭本性之光学会了世俗学问所不能教的事，而他们自己虽努力长训却未能学会。关于贪饕本身，这就够了。现在让我们回到我们最初开始讨论的毛病相互关系的一般论述。\par

% structure: p0066 source=h2 target=subsection reason=relative-heading-rank; base=h2

\subsection{第二十二章}

\Emph{天主如何向亚巴郎预言以色列必须驱逐十个民族。}\par

当主与亚巴郎谈论未来时（这是你们没有问的一点），我们发现祂没有列举七个民族，而是十个，祂应许将他们的土地赐给他的后裔。\BibleRef{创 15:18} 这个数目显然是由加上崇拜偶像和亵渎而构成的，在认识天主和圣洗的恩宠之前，外邦人不敬神的群众和犹太人亵渎的群众都受这两者的统治，因为他们居住在精神的埃及。但当一个人作了弃绝，从那里出来，并藉天主的恩宠战胜了贪饕，进入精神的旷野时，他就摆脱了这三种的攻击，只需与梅瑟所列举的那七个民族作战。\par

% structure: p0069 source=h2 target=subsection reason=relative-heading-rank; base=h2

\subsection{第二十三章}

\Emph{占有他们的土地对我们有何益处。}\par

但命令我们为自身的益处去占据那些最邪恶民族的国土，其含义可这样理解。每一种恶习在心底都有自己特定的角落，它在灵魂的深处为自己占据位置，并驱赶以色列——即神圣与天上事物的默观——且从不停止与之敌对。因为美德不可能与恶习并存。\Quote{因为正义与不义有什么相通之处？光明与黑暗又有什么联系？} \BibleRef{格后 6:14} 然而，一旦这些恶习被以色列人民——即那些与它们作战的美德——所战胜，那么，贪欲和邪淫之灵在我们心中所占据的位置，立刻就会被贞洁所充满。愤怒所占据的，将由忍耐所占领。那曾由致死的悲伤所占据的，将被属神的、充满喜乐的忧伤所取代。那曾被懈怠所荒废的，将立即被勇敢所耕耘。那曾被骄傲践踏的，将被谦卑所尊崇。如此，当每一种恶习被驱逐后，它们的位置（即对这些恶习的倾向）将被相反的美德所充满，这些美德恰当地被称为以色列——即看见天主的灵魂——的子女；当这些美德将一切情欲从心中驱逐出去后，我们可以相信它们恢复了自己的产业，而非侵占了别人的产业。\par

% structure: p0072 source=h2 target=subsection reason=relative-heading-rank; base=h2

\subsection{第二十四章}

\Emph{迦南人被驱逐之地如何已分配给了闪的后裔。}\par

因为，据古代传说所载，以色列子民被带入的这些迦南人的土地，原是在世界分划时分配给闪的子孙的，后来含的后裔以暴力不义地入侵并占据了它们。在此显示了天主的公义审判：祂将那些不正当地占据他人土地的人从别人的土地上驱逐出去，并将祖先在世界分划时分配给他们的古老产业归还给另一些人。我们完全可以清楚看到，这一预表在我们身上同样成立。因为按本性，天主的旨意将我们心灵的占有权不是分配给恶习，而是分配给美德；在亚当堕落后，这些美德被逐渐兴起的罪——即迦南人——从自己的家乡驱逐出去；因此，靠着天主的恩宠，经由我们的努力和劳作，它们再次恢复到心灵中时，我们可以认为它们没有占据别人的领土，而是恢复了自己的家乡。\par

% structure: p0075 source=h2 target=subsection reason=relative-heading-rank; base=h2

\subsection{第二十五章}

\Emph{关于八种恶习含义的不同圣经段落。}\par

关于这八种恶习，我们在福音中也有以下记载：\Quote{但不洁之灵从人身上出去后，它走过干旱之地，寻找安息，却寻不着。于是它说：我要回到我出来的那房屋去；及至来到，看见里面空着，打扫干净，装饰整齐；它就去带了七个比自己更恶的灵来，一同进去住在那里；那人最后的处境就比先前更坏了。} \BibleRef{玛 12:43} 看，正如前面段落中我们读到，除了以色列子民所出来的埃及人之外，还有七个民族；同样，这里也说到七个不洁之灵返回，除了我们最初听到从那人身上出去的那一个。关于这种七重的罪过刺激，撒罗满在《箴言》中这样记载：\Quote{如果你的敌人高声对你说话，不要同意他，因为他心里有七样恶事；} 也就是说，如果贪饕之灵被战胜，便开始以羞辱了它为借口来谄媚你，仿佛请求你放松一点起初的热忱，在正常的禁欲限度与严厉克苦的标准之外向它让步一点；你不要被它的顺服所胜，也不要因为似乎暂时摆脱了肉欲，就自以为安全地从它的攻击中返回你先前的大意或对美好事物的喜好。因为被你战胜的这个灵正是说：\Quote{我要回到我出来的那房屋去；} 于是，由它产生的七种罪的灵会比你在最初战胜的那个情欲更加有害，并将立刻把你拖向更坏种类的罪过。\par

% structure: p0078 source=h2 target=subsection reason=relative-heading-rank; base=h2

\subsection{第二十六章}

\Emph{当我们战胜了贪饕的情欲后，必须努力获得其他一切美德。}\par

因此，在我们操练禁食与守斋时，一旦战胜了贪饕的情欲，必须小心切勿让我们的心灵空着，没有我们所必需的美德；反而当更热切地以它们充满我们内心的最深处，以防贪欲之灵返回时，发现我们空虚、毫无美德，它不只满足于单独为我们进入，还要将这七重罪过的刺激一同带入我们的心，使我们的最后处境比先前更坏。因为那自夸舍弃了这世界以及其中掌权的八种恶习的灵魂，后来会比从前尚在世间、既未接受隐修规则亦未接受隐修士之名时更加污秽不洁，并受到更严厉的惩罚。这七个灵之所以被说成比先前出去的那个更坏，原因在于：对美食的贪爱（即贪饕）本身并不有害，若非它为其他情欲打开了门——即邪淫、贪婪、忿怒、沮丧和骄傲，这些都显然对灵魂有害并主宰灵魂。因此，一个人若希望仅靠禁食——即身体的斋戒——来获得完美的洁净，他永远不能达到；除非他知道，他应当为此目的而操练禁食：即通过禁食使肉体降伏，以便更轻易地进入与其他恶习的战斗，因为肉体尚未习惯于贪饕和暴食。\par

% structure: p0081 source=h2 target=subsection reason=relative-heading-rank; base=h2

\subsection{第二十七章}

\Emph{我们与过失的战斗并非照列表次序展开。}\par

但你们必须知道，我们的战斗并非都以同样的顺序进行；因为，正如我们提到过，攻击并非总以同样的方式临于我们，我们每个人都应当根据特别攻击他的攻击特性来开始战斗；这样，一个人必须首先与列表中排在第三位的过失作战，另一个人则与第四或第五位的过失作战。而且，根据过失对我们的控制程度以及它们攻击的特性可能要求的那样，我们也应当调整我们战斗的顺序，以致胜利和凯旋的喜乐结果能确保我们获得心灵的纯洁和全然的成全。\par

\Person{塞拉皮翁院长}就这样向我们讲述了八种主要过失的本性，并且如此清楚地阐明了潜藏在我们内的各种私欲偏情——它们的起源和联系，尽管我们每日受其折磨，却从未能透彻理解和领悟——以致我们几乎就像在镜中看见它们陈列在我们眼前。\par

\SourceRecord{Translated by C.S. Gibson.；Nicene and Post-Nicene Fathers, Second Series；Vol. 11.；Edited by Philip Schaff and Henry Wace.；Buffalo, NY: Christian Literature Publishing Co.,；1894.；<http://www.newadvent.org/fathers/350805.htm>.}

% structure: p0001 source=h1 target=section reason=page-title

\section{第六次会晤}

\Emph{圣若望·卡西安}\par

院长德奥道罗的会晤：论圣人的死亡。\par

% structure: p0004 source=h2 target=subsection reason=relative-heading-rank; base=h2

\subsection{第一章}

\Emph{旷野的描述，以及关于圣人死亡的问题。}\par

在\Place{巴勒斯坦}地区，靠近\Place{特科亚}村——这村有幸生出了先知\Person{亚毛斯}\BibleRef{亚 1:1}——有一片广大的旷野，远远延伸，直达\Place{阿拉伯}和\Place{死海}，\Place{约旦河}的河流注入其中便消失了，那里还有\Place{索多玛}的灰烬。在这地区，曾长期住着一些修行最完美、生活最圣洁的隐修士，他们突然遭\Place{撒拉逊}强盗袭击而被杀害。我们知道，他们的遗体被邻近的主教们和整个\Place{阿拉伯}的百姓以极大的敬意抢去，安放在殉道者的圣髑中，以致两个城镇的人群聚集起来，彼此激烈相争，甚至为争夺这神圣的遗骸而动手打起来。他们怀着虔诚的热忱相互争夺，看谁更有权埋葬他们并保存其遗髑——一方夸耀自己离他们居住的地方最近，另一方则夸耀自己靠近他们的出生地。但我们对此感到不安，为自己或为一些因此事而颇感惊愕的弟兄们所困扰，便探问为何那些功德如此卓著、德行如此伟大的人竟会遭强盗杀害，为何主竟允许如此严重的罪行加于祂的仆人身上，以至将那些为众人所钦佩的人交到恶人手中。于是我们满怀忧伤地来到圣\Person{德奥道罗}面前——他是一位在实践常识方面出类拔萃的人。因为他住在\Place{塞拉}，一个介于\Place{尼特里亚}和\Place{塞特}之间的地方，距\Place{尼特里亚}的隐修院五英里，并被八十英里的旷野与我们居住的\Place{塞特}旷野隔开。当我们向他倾诉上述那些人的死亡，并对天主的极大忍耐表示惊讶——因祂竟容许如此有德之人被这样杀害，以致那些本应凭其圣德之分量解救他人脱离此类考验的人，却不能自救脱离恶人之手——并问他为何天主允许如此严重的罪行加于祂的仆人身上时，真福\Person{德奥道罗}回答说。\par

% structure: p0007 source=h2 target=subsection reason=relative-heading-rank; base=h2

\subsection{第二章}

\Emph{院长德奥道罗对向他提出之问题的答复。}\par

这个问题常困扰那些信心或知识不足的人，他们想像圣人的赏报和报酬——这些并非在今世赐予，而是为将来存留的——是在这短暂必死的一生中赐予的。但我们的希望寄于基督，并不仅限于今生，以免如宗徒所说，我们\Quote{就成了众人中最可怜的了}\BibleRef{格前 15:19}，因为我们既在今世未得任何应许，若因不信而连来世的应许也失掉，那就不对了。我们不应错误地追随他们的想法，以免因对真实解释的无知，而在受试探时犹豫、颤抖、跌倒，如果发现自己落入这般人手中；并归咎于天主不义或漠不关心人间事务——提这事几乎就是犯罪——因为祂没有在试探中保护那些正直圣洁生活的人，也没有在今世以善报善、以恶报恶；这样我们便该受先知\Person{索福尼亚}所斥责之人的审判，他说：\Quote{他们在心里说：上主不降福，也不降祸。}\BibleRef{索 1:12}或者至少我们会被列入那些以如下抱怨亵渎天主的人当中：\Quote{凡行恶的，在上主眼中看为善，并且祂喜悦他们；不然，审判的天主在哪里呢？}\BibleRef{拉 2:17}此外还加上随后同样记载的亵渎话：\Quote{事奉天主是徒劳的，遵守祂的规诫，在上主面前忧愁行走，有什么益处呢？如今我们称骄傲的人为有福，因为行恶的人得以富足，他们试探天主，竟然得免。}\BibleRef{拉 3:14}因此，为避免这种无知——它是这最致命错误的根源和原因——我们首先应当知道什么是真正的善，什么是真正的恶；这样，如果我们最终把握了这些词语在圣经中的真实含义，而非虚假的通俗含义，我们就能避免被不信者的错误所欺骗。\par

% structure: p0010 source=h2 target=subsection reason=relative-heading-rank; base=h2

\subsection{第三章}

\Emph{论世间三种事物：善、恶与无记。}\par

大体上，世上有三种事物：即善、恶与中性者。因此，我们应当知道什么是真正的善，什么是恶，什么是中性者，好使我们的信德能以真知灼见为支撑，并在一切诱惑中屹立不摇。那么，我们必须相信，在纯属人的事物中，除了灵魂的德行之外，没有真正的善；这德行以无伪的信德引领我们趋向神圣之事，并使我们恒常持守那不变的善。另一方面，我们不应称任何事物为恶，唯独罪除外，因为它使我们与良善的天主分离，而与邪恶的魔鬼结合。但那些事物是中性者，它们可以根据拥有者的意向或意愿而用于任何一方，例如：财富、权能、荣誉、体力、健康、美貌、生命本身与死亡、贫穷、身体软弱、伤害，以及其他同类事物；它们可以按照拥有者的品性和意向而趋向善或恶。因为财富常有益于我们的善，正如宗徒所说，他嘱咐\BibleQuote{今世的富人，要乐意施舍，分给穷人，为自己积储美好的根基，以准备将来}，借此\BibleQuote{他们能把握真正的生命。}\BibleRef{弟前 6:17} 并且根据福音，财富对那些\BibleQuote{藉不义的钱财结交朋友}的人而言是好事。\BibleRef{路 16:9} 然而，它们也可以被引向恶的方向，若只是为了囤积或过奢华的生活而积聚，而没有用来满足穷人的需要。至于权能、荣誉、体力、健康也是中性者，可用于善恶两边，这一点很容易证明：许多旧约的圣人享有这一切，他们拥有巨大的财富和极高的荣誉，且蒙赐体力，却仍被天主所悦纳。相反，那些不正当地滥用这些事物、为自身目的而扭曲它们的人，并非无缘无故地受惩罚或被毁灭，正如\WorkTitle{列王纪}向我们显示的那样屡见不鲜。至于生命与死亡本身也是中性者，这可由\Person{圣若望}和\Person{犹达斯}的诞生证明。因为在前者，他的生命如此有益于他自己，甚至我们得知他的诞生也给别人带来喜乐，正如我们读到：\BibleQuote{许多人要因他的诞生而喜乐；}\BibleRef{路 1:14} 但至于后者，论到他的生命则说：\BibleQuote{那人若没有出生，为他更好。}\BibleRef{玛 26:24} 此外，论到\Person{若望}和所有圣人的死亡，经上说：\BibleQuote{上主圣人的去世，在主眼中十分珍贵。} 但至于\Person{犹达斯}和像他一样的人，则说：\BibleQuote{恶人的死亡极凶恶。} 身体疾病有时何等地有益，这由满身疮痍的乞丐\Person{拉匝禄}所受的祝福显示给我们。因为圣经没有提到他别的善行或功绩，而单单因这一点：即他以极大的忍耐承受了贫乏和身体的疾病，从而被相称于在\Person{亚巴郎}怀中得一席之地的有福命运。\BibleRef{路 16:20} 至于贫乏、迫害和伤害，人人都以为它们是恶，但它们何等地有用和必要，这一点清楚地由下面的事实证明：圣人不仅从未试图躲避它们，反而要么全力追求它们，要么勇敢地忍受它们，从而成为天主的朋友，并获得永生的赏报，正如真福宗徒所颂唱：\BibleQuote{为此，我喜爱我的软弱、凌辱、急难、迫害和困苦，为基督的缘故。因为我几时软弱，正是我有能力时，因为能力在软弱中得以成全。}\BibleRef{格后 12:9} 因此，那些在今世拥有最伟大财富、荣誉和权能的人，不应被看作已从其中获得了他们的至善（因为至善只在于德行），而只是获得了中性者；因为正如对善用它们的好人而言，它们是有用和便利的（因为它们为他们提供了善工和可存留到永生的果实的机遇），同样对不正当地滥用财富的人而言，它们是无用和不合适的，并提供了犯罪和死亡的机会。\par

% structure: p0013 source=h2 target=subsection reason=relative-heading-rank; base=h2

\subsection{第四章}

\Emph{如何不能违反他人的意愿将恶强加于人。}\par

因此，在保持这些区分清楚而固定，并知道除了德行之外无善，除了罪与离开天主之外无恶之后，让我们现在仔细考虑：天主是否会让自己或其他人将恶强加于其圣人？你必定会发现绝无此事。因为若一个人不同意并抗拒，他人绝不可能将罪的恶强加于他；但若一个人因懒惰和心中的邪欲而接纳它，便会被强加。最后，当魔鬼用尽所有邪恶计谋，试图将罪的恶强加于有福的\Person{约伯}，不仅夺尽他所有财物，更在七个儿女猝死的可怕且全然意外的灾祸之后，又使他从头到脚遭受可怕的创伤和难以忍受的痛苦，却仍不能将罪的污点粘在他身上，因为他在一切中坚守，从未同意亵渎。\par

% structure: p0016 source=h2 target=subsection reason=relative-heading-rank; base=h2

\subsection{第五章}

\Emph{一个异议，天主自己如何可以说是创造恶的。}\par

\Person{日耳曼努}：我们常在圣经中读到天主创造了恶或将恶降于人，如这段：\BibleQuote{在我以外没有别的。我是上主，再没有别的；我造光，也造暗；我造和平，也造恶。}\BibleRef{依 45:6} 又说：\BibleQuote{城中若有灾祸，岂非上主所为？}\BibleRef{亚 3:6}\par

% structure: p0019 source=h2 target=subsection reason=relative-heading-rank; base=h2

\subsection{第六章}

\Emph{对所提问题的答复。}\par

特奥多雷：有时，圣经惯于以不恰当的用词，称\Quote{灾祸}为\Quote{磨难}；并非这些事物本身并按其本性是灾祸，而是因为受它们为得到益处而承受的人认为它们是灾祸。因为当天主的审判与人对话时，必须用人的语言和感受来说话。当一位医生为了健康，有充分理由给患溃疡炎症的人开刀或烧灼时，这对承受者来说被视为灾祸。马刺和鞭子对倔强的马也不愉快。再者，一切惩戒在受责者看来，当时似乎是痛苦的事，正如宗徒所说：\BibleQuote{固然各种惩戒，在受的当时，似乎不是乐事，而是苦事；可是以后，为那些受过它锻炼的人，却结出正义与平安的果实。}\BibleQuote{主爱谁，就惩戒谁；凡祂接纳的儿子，祂就鞭打。因为哪里有儿子不被父亲惩戒的呢？}\BibleRef{希 12:6}所以，灾祸有时也用以指磨难，例如我们读到：\BibleQuote{天主就后悔了，没有降下祂所说的要对他们施行的灾祸。}又说：\BibleQuote{因为祢，上主，是宽仁和慈悲的，极有耐心，无限仁慈，且乐于后悔降灾，}即祢不得不因我们的罪过而加给我们作为报应的痛苦和损失。另一位先知知道这些对某些人有益，绝不是出于对他们安全的嫉妒，而是着眼于他们的好处，这样祈祷：\Quote{上主，求祢把灾祸加给他们，把灾祸加给地上的骄傲人吧！}上主自己说：\BibleQuote{看，我要给他们降下灾祸，}\BibleRef{耶 11:11}即痛苦和损失，他们现在要为灵魂的健康而受这些的惩戒，如此终将被迫回心转意，奔向那在顺境中被他们轻视的我。因此，这些事物原本是灾祸，我们绝不可能断言：因为对许多人来说，它们促成他们的益处，并提供永恒福乐的机会。所以（回到所提的问题），所有那些被认为是由我们的仇敌或任何其他人加在我们身上的灾祸，不应算作灾祸，而应算作中性之物。因为结局不是由那个愤怒狂暴地加给我们的人所认为的那样，而是由那忍受的人所变成的那样。因此，当死亡临到一位圣人时，我们不应认为他遭遇了灾祸，而是一件中性之物；对恶人它是灾祸，对善人则是安息和脱离灾祸。\Quote{因为死亡是道路被隐藏的人的安息。}所以，善人不会因死亡而受任何损失，因为他没有遭受任何异常的事；他只不过因仇敌的罪行而得到（且不乏永生的赏报）那在自然过程中也会发生的事，并以最有成果的苦难和巨大赏报的补偿，偿还必须由不可避免的法则偿还的人类死亡的债务。\par

% structure: p0022 source=h2 target=subsection reason=relative-heading-rank; base=h2

\subsection{第七章}

\Emph{一个疑问：如果善人因死亡而获益，那么导致善人死亡的人是否有罪？}\par

革尔玛诺：那么，如果善人被杀不但不受害，反而从苦难中获益，我们怎能指控那个杀死他而对他无害反益的人呢？\par

% structure: p0025 source=h2 target=subsection reason=relative-heading-rank; base=h2

\subsection{第八章}

\Emph{对以上问题的回答。}\par

特奥多雷：我们讨论的是善恶事物本身的性质，以及我们称之为中性之物的事物；而不是行这些事的人的品性。任何恶人或坏人都不应因其恶行未能加害于善人而免受惩罚。因为义人的忍耐和善良，对导致他死亡或受苦的人毫无益处，只对那耐心承受所加给他的人有益。因此，前者因残忍暴行受到公正惩罚，因为他有意伤害他；而后者却未遭受任何灾祸，因为他以心中的善良耐心承受了他的考验和苦难，从而使得所有那些以恶意加给他的事，都转变为对他的益处，并有助于永生的福乐。\par

% structure: p0028 source=h2 target=subsection reason=relative-heading-rank; base=h2

\subsection{第九章}

\Emph{约伯受魔鬼诱惑及主被犹达斯出卖的案例；以及顺境和逆境如何对善人都有益处。}\par

约伯的忍耐并没有给魔鬼带来任何好处，魔鬼藉着试探并没有使他变得更好，而只给勇敢承受试探的约伯本人带来益处；犹大也没有因他的背叛行为有助于人类的救恩而免于永罚。因为我们不应看行为的结果，而应看行者的意图。因此，我们应始终持守这一主张：除非人因自己的怠惰或心灵软弱而接纳了恶，否则恶不能由他人加于人；正如蒙福的宗徒用圣经中的一节经文证实了我们的这一看法：\BibleQuote{但我们知道，一切事情都使爱天主的人获得益处。}\BibleRef{罗 8:28} 但他说\BibleQuote{一切事情都获得益处}时，他包括了一切事物，不仅是幸运的事，也包括那些看似不幸的事；宗徒在另一处告诉我们他自己也经历过这些，他说：\BibleQuote{用正义的武装在左右手，即透过荣耀和羞辱，透过恶名和美名，似欺骗者却是真诚的，似忧愁却常喜乐，似贫穷却使许多人富足：}\BibleRef{格后 6:7} 那么，所有那些被认为是幸运的、被称为\Quote{在右手边的}（圣宗徒以荣耀和美名来指称），以及那些被视为不幸的（他明显以羞辱和恶名来意指，并称之为\Quote{在左手边的}），若完美之人在它们临到时勇敢地承受，就成为他的\Quote{正义的武装}；因为，当他与这些对抗，并利用那些看似攻击他的武器来保护自己，如同用弓、剑和坚固的盾牌来抵挡那些将这些加于他的人时，他就获得了忍耐和良善的利益，并藉着敌人用来杀害他的那些武器，取得了坚毅的伟大胜利；只要他不因成功而得意，也不因失败而沮丧，而是始终在君王的坦途上直行，不在喜乐笼罩他时转向右手的安宁状态，也不在不幸压倒他、悲伤控制他时被驱向左手。因为\Quote{爱慕祢法律的人，大享平安，他们没有任何绊脚石。} 但对于那些随着所遇各种际遇的性质和变化而摇摆不定的人，我们读到：\BibleQuote{愚人却像月亮一样变化。}\BibleRef{德 27:11} 因为正如对完美智慧的人所说：\BibleQuote{一切事情都使爱天主的人获得益处，}\BibleRef{罗 8:28} 同样，对软弱愚昧的人则宣告：\Quote{一切事情都与愚人作对，}因为他不能从顺境中获得益处，逆境也不能使他变好。因为勇敢承受悲伤所需要的良善，与在顺境中节制是一样的；并且可以肯定，在其中一项上失败的人，也无法承受另一项。但人更容易被顺境而非不幸所击败：因为不幸有时会违反人的意愿约束他，使他谦卑，并通过极有益的悲伤使他少犯罪，变得更好；而顺境则以悦人但最有害的奉承使人心高气傲，当人们在自己的幸福前景中自认为安全时，却以更大的毁灭将其摔到地上。\par

% structure: p0031 source=h2 target=subsection reason=relative-heading-rank; base=h2

\subsection{第十章}

\Emph{论完美之人（以左右手皆能者比喻）的卓越。}\par

这些人就是圣经中比喻性地称为\OriginalTerm{左右手都灵活的人}{\GreekText{ἀμφοτεροδέξιον}}，即左右手都灵活的人，正如在\BibleBook{民长}{纪}中描述的厄胡得，\BibleQuote{他使用左手如同右手一样。}这种能力我们也能在精神上获得，如果我们正确而恰当地使用那些幸运且看似\Quote{在右手边的}事物，以及那些不幸且我们称之为\Quote{在左手边的}事物，使两者都归属于右手边，那么，按照宗徒的话，无论发生什么，在我们身上都会成为\BibleQuote{正义的武器。}因为我们看到内心的人由两部分组成，如果我可以这样表达，两只手，任何圣人都不能缺少我们称之为左手的那部分：但通过它，德行的圆满得以显现，当一个人通过巧妙运用能将两只手都变成右手时。为了使我们的意思更清楚，圣人右手是他神修上的成就，当他以热切的精神战胜欲望和情欲时，当他摆脱魔鬼的一切攻击，毫不费力地拒绝和切断所有肉性的罪时，当他超脱尘世，视一切现世和属世之物如轻烟或虚幻的影，并因它们即将消失而蔑视它们时，当他的心满溢，不仅极其渴望未来，而且更清晰地看见它时，当他更有效地以神修默观为食时，当他更明亮地看到天上的奥秘向他敞开时，当他更纯洁和更迫切地向天主倾心祈祷时，当他被精神的热情燃烧，以极度的灵魂准备转向不可见和永恒的事物，以致几乎不相信自己仍处于肉身之中时。他也有左手，当他陷入诱惑的罗网，被肉欲之火的热情燃烧，被愤怒和怒气的情绪点燃，被骄傲或虚荣所膨胀而征服，被致死的悲伤压迫，被懈怠的诡计和攻击撕裂，当他失去了所有神修的热度，并以某种冷淡和不合理的悲伤变得冷漠，以致不仅善良和点燃的思绪离开了他，而且甚至圣咏、祈祷、诵读和独居于斗室都使他厌倦，所有修德的操练因一种难以忍受和恐怖的厌恶而失去了味道。当隐修士这样被困扰时，他就知道自己受到\Quote{左手的}攻击。因此，任何人若因我们提到的那些右手之事而不被虚荣所膨胀，并勇敢地与左手之事抗争，不屈服于绝望，反而抓住忍耐的盔甲来锻炼自己的德行，这样他就能将两手都用成右手，在每种行动中都证明自己是胜利者，并从左手和右手的状况中获得胜利的奖赏。我们读到，有福的约伯就是这样得到了奖赏。他当然在右手上获得了（胜利的）冠冕，当他有七个儿子，作为富有和富裕的人行走时，每天都为他们献祭给主以清洁他们，焦虑地希望他们讨天主喜悦而非自己，当他的门向每一个陌生人敞开，当他成为\BibleQuote{瘸子的脚，瞎子的眼睛}\BibleRef{约 29:15}，当受苦者的肩膀被他羊的毛所温暖，当他成为孤儿的父亲和寡妇的丈夫，当他甚至在心中也不因仇敌的跌倒而欢喜时。同样，也是这个人，以更大的德行在左手的逆境中得胜，当他瞬间被夺去七个儿子时，不是作为父亲被苦涩的悲伤所压倒，而是作为天主真正的仆人，欣喜于造物主的旨意；当他从富人变为穷人，从富足变为赤贫，从强壮变为憔悴，从有名望可敬变为被鄙视和轻贱，却保持灵魂坚不可摧时；最后，当他失去所有财富与财产，住在粪堆上，像自己身体的严厉刽子手，用陶片刮去破裂流出的脓液，将手指深深插入伤口，从四肢各处拖出一团团蛆虫时。在这一切中，他从未陷入绝望和亵渎，也绝未抱怨造物主。不仅如此，如此沉重苦痛的重担也没有压倒他，以至于他仅存的那件用以遮盖身体的外衣——那是魔鬼无法毁灭的，因为他正穿着它——他撕下并丢弃，用它遮盖他自愿忍受的赤身，这是那可怕的强盗加给他的；他头上的头发，也是从前荣耀所存一切中唯一未被触及的东西，他剃下并丢给折磨他的人，甚至割断他那凶残的敌人留给他的那一点，并向他夸耀，用那天上的呼喊嘲弄他：\BibleQuote{我们从主手中接受了福，难道不接受祸吗？我赤身从母胎出生，也要赤身回到那里。主赐予，主收回；主所喜悦的，就这样成就了；愿主的名受赞颂。}我也完全有理由称若瑟为左右手都灵活的人，因为在顺境中他极其爱父亲，对兄弟亲切，蒙天主悦纳；在逆境中他贞洁，忠于主，在监狱中对囚犯最仁慈，不记恨，对仇敌宽厚；对他的兄弟——那些嫉妒他并尽力杀害他的人——他不仅慈爱，而且实际上慷慨。这些人以及像他们一样的人，正当地称为\OriginalTerm{左右手都灵活的人}{\GreekText{ἀμφοτεροδέξιον}}，即左右手都灵活的人。因为他们能用任何一只手作为右手，经过宗徒所列举的那些事，可以公正地说：\BibleQuote{借着左右两手的正义武器，借着尊荣和羞辱，借着恶名和美名，等等。}关于这右手和左手，撒罗满在\BibleBook{}{雅歌}中以新娘的口吻这样说：\BibleQuote{他的左手在我头下，他的右手必拥抱我。}\BibleRef{歌 2:6}这段经文表明两者都有用，但将一只手放在头下，因为不幸应当受心灵的控制，因为它们只对此有益：即暂时训练我们，为我们的救恩而管教我们，使我们在忍耐上得以成全。但她希望右手永远紧贴她，疼爱她，在蒙福的新郎的拥抱中紧紧握住她，并不可分割地与她结合。当富足和匮乏都不影响我们，前者不引诱我们陷入危险疏忽的奢侈，后者不使我们陷入绝望和抱怨，而是我们对两者同样感恩，从好与坏的命运中获得同样的益处时，我们就是左右手都灵活的人了。那位真正左右手都灵活的人，外邦人的导师，证明他自己就是这样，他说：\BibleQuote{因为我已经学会了无论在什么境况下都可以知足。我知道怎样处卑贱，也知道怎样处丰富；在任何事上，在一切事上，我都得到了秘诀：既知饱足，也知饥饿；既知富足，也知匮乏。我在那加给我力量的那一位内，凡事都能做。}\BibleRef{斐 4:11}\par

% structure: p0034 source=h2 target=subsection reason=relative-heading-rank; base=h2

\subsection{第十一章}

\Emph{论两种考验临于我们的三重方式。}\par

那么，虽然我们说试炼是双重的，即在顺境和逆境中，但你必须知道，所有人都有三种不同的试炼方式。有时是为了考验他们，有时是为了改进他们，有时则是因为他们的罪过应受惩罚。果真，是为了考验他们，正如我们读到，有福的\Person{亚巴郎}、\Person{约伯}和许多圣人忍受了无数的患难；或者如\Person{梅瑟}在\WorkTitle{申命纪}中对百姓所说的：\BibleQuote{你要记念上主你的天主在这四十年里引领你走过旷野的所有道路，为磨难你、考验你，好知道你心中的事，是否愿意遵守祂的诫命。}\BibleRef{申 8:2} 以及我们在\WorkTitle{圣咏}中读到的话：\Quote{我在默黎巴的水旁考验了你。} 对\Person{约伯}也说：\Quote{你以为我说这些话是为别的缘故，不是为使你显为正义吗？} 但为了改进，当天主为了他们微小的小罪而惩戒祂的义人，或为了提升他们达到更高的纯洁状态，就让他们经历各种试炼，为炼净他们一切不洁的思想，并用先知的话来说，就是除掉他们隐秘部分所积聚的\Quote{渣滓}，好使他们如同纯金一样被送达将来的审判，因为祂不允许在他们身上留下任何东西，让审判之火在日后用刑罚折磨他们时发现。正如这句话所说的：\Quote{义人的患难虽然繁多。} 以及：\BibleQuote{我儿，不可轻视上主的管教，受祂责备时也不要厌倦；因为上主所爱的，祂必管教，祂鞭打所收纳的每一个儿子。哪有儿子不受父亲管教的呢？如果你们不受管教，就是私生子，而不是儿子。}\BibleRef{希 12:5} 在\WorkTitle{默示录}中：\BibleQuote{凡我所爱的，我就指责并管教。}\BibleRef{默 3:19} 对耶路撒冷以拟人法，上主借着\Person{耶肋米亚}说了以下的话：\BibleQuote{因为我必完全消灭我所分散你们到的各国，但我不完全消灭你，只按正义管教你，使你不再自以为无罪。}\BibleRef{耶 30:11} \Person{达味}为这种赋予生命的洁净祈祷说：\Quote{上主，求你考验我，试炼我，煅炼我的肾和我的心。} \Person{依撒意亚}深知这种试炼的价值，也说：\Quote{上主，求你按正义管教我们，不要发怒。} 又说：\BibleQuote{上主，我要称谢你，因为你曾向我发怒，你的怒气已转消，并且安慰了我。}\BibleRef{依 12:1} 但作为对罪的惩罚，试炼的打击也会降临，例如上主威胁要向以色列百姓降下瘟疫：\BibleQuote{我要打发野兽的牙齿到他们中间，以及地上爬行的毒蛇的狂怒。}\BibleRef{申 32:24} 以及\BibleQuote{我白白击打了你们的儿女，他们仍不受管教。}\BibleRef{耶 2:30} 在\WorkTitle{圣咏}中也有：\Quote{罪人的鞭笞众多。} 在\WorkTitle{福音}中：\BibleQuote{你已经痊愈了，不要再犯罪，免得你遭遇更坏的事。}\BibleRef{若 5:14} 的确，我们还发现第四种方式，根据圣经的权威，我们知道有些苦难仅仅是为了彰显天主和祂工程的荣耀而加在我们身上，正如福音的这些话：\BibleQuote{也不是这人犯了罪，也不是他的父母，而是为叫天主的工程在他身上显扬出来。}\BibleRef{若 9:3} 又说：\BibleQuote{这病不至于死，而是为光荣天主，叫天主子因此得光荣。}\BibleRef{若 11:4} 此外还有其他种类的惩罚，有些越过了邪恶界限的人在此生就受到打击，正如我们读到\Person{达堂}、\Person{阿彼兰}或\Person{科辣黑}所受的惩罚，尤其是宗徒所说的那些人：\Quote{因此，天主任凭他们陷入可耻的情欲和败坏的心。} 这应被视为比所有其他惩罚更严重。因为关于这些人，圣咏作者说：\Quote{他们不像别人那样受苦，也不像别人那样受鞭笞。} 因为他们不配通过天主的生命之访和现世的瘟疫得到医治，正如他们\BibleQuote{既已绝望，就纵情恣欲，行各种不洁之事，贪求一切污秽。}\BibleRef{弗 4:19} 并且由于心硬，习惯于犯罪，他们在这短暂的生命和现世的惩罚中已无法得到洁净。这些人受到先知圣言的责备：\BibleQuote{我毁灭了你们中的一些人，如同天主毁灭了\Place{索多玛}和\Place{哈摩辣}，你们像从火中抽出的火把，但你们仍不归向我——上主说。}\BibleRef{亚 4:11} 以及\Person{耶肋米亚}：\BibleQuote{我击杀了你们的百姓，毁灭了他们，但他们仍不离开自己的道路。}\BibleRef{耶 15:7} 又说：\BibleQuote{你击打了他们，他们却不悲伤；你破碎了他们，他们拒绝接受管教；他们使脸面坚硬于磐石，拒绝回头。}\BibleRef{耶 5:3} 先知看到今生所有救治方法对医治他们都已徒劳，几乎已对他们的生命绝望，便宣告：\BibleQuote{风箱在火中已失败，炼匠白白地熔炼；因为他们的恶行并未被除去。称他们为废弃的银子，因为上主已弃绝了他们。}\BibleRef{耶 6:29} 上主也遗憾地对因罪孽而满身锈垢的耶路撒冷说，祂徒然用这有益的炼火对待那些心硬于罪恶的人：\BibleQuote{将空锅放在炭火上，烤热它，熔化它的铜，让它的污秽在其中熔化。费了大力，但它的厚锈仍未除去，甚至用火也不行。你的污秽是可憎的，因为我曾洁净你，你却未从你的污秽中得洁净。}\BibleRef{则 24:11} 因此，像一位熟练的医生，试过了一切拯救疗法，却看到没有药物可治其病时，上主似乎被他们的不义所征服，不得不停止祂那慈爱的管教，并这样谴责他们：\BibleQuote{我不再向你发怒，我的妒火也离开了你。}\BibleRef{则 16:42} 但至于另一些人，他们的心并未因持续犯罪而变硬，也不需要那最严厉的（如果可以这样称之）灼伤性疗法，而是只需生命之言的教导就足以使他们得救——关于他们，经文说：\Quote{我要因听见他们的苦难而改善他们。} 我们很清楚，对犯大罪者所施的惩罚和报应还有其他原因——不是为了赎他们的罪，也不是为了抵消他们应得的惩罚，而是为了使活着的人受到警戒而改过自新。我们看到这些惩罚明确地加在了\Person{雅洛贝罕}（\Person{乃巴特}之子）、\Person{巴厄沙}（\Person{阿希雅}之子）、\Person{阿哈布}和\Person{依则贝尔}身上，天主的责备如此宣称：\BibleQuote{看，我要降祸于你，剪除你的后裔，杀死\Person{阿哈布}家中所有男子，无论被囚禁的还是在以色列中最后的。我要使你的家像\Person{乃巴特}之子\Person{雅洛贝罕}的家，又像\Person{阿希雅}之子\Person{巴厄沙}的家，因为你所做的一切招惹我发怒，并使以色列犯罪。狗在\Place{依次勒耳}的田间要吃\Person{依则贝尔}。如果\Person{阿哈布}死在城里，狗要吃他；如果死在田野，空中的飞鸟要吃他。}\BibleRef{列上 21:21} 以及作为最大威胁的这些话：\BibleQuote{你的尸体不得安葬在你祖先的坟墓中。}\BibleRef{列上 13:22} 这短暂的惩罚并不足以洁净那人首先铸造金牛犊、导致百姓长期犯罪和离弃上主的亵渎发明，也不足以洁净其他人的无数可耻亵渎，而是要以他们为榜样，使那些不重视将来、甚至全然不信的人，因畏惧这些惩罚而受到影响；并借此祂严厉的证明，他们可能会承认至高天主的威严并未疏于照管人间事务和日常行为，从而通过他们所极度畏惧的，更清楚地看到天主是万有的报应者。的确，我们发现即使较轻的过失，有些人在今生也受到了同样的死刑判决，就如我们之前说的那些亵渎背道之人的惩罚一样。例如在安息日捡柴的人，\BibleRef{户 15:32} 以及\Person{阿纳尼雅}和\Person{撒斐辣}，因不信之罪隐瞒了一部分财产。并非他们的罪责相当，而是因为他们是最早在新种类过犯中被发现的，因此他们既给了别人犯罪的榜样，也应当给他人惩罚和畏惧的榜样，使任何试图效法他们的人知道（即使惩罚在今生被延迟），在将来审判的考验中，他们必受同样的惩罚。由于我们在急于列举各种试炼和惩罚时，似乎与我们之前说完全人在任何试炼中都会保持坚定的主题有些偏离，现在让我们再回到主题上。\par

% structure: p0037 source=h2 target=subsection reason=relative-heading-rank; base=h2

\subsection{第十二章}

\Emph{正直的人应当如何像硬钢的印记，而非蜡做的印记。}\par

因此，正直之人的心志不应像蜡或其他柔软的材料，它们总是屈服于压在其上的形状，被其形态和印记所烙印，并保持该印记，直到另一个印记盖在上面而改变形状；结果它从不保持自己的形状，而是随着压在其上的东西扭转变化。相反，他应当像一块硬钢的印记，使心志始终保持其固有的形态和形状不变，并将自身的状态印记在发生的一切事物上，而任何发生的事都不能在它上面留下痕迹。\par

% structure: p0040 source=h2 target=subsection reason=relative-heading-rank; base=h2

\subsection{第十三章}

\Emph{一个疑问：心志能否恒常保持同一状态。}\par

\Person{革尔玛诺}：但我们的心志能否恒常保持其不变的状态，并始终处于同一境界？\par

% structure: p0043 source=h2 target=subsection reason=relative-heading-rank; base=h2

\subsection{第十四章}

\Emph{对提问者所提问题的答复。}\par

\Person{德奥多}：人必须像宗徒所说的那样，\BibleQuote{在心灵的精神上更新}\BibleRef{弗 4:23}，并每日\BibleQuote{向着前面的目标奋力进取}\BibleRef{斐 3:13}而进步；否则，若忽略此事，就必定后退，变得更糟。因此，心志不可能保持同一不变。正如一个人奋力逆流划船，必须用臂力抵抗急流以向上游前进，否则若放松双手，就会被冲向下游。因此，如果我们发现没有获得更多进步，那就是失败的明显证据；而且每当我们发现没有向上前进时，则应当毫不怀疑我们已经完全倒退了。因为如前所述，人的心志不可能保持同一状态，只要人仍在肉身内，没有一位圣人能达到所有德行的顶峰而恒定不变。因为或增或减必然发生，任何受造物都不可能完美到不受变化的影响。正如我们在\WorkTitle{约伯传}中所读：\BibleQuote{人是什么，竟能洁净？妇人所生的，怎能显示为义？看，在他圣者中，没有不变者；在他眼中，诸天也不纯洁。}\BibleRef{约 15:14}因为我们承认唯有天主是不变的，圣人先知在祈祷中这样向祂呼求：\BibleQuote{但祢却永存不变}，祂也自称：\BibleQuote{我是天主，我从不改变}\BibleRef{拉 3:6}，因为唯有祂在本性上常是善的、圆满而完美，无可增加，也无可减少。因此，我们应当以不懈的关切和焦虑，恒常献身于德行的取得，并持续操练之，以免我们一旦停止前进，就立刻后退。因为如前所述，心志不能保持同一状态，我是说，若不领受其善质的增加或减少。因为不能获得新善就是失去已有之善，因为当进步的愿望停止时，后退的危险便已临在。\par

% structure: p0046 source=h2 target=subsection reason=relative-heading-rank; base=h2

\subsection{第十五章}

\Emph{离开小室导致的损失。}\par

因此，我们应当始终闭居在小室内。因为每当人离开它，再返回并重新开始居住时，便会心神不宁。因为如果人离开了小室，便难以毫不费力地恢复那在他安居小室时所获得的坚定心志；由于他因此退步了，他不会看重那因离开小室而错过的进步——若他不允许自己离开小室，本可保有这进步——反而会庆幸自己重新获得了那已失落的状态。因为正如失去的时间无法追回，同样，错过的利益也无法恢复；因为无论日后心志多么专注，那只是当天的收益和当时的好处，并不能弥补那已永远失去的收益。\par

% structure: p0049 source=h2 target=subsection reason=relative-heading-rank; base=h2

\subsection{第十六章}

\Emph{连天上的神体也是可变的。}\par

然而，我们所说连天上的神体也是可变的，这由那些因堕落意志的过错而从其等级中坠落者所证实。因此，我们不应认为那些留在其受造时的有福状态中的神体本性不变，仅仅因为他们没有同样地被误导选择较差的部分。因为本性不变与一个人通过德行的努力并因不变天主的恩宠守护善而免于变化，是两回事。因为凡靠谨慎守护的一切，也能因疏忽而失去。因此我们读到：\BibleQuote{人未死以前，不要称他有福}\BibleRef{德 11:30}，因为只要人仍在奋战中——容我借用一词，仍在角力——即使他通常得胜并取得许多胜利的奖品，他也不能免于恐惧和对不确定结局的猜疑。因此，惟有天主被称为不变和善，因为祂的善不是努力的结果，而是本性的拥有，因此祂除了善以外，不能是别的。所以，人所能获得的任何德行都难免变化，但为使它一旦存在便能持续保存，必须用取得它时同样的谨慎和勤奋来守护。\par

% structure: p0052 source=h2 target=subsection reason=relative-heading-rank; base=h2

\subsection{第十七章}

\Emph{没有人是突然跌倒的。}\par

\BibleQuote{败坏以先，人心高傲；跌倒之前，人心邪恶。}正如房屋不会因突然坍塌而坠地，除非地基存在久远的裂缝，或因居住者长期疏忽，起初只是一点渗水，逐渐侵蚀，保护性的墙壁随之慢慢毁坏，因长期忽视，缺口变大，最终断裂，一旦暴风雨和雨水如河流般涌入：因为\BibleQuote{懒惰使房屋倾颓，双手无力，屋漏滴穿。}同样的事发生在灵魂的灵性层面，同一撒罗满用另一些话告诉我们：\BibleQuote{滴漏的水，在风雨之日，使人离家。}因此，他巧妙地将心灵疏忽比作屋顶和未加照管的瓦片，起初（姑且如此说）只有极细微的情欲滴漏渗入灵魂；但若因它们微小而忽视，那么德行的梁木便会腐朽，并被罪过的大风暴卷走，藉此\BibleQuote{在风雨之日}，即在诱惑之时，魔鬼的攻击将侵袭我们，灵魂将从德行的居所被驱逐出去；只要它保持一切警醒的勤奋，它本应如同在自己的房屋中一样安居。\par

我们听完这话，对这场灵性盛宴如此欣喜，以致这次会谈带给我们的心灵愉悦，超过了先前因圣人去世而经历的悲伤。因为我们不仅在我们曾困惑的事上得到了教导，而且通过提出那个问题，我们还学到了一些我们理解力太小以致无法询问的事。\par

\SourceRecord{Translated by C.S. Gibson.；Nicene and Post-Nicene Fathers, Second Series；Vol. 11.；Edited by Philip Schaff and Henry Wace.；Buffalo, NY: Christian Literature Publishing Co.,；1894.；<http://www.newadvent.org/fathers/350806.htm>.}

% structure: p0001 source=h1 target=section reason=page-title

\section{\WorkTitle{第七次会议}}

\Signature{圣若望·卡西安}\par

\Person{塞林}院长的第一次会议。论心灵的反复无常及属灵的邪恶。\par

% structure: p0004 source=h2 target=subsection reason=relative-heading-rank; base=h2

\subsection{第一章}

\Emph{论\Person{塞林}院长的贞洁。}\par

我们渴望向热心之人介绍\Person{塞林}院长，一位极其圣洁与节制的人，他名副其实，如一面镜子，我们以特有的尊崇将他置于众人之上。我们认为，藉着将他的会议收录于本书，便实现了我们的愿望。此人除了在其他德行上出类拔萃——这些德行不仅在他的言行中，而且藉着天主的恩宠也在他的面容上闪耀——还因特殊的祝福获赐节制的恩赐，以致即使在睡眠中，他也从未感到被自然的冲动所扰乱。他是如何藉着天主恩宠的助佑，达到了如此奇妙的内身纯洁，似乎超越了人性的条件，我认为应当首先加以说明。\par

% structure: p0007 source=h2 target=subsection reason=relative-heading-rank; base=h2

\subsection{第二章}

\Emph{上述长者关于我们思想状况的提问。}\par

于是，这位长者日夜祈祷、禁食守夜，不知疲倦地恳求内心与灵魂的贞洁；当他看见自己已获得所愿所求的，心中一切肉欲的情欲都已死去，便仿佛因纯洁之甘美而被唤醒，对贞洁的热忱点燃了更炽烈的渴望，开始致力于更严格的禁食和祈祷，好使天主恩宠已赐予他内在之人的这种情欲的克制，也能延伸到外在的纯洁，以至于他不再受任何简单自然之动的影响，即使是在孩童和婴儿身上也会激起的动。他因所获恩赐的经验——他知道这恩赐并非靠自己劳苦的功德，而是靠天主的恩宠——而更加热切地渴望同样获得这外在的纯洁，因为他相信：既然天主白赐了那更伟大的、非人力所能获得的内心纯洁，那么祂必更易将肉体的这种刺激连根拔除（这刺激有时甚至凭人的技艺和技能藉药剂、疗法或手术也能根除）。当他以不断的恳求和眼泪不知疲倦地致力于已开始的祈求时，有一位天使在夜间显现于他，似乎剖开他的腹部，从他的内脏中取出一团火一般的肉欲体液，抛掉，然后一切复归原状；他说：\Quote{看}，\Quote{你的肉欲刺激已被除去，务要确信，你今天获得了你忠实求告的那持久的身体纯洁。}关于赐予这位杰出之人的特殊恩宠，如此简略地叙述便已足够。但我认为无需提及他与其他善人所共有的那些德行，惟恐那段关于此人之名的特别叙述，会剥夺他人所独有的特点。因此，我们怀着与他会议并受他教导的极大热忱，安排在大斋期拜访他；当他极其平静地询问我们思想的特性、我们内心之人的状况，以及我们在旷野久居对其纯洁有何助益时，我们向他倾诉了这些抱怨：\par

% structure: p0010 source=h2 target=subsection reason=relative-heading-rank; base=h2

\subsection{第三章}

\Emph{我们关于思想反复无常的答复。}\par

在这里度过的时光、独处的居所和默想——您认为我们藉此应当已获得内心之人的完美——只为我们成就了这一点：即教导我们无法成为什么，却未使我们成为我们所努力要成为的。我们并不觉得因这知识，我们获得了所渴望的纯洁的任何确定稳固，或任何力量与坚定；只觉增添了困惑与羞愧：因为我们所有修行的默想，在日常学习中旨在达到这一点，并努力从颤抖的起点达到确定无误的技能，开始认识起初只是模糊知晓或全然不知的事物，并通过（可以说）稳步迈向那种修习状态，使自己毫无困难地完全熟悉它；然而我发现，正相反，在我为这纯洁渴望而奋斗时，我只达到了能认识到自己不能成为什么的地步。因此我感受到，这一切心灵的痛悔带给我的只有苦恼，以致眼泪从不缺乏，然而我并未停止成为我不应当成为的。那么，知道了什么是最好的，却又无法在知道后达到，又有何益？因为当我们感到内心的目标指向我们所意图的时，心神不知不觉又回到先前游荡的思绪，更猛烈地滑回去，被日常的纷扰占据，被无数俘获它的事物不断牵引，以至于我们几乎对所渴望的改进绝望，这一切修行似乎都无用。因为那每一刻都漫无目标游荡的心神，当被带回对天主的敬畏或属灵默观时，还未在其中稳固，便又飞离和迷失；当我们被唤醒，发现它已偏离设定的目标，想要把它召回它所偏离的默想，并用最坚定的心意（如同用锁链）将其束缚时，在我们尝试之际，它却比蛇更快地从心底深处溜走。因此，我们虽被这种日常修行所激励，却未见从中获得任何内心的力量和稳定，便被打败，绝望中得出这一结论：即相信这些心神游荡出现在人身上，并非出于我们自己的过错，而是出于本性的过错。\par

% structure: p0013 source=h2 target=subsection reason=relative-heading-rank; base=h2

\subsection{第四章}

\Emph{长者关于灵魂状态及其卓越的论述。}\par

\Person{塞雷努斯}：轻率下结论，尚未恰当探讨主题、考虑其真实性质之前便对某事的性质急于定论，是危险的。你也不应只看见自己的软弱，凭猜想冒险，而应根据实践本身的性质和价值以及他人的经验作出判断。倘若有人不识水性，却知水的浮力无法支撑自身重量，便想凭其无经验提供的证据断言，绝无可能有人能以血肉之躯漂浮于液体之上；我们不应因此便以为他的观点正确——他似乎根据自己的经验提出，但这不仅被他人证明绝非不可能，而且实际上极易做到，有最清晰的证据和目击为证。因此，\OriginalTerm{心神}{\GreekText{νοῦς}}，即理智，被定义为\OriginalTerm{永远变动且多变的}{\GreekText{ἀεικίνητος} \GreekText{καὶ} \GreekText{πολυκίνητος}}，正如在\WorkTitle{智慧篇}中以另一方式所描述的那样：\OriginalTerm{属地的帐幕压下了多虑的心神}{\GreekText{καὶ} \GreekText{γεῶδες} \GreekText{σκῆνος} \GreekText{βρίθει} \GreekText{νοῦν} \GreekText{πολυφρόντιδα}}，即\BibleQuote{属地的帐幕压下了多虑的心神。}\BibleRef{智 9:15}因此，心神按其本性永不能静止，除非提供地方使其运动，并常有使其忙碌之事，否则它必因自身善变而游荡，流连于各类事物，直至它藉长期实践和日常运用——你说你在这其中徒劳挣扎——尝试并学会当预备何种记忆食粮，使它能向之回归其不懈的飞行，并获得持留的力量，从而得以驱散敌手那使之分心的恶意暗示，并坚持于它所渴望的状态与境况。因此，我们不应将心的这种游荡倾向归咎于人性或造物主天主。因为圣经有真确之言：\Quote{天主造人原为正直；但他们自己寻出了许多思念。}这些思念的性质便在于我们自己，因为经上说：\Quote{善念亲近认识它的人，但明智的人必会找到它。}因为凡属我们审慎和勤勉所能发现之处，若\Emph{未}能发现，我们定当归咎于己的懒惰或疏忽，而非本性的缺陷。圣咏作者也同意此意，他说：\Quote{那以祢为助佑的人是有福的：他在心中安排了他的上升。}你看，那么，我们有权在心中安排\Emph{上升}，即属于天主的思念，或\Emph{下降}，即沉于血肉和属世之事的思念。如果这不在于我们权下，主就不会责斥法利塞人说：\BibleQuote{你们为什么心中思念恶事？}\BibleRef{玛 9:4}也不会藉先知吩咐：\Quote{从我眼前除去你们思念的邪恶；}以及\Quote{邪恶的思念在你们内要到几时？}在审判之日，我们思念的性质如同我们的行为一样会受审，正如主藉依撒意亚恐吓说：\BibleQuote{看，我来聚集他们的行为和思念，连同万国万民；}\BibleRef{依 66:18}在那可怕可畏的审查中，我们也不该因它们的见证而被定罪或辩护，正如蒙福的宗徒所说：\BibleQuote{他们的思念彼此互相控告或辩护：在天主按我的福音审判人心的隐秘之日。}\BibleRef{罗 2:15}\par

% structure: p0016 source=h2 target=subsection reason=relative-heading-rank; base=h2

\subsection{第五章}

\Emph{论灵魂的完美，由福音中百夫长的比较引申而来。}\par

论这完美的精神，在福音中百夫长的事迹里有一个极好的比喻画像。他的德行和坚定，使他不会被思想的洪流冲走，而是根据自己的判断，要么接纳善的思想，要么轻易驱走相反性质的思想，这些都以比喻的形式描述如下：\BibleQuote{因为我也属人权下，也有兵士在我以下：我对这个说：你去，他就去；对另一个说：你来，他就来；对我的奴仆说：你做这个，他就做。}\BibleRef{玛 8:9} 如果我们也能勇敢地与纷扰和罪恶搏斗，并将它们置于自己的控制和判断之下，与肉体中的情欲争战并消灭它们，使理智统驭思想的蜂群，并借着主十字架赋予生命的旗帜，从胸中击退敌对我们的可怕权势的军队，那么作为这些胜利的赏报，我们就会被提升到那位在属神意义上理解的百夫长的地位；正如我们在出谷纪中读到，梅瑟曾神秘地指向他：\BibleQuote{你要为自己指派千夫长、百夫长、五十夫长、十夫长。}\BibleRef{出 8:21} 同样，当我们被提升到这种尊荣的高度时，也将拥有同样的命令权与能力，使我们不会被思想违背意愿地冲走，而是能够持守并紧贴那些在属神上令我们喜悦的思想，命令邪恶的暗示离开，它们就离开；而对善的暗示，我们则说\Quote{来}，它们就来；对我们的仆人，即身体，我们也同样命令与贞洁和节制相关的事，它就会毫无反对地服侍我们，不再在我们内激起私欲的敌对冲动，反而完全顺服于神魂。这位百夫长的武器有何特性，在战斗中有何用途？请听蒙福的宗徒宣告：他说：\BibleQuote{我们的作战武器不是血肉的，而是在天主前有力的。}他告诉我们它们的特性：即不是血肉的或软弱的，而是属神的，在天主前有力的。接着他指示它们将在哪些战斗中被使用：\BibleQuote{用以攻破堡垒，清除思念，以及每一种高抬自己反对认识天主的骄傲，并俘虏一切理性，使其服从基督，并且准备好惩罚一切悖逆，当你们的服从首先被成全时。}\BibleRef{格前 10:4} 虽然详述这些武器逐一有益，但应当留待其他时机；我只愿你们看到这些武器的不同种类及其特性，因为如果我们要打主的战争，并在福音的百夫长中服务，我们也当始终佩戴它们行走。他说：\BibleQuote{拿起信德的盾牌，使你们能扑灭恶者一切的火箭。}\BibleRef{弗 6:16} 信德就是如此，它拦截情欲的火焰之箭，并以对未来审判的恐惧和对天国的信仰摧毁它们。他说：\BibleQuote{以及爱德的胸甲。}\BibleRef{得前 5:8} 这胸甲环绕胸口的要害部位，保护易受骄傲思想致命创伤之处，抵挡攻击的打击，不容魔鬼的箭穿透到我们的内里。因为它\BibleQuote{凡事忍受，凡事受苦，凡事担当。}\BibleRef{格前 13:7} \BibleQuote{以及救恩之望的头盔。}\BibleRef{得前 5:8} 头盔是保护头部的。既然基督是我们的头，我们就应在一切诱惑和迫害中，用对未来美善事物的希望来保护它，并尤其要持守对祂的信德纯全无玷。因为失去身体其他部分的人，尽管软弱，仍可能存活；但没有头，无人能多活片刻。\BibleQuote{以及圣神的剑，即是天主的话。}\BibleRef{弗 6:17} 因为它\BibleQuote{比任何双刃剑更锐利，甚至穿透灵魂与神魂的分开，以及关节与骨髓的分开，且能辨别心中的思念和意向。}\BibleRef{希 4:12}它分割并切除在我们内发现的一切血肉和属世之物。凡受这些武器保护的人，必常得保卫，免于敌人的武器和蹂躏，也不会被掠夺者的锁链捆绑，像俘虏和囚犯一般，被带到虚妄思念的敌对之地，也不会听到先知的话：\BibleQuote{你为什么在异域变老了？}\BibleRef{巴 3:11} 但他会像一个得胜的征服者，站在他所选择的思念之地。你也愿意了解这位百夫长的力量和勇气，他正是靠着这些武器——我们之前所说不是血肉的，而是在天主前有力的——来背负它们吗？请听君王在召唤勇士参加属神战斗时，如何挑选并认可他们。祂说：\Quote{让软弱者说我是刚强的；}以及：\Quote{让受苦者成为战士。} 那你们看，只有受苦者和软弱的人才能打主的战争，且正是那种软弱，我们的福音百夫长曾凭着它自信地说：\BibleQuote{因为我软弱时，正是我刚强的时候，并且因为能力在软弱中得以成全。}\BibleRef{格后 12:9} 论这种软弱，有一位先知说：\BibleQuote{他们中软弱的，必如达味家。}\BibleRef{匝 12:8} 因为忍耐的受苦者必打这些仗，带着那忍耐，正如经上说：\BibleQuote{你们需要忍耐，以便承行天主的旨意，而获得赏报。}\BibleRef{希 10:36}\par

% structure: p0019 source=h2 target=subsection reason=relative-heading-rank; base=h2

\subsection{第六章}

\Emph{论关于思念之照顾的恒心。}\par

但我们将会从自己的经验中发现，如果我们克制自己的意愿，断绝今世的欲望，我们就能且应该紧紧依附于主；并且我们将从那些与主交谈并满怀信心说\Quote{我的灵魂紧紧依附于祢；}和\Quote{主啊，我依附于祢的训言；}以及\Quote{依附天主为我是好的；}和\Quote{那与主结合的，便成为一灵。}之人的权威中得到教导。因此，我们不应因心思的游荡而疲倦，放松我们的热忱；因为\Quote{那耕种自己田地的，必得饱食；那追求闲逸的，必充满贫穷。}\BibleRef{箴 28:19}我们也不应因危险的绝望而偏离对这警醒的专注，因为\Quote{凡忧虑的人，必有余；而那安逸无忧的，必有所缺；}又：\Quote{一个忧伤的人为自己劳作，并强行使自己毁灭。}此外：\Quote{天国遭受暴力，暴力的人以夺取它，}\BibleRef{玛 11:12}因为没有辛劳就不能获得任何德行，没有极大的内心忧伤就不能达到所期望的心灵稳定，因为\Quote{人生来就为受苦，}\BibleRef{约 5:7}并且为了能够达到\Quote{圆满的人，达到基督圆满年龄的程度}\BibleRef{弗 4:13}，他必须更加专注地警醒，以不懈的谨慎努力。但没有人能达至这圆满的程度，除非他预先思考并接受训练，如今还在世上便预先尝到一些滋味，并被标为基督最宝贵的肢体，在肉身中拥有那\Quote{联络}\BibleRef{弗 4:13}的凭据，借此得以与祂的身体结合：只渴望一件事，只渴求一件事，始终不仅将他的行为，甚至连他的思想都集中于一件事；即，他甚至现在就能保有那对将来圣者幸福生活的预言作为凭据，即：\Quote{天主成为}他的\Quote{万有中的万有。}\BibleRef{格前 15:28}\par

% structure: p0022 source=h2 target=subsection reason=relative-heading-rank; base=h2

\subsection{第七章}

\Emph{关于心思游荡的倾向与邪魔攻击的提问。}\par

\Person{革尔玛努斯}：也许，这种心思游荡的倾向在某种程度上可以被抑制，如果不是因为如此大群的敌人包围着它，并不断驱使它走向它不愿去的地方，或者更确切地说，走向其本性游荡的倾向所驱使的地方。既然如此无数的敌人，而且如此强大可怕，包围着它，我们就不应认为，特别是凭我们这脆弱的肉身，能够抵挡它们，若不是被你的话如同来自天上的神谕所鼓励，我们不会这样想。\par

% structure: p0025 source=h2 target=subsection reason=relative-heading-rank; base=h2

\subsection{第八章}

\Emph{关于天主的帮助与自由意志能力的回答。}\par

\Person{塞雷努斯}：凡经历过内在人争战的人，都不会怀疑我们的敌人不断埋伏等待我们。但我们认为，他们对我们的进步进行阻碍的方式，使我们只能认为他们只是在\Emph{煽动}恶事，而不是\Emph{强迫}。但是，如果他们有一种强烈的冲动，像建议那样强迫（恶事）进入我们心中，那么没有人能够完全避免他们倾向于印在我们心里的任何罪。因此，正如他们里面有充分的煽动能力，我们里面也有拒绝的能力和同意的自由。但是，如果我们害怕他们的能力和攻击，我们也可以祈求天主的保护和帮助来对抗他们，正如我们读到：\Quote{因为那在我们内的，比那在世界上的更大；}\BibleRef{若一 4:4}而祂的帮助以更大的能力为我们争战，胜过他们的大军攻击我们；因为天主不仅是善事的倡导者，也是善事的维持者和坚持者，以至于有时祂甚至在我们不愿意且不知觉的情况下将我们引向救恩。由此可见，没有人会被魔鬼欺骗，除非他选择同意自己的意志；正如训道篇用这些话清楚地说明：\Quote{因为恶人迅速行恶而无人反对，因此世人的心充满了行恶的意念。}因此很清楚，每个人从这一点上犯错；即，当恶念攻击他时，他没有立即以拒绝和反对来应对，因为经上说：\Quote{抵挡他，他必从你们面前逃去。}\BibleRef{雅 4:7}\par

% structure: p0028 source=h2 target=subsection reason=relative-heading-rank; base=h2

\subsection{第九章}

\Emph{关于灵魂与魔鬼结合的问题。}\par

\Person{革尔玛努斯}：请问，那种灵魂与邪魔之间无区别且普通的结合是什么，以至于它们能够（我不说结合，而是）与它联合，以致它们能够不知不觉地与它谈话，进入它内部，向它建议它们想要的一切，煽动它去做它们喜欢的事，查看并看到它的思想和活动；其结果使它们与灵魂之间如此紧密，以至于没有天主的恩宠几乎无法分辨哪些源于它们的煽动，哪些出于我们的自由意志？\par

% structure: p0031 source=h2 target=subsection reason=relative-heading-rank; base=h2

\subsection{第十章}

\Emph{关于不洁之灵如何与人的灵魂结合的回答。}\par

\Person{塞雷努斯}：灵体能够不知不觉地与灵体结合，并对其施加看不见的劝说力量，这并不奇怪。因为它们之间（如同人与人之间）有某种相似性和本质上的亲缘关系，因为对灵魂本性的描述也同样适用于它们的实体。但是，灵体不可能内在地植入灵体，或以这样一种方式结合，以致一个能抓住另一个；因为只有天主才有这种特权，祂是唯一单纯和非形体的本性。\par

% structure: p0034 source=h2 target=subsection reason=relative-heading-rank; base=h2

\subsection{第十一章}

\Emph{一个异议：不洁之灵是否能临于或结合于他们所充满之人的灵魂中。}\par

\Person{Germanus}：我们觉得，在附魔者身上所见的现象足以反驳此观点，即当他们被不洁邪灵影响时，所说所作皆非己所知。那么，我们怎能拒绝相信他们的灵魂与那些不洁邪灵结合呢？因为我们看见他们成为不洁邪灵的工具，（放弃天然状态）顺从不洁邪灵的动作与情绪，以致不再表达自己的言语、行动和意愿，而表达不洁邪灵的言语、行动和意愿？\par

% structure: p0037 source=h2 target=subsection reason=relative-heading-rank; base=h2

\subsection{第十二章}

\Emph{答案：不洁邪灵如何能主宰附魔者。}\par

\Person{Serenus}：你所说附魔者的情况并不反对我们的主张，即被不洁邪灵附身的人所说所作并非出于自愿，被迫说出不知之事。因为明显可见，他们并非以同一方式受不洁邪灵侵入：有些受其影响，对自己所言所行毫无意识；另一些则能知晓并事后回忆。但不可想象，不洁邪灵的注入能渗透到灵魂的实体中，与之结合，仿佛穿上灵魂，通过患者的口说出话语。因为不应相信此事可能。我们清楚看出，这并非由于灵魂的缺失，而是由于身体的软弱：当不洁邪灵抓住灵魂力量所在的肢体，施以极大而不可忍受的重压，以极深的黑暗淹没灵魂，干扰其理智能力。正如有时因饮酒、发烧、过度寒冷或其他外在失调而见的情况。这正是魔鬼被禁止加于约伯身上的事：虽然魔鬼已获得对其肉身的权柄，但主命他说：\Quote{看，我将他交在你手中；只要保全他的灵魂}，即，不要削弱他灵魂的座所而使他疯狂，不要用你的重量压制他心中的统治权，从而压垮他残余的理解与智慧。\par

% structure: p0040 source=h2 target=subsection reason=relative-heading-rank; base=h2

\subsection{第十三章}

\Emph{不洁邪灵如何不能渗透不洁邪灵，以及唯天主是非物质的。}\par

因为即使不洁邪灵能与粗糙而坚实的物质即肉体混合（这极易发生），难道我们应相信它能与同样是灵体的灵魂结合，使灵魂同样接纳其本性？这是唯有天主圣三才能做到的：祂能渗透一切理智本性，不仅拥抱环绕，而且插入其中，无形体却能注入形体？因为我们虽承认有些灵体存在，如天使、总领天使及其他权能，以及我们自己的灵魂和稀薄空气，但绝不应视之为非物质的。它们各有自己的形体存在，虽远较我们的身体精细，如宗徒所说：\BibleQuote{有属天的形体，也有属地的形体}；又说：\BibleQuote{所种的是属生灵的身体，复活的是属神的身体}\BibleRef{格前 15:40}；由此清楚可知，唯有天主才是非物质的。因此，只有祂能渗透一切灵性和理智实体，因为祂独一、全在、在万有之中，得以鉴察人的思想、内在活动及灵魂的一切隐秘。对此，蒙福的宗徒说：\BibleQuote{天主的话是生活的，有效的，比双刃的剑更锐利，直穿入灵魂和神魂、关节与骨髓之间，且能辨明心中的思想和意念。没有受造物在祂面前不是显明的，万物在祂眼前都是袒露敞开的}\BibleRef{希 4:12}。蒙福的达味也说：\BibleQuote{祂一一塑造了他们的心}；又说：\BibleQuote{因为祂知道人心的隐秘}；约伯也说：\BibleQuote{唯有祢知道人的心}\BibleRef{编下 6:30}。\par

% structure: p0043 source=h2 target=subsection reason=relative-heading-rank; base=h2

\subsection{第十四章}

\Emph{质疑：应如何相信魔鬼能洞察人的思想。}\par

\Person{Germanus}：照你所描述的方式，那些不洁邪灵绝不可能看出我们的思想。但我们认为这观点极其荒谬，因为圣经说：\BibleQuote{若有掌权者的灵上升到你身上}\BibleRef{训 10:4}；又说：\BibleQuote{魔鬼已将出卖主的意思放入西满依斯加略的心}\BibleRef{若 13:2}。我们怎能相信我们的思想不向他们显露，当我们感觉多数情况下，思想是因他们的建议和煽动而生发并滋长？\par

% structure: p0046 source=h2 target=subsection reason=relative-heading-rank; base=h2

\subsection{第十五章}

\Emph{答案：关于人的思想，魔鬼能做什么和不能做什么。}\par

\Person{色勒诺}：无人怀疑邪魔能影响我们思念的特质，但这是通过从外部施加可感知的影响，即要么来自我们的倾向，要么来自我们的言语以及那些它们看到我们特别偏爱的喜好。但它们不可能接近那些尚未从灵魂最深处涌出的事物。此外，它们所建议的思念，无论是否被实际或以某种方式接纳，它们并非从灵魂本身的本质——即那潜藏在骨髓中的内在倾向——来发现，而是从外在之人的动作和迹象来发现。例如，当它们怂恿贪食时，若它们看到一位隐修士焦急地抬眼望向窗户或太阳，或急切地询问几点钟，它们就知道他已接纳了贪婪之情。当它们怂恿邪淫时，若它们发现他冷静地屈从于情欲的攻击，或看到他的身体骚动不安，或至少没有像理应的那样在污秽建议的放纵下叹息，它们就知道情欲之箭已刺入他的灵魂。若它们挑动悲伤、愤怒或暴怒的刺激，它们便能通过身体的动静和可见的骚乱来判断这些是否已在心中扎根，例如，当它们注意到他要么默默叹息，要么因愤慨而喘息，或脸色变化；它们便狡猾地发现他所陷的过错。因为它们知道，我们每个人都会被那一种诱惑以惯常方式所诱，而它们看到他以一种身体同意的动作，已向那刺激献出了他的同意和认可。这被那些空中权势发现并不足为奇，因为甚至聪明人往往也能从一个人的神情、面容和外在举止发现其内心状态。那么，那些身为属灵体、无疑比人精细聪明得多的存在，岂不更能肯定地发现这一点吗？\par

% structure: p0049 source=h2 target=subsection reason=relative-heading-rank; base=h2

\subsection{第十六章}

\Emph{一个说明，显示我们如何被教导：邪魔知道人的思念。}\par

正如有些盗贼习惯于探查他们打算抢劫的那些房屋里隐藏的财宝，在漆黑的夜色中小心翼翼地撒下细沙，借着沙子落下时发出的叮当声来发现他们看不见的隐藏宝藏，从而确知每件东西和金属，它们仿佛通过所发出的声音暴露了自己；同样，这些邪魔为了探查我们心中的宝藏，也向我们撒下某些邪恶建议的沙子，当它们看到某种与这些建议性质相应的身体情感出现时，便仿佛通过发自内心深处的一种叮当声，认出内室中储存的是什么。\par

% structure: p0052 source=h2 target=subsection reason=relative-heading-rank; base=h2

\subsection{第十七章}

\Emph{论并非每个魔鬼都有能力向人注入每一种情欲。}\par

但我们应当知道，并非所有魔鬼都能向人注入所有情欲，而是某些魔鬼专司某一罪恶，有些喜爱不洁与污秽的淫欲，有些喜爱亵渎，有些特别投身于愤怒与暴怒，有些在忧郁中生长，有些以虚荣与骄傲得慰藉；每一个魔鬼都将自己所沉溺的罪恶注入人心，它们不能同时注入它们特有的恶行，而是轮流地，按照时间、地点或一个易于接受它们建议的人所引发的机会来施加。\par

% structure: p0055 source=h2 target=subsection reason=relative-heading-rank; base=h2

\subsection{第十八章}

\Emph{一个问题：在魔鬼中间，攻击时是否遵循某种秩序，或其变化中有某种体系？}\par

\Person{革尔玛奴斯}：那么，我们是否必须相信，邪恶在他们中间被如此安排，可以说被系统化，以至于他们的变化中有某种秩序，执行一种有规律的攻击计划？尽管显然方法和体系只能存在于良善正直的人中间，正如圣经上说：\BibleQuote{你们在恶人中寻求智慧，必找不到；}并且：\BibleQuote{我们的仇敌是无知的；}以及：\BibleQuote{在恶人中没有智慧，没有勇气，也没有谋略。}\par

% structure: p0058 source=h2 target=subsection reason=relative-heading-rank; base=h2

\subsection{第十九章}

\Emph{回答：关于攻击及其变化，魔鬼之间存有何种一致？}\par

\Person{色勒诺}：确实有一种真实的说法，即恶人中没有持久的和睦，甚至对那些对他们都有吸引力的特殊过犯，也不可能有完美的和谐。因为，正如你所说，在无纪律的事物中永远不可能保持体系和纪律。但在某些事情上，当利益共同体、必要性强加之，或分享某种收益所建议时，他们必须安排暂时的协议。我们非常清楚地看到，在这场属灵邪恶的战争中也是如此：因此，它们不仅在彼此之间遵守时间和变化，而且实际上还特别占据某些特定地点并持续出没：因为既然它们必须通过某些固定的诱惑和明确的罪恶，并在特定的时间进行攻击，我们从中明确推断，没有一个人能在同一时刻被虚荣的空虚所迷惑，又被邪淫的欲望所点燃；也没有人能同时被属灵骄傲的狂妄所膨胀，又受制于肉体贪食的屈辱。也没有人能同时被愚蠢的咯咯笑和笑声所征服，同时又因愤怒的刺痛而兴奋，或至少被啃噬悲伤的痛苦所充满：相反，所有魔鬼必须一个接一个地前进，攻击灵魂，其方式是一个被打败退却后，必须让位于另一个魔鬼更猛烈地攻击它；或者如果一个魔鬼获胜利，它仍会将灵魂交给另一个魔鬼去欺骗。\par

% structure: p0061 source=h2 target=subsection reason=relative-heading-rank; base=h2

\subsection{第二十章}

\Emph{论对立的权势并非具有同等的胆量，而且诱惑的时机并不受他们控制。}\par

我们也不该不知道这一点：并非所有魔鬼都有同样的凶暴和能量，也没有同样的胆量和恶意。对于初学者和软弱之人，只有较弱的属灵邪恶势力与之交战；当这些属灵的邪恶势力被打败后，攻击基督运动员的乃是逐渐增强的更强者的冲击。因为斗争的难度是随着人的力量和进步而加大的：因为没有一个圣人能独自承受如此众多且强大的敌人的恶意，或抵挡他们的攻击，甚至忍受他们的残酷和野蛮，若非我们竞赛的仁慈审判者、赛场的主持人——基督本人——均衡了交战者的力量，击退并遏制了它们过分的攻击，且在诱惑中为我们开了出路，使我们能承受得住。\BibleRef{格前 10:13}\par

% structure: p0064 source=h2 target=subsection reason=relative-heading-rank; base=h2

\subsection{第21章}

\Emph{论魔鬼与人的斗争并非没有它们的努力。}\par

但我们相信，他们所进行的这场战斗并非毫无自己的努力。因为在他们自己的冲突中，他们也有某种焦虑和沮丧，特别是当他们与更强的对手，即圣人和完美的人相比时。否则，就不会有竞赛或战斗，而只会将一种简单的、对他们而言毫无焦虑的欺骗归于他们。那么宗徒的话如何成立，他说：\Quote{我们作战，不是对抗血和肉，而是对抗率领者，对抗掌权者，对抗这黑暗世界的霸主，对抗天界里邪恶的鬼神；} 以及：\Quote{我战斗，不像打空气的；} 并且：\Quote{我打了美好的仗}？因为凡是说到战斗、冲突和战争的地方，双方必然都有努力、辛劳和焦虑，同样，他们要么因失败而面临懊恼和混乱，要么因胜利而享受喜悦。但是，如果一方轻松而有把握地对抗另一奋力挣扎的人，并且仅凭自己的意愿作为力量来压倒对手，那就不应称为战斗、搏斗或争执，而是一种不公平、不合理的袭击和攻击。但是，他们确实必须劳苦，当他们攻击人时，也付出同样大的努力，以便从每个人身上取得他们想要获得的胜利，并且，如果我们被他们击败，那原本等待我们的混乱就会返回到他们身上；正如经上所说：\Quote{他们围困我的头，他们嘴唇的劳苦必压倒他们；} 以及：\Quote{他的愁苦必归到他自己的头上；} 并且：\Quote{愿那他不知道的网罗临到他，愿他隐藏的网捉住他，愿他掉进那网罗里；} 即，他为欺骗人所设的计谋。他们自己也遭受灾祸，当他们在伤害我们时，也同样被我们伤害，当他们战败时，也并非毫无混乱地离开。那些内在之眼明亮的人，看见他们的失败和挣扎，又看见他们因个人的跌倒和不幸而幸灾乐祸，害怕自己的情况也给他们带来这种快乐，便向主祈祷说：\Quote{求你明亮我的眼睛，免得我沉睡至死，免得我的仇敌说：我战胜了他。倘若我动摇，那烦扰我的人就要欢喜。} 以及：\Quote{我的天主，求你不要让他们因我而欢喜；不要让他们心里说：啊哈，啊哈，我们的愿望！也不要说：我们吞灭了他。} 以及：\Quote{他们向我咬牙切齿。上主，你旁观要到何时？} 因为：\Quote{他像狮子在洞穴中埋伏，他埋伏要抢夺穷人；} 以及：\Quote{他从天主那里求食。} 而且，当他们用尽一切努力，仍未能诱骗我们时，他们必然因努力失败而\Quote{感到困惑和脸红}，\Quote{那些寻索我们灵魂要毁灭他们的人，愿他们披上羞耻和混乱，因为他们图谋恶事攻击我们。} 耶肋米亚也说：\Quote{愿他们感到困惑，而不要使我感到困惑；愿他们恐惧，而不要使我恐惧；求你使你的愤怒降在他们身上，以双重的毁灭毁灭他们。} \BibleRef{耶 17:18} 因为无人怀疑，当他们被我们战胜时，将遭受双重毁灭：首先，因为当人在追求圣德时，他们虽曾拥有却失去了圣德，并成为人毁灭的原因；其次，因为他们作为属灵的存在，竟被属肉体和属地的存在所战胜。于是，每一位圣人看到仇敌的毁灭和自己的凯旋，便喜悦地高呼：\Quote{我要追赶我的仇敌，追上他们；不消灭他们，决不回转。我要击打他们，使他们不能起来：他们必倒在我脚下，} 在同一位先知对抗他们的祈祷中说到：\Quote{上主，求你判断那错待我的人，推翻那攻击我的人。拿起武器和盾牌，起来帮助我。抽出剑，拦住那追赶我的人的道路；对我的心说：我是你的救恩。} 当我们通过征服和消灭我们所有的情欲而战胜了这些仇敌时，我们将得以听到这些祝福的话：\Quote{你的手要举起在仇敌之上，你所有的仇敌都要灭亡。} \BibleRef{米 5:9} 因此，当我们阅读或咏唱圣经中所有这一切和类似的段落时，除非我们将它们理解为是针对那些日夜埋伏我们的邪恶鬼神而写的，否则我们不仅无法从中得到任何使我们温柔和忍耐的教诲，反而会遭遇某种可怕的后果，与福音的成全完全相反。因为我们将不仅学不会为仇敌祈祷或爱仇敌，反而会被激起对他们怀有刻骨的仇恨，诅咒他们，并不断向他们倾泻祈祷。认为圣人们，这些天主的朋友，在基督来临之前就已经不受法律的约束，因为他们超越了法律的命令，选择遵守福音的训诫并追求成全，尽管他们生活在时间的规定之前，却以这样的精神说出这些话，是极其错误和亵渎的。\par

% structure: p0067 source=h2 target=subsection reason=relative-heading-rank; base=h2

\subsection{第二十二章}

\Emph{论伤害的能力并不取决于魔鬼的意志}\par

但是，那些恶者并无伤害任何人的能力，这由真福约伯的例子非常清楚地显明出来：敌人不敢超过天主许可的范围试探他。这也由福音记载中同一恶灵的告白所证实，他们说：\BibleQuote{你若驱逐我们，就准我们进入猪群。}\BibleRef{玛 8:31} 何况我们更应坚信，他们若没有天主的许可，连进入无言的、不洁的动物都不能，就更不能随意进入任何按天主肖像受造的人。但是，任何人——我不说那些我们看见在此旷野中极其坚定地生活的年轻人，就连那些成全的人——若他们拥有无限的权力和自由来伤害并试探我们，就不可能独自生活在旷野，被如此众多的这类仇敌包围。这一点更清楚地由我们的主和救主的话所支持：祂在祂所取的人性的卑微中，对比拉多说：\BibleQuote{你对我没有任何权柄，除非是从上赐给你的。}\BibleRef{若 19:11}\par

% structure: p0070 source=h2 target=subsection reason=relative-heading-rank; base=h2

\subsection{第二十三章}

\Emph{论魔鬼权能的减弱。}\par

但是，我们凭借自己的经验以及长老们的见证，已经充分发现：如今魔鬼不再拥有像隐修士初期——那时还只有少数隐修士住在旷野——那样的权能。因为起初他们的凶猛如此强烈，以至于只有少数极其坚定且年长的人才能忍受独修的生活。在那些八、十人一起居住的隐修院里，他们的攻击如此猛烈，他们的侵袭如此频繁且显著，以至于众人不敢全部在夜间同时就寝，而是轮流：一些人稍睡片刻，另一些人则保持警醒，专务圣咏、祈祷和诵读。当自然的需要迫使他们睡眠时，他们唤醒其他人，同样把看守那些正要就寝者的职责交给他们。由此我们无法怀疑：有两种可能性之一造成了这一结果，不仅对于我们这些由于年龄给予经验而看似相当强壮的人，而且对于更年轻的人也是如此。因为，要么是魔鬼的恶意被十字架的力量所击退——这十字架的力量甚至渗透到旷野，并以其恩宠照耀各处；要么是我们的疏忽使他们放松了最初的猛攻，因为他们不屑于以他们从前攻击那些最可赞美的基督士兵同等的精力来攻击我们；并且通过这种欺骗和停止公开攻击，他们对我们造成了更大的损害。因为我们看到有些人已经陷入如此懒散的状态，以至于他们必须被极其温和的劝诫所哄骗，惟恐他们离开自己的斗室，陷入更危险的麻烦，到处游荡迷失，并被缠住在我所谓更粗俗的罪中；人们认为，如果他们即使有些懈怠地留在旷野，就已从他们那里得到了很大的收获，长老们常常对他们说，作为一种巨大的解脱：\Quote{待在你们的斗室里，尽量吃喝睡眠，只要你们总是留在其中。}\par

% structure: p0073 source=h2 target=subsection reason=relative-heading-rank; base=h2

\subsection{第二十四章}

\Emph{论魔鬼为自己预备进入他们所要附身之人身体的途径。}\par

显然，不洁之神除了先占据他们的心神和思想之外，不能以其他方式进入那些他们将要攫取身体的人。当他们夺去了这些人的敬畏、对天主的记忆和属神的默想之后，他们便大胆地逼近他们，好像他们被剥夺了一切保护和神圣的守卫，可以轻易被束缚，然后便在他们内安家，如同在一处交给他们的财产中。\par

% structure: p0076 source=h2 target=subsection reason=relative-heading-rank; base=h2

\subsection{第二十五章}

\Emph{论那些被罪恶附身的人比那些被魔鬼附身的人更为不幸。}\par

虽然事实是，那些在身体上似乎很少受影响、但在精神上以更恶劣的方式被占有的人——即被他们的罪恶和情欲所缠住——遭受着更严重更剧烈的困扰。因为正如宗徒所说：\Quote{人被谁制伏，就是谁的奴隶。}不过，在这方面他们处于更危险的病态，因为尽管他们是奴隶，却不知道被它们所攻击，并处于其统治之下。但我们知道，甚至圣善的人也因一些极其轻微的过失，被交付于肉身给撒殚和重大的苦难，因为神圣的仁慈不愿在审判之日他们身上有丝毫斑点或污秽，并在此世炼净他们的一切污秽，正如先知——或者更确切地说，天主自己——所说，为要把他们如同炼净的金银、无需惩罚性的净化一样交予永恒。祂说：\Quote{我必清洁炼尽你的渣滓，除去你一切的罪；此后你必被称为正义之城，忠信之城。}又说：\Quote{如同银子和金子在炉中锻炼，主也选择人心。}又说：\Quote{火试炼金银；但人在卑微的炉中被试炼。}还有：\Quote{主所爱的，祂必管教，又鞭打凡所收纳的儿子。}\par

% structure: p0079 source=h2 target=subsection reason=relative-heading-rank; base=h2

\subsection{第二十六章}

\Emph{论那被误导的先知的死亡，以及院长保禄的软弱——他为炼净之故所受的苦难。}\par

我们在\BibleBook{列王}{纪上}中看到那位先知和天主的人的清晰实例，他因一次不顺从的过失立刻被狮子消灭，他并非有意或出于己意之过，而是受他人引诱而卷入，就如圣经论及他所说：\BibleQuote{这是天主的人，他违背了主的口谕，主就把他交给狮子，狮子撕裂了他，正如主所说的话。}\BibleRef{列上 13:26} 在此事例中，对眼前过犯和疏忽的惩罚，连同他正义的赏报——主为此在世上将自己的先知交付给毁灭者——都由那猛兽的节制和禁食显示出来，因为那最凶残的活物甚至不敢尝那交付给它的尸体。关于同样的事，在我们这个时代，我们亲眼看到极其清楚明白的证据，发生在住在这一旷野中一处名为\Place{卡拉穆斯}的\Person{保禄}院长和\Person{梅瑟}院长身上。前者曾居住在靠近\Place{帕内菲西斯}城的荒野中，我们知道那里最近因盐水泛滥而成为荒野；每当北风吹起，水从沼泽被驱赶出来，蔓延到邻近田野，覆盖整个地区的表面，使得古老的村庄——正因如此，所有居民已弃之而去——看起来像岛屿。在那里，\Person{保禄}院长在旷野的寂静和沉默中，在心灵的纯洁上取得了如此大的进步，以致他不容忍——我不说女人的面容——甚至任何一个女性性别的衣服出现在他眼前。因为当他与住在同一旷野的\Person{阿尔赫比乌}院长一起前往一位长老的静舍时，偶然有一个女人遇见了他，他对此相遇如此厌恶，以致他放弃了已着手进行的那次友好访问的事务，并以比人逃避狮子或可怕巨龙之面更快的速度冲回自己的隐修院；甚至没有因前述\Person{阿尔赫比乌}院长的喊叫和祈求而动摇——后者呼唤他继续他们已开始的旅程，去请教那位长老他们原本打算做的事。但是，尽管这是出于他对贞洁的热忱和对纯洁的渴望，然而因为这不是按照知识做的，而且因为纪律的遵守和适当严厉的方法被过度强调——因为他想象不仅与女人亲近（这是真正的危害），而且甚至那个性别的形像本身也应被诅咒——他立刻受到如此惩罚，以致全身瘫痪，没有一根肢体能够执行其正常功能，因为不仅手和脚，而且使我们形成言语的舌头动作以及他的耳朵本身也失去了听觉，以致在他身上除了一个不动不感的人形外，什么也没有留下。他被降到这样的状态，以致男人最大的护理也无法帮助他的虚弱，只有女人的温柔服事才能照顾他的需要：因为当他被带到一座圣洁童贞女的隐修院时，由女性看护者给他送上食物和饮料——他连用手势都不能要求；并且为了执行本性所需的一切，他由同样的服事照顾了将近四年，即直到他生命的终结。尽管他因所有肢体的软弱而受苦，以致没有一根肢保留敏锐的运动和感觉能力，然而从他身上发出如此之大的恩宠，以致当病人用沾过那应称为尸首而非身体的油涂抹时，他们立刻痊愈一切疾病，以致关于他自己的疾病，甚至对不信者也清楚明白地表明，他所有肢体的软弱是因主的眷顾和慈爱所致，而这些治病的恩宠是由圣神的能力所赐，作为他纯洁的见证和功绩的显示。\par

% structure: p0082 source=h2 target=subsection reason=relative-heading-rank; base=h2

\subsection{第二十七章}

\Emph{论院长梅瑟的诱惑}\par

但是，我们提到的住在这一旷野中的第二个人，虽然他也是杰出而引人注目的人，然而，为惩罚他在与\Person{玛加略}院长的一次辩论中，因在某意见上被抢先而说得过于尖锐的一句话，他立刻被交付给一个如此可怕的魔鬼，以致魔鬼往他嘴里塞满污秽之物，而主通过他迅速痊愈和治愈的来源显示，祂将这灾祸降在他身上是为洁净他，使他不留下任何来自瞬间过失的污点。因为\Person{玛加略}院长一投入祈祷，话还不及出口，那恶灵就从他那里逃走了。\par

% structure: p0085 source=h2 target=subsection reason=relative-heading-rank; base=h2

\subsection{第二十八章}

\Emph{我们不应轻视那些被交付给不洁之神的人}\par

由此明显地得出结论：我们不应仇恨或轻视那些我们看见被交付给各种诱惑或不洁之神的人，因为我们应当坚定地持守两点：第一，没有天主的允许，他们中任何人都不可能被它们诱惑；第二，天主加于我们的一切，无论此刻在我们看来是悲伤还是喜乐，都是如最慈爱的父亲和最怜悯人的医生为我们的益处而施加的，因此，人就好像被交付给老师的管教，并被贬抑，好让他们在离开此世时，能以更大的纯洁被提升到另一个生命，或受到更轻的惩罚，因为他们如宗徒所说，现时被交付给了\BibleQuote{撒殚，败坏他的肉体，为使他的灵魂在主耶稣的日子可以得救。}\BibleRef{格前 5:5}\par

% structure: p0088 source=h2 target=subsection reason=relative-heading-rank; base=h2

\subsection{第二十九章}

\Emph{质疑：为什么受不洁之神折磨的人被隔绝于主的共融之外？}\par

\Person{Germanus}：我们怎能看见他们不仅被众人鄙视和回避，而且在我们各省中实际上一直被禁止领受主的圣体，正如福音所言：\BibleQuote{不要把圣物给狗，也不要把你们的珍珠丢在猪前}；\BibleRef{玛 7:6}而您却告诉我们，不知为何我们应当认为这种诱惑的羞辱临到他们，是为着他们的净化和益处？\par

% structure: p0091 source=h2 target=subsection reason=relative-heading-rank; base=h2

\subsection{第三十章}

\Emph{对所提问题的回答。}\par

\Person{Serenus}：倘若我们拥有上述的知识，或更确切地说，拥有这样的信德，即相信万事皆由天主安排，为使我们灵魂获益而安排，那么我们不仅绝不会鄙视他们，反而会不断为他们祈祷，如同为自己肢体祈祷，并以全副心神和最深的爱怜同情他们（因为\BibleQuote{若一个肢体受苦，所有的肢体都一同受苦} \BibleRef{格前 12:26}），因为我们知道，既然他们是我们的肢体，没有他们我们便不可能达至成全；正如我们读到，我们的先辈没有我们也不能得到许诺的圆满，正如宗徒论及他们说：\BibleQuote{这些人虽然都因信德得了褒扬，却没有获得那所应许的，因为天主为我们预备了更好的事，以致他们若没有我们，不得成全。} \BibleRef{希 11:39} 但我们从未记得圣体圣事曾被禁止给他们；相反，倘若可能，他们认为应当\Emph{每日}赐予他们；而且，按照您在此不当引用的福音经句：\BibleQuote{不要把圣物给狗，} \BibleRef{玛 7:6} 我们不应相信圣体圣事成了魔鬼的食物，而非灵魂和身体的净化与保护；因为当人领受它时，它仿佛燃烧并驱逐那盘踞在肢体中或试图潜伏其中的恶神。最近我们就这样看到\Person{安当尼各}院长和许多其他人被治愈。因为若敌人看见那附魔之人被截断于天上之药，便会更加虐待他，而且一旦见他越发远离这属神的疗法，便会更频繁、更可怖地诱惑他。\par

% structure: p0094 source=h2 target=subsection reason=relative-heading-rank; base=h2

\subsection{第三十一章}

\Emph{论那些不被允许遭受这些暂时诱惑的人更值得怜悯。}\par

但我们应当视那些人为真正可怜和悲惨的，他们虽然以各种罪恶和邪恶玷污自己，但身上不仅没有出现魔鬼附身的明显迹象，也没有与他们的行为相称的诱惑，更没有施加任何惩罚的鞭笞。因为他们在今世得不到任何迅速及时的疗法，他们的\BibleQuote{硬心和不知悔改的心}，由于超过今世的惩罚，\BibleQuote{为自己积蓄忿怒，直到忿怒和启示天主正义审判的日子}，\BibleQuote{那里的虫子不死，火也不灭。} 先知仿佛对圣人们的苦难感到困惑，当他看见他们遭受各种损失和诱惑，而另一方面又看见罪人不仅毫无羞辱的鞭笞而通过此世，甚至因巨大的财富和万事亨通而欢喜，便充满无法抑制的义愤和热忱，喊道：\BibleQuote{至于我，我的脚几乎跌倒，我的步伐险些滑倒。因为我嫉妒恶人，看见罪人享平安。因为他们没有死前的痛苦，他们的力量也强壮。他们不像别人受劳苦，也不像别人遭鞭打。} 因为他们今后将与魔鬼一同受罚，而在今世他们不配与世人一同承受儿子份内的管教和鞭笞。\Person{耶肋米亚}在与天主谈论恶人的兴盛时，虽然从未声称怀疑天主的公义，正如他说\BibleQuote{上主，若我与祢争辩，祢是公义的}，但在探究这种不平等的缘由时，继续说：\BibleQuote{但我仍要向祢说公义的话：恶人的道路为何亨通？一切背信作恶的人为何平顺？祢栽植了他们，他们就扎根，长大，结出果实。祢的口唇离他们近，他们的心却离祢远。} \BibleRef{耶 12:1} 当上主藉先知哀叹他们的毁灭，并急切地指示医生和医师去医治他们，并好似催促他们作类似的哀叹说：\BibleQuote{巴比伦忽然倾覆，毁坏了。为她哀号吧！拿乳香来治她的创伤，或许可治愈。} 于是，无所指望的受托照管人类救恩的天使们回答；或者至少先知以宗徒、属神的人和看见他们心灵刚硬、不知悔改的医生的口吻说：\BibleQuote{我们医治了巴比伦，她却不痊愈。我们离开她吧！各人归回本地，因为她的审判已上达于天，高耸入云。} \BibleRef{耶 51:8} 关于他们绝望的软弱，\Person{依撒意亚}遂以天主的位格对耶路撒冷说：\BibleQuote{从脚掌到头顶，没有一处完好，尽是伤口、青肿与流血的创伤；没有挤净，没有包扎，也没有用油润敷。} \BibleRef{依 1:6}\par

% structure: p0097 source=h2 target=subsection reason=relative-heading-rank; base=h2

\subsection{第三十二章}

\Emph{论空中权势中存在的各种欲望和意愿。}\par

然而，有充分证据显示，在污秽的鬼神中存在着与人类同样多的欲望。因为其中有些鬼，通常被称为\OriginalTerm{游荡者}{Plani}，表现得如此诱人和嬉戏，以至于当它们持续占据某些地方或道路时，它们并非以折磨那些被它们欺骗的过路人为乐，而仅仅满足于嘲笑和戏弄他们，试图使他们疲惫不堪而非伤害他们；而另一些鬼则只在夜间无害地附身于人，但还有一些鬼却如此暴怒和凶残，不仅满足于可怕地撕裂被附之人的身体，甚至还急于扑向远处路过的人，以最野蛮的屠杀攻击他们——正如福音中所描述的那些人，因惧怕他们，无人敢走那条路。毫无疑问，这些以及类似之鬼在其无法满足的狂怒中以战争和流血为乐。我们发现其他鬼会用虚妄的骄傲冲击被它们附身之人的心（这些鬼通常被称为\OriginalTerm{浮夸者}{Bacucei}），以至于它们伸展身体超出其正常高度，时而自高自大、趾高气扬，时而以平常和温和的方式屈就于平静和友善的状态；而由于它们自认为是伟人和众人惊叹的对象，时而弯腰鞠躬表示崇拜更高的权能，时而又认为受到他人崇拜，于是以骄傲或卑微的方式做出所有表示真正服侍的动作。我们还发现有些鬼不仅热衷于谎言，还煽动人亵渎神明。对此我们亲眼见证，曾听到一个魔鬼公开承认它通过\Person{亚流}和\Person{欧诺米}的口宣扬了邪恶和不敬神的教义。同样地，我们在列王纪第四卷中读到，其中一个魔鬼公开宣称：\Quote{我要出去，}他说，\Quote{并要在他众先知的口中作说谎的灵。}\BibleRef{列上 22:22}因此，当宗徒责备那些被它们欺骗的人时，接着补充说：\Quote{听从诱惑人的神和鬼魔的教导，在虚伪中说谎。}\BibleRef{弟前 4:1}还有许多别的聋哑魔鬼，福音书可以作证。有些鬼煽动情欲和放荡，先知也证实说：\Quote{淫乱的神迷惑了他们，他们偏离了他们的天主。}\BibleRef{欧 4:12}同样地，圣经的权威教导我们，有黑夜、白昼和晌午的魔鬼。但要通览整部圣经并一一列举它们的种类，则耗时过长，因为这些种类在先知著作中被称为\OriginalTerm{半驴半人}{onocentaurs}、\OriginalTerm{半羊人}{satyrs}、\OriginalTerm{塞壬女妖}{sirens}、\OriginalTerm{女巫}{witches}、\OriginalTerm{嚎叫者}{howlers}、\OriginalTerm{鸵鸟}{ostriches}、\OriginalTerm{刺猬}{urchins}；在圣咏集中称为\OriginalTerm{蝮蛇}{asps}和\OriginalTerm{蜥蜴}{basilisks}；在福音中被称为\OriginalTerm{狮子}{lions}、\OriginalTerm{龙}{dragons}、\OriginalTerm{蝎子}{scorpions}；宗徒则称它们为\OriginalTerm{今世的首领}{the prince of this world}、\OriginalTerm{这黑暗的统治者}{rulers of this darkness}和\OriginalTerm{邪恶的鬼神}{spirits of wickedness}。所有这些名称，我们不应视为随意或偶然的称呼，而是藉着那些在我们中间或多或少有害或危险的野兽的标记，暗指它们的凶残和疯狂，并通过将它们与那些在野兽或蛇类中因邪恶的优越性而具有特别毒害能力的东西相比较，依其名称称呼它们，以至于给这一个取了狮子的名，因为它狂暴的怒气；给那一个取了蜥蜴的名，因为它致命的毒液，在人察觉之前就已致死；又给另一个取了半驴半人或刺猬或鸵鸟的名，因为它迟钝的恶毒。\par

% structure: p0100 source=h2 target=subsection reason=relative-heading-rank; base=h2

\subsection{第三十三章}

\Emph{关于天空中邪恶势力这些差异之起源的提问。}\par

\Person{格尔玛努斯}：我们当然不怀疑宗徒所列举的那些阶层指的是它们：\Quote{因为我们战斗不是对抗血肉，而是对抗首领、对抗掌权者、对抗这黑暗世界的霸主、对抗天界中邪恶的鬼神。}\BibleRef{弗 6:12}但我们想知道，它们之间的这种差异从何而来，或者这些邪恶的等级是如何存在的？它们被造是为了这些吗？为了遇见这些邪恶的阶层，并以某种方式服务于这种邪恶？\par

% structure: p0103 source=h2 target=subsection reason=relative-heading-rank; base=h2

\subsection{第三十四章}

\Emph{对所提问题的回答的推迟。}\par

\Person{塞雷努斯}：虽然你们的提议会使我们失去整夜的休息，以致我们不会注意到黎明的临近，并会贪婪地将我们的谈话一直延长到日出，然而，因为如果开始探究所提问题的解决，它将把我们投入一个广阔而深奥的问题之海，而我们所拥有的短暂时间不允许我们横渡它，我认为更合适的做法是将其留到另一个夜晚来考虑，那时，藉着提出这个问题，我将从你们极其愉快的交谈中获得一些属灵的喜悦和更丰硕的果实，并且如果圣神赐给我们顺风，我们将更能自由地深入所提问题的复杂性。因此，让我们稍作睡眠，在黎明临近时驱散涌上眼帘的睡意，然后我们一起去教堂，因为遵守主日要求我们这样做；礼仪结束后，我们再回来，按你们所愿，以加倍的喜悦讨论主可能为我们的共同益处而赐予我们的一切。\par

\SourceRecord{Translated by C.S. Gibson.；Nicene and Post-Nicene Fathers, Second Series；Vol. 11.；Edited by Philip Schaff and Henry Wace.；Buffalo, NY: Christian Literature Publishing Co.,；1894.；<http://www.newadvent.org/fathers/350807.htm>.}

% structure: p0001 source=h1 target=section reason=page-title

\section{第八次会谈}

\Signature{圣若望·卡西安}\par

塞雷诺院牧的第二次会谈。论宰制者。\par

% structure: p0004 source=h2 target=subsection reason=relative-heading-rank; base=h2

\subsection{第一章}

塞雷诺院牧的款待。\par

我们完成了一天的职责，会众从教堂散去后，我们回到老者的单人房间，享受了一顿极其丰盛的晚餐。因为平时他日常用餐的桌上只放着加了几滴油的酱汁，这次他却调了一点汤汁，并比平时更为慷慨地多倒了些油。原来，每位隐修士在准备用日常餐食时，都会倒上几滴油，并非要品尝其滋味（因为油量极少，我甚至不能说足以润滑喉咙和上下颚的通道，更不用说咽下去了），而是借以抑制心灵的骄傲（倘若他更加克己，这种骄傲必然会悄悄潜入）以及虚荣心的刺激——因为他的克己越是隐秘、无人看见，就越是狡诈地不断诱惑那隐藏它的人。随后，他摆上了餐桌上的盐和三颗橄榄；接着又拿出一个篮子，里面装有烘干的野豌豆，他们称之为\OriginalTerm{烤豆}{trogalia}，我们从每人取了五粒，两颗李子以及一个无花果。因为在那个旷野里，任何人超过这个数量都被视为不当。用餐完毕，我们再次请求他履行承诺，解答问题。老者说：\Quote{让我们听听你们的问题，我们将其推迟到现在考虑。}\par

% structure: p0007 source=h2 target=subsection reason=relative-heading-rank; base=h2

\subsection{第二章}

关于各类属灵邪恶的陈述。\par

日耳曼纽斯说：我们想知道，敌对人类的众多权势的起源以及它们之间的区别是什么？蒙福的宗徒概括如下：\BibleQuote{我们战斗不是对抗血肉，而是对抗宰制者、对抗掌权者、对抗这黑暗的世界霸主、对抗天界中的属灵邪恶。} \BibleRef{弗 6:12} 又说：\BibleQuote{无论是天使、宰制者、掌权者，还是其他任何受造物，都不能使我们与天主在基督耶稣我们的主内的爱隔绝。} \BibleRef{罗 8:38} 那么，所有这些嫉妒我们的恶意的敌意从何而来？我们是否应该相信这些权势是主为此目的而创造的，即按这些等级和秩序与人争战？\par

% structure: p0010 source=h2 target=subsection reason=relative-heading-rank; base=h2

\subsection{第三章}

关于圣经中多种食物的答案。\par

\Person{塞雷努斯}说：\Quote{圣经的权威清楚地说明了某些事，甚至对完全无知的人也如此浅显明了，以至于不仅不为任何隐秘含义的晦涩所遮蔽，甚至不需要任何解释的帮助，而是将其意义和感觉展现在词句和字母的表面；但有些事却如此隐秘地藏在奥秘中，以至于为我们在讨论和解释它们时提供了极大的技巧和谨慎的领域。显然，天主这样安排有许多原因：首先，惟恐神圣的奥秘若不被某种属神意义的帷幕所遮盖，便会同样暴露于所有人的知识和理解，即亵渎者和信徒，从而在善与智慧方面，懒惰者和热忱者之间便没有区别；其次，在那些确实是信仰之家的人中，当智力能力的巨大差异展现在他们面前时，就有机会责备懒惰者的懈怠，并证明热忱者的敏锐和勤奋。因此，圣经被恰当地比作一块富饶肥沃的田地，它虽然产生并结出许多不经火煮就适宜人食用的好东西，但也产生一些除非经过火的加热而调和软化，否则便不适于人用甚至有害的东西。但有些东西在两种状态下都自然适合使用，以至于即使未经煮熟，其生涩状态也并非令人不快，但经过火烹加热后更加可口。还有许多东西只适合非理性受造物、牲畜、野兽和鸟类的食物，而完全不适合作为人的食物，这些东西在未经过任何火处理的生涩状态下，却有助于牲畜的健康和生命。我们可以清楚地看到，同样的原则也适用于那最丰饶的属神圣经园地。在这园地中，有些事物按其字面意义清晰明亮地照耀出来，以至于它们不需要任何更高层次的解释，仅凭字词的简单声音就为听众提供了丰富的食物和营养，例如这段：\BibleQuote{以色列，你要听，上主你的天主是唯一的上主；}以及：\BibleQuote{你当全心、全灵、全力爱上主你的天主。}\BibleRef{申 6:4}但有些事物，除非通过\OriginalTerm{寓意解经}{allegorical interpretation}而减弱，并通过圣神的火的考验而软化，否则无法成为内在之人的健康食物，反而会对他造成损害和损失；接受它们不仅无益，反而有害，例如这段：\BibleQuote{你们的腰要束上，灯要点着；}以及：\BibleQuote{没有剑的，要卖了他的外衣买剑；}以及：\BibleQuote{凡不背起自己的十字架跟随我的，不配作我的门徒。}这段经文被一些最热忱的隐修士（他们\BibleQuote{对天主确有热心，但不按真知}\BibleRef{罗 10:2}）按字面理解，以至于他们做了木制十字架，经常背在肩上，结果在所有看见的人眼中，不是得到造就，而是成为笑柄。但有些经文可以按两种方式——即历史和寓意——正确而适当地理解，以至于任何一种解释都为灵魂提供疗愈之饮，例如这段：\BibleQuote{有人打你的右脸，连左脸也转过来由他打；}以及：\BibleQuote{有人在这城迫害你们，就逃到另一城去；}以及：\BibleQuote{你若愿意是成全的，去，变卖你所有的，施舍给穷人，你必有宝藏在天上，然后来跟随我。}它也确实产生了\BibleQuote{牲畜的青草}（圣经所有田地都充满了这种食物），即清楚简单的历史叙述，借此，简单的人和那些不能完全和正确理解的人（关于他们，经上说：\BibleQuote{主，祢要拯救人和牲畜}）可以按照他们能力的状况和度量，变得更强壮、更有活力，以应对他们艰苦的工作和实际生活的劳动。}\par

% structure: p0013 source=h2 target=subsection reason=relative-heading-rank; base=h2

\subsection{第四章}

\Emph{论圣经可能具有的双重意义。}\par

因此，对于那些带有清晰解释而提出的段落，我们也可以经常确定其意义，并大胆陈述我们自己的意见。但那些圣神为留给我们默想和操练而插入圣经中、具有隐秘意义的段落——祂希望其中一些通过不同的证据和推测来收集——应当逐步而谨慎地汇集起来，以便论述或支持它们的人能凭其判断安排其断言和证据。因为有时对同一主题表达不同意见时，两种观点都可以被视为合理，并且在不损害信仰的情况下，可以坚定地或存疑地（即既不完全相信也不完全拒绝）予以保持，而且第二种观点不必妨碍第一种，只要两者都不与信仰相悖：例如在此情形中：厄里亚以若翰的身份而来\BibleRef{玛 11:14}，并又将作为主来临的前驱；又如关于\BibleQuote{那行毁坏可憎之物}\BibleQuote{站在圣地}，即藉由我们读到立在耶路撒冷圣殿中的朱庇特偶像，这又将藉着假基督的来临而站在教会中，以及福音中随后的一切事物——我们视这些事物已在耶路撒冷被俘之前应验，又将在世界末日应验。在这些事上，两种观点互不抵触，第一种解释也不妨碍第二种。\par

% structure: p0016 source=h2 target=subsection reason=relative-heading-rank; base=h2

\subsection{第五章}

\Emph{论所提出的问题应属于以中立或存疑方式持有的那一类事物。}\par

因此，由于我们所提出的问题，似乎尚未在人们中间得到充分或经常的讨论，对大多数人来说是清晰的，而正因为此，我们所提出的内容或许在有些人看来是可怀疑的，我们应当调整自己的观点（因为它并不干涉对天主圣三的信德），使之被归入那些应被持为可疑的事物之列；尽管它们所依据的并非仅仅是通常给予猜测和臆断的见解，而是清晰的圣经证据。\par

% structure: p0019 source=h2 target=subsection reason=relative-heading-rank; base=h2

\subsection{第六章}

\Emph{论天主未造任何邪恶之物}\par

我们切不可承认天主曾创造了任何本质上是恶的事物，正如经上记载：\BibleQuote{天主看了祂所造的一切，认为样样都好。} \BibleRef{创 1:31} 因为如果它们被造时便如今日之状态，或被造为此目的，即占据这些恶意的位置，并时刻准备欺骗和毁灭人类，那么我们便与上述圣经的观点相悖，诽谤天主是邪恶的创造者和根源，认为祂亲自形成了完全邪恶的意志和本性，并为此目的创造了它们，即让它们永远固守于邪恶，永不转向善意。我们随后从教父传统中接受了关于这种差异的理由，这传统源自圣经的泉源。\par

% structure: p0022 source=h2 target=subsection reason=relative-heading-rank; base=h2

\subsection{第七章}

\Emph{论掌权者或权势者的起源}\par

没有信友会质疑这一事实：在这有形受造界形成之前，天主就已创造了属神和天上的权势，目的是为了使它们因深知自己是由造物主的良善从无中造出、并承受如此荣耀和福乐，而向祂不断感恩、永不停息地赞美祂。因为我们也不可想象，天主是在这世界的形成中才开始祂的创造和工作，仿佛在之前那无数世代中，祂未曾考虑过眷顾和神圣的安排，又仿佛我们可以相信，祂因没有对象可展示其良善的恩赐，而处于孤独、与一切慷慨无缘的状态——这对于那无限、永恒、不可理解的尊威而言，是一种太过贫乏和不合适的想象。正如主亲自论及这些权势时所说：\Quote{当星辰一同形成时，我的众天使都高声赞美我。} 那些在星辰创造时在场者，最清楚地被证明是在那\Quote{起初}之前被创造的，即天地被造的那个\Quote{起初}，因为他们被说成以高声和惊叹赞美造物主，因他们所见的一切有形受造物从无中生出。所以，在那按时间而言的起初（即梅瑟所讲述的，按历史和犹太人的解释，指这世界的时代；但这并不妨碍\Emph{我们的}解释，即我们说明万物的\Quote{起初}就是基督，圣父在祂内创造了万物，正如经上所说：\BibleQuote{万物是藉着祂造的；凡受造的，没有一样不是藉着祂造的。} \BibleRef{若 1:3}）——我说，在那创世纪的时间起初之前，毫无疑问，天主早已创造了所有这些权势和天上的德能；宗徒按次序列举并描述说：\BibleQuote{因为在天上和地上的一切，无论看得见的、看不见的，或是上座的、统治的、率领的、掌权的，都是在基督内受造的。万物都是藉着祂，并且是为了祂而受造的。} \BibleRef{哥 1:16}\par

% structure: p0025 source=h2 target=subsection reason=relative-heading-rank; base=h2

\subsection{第八章}

\Emph{论魔鬼和天使的堕落}\par

这样，我们清楚地看到，在他们当中有些首领由厄则克耳和依撒意亚的哀歌中坠落，我们得知提洛的君王或那清晨升起的\OriginalTerm{路济弗尔}{Lucifer}以悲伤的挽歌被哀悼：主对厄则克耳这样说：\BibleQuote{人子，你应对提洛的君王唱哀歌，向他说：上主天主这样说：你原是完美的印记，充满智慧，十分美丽。你曾在伊甸乐园中，佩戴各种宝石：红玉髓、黄玉、碧玉、橄榄石、缠丝玛瑙、绿柱石、蓝宝石、红宝石和祖母绿；你的饰物是金，又是笛子，在你受造之日便已预备。你是伸展而护卫的\OriginalTerm{革鲁宾}{cherub}，我将你安置在天主圣山上，你在烈火石中行走。从你受造之日，你的道路就完备无缺，直到在你身上发现罪孽。因你交易繁多，你内心充满横暴，以至犯了罪，我将你从天主圣山上抛下，消灭你，护卫革鲁宾，从烈火石中。你因美丽而心高气傲，你因华美而丧失智慧，我将你摔到地上，把你放在君王面前，使他们目睹你。因你罪孽众多，因你贸易的不义，你亵渎了你的圣所。}\BibleRef{见：则 28:11}依撒意亚也对另一个这样说：\BibleQuote{啊，\OriginalTerm{路济弗尔}{Lucifer}，清晨之星，你怎么从天坠落？你怎样被摔到地上，伤害了万民？你曾在心中说：我要升到天上，我要高举我的宝座在天主星辰之上，我要坐在集会之山上，在极北之处。我要升到云层之上，我要与至高者相似。}\BibleRef{见：依 14:12}但圣经告诉我们，这些坠落不是单独从他们福乐的地位顶峰，因为它告诉我们那龙拖走了三分之一的星辰。\BibleRef{见：默 12:4}一位宗徒还更明确地说：\BibleQuote{至于那些没有保持自己原本地位，而离开自己居所的天使，主用永久的锁链将他们拘留在黑暗中，直到大审判的日子。}对我们所说的：\BibleQuote{你们却要像人一样死去，像一位王子一样跌倒。}这暗示什么，难道不是许多王子坠落了吗？通过这些见证，我们可以推断这种差异的原因：即他们要么仍然保持那由他们各自受造时的原初等级的地位，如同敌对势力据有与圣洁和天上的美德相似的地位；要么他们自己从天上被摔下，却作为他们恶行的报应，他们在邪恶中步步升高，以模仿那些坚持下来的美德的方式，这些等级和头衔在悖逆中自称。\par

% structure: p0028 source=h2 target=subsection reason=relative-heading-rank; base=h2

\subsection{第九章}

\Emph{一个反对意见，声称魔鬼的堕落起源于天主的欺骗。}\par

杰尔玛诺：到目前为止，我们曾相信魔鬼败坏和坠落的理由与开端，即他被逐出天上地位，更特别是出于嫉妒，当他怀着恶意的诡计欺骗了亚当和厄娃时。\par

% structure: p0031 source=h2 target=subsection reason=relative-heading-rank; base=h2

\subsection{第十章}

\Emph{魔鬼堕落之始的答覆}\par

塞雷诺：《创世纪》的段落表明，那不是他堕落和败坏的开端，因为在欺骗他们之前，它已经认为他受到了蛇这一名称的羞辱，那里说：\BibleQuote{蛇是更聪明的}，或如希伯来抄本所说：\BibleQuote{比上主天主所造的一切野兽更狡猾。}\BibleRef{见：创 3:1}那么你看，他在欺骗第一个人之前就已经从天使的圣洁中坠落了，以致他不仅配得上这一称号的羞辱，而且实在在邪恶的托词中超越地上所有野兽。因为圣经不会用一个好天使来称呼这样的名字，也不会对那些仍处于福乐状态中的天使说：\BibleQuote{蛇比地上一切野兽更聪明。}因为这一称号不可能适用于加俾额尔或弥额尔，甚至对任何一个善人也不合适。因此，蛇这一称号和与野兽的比较最清楚地暗示的不是天使的尊严，而是背教者的耻辱。最后，导致他欺骗人的嫉妒和诱惑的机缘源于他先前堕落的基础，因为他看到那个人，刚刚从尘土中被塑造出来，将被召唤到那个荣耀中，而他记得自己，虽然仍是一位首领，却从那里坠落了。因此，他因骄傲而第一次堕落，这使他得到了蛇的名字，随后又因嫉妒而第二次堕落；由于这次堕落，他仍然拥有某种正直，可以与人进行一些交流和商议，因此主非常合理地将他降到最深的深渊，使他不再像以前那样直立和仰望高处，而应匍匐于地爬行，用腹部爬行，以地上的食物和罪的行为为食，从此宣布他的秘密敌意，并在自己与人之间建立对我们有利的敌意，对我们有益的纷争，以便人们对他保持警惕，将他视为危险的敌人，他便不能再以欺骗性的友谊伤害他们。\par

% structure: p0034 source=h2 target=subsection reason=relative-heading-rank; base=h2

\subsection{第十一章}

\Emph{欺骗者与被欺骗者的惩罚}\par

但在此事上，为了躲避恶谋，我们应从这一事实学得一个特别的教训：虽然欺骗的始作俑者受了恰当的惩罚与定罪，然而被引入歧途者却也未完全免受刑罚，尽管比那欺骗的始作俑者稍轻。这一点我们看得很清楚。因为被欺骗的亚当——或者更确切地说，按宗徒的话，\Quote{他\Emph{并非}被欺骗}——而是顺应那被欺骗者的意愿，似乎同意了致命的服从——仅被判处劳苦和额上的汗水，这并非直接诅咒他本人，而是诅咒土地及其荒芜。但那诱使他如此做的女人，却增加了苦楚、疼痛与悲伤，并被置于永久臣服的轭下。而那首先唆使他们犯此过犯的蛇，则受永久的诅咒。为此，我们当以极大的谨慎和小心防备恶谋，因为它们既给其始作俑者带来惩罚，也不让被它们所欺骗者免于罪责和惩罚。\par

% structure: p0037 source=h2 target=subsection reason=relative-heading-rank; base=h2

\subsection{第十二章}

\Emph{论魔鬼的群集，以及它们在我们的大气中不断引起的骚动。}\par

但那介于天地之间的大气，经常充满着一大群浓密的魔鬼，它们并非安静或懒散地在其中飞翔；因此，天主的眷顾极为幸运地将它们从人的眼前撤去。因为若人因惧怕它们的攻击，或因它们随意变化转变形态的形状而惊骇，就必因难以忍受的恐惧而失魂落魄，或因肉眼无法目睹此类事物而晕厥；否则，人会因不断的榜样和效法而日益堕落腐败，从而在人与空中不洁的权势之间产生一种危险的熟稔和致命的交往；如今在人间所犯的罪行，尚可借墙壁、围栏、距离或空间，或因某种羞耻和混乱而被遮掩；但若能常用明眼看它们，人便会受到更大的疯狂刺激，因为无时无刻不见它们行恶，而它们从不停歇，因为不像我们，身体的疲乏、事务的忙碌或饮食的忧虑，有时迫使它们甚至违背意愿，放弃已开始实施的目的。\par

% structure: p0040 source=h2 target=subsection reason=relative-heading-rank; base=h2

\subsection{第十三章}

\Emph{论敌对势力将针对人的攻击甚至施于彼此。}\par

因为十分明显，它们也将用来攻击人的那些打击施于彼此；同样，它们因某种天生的喜爱邪恶，不知疲倦地以争斗来促进它们为某些民族所挑起的纷争和冲突；这点我们在达尼尔先知的异象中读得很清楚，当时天使加俾额尔说：\BibleQuote{不要害怕，达尼尔！因为自从你用心去领悟，在你天主前自卑的第一天，你的祈祷已蒙应允，我是因你的祈祷而来的。但是波斯国的君侯反对了我二十一天；看，总领天使弥额尔来援助我，我便停留在波斯国的君王那里。现在我来，是告诉你将来你的百姓要遭遇的事。} \BibleRef{达 10:12} 我们绝不可能怀疑那波斯国的君侯是一个敌对势力，它袒护与天主子民为敌的波斯民族；因为出于嫉妒，它阻碍天使总领加俾额尔为先知祈祷所求解的疑问而带来的益处，阻止那救慰的天使过快到达达尼尔那里，并坚强天主所眷顾的子民；后者说，即使如此，若不是总领天使弥额尔前来援助他，与波斯国的君侯作战并干预，保护他免受攻击，他仍不能前来，故二十一天后才得以到达先知那里。稍后不久又说：\BibleQuote{他说：\InnerQuote{你知道我为什么来吗？现在我要回去与波斯的君侯交战；我一出现，希腊的君侯就来。但我要告诉你那载在真理书上的事；除了你们的君侯弥额尔外，没有支持我的人。}} \BibleRef{达 10:20} 又说：\BibleQuote{那时，那保护你百姓的伟大的君侯弥额尔就起来。} \BibleRef{达 12:1} 因此我们读到，同样，另一位被称为希腊国的君侯，他既是那属于他的民族的护守，似乎既反对波斯民族，也反对以色列子民。由此我们清楚看到，敌对势力在民族之间挑起那些纷争和冲突，它们自己在唆使下也互相攻击，或因胜利而欢欣，或因失败而沮丧，彼此之间不能和睦共处，因为每个势力都怀不安的嫉妒，为其所管辖者，对抗另一民族的护守。\par

% structure: p0043 source=h2 target=subsection reason=relative-heading-rank; base=h2

\subsection{第十四章}

\Emph{论属灵的邪恶如何获得\Quote{权力}或\Quote{统权者}的名称。}\par

我们于是可以看到清楚的理据，除了我们上文所阐述的那些观念之外，为何它们被称为宰制者或掌权者；即，因为它们统治并治理不同的民族，并至少掌管较低级的灵体和魔鬼，关于这些，福音书借它们自己的自认给了我们证据，说明存在军团。因为它们不能被称为主，除非它们有某些可以行使主权管辖的对象；也不能被称为掌权者或宰制者，除非有某些对象可对之行使权柄：这一点我们在福音书中，从法利塞人的亵渎之言可以非常清楚地看出：\BibleQuote{祂藉贝耳则步，魔鬼的首领，驱除魔鬼} \BibleRef{路 11:15}；因为我们也发现它们被称为\BibleQuote{黑暗的统治者} \BibleRef{弗 6:12}，并且它们中有一个被称为\BibleQuote{这世界的首领} \BibleRef{若 14:30}。但蒙福的宗徒宣布，将来当万物都屈服于基督时，这些等级将被消灭，说道：\BibleQuote{当祂把国度交给天主父时，当祂消灭了一切宰制者、一切掌权者和一切异能者之后} \BibleRef{格前 15:24}。这当然只有在它们从那些我们知道在此世有掌权者、异能者和宰制者负责管辖的对象之管辖中被移除时，才会发生。\par

% structure: p0046 source=h2 target=subsection reason=relative-heading-rank; base=h2

\subsection{第十五章}

\Emph{论圣洁天上的掌权者得名天使和总天使并非无故。}\par

无人怀疑这些相同的品级称号并非无缘无故地赋予较善的种类，并且它们是职衔和功绩或尊严的名号，因为显然它们被称为天使，即来自传递消息的职务的使者，而名称的适当性教导它们被称为\Quote{总天使}，因为它们统辖天使；\Quote{异能者}，因为它们对某些人行使宰制权；\Quote{宰制者}，因为它们有某些对象可作其首领；\Quote{宝座者}，因为它们如此接近天主，如此亲密贴近祂，以致天主的威严特别安息于它们，如同在神圣宝座上，而且以某种方式必然安歇于它们之上。\par

% structure: p0049 source=h2 target=subsection reason=relative-heading-rank; base=h2

\subsection{第十六章}

\Emph{论魔鬼对其首领所显示的服从，正如在某位弟兄的受害者中所见。}\par

但不洁之灵被更恶的掌权者统治并服从于他们，我们不仅从圣经的那些段落中发现，福音书记载当法利塞人诽谤主时，祂回答说：\BibleQuote{我若藉贝耳则步，魔鬼的首领，驱除魔鬼} \BibleRef{路 11:19}；而且我们也通过清晰的异象和圣人们的许多经验而学到这一点。因为当我们的一个弟兄行走在这旷野中时，天色渐晚，他找到一个洞穴并停在那里，打算在里面念晚课，他还在咏唱圣咏时，午夜已过。当他做完课业后，在恢复疲惫身体之前稍坐片刻，突然他开始看见无数成群的魔鬼从四面八方聚集而来，他们以庞大的队伍和长长的行列前进，有些在前，有些跟随他们的首领；首领终于来到，身材比所有其他魔鬼更高大、更可怕；一个宝座被放置后，他如同坐在某种崇高审判台上，开始通过彻底调查来审查他们中每一个的行为；那些说尚未能诱骗对手的，他命令被羞愧和耻辱地赶出他的视线，作为懒惰和怠惰的，用愤怒的怒火斥责他们浪费了如此多的时间，徒劳无功；但那些报告说已欺骗了指派给他们的对象的，他在所有魔鬼的欢腾和鼓掌中，以最高的赞扬释放他们，作为最勇敢的战士，并以最著名的榜样给所有其余者；当在这个数目中，有一个最邪恶的灵出现，因要报告某种辉煌胜利而欢喜时，他提到一个非常著名的隐修士的名字，并宣称在连续攻击他十五年后，终于打败了他，就在当晚因邪淫之罪毁灭了他，因为他不仅促使他与某个献身于主的少女犯奸淫，还实际上说服他娶了她并留住她。当这消息引起欢呼时，他被黑暗之首以最高赞扬所称颂，并带着极大荣誉离开。因此，当在破晓时分，整个魔鬼群已从他眼前消失，这位弟兄对不洁之灵的断言感到怀疑，反而认为他想用古老的惯常欺骗来引诱他，并以乱伦之罪玷污一个无辜的弟兄，他记起福音中的那些话，即\BibleQuote{他不持守真理，因为在他内没有真理。他说谎时，是出于自己说的，因为他是说谎者，也是说谎之父} \BibleRef{若 8:44}；于是他前往\Place{培琉西雍}，他知道那恶灵宣称已被毁灭的人住在那里；因为这弟兄对他非常熟悉，当他询问后，他发现就在那污秽之魔向他的同伴和首领宣布他的坠落的那同一个晚上，他离开了他原来的隐修院，去了那个城镇，并因与所述少女的堕落而误入歧途。\par

% structure: p0052 source=h2 target=subsection reason=relative-heading-rank; base=h2

\subsection{第十七章}

\Emph{论及两个天使常依附于每个人。}\par

因为\WorkTitle{圣经}证实，二位天使，一善一恶，依附于我们每一个人。关于善的天使，救主说：\BibleQuote{你们切勿轻视这些小子中的一个；因为我告诉你们，他们的天使在天上，常常看见我天父的面。}\BibleRef{玛 18:10}又说：\BibleQuote{上主的天使在敬畏祂的人四围安营，并拯救他们。}此外，\WorkTitle{宗徒大事录}中论及\Person{伯多禄}时也说：\BibleQuote{是他的天使。}\BibleRef{宗 12:15}但\WorkTitle{牧者书}关于这两类天使教导得十分详尽。然而，若我们思考那攻击真福\Person{约伯}的，便会清楚得知：那常谋算反对他、却从未能引诱他犯罪的，正是他；因此他向主求权能，因他不是被\Person{约伯}的德行所胜，而是被那常护佑他的上主的保护所击败。关于\Person{犹达斯}也记载说：\BibleQuote{愿魔鬼站在他的右边。}\par

% structure: p0055 source=h2 target=subsection reason=relative-heading-rank; base=h2

\subsection{第十八章\newline{}}

\Emph{论存在于敌对诸灵中的邪恶程度，如两位哲学家的案例所示。}\par

但关于魔鬼之间的区别，我们通过那两位哲学家学到了很多；他们从前藉魔法行为，对魔鬼的懒惰、勇气和凶恶的邪恶常常大有经验。因为他们轻视\Person{圣安当}，视之为粗野无文之辈，且想，即使不能再进一步伤害他，至少也要用魔法幻象和魔鬼的伎俩将他逐出他的隐修院，就派遣了最污秽的邪魔攻击他；他们之所以被迫发动此攻击，乃是出于嫉妒的刺痛，因为每天都有海量人群到他那里，视他为天主的仆人。当这些最凶恶的魔鬼甚至不敢接近他时（他正一面在自己的胸和额上画十字圣号，一面致力于祈祷和恳求），它们便毫无结果地返回到那些派遣它们的人那里；这些人又派遣了其他在邪恶上更绝望的魔鬼去攻击他，而当它们也徒耗力量、毫无所成地返回，仍有更强大的魔鬼被派去对抗这位得胜的基督的战士，却无法胜过他；他们费尽心机用尽魔法的这一切大阴谋，只证明了基督徒信仰的极大价值，以致那些凶猛强大的阴影（他们以为如果指向太阳月亮就能遮蔽它们），不仅不能伤害他，甚至连片刻都不能将他引离他的隐修院。\par

% structure: p0058 source=h2 target=subsection reason=relative-heading-rank; base=h2

\subsection{第十九章\newline{}}

\Emph{论魔鬼除非先占据人的心，否则绝不能胜过人的事实。}\par

他们对此感到惊异，便直接来到\Person{安当}院长那里，揭露了他们攻击的程度、原因和密谋，并掩饰了他们的嫉妒，请求立刻成为基督徒。但当他询问他们发动攻击的那一天时，他声称那时他正被极痛苦的思维煎熬。藉此经验，\Person{圣安当}证实并确立了我们在昨天的会议上发表的意见，即魔鬼不可能进入任何人的心灵或身体，也没有能力压倒任何人的灵魂，除非它们先剥夺了它所有的神圣思想，使它空洞，远离属灵默想。但你们必须知道，不洁之神以两种方式服从于人：要么因天主的恩宠和权能而服从于信者的圣德，要么被罪人的祭献和某种咒语所俘获，并被他们奉为敬拜者而谄媚。\Person{法利塞}人也为这观念所误导，幻想甚至主救主也以这种方式命令魔鬼，于是说：\BibleQuote{祂是靠魔鬼之王贝耳则步驱魔。}这是根据他们所知的法则：他们自己的巫士和行巫术的人——藉呼求他的名并献上他所喜悦和满足的祭品——作为他的仆人，甚至有权管辖那顺从他的魔鬼。\par

% structure: p0061 source=h2 target=subsection reason=relative-heading-rank; base=h2

\subsection{第二十章\newline{}}

\Emph{关于\WorkTitle{创世纪}所载堕落天使与人的女儿结合的问题。}\par

\Person{革尔玛诺}说：既然刚才因天主的眷顾，在我们中间提出了\WorkTitle{创世纪}的一段经文，并愉快地提醒我们现在可以方便地询问我们一直渴望知道的一点，我们想知道关于那些所谓与人的女儿结合的堕落天使，我们应持什么看法；这种事是否能与一个精神性的本性字面地发生。此外，关于你刚才引用的福音中关于魔鬼的那段经文：\BibleQuote{因为他是说谎者，也是他的父亲，}我们同样愿意听一听，谁是\Quote{他的父亲}所指的。\par

% structure: p0064 source=h2 target=subsection reason=relative-heading-rank; base=h2

\subsection{第二十一章\newline{}}

\Emph{对所提问题的回答。}\par

\Person{Serenus}：您提出了两个并非不重要的问题。我将尽己所能，按照您提出的顺序予以答复。我们不可能相信属灵的存在能与女人有肉体的交合。但如果这曾确实发生过，何以现在不再发生？我们为何看不到类似的由魔鬼通过女人而非通过男人所生的人？尤其明显的是，他们喜爱情欲的玷污，若能办到，他们必定更愿通过自己而非通过男人来实现。正如《训道篇》所言：\BibleQuote{曾有过的，将来仍有。已做成的，将来仍要做。太阳之下绝无新事，以致有人能说：看，这是新的；因为在我们以前的世代早已有了。}\BibleRef{训 1:9}但所提出的问题可以这样解决：义人亚伯死后，为使整个人类不从一个邪恶的杀人犯而生，塞特被生下来代替他被杀之兄，不仅在后裔上，也在正义和良善上代替他。他的后裔仿效其父的良善，始终不与从恶人加因所出的亲属交往和共处，正如族谱的差异所清楚表明的，其中有载：\BibleQuote{亚当生塞特，塞特生厄诺士，厄诺士生刻南，但刻南生玛拉肋耳，但玛拉肋耳生雅勒特，雅勒特生哈诺客，哈诺客生默突舍拉，默突舍拉生拉默客，拉默客生诺厄。}\BibleRef{创 5:4}而加因的族谱则另列如下：\BibleQuote{加因生哈诺客，哈诺客生刻南，刻南生玛拉肋耳，玛拉肋耳生默突舍拉，默突舍拉生拉默客，拉默客生雅巴耳和犹巴耳。}\BibleRef{创 4:17}因此，从义人塞特后裔所出的支系始终与本族相交，长久保持其祖先的圣洁，未受邪恶后裔的亵渎和恶行沾染，那后裔植入了仿佛祖先传递的罪种。只要两支系之间仍有分离，塞特的后裔因其出自佳美的根，便因其圣洁而被称为\OriginalTerm{天主的使者}{angels of God}，或如有些抄本所称的\OriginalTerm{天主的儿子}{sons of God}；反之，其他人因其自身和祖先的邪恶及属世行为，被称为\OriginalTerm{人子}{children of men}。那时虽有这神圣有益的分离，但此后塞特的儿子们，即天主的儿子们，看见从加因支系所生的女儿们，便贪恋她们的美貌，娶她们为妻。这些妻子将她们父亲的邪恶教导给丈夫，立刻领他们背离了天生的圣洁和祖先的纯朴。以下言论对此说得很准确：\BibleQuote{我说过：你们是神，全都是至高者的儿子。但你们必要像人一样死去，像一位首领一样跌倒。}他们背离了祖先所传授的真正自然哲学研究。那第一个人，曾立即追溯全部自然的研究，能清楚掌握并将其确凿地传于后裔，因为他看见这世界的幼年，那时仍如柔软、悸动且未成形；他不仅有如此圆满的智慧，还有神圣灵感所赐的先知恩宠，以致尚在这世上未受教导时，就给一切生物命名；不仅知道所有野兽和毒蛇的狂怒与毒液，还辨别植物与树木的功效、石头的性质、以及他尚未经历的季节变化，以致他能说道：\BibleQuote{上主赐给了我认识万物的真实知识，使我得以知晓世界的构造和元素的效能，时间的开始、终结与中间，时节的运转和季节的变迁，年岁的循环和星辰的方位，生物的性情和野兽的凶猛，风的力量和人的思虑，植物的多样和根茎的功效——凡隐晦和显明的事我都学会了。}\BibleRef{智 7:17}塞特的后裔世代相传，从祖先继承了这全部自然的知识，只要他们与邪恶支系保持分离，便圣洁地领受并用于光荣天主和日常生活的需要。但一旦与邪恶世代混合，他们便在魔鬼的唆使下，将所无知学会的东西用于亵渎和有害的用途，并胆敢用它教导巫士、咒法和魔幻迷信的奇术，教其后代离弃对天主的神圣敬拜，转而尊崇崇拜元素、火或空中的魔鬼。至于我们所谈的这奇术知识何以未在洪水中灭亡，反而为后世所知，我认为，按此次讨论的契机，虽对所提问题的回答几乎无需此点，但仍应简要说明。据古代传说，诺厄的儿子含，曾受教于这些迷信和邪恶亵渎的技艺。他知道自己无法将任何这方面的手册带入与他善父圣兄同进的方舟。于是，他将这些邪术和亵渎器具铭刻于各种洪水无法毁灭的金属版和坚硬岩石上。洪水过后，他带着当初隐藏它们时同样的好奇去搜寻，从而将亵渎和永罪的温床传于子孙。这样，那普遍观念——人们相信天使将咒法和各种技艺传给人类——便确实应验了。从这些塞特之子与加因之女，如我们所说，生下了更恶的子孙，他们成为强悍的猎人、凶暴残忍的人，因其身体庞大、残暴邪恶而被称作\OriginalTerm{巨人}{giants}。这些人首先开始侵扰邻人，在人间实行掠夺，以抢掠为生，而不满足于辛劳和汗水。他们的恶行增至如此地步，以致世界只能通过洪水大灾得以净化。因此，当塞特之子在情欲的煽动下，违犯了那从世界之初由自然本能长期遵守的命令时，后来必须通过法律明文来恢复：\BibleQuote{不可将你的女儿嫁给他的儿子，也不可将他的女儿娶给你的儿子；因为他们必要引诱你们的心离开你们的天主，去事奉别的神，叩拜他们。}\par

% structure: p0067 source=h2 target=subsection reason=relative-heading-rank; base=h2

\subsection{第二十二章}

\Emph{一个异议：在法律禁止之前，如何能将舍特的后裔与加音的女儿的非法混杂归咎于他们？}\par

\Person{杰尔马努斯}说：如果那条命令曾赐给他们，那么因胆敢如此婚配而违反它的罪过就可以合理地归咎于他们。但是，既然那种分离的遵守尚未由任何规则确立，那么这种种族的混杂，既未受任何命令禁止，又如何能算作他们的过错呢？因为法律通常不会禁止已犯的罪，而是禁止将来的罪。\par

% structure: p0070 source=h2 target=subsection reason=relative-heading-rank; base=h2

\subsection{第二十三章}

\Emph{答：人从起初就因自然法而应受审判与惩罚。}\par

\Person{塞雷努斯}说：天主在创造人时，就将完整的法律知识自然种植在人里面；如果人按主的旨意，如同起初一样保持这知识，就不需要另外颁布一条后来以文字宣示的法律；因为若内在的药物仍被种植且充满活力，再提供外在的治疗就是多余的。但是，由于这内在的法律，如我们所说，被自由和犯罪的机会彻底腐蚀了，于是附加了梅瑟法律的严厉限制作为此（较早法律）的执行者和维护者，并借用圣经的表达，作为其辅助者，使人因惧怕即刻的惩罚，而不致完全丧失自然知识的美好。正如先知所说的：\Quote{祂颁布法律以帮助他们。} 宗徒也描述它被赐下作小孩子的启蒙师\BibleRef{迦 3:24}，因为它教导和保护他们，防止他们因纯粹健忘而偏离他们曾由自然之光所受的教导。因为完整的法律知识在最初受造时就已种植在人里面，这从以下事实清楚证明：我们知道在法律之前，甚至在洪水之前，所有圣人都遵守法律的诫命，而并未阅读法律的文字。因为\newline{}\newline{}亚伯尔如果没有由自然种植于他内的法律的教导，如何能在没有法律命令的情况下，就知道应将羊群中首生的和最好的脂肪祭献给天主呢？\BibleRef{创 4:4}\newline{}\newline{}诺厄如何能区分哪些动物是洁净的，哪些是不洁净的，\BibleRef{创 7:2} 而那时法律的命令尚未做出区分，除非他受了自然知识的教导？\newline{}\newline{}哈诺客从未从别人那里获得任何法律的亮光，他又从何处学会\BibleQuote{与天主同行}？\BibleRef{创 5:22}\newline{}\newline{}闪和耶斐特在哪里读到\Quote{不可揭露你父亲的裸体}，以致他们倒退着遮盖了父亲的羞耻？\newline{}\newline{}亚巴郎如何被教导要戒绝敌人献给他的战利品，以免接受任何劳苦的报偿，并将什一之物献给司祭默基瑟德，而这正是梅瑟法律所规定的？\newline{}\newline{}同样，亚巴郎和罗特又如何谦卑地向过路人和陌生人提供好客的服务和洗脚，而此时福音的命令尚未发出？\newline{}\newline{}约伯从哪里获得了如此热切的信德、如此纯洁的贞洁、如此谦逊、温和、怜悯和仁慈的知识，甚至我们现在在那些熟记福音的人身上也见不到这些？\newline{}\newline{}我们读到哪位圣人，在法律颁布之前没有遵守某条法律的诫命？有谁没有遵守这一条：\BibleQuote{以色列，你要听！上主我们的天主是唯一的上主。} \BibleRef{申 6:4} 有谁没有履行这一条：\BibleQuote{不可为你制造任何雕像，或天上、地上、地下水中之物的任何肖像？} 有谁没有遵守这一条：\BibleQuote{应孝敬你的父亲和你的母亲，} 或十诫中接下来的：\BibleQuote{不可杀人；不可奸淫；不可偷盗；不可作假见证；不可贪恋你近人的妻子。} \BibleRef{出 20:4} 以及其他许多事，在这些事上，他们不仅预见了法律的命令，甚至预见了福音的命令？\par

% structure: p0073 source=h2 target=subsection reason=relative-heading-rank; base=h2

\subsection{第二十四章}

\Emph{论在洪水前犯罪的人受罚是公义的。}\par

因此，我们看到，天主起初创造了完美的一切，如果万物都保持在祂所创造的状态和条件下，就不需要再给祂最初的安排添加任何东西——好像它短视而不完美似的。所以，对于在法律之前甚至在洪水之前就犯了罪的人，我们看到天主以公义的审判惩罚了他们，因为他们违背了自然法律，无可推诿地该受惩罚；我们也不应陷入那些不明此理之人的亵渎诽谤，他们以此贬低旧约的天主，诋毁我们的信仰，并嘲讽说：那么，为什么你们的天主愿意在数千年之后才颁布法律，而让如此漫长的时代没有任何法律呢？但若祂后来发现了更好的东西，那就说明在世界之初，祂的智慧是较低和较贫乏的，而后仿佛受经验教导才开始提供更好的东西，并修正和改进祂最初的安排。这种事情当然不可能发生在天主无限的预知上，异端者的疯狂愚昧也不能不犯严重亵渎而如此论断，正如\BibleBook{}{训道篇}所说：\BibleQuote{我已知道，天主从起初所造的一切言语将永远存续：不能加添，也不能删除，}\BibleRef{训 3:14}因此\BibleQuote{法律不是为义人设立的，乃是为不法和不服的，不虔敬和犯罪的，不圣洁和恋世俗的。}\BibleRef{弟前 1:9}因为他们拥有健全完整的自然法律系统植入内心，就不需要外加这成文的法律，就是作为自然法律辅助而赐下的法律。由此我们得出最清楚的推理：那成文的法律不必在起初就赐下（因为当自然法律仍然存在且未被完全违反时，这样做是不必要的），而且福音的完美也不能在法律被遵守之前就赐予。因为他们不能听从这句话：\BibleQuote{若有人打你的右脸，连另一面也转给他，}\BibleRef{玛 5:39}这些人不满足于以\OriginalTerm{同态复仇法}{lex talionis}的公平正义来报复所受的伤害，反而用致命的拳打脚踢和刀剑还击极轻微的触碰，并且为了一句真话就想杀死打他们的人。也不能对他们说：\Quote{爱你们的仇敌，}在他们当中，爱自己的朋友被视为一件大事和最重要的事，却避开仇敌，仅以仇恨与他们分离，而不急于压迫和杀害他们。\par

% structure: p0076 source=h2 target=subsection reason=relative-heading-rank; base=h2

\subsection{第二十五章}

\Emph{如何理解福音中论及魔鬼的话，即\Quote{他是说谎者，也是说谎者之父。}}\par

但至于这件事，即那有关魔鬼并使你们不安的，就是\BibleQuote{他是说谎者，并且是说谎者之父，}\BibleRef{若 8:44}——因为似乎主曾宣称他和他的父亲都是说谎者，——即便粗略地这样想，也是十分可笑的。因为我们不久前说过，灵魂不能生出灵魂，正如灵魂不能产生灵魂一样，虽然我们不怀疑肉身的结合是由人的种子形成的，正如宗徒在两种实体——即肉身和灵魂——的情况下清楚地区分了谁应当被归于谁作为其创造者，并且说：\BibleQuote{此外，我们曾有过我们肉身的父亲管教我们，尚且尊敬他们；难道我们不是更要服从灵魂的父亲而生活吗？}\BibleRef{希 12:9}有什么比这个区分更清楚的呢？他（宗徒）确定人是肉身之父，但始终教导惟有天主是灵魂之父。虽然即使在这个身体的实际结合中，也必须只将辅助的服务归给人，但将其形成的主要部分归给万有的创造者天主，正如\Person{达味}说：\BibleQuote{祢的双手造了我，并塑造了我。}以及真福\Person{约伯}说：\BibleQuote{祢不是将我倒出如奶，使我凝结如酪吗？祢以骨和筋把我连络起来；}\BibleRef{约伯 10:10}以及上主对\Person{耶肋米亚}说：\BibleQuote{我尚未在母腹中形成你，我已认识了你。}\BibleRef{耶 1:5}但是\BibleBook{}{训道篇}通过考察每种实体的来源和开端，以及考虑它们各自趋向的终结，非常清楚而精确地总结了两种实体的本质及其起始，并且也决定了这个身体和灵魂的分离，论述如下：\BibleQuote{在尘土归于它所出的地，灵魂归于赐它的天主之前。}\BibleRef{训 12:7}还有什么比这说得更明白呢？它宣称肉身的物质——它称之为尘土，因为它来自人的种子，似乎是由人的协助撒下的——既从土而出，就必须再归回土；同时它指出，那并非由两性交合而生，而是唯独以特殊方式属于天主的灵魂，归回它的创造者。这一点也清楚地隐含在天主的气息中，\Person{亚当}就是藉此初次获得了生命。因此，从这些段落我们清楚推知，除天主以外，没有人可被称为灵魂之父，祂随己意从无中创造灵魂，而人只能被称为我们肉身之父。所以，魔鬼既然被造为灵魂或天使，且是善的，除天主创造者外，他没有其他父亲。但当他因骄傲而自高自大，心中说：\BibleQuote{我要升到云层高处，我要与至高者相似，}\BibleRef{依 14:14}他就成了说谎者，并且\BibleQuote{不居于真理中；}\BibleRef{若 8:44}但从自己邪恶的宝库中生出谎言，因此不仅成了说谎者，而且成了谎言之父；他藉此向人许诺神性，说：\BibleQuote{你们将如同神一样，}\BibleRef{创 3:5}他不住在真理中，反而从起初就成了杀人者，既使\Person{亚当}陷入必死的状态，又藉其兄弟之手杀害\Person{亚伯尔}，由于他的怂恿。但是黎明的临近即将结束我们的讨论，它已经几乎占用了整整两个夜晚；我们简短而朴素的言语已将这次会谈的小舟从问题的大海带进了安全的静默港；在这深海中，确实，随着天主圣神的气息驱使我们更深入，总是展现出更广阔无垠的空间，超越我们眼目所及，并且，正如\Person{撒罗满}说：\BibleQuote{它将离我们比从前更远，且是一个深渊；谁能找出它呢？}\BibleRef{训 7:25}因此，让我们祈求主，使祂的敬畏和祂永不失败的爱，能坚定地存留在我们内，并使我们在一切事上有智慧，永远保护我们免受魔鬼的箭矢伤害。因为有了这些护卫，没有人可能陷入死亡的罗网。但在成全者与不成全者之间有这样一个区别：在前者，爱情是坚定的，可以说是更成熟的，且更持久存在，因此使他们更坚定且更容易地持守圣德；而在后者，爱情的地位更软弱，且更容易冷却，因而更快、更频繁地让他们被罪恶的罗网缠住。当我们听到这些时，这次会谈的话如此点燃了我们，以致当我们离开老人的小屋时，我们比以前初来时怀着更炽热的灵魂热忱，渴望实现他的教导。\par

\SourceRecord{Translated by C.S. Gibson.；Nicene and Post-Nicene Fathers, Second Series；Vol. 11.；Edited by Philip Schaff and Henry Wace.；Buffalo, NY: Christian Literature Publishing Co.,；1894.；<http://www.newadvent.org/fathers/350808.htm>.}

% structure: p0001 source=h1 target=section reason=page-title

\section{第九讲}

\Emph{\Emph{\Person{圣若望·卡西安}}}\par

院长\Person{依撒格}的第一场讲道。论祈祷。\par

% structure: p0004 source=h2 target=subsection reason=relative-heading-rank; base=h2

\subsection{第一章}

\Emph{讲道引言}\par

在第二部\WorkTitle{隐修制度}中，论及持续不断祈求于祈祷的许诺，将藉\Person{主}的助佑，由这位长者的讲道——即院长\Person{依撒格}——来实现；待这些讲道陈述完毕后，我想我已满足了真福教宗\Person{卡斯托}的嘱咐，以及您，真福教宗\Person{良定}和神圣的弟兄\Person{赫拉迪乌斯}的愿望；本书前部篇幅过长，虽已尽力将所述内容压缩为简短论述，并略去甚多要点，仍超出原定篇幅。因为从对各种规诫的完整论述开始——我们为简练起见已将其删节——最后，真福\Person{依撒格}说了以下的话。\par

% structure: p0007 source=h2 target=subsection reason=relative-heading-rank; base=h2

\subsection{第二章}

\Emph{院长\Person{依撒格}论祈祷本质的言论}\par

每位隐修士的目标及其内心的成全，在于持续不断、无间断地坚持祈祷，并尽可能在人性软弱允许的范围内，力求获得不可动摇的心灵宁静与永久的纯洁；为达此目的，我们不知疲倦、恒常不懈地操练一切身体的劳作和心灵的痛悔。此二者之间有一种相互的、不可分离的联合。因为正如一切德行的建筑之冠是祈祷的成全，同样，除非万事万物因此而联合、固结于此冠冕之下，它就不可能持久巩固。我们所谈论的祈祷之持久与恒常的安宁，若无它们便不能获得或完成；同样，奠基于祈祷的诸德行，若不坚持祈祷，也不能完全获得。因此，我们既不能恰当地论述祈祷的效果，也无法通过快速的言谈深入其最终目标——这目标是通过操练一切德行而获得的——除非首先逐一列举并讨论那些为祈祷之故而必须弃绝或确保的一切事物，并如福音中的比喻教导的\BibleRef{路 14:28}，凡事关建造那属灵的、至高的塔的，都事先计算并仔细考虑。然而，这些准备若未先清除一切过犯，挖掘已腐朽的、死去的私欲之垃圾，并在我们胸怀中坚实而（可说是）活的土壤上，或更确切地说，在那福音的磐石上\BibleRef{路 6:48}奠定单纯与谦卑的坚固根基，便毫无用处，也无法使至高的德行之巅恰当地置于其上；如此建造，这属灵德行之塔必将升起，屹立不动，并在其稳固的保证下，被提升至天空的至高之处。因为若它奠基于这样的根基，那么即使激烈的私欲风暴冲击它，即使迫害的洪流如撞锤般击打它，即使属灵仇敌的猛烈暴风雨猛冲并攻击它，它不仅不会倾覆，甚至连丝毫损伤也不会受到。\par

% structure: p0010 source=h2 target=subsection reason=relative-heading-rank; base=h2

\subsection{第三章}

\Emph{如何获得纯洁真诚的祈祷}\par

因此，为使祈祷能以应有的热忱和纯洁献上，我们必须切实遵守以下规则。首先，必须完全消除一切关于肉身事物的焦虑；其次，不仅不能有任何事务的挂虑，甚至不能再想起它们；同样，必须放下一切诽谤、虚浮不断的闲聊和戏谑；特别是愤怒和令人不安的愁苦必须完全毁灭，致命的肉欲和贪欲之劣根必须连根拔除。这样，当这些以及类似在世人眼前也可见的过失被完全清除和剪除后，当我们所说的那种起初的净化与洗涤——由纯朴单纯与清白无瑕所带来——完成后，就必须先奠定深沉的谦卑之稳固根基，它能支撑那直达天际的高塔；其次，必须在它们之上建造德行之属灵结构，使灵魂远离一切交谈和飘忽的思绪，从而逐步开始上升到对\Person{天主}的默观和属灵洞见。因为无论我们的心灵在祈祷前思虑什么，在祈祷时，借由记忆的活动，它必定会呈现给我们。因此，我们希望在祈祷时发现自己处于何种状态，就应当在祈祷前准备好自己。因为祈祷时的心灵由先前的状态所形成；当我们致力于祈祷时，相同行为、言语和思想的影像会在眼前舞动，使我们或如先前那样愤怒，或忧郁，或回忆起先前的欲望和事务，或为某个愚蠢的笑话而颤抖着发出愚蠢的笑声（我羞于提及），或对某个行动微笑，或回到先前的谈话中。因此，若我们不想在祈祷时有任何事物纠缠我们，就应在祈祷前小心地将它排除出心灵的圣所，如此我们才能实现宗徒的训令：\Quote{不断祈祷；}和\Quote{在每处举起圣洁的手，不发怒，不争论。}否则，除非我们的心灵从一切罪恶的污秽中净化，并归向德行如同归向它自然的善，以不断默观\Person{全能天主}为食粮，否则我们将无法履行那命令。\par

% structure: p0013 source=h2 target=subsection reason=relative-heading-rank; base=h2

\subsection{第四章}

\Emph{论灵魂的轻灵，可比作羽翼}\par

灵魂的本性，若比作一根极轻的羽毛或极轻的翅膀，并非不恰当：它若没有因外界落在它上面的任何湿气而受损或受影响，便几乎自然而然地藉其本性的轻飘和微风的助力，升到天空的高处；但若因落在它上面并渗入它里面的湿气而加重，那么它不仅不会因其本性的轻飘而飘向任何空中的飞行，反而会因所吸收的湿气之重量而被拖到地的深处。同样，我们的灵魂，若没有因触及它的过错和世俗的挂虑而加重，或因有害情欲的湿气而受损，便会因其自身纯洁的本性之祝福而上升，并由属灵默想的微风带到高处；离开低下和属地的事物，被提升到属天和不可见的事物。为此，主的命令很好地警告了我们：\Quote{你们要谨慎，免得你们的心因贪食、醉酒和世俗的挂虑而沉重。}\BibleRef{路 21:34}因此，如果我们希望我们的祈祷不仅能达到天空，而且能达到天空之上，我们就要小心，使我们的灵魂从一切属地的过错中净化，从一切污点中涤净，恢复到其本性的轻飘，这样我们的祈祷就能不受任何罪恶重量的阻碍而上升到天主。\par

% structure: p0016 source=h2 target=subsection reason=relative-heading-rank; base=h2

\subsection{第五章}

\Emph{论灵魂被沉重的方式。}\par

\Emph{论灵魂被沉重的方式。}但我们应当注意主所指出的灵魂被沉重的方式：因为祂没有提奸淫、或淫乱、或谋杀、或亵渎、或抢劫，这些人都知道是致死的、该受惩罚的；而是贪食、醉酒和世俗的挂虑或焦虑——世俗之人不但不避免这些，不认为该受惩罚，反而（我羞于说）有些自称隐修士的人竟把这些职责当作无害或有用的而缠绕其中。虽然这三样东西，如果放任，会沉重灵魂，使它远离天主，拖向属地的事物，但要避免它们却很容易，尤其是对于我们这些因如此巨大的距离而与世俗的一切交往隔绝、并且在任何情况下都与那些可见的挂虑、醉酒和贪食无关的人。但还有另一种贪食，同样危险；一种属灵的醉酒，更难避免；一种世俗的挂虑和焦虑，常常在我们彻底舍弃一切财物、戒绝酒宴、甚至独居时也缠住我们——关于这样的醉酒，先知说：\Quote{醒来，你们这些醉酒却不是因酒的；}\BibleRef{岳 1:5}另一位说：\Quote{你们要惊骇、奇怪、蹒跚：醉酒却不是因酒；动摇，却不是因醉酒。}\BibleRef{依 29:9}因此，这醉酒的酒必然就是先知所说的\Quote{龙之狂怒}；至于这酒出自什么根，你可以听到：\Quote{他们的葡萄树是出自索多玛的葡萄园，}他说，\Quote{他们的枝条来自哈摩辣。}你也要知道那葡萄树的果实和那枝条的种子吗？\Quote{他们的葡萄是苦胆的葡萄，他们的串串是苦的。}\BibleRef{申 32:32}因为除非我们完全净化一切过错，戒除一切情欲的饱足，我们的心，即使没有醉酒和宴席的过度，也会被一种更危险的醉酒和贪食所沉重。因为世俗的挂虑有时会落在我们这些与世俗行为无涉的人身上，这一点根据长老们的规矩已清楚表明，他们规定：任何超出日常食物需要和肉身不可避免的需要的东西，都属于世俗的挂虑和焦虑，例如，如果一天赚一分钱足以满足身体需要，我们却试图通过更长时间的辛劳和工作来赚两分或三分钱；如果两件外衣的遮盖足以在日夜使用，我们却设法拥有三件或四件；如果一间有一两个房间的小屋就足够了，我们却因世俗的野心和伟大的骄傲而建造四五个房间，而且装饰华丽，比需要的更大——这样，每当我们能够，就显出了世俗情欲的激情。\par

% structure: p0019 source=h2 target=subsection reason=relative-heading-rank; base=h2

\subsection{第六章}

\Emph{论一位长老所见的关于一位不安工作的弟兄的异象。}\par

此事非因魔鬼的煽动而不能成就，我们由最确凿的证据得知：一位极为尊贵的长老经过某位罹患此心智疾病（我们所论及的）的弟兄的小室时，见他不懈地劳碌于日常修建与修补无用之物，从远处看到他正用重锤敲打极坚硬的石头，且见一埃塞俄比亚人立于其旁，双手交握，与他一同击锤，并以火热的刺激催逼他勤于工作。长老惊愕于这凶猛魔鬼的力量与如此幻象的欺骗性，久久伫立。当那弟兄疲惫欲歇、想停止劳作时，却被那灵的刺激鼓动，重又拿起锤子，不停歇地致力于已开始的工作，以致在它的催逼下不觉如此劳苦之害。最后，那长老厌恶魔鬼如此可怖的迷惑，转向弟兄的小室问候他，问：\Quote{兄弟，你在做什么工作？}他回答：\Quote{我们在敲这块极其坚硬的石头，几乎敲不碎。}长老便说：\Quote{你说\InnerQuote{我们}能，这话不错，因为你敲的时候不是独自一人，还有另一个人在你旁边，你看不见他，他与你同立，与其说是帮助你，不如说是全力催逼你。}因此，摆脱世俗虚荣之病的心，不是仅凭戒绝那些即便想做也无法实行的事务，也不是轻看那些若追求便会在属神的人与世俗人中名列前茅的事，而是以不可动摇的坚毅之心拒绝一切助长我们权势、看似披着正当外衣的事物。事实上，这些看似微不足道、无足轻重，且被我们同圣召的人视为无妨之事，其性质对灵魂的影响，并不亚于那些按状况常使世俗人迷醉、且不容隐修士摒弃尘世污染而渴慕天主（其注意力当常凝注于祂）的重大事物；因为对隐修士而言，即使与至高之善稍有分离，亦当视为现世的死亡与最危险的毁灭。当灵魂被安置于如此平安的状态，摆脱一切肉欲的网罗，心灵的意向稳固地定于那唯一的至高之善时，他便将履行这宗徒的训令：\BibleQuote{你们要不断祈祷}；并：\BibleQuote{在每一个地方举起圣洁的双手，不发怒，不争论}。因为当灵魂的思想被此纯洁（若可如此说）所占据，并从尘世状态被重塑为属神与天使的样式时，无论它接受什么、着手什么、做什么，祈祷都将完全纯洁而真诚。\par

% structure: p0022 source=h2 target=subsection reason=relative-heading-rank; base=h2

\subsection{第七章}

\Emph{一个问题：为何保持善念比产生善念更难。}\par

格曼努斯说：但愿我们能以同样方式、同样轻易地持久保有那些属神的思想，就像我们一般萌生其萌芽一样！因为当这些思想通过回忆圣经或记忆某些属神行动或默观天上奥迹而在心中萌生时，它们却转瞬即逝，不知不觉地飞逝无踪。当我们的灵魂发现其他引发属神情感的机会时，又有不同的思想蜂拥而至，我们所抓住的那些思想便散失轻逸而去，以致心神毫无恒心，凭己力亦不能牢牢把握圣洁的思想，即使一时似乎持有它们，也必被视为偶然萌生而非有意为之。因为若它们不持久存留，我们怎能认为其兴起应归因于我们自己的意志？但为避免因思考此问题而偏离已开始讨论的计划，或再推迟对祈祷本质的承诺解释，我们将此留待适当时候，并请立即告知祈祷的特性，尤其蒙福的宗徒劝勉我们不可停止祈祷，说：\BibleQuote{你们要不断祈祷}。所以我们希望先被教导祈祷的特性，即祈祷应当如何被\Emph{常常}举起，然后我们如何确保这一点，无论是什么，并不断地实行它。因为此事不能由任何轻浮的心意达成，日常经验与四位圣人的解释都向我们指明，正如您所规定的：隐修士的目标与一切成全的顶峰在于祈祷的完满。\par

% structure: p0025 source=h2 target=subsection reason=relative-heading-rank; base=h2

\subsection{第八章}

\Emph{论祈祷的不同特性。}\par

\Person{Isaac}说：我认为，若没有极其纯净的心灵和灵魂以及圣神的照亮，便无法领会各种祈祷。因为在一灵魂内——或可说在众灵魂内——可能产生的状况和性情有多少，祈祷的种类就有多少。因此，尽管我们知道自己因心灵的迟钝无法看见所有种类的祈祷，但我们仍将尽力按照我们浅薄的经历所能及的，以某种秩序来讲述它们。因为根据每一灵魂所达到的纯净程度，以及因所遇之事而沉沦或自身努力而更新的状态，其祈祷本身每时每刻都在变化；因此显而易见，没有人能始终献上相同的祈祷。因为每个人在精神振奋时以一种方式祈祷，在悲伤或失望的重压下以另一种方式祈祷，在因灵性成就而振作时又以另一种方式祈祷，在受攻击的重担下沮丧时又以另一种方式祈祷，在祈求罪过的赦免时以一种方式祈祷，在祈求获得恩宠或某种德行或为消灭某种罪而祈祷时又以另一种方式祈祷，在因地狱的念头和对将来审判的恐惧而心痛时以一种方式祈祷，在对未来美善的希望和渴望中炽热时又以另一种方式祈祷，在处理事务和危险时以一种方式祈祷，在平安与稳妥中又以另一种方式祈祷，在接受天上奥秘的启示而受光照时以一种方式祈祷，在因德行的贫乏和情感上的干枯而沮丧时又以另一种方式祈祷。\par

% structure: p0028 source=h2 target=subsection reason=relative-heading-rank; base=h2

\subsection{第九章}

\Emph{论祈祷的四重性质。}\par

因此，当我们对于祈祷的性质作了这样的规定——虽未如主题要求的那样详尽，却已如时间所允许，且至少如我们浅薄的才能所及，心灵的迟钝所容——一个更大的困难现在等待着我们；即逐一阐述各种不同的祈祷，宗徒将它们分为四重，如下所说：\Quote{所以我最劝勉的是，首先要为一切人、为君王及一切有权位者，恳求、祈祷、转求、感恩。} \BibleRef{弟前 2:1} 对于宗徒的这种划分绝非随意所作，我们丝毫不能怀疑。首先，我们必须研究什么是恳求、什么是祈祷、什么是转求、什么是感恩。其次，我们要追问，行祈祷者是否应当同时进行这四种，即是否在每次祈祷中都把它们结合在一起——或者它们应当轮流并逐一献上，例如，有时献上恳求，有时祈祷，有时转求，有时感恩；或者，是否一个人向天主献上恳求，另一个人献上祈祷，第三个人献上转求，第四个人献上感恩，这与每一灵魂因着志向的诚挚而长进的年龄度量相符。\par

% structure: p0031 source=h2 target=subsection reason=relative-heading-rank; base=h2

\subsection{第十章}

\Emph{论就祈祷的性质所定的各种顺序。}\par

因此，首先我们必须考虑这些名称和词语本身的含义，讨论祈祷、恳求和转求之间的区别；然后，同样地，我们要研究它们是应当分开还是全部一起献上；第三，要考察宗徒权威这样排列的特定顺序是否还向听者传达了一些什么，还是仅仅需要作出区分，并且应该认为他是这样随意排列的：在我看来，这是完全荒谬的。因为不可相信圣神会通过宗徒随意或无理由地说出任何事。因此，我们将按主所赐予我们的，按照我们开始的顺序来考虑它们。\par

% structure: p0034 source=h2 target=subsection reason=relative-heading-rank; base=h2

\subsection{第十一章}

\Emph{论恳求。}\par

\Quote{所以我最劝勉的是，首先要作恳求。} 恳求是一种关于罪的哀求或祈求，懊悔自己当前或过去行为的人在此祈求宽恕。\par

% structure: p0037 source=h2 target=subsection reason=relative-heading-rank; base=h2

\subsection{第十二章}

\Emph{论祈祷。}\par

祈祷是我们向天主奉献或许愿的那些事，希腊人称之为\OriginalTerm{誓愿}{\GreekText{εὐκή}}，即一个誓愿。因为在希腊文\Quote{\OriginalTerm{我将向主还我的誓愿}{\GreekText{ἰὰς} \GreekText{ἐυκάς} \GreekText{μου} \GreekText{τῶ} \GreekText{κυρίῶ} \GreekText{ἀποδώσω}}}之处，拉丁文读作：\Quote{我将向主还我的誓愿}；按字面精确含义可如此表达：\Quote{\BibleQuote{我将向主还我的祈祷。}}而我们在\WorkTitle{训道篇}中读到：\BibleQuote{你向主许了愿，不可迟延不还}，希腊文同样写作\Quote{\OriginalTerm{你若向主许愿}{\GreekText{ἐάν} \GreekText{ἐύξῃ} \GreekText{ἐυχὴν} \GreekText{τῶ} \GreekText{κυρίῶ}}}，即\BibleQuote{你若向主祈祷，不可迟延不还}。\BibleRef{训 5:3} 我们每人将以这样的方式履行它。当我们弃绝此世，并许诺在一切世俗行为及此世生命上死去，以全心服事主时，我们是在祈祷。当我们许诺轻视世俗的尊荣，蔑视地上的财富，以全心的忧伤和心灵的谦卑紧紧依随主时，我们是在祈祷。当我们许诺永远保持身体最完全的纯洁和坚定的忍耐，或者当我们发誓要从心中完全根除愤怒或致死的忧伤之根时，我们是在祈祷。如果我们因懈怠而软弱，回到从前的罪，未能做到这一点，我们就将因我们的祈祷和誓愿而有罪，并且这些话将适用于我们：\Quote{不更好地不祈祷，就是不祈祷而不还}，按希腊文可译为：\BibleQuote{不更好地不祈祷，祈祷而不还。}\par

% structure: p0040 source=h2 target=subsection reason=relative-heading-rank; base=h2

\subsection{第十三章}

\Emph{论转求。}\par

第三位的是转求，我们在充满热忱时也常为他人献上转求，或为我们所爱的人，或为全世界的和平，并用宗徒自己的话说，我们为\BibleQuote{一切人、为君王及一切有权位者}祈祷。\BibleRef{弟前 2:1}\par

% structure: p0043 source=h2 target=subsection reason=relative-heading-rank; base=h2

\subsection{第十四章}

\Emph{论感恩。}\par

然后，第四种是感恩，心灵在不可言喻的激动中向天主献上感恩，或是回忆天主过去的恩惠，或是默观祂现今的恩惠，或是展望祂为爱祂之人所预备的将来的伟大恩惠。为此目的，我们有时也会倾注更丰盛的祈祷，当我们以纯洁的眼凝视将来为圣人们所积存的赏报时，我们的心神被激励，以无边的喜乐向天主献上不可言喻的感恩。\par

% structure: p0046 source=h2 target=subsection reason=relative-heading-rank; base=h2

\subsection{第十五章}

\Emph{这四种祈祷是否必须每个人同时全部献上，还是分别轮流献上。}\par

在这四种之中，虽然有时会出现更丰盛、更圆满的祈祷（因为我们知道，从因悔罪而产生的恳求，从基于纯洁良心而献上祭献和履行誓愿的信赖所发出的祈祷，从爱德热情而生的代祷，以及从默观天主的恩赐、伟大和良善所生的感恩中，常常产生最热切、最炽热的祈祷，以至于显然，我们所说的这几种祈祷对所有人都是有用且必要的，以致在同一个人身上，他变化的情感会发出纯洁而热切的祈求，时而恳求，时而祈祷，时而代祷），然而，第一种似乎特别属于初学者，他们仍受罪恶刺痛和回忆的困扰；第二种属于那些在灵性进步和追求德行中已达到某种崇高心灵的人；第三种属于那些以行动完成誓愿，并因考虑他人的软弱和自身爱德的热情而激发为他人代祷的人；第四种属于那些已从良心中拔除有罪荆棘，因而无忧无虑，能以纯洁之心默观天主过去的恩赐、现今的恩惠，以及为未来所预备的慈悲，从而以炽热之心奔向那种人嘴无法容纳或表达的炽热祈祷。然而，有时那迈向完全纯洁境界、并已开始在其中稳固的心灵，会同时涵盖所有这些，如同某种不可理解、吞噬一切的火焰，贯穿它们，向天主献上最纯洁力量的无言祈祷，这祈祷由圣神亲自以无可言喻的叹息代祷，连我们自己也不明白，就在那一刻，在恳求中如此伟大而不可言喻地倾泻，以致无法用口说出，甚至日后也无法被心灵回忆。因此，无论一个人处于何种程度，他有时都会献上纯洁而虔诚的祈祷；即使在那与回忆未来审判相关的第一和卑微阶段，那仍处于恐怖惩罚和审判恐惧之下的人，也因当时的悲痛而深受打击，以至于他从恳求的丰盛中所获得的灵性敏锐度，并不亚于那因心灵纯洁而默观并思考天主的恩惠、被不可言喻的喜乐和欢愉所淹没的人。因为，正如主亲自所说，那知道自己被赦免得更多的人，开始爱得更多。\BibleRef{路 7:47}\par

% structure: p0049 source=h2 target=subsection reason=relative-heading-rank; base=h2

\subsection{第十六章}

\Emph{我们应致力于何种祈祷。}\par

然而，我们在生活中前进并修德时，应更致力于那些从默观未来美善或从爱德热情而发出的祈祷，或者，至少更谦卑地、按照初学者的尺度，为获得某些德行或消除某些过失而发出的祈祷。否则，除非我们的心灵通过那些代祷的常规次序逐渐攀升，我们不可能达到前述那些更崇高的恳求方式。\par

% structure: p0052 source=h2 target=subsection reason=relative-heading-rank; base=h2

\subsection{第十七章}

\Emph{主如何创立了四种恳求。}\par

主自己以祂的榜样恩赐我们开始这四种恳求，为叫祂也在此事上成就那关于祂所说的话：\Quote{耶稣开头所行所教训的一切事。}\BibleRef{宗 1:1} 因为祂在说\Quote{父啊！祢若愿意，请免去我这杯罢！}时，便使用了\Emph{恳求}的类别；或者那在圣咏中以祂的身份所咏唱的：\Quote{我的天主，我的天主，求祢眷顾我，为什么舍弃了我？}以及类似者。这是\Emph{祈祷}，当祂说：\Quote{我在地上已经光荣了祢，我完成了祢所托付我的工作，}以及这句：\Quote{我为他们的缘故，祝圣我自己，为叫他们也在真理中被祝圣。}\BibleRef{若 17:4} 这是\Emph{转求}，当祂说：\Quote{父啊！祢赐给我的人，我愿他们在那里，同我在一起，为叫他们看见祢赐给我的光荣；}或者至少当祂说：\Quote{父啊！宽恕他们，因为他们不知道他们做的是什么。}这是\Emph{感恩}，当祂说：\Quote{父啊！天地的主宰，我称谢祢，因为祢将这些事瞒住了智慧及明达的人，而启示给了小孩子。是的，父啊！祢原来喜欢这样。}或者至少当祂说：\Quote{父啊！我感谢祢，因为祢俯听了我。我本来知道祢常常俯听我。}不过，我们的主虽然按我们所知的模式区分了这四种祈祷，使之可以分别地、一个一个地献上，然而，祂又藉着祂自己的榜样，在那篇我们在圣若望福音末尾读到祂如此丰满地献上的祈祷中，显明它们全部可以在同一时间被包含在一个完美的祈祷中。从这篇祈祷的言词（因为太长不能全部引述），细心的研究者可以从经文的次序发现这是如此。宗徒在写给斐理伯人的书信中也表达了同样的意思，以略为不同的顺序列出这四种祈祷，并表明它们有时应当在单独一次祈祷的热忱中一起献上，如此说道：\Quote{一切事上，藉着祈祷和恳求，以感恩，将你们的请求呈报给天主。}\BibleRef{斐 4:6} 为此，他特别要我们明白：在祈祷和恳求中，感恩应当与我们的请求混合在一起。\par

% structure: p0055 source=h2 target=subsection reason=relative-heading-rank; base=h2

\subsection{第十八章}

\Emph{论主祷文。}\par

因此，在这几种不同的恳求之后，紧接着有一种更崇高、更超越的境界，那是经由唯独默观天主和炽热的爱情而产生的；藉此境界，心灵将自己投入并抛向对祂的爱，以自己特有的孝爱之情，最亲昵地称呼天主为自己的父亲。主祷文的格式教导我们应当热切寻求这境界，说：\Quote{我们的天父。}当我们亲口承认宇宙的天主和主是我们的父亲时，我们当即宣告自己已从奴隶的地位被召回成为义子的身份；接着加上\Quote{祢在天上，}为叫我们以极致的恐惧逃避所有滞留于现世生命（我们在此世如同过客），以及那使我们与父亲相隔甚远的一切，反倒以全力急切地赶赴我们承认父亲所居的国度，并且不容许任何这类事物，使我们不配于这身份和如此义子的尊严，以致作为父亲继承的耻辱而被剥夺，并招致祂公义和严厉的愤怒。及至我们达到这种儿子身份的状态和情况时，我们将立刻被善生儿子应有的热忱所燃烧，以致我们所有精力都用于推进父亲的光荣，而不求自己的利益，向祂说：\Quote{愿祢的名被尊为圣。}以此证明我们的渴望和喜乐是祂的光荣，效法那一位所说的话：\Quote{由自己发言的，寻求自己的光荣；但那寻求派遣祂者光荣的，祂是真实的，在祂内没有不义。}\BibleRef{若 7:18} 最后，那充满了这情感的\OriginalTerm{特选的器皿}{vas electionis}希望自己能被诅咒而脱离基督\BibleRef{罗 9:3}，只愿属于祂的子民能增多和繁衍，整个以色列民族的得救能归功于祂父亲的光荣；因为他知道没有人会为了生命而灭亡，所以他可以完全确信地愿为基督而死。他又说：\Quote{当我们软弱而你们强壮时，我们喜欢。}\BibleRef{格后 13:9} 那特选的器皿愿为基督的光荣、自己弟兄的归化及民族的殊恩而与基督隔绝，这有什么奇怪呢？因为米该亚先知愿意自己成为说谎者，并失去圣神的灵感，只愿犹太百姓能逃脱他所预言的那些灾祸和被掳，说：\Quote{但愿我不是有圣神的人，我宁愿说谎言。}\BibleRef{米 2:11}——且不提那立法者（梅瑟）的心愿，他并不拒绝与注定死亡的弟兄一同死亡，说：\Quote{我求祢，上主！这百姓犯了大罪；求祢赦免他们的过犯，不然，求祢从祢所写的册子上抹去我的名字。}\BibleRef{出 32:31} 不过，当说\Quote{愿祢的名被尊为圣}时，也很可以这样理解：\Quote{天主的圣化就是我们的成全。}所以，当我们向祂说\Quote{愿祢的名被尊为圣}时，我们就是说：父啊，求祢使我们能如此，以致我们能领悟并懂得祢的圣化是什么，或者至少，祢在我们灵性的交往中被看见受尊崇。这在我们身上有效地实现，是当\Quote{人们看见你们的善行，而光荣你们的在天之父。}\BibleRef{玛 5:16}\par

% structure: p0058 source=h2 target=subsection reason=relative-heading-rank; base=h2

\subsection{第十九章}

\Emph{论\Quote{愿祢的国来临}这一句。}\par

纯洁的心灵的第二个恳求是渴望其父的王国立刻来临；即，要么是基督日复一日在圣人中为王（当魔鬼的统治因可耻的罪恶被消灭而逐出我们心中，天主开始以德行的馨香掌管我们，当淫乱被克服，爱德连同宁静在我们心中为王，当愤怒被征服；当骄傲被践踏时，谦卑得胜），要么是在适当时刻应许给所有成全者和所有天主子民的那个，那时基督会对他们说：\BibleQuote{我父所祝福的，你们来吧！承受自创世以来为你们预备的国度吧！}\BibleRef{玛 25:34}（而心灵），可以说是以坚定稳定的目光，渴慕它并对祂说：\BibleQuote{愿祢的国来临。}因为良心见证它知道当祂显现时，它将立即分享祂的命运。因为没有罪人敢于这样说或盼望这事，因为没有人愿意面对审判官的法庭，他知道在祂来临时，他将立即得到的不是奖赏或功绩的赏报，而只是惩罚。\par

% structure: p0061 source=h2 target=subsection reason=relative-heading-rank; base=h2

\subsection{第二十章}

\Emph{论\BibleQuote{愿祢的旨意奉行}这一句。}\par

第三项恳求是子女的：\BibleQuote{愿祢的旨意奉行在人间，如同在天上。}没有比希望地上的事物与天上的事物相等更大的祈祷了：因为说\BibleQuote{愿祢的旨意奉行在人间，如同在天上}，除了祈求人能像天使一样，并且如天主的旨意在天上永远是满全的那样，所有地上的人也能行祂的旨意而不是自己的旨意之外，还有什么呢？这没有人能发自内心地这样说，除非是相信天主为我们命定一切所见之事，无论幸运或不幸运，都是为我们的益处，而且祂比我们自己更加关心和眷顾我们的益处和救恩的人。或者至少可以这样理解：天主的旨意是所有人的得救，正如圣保禄的话：\BibleQuote{祂愿意所有的人都得救，并得以认识真理。}\BibleRef{弟前 2:4}关于这个旨意，依撒意亚先知以天主圣父的位格说：\BibleQuote{我的一切旨意都要成就。}\BibleRef{依 46:10}所以当我们说\BibleQuote{愿祢的旨意奉行在人间，如同在天上}时，我们求的是这个：如同天上的人，同样地上所有的人都因认识祢而得救，圣父啊。\par

% structure: p0064 source=h2 target=subsection reason=relative-heading-rank; base=h2

\subsection{第二十一章}

\Emph{论我们的超本质的或每日的食粮。}\par

接着：\BibleQuote{求祢今天赐给我们日用的食粮，这食粮是\OriginalTerm{超本质的}{\GreekText{ἐπιούσιον}}}，另一位福音作者称之为\Quote{日常的}。前者表明其高贵和本质的性质，因此它超乎万有之上，其伟大和神圣的崇高超越一切受造物；而后者则表明其用途和价值的目地。因为当它说\Quote{日常的}时，表明没有它我们一天也不能过属灵的生活。当它说\Quote{今天}时，表明它必须每天领受，昨日的供应不够，但今天也必须同样赐予我们。我们每日需要它提醒我们，应当时刻献上这个祈祷，因为没有一天我们不需要通过吃和领受它来加强我们内在之人的心，尽管所用的表达\Quote{今天}可应用于他现世的生活，即当我们生活在这世界上时，求祢将这食粮供应我们。因为我们知道它将来会由祢赐给那些配得的人，但我们求祢今天赏给我们，因为除非一个人在这生命中蒙恩领受它，否则他将来永远不会成为其分享者。\par

% structure: p0067 source=h2 target=subsection reason=relative-heading-rank; base=h2

\subsection{第二十二章}

\Emph{论这句：\BibleQuote{免我们的债，等等。}}\par

\BibleQuote{宽恕我们的过犯，如同我们也宽恕那些亏负我们的人。} 啊，无法言说的天主的仁慈，祂不仅赐给我们一种祈祷的形式，教导我们一种蒙祂悦纳的生活方式，并透过所赐形式的要求（在其中祂命令我们时常祈祷），连根拔起了忿怒和忧愁的根源，而且还给祈祷的人提供一个机会，并向他们启示一条途径，藉此他们能促使天主对他们作出仁慈和善意的判决，并且以某种方式赐给我们能力，藉此我们能缓和审判官的判决，以我们宽恕的榜样吸引祂宽恕我们的过犯：当我们向祂说：\BibleQuote{宽恕我们，如同我们也宽恕。} 因此，一个人可以毫无忧虑、满怀信心地从这祈祷中请求赦免自己的过犯，如果他已宽恕了自己的债主，而非主的债主。因为有些人（这是极恶劣的）对于损害天主的事，无论罪行多大，都表现得平静、极其仁慈；但对自己有亏欠的事，即便最微小的不当，却成为最苛刻、毫不留情的追债者。所以，凡不存心宽恕得罪他的弟兄的人，藉此祈祷不是召来宽恕，而是定罪，并凭他自己的宣告，请求更严厉地审判自己，说：宽恕我，如同我也宽恕了。如果照他自己的请求受报应，那么除了照他自己的榜样，以不可遏止的义怒和不可撤销的判决受惩罚，还能有什么呢？因此，如果我们想得到仁慈的审判，我们也应当对得罪我们的人仁慈。因为我们宽恕了多少，也只有多少会被赦免，无论他们多么恶毒地伤害过我们。有些人畏惧这一点，当这祷文在教堂里被全体会众咏唱时，他们默默地略过这一句，唯恐他们自己的言词似乎是在捆绑自己而非宽恕自己，因为他们不明白，他们试图向所有人的审判者提出这些狡辩是徒劳的。祂愿意预先向我们表明祂将如何审判那些向祂哀求的人。因为祂不愿意让自己显得对他们严厉和无情，祂标示了祂审判的方式：正如我们渴望被祂审判一样，我们也应当审判我们的弟兄，如果他们有什么亏负我们的事，因为\BibleQuote{他没有施怜悯，必受无怜悯的审判。}\BibleRef{雅 2:13}\par

% structure: p0070 source=h2 target=subsection reason=relative-heading-rank; base=h2

\subsection{第二十三章}

\Emph{论\BibleQuote{不要让我们陷于诱惑}一句。}\par

接下来是：\BibleQuote{不要让我们陷于诱惑}，由此产生一个重要的问题，因为如果我们祈求不让我们受试探，那么我们考验的力量如何被证明呢？按照这段经文：\BibleQuote{凡未受试探的人，不被证明}；\BibleRef{德 34:11} 再次：\BibleQuote{忍受试探的人是有福的}？\BibleRef{雅 1:12} 那么，\BibleQuote{不要让我们陷于诱惑}这句，意思并非如此；即，不要允许我们永远不受试探，而是不要允许我们陷入试探时被征服。因为约伯受了试探，但并没有被引入试探。因为他没有归咎于天主的愚昧或亵渎，也没有用不敬的口屈服于诱惑者引他向的那个意愿。亚巴郎受了试探，若瑟受了试探，但他们都没有被引入试探，因为他们都没有向诱惑者屈服。接下来是：\BibleQuote{但救我们免于凶恶}，即，不要允许我们受魔鬼诱惑超过我们所能承受的，而是\BibleQuote{在试探中也给我们开一条出路，使我们能忍受得住。}\BibleRef{格前 10:13}\par

% structure: p0073 source=h2 target=subsection reason=relative-heading-rank; base=h2

\subsection{第二十四章}

\Emph{我们不应祈求其他事物，只应祈求包含在主祷文界限内的事物。}\par

那么你看出，审判者亲自向我们提议的祈祷方法和形式是什么，祂是该受我们以它祈祷的那一位。这形式中不包含对财富的祈求，没有荣誉的想法，没有对权力和能力的请求，没有提及身体健康和现世生命。因为永恒的作者愿意人向祂祈求任何不确定、微不足道、暂时的事物。所以，如果一个人撇开这些永恒的祈求，反而选择向祂祈求暂时和不确定的东西，那将是对祂的威严和仁慈的最大侮辱；并且，由于他祈祷的琐碎，他将招致审判者的愤怒而非平息。\par

% structure: p0076 source=h2 target=subsection reason=relative-heading-rank; base=h2

\subsection{第二十五章}

\Emph{论更崇高祈祷的特质。}\par

这祈祷虽然似乎包含了一切圆满的丰满，因为是主亲自发起并指定的，但它将属于它的人提升到我们上面所说的那个更高境界，并通过一个更崇高的阶段将他们带向那炽热的祈祷，这祈祷只有极少数人知道并经历过，更实在地说，是无可言喻的；它超越所有人类思想，其特点，我并非说有任何声音，而是没有舌头或言语的动作，而是心灵受到天上光芒的注入而开明，用非人类和有限的语言描述，却如从丰沛的泉源涌出，在思想的积累中丰富地倾泻，并不可言喻地向天主述说，在极其短暂的时间内表达如此伟大的事物，以至于当心灵回到平常状态时，无法轻易说出或叙述。我们的主也通过那些恳求的形式预先描绘了这种状态，据说当祂独自退到山上时，祂曾默默地倾泻这些恳求，并且在祈祷的苦痛中，祂甚至流出血滴，作为一个难以模仿的决心榜样。\par

% structure: p0079 source=h2 target=subsection reason=relative-heading-rank; base=h2

\subsection{第二十六章}

\Emph{论各种不同的定罪原因。}\par

然而，无论一个人天生拥有何等经验，谁又能充分描述使心灵受到激励、点燃和推动去作纯洁而最热切祈祷的各种感悟、理由和信服依据呢？以下我们仅凭天主的启迪，就所能回忆者，略举数例。因为有时在咏唱圣咏时，某篇圣咏中的一节会给我们带来热切祈祷的契机。有时弟兄声音的和谐旋律会激发懒散者的心灵去作恳切的祈求。我们也知道，咏唱者的发音和虔诚会大大增添旁听者的热忱。此外，一个成全之人的劝勉和一场灵性交谈，往往能将在场者的情感提升至最丰富的祈祷。我们还知道，一位兄弟或亲人的去世，同样使我们被带入完全的信服。对自己冷淡和疏忽的回忆，有时也会在我们心中激起一种健康的灵性热忱。如此，无人能怀疑，借着天主的恩宠，有无数的机会可以驱散我们心灵的冷淡和昏睡。\par

% structure: p0082 source=h2 target=subsection reason=relative-heading-rank; base=h2

\subsection{第二十七章}

\Emph{论不同的感悟种类。}\par

然而，那些感悟本身如何以及以何种方式从灵魂的最深处产生，同样难以追溯。因为有时，借着某种不可言喻的喜悦和心灵的锐利，一种极为有益的感悟之果会如此产生，以致它实际上因那不可控制的喜乐之宏大而爆发为呼喊；这时，内心的喜悦和极大的欢欣甚至在邻人的小室中也听得见。但有时，心灵在深沉宁静的奥秘中完全隐藏自己，以致突然光照的惊奇窒息了一切言语的声音，而敬畏的心灵要么将所有感受保留给自己，要么失去它们，并以不能用言语表达的叹息向天主倾吐它的愿望。但有时，它被如此压倒性的感悟和悲伤所充满，以致除了借泪水的洪流外，无法表达。\par

% structure: p0085 source=h2 target=subsection reason=relative-heading-rank; base=h2

\subsection{第二十八章}

\Emph{关于流泪的丰沛供应并不在我们自己能力之内的问题。}\par

杰尔马诺：我自己这卑微之人确实并非全然不知这感悟之情。因为当泪水因回忆我的过失而涌出时，我常被主的眷顾因你所描述的这种不可言喻的喜乐而如此更新，以至于那喜乐的宏大使我确信我不应对其宽恕感到绝望。我认为，如果这种心境能随我们自己的意愿而被唤起，那么没有比这更崇高的了。因为有时，当我渴望尽全力激励自己达到同样的感悟和眼泪，并将我所有的过失和罪过置于眼前时，我却无法带回那丰沛的泪水，我的眼睛干燥坚硬如最硬的燧石，以致一滴泪也流不出来。因此，当我为自己那丰沛的泪水而庆幸时，也同样哀叹自己无法在愿意时将其唤回。\par

% structure: p0088 source=h2 target=subsection reason=relative-heading-rank; base=h2

\subsection{第二十九章}

\Emph{关于由眼泪产生的各种感悟的答案。}\par

依撒格：并非每一种流泪都是由同一种感受或同一种德行产生的。因为有一种哭泣源于我们罪恶的刺伤心灵，如经上所载：\BibleQuote{我因呻吟而疲惫；每夜我洗涤我的床铺，以眼泪浸透我的卧榻。}\BibleRef{咏 6:7}又说：\BibleQuote{愿眼泪日夜如急流流下，不要让你自己休息，也不要让你的眼目歇息。}\BibleRef{哀 2:18}另一种哭泣源于对永恒美善的默观和对未来荣耀的渴望，由此，从不可控制的喜乐和无限的欢欣中涌出更丰富的泪泉，同时我们的灵魂渴慕那大能生活的天主，说：\BibleQuote{我何时来，出现在天主面前？我的眼泪昼夜成了我的食物，}\BibleRef{咏 42:3-4}每日哭泣哀叹说：\BibleQuote{我有祸了，因我的寄居延长了；}\BibleRef{咏 120:5}以及：\BibleQuote{我的灵魂寄居太久。}\BibleRef{咏 120:6}又有一种流泪的方式，虽然没有致命之罪的意识，却仍出于对地狱的恐惧和对那可怕审判的回忆，先知因这恐惧而受击打，向天主祈祷说：\BibleQuote{求祢不要与祢的仆人进入审判，因为在祢面前凡活着的人没有一个是义的。}\BibleRef{咏 143:2}还有另一种眼泪，并非由对自己的认识引起，而是由他人的硬心和罪恶引起；撒慕尔即为撒乌耳哭泣，主在福音中为耶路撒冷城哭泣，耶肋米亚昔日也为她哭泣，后者这样说：\BibleQuote{但愿我的头是水，我的眼是泪水的泉源！我要为我的百姓中被杀的人昼夜哭泣。}\BibleRef{耶 9:1}或者还有如我们在第一百零一篇圣咏中所听到的泪水：\BibleQuote{因为我吃灰烬如同吃饭，我以眼泪调和我的饮料。}\BibleRef{咏 102:10}这些泪水肯定不是由与第六篇圣咏中忏悔者所流之泪相同的感受引起的，而是由于今生的忧虑、困苦和损失，这些压制着活在这世界上的义人。这不仅从圣咏本身的话语中清楚显示，也从其标题中显示，该标题以那穷人的身份写道——福音记载说：\BibleQuote{神贫的人是有福的，因为天国是他们的。}\BibleRef{玛 5:3}——\BibleQuote{穷人的祈祷，当他困苦时，向天主倾吐他的祈求。}\BibleRef{咏 102:1}\par

% structure: p0091 source=h2 target=subsection reason=relative-heading-rank; base=h2

\subsection{第三十章}

\Emph{论眼泪不应在不能自然流出时强行挤出。}\par

由这些眼泪大不相同的是那些从干涸的眼睛中挤出来的泪水，而心却是硬的：虽然我们不能认为这些泪水完全无果（因为流出它们的尝试是出于善意，尤其是那些尚未能获得完美知识或尚未彻底洗净过去或现在罪污的人），但那些已前进至对德行之爱的人，当然不应这样强行挤出泪流，也不应费大力气去追求外在人的哭泣，因为即使产生了，也决不会达到自发眼泪的丰沛。相反，它会因人的努力而使祈祷者的灵魂沮丧，并羞辱他，使他陷入世俗事务，远离天上的高峰——祈祷者敬畏的心本应坚定不移地定在那里——并迫使他放松对祈祷的专注，因贫乏而强挤的泪水而变得衰弱。\par

% structure: p0094 source=h2 target=subsection reason=relative-heading-rank; base=h2

\subsection{第三十一章}

\Emph{院长安当论祈祷状况的意见。}\par

为使你们看见真祈祷的特性，我要给你们的不只是我自己的意见，而是有福的安当的意见：我们曾见过他有时在祈祷中如此持久，以致常常在他心神超拔祈祷时，当太阳开始升起，我们听到他在神魂激昂中喊道：\Quote{啊，太阳，你为何妨碍我？你升起正是为这个目的，即把我从这真光的光明中拉走？} 关于祈祷的结局，也是他这句属天且超越凡人的名言：\Quote{一个隐修士了解自己和所祈祷的言语，那还不是完美的祈祷。} 如果我们，按照我们微薄的能力，也敢在这辉煌的言论上补充一点，我们将按照我们所经验过的，提出主所垂听的祈祷的标志。\par

% structure: p0097 source=h2 target=subsection reason=relative-heading-rank; base=h2

\subsection{第三十二章}

\Emph{关于祈祷蒙垂听的证明。}\par

当我们祈祷时，没有犹豫介入，并以某种绝望击溃我们恳求的信心，反而感到我们通过倾心祈祷已获得所祈求的，我们就不怀疑我们的祈祷已有效地达到天主。因为主的话不会收回：\BibleQuote{你们祈祷，不论求什么，只要相信必得，必给你们成就。} \BibleRef{谷 11:24}\par

% structure: p0100 source=h2 target=subsection reason=relative-heading-rank; base=h2

\subsection{第三十三章}

\Emph{异议：如此蒙垂听的信心只属于圣人。}\par

杰尔曼：我们确实相信，这种蒙垂听的信心来自良心的纯洁，但对我们这些心仍被罪恶刺痛的人，怎能拥有它呢？我们没有功德为我们代求，使我们能自信地认为我们的祈祷会被垂听。\par

% structure: p0103 source=h2 target=subsection reason=relative-heading-rank; base=h2

\subsection{第三十四章}

\Emph{回答：关于祈祷蒙垂听的不同原因。}\par

\Person{依撒格}：福音和先知的话教导我们，祈祷被俯听的原因，依照灵魂不同变化的状态而各不相同。因为我们有主的话指出，两人同心合意时祈求的回答之果实；正如所说：\BibleQuote{如果你们两人在地上同心合意，无论为任何事祈求，我在天之父必要为他们成就。}\BibleRef{玛 18:19} 另一个是在信德的圆满内，这信德被比作一粒芥菜子。祂说：\BibleQuote{因为，如果你们有信德像一粒芥菜子，你们要对这座山说：从这边移到那边去，它必移过去；在你们没有不可能的事。}\BibleRef{玛 17:19} 另一个是在祈祷中持续不断，主的话称之为，由于祈求时不倦的坚忍，即强求：\BibleQuote{因为，我切实告诉你们：纵然不因为他是朋友，而因为他的强求，他必起来，按他所需要的给他。}\BibleRef{路 11:8} 另一个是在施舍的果实上：\BibleQuote{要把施舍物藏在穷人的心里，它必在患难时为你祈祷。}\BibleRef{德 29:15} 另一个是在生活的净化与慈善工作中，正如所说：\BibleQuote{解开邪恶的锁链，解除压迫的重轭；} 并且在责备无益禁食的不结果实之后，他说：\BibleQuote{那时，你呼求，上主必要俯听；你哀号，祂必说：我在这里。}\BibleRef{依 58:6} 有时过多的苦难也促使祈祷被俯听，正如所说：\Quote{我在患难中呼求了上主，祂便俯听了我；} 又说：\Quote{不要虐待外方人，因为他若向我呼求，我必俯听他，因为我是仁慈的。} 因此你看见，得到回答的恩赐有多少种方式，以致无人需要被良心的绝望压垮，为了获取那些有益于救恩和永恒的事物。因为如果在默想我们可怜的状况时，我承认我们完全缺乏上述所有那些德行，既没有那值得称赞的两人同心，也没有那比作一粒芥菜子的信德，也没有先知描述的那些虔敬的工作，那么我们肯定不会缺少那强求，祂供给所有渴望它的人，单凭这强求，主就应许祂会赐予凡被祈求赐予的一切。因此我们应当毫不怀疑地坚持，毫不怀疑地确信，借着在这些事上持续，我们必会获得那些我们按天主意愿所祈求的一切。因为主渴望赐予天上永恒之事，敦促我们好像用我们的强求来约束祂，因为祂不仅不轻视或拒绝强求的人，反而欢迎并赞美他们，且最仁慈地应许赐予他们坚忍希望的一切；说：\BibleQuote{你们求，必要给你们；你们找，必要找着；你们敲，必要给你们开。因为凡求的，就必得到；找的，就必找到；敲的，就必给他开。}\BibleRef{路 11:9} 又说：\BibleQuote{凡你们在祈祷中祈求的一切，只要相信，你们必得到，在你们没有不可能的事。} 因此，即使我们提到的所有被俯听的理由全然缺乏，至少强求的热诚可以激励我们，因为这置于任何愿意的人的能力之内，无需任何功德或劳苦的困难。但任何祈求者不要怀疑他一定不会被俯听，只要他怀疑自己是否被俯听。但我们也被真福达尼尔的榜样教导，也要不懈地向主祈求这一点，因为他虽然自开始祈祷的第一天就被俯听，但他只在二十一天后才得到祈求的结果。因此，我们也不应该在已开始的祈祷的热诚中松懈，如果我们觉得回答来得缓慢，害怕也许回答的恩赐在天主的上智安排中被延迟，或是那将带天主祝福给我们的天使，当祂从全能者的面前出来时，可能被魔鬼的抵抗阻碍，因为确实，如果祂发现我们松懈了所祈求的热诚，祂就不能传递并带给我们所渴望的恩赐。而这事无疑也会发生在上述先知身上，除非他以无比的坚定延长并坚持祈祷直到第二十一天。那么，让我们绝不被绝望从我们这信德的确信上击倒，即使当我们觉得远未获得我们所祈求的，也不要对主的应许有任何怀疑，祂说：\BibleQuote{凡你们在祈祷中祈求的一切，只要相信，你们必得到。}\BibleRef{玛 21:22} 因为我们最好考虑真福福音作者若望的这句话，藉此，这个问题的含糊之处被清楚解决：他说：\BibleQuote{这就是我们向祂所怀有的信赖：无论我们按祂的旨意求什么，祂必俯听我们。}\BibleRef{若壹 5:16} 因此，祂命令我们，只在那些不是为了我们自己利益或现世舒适，而是符合主的旨意的事上，对回答有完全而无疑的信赖。我们也由主祷文被教导将此放入我们的祈祷中，在那里我们说\Quote{愿祢的旨意奉行}——\Emph{祢的}，而非我们的。因为如果我们也记得宗徒的话：\BibleQuote{我们不知道该如何祈祷才对}\BibleRef{罗 8:26}，我们将看见我们有时祈求相反我们救恩的事物，并且我们被那位比我们更正确、更真实地看到什么对我们有益的天主，以最上智的方式拒绝了我们的请求。显然，这事也发生在外邦人的导师身上，当他祈祷那出于他益处、由主所允许的、打击他的撒旦使者会被移去时，说：\BibleQuote{为了这事，我曾三次求主使它离开我。但祂对我说：我的恩宠够你用，因为我的能力在软弱中才全显出来。}\BibleRef{格后 12:8} 这情感甚至我们的主也表达了，当祂以所取的人性祈祷时，为能给我们祈祷的样式，如同其他事物也由祂的榜样所赐予；这样说：\BibleQuote{父啊！如果可能，就让这杯离开我，但不要照我的意愿，而要照祢的意愿。}\BibleRef{玛 26:39} 尽管祂的意愿肯定与圣父的意愿不相违背，因为祂来是为拯救失落的，并交出祂的生命为大众作赎价；如祂自己说：\BibleQuote{没有人夺去我的生命，是我自己舍掉的。我有权舍掉，也有权再取回。}\BibleRef{若 10:18} 在这种身份中，在第三十九篇圣咏中，真福达味歌唱了祂始终与圣父保持的旨意合一：\Quote{我的天主，我乐意承行祢的旨意。} 因为即使我们读到有关圣父的话：\BibleQuote{天主如此爱了世界，竟赐下祂的独生子}\BibleRef{若壹 3:16}，我们也同样找到有关圣子的话：\BibleQuote{祂为我们舍弃了自己。}\BibleRef{迦 1:4} 正如论及圣父说的：\BibleQuote{祂没有怜惜自己的儿子，反而为我们众人交出了祂}\BibleRef{罗 8:32}，同样论及圣子记载：\Quote{祂被献上，因为祂自己愿意。} 并且显示圣父与圣子的旨意在所有事上都合一，以致即使在主复活的实际奥迹中，我们也被教导运作没有不一致。因为正如真福宗徒声明圣父使祂的身体复活，说：\BibleQuote{天主圣父，那使祂从死者中复活的}\BibleRef{迦 1:1}，同样圣子也证明祂自己要重建祂身体的圣殿，说：\BibleQuote{你们拆毁这座圣殿，三天内我要把它重建起来。}\BibleRef{若 2:19} 因此，我们被所有这些列举的主的榜样所教导，也应当以同样的祈祷结束我们的恳求，并总是在我们所有的祈求中附加这句：\BibleQuote{但是不要照我的意愿，而要照祢的意愿。}\BibleRef{玛 26:39} 但足够清楚的是，一个不专心祈祷的人不能遵守通常在弟兄集会中、在礼仪结束时所做的三重敬礼。\par

% structure: p0106 source=h2 target=subsection reason=relative-heading-rank; base=h2

\subsection{第三十五章}

\Emph{论在内室并关上门祈祷}\par

首先，我们必须最谨慎地遵守福音的诫命，它告诉我们进入内室，关上门，向我们的圣父祈祷。我们可以如下方式实现这一诫命：当我们从内心除去一切思想和焦虑的喧闹，在隐秘和亲密的交往中向主倾诉我们的祈祷时，我们就在内室祈祷。当我们闭口不言，完全沉默地向那位不是审阅言语而是审阅心灵的主祈祷时，我们就关上了门。当我们从心中和热切的意念中唯独向天主披露我们的恳求，以致敌对的权势甚至无法察觉我们祈祷的性质时，我们就在暗中祈祷。因此，我们应当在完全的沉默中祈祷，这不仅要避免我们低语或较大声的说话打扰站在附近的弟兄，干扰正在祈祷者的思想，而且也要使我们的恳求的意图不被那些在我们祈祷时特别窥探我们的敌人所察觉。因为这样我们将实现这条诫命：\BibleQuote{谨守你口中之门的，是从那睡在你怀中的妇人。} \BibleRef{米 7:5}\par

% structure: p0109 source=h2 target=subsection reason=relative-heading-rank; base=h2

\subsection{第三十六章}

\Emph{论简短而无声祈祷的价值}\par

因此，我们应当常常祈祷，但祈祷要简短，免得我们祈祷过长时，狡猾的敌人可能得以在我们心中植入什么。因为真正的祭献就是\Quote{天主所悦纳的祭献是破碎的心灵}。这是有益的奉献，这些是纯洁的奠祭，即\Quote{正义的祭献}，\Quote{赞颂的祭献}，这些是真实而肥美的牺牲，\Quote{充满骨髓的全燔祭}，是由忏悔和谦卑之心奉献的，那些实行我们所谈论的这种对心灵的控制和热忱的人，能以有效的力量歌唱：\Quote{愿我的祈祷如馨香上达于你面前，愿我的双手高举如同晚间的祭献。}\par

\Person{圣依撒格}这些话使我们眼花缭乱，而非满足；晚祷结束后，我们休息了片刻，并打算在黎明时分再次返回，以得到更充分的讨论，我们离开了，为获得这些训诫以及他承诺的保证而欢喜。因为我们觉得，虽然祈祷的卓越性已经向我们展示，但我们仍未从他的讲论中理解其本质，以及藉以获得并保持持续祈祷的能力。\par

\SourceRecord{Translated by C.S. Gibson.；Nicene and Post-Nicene Fathers, Second Series；Vol. 11.；Edited by Philip Schaff and Henry Wace.；Buffalo, NY: Christian Literature Publishing Co.,；1894.；<http://www.newadvent.org/fathers/350809.htm>.}

% structure: p0001 source=h1 target=section reason=page-title

\section{第十次会谈}

\Emph{圣若望·卡西盎著}\par

院长依撒格第二次会谈。论祈祷。\par

% structure: p0004 source=h2 target=subsection reason=relative-heading-rank; base=h2

\subsection{第一章}

\Emph{引言。}\par

在隐修士们崇高的习俗中，藉天主的助佑，我们虽以朴素无华的文体加以阐述，但我们叙述的进程迫使我们插入并安排一些内容；这些内容，可以说，在美丽的躯体上造成了瑕疵。然而我毫不怀疑，藉此将为一些较为单纯的人，就我们在\WorkTitle{创世纪}中读到的全能天主的肖像赋予不少教导，尤其是当考虑到一个如此重要的教义的理由时，若对此无知，必会导致可怕的亵渎和对天主教信仰的严重损害。\par

% structure: p0007 source=h2 target=subsection reason=relative-heading-rank; base=h2

\subsection{第二章}

\Emph{论埃及省通报复活节时间的习惯。}\par

在埃及地区，自古相传一种习俗：主显节过后——该省的司铎视此日为主受洗及祂降生成人的日子，因此不像西方各省那样分别庆祝两个奥迹，而是在这同一庆节中一同纪念——亚历山大里亚主教致函埃及各地教会，藉此不仅在各城市，也在各隐修院中指明四旬期的开始和复活节的日子。依照此习俗，在先前与院长依撒格举行会谈后的几天，上述城市的主教德奥斐洛的节庆信函到达了。信中在宣告复活节的同时，他也仔细考虑了\OriginalTerm{拟人论者}{Anthropomorphites}的愚昧异端，并予以充分驳斥。这封信由于他们的单纯和错误，几乎被所有居住在全埃及省的隐修士团体以如此苦毒的态度接受，以致大部分长老反其道而行之，决定上述主教应被全体弟兄憎恶，视之为最恶劣异端之沾染者，因为他似乎否认全能天主具有人形，从而攻击圣经的教导，而圣经明示亚当是照天主的肖像受造的。最后，连那些住在斯刻提旷野、在成全和知识上超越埃及所有隐修院修士的人，也拒绝这封信。除了我们团体的司铎帕夫努提院长外，其他三位在同一旷野领导其他三座教堂的司铎，没有一个允许在他们的聚会中诵读或传阅这封信。\par

% structure: p0010 source=h2 target=subsection reason=relative-heading-rank; base=h2

\subsection{第三章}

\Emph{论院长撒拉丕翁及其因无知而陷入的拟人论异端。}\par

当时被这种错误见解所迷惑的人中，有一位名叫撒拉丕翁，他是一位长期恪守苦修生活的人，在实际操练上全然成全。他关于前述教义观点的无知，成为所有持守真信仰者的绊脚石，恰如他本人无论在生活的功德上还是在他所居留的时间长度上都超越几乎所有的隐修士。当圣司铎帕夫努提多次劝诫他，仍不能使他回归正道，因为他认为这种见解是新奇的，是前人从未知晓或传授的。恰巧有一位来自卡帕多细亚地区的执事，名叫佛提诺，博学多才，他来访同一旷野中的弟兄。有福的帕夫努提极为热忱地接待了他，为了证实上述主教信函中所阐述的信仰，便让他站在中间，当着全体弟兄的面问他：东方各天主教会如何解释\WorkTitle{创世纪}中\BibleQuote{让我们按我们的肖像和模样造人}\BibleRef{创 1:26}的那段经文。他解释说，天主的肖像和模样被所有教会领袖不是按照字面的粗俗含义，而是以属灵的方式理解，并且用许多圣经章节非常充分地支持这一点。他指出，那无限、不可理解、不可见的光荣不可能被容纳于人的形态和模样中，因为祂的本性是无形体、无组合、纯一的，既不能被眼睛看见，也不能被心智领会。最后，这位老人被这位博学之士众多且极其有力的论证所摇动，被吸引到天主教传统的信仰中。帕夫努提院长和我们全都因他的归从而充满极大的喜悦，因为主没有允许这样一位年迈且被如此多美德加冕的人（他仅因无知和乡村般的单纯而误入歧途）直到最后仍偏离正道。当我们起身感恩，并一同向主献上祈祷时，这位老人在祈祷中思绪混乱，因为他觉得他惯于在祈祷时摆在自己面前的神性拟人形象从他心中被驱散了。突然，他泪如泉涌，不断啜泣，扑倒在地，大声呻吟道：\Quote{哀哉！我这可怜的人！他们从我这里夺去了我的天主，我现在无从把握，不知道该朝拜谁、向谁说话。}对此情景我们大为惊骇，加之先前会谈的效果仍留在我们心中，我们便回到院长依撒格那里。我们见他临近，便以这些话向他说话。\par

% structure: p0013 source=h2 target=subsection reason=relative-heading-rank; base=h2

\subsection{第四章}

\Emph{论我们回到院长依撒格处，并询问关于前述老人陷入的错误。}\par

尽管除了近日新起的这些事以外，我们从前一次关于祈祷本质的讲道中所获得的快乐，也会催促我们放下一切，回到你的圣德那里，然而撒拉彼翁院长的这一严重错误——我们认为是由于最卑劣的魔鬼的诡计而产生的——却更增添了我们的这种渴望。因为当我们想到，由于这一无知，他不仅彻底丧失了他在这旷野中五十年里如此值得称赞地所做的一切劳苦，而且还冒了永死的危险，这使我们陷入了不小的沮丧。因此，我们首先想知道，为什么以及因何缘故如此严重的错误潜入了他。其次，我们愿意被教导，如何才能达到你在不久前不仅充分而且精彩地论述过的祈祷状态。因为那令人钦佩的讲道对我们产生了这样的效果：它只是让我们眼花缭乱，却没有告诉我们如何实施或获得它。\par

% structure: p0016 source=h2 target=subsection reason=relative-heading-rank; base=h2

\subsection{第五章}

\Emph{对于上述异端的答复。}\par

我们不必惊讶，一个真正纯朴的人，从未受过关于天主性本体和本质的任何教导，仍可能被一种纯朴的错误和长期错误的习惯所纠缠和欺骗，并且（更真实地说）继续停留在原始的谬误中，这种谬误并非如你所猜想的那样是由于魔鬼的新幻觉，而是由于古代外邦世界的无知所造成的；因为他们遵守那种错误观念的习惯，即他们曾用来崇拜以人形象形成的魔鬼，如今甚至认为，真正的天主那不可理解、不可言喻的光荣应该以某种形象的局限性来被崇拜，因为他们相信，如果面前没有某种形象，他们就无法把握和持守任何东西，这形象可以在他们祈祷时不断向其说话，并在心中携带，常置于眼前。针对他们的这种错误，可以用这段经文：\BibleQuote{他们将不可朽坏的天主的光荣，变成了可朽坏的人的图像的模样。} \BibleRef{罗 1:23} 耶肋米亚也说：\BibleQuote{我的子民将自己的光荣换成了偶像。} \BibleRef{耶 2:11} 这种错误，虽然由于我们所说的这一起源，在某些人的观念中根深蒂固，但同样地，它也沾染在那些从未受过外邦世界迷信玷污的人心中，其借口是这段经文：\BibleQuote{让我们按我们的肖像和模样造人，} \BibleRef{创 1:26} 无知和纯朴是它的根源，以至于实际上由于这种可憎的解释，产生了一种称为\OriginalTerm{拟人论者}{Anthropomorphites}的异端，它固执地坚持认为无限而单纯的天主性本体是按照人的形体和外貌塑造的。然而，任何受过天主教教义教导的人都会将其视为外邦人的亵渎而憎恶，从而在祈祷中达到完全纯洁的状态，这种状态不仅不会在其祈祷中结合任何天主性的形象或身体轮廓（连提说都是罪），而且甚至不允许自身存有名字的记忆、行动的出现或任何特征的轮廓。\par

% structure: p0019 source=h2 target=subsection reason=relative-heading-rank; base=h2

\subsection{第六章}

\Emph{关于耶稣基督以祂的谦卑或祂的光荣方式显现给我们各人的原因。}\par

因为按照其纯洁的程度，正如我在上一次讲道中所说的，每个心灵在祈祷中都被提升和塑造，只要它放弃对尘世和物质事物的思虑，如其纯洁程度所能推动它前进，并使其能用灵魂的内眼看到耶稣，要么仍处于祂的谦卑和肉身中，要么在祂光荣和威严的荣耀中来临。因为那些仍然停留在某种犹太人的软弱状态中，并且不能与宗徒一同说：\BibleQuote{我们虽然曾按肉体认识基督，但如今不再这样认识祂了；} \BibleRef{格后 5:16} 的人，不能看见耶稣在祂的王国中来临；唯有那些与祂一同从低贱和属世的工作和思想中升起，登上高处的孤寂之山（这山远离一切属世思想和烦恼的搅扰，免受一切罪恶的干扰，并由纯洁的信德和美德的高峰所高举，将祂面容的光荣和祂荣耀的形象启示给那些能以灵魂的纯洁眼睛注视祂的人），才能以最纯洁的眼睛注视祂的天主性。但耶稣也被那些住在城镇、村庄和乡间的人看见，即那些从事实际事务和劳作的人，但亮度不及那些能与祂一同登上上述美德之山的人（即伯多禄、雅各伯和若望）。因为如此，祂在旷野中向梅瑟显现，并与厄里亚说话。我们的主既然想确立这一点，并为我们留下完全纯洁的榜样，虽然祂自己，那不可侵犯的圣德的泉源，不需要外部的帮助和独处的辅助来获得它（因为纯洁的圆满不能被人群的任何污点所玷污，也不能被与人的交往所污染，因为祂洁净和圣化一切被污染的事物），但祂仍独自退到山上祈祷，从而以祂退隐的榜样教导我们，如果我们也想以纯洁无瑕的心灵情感接近天主，我们也应当远离一切人群的搅扰和混乱，以便在仍活在肉身中时，在某种程度上设法适应那以后应许给圣人的福乐的模样，并且\BibleQuote{天主成为对我们来说万物中的万有。} \BibleRef{格前 15:28}\par

% structure: p0022 source=h2 target=subsection reason=relative-heading-rank; base=h2

\subsection{第七章}

\Emph{什么是我们的终向和完美的福乐。}\par

那时，在我们身上将完美地实现救主在向圣父为祂的门徒祈祷时所说的祈祷：\Quote{祢爱我的爱在他们内，他们也在我们内}；又说：\Quote{愿众人合而为一，父啊！如同祢在我内，我在祢内，使他们在我们内也合而为一。} 那时，天主那完美的爱——祂首先爱了我们\BibleRef{若一 4:16}——也渗入了我们内心的感情，藉着主这个我们认为绝不能无效的祈祷的应验。这将发生：当天主成为我们全部的爱，我们所有的欲望、意愿、努力、思想、生命、言语和呼吸，以及那已存在于圣父与圣子、圣子与圣父之间的合一，倾注在我们心中和意念里，以致祂以纯洁、真诚、不可分解的爱爱我们，我们也能以持久、不可分离的深情与祂结合，因为我们与祂如此合一，以致我们所呼吸、思想或言说的都是天主。如我所说，我们达到了之前谈到的那个目的，同一主在祈祷中希望在我们身上实现：\BibleQuote{愿众人都合而为一，就如我们原为一体，我在他们内，祢在我内，使他们完全合而为一}\BibleRef{若 17:22}\BibleQuote{父啊！祢所赐给我的人，我愿我在哪里，他们也同我在那里。}\BibleRef{若 17:22}这应当是隐修士的目标，他全部的目的就是蒙恩即使在肉身内也能拥有未来福乐的图像，并在今世开始预尝天上生命与荣耀的某种凭据。我说，这就是一切完美的终向：心灵从一切肉欲中净化后，日日被提升至神性事物，直到整个生命和心内的一切思念成为一篇连续的祈祷。\par

% structure: p0025 source=h2 target=subsection reason=relative-heading-rank; base=h2

\subsection{第八章}

\Emph{关于成全修行的问题，我们藉此能达到对天主的恒久纪念。}\par

\Person{杰尔曼诺}：我们对前一次会谈的惊叹与敬畏所导致的困惑，因之我们再次返回，如今更加深。因为正如此次教导的激励，我们渴望成全的福乐，我们却更深地陷入绝望，因为我们不知如何寻求或获得达到如此崇高境界的修行。因此，我们恳求你耐心容许我们（这也许需要相当多的讲述来展开）解释我们在静坐于斗室时开始长时间默想的内容，尽管我们知道你的圣德丝毫不受弱者软弱之扰，而正是为此，这些软弱应当被公开陈述，以便其中不当之处得到纠正。我们的观念是：任何艺术或修行体系的成全都必须从一些简单的入门开始，首先要习惯于稍微容易且幼稚的开端，以便藉着一种理性的奶水被滋养和训练，逐渐成长，并逐步从最底层攀升至高境。当藉此方法进入更显明的原则，并可以说通过了职业的入门之门后，它便毫不困难地抵达内在的圣殿和成全的高峰。因为一个孩子若不先仔细学习字母，怎能发出最简单的音节组合？或者，一个尚不能连接简短句子的人，怎能学会快速阅读？一个在语法学上缺乏训练的人，又怎能达到修辞学的雄辩或哲学的知识？因此，对于这至高的学问——我们藉之被教导甚至与天主结合——我毫不怀疑有一些系统的基础，必须首先牢固奠定，然后成全的高塔才能在其上建造和竖立。我们稍有些想法，即其首要原则如下：首先，我们应学习藉何种默观可以把握并默想天主；其次，我们应设法牢牢抓住这个主题——无论它是什么，我们毫不怀疑这是所有成全的高峰。因此，我们希望你向我们展示一些关于此纪念的材料，藉此我们可以在心中构思并恒久保持对天主的观念，以便时常将它置于眼前，当我们发现已从祂那里失落时，能立刻恢复并返回，并在无需因四处徘徊和寻找而耽搁的情况下成功再次抓住它。因为常发生这样的情况：当我们从神性默观中迷失，如同从致命的沉睡中醒来，彻底振作，寻找那能使已毁的神性纪念复兴的主题时，我们被实际寻找的耽搁所阻碍，在找到之前再次被从努力中拉走，而心灵的目的在神性洞见实现之前已经消失。这种困扰必然发生在我们身上，因为我们没有将某种特别的东西稳固地置于眼前，如同一原则，使飘忽的思想经过多次离题和多样漫游后能被召回；如果我可以这样表达，就是在长久风暴后进入安宁的港湾。因此，由于这种知识的匮乏和困难，心灵总是盲目地、如同醉酒般在各种事物中飘荡，甚至不能长时间牢固把握任何偶然而非有意发生的神性事物；同时，由于它不断接受一个又一个事物，它既未注意到它们的开始和起源，也未注意到它们的结束。\par

% structure: p0028 source=h2 target=subsection reason=relative-heading-rank; base=h2

\subsection{第九章}

\Emph{关于由经验获得的领悟效力的答复。}\par

\Person{依撒格}：您这种精细而深入的探究表明，您已经非常接近纯洁了。因为除非一个人被心灵的谨慎而有效的勤奋和警醒的焦虑所催促，去探究这些问题的深处，除非一个人通过恒常追求节制生活的目标，并经实际经验学会尝试进入这种纯洁并叩击其门，否则他根本不可能就这些问题进行探究，更不用说内省和分辨了。因此，我看到您——我不是说您站在我们所谈论的那种真正祈祷的门前，而是说您用经验之手触摸了它的内室和内在部分，并且已经抓住了它的一些部分——我不认为我会有任何困难将您引入它的大厅，让您在其深处漫步，正如主所引导的那样；我也不认为您会被任何障碍或困难阻挡，无法探究将要向您展示的事。因为一个谨慎地认识到自己应当询问什么的人，离理解只有一步之遥；而一个开始明白自己多么无知的人，离知识也不远了。因此，如果我在先前关于祈祷完美的论述中曾经保留没有讨论的内容，现在却予以揭示，我并不害怕被指控泄露秘密或轻率，因为我认为，即使没有我言辞的帮助，天主恩宠也会向我们这些致力于这个主题和兴趣的人解释其力量。\par

% structure: p0031 source=h2 target=subsection reason=relative-heading-rank; base=h2

\subsection{第十章}

\Emph{论持续祈祷的方法}\par

因此，依照你恰当地比作教导孩童的那套方法（孩童只能接受字母表的最初几课，辨识字母的形状，并用稳健的手描画其笔画，若他们借助一些精心刻在蜡板上的样本和形状，通过不断注视和每日模仿而习惯于表达其图形，便能如此），我们也必须将此属灵默观的形式交付与你，使你得以始终以最大的定力注视它，并在不间断的持续中学习有益地思考它，且经由实践与默想而攀升至更崇高的洞见。如此，我们将向你提出这套你所渴求的系统的公式，以及祈祷的公式——每一位隐修士在迈向不断记念天主的进程中，惯于在心中不断反复思量它，摆脱所有其他类型的思绪；因为若他未摆脱一切肉体的挂虑和焦虑，便无法持久把握它。这公式由少数几位尚存的最古老教父传授给我们，我们也只透露给极少数真正热忱者。因此，为保持不断记念天主，这虔敬的公式应常摆在你面前：\Quote{天主，求你快来拯救我；上主，求你速来扶助我。}因为这一短诵从整部圣经中被选出来用于此目的，并非无理。它涵盖人性中能被植入的一切情感，并能恰当而圆满地适应各种境况与一切攻击。因为它包含面对一切危险时对天主的呼求，包含谦卑而虔诚的宣认，包含焦虑的警觉与不断的敬畏，包含对自身软弱的认知、对垂允的信赖，以及对那临在且随时可得的助佑的确证。因为不断呼求其保护者的人，确知他常近在咫尺。它包含爱德的热忱，包含对仇敌阴谋的洞察与畏惧——那昼夜被其包围的人，承认若无保护者的援助便无法解脱。此短诵对于所有受魔鬼攻击者，是坚不可摧的城墙，也是不可穿透的铠甲与坚固的盾牌。它不让那些处于愁闷、焦虑、或被悲伤及各种思绪压迫的人对拯救的良方绝望，因为它显示那被呼求者常注视我们的奋战，且不远离祈求者。它警告那些蒙受属灵成就与心喜之福的人，不应因幸福境况而自满或骄傲，因为若无天主作保护者，这境况无法持久，而它呼求他不仅时常，且迅速帮助。我说，这短诵对我们每个人，无论处于何种境况，都有益且有用。因为那常并凡事求帮助的人，显示他需要天主的助佑不仅在忧愁或艰难的事上，也同样在顺利和幸福的事上，以便从前者得拯救，在后者得持续，因为他知道在这两种境况中，人的软弱若无他的助佑便无法持久。我被贪饕之情所扰。我渴望旷野所无的食物，在荒凉的旷野中，王室的佳肴香气向我飘来，我发现自己甚至违愿地被牵引去渴望它们。我须立刻说：\Quote{天主，求你快来拯救我；上主，求你速来扶助我。}我被诱惑要提前晚饭的时间，或我怀着极大忧苦试图坚守正确且惯常的清淡饮食的限度。我须以呻吟呼喊：\Quote{天主，求你快来拯救我；上主，求你速来扶助我。}胃弱阻碍我，当我因肉身的攻击而渴望更严苛的禁食时，或腹部干燥与便秘令我恐惧。为使我的愿望得以实现，或为让肉欲之火不经更严格禁食的补救而熄灭，我须祈祷：\Quote{天主，求你快来拯救我；上主，求你速来扶助我。}当我按适当时辰的召唤来到晚餐时，我厌恶进食，无法吃任何东西以满足本性的需求：我须叹息呼喊：\Quote{天主，求你快来拯救我；上主，求你速来扶助我。}当我为持守心灵的坚定而欲投身阅读时，头痛干扰并阻止我，在第三时辰睡眠将我的头粘在圣经页上，我被迫要么超出、要么提前指定的休息时间；最终，一股不可抗拒的睡意迫使我缩短诵唱圣咏的规范规则：同样，我须呼喊：\Quote{天主，求你快来拯救我；上主，求你速来扶助我。}睡眠从我眼中被夺去，许多夜晚我被魔鬼引起的失眠折磨得疲惫不堪，夜间一切安息与休息远离我的眼皮；我须叹息并祈祷：\Quote{天主，求你快来拯救我；上主，求你速来扶助我。}当我仍处于与罪的争斗中时，突然肉身的激惹触动我，试图以快感引诱我在睡梦中同意。为使外在的怒火不致烧毁贞洁的芬芳花朵，我须呼喊：\Quote{天主，求你快来拯救我；上主，求你速来扶助我。}我感受到情欲的冲动被移除，激情的燥热在我肢体中消逝：为使这已获得的善境，或更确切地说，天主的恩宠得以更持久或永远与我同在，我须恳切地说：\Quote{天主，求你快来拯救我；上主，求你速来扶助我。}我被愤怒、贪婪、阴郁的刺痛所扰，被迫搅乱我原先珍视的平安状态：为使我不致被狂怒之情卷入苦胆之苦，我须以深深的呻吟呼喊：\Quote{天主，求你快来拯救我；上主，求你速来扶助我。}我受试探，因懈怠、虚荣、骄傲而自满，我的心思以微妙的思绪，因他人的冷淡和疏忽而暗自沾沾自喜：为使这危险的敌人建议不致控制我，我须以全心的痛悔祈祷：\Quote{天主，求你快来拯救我；上主，求你速来扶助我。}我获得了谦卑与纯朴的恩宠，并借不断克制心神，摆脱了骄傲的膨胀：为使\Quote{骄傲的脚}不再\Quote{攻击我}，\Quote{罪人的手不再搅扰我}，并使我不致因成功而自满受到更严重的损害，我须全力呼喊：\Quote{天主，求你快来拯救我；上主，求你速来扶助我。}我被无数种思想游荡和心绪不定所灼烧，无法收聚涣散的思绪，甚至无法不间断地倾吐祈祷而不受虚妄形象的干扰、谈话和行为的回忆所困，且我觉自己被如此干枯贫乏所束缚，以致感到不能生出任何属灵思想的果实：为使我能被解救脱离这可怜的心境——我无法以任何叹息呻吟自拔——我须确实呼喊：\Quote{天主，求你快来拯救我；上主，求你速来扶助我。}反之，我感觉到因圣神的眷顾，我获得了心志的坚定、思想的稳固、心灵的敏锐，连同一种不可言喻的喜乐与心灵的飞跃，在属灵情感的洋溢中，我因主忽然的光照，感知到那先前全然向我隐藏的最神圣思想的丰沛启示：为使我能更长久地停留在其中，我须经常而焦切地呼喊：\Quote{天主，求你快来拯救我；上主，求你速来扶助我。}被魔鬼的夜间恐怖所包围，我激动不安，被不洁神灵的显现所扰动，我生活的希望与救恩本身被恐惧的恐怖所剥夺。我逃向这短诵的安全避难所，全力呼喊：\Quote{天主，求你快来拯救我；上主，求你速来扶助我。}反之，当我因主的安慰而恢复，因他的来临而欢慰，感觉如同被无数天使环绕，以致我忽然敢于寻求冲突，向那些片刻前我惧之甚于死亡、连其触摸或接近都令我心身战栗者挑战：为使这勇气的力量因天主的恩宠更持久地存于我内，我须全力呼喊：\Quote{天主，求你快来拯救我；上主，求你速来扶助我。}因此，我们必须不停息地、持续地倾吐这短诵的祈祷：在逆境中，为得救赎；在顺境中，为得保存而不致骄傲。我告诉你，让这短诵的思想毫不间断地在你的胸中反复默念。无论你在做何事、担任何职、行何路，不要停止咏唱它。在你就寝、进食、以及满足最后的本性需求时，思念它。这思想在你心中或许成为你的拯救公式，不仅保护你免受魔鬼的一切攻击，也净化你脱离一切过犯与尘世的污秽，并引领你达至那不可见的、属天的默观，将你载至那极少数人经验过的、不可言喻的祈祷热火。愿你在仍思量这短诵时入睡，直到因不断使用它而被塑造，以至在睡梦中习惯重复它。若你醒来，愿它首先进入你的思想，愿它先于一切醒来的意念，愿它在你起身时令你跪下，然后从那里差遣你去往一切工作与事务，并终日跟随你。按立法者的吩咐，这应是你所思念的：\Quote{在家中和行在路上，}\BibleRef{申 6:7}睡时和醒时。这应写在你的口之门槛与门上，这应置于你房屋的墙壁和你心灵的内室，以致当你跪下祈祷时，这短诵是你跪着时的咏唱；当你起身去处理一切必要的生活事务时，它成为你站立时的恒久祈祷。\par

% structure: p0034 source=h2 target=subsection reason=relative-heading-rank; base=h2

\subsection{第十一章}

\Emph{论藉上述方法所能达到的祈祷之圆满。}\par

这便是那心神当始终持守的公式，直至因恒常使用及不断默想而坚强，它便抛却并拒绝一切思想之丰盈内容，局限于这一短句的贫乏，从而轻易抵达福音真福中居首的那一端：因为祂说：\BibleQuote{神贫的人是有福的，因为天国是他们的。}\BibleRef{玛 5:3}如此，藉此类贫穷而成为伟人般贫穷的人，必将实现先知的话：\Quote{贫穷困苦的人将赞美上主的名。}的确，有何种贫穷比这更大或更神圣呢？即那人自知毫无防卫与能力，每日乞求他人恩赐的援助，且意识到自己每一刻的生命与存留皆赖天主的眷顾，便不无理由地自称为主的乞丐，每日在祈祷中向祂呼喊：\Quote{但我贫穷困苦：主帮助了我。}于是，藉天主自身的启迪，他攀升至对祂的多样认识，并从此开始以更高深、更神圣的奥秘为养料，符合先知的话：\Quote{高山是鹿的避难所，岩石是刺猬的藏身处。}按我们所述之意，这十分适用，因为凡持守纯朴与无辜者，不伤害或冒犯任何人，反满足于自身简单状况，只图单纯自卫以免遭仇敌掠夺，便成为某种灵性刺猬，被那福音之石的永恒盾牌所保护，即藉回忆主的苦难并不断默想上述短句，逃避敌对仇敌的陷阱。关于这些灵性刺猬，我们在\BibleBook{}{箴言}中读到如下：\Quote{刺猬是柔弱之民，却在岩石中安家。}的确，有何比基督徒更柔弱？有何比隐修士更软弱？他不仅不许为所受的委屈报仇，甚至不许心中生起一丝微弱的恼怒。但谁若从此状况前进，不仅保有无辜的纯朴，更以明辨之德护身，便成为毒蛇的消灭者，将撒殚踏于脚下，并以其敏捷之心符合明智之鹿的形象，此人将以先知和宗徒们的山岭为食，即他们最高最深的奥秘。不断以此牧草滋养，他将汲取\BibleBook{}{圣咏}中一切思想，并以如此方式歌唱，以致他不仅以最深切的心情感动地诵念，仿佛它们不是圣咏作者的作品，而是他自己的言语和他自己的祈祷；他必视它们为针对自己，并认出它们的话语不仅昔日藉先知或在其身上成就，且每日在他自己身上成就并实行。因为那时圣经对我们更清楚地开启，仿佛其脉络和骨髓暴露无遗，因我们的经验不仅领悟其意义，实则预先领会；而经文的意义并非藉注释向我们揭示，而是藉实际证明。因为若我们亲身体验了每篇圣咏所歌唱和写作时的心境，我们便似其作者，预先领会其意义而非追随，即在我们真正认识话语之前便聚集其力量，回忆我们所遭遇之事，以及每日攻击中发生之事，当这些思想临于我们时，我们歌唱并记起所有因疏忽所致的、或因热忱所得的、或因天主眷顾所赐的、或因仇敌的引诱所失的、或因狡猾的遗忘所带走的、或因人性软弱所致的、或因无知所欺骗的。这一切情感我们都在\BibleBook{}{圣咏}中找到表达，以致如同在极其清晰的镜子中看见所发生之事，从而更清楚领悟，并如此以情感为师，抓住它，视之为不仅听闻，实则眼见，仿佛不是记在记忆中，而是植于事物本性中，我们从心底被感动，以致我们不是通过阅读文字，而是通过经验预先领会其意义。如此，我们的心神将达到那不朽坏的祈祷，在我们前次论述中，如主恩准，我们会议的大纲便已升至这一祈祷；它不仅不注视任何形象，而且实际上不用任何言语或发声；但心神目的燃起火焰，藉某种不可言喻的精神敏锐，经由心的出神而产生，而心神如此受影响，无需感官或任何可见物质，便以不可言喻的叹息和呻吟向天主倾吐出来。\par

% structure: p0037 source=h2 target=subsection reason=relative-heading-rank; base=h2

\subsection{第十二章}

\Emph{一个疑问：如何能持守灵性思想而不失去它们。}\par

杰尔马努斯说：我们认为您不仅向我们描述了所问的这灵性操练的体系，而且描述了圆满本身，且极为清晰明了。因为有什么比以如此简短的默想拥抱对天主的记忆，并藉专注于一句短句，摆脱一切可见之物的限制，将所有祈祷的思想概括于一个短词中更圆满、更崇高呢？因此我们恳请您解明仍存的一点：即我们如何能稳固持守您赐予我们作为公式的这一短句，以致我们既已藉天主的恩宠从世俗思想的琐事中解脱，也能稳固把握一切灵性思想。\par

% structure: p0040 source=h2 target=subsection reason=relative-heading-rank; base=h2

\subsection{第十三章}

\Emph{论思想的轻浮。}\par

因为当心灵领会了某篇圣咏中一段经文的意义时，这意义不知不觉便从心灵中溜走，而它又无知无虑地转向另一段圣经经文。当它开始独自思量这段经文时，尚未彻底探究，另一段经文的回忆又浮现，挤走了先前主题的思量。由此，它又被某个新思虑的进入转至别处，灵魂总是从一篇圣咏转向另一篇，从福音段落跳到书信段落来读，又从书信转到先知书，再从那里被带到某段灵修史书，如此在全本圣经中漫无目的地游荡，既不能自主地拒绝或抓住任何东西，也无法通过充分反思和审视来完成任何事，于是只成了灵性意义的触摸者或品尝者，而非其作者和拥有者。因此，心灵总是轻浮游荡，甚至在礼仪时间也被各种事物分心，犹如醉酒，无法妥善履行任何职责。例如，祈祷时，它回忆某篇圣咏或某段经文；咏唱时，它想的是圣咏正文内容之外的东西；诵读经文时，它却想着该做的事或已做的事。如此，它无法有纪律且恰当地接纳或拒绝任何事物，仿佛被随机的侵扰驱赶，既无力保留所喜好的，也无法逗留其上。因此，我们首先应当懂得如何妥善履行这些灵修职责，并牢牢持守你们赐给我们作为公式的这一特定经句，使我们的情感起伏不因自身的轻浮而波动，而是置于我们自己的掌控之下。\par

% structure: p0043 source=h2 target=subsection reason=relative-heading-rank; base=h2

\subsection{第十四章}

\Emph{回答：我们如何获得心灵的稳定或思想的稳定。}\par

依撒格说：虽然在我们先前关于祈祷特质的讨论中，我认为对这个问题已说得足够，但既然你们要求我再次重复，我就对心灵的坚定作一简短教导。使游移的心灵坚定不移的有三件事：守夜、默想和祈祷，持之以恒和持续专注将产生心灵的坚定稳固。但这只能通过首先以不懈的勤奋摆脱今世的一切牵挂和焦虑才能达成，而此勤奋应投入不为贪财而为隐修院圣用的工作，如此我们才能履行宗徒的命令：\BibleQuote{不断祈祷。} \BibleRef{得前 5:17} 因为那只在屈膝祈祷时才祈祷的人，祈祷得\Emph{太少}；而那在屈膝时仍被种种心灵游荡所分心的人，\Emph{从未}祈祷。因此，我们在祈祷时当成为怎样的人，就应当在祈祷前成为怎样的人。因为祈祷时，心灵无法不受先前状态的影响，它在祈祷中，要么被之前徘徊的思想带入天上之事，要么被拖入地上之事。\par

依撒格院长就这样向我们这些惊奇的听众进行了他关于祈祷特质的第二篇谈话，直到这里。我们非常钦佩他对上述经句（他作为初学者持守的大纲所赐予我们的）的教导，并渴望紧紧追随，因为我们原以为这会是简短易行的方法；但我们发现，它甚至比我们以前不受任何恒心束缚、在全本圣经中漫无目的地游荡于各种默想中的方法更难遵守。因此，可以肯定，无人因不识字而被拒于心灵的成全之外，也无人的粗鄙无知会妨碍获得心灵和思想的纯洁，这纯洁对所有人都是近在咫尺，只要他们藉着不断默想这节经句，将心灵的思想稳健地持守于天主。\par

\SourceRecord{Translated by C.S. Gibson.；Nicene and Post-Nicene Fathers, Second Series；Vol. 11.；Edited by Philip Schaff and Henry Wace.；Buffalo, NY: Christian Literature Publishing Co.,；1894.；<http://www.newadvent.org/fathers/350810.htm>.}

% structure: p0001 source=h1 target=section reason=page-title

\section{会议十一}

\Emph{圣\Person{若望·卡西安} 著}\par

院长\Person{赫肋孟}的第一场会议。论成全。\par

虽然许多圣人受你们榜样的教导，几乎无法企及你们成全的伟大，你们如同伟大的光体，以惊人的光芒在这世上闪耀；然而，圣兄弟们\Person{霍诺拉图斯}和\Person{欧凯里乌斯}，你们被那些杰出之人的伟大荣耀所激励，我们从他们接受了隐修制度的最初原则；其中一人主持着一个大型弟兄团体，希望他的会众——通过每日目睹你们圣善生活的榜样——能从那些教父的训导中受教；而另一人则渴望前往埃及，亲见这些圣人而获益，以便离开这如同被高卢寒气冻结的省份，像一只纯洁的斑鸠飞向那义德太阳所照耀且最接近之地，那里充满成熟的德行之果实。当然，我深厚爱德促使我做这件事：即考虑到一人的渴望和另一人的辛劳，我不应拒绝写作的危险和艰辛，哪怕只为给前者增添在其子女中的权威，并为后者免除如此危险旅程的必要性。此外，既然我们为可纪念的\Person{卡斯特}主教所写的十二卷《隐修院规章》，以及我们应圣\Person{赫拉迪乌斯}和\Person{莱昂提乌斯}主教的命令创作的在斯刻特旷野生活的教父的十次会议，未能满足你们的信德和热忱，如今为使我们的旅程缘由也为人所知，我认为应将我们最初在另一旷野所见的三位教父的七次会议，以同样风格写成并献给你们；在这些会议中，凡是我们先前作品中对成全问题或许解释隐晦或甚至遗漏的，都将得到补充。但若这仍不足以满足你们渴望的神圣干渴，另外七场会议——将寄给住在斯多沙德群岛的圣兄弟们——我想将满足你们的需求和热忱。\par

% structure: p0005 source=h2 target=subsection reason=relative-heading-rank; base=h2

\subsection{第一章}

\Emph{\Place{滕内苏斯}城的描述。}\par

当我们信德初萌后住在叙利亚的一所隐修院里，并渐渐长大，开始渴望更大的成全恩宠时，就决定立即前往埃及，甚至深入特拜德最偏远的旷野，拜访许多圣人——他们的荣耀和声誉已传遍各处——为的是即使不能效法他们，至少也要认识他们。于是我们经过一次非常漫长的航行，来到一个名为\Place{滕内苏斯}的埃及城镇。那里的居民要么被海包围，要么被咸水湖环绕，因此他们只从事商业，并靠海上贸易获得财富和生计，因为土地不给他们提供这些；以致当他们想建造房屋时，没有足够的土壤，除非从远处用船运来。\par

% structure: p0008 source=h2 target=subsection reason=relative-heading-rank; base=h2

\subsection{第二章}

\Emph{关于\Person{阿尔赫比乌斯}主教。}\par

当我们到达那里时，天主满足了我们的愿望，使那位最蒙福、最杰出的主教\Person{阿尔赫比乌斯}也恰好抵达。他是从独修士团体中被选出来，被任命为\Place{帕内菲西斯}城的主教；他毕生如此严格地持守独修志向，以致丝毫没有放松他从前谦卑的性情，也没有因加给他的尊荣而自满（因为他发誓自己并非因其胜任而被召担任此职，反而抱怨自己是不配地被迫离开隐修制度，因为尽管在其中度过了三十七年，却从未能达到如此崇高修道的纯洁要求）。他在上述的\Place{滕内苏斯}亲切而最和善地接待了我们——他因在那里选举一位主教的事务而来到此地——在听到我们渴望和意愿去请教埃及更偏远地区的圣父们之后，他说：\Quote{来吧，暂且看看那些住在我隐修院不远处的老人；他们服务的长度由他们弯曲的身体显示，而他们的圣德在他们外貌上闪耀，以致仅仅看到他们就能给见者带来极大的教训。从他们那里，你们不仅能通过言语，更通过他们圣善生活的实际榜样，学到我所痛惜失去、并因失去而无法给予你们的东西。但我想，若你们在寻求我所没有的福音珍珠时，我至少提供一个你们能方便获得它的地方，那么我的贫穷将因这份热忱而有所减轻。}\par

% structure: p0011 source=h2 target=subsection reason=relative-heading-rank; base=h2

\subsection{第三章}

\Emph{关于\Person{赫肋孟}、\Person{内斯特罗斯}和\Person{若瑟}居住的旷野的描述。}\par

于是他拿起手杖和行囊——那是当地所有隐修士出发旅行的惯例——亲自作为向导领我们前往他自己的城，即\Place{帕内菲西斯}。该城的土地，以及邻近地区的大部分（从前是一片极其富饶的土地，因为据说从那里供应皇家餐桌的一切），都因一场突如其来的地震使海水翻腾、溢出河岸而被淹盖；如此（几乎所有村庄都成了废墟），先前肥沃的土地变成了盐沼泽，以致你会认为圣咏中那首灵歌是对该地区的字面预言：\BibleQuote{祂使河流变为旷野，使水泉变为干地；使肥沃之地变为咸地，因为居民的邪恶。}在这些地区，许多这样矗立在高丘上的城镇，因洪水而被居民遗弃，变成了岛屿，为神圣的独修士提供了所渴望的独处之所。在这些独修士中，有三位老人——\Person{赫肋孟}、\Person{内斯特罗斯}和\Person{若瑟}——作为最资深的独修士而出类拔萃。\par

% structure: p0014 source=h2 target=subsection reason=relative-heading-rank; base=h2

\subsection{第四章}

\Emph{关于院长\Person{赫肋孟}和他对我们所求教导的推辞。}\par

于是有福的\Person{阿尔赫比乌斯}认为最好先带我们去见\Person{赫勒蒙}，因为他离他的隐修院较近，而且他在年龄上也比另外两人更年长：因为他已度过人生中的第一百个年头，仅在心志上充满活力，而他的脊背因年老和不断祈祷而弯曲，以至于他仿佛重归孩童时代，双手下垂、手掌贴地爬行。当我们同时凝视这人奇妙的面容和步伐时（因为他虽然所有的肢体都已衰败、如同死亡，但他丝毫没有失去从前严谨生活的严厉），我们谦卑地恳求言语和道理，并声明渴求灵性教导是我们前来的唯一原因。他深深叹息说：我能教导你们什么道理呢？我已年老体衰，松弛了从前的严谨，也丧失了说话的自信心？因为我怎敢教导自己所不做的，或指示他人那些我知道自己现在实践得软弱又冷淡的事理？因此，如今我已如此高龄，我不允许任何年轻人与我同住，唯恐他人因我的榜样而放松了严谨。因为教师的权威永远不会强大，除非他通过实际行动将其建立在听者的心中。\par

% structure: p0017 source=h2 target=subsection reason=relative-heading-rank; base=h2

\subsection{\Emph{第五章}}

\Emph{论我们对他推辞的回答。}\par

听到这话，我们感到不小的困惑，于是这样回答说：虽然此地的艰难以及独居生活本身——即使年富力强的人也几乎难以忍受——足以教导我们一切（并且，无需您言语，它们\Emph{确实}教导并使我们深受感动），但我们仍恳求您暂时放下沉默，以更相称的方式在我们心中植入那些原则，使我们能够——与其说通过模仿，不如说通过赞叹——拥抱我们在您身上所见的良善。因为，即使我们的冷淡为您所知，不配获得我们所求的，但至少如此长途跋涉的辛劳应因它而得到报偿，因为我们出于对您教导的渴望和对自身益处的向往，在\Place{贝特莱恒}隐修院初学后，急忙赶到了这里。\par

% structure: p0020 source=h2 target=subsection reason=relative-heading-rank; base=h2

\subsection{\Emph{第六章}}

\Emph{院长赫勒蒙的陈述：如何以三种方式克服过犯。}\par

于是有福的\Person{赫勒蒙}说：使人控制过犯的有三样事物，即：或是对地狱或现已加诸的法律的恐惧；或是对天国的盼望与渴望；或是对良善本身的好感与对德行的爱。因为我们读到，对邪恶的恐惧憎恶玷污：\BibleQuote{敬畏上主憎恨邪恶。}\BibleRef{箴 8:13} 盼望也阻挡所有过犯的攻击：因为\Quote{凡仰望祂的必不蒙羞。} 爱也不怕被罪恶毁灭，因为\BibleQuote{爱永不止息；}\BibleRef{格前 13:8} 又说：\BibleQuote{爱遮盖许多罪过。}\BibleRef{伯前 4:8} 因此，有福的宗徒将全部救恩的总和归于这三样德行的获得，说：\BibleQuote{现今存在的，有信德、望德、爱德，这三样。}\BibleRef{格前 13:13} 因为信德使我们因惧怕将来的审判和惩罚而躲避罪的污点；望德使我们的心思脱离现世事物，因期待天上的赏报而蔑视一切肉体的快乐；爱德则使我们心中炽烈地爱慕\Person{基督}和灵性良善的果实，并以完全的憎恨憎恶反对这一切的事物。这三样虽然似乎都指向同一目标（因为它们激励我们戒绝不当之事），但在其卓越的程度上相互之间却大不相同。因为前两者专属于那些在追求良善上尚未获得对德行之爱的人，而第三样则专属于天主以及那些已将天主的肖像和模样接受到自己内的人。因为唯独祂行善，不为恐惧，不为感谢或酬报所激发，而纯粹出于对良善的爱。因为，正如\Person{撒罗满}所说：\BibleQuote{上主为了自身创造了万物。}\BibleRef{箴 16:4} 因为祂在自己良善的荫庇下，将一切美善的丰盈赐予配得和不配得的人，祂不会因过犯而疲惫，也不被人间的罪所激动，因为祂永远保持为完美的良善，本性不变。\par

% structure: p0023 source=h2 target=subsection reason=relative-heading-rank; base=h2

\subsection{\Emph{第七章}}

\Emph{论我们如何攀登至爱的高峰，以及其中的恒久性。}\par

若有人追求成全，从最初我们恰当地称之为奴性的恐惧阶段（关于此阶段，经上说：\BibleQuote{你们做完了一切，要說：我们是无用的仆人\BibleRef{路 17:10}}）开始，就应前进，登上更高的希望之路——这希望不是与奴隶相比，而是与雇工相比，因为它期待报酬的支付；而且它似乎对罪过的赦免和惩罚的恐惧漠不关心，自知善行，虽看似期待应许的赏报，却未能达到那信赖父亲仁慈与慷慨的儿子的爱，这儿子毫不怀疑父亲所有的一切都属于他。那个挥霍者连同父亲的财产一起失去了儿子的名分，也不敢企及这爱，当他说：\BibleQuote{我不配再称为你的儿子\BibleRef{路 15:19}}；因为在吃了猪吃的豆荚（即令人作呕的罪恶食物）之后，他连这也得不到，当他\BibleQuote{醒悟过来\BibleRef{路 15:17}}，被一种有益的恐惧所征服，开始厌恶猪的不洁，惧怕饥饿的折磨；仿佛已成了仆人，他渴望雇工的身份，思虑报酬，说：\BibleQuote{我父亲有多少雇工，口粮有馀，我倒在这里饿死！我要起身到我父亲那里去，向他说：\InnerQuote{父亲！我得罪了天，也得罪了你，我不配再称为你的儿子，请把我当作你的一个雇工吧！}\BibleRef{路 15:18-19}}。但那位跑去迎接他的父亲，以比这些话更大的慈爱接纳了他谦卑的忏悔，并不满足于让他享受次要的东西，而是立即跨过两个阶段，恢复了他原先儿子的尊严。我们也应当立刻努力，借着不可分离的爱德恩宠，登上那第三阶段——儿子的阶段，这阶段相信父亲所有的一切都是自己的；这样我们才配领受天上父亲的肖像和模样，并能像真正的儿子那样说：\BibleQuote{凡父所有的，都是我的\BibleRef{若 16:15}}。蒙福的宗徒也论到我们说：\BibleQuote{一切都是你们的，无论是保禄，或是阿颇罗，或是刻法，或是世界，或是生命，或是死亡，或是现今的，或是将来的；一切都是你们的\BibleRef{格前 3:22}}。我们的救主也以这样的模样召叫我们，说：\BibleQuote{你们应当是成全的，如同你们的天父是成全的一样\BibleRef{玛 5:48}}。在这些人的身上，有时对良善的爱会受到中断，当灵魂的力量因某种冷淡、快乐或喜悦而松懈时，就会暂时失去对地狱的恐惧或对将来福乐的渴望。在这些状态中，确实有一个通向某种进步的阶段，它影响我们，使我们从对惩罚的恐惧或对赏报的希望开始避免罪恶，从而得以进入爱的阶段。因为经上说：\BibleQuote{在爱内没有恐惧，反之，圆满的爱把恐惧驱逐于外，因为恐惧含有惩罚；那恐惧的，在爱内还没有圆满。我们所以爱，因为天主先爱了我们\BibleRef{若一 4:18}}。只有当祂先爱了我们（为我们的救恩，别无他恩），我们也因祂自己的爱而爱祂（别无他因）时，我们才能登上那真正的成全。因此，我们必须尽最大努力，以完全热忱的意念，从这恐惧上升到希望，从希望上升到对天主的爱，以及对德性本身的爱，以便当我们稳步走向对良善本身的爱时，能够尽可能按人性所及，牢牢持守那良善。\par

% structure: p0026 source=h2 target=subsection reason=relative-heading-rank; base=h2

\subsection{第八章}

\Emph{那些因爱而远离罪恶的人，是何等卓越。}\par

因为，靠着对地狱的恐惧或对未来赏报的希望而扑灭内心罪火的人，与因感到天主的爱而憎恨罪恶本身和不洁、纯粹出于对贞洁的喜爱和渴望而持守贞洁美德、不期待未来应许的赏报、反而因认识当前的美好事物而欢喜、行事并非出于对惩罚的顾虑而是出于对美德之喜悦的人之间，存在着巨大的差别。因为后一种状态，即使所有人为的见证都不在场，也不会滥用犯罪的机会，也不会被秘密的思绪诱惑所败坏；相反，由于它将美德之爱本身保持在自己的骨髓里，不仅不让任何与之相悖的事物进入内心，而且以极度的憎恶憎恨它。因为一个人因某种当前美好的喜悦而憎恨罪恶和肉体的污秽，与另一个人通过默想未来的赏报而制止不当欲望，是两回事；害怕当前的损失与害怕未来的惩罚也是两回事。最后，为了善本身而不愿离弃善，比因害怕恶而拒绝同意恶，要伟大得多。因为前者是自愿的善，而后者则是被迫的，仿佛是因害怕惩罚或贪图赏报而从勉强的人那里强行逼出来的。因为一个因恐惧而远离罪之诱惑的人，一旦恐惧的障碍被除去，就会重新回到他所喜爱的事物中，因此不会在善中持续获得任何稳定性，也永远不会摆脱攻击而安息，因为他不会获得确定的、持久的贞洁和平。因为哪里有战争的纷扰，哪里就难免有受伤的危险。因为一个身处战斗之中的人，即使他是一个战士，并且通过勇敢作战频频给敌人造成致命伤，有时仍会被敌人的剑尖刺中。但是，一个已经击溃罪的攻击、现在享受着和平的保障、并已过渡到对美德之爱本身的人，将持续保持这种善的状态，因为他完全沉浸其中，因为他相信没有什么比失去内心的贞洁更糟糕的了。因为他认为没有什么比当前的纯洁更宝贵、更珍贵；对他来说，危险地偏离美德或罪恶的有毒污点是一种严厉的惩罚。对这样的人，我说，别人的在场不会给他的善行增添任何东西，独处也不会夺走任何东西；相反，因为他随时随地都带着他的良心作为他行为甚至思想的审判者，他将特别努力取悦它，因为他知道良心不能被欺骗或蒙蔽，而且他无法逃避它。\par

% structure: p0029 source=h2 target=subsection reason=relative-heading-rank; base=h2

\subsection{第九章}

\Emph{爱不仅使仆人成为儿子，而且还赋予天主的肖像和模样。}\par

若有人依靠天主的助佑而非自身的努力，得蒙恩赐达到这种境界，他便开始从奴仆的境况（其中有恐惧）和佣工式的希望之贪（此贪求所寻求的不是施予者的美善，而是报偿的酬劳）转向义子的名分，在那里不再有恐惧，也不再有贪求，而是那永不失败的爱恒久长存。主在责备某些人时，指出了这恐惧和爱各自所应得的：\BibleQuote{儿子认识自己的父亲，仆人畏惧自己的主人：我若是父亲，我的尊敬在哪里？我若是主人，我的敬畏在哪里？}\BibleRef{拉 1:6} 因为身为仆人者必须畏惧，因为\BibleQuote{明知主人的意愿，却做了当受鞭打的事，必受多番鞭打。}\BibleRef{路 12:47} 凡藉此爱达到天主的肖像和模样的人，如今必因善本身而喜悦于善，且以某种相似忍耐和温和的情感，从此不再为罪人的过失而发怒，反以怜悯和同情为他们的软弱祈求宽恕，且记着自己也曾长久被类似情欲的刺所折磨，直至因主的仁慈得救，便会感到对于那些迷失者，不应是愤怒，而是怜悯；在完全的平安中，他将向天主歌唱以下诗句：\BibleQuote{祢斩断我的锁链，我要向祢献上赞颂的祭献；}以及：\BibleQuote{若非主帮助了我，我的灵魂几乎已居于地狱。}当他怀着这种谦卑之心时，甚至能完成这福音的完美诫命：\BibleQuote{爱你们的仇敌，善待恨你们的人，为迫害你们和诽谤你们的人祈祷。}\BibleRef{玛 5:44} 如此，我们也将得蒙赐予那随之而来的赏报，使我们不仅带有天主的肖像和模样，甚至被称为儿子：\BibleQuote{好成为你们在天之父的儿子，}祂说，\BibleQuote{祂使太阳升起，光照恶人也光照善人，降雨给义人也给不义的人；}圣若望知道自己已达到这种情感，他说：\BibleQuote{这样，我们在审判的日子就可以放心大胆，因为祂怎样，我们在这世上也怎样。}\BibleRef{若一 4:17} 软弱脆弱的人性如何能像祂，除非常以平静的爱在心中对待善与恶、义人与不义之人，效法天主，并为了善本身而行善，从而确实达到天主儿女的真义子名分，圣宗徒也这样说：\BibleQuote{凡由天主生的，就不犯罪，因为天主的种子存留在他内，他不能犯罪，因为他是天主生的；}又说：\BibleQuote{我们知道，凡由天主生的，就不犯罪，而且由天主所生保守了他，那恶者不能触动他。}这应被了解为不是指所有种类的罪，而仅指大罪：若有人不愿从中解脱和洁净自己，宗徒在另一处告诉我们甚至不应为他祈祷，说：\BibleQuote{若有人看见自己的弟兄犯了不至于死的罪，就当祈求，天主必将生命赐给他，就是给那些犯不至于死的罪的人。有至于死的罪：我不说当为他祈求。}至于那些他称为不至于死的罪，连忠诚事奉基督的人，无论怎样谨慎自守，也不能免于这些罪，他说：\BibleQuote{若我们说我们没有罪，便是欺骗自己，真理也不在我们内；}又说：\BibleQuote{若我们说我们没有犯过罪，便是以祂为说谎者，祂的话不在我们内。}\BibleRef{若一 1:8} 因为任何一位圣人不可能不落入那些因言语、思想、无知、遗忘、必然、意愿和突然而来的微小过失中；这些过失虽与那称为至于死的罪大不相同，但仍不能免于过错和责罚。\par

% structure: p0032 source=h2 target=subsection reason=relative-heading-rank; base=h2

\subsection{第十章}

\Emph{爱德如何达到为仇敌祈祷的完美境界，以及用什么记号可以辨认尚未净化的心灵。}\par

因此，当一个人获得了我们所说的对善的爱和效法天主时，他便拥有了主的怜悯之心，也会为迫害他的人祈祷，同样说：\BibleQuote{父啊，宽恕他们，因为他们不知道他们做的是什么。}\BibleRef{路 23:34} 然而，一个尚未完全清除罪恶渣滓的灵魂，明显的记号是不以怜悯之情哀叹他人的过犯，而坚持法官的严厉责难：因为若没有那使法律得以圆满的要素，他怎能获得心灵的完美？宗徒指出：\BibleQuote{你们要彼此分担重担，这样就满了基督的法律，}\BibleRef{迦 6:2} 而他没有那爱德，这爱德\BibleQuote{不动怒，不夸张，不图谋恶事，}它\BibleQuote{凡事忍耐，凡事包容。}\BibleRef{格前 13:4} 因为\BibleQuote{义人怜惜自己牲畜的生命；恶人的心却毫无怜悯。}所以，隐修士若以冷酷无情的严厉谴责他人的过错，自己也必定会陷入同样的罪中，因为\BibleQuote{严厉的君王必陷于灾祸，}而\BibleQuote{塞耳不听弱者哀求的，他自己呼求时，也无人应允。}\par

% structure: p0035 source=h2 target=subsection reason=relative-heading-rank; base=h2

\subsection{第十一章}

\Emph{一个问题：为什么他称恐惧和希望的情感为不完美。}\par

格曼诺：你确实强而有力且恢宏地论及了天主完美的爱德。但仍有件事困扰我们：即当你以如此赞美高举爱德时，却说敬畏天主和对永赏的希望是不成全的，尽管先知似乎对它们持有截然不同的看法，他说：\Quote{你们敬畏上主，你们的一切圣民，因为敬畏祂的人一无所缺。} 又，在遵行天主的公义行为上，他承认自己是出于对赏报的考虑而行的，说：\Quote{我倾心行祢的律例，直到永远，因有赏报。} 宗徒也说：\Quote{因着信德，梅瑟长大以后，拒绝被称为法郎公主的儿子；宁愿与天主的百姓一同受苦，也不愿暂时享受罪中之乐；他以为基督所受的凌辱比埃及的宝物更宝贵，因为他注目于赏报。} \BibleRef{希 11:24} 那么，我们怎能认为它们是不成全的呢？既然有福的达味夸耀他是在盼望酬报中遵行了天主的公义行为，而颁赐法律者据说也寄望于未来的赏报，因此轻视了王室的收养，并将最可怕的苦难以作比埃及的宝物更可取？\par

% structure: p0038 source=h2 target=subsection reason=relative-heading-rank; base=h2

\subsection{第十二章}

\Emph{关于不同种类成全的答复。}\par

赫隆：按照每一个心灵的状况和尺度，圣经召唤我们的自由意志走向不同等级的成全。因为不能向所有人提供同一的成全冠冕，因为众人并非都具有同样的德行、意向或热忱，因此天主圣言在某种程度上为成全本身指定了不同的等级和不同的尺度。这一点在福音中真福品类的多样性上清楚显示。因为虽然那些拥有天国的人被称为有福的，那些将拥有大地的人被称为有福的，那些将接受安慰的人被称为有福的，那些将获得饱饫的人被称为有福的，但我们相信，天国的居所与（不论怎样的）大地的拥有之间，以及接受安慰与义德的饱饫满足之间，有巨大的差异；那些获得怜悯的人与那些被堪当享受天主最光荣的显见的人之间也有巨大的区别。\Quote{因为太阳有一种光辉，月亮有另一种光辉，星辰有另一种光辉；而且星辰与星辰的光辉也有区别；死人的复活也是如此。} \BibleRef{格前 15:41} 因此，按照这一规则，圣经赞美敬畏天主的人，并说\Quote{凡敬畏上主的人，真是有福的，} 并为此应许他们圆满的福乐；然而又说：\Quote{爱内没有恐惧，反而圆满的爱德驱除恐惧，因为恐惧含有惩罚；那恐惧的人，在爱内还不成全。} \BibleRef{若一 4:18} 并且，虽然事奉天主是伟大的事，经上说：\Quote{当以敬畏事奉上主；} 又说：\Quote{你被称为我的仆人，这是一件大事；} 又说：\Quote{那仆人在主人来时，发现他如此行事，便是有福的，} 然而却对宗徒们说：\Quote{我不再称你们为仆人，因为仆人不知道他主人所做的事；我称你们为朋友，因为我从父那里听到的一切，我都告诉了你们。} \BibleRef{若 15:14} 又说：\Quote{你们如果作我所命令你们的事，就是我的朋友。} \BibleRef{若 15:13} 由此可见，存在着不同的成全阶段，主如此召唤我们由高处走向更高处，以致那在敬畏天主中已成为有福和成全的人，如经上所记\Quote{从力量到力量}，从一种成全走向另一种成全，即以心灵的锐敏从恐惧攀升至希望，最终被召至那更蒙福的阶段，即爱德；那曾是\Quote{忠信而精明的仆人} \BibleRef{玛 24:45} 的，将进入友谊的同伴和义子的名分。因此，我们的言论也须按此含义来理解：并非说对那永久的惩罚或应许给圣人的蒙福赏报的考虑没有价值，而是因为，虽然它们有用并引领追求它们的人进入蒙福的最初开端，然而爱德，其中已有更圆满的信心和持久的喜乐，却将他们从奴仆式的恐惧和雇佣式的希望中移除，导向对天主的爱德，并将他们带入义子的名分，并使他们由成全者变得更成全。因为救主说，在祂父的家里有\Quote{许多住所，} \BibleRef{若 14:2} 并且虽然所有的星辰似乎都在天上，但太阳和月亮的光辉之间，以及晨星与其他星辰之间，却有极大的区别。因此，有福的宗徒不仅将爱德置于恐惧和希望之上，也置于那些被视为伟大而奇妙的恩赐之上，并指出爱德的道路比一切更卓越。因为在他列举完神恩与德行的清单后，想要描述其组成部分时，他如此开始：\Quote{但我还指示你们一条更高超的道路。我若能说人间的语言和天使的语言，并拥有先知之恩，而又明白一切奥秘和一切知识，并有全备的信德，甚至能移山，并将我所有的财产全施舍给人，且舍身被火焚烧，若没有爱德，于我毫无益处。} 由此可见，再没有比爱德更珍贵、更成全、更崇高，若我可这么说，更持久的东西了。因为\Quote{先知之恩终必消失，语言之恩终必停止，知识之恩终必消逝，} 但\Quote{爱德永不消逝，} 若没有爱德，不仅那些最卓越的恩赐种类，甚至殉道的荣光本身也将消逝。\par

% structure: p0041 source=h2 target=subsection reason=relative-heading-rank; base=h2

\subsection{第十三章}

\Emph{论出于最大爱德的恐惧。}\par

凡已在这成全之爱上站稳的人，必定会登上更高境界，达致那属于爱的更崇高的敬畏。这敬畏并非出自对惩罚的恐惧或对赏报的贪求，而是出于最大的爱。因此，子以热切的情感敬畏最宽厚的父亲，兄弟敬畏兄弟，朋友敬畏朋友，妻子敬畏丈夫；其间并不是惧怕他的责打或斥责，而只是唯恐伤害到他的爱；并且在每一言每一行中，总是以焦虑的爱惜小心翼翼，唯恐自己对所爱之人的热情有丝毫减退。有一位先知曾美妙地描述这敬畏的伟大，说：\Quote{智慧与知识是救恩的财富：敬畏上主是他的宝藏。}\BibleRef{依 33:6} 他再也无法更清楚地说明那敬畏的价值与意义，竟说我们救恩的财富——即在于对天主的真智慧与真知识——只能藉着敬畏上主来保存。因此，先知的话邀请的不是罪人，而是圣人们来拥有这敬畏，圣咏作者说：\Quote{上主的圣徒，你们当敬畏上主，因为敬畏祂的人一无所缺。}因为一个人若以此敬畏敬畏上主，他的成全必定一无所缺。因为显然地，若望宗徒论及那另一种惩罚性的敬畏时，曾说过：\Quote{那畏惧的，在爱内还未成全，因为畏惧含有惩罚。}\BibleRef{若一 4:18} 由此可见，这无所缺、且是智慧与知识之宝藏的敬畏，与那不完美的、被称为\Quote{智慧的开端}、且含有惩罚、因而从成全者心中被爱的圆满所驱逐的敬畏，二者之间有天壤之别。因为\Quote{爱内无恐惧，但圆满的爱驱逐恐惧。}\BibleRef{若一 4:18} 其实，若智慧的开端在于敬畏，那么成全除了在于基督的爱——它因包含属于成全之爱的敬畏而被称为智慧与知识的宝藏，而非开端——还能是什么呢？因此，敬畏有两个层次。一个是初学者的，即那些仍受轭与奴役性恐惧支配者的敬畏；关于此，我们读到：\Quote{仆人当敬畏他的主；}在福音中：\Quote{我不再称你们为仆人，因为仆人不知道他主人所做的事；}因此祂告诉我们：\Quote{仆人不能永远住在家里，但儿子永远住着。}因为祂是在教导我们从那惩罚性的敬畏过渡到爱的完全自由，以及天主的朋友和儿子的信赖。最后，蒙福的宗徒，因主爱的力量已超越了奴役性敬畏的阶段，轻视较低的事物，声明自己已由主以美好事物所充盈，\Quote{因为天主赐给我们的，不是怯懦之神，而是大能、爱德和健全之神。}\BibleRef{弟后 1:7} 那些为对天上父亲的完美爱火所燃、并藉神圣义子之恩已不再是仆人而是儿子的人，他这样对他们说：\Quote{因为你们没有领受奴役之神，以致再畏惧，而是领受了义子之神，使我们呼号：阿爸，父啊！}\BibleRef{罗 8:15} 先知论及同一种敬畏时，描述了七重神恩，根据降生成人的奥迹，这七重神恩确实降临在那位天主人身上：\Quote{上主的神将住在祂内：智慧和聪敏之神，超见和刚毅之神，知识和真虔敬之神；}最后，他特别加上这些话：\Quote{敬畏上主之神将充满祂。}\BibleRef{依 11:2} 在此，我们首先要注意，他没有像前几次那样说\Quote{敬畏之神将\Emph{住在祂内}}，而是说\Quote{敬畏上主之神将\Emph{充满祂}}。因为其丰富的伟大之处在于，一旦它以其力量抓住了人，它所占据的就不是一部分，而是整个心灵。这并非没有充分的理由。因为它与那\Quote{永不消亡}的爱紧密相连，不仅充满人，而且对它所开始占据的人，施行持久、不可分离且持续不断的占有；它不会像那被驱逐的敬畏有时那样，因任何短暂喜乐或欢愉的诱惑而减弱。这就是属于成全的敬畏，据说那位天主人——祂来不仅是为了救赎人类，也是为了赐给我们成全的模范和善的榜样——被这敬畏所充满。因为天主的真子\Quote{没有犯过罪，口中也从未出过诡诈，}\BibleRef{伯前 2:22} 不可能感受到那种惩罚性的奴役敬畏。\par

% structure: p0044 source=h2 target=subsection reason=relative-heading-rank; base=h2

\subsection{第十四章}

\Emph{关于完全贞洁的问题。}\par

杰尔马诺：现在您已结束关于完全贞洁的论述，我们还想更自由地问一下有关贞洁的终向。因为，我们毫不怀疑，那些爱的崇高境界——藉着它们，我们得以登上天主的肖像与模样——绝不可能没有完全的纯洁。但我们想知道，是否能够确保永久持有它，以致没有任何情欲的刺激会骚扰我们内心的宁静，从而我们得以在肉体寄居的时期免于这种肉欲之情，永不被兴奋之火点燃。\par

% structure: p0047 source=h2 target=subsection reason=relative-heading-rank; base=h2

\subsection{第十五章}

\Emph{对所请求的阐释的延迟。}\par

\Person{凯勒蒙}: 确实，不断学习并教导那使我们与主相连的爱，是至福和非凡善德的标志，以致对祂的默想，如圣咏作者所说，占据我们生命的所有昼夜，并通过不断反刍这天上的食粮，滋养我们那如饥似渴地追求正义的灵魂。但我们也必须按照救主的仁慈远见，为身体的食粮做些准备，免得我们在路上昏厥\BibleRef{玛 15:32}，因为\BibleQuote{精神固然刚强，但肉体却软弱。}\BibleRef{玛 26:41} 现在我们必须摄入少许食物，以便晚饭后，心智能更加专注，以便仔细探究你所想要的内容。\par

\SourceRecord{Translated by C.S. Gibson.；Nicene and Post-Nicene Fathers, Second Series；Vol. 11.；Edited by Philip Schaff and Henry Wace.；Buffalo, NY: Christian Literature Publishing Co.,；1894.；<http://www.newadvent.org/fathers/350811.htm>.}

% structure: p0001 source=h1 target=section reason=page-title

\section{第十三讲道集}

\Signature{圣若望·卡西安}\par

\Person{卡勒蒙}院长的第三次讲道集。论天主的保护。\par

% structure: p0004 source=h2 target=subsection reason=relative-heading-rank; base=h2

\subsection{第一章}

\Emph{引论。}\par

当我们在短暂的睡眠后返回参加晨祷，等候那位老人时，院长\Person{日耳曼努斯}因极大疑虑而困扰，因为在之前的讨论中，其力量激发了我们对此前尚不认识的贞洁的最大渴望，而那位有福的老人却仅凭增加一句话，就摧毁了人自身努力的价值，加上一句：即使人为了善果而竭尽全力，除非仅凭天主恩赐的恩惠而非自身劳苦的努力获得善，否则他无法成为善的主宰。当我们正为这个问题困惑时，有福的\Person{卡勒蒙}来到斗室，见我们低声议论某事，他便比往常更早地缩短了祈祷和圣咏的礼仪，问我们发生了什么事。\par

% structure: p0007 source=h2 target=subsection reason=relative-heading-rank; base=h2

\subsection{第二章}

\Emph{一个问题：为何善行的功绩不可归于行善者的努力？}\par

\Person{日耳曼努斯}说：我们几乎被昨晚讨论中描述的那杰出美德之伟大所排斥，以致无法相信它的可能性；同样，恕我直言，我们的努力之酬报，即完美的贞洁——它是通过自身努力的勤勉获得的——却不主要归于努力者的付出，这在我们看来是荒谬的。例如，当我们看见一个农夫以最大的苦心经营土地，却不对收成归功于他的努力，这是愚蠢的。\par

% structure: p0010 source=h2 target=subsection reason=relative-heading-rank; base=h2

\subsection{第三章}

\Emph{回答：没有天主的帮助，不仅不能成就完美的贞洁，连一切善事都无法成就。}\par

\Person{卡勒蒙}说：通过你提出的这个实例，我们可以更清楚地证明，没有天主的助佑，工人的努力一无所能。因为农夫即使在最细心地耕作土地之后，也不能立即将庄稼的收成和丰硕的果实归功于自己的努力，因为他发现这些常常是徒劳的，除非有了及时的雨水和平静温和的冬季相助；因此我们时常看见已经成熟、结出并完全成熟的果实，仿佛从那些正欲收获的人手中被夺走；他们持续而认真的努力对劳动者毫无用处，因为它们没有受到主助佑的引导。既然天主的美善不将这些丰富的收成赐给那些不勤奋耕作、不频繁犁地的懒散农夫，同样，整夜劳苦对工人也无用，除非主的仁慈使之顺利。然而，人的骄傲绝不应该试图把自己与天主的恩宠等同，或与它混淆，从而幻想着自己的努力是天主恩赐的原因，或者夸耀非常丰硕的果实是对自己努力功劳的回应。人应当以最审慎的斟酌思考并权衡这一事实：他凭自己的力量无法付出那些他热切用于求财的努力，除非主的保护和怜悯赐予他完成所有农业劳动的力量；他自己的意志和力量将毫无能力，除非天主的仁慈提供了完成它们的条件，因为它们有时因雨水过多或过少而失败。因为当主赐予牛力量、身体健康、完成所有工作的能力以及事业的成功时，人仍必须祈祷，免得圣经所说的\BibleQuote{铜天铁地}和\BibleQuote{蝗虫所剩下的，蛀虫吃了；蛀虫所剩下的，蝗虫吃了；蝗虫所剩下的，霉病毁坏了}降临于他。不仅在此，农夫在劳作中的努力需要天主的帮助，除非它阻止了意外的不幸，即使田地里充满了预期的丰收果实，人不仅被剥夺了他徒劳希望和期待的东西，而且实际上失去了他已经收获并储存在打谷场或粮仓里的丰盛果实。由此我们清楚推断，不仅我们行动的起始，而且良善思想的起始也来自天主，祂首先激励我们有良好的意志，并为我们提供机会去完成我们正确渴望的事：因为\BibleQuote{一切美好的赠予和一切完善的恩赐都是从上而来，来自光明之父} \BibleRef{雅 1:17}，祂在我们内开始、继续并完成善事，正如宗徒所说：\BibleQuote{那赐种子给撒种人的，必将提供食物给吃，并增多你们的种子，使你们义德的果实增加} \BibleRef{格后 9:10}。但我们需要卑微地追随每日吸引我们的天主恩宠，否则如果我们以\BibleQuote{硬着颈项}和（用圣经的话说）\BibleQuote{心与耳未受割礼} \BibleRef{宗 7:51}来抵抗，我们将配得听见耶肋米亚的话：\BibleQuote{跌倒的，岂不再起来？转离的，岂不转回？为何这耶路撒冷的子民执著地转离？他们硬着颈项，拒绝归回} \BibleRef{耶 8:4}。\par

% structure: p0013 source=h2 target=subsection reason=relative-heading-rank; base=h2

\subsection{第四章}

\Emph{一个异议：问外邦人怎能在没有天主恩宠的情况下拥有贞洁。}\par

\Person{格尔玛努斯}：对于这个解释——我们不能仓促地否认其卓越——似乎有一个困难，即它倾向于摧毁自由意志。因为我们看见许多外教人，天主的恩宠的协助肯定未曾赐予他们，却不仅在节俭和忍耐的美德上，而且在（更显著的是）贞洁的美德上超群出众，那么我们怎能认为他们的意志自由被掳获，并且这些美德是由天主的恩赐赋予他们的呢？尤其是，在追求今世的智慧，并且不仅对天主的恩宠，甚至对真天主的存在（正如我们通过阅读的历程和他人的教导所认识的那样）全然无知的情况下，据称他们凭自己的努力和苦修获得了最完美的贞洁纯洁呢？\par

% structure: p0016 source=h2 target=subsection reason=relative-heading-rank; base=h2

\subsection{第五章}

\Emph{关于哲学家们所谓贞洁的答复。}\par

\Person{赫雷蒙}：我很高兴，尽管你怀有探求真理的极大热忱，却提出了一些愚蠢的观点；因为通过你的这些异议，天主教信仰的价值似乎更能确立，并且，如果我可以这样说，更能被彻底探究。因为有哪个智者会作出如此自相矛盾的陈述：昨天还坚持贞洁属天的纯洁性甚至不可能因天主的恩宠赐予任何凡人，现在却主张它甚至被外教人凭自身力量获得呢？但是既然你确实，如我所说，出于求真的渴望提出了这些异议，那就请思考一下我们关于这些要点的立场。首先，我们当然不能认为哲学家们达到了我们所要求的灵魂的贞洁——我们被命令不仅淫乱，连不洁在我们中间都不可提及——他们拥有一种\OriginalTerm{部分的}{\GreekText{μερική}}贞洁，即肉身的节制，借此他们能约束自己免于肉欲的性行为；但内心的纯洁和持续的身体纯洁，我不敢说在行为上，甚至在思想上他们都无法达到。最后，他们中最著名的\Person{苏格拉底}——正如他们自己所评价的——也不羞于公开承认这一点。因为当一位通过相貌判断性格的人（\OriginalTerm{面相学者}{\GreekText{ψυσιογνώμοιν}}）注视着他，并说\Quote{\OriginalTerm{败坏少年者的眼睛}{\GreekText{ὄμματα} \GreekText{παιδ} \GreekText{εραστοῦ}}，即\InnerQuote{败坏少年者的眼睛}}，他的弟子们冲向他，把他带到他们的老师面前，想要为他报复这一侮辱。据说\Person{苏格拉底}用这些话平息了他们的愤慨：\Quote{\OriginalTerm{安静，朋友们}{\GreekText{παύσαοθε}, \GreekText{ἐταῖροι}}}。\par

\SourceRecord{Translated by C.S. Gibson.；Nicene and Post-Nicene Fathers, Second Series；Vol. 11.；Edited by Philip Schaff and Henry Wace.；Buffalo, NY: Christian Literature Publishing Co.,；1894.；<http://www.newadvent.org/fathers/350813.htm>.}

% structure: p0001 source=h1 target=section reason=page-title

\section{第十四次会谈}

\Emph{由圣若望·卡西安著}\par

院长内斯特罗斯第一次会谈。论灵性知识。\par

% structure: p0004 source=h2 target=subsection reason=relative-heading-rank; base=h2

\subsection{第一章}

\Emph{院长内斯特罗斯论修道者知识的话语。}\par

我们的诺言与进程之次序要求接着有院长内斯特罗斯的训导，他是一位在各方面杰出且极有知识的人。当他看到我们将圣经的一些部分背下并渴望理解它们时，便对我们说了以下的话。事实上，此世有许多不同种类的知识，因为其种类之多如同技术与科学一般。然而，尽管它们全都要么完全无用，要么只对此生的利益有用，但没有一种知识没有其自己的学习体系和方法的，藉此寻求者可以掌握它。如果那些技术有特定的规则来传播它们，那么我们的信仰的体系和表达，它倾向于默观无形奥秘的隐秘，并寻求永恒的赏报而非现世的利益，就更加依赖于固定的秩序与计划了。而此知识是双重的：首先，\OriginalTerm{实践知识}{\GreekText{πρακτική}}，即实践的，它透过道德改进和净化过失而实现；其次，\OriginalTerm{理论知识}{\GreekText{θεωρητική}}，它在于默观神圣事物和认识最神圣的思想。\par

% structure: p0007 source=h2 target=subsection reason=relative-heading-rank; base=h2

\subsection{第二章}

\Emph{论把握属灵事物的知识。}\par

凡愿意达到此理论知识的人，必须首先竭尽全力追求实践知识。因为无理论知识亦可获得实践知识，但无实践知识则绝不可能获得理论知识。因为有些阶段如此分明、安排有序，以便人的谦卑能攀至高处；若这些阶段按我们所说的次序依次连接，人便能达到若无第一步则无法飞升的高度。因此，那不清除罪的污秽却妄图看见天主的人，是徒劳的：\Quote{因为天主的圣神厌恶欺诈，不住在罪恶的身体内。} \BibleRef{智 1:4}\par

% structure: p0010 source=h2 target=subsection reason=relative-heading-rank; base=h2

\subsection{第三章}

\Emph{实践成全如何依赖双重系统。}\par

然而，此实践成全依赖于双重系统：其第一方法是认识一切过犯的本性及其治疗方式；其第二是发现德行的秩序，并以德行的成全塑造我们的心思，使之服从它们，并非如被迫而屈从于某种残暴的统治，而是因其自然之善而喜悦，并以此为乐，真正愉快地攀登那条陡峭狭窄的道路。因为若一个人既未成功认识自身过犯的本性，也未努力加以根除，他又如何能获得德行的知识——即实践训练的第二阶段——或属灵与属天事物的奥秘——它们存在于更高阶段的理论知识中？因为必然的推论是，未能超越较低境界的人无法攀登更高境界，而不能理解自身领悟范围内之事的人，更无法把握外在之事。但你们应当知道，我们必须以双重目的努力：既为驱逐恶习，也为获得德行。而我们并非从自己的猜测得出此点，而是由唯独知道祂工作力量和方法的那一位的话所教导：\Quote{看，}祂说：\Quote{我今天派你管理列邦和万国，要拔出、拆毁、毁灭、推翻、建造和栽种。} \BibleRef{耶 1:10} 祂指出，为除去有害之物需要四件事：即拔出、拆毁、毁灭、推翻；而为行善并获得正义之事，只需建造和栽种。由此完全显明，拔除并根除肉体与灵魂中根深蒂固的情欲比引入并栽种属灵的德行更为艰难。\par

% structure: p0013 source=h2 target=subsection reason=relative-heading-rank; base=h2

\subsection{第四章}

\Emph{实践生活如何分配在许多不同职业和兴趣之中。}\par

于是，此实践生活——如上所述基于双重系统——分配在许多不同职业和兴趣之中。因为有些人以隐修者的独处和心灵的纯洁为其全部目标，如同我们知晓古时的厄里亚和厄里叟，以及我们时代的真福安当和其他以同样目标追随的人，藉着独处的静默与天主紧密结合。有些人则将全部努力和兴趣放在弟兄们的团体生活和修院生活的谨慎照顾上；如我们记得近期的院长若望，他曾主持特穆伊斯城附近一所大隐修院，以及其他一些具有同样功绩的人，以宗徒的标记而卓越。有些人喜欢客舍的慈善服务和接待，古时的圣祖亚巴郎和罗特曾以此取悦主，近期则有真福玛卡略——一位极其谦和且有耐心的人，他主持亚历山大城的客舍，其方式被认为不亚于任何追求沙漠隐居的人。有些人选择照顾病人，有些人致力于为受压迫和受苦的人代祷，或专务教导，或向穷人施舍，并因他们的爱和善良而在杰出和有名望的人中昌盛。\par

% structure: p0016 source=h2 target=subsection reason=relative-heading-rank; base=h2

\subsection{第五章}

\Emph{论在所选择的道路上的恒心。}\par

因此，人人都应竭力、尽力追求在已开始的工作上达致成全，按照他所选择的路线，以及他所领受的恩宠；并且当他称赞并钦佩他人的德行时，不偏离他已一劳永逸选择的那条路线，因为他知道，正如宗徒所说，教会身体只有一个，但肢体有许多，并且它有\BibleQuote{恩赐各有不同，是按着赐给我们的恩宠：或是预言，应按信德的比例而行；或是服务，应专务服务；或是教导的，应专务教导；或是劝勉的，应专务劝勉；施与的，要慷慨；治理的，要殷勤；怜悯人的，要喜乐。}\BibleRef{罗 12:4}因为没有肢体可以要求其他肢体的职分，因为眼睛不能代替手的工作，鼻孔也不能代替耳朵的工作。所以，不都是宗徒，不都是先知，不都是教师，不都是有治病的恩赐，不都是说各种语言的，不都是解释的。\BibleRef{格前 12:28}\par

% structure: p0019 source=h2 target=subsection reason=relative-heading-rank; base=h2

\subsection{第六章}

\Emph{软弱者如何容易被动摇。}\par

因为那些尚未在自己已选定的路线中安定下来的人，常常当他们听到有人因不同的志趣和德行而受称赞时，就会被这些称赞如此激动，以致立即试图模仿他们的方法：而在这件事上，人的软弱注定要徒劳无功。因为同一个人同时擅长我上面列举的所有善行是不可能的。如果有人渴望同样地影响所有，他必定会落到这种地步：即当他追求所有时，他将不能彻底成功于任何一项，并且从这种改变和动摇中，他失去的将比获得的更多。因为人走向天主的途径有很多，所以每个人都应完成他一旦选定的那一条，永不改变他志向的路线，以便他能在自己所行的任何生活道路上达致完善。\par

% structure: p0022 source=h2 target=subsection reason=relative-heading-rank; base=h2

\subsection{第七章}

\Emph{一个有关贞洁的例子，教导我们并非所有的人都应羡慕所有的事。}\par

除了我们所说的一个因轻浮想从一种追求转到另一种追求的隐修士所遭受的那种损失之外，还有一种由此导致的死亡危险，因为有时一些别人做得正确的事却被另一些人错误地当作榜样，而一些对某些人结果良好之事，却发现对另一些人有害。例如，若有人想模仿那个人的善行——若望院长常提起这善行，并非为了模仿他，而只是出于对他的钦佩——因为有一个人穿着世俗的衣服来到上述长者那里，当他带来一些初收的庄稼时，发现那里有一个人被一个极其凶猛的魔鬼附身。这个魔鬼藐视若望院长的驱魔命令，并宣称绝不会听从他的命令离开所附的身体，然而此人的来到却吓坏了魔鬼，魔鬼以极其谦卑的呼出他的名字而离开了。这位长者对此人如此明显的恩宠大为惊讶，更加惊奇，因为他看到此人穿着世俗的衣服；于是开始仔细询问他的生活方式和追求。当他说自己生活在世俗中，并受婚姻之绑时，真福若望心里思量他德行和恩宠的伟大，更加仔细地探查他的生活方式。他声称自己是一个农夫，靠每日双手劳动寻求食物，除了在早晚去田间劳动和回家时，在没有先在教会为日常生活的食物感谢赐予食物的天主，他并不觉得自己有什么好处；而且，他从未使用过任何庄稼而不先向天主奉献初果和十分之一；他也从未驱牛越过他人的庄稼地界而不先给牛戴上口套，以免因他的疏忽使邻人遭受丝毫损失。当这些似乎不足以使若望院长认为他获得如此恩宠的原因时，他询问并调查了什么可能与如此恩宠的功德相连，由于尊重这样焦心的询问，他被迫承认，十二年前，当他想要宣发成为隐修士时，他被迫于父母之命而娶了一个妻子，而这位妻子至今无人知晓，他一直当作姐妹般保持童贞。当长者听到这话，钦佩之情深深打动了他，他当众宣布，那蔑视他本人的魔鬼不能忍受此人的临在，是很有理由的；对于此人的德行，他本人不仅在青春的热火中，甚至现在，也不敢冒丧失贞洁的危险去追求。尽管若望院长怀着极大的钦佩讲述这个故事，但他从未建议任何隐修士尝试这个计划，因为他知道，许多别人做得正确的事，对于模仿者却会陷入巨大危险，而且那些主作为特殊恩赐赋予少数人的事，不是所有人都能尝试的。\par

% structure: p0025 source=h2 target=subsection reason=relative-heading-rank; base=h2

\subsection{第八章}

\Emph{论属灵的知识。}\par

但回到我们论述所由起的知识之解释。正如我们以上所说，\Emph{实用}知识分散在众多学科和兴趣中，而\Emph{理论}知识分为两部分，即历史解释和灵性意义。因此，撒罗满在总结教会丰厚的恩宠时，也补充说：\Quote{因为凡与她在一起的，都穿着双层的衣裳。}但灵性知识有三种：\OriginalTerm{伦理诠释}{tropological}、\OriginalTerm{寓意诠释}{allegorical}、\OriginalTerm{奥秘诠释}{anagogical}；我们在\WorkTitle{箴言}中读到如下记载：\Quote{但你要按你心之宽广，以三种方式将这一切向你描述。}因此，历史包含对过去和可见事物的知识，正如宗徒这样重复说：\BibleQuote{因为经上记载：亚巴郎有两个儿子，一个由婢女所生，一个由自由妇人所生；但那生于婢女的是按血肉而生的，那生于自由妇人的却是因恩许而生的。}但属于寓意的是接下来的部分，因为实际发生的事据说预表了某种奥秘的形式：\BibleQuote{\InnerQuote{因为这两个，}他说，\InnerQuote{是两个盟约，一个来自西乃山，生子为奴，那就是哈加尔。原来西乃山是阿拉伯的一座山，它与现时的耶路撒冷相比，她和她的子女都处于为奴状态。}}但奥秘诠释从灵性奥秘上升到天上更崇高神圣的奥秘，宗徒接着用这些话：\BibleQuote{但上面的耶路撒冷是自由的，她是我们的母亲。因为经上记载：\InnerQuote{不生育、不生产的人，欢呼吧！未经产痛的人，高呼吧！因为被遗弃的比有丈夫的儿女更多。}} \BibleRef{迦 4:22}伦理诠释是道德解释，涉及生活的改善和实践教导，就好比我们透过这两个盟约理解实用和理论的教导，或者至少我们愿意将耶路撒冷或熙雍理解为人的灵魂，正如这样：\Quote{耶路撒冷啊，请赞美主；熙雍啊，请赞美你的天主。}因此，上述四种形象若我们愿意，可汇聚于一个主题，使同一个耶路撒冷可以有四种含义：历史地，作为犹太人的城；寓意地，作为基督的教会；奥秘地，作为天上的天主之城\Quote{她是我们众人的母亲}；伦理地，作为人的灵魂，它常受主以这名号予以赞美或责备。关于这四种解释，蒙福的宗徒如此说：\BibleQuote{但如今，弟兄们，我若到你们那里去，说方言，除非我将启示、知识、预言或教训说给你们，于你们有什么益处呢？} \BibleRef{格前 14:6}因为\Quote{启示}属于寓意，借此隐藏于历史叙事之下的内容得以在灵性意义和解释中显明，例如，我们若试图解释：\BibleQuote{我们的祖先都在云下，都从云中海中受了洗归于梅瑟，并且他们\Quote{都吃了同样的神粮，都喝了同样的神饮，因为他们从跟随他们的磐石中喝，而那磐石就是基督。}} \BibleRef{格前 10:1}这个解释将我们每日领受的基督体血之形象相比，包含寓意意义。但同样由宗徒提及的知识，是伦理诠释，借此我们通过仔细研究，可以看到所有关乎实用辨识的事物是否有益和良善，就如在此情况中，我们被告知自己判断：\BibleQuote{女人是否适宜不蒙头祈祷天主？} \BibleRef{格前 11:13}这个体系，如前所述，包含道德意义。所以宗徒放在第三位的\Quote{预言}，指向奥秘诠释，借此话语应用于未来和不可见之事，就如这里：\BibleQuote{弟兄们，关于安眠的人，我们不愿意你们不知道，免得你们忧伤，像那些没有希望的人一样。因为我们若信耶稣死了，也复活了，那么，那些在耶稣内安眠的人，天主也必将他们与耶稣一同带来。我们照主的话告诉你们这件事：我们活着存留到主来临的人，决不会在安眠的人以前。因为主必亲自从天降下，有呼喊声、天使长的声音和天主的号声；那在基督内死去的人要先复活。} \BibleRef{得前 4:12}在这种劝诫中，奥秘诠释的形象被提出。但\Quote{教训}展开历史解释的简单进程，其中不再包含更隐秘的意义，而只是由字句本身所宣告的：就如这段经文：\BibleQuote{我当日所领受又传给你们的，第一是：基督照经上所说，为我们的罪死了，而且埋葬了，又照经上所说，第三天复活了，并且显现给刻法；} \BibleRef{格前 15:3}以及：\BibleQuote{天主派遣了自己的儿子，生于女人，生于法律之下，为把在法律之下的人赎出来；} \BibleRef{迦 4:4}或这段：\BibleQuote{以色列！你要听：上主我们的天主是唯一的主。} \BibleRef{申 6:4}\par

% structure: p0028 source=h2 target=subsection reason=relative-heading-rank; base=h2

\subsection{第九章}

\Emph{如何从实用知识进向灵性知识}\par

因此，若你渴望获得属神知识的光明，不是为了虚妄的夸耀，而是为成为更好的人，你首先被那真福的渴望所激发，对此我们读到：\BibleQuote{心里洁净的人是有福的，因为他们要看见天主。}\BibleRef{玛 5:8} 如此你也可能获得天使对\Person{达尼尔}所说的那真福：\BibleQuote{智者必发光如穹苍的光辉；那使多人归义的人，必如星辰，直到永远。}又在另一位先知书中：\Quote{趁着还有时间，用知识的光照亮自己吧。}所以，要保持我在你身上所见到的那股阅读的勤勉，要竭力首先彻底掌握实践知识，即伦理知识。因为若没有它，我们前面所说的那种理论上的纯净就无法获得；只有那些因自己行为的卓越、而非因教导他们之人的言语而成全的人，才能在付出大量辛劳与努力之后，将其作为奖赏而获得。因为他们获得知识，并非借由默思法律，而是借由自己劳作的果实；因此他们与圣咏作者同唱：\Quote{从你的诫命我获得了明悟。}在战胜了所有私欲偏情之后，他们自信地说：\Quote{我要歌唱，并在\OriginalTerm{无玷的道路}{undefiled way}上理解。}因为那在纯洁心灵的路径上努力行走于无玷之路的人，当他咏唱圣咏时，便能理解所咏唱的话语。因此，若你愿意在心灵中为属神知识预备一座神圣的帐幕，就当涤除一切罪恶的污秽，并摆脱对此世的挂虑。因为一个灵魂，即使只稍微被世俗的烦恼所占据，要获得知识的恩赐、成为属神诠释的作者并勤于阅读圣书，是不可能的事。因此，首先要谨慎，尤其是你，\Person{若望}，因你较为年轻更需要留意我即将说的话——你要严加管束嘴唇，保持绝对静默，以免你阅读的热忱和志向的努力被虚妄的骄傲摧毁。因为这是迈向学问的第一步实践步骤：以热诚的心和缄默的嘴唇，领受所有长老的规条和意见，并殷勤地将它们珍藏于心中，努力去实践而非教导它们。因为从教导中，生出虚荣的危险傲慢；而从实践中，则结出属神知识的果实。因此，你决不可在长老们的集会中贸然发言，除非某种可能造成伤害的无知，或某种重要知晓的事情促使你发问；因为有些被虚荣冲昏头脑的人，假装发问，实则为了炫耀他们所完全掌握的知识。因为一个以获取人的赞誉为目的从事阅读追求的人，不可能获得真知识的恩赐作为赏报。因为被这种私欲的锁链束缚的人，也必然被其他缺点所捆绑，尤其是骄傲；因此，若他在实践与伦理知识面前受挫，必不能获得由此生发出的属神知识。所以，在一切事上要\BibleQuote{敏于听闻，迟于发言}\BibleRef{雅 1:19}，免得那时常记载于\Person{撒罗满}口中的话临到你：\Quote{你若见一个说话急促的人，要知道愚人比他更有希望。}并且不要妄图以言语教导他人你尚未以行动实践的事。因为我们的主以自身的榜样教导我们应遵循这个次序，正如论及祂所说的：\BibleQuote{耶稣所开始行和教导的一切}\BibleRef{宗 1:1}。那么，你要小心，不可在行动之前急于教导，以免被列入主在福音中对门徒所提到的那类人：\BibleQuote{凡他们吩咐你们的，你们要遵守实行；但不要效法他们的行为，因为他们只说不做。他们把沉重而难负的担子捆好，放在人的肩上，自己却连一个指头也不肯动。}\BibleRef{玛 23:3} 因为若有人\BibleQuote{违犯这些诫命中最小的一条，并且这样教训人，将在天国里称为最小的}\BibleRef{玛 5:19}，那么，那胆敢轻视许多更大诫命并这样教训人的人，理所当然不被视为天国里最小的，而应被视为在地狱惩罚中最大的。因此，你必须小心，不要被那些在议论上颇有技巧、口才流利、能优雅而充分地谈论自己喜爱之事物的人的榜样所引导而教导别人，他们被那些不懂得分辨属神知识真实力量与特质的人想象为拥有属神知识。因为拥有伶俐的口舌与优雅的言辞是一回事，而穿透天上言语的核心与精髓、以灵魂纯洁的眼睛凝视深邃而隐藏的奥迹，则是另一回事；因为后者靠人的学问或今世的条件无法获得，只能靠灵魂的纯洁，借由圣神的启迪获得。\par

% structure: p0031 source=h2 target=subsection reason=relative-heading-rank; base=h2

\subsection{第十章}

\Emph{如何拥抱真知识的体系。}\par

那么，你若想要获得圣经的真知识，首先当致力于获得内心坚定的谦卑，凭完美的爱德引导你，不是走向那使人自高自大的知识，而是走向那使人启迪的知识。因为不洁的头脑不可能获得属神知识的恩赐。因此，你要以最大的谨慎避免这一点，免得因你阅读的热忱，在你内产生的不是知识的明光，也不是因学习之光所应许的永恒荣耀，而是因虚妄的骄傲而成为毁灭你的工具。其次，你必须竭力摆脱一切焦虑和世俗思想，专心致志地，或不间断地投入神圣阅读，直到持续默想充满你的心，并按照它自己的肖像塑造你，在某种意义上使之成为约柜，\BibleRef{希 9:4}，柜内有两块石版，即两约的坚定保证；一个金罐，即纯洁无瑕的记忆，以恒常的持守保存其中所储存的玛纳，即属神意义永恒而属天的甜蜜和天使的食粮；此外还有亚郎的棍杖，即我们真正的大司祭耶稣基督的救恩旗帜，它永远发芽，蕴含着不朽记忆的新鲜活力。因为这根棍杖从叶瑟的根砍下后，死去而又复生，焕发更旺盛的生命。而这一切都由两个革鲁宾守护，即历史和属神知识的圆满。因为革鲁宾意味着丰盛的知识；它们不断守护天主的赎罪盖，即你内心的平安，并荫庇它免受一切邪恶势力的攻击。这样，你的灵魂不仅将前进到天主盟约的约柜，也将进入司祭王国，由于对纯洁之爱的不间断热爱，如同沉浸于属神的研究，将履行那赋予司祭的、由立法者所吩咐的命令：\BibleQuote{他不可离开圣所，免得玷污天主的圣所，}\BibleRef{肋 21:12}即他的心，主曾应许永远居住在其中，说：\BibleQuote{我要在他们内居住，并要在他们中间行走。}\BibleRef{格后 5:16}因此，应将全部圣经勤奋地铭记于心，并不断背诵。因为这持续不断的默想将带给我们双重的果实：第一，当心思专注于阅读和预备功课的时候，它不可能被任何恶念的罗网所俘虏；其次，那些我们反复背诵、反复记忆的经文，当我们试图记住它们时，由于当时心思被占据而未能理解，日后当我们摆脱一切行动和景象的纷扰，尤其在夜间静默默想时，就能更清晰地领悟。这样，当我们安息，仿佛沉睡于迷寐之中时，那些我们清醒时毫无概念的最隐秘含义的理解，就会向我们启示。\par

% structure: p0034 source=h2 target=subsection reason=relative-heading-rank; base=h2

\subsection{第十一章}

\Emph{论圣经的多重含义}\par

但借着这种研习，我们灵魂的更新成长时，圣经也将开始呈现新的面貌，神圣含义的美会随着我们的成长而增长。因为其形式适应人心智的能力，对属血肉的人显得属世，对属神的人显得属神，以致那些先前见它似乎被密云笼罩的人，不能领会其精妙，也无法承受其光明。但为使我们所追求的目标通过一个实例更加清楚，只举法律中的一个段落就足够了，借此我们可以证明，所有天上的诫命也是按我们状态的程度应用于人。因为法律上写着：\BibleQuote{不可奸淫。}\BibleRef{出 20:14} 这话按照字面的简单意义，被一个仍受制于卑劣情欲的人正确地遵守。但一个已经抛弃这些肮脏行为和不洁情感的人，则必须在精神上遵守它，使他不仅抛弃偶像崇拜，而且抛弃所有外邦迷信、占卜、征兆、所有迹象、日子和时辰的遵守，无论如何也不被那些破坏我们信仰纯朴的言语和名字的猜测所缠累。因为这样的奸淫，我们读到耶路撒冷被玷污了，她曾\BibleQuote{在各高山上和每棵绿树下}行奸淫，\BibleRef{耶 3:6} 主也曾藉先知责备她说：\BibleQuote{让观星象的站起来拯救你，他们观看星辰，数算月份，好从其中告诉你那将要临到你的事，}\BibleRef{依 47:13} 对于这种奸淫，主在别处责备他们时也说：\BibleQuote{行淫的神欺骗了他们，他们就离开他们的天主去行淫。}\BibleRef{欧 4:12} 但一个已经抛弃了这两种奸淫的人，将会有第三种要避免的，它包含在法律和犹太教的迷信中；关于这，宗徒说：\Quote{你们遵守日子、月份、节期和年份；}又说：不可摸，不可尝，不可碰。无疑这是指着法律的迷信说的，陷入其中的人确实已离开基督行淫，也不配听宗徒这样说：\BibleQuote{因为我已把你们许配给一个丈夫，把你们当作贞洁的童女献给基督。}\BibleRef{格后 11:2} 但以下的话将由同一宗徒的话对他说：\Quote{但我怕，正如蛇用狡猾欺骗了厄娃，你们的心意也会被败坏，失去那在基督耶稣内的纯朴。}但如果有人甚至逃脱了这种奸淫的污秽，还有第四种，就是通过与异端教导的淫乱交合所犯的。对此，蒙福的宗徒也说：\BibleQuote{我知道在我离开之后，将有凶暴的豺狼进入你们中间，不顾惜羊群，从你们中间也要起来说悖谬话的人，引诱门徒跟从他们。}\BibleRef{宗 20:29} 但如果一个人成功避免了这一种，他还得小心，以免因更微妙的罪陷入奸淫的罪过。我指的是那种由游荡思绪构成的奸淫，因为每一个不仅是可耻的、甚至是闲散的思想，以及无论多么微小的偏离天主，都被成全者视为最丑陋的奸淫。\par

% structure: p0037 source=h2 target=subsection reason=relative-heading-rank; base=h2

\subsection{第十二章}

\Emph{一个关于我们如何能达到忘却今世挂虑的问题。}\par

听到这话，我起初被一种隐秘的情感所动，然后深深叹息说：您所详尽阐述的这些事，使我比先前所承受的更加绝望了：因为，除了那些外部袭击软弱之人的普遍灵魂被掳之外，我所似乎已经获得的些许文学知识，又增添了特别的救恩障碍，我教师的努力，或我对持续阅读的专注，已经如此削弱了我，以致现在我的心中充满了那些诗人的诗歌，甚至在祈祷的时刻，还在想着那些无聊的寓言和战斗故事，这些从最早婴儿时期就通过童年的功课储存在我心里；因此，在咏唱圣咏或祈求罪赦时，要么有某些放荡的诗歌回忆闯入，要么战斗的英雄形象出现在眼前，这种幻象的想象总是欺骗我，不容我的灵魂升华到对至高之事的洞察，以至于藉着我每天的哀叹也无法摆脱它。\par

% structure: p0040 source=h2 target=subsection reason=relative-heading-rank; base=h2

\subsection{第十三章}

\Emph{关于我们如何能除去记忆中的渣滓的方法。}\par

\Person{内斯特罗斯}：从这一事实本身——它为您带来对您净化的极度绝望——反而可能产生快速而有效的补救，只要您愿意将您所说的在世俗研究中表现出的那种勤勉和热忱，转移到对圣神著作的阅读和默想上来。因为您的心灵必然会被那些诗歌所占据，直到它以同样的热忱和勤奋获得其他可供思考的事物，并代之以属神和神圣的事物，而不是无益的尘世事物。但当这些事物被完全彻底地领会，并且心灵被它们滋养之后，那么早先的思想就能逐渐被驱逐并彻底消除。因为人的心灵不可能被完全清空所有思想；因此，只要它不被属神的兴趣所占据，它就必然会沉浸于很久以前学到的东西。因为只要它没有东西可反复思考和不知疲倦地操练，它就必然会回到童年所学的东西，并不断思考那些通过长期使用和默想而吸收的东西。因此，为了使这种神性知识在您内持久稳固地得到加强，并使您不仅像那些只是通过他人复述而非自身努力接触它的人那样暂时享受它——如果我可以这样表达的话，在空中闻到它的气味——而是让它被安置在您的心中，并深深地铭刻在它上面，被透彻地看见和触摸，您最好尽最大努力确保：即使您也许在会谈中听到一些您非常熟悉的内容，也不要因为已经知道而轻蔑和鄙视地看待它们，而是以同样的热忱将它们放在心里——这种热忱正是我们渴望的救恩之言应当不断倾注到我们耳中或从我们口中发出的热忱。因为虽然圣言的讲述常常重复，但在一个渴求真知识的心灵中，饱足永远不会引起厌恶；相反，它每天如同接受某种新的、渴望的东西那样接受它，无论已经吸收了多少次，它都会更加热切地要么听要么说，并藉着这些重复而巩固已有的知识，而不是因频繁的会谈而感到任何厌倦。因为这是一个冷漠而骄傲的心灵的确凿标志：如果它轻蔑和粗心地接受救恩之言的药物，即使这药物是以过分坚持的热忱提供的。因为\BibleQuote{一个饱足的心灵轻蔑蜂房；但一个缺乏的心灵，即使微小的事物也显得甘甜。}\BibleRef{箴 27:7}因此，如果这些事物被仔细地吸收，储存在灵魂的深处，并以沉默的封印封存，那么后来，像某种使人心喜悦的香甜美酒，它们会因思想的古老和长久的忍耐而变得醇厚，从您心灵之罐的深处带着极大的香气被取出，又像某种长流不息的泉水，从经验的脉络和灌溉的德行渠道丰沛地涌流，从您心灵的深井中倾泻出丰富的溪流。因为您身上将发生箴言中对一位在工作中实现这一点的人所说的事：\BibleQuote{您要从自己的水池中饮水，并从自己井泉的源头取水。让您自己泉源的水向您丰盛地流淌，但让您的水流经您的街市。}\BibleRef{箴 5:15}并且根据先知依撒意亚：\BibleQuote{您将成为一座灌溉的花园，像一股永不干涸的水泉。在您内，荒废已久的处所将被重建；您将立起一代又一代的根基；您将被称为围墙的修补者，使路径转向安息。}\BibleRef{依 58:11}而且那位先知所应许的福乐将降临到您身上：\BibleQuote{上主必不再使您的导师远离您，您的眼睛必看见您的导师。您侧耳将听到背后告诫的声音说：\InnerQuote{这是正路，你们走在这路上，不要偏左或偏右。}}\BibleRef{依 30:20}因此，不仅您内心的一切目的和思想，而且您想象的一切游荡和漫游，都将成为您对神圣法律的一种神圣而不断的沉思。\par

% structure: p0043 source=h2 target=subsection reason=relative-heading-rank; base=h2

\subsection{第十四章}

\Emph{不洁的灵魂不能给予或接受神性知识。}\par

但正如我们所说，一个初学者不可能理解或教导这一点。因为如果一个人不能接受它，他怎能适合传给另一个人呢？但如果他胆敢教导任何关于这些事情的内容，那么他的话必定是空洞无用的，只能传到听者的耳朵，却无法触动他们的心，因为这些话纯粹是出于懒惰和不结果实的虚荣，并非来自良善良心的宝藏，而是来自夸夸其谈的虚妄无礼。因为一个不洁的灵魂（无论它多么热心地投入阅读）不可能获得神性知识。因为没有人会把珍贵的香油、精细的蜂蜜或任何宝贵的液体倒入肮脏发臭的容器。一个曾经充满恶臭的罐子，更容易破坏最甜美的没药，而不是从它那里接受任何甜蜜或恩宠，因为纯洁的东西被玷污比污秽的东西被洁净要快得多。因此，我们胸怀的容器，除非首先从一切罪恶的污秽斑点中被净化，否则将不配接受那有福的油膏——先知论及此油膏说：\Quote{像头上的油膏，流到亚郎的胡须上，又流到他衣领的边缘上}——也不能保持神性知识和圣经言语的纯洁，这些言语\Quote{比蜂蜜和蜂房更甘甜}。\BibleQuote{因为正义与不法有什么相通？光明与黑暗有什么联系？基督与贝利亚尔有什么和谐？}\BibleRef{格后 6:14}\par

% structure: p0046 source=h2 target=subsection reason=relative-heading-rank; base=h2

\subsection{第十五章}

\Emph{一个异议：因为许多不洁的人有知识，而圣人却没有。}\par

\Person{Germanus}：这个断言在我们看来既不是建立在真理上，也不是建立在可靠的推理上。因为如果显然所有从未接受基督的信德或以邪恶的异端之罪败坏它的人，都有不洁净的心，那么，许多犹太人、异端者以及陷于各种罪恶中的天主教徒，却获得了对圣经的完美知识，并夸耀其灵性学问的伟大；另一方面，无数圣善的人，他们的心已从一切罪的玷污中洁净，却满足于纯朴信德的虔诚，对更深知识的奥秘一无所知——这怎么解释呢？那么，那个将灵性知识完全归于心灵洁净的观点，又如何站得住呢？\par

% structure: p0049 source=h2 target=subsection reason=relative-heading-rank; base=h2

\subsection{第十六章}

\Emph{答案是：恶人不能拥有真知识。}\par

\Person{Nesteros}：不仔细权衡所发表见解中每一个词的人，无法正确发现断言的价值。因为我们首先说过，这种人仅拥有辩论的技巧和言语的装饰，却不能深入圣经的核心及其属神意义的奥秘。因为真正的知识只有通过真正的天主崇拜者才能获得；而这些人当然不拥有它，关于他们曾这样说：\Quote{听啊，愚昧的人民，你们没有心：你们有眼却看不见，有耳却听不见。}又说：\Quote{因为你抛弃了知识，我也必抛弃你，不让你作我的司祭。}因为正如经上所说，在基督内\BibleQuote{蕴藏着一切智慧和知识的宝藏}，\BibleRef{哥 2:3}那么，我们怎能认为，轻视寻找基督，或找到后以不敬之口亵渎祂，或至少以不洁行为玷污天主教信仰的人，获得了属神的知识呢？\BibleQuote{因为天主的圣神必远避诡诈，不居住于屈从罪恶的身体内。}\BibleRef{智 1:4}因此，除了先知所精妙描述的这一途径，别无他法达到属神的知识：\BibleQuote{你们要为自己播种正义，收割生命的希望；用知识之光光照自己。}\BibleRef{欧 10:12}首先我们必须播种正义，即通过正义的行为扩展实际成全；其次我们必须收割生命的希望，即通过驱逐肉身的罪恶，收集属神德行的果实；这样我们才能以知识之光成功地光照自己。圣咏作者也认为应遵循这一秩序，他说：\Quote{品行无缺、遵行上主法律的人，是有福的。寻求祂的证言的人，是有福的。}他并非先说：\Quote{寻求祂证言的人是有福的，}然后加上：\Quote{品行无缺的人是有福的；}而是以\Quote{品行无缺的人是有福的}开始；由此清楚表明，除非人首先通过实际生活，在基督的道路上品行无缺地行走，否则无人能正确地来寻求天主的证言。因此，你提到的那些人并不拥有那种不洁者无法获得的知识，而是\OriginalTerm{假知识}{\GreekText{ψευδώνυμον}}，即可称为伪知识的东西，蒙福的宗徒曾说：\BibleQuote{啊，\Person{弟茂德}，你要保管所受的寄托，远避亵渎的空谈和所谓知识的虚妄反对；}\BibleRef{弟前 6:20}这在希腊文中是\OriginalTerm{所谓知识的虚妄反对}{\GreekText{τάς} \GreekText{ἀντιθέσεις} \GreekText{τῆς} \GreekText{ψευδωνύμου} \GreekText{γνώσεως}}。对于那些似乎获得某种知识的外表或那些虽勤于阅读圣经并熟记经文，却不放弃肉身罪恶的人，箴言中说得很好：\BibleQuote{女人貌美而无见识，就如金环戴在猪鼻上。}\BibleRef{箴 11:22}因为人获得天上雄辩的装饰和圣经最宝贵的美丽，却又因执着于污秽的行为和思想，将其埋葬在最肮脏的土地中，或因自身情欲的污秽翻滚而玷污它，这于人何益呢？结果将是：那本为正当使用者之装饰的东西，不仅不能装饰他们，反而因着增加的污秽和泥浆而变得肮脏。因为\BibleQuote{罪人口里的赞美，实在不雅；}\BibleRef{德 15:9}于是先知对他说：\Quote{你为何要宣扬我的正义，你的嘴为何提及我的盟约？}对于这样的灵魂，他们从不会持久地拥有那经上所言的敬畏之德：\BibleQuote{敬畏上主是教诲和智慧，}\BibleRef{箴 15:33}却试图通过持续沉思来理解圣经的含义，箴言中恰当地问道：\BibleQuote{愚昧人手里拿着钱财，有什么用处呢？他不能买得智慧。}\BibleRef{箴 17:16}然而，这种真正的属神知识与被肉身罪恶的污点所玷污的世俗学问相距甚远，以至于我们知道，它有时在那些不具口才、几乎未受教育的人身上极其璀璨地兴盛。这一点在宗徒和许多圣人的事例中表现得很清楚，他们并未以虚荣的枝叶来炫耀，而是因着真正属神知识的果实之重量而俯首：关于他们，宗徒大事录中写道：\BibleQuote{他们看见\Person{伯多禄}和\Person{若望}的胆量，并得知他们是没有读过书的平常人，就十分惊讶。}\BibleRef{宗 4:13}因此，如果你渴望获得那永不消失的芬芳，你必须首先全力以赴，从主那里求得贞洁的纯洁。因为没有一个仍受肉身情欲，特别是邪淫之爱支配的人，能够获得属神的知识。因为\Quote{智慧安居于善心，}又说：\Quote{敬畏上主的人，必找到正义的知识。}但我们必须在已谈过的秩序中获得属神的知识，蒙福的宗徒也教导了我们。因为当他不仅想列出自己所有德行的清单，而且更想描述它们的秩序，以便说明哪个紧随哪个，哪个产生哪个时，他在列出某些德行后接着说：\BibleQuote{借着警醒、斋戒、贞洁、知识、坚忍、仁慈、圣神、无伪的爱德。}\BibleRef{格后 6:5}通过列举这些德行，他显然想教导我们，我们应从警醒和斋戒达到贞洁，从贞洁达到知识，从知识达到坚忍，从坚忍达到仁慈，从仁慈达到圣神，从圣神达到无伪爱德的赏报。因此，当你按照这一方法和秩序也达到了属神的知识时，你必定会像我们所说，拥有并非荒芜或空洞的学问，而是充满活力和果实的；由你播种在听众心中救恩的言语之种子，将因随之而来的圣神丰沛的雨水而得到灌溉；并且，按照先知所应许的，\BibleQuote{雨必降在你所播的种子上，你土地上所产的粮食必丰饶肥美。}\BibleRef{依 30:23}\par

% structure: p0052 source=h2 target=subsection reason=relative-heading-rank; base=h2

\subsection{第十七章}

\Emph{\OriginalTerm{成全之道}{method of perfection}应启示给谁}\par

也要小心，当你年岁成熟、引人教诲时，不要被\OriginalTerm{爱慕虚荣}{love of vainglory}引入歧途，而将这些你并非借由阅读，而是借由经验效果所学到的东西，随意教授给最不洁的人，从而遭受那最智慧的人所罗门谴责的事：\Quote{不要将恶人附于义人的牧场，也不要因肚腹饱足而被引入歧途。}因为\Quote{美味对愚人无益，在缺乏感官之处，智慧也无立足之地。因为愚昧更易被引，一个倔强的仆人不能因言语而改善，即使他理解，他也不会服从。}并且\Quote{不要在不谨慎的人耳边说什么，以免他嘲笑你的智慧言辞。}并且\BibleQuote{不要把圣物给狗，也不要把你们的珍珠投在猪前，以防它们用脚践踏，又转过来咬伤你们。}\BibleRef{玛 7:6}因此，将\OriginalTerm{灵性意义的奥秘}{mysteries of spiritual meanings}向这种人隐藏是对的，以便你能有效地歌唱：\Quote{我将你的话藏在心里，免得我得罪你。}但你会说：那么圣经的奥秘要分施给谁呢？最智慧的人所罗门将教导你：\BibleQuote{将烈酒给那些忧伤的人，将酒给那些痛苦的人，让他们忘记自己的贫穷，不再记念自己的痛苦。}\BibleRef{箴 31:6}即，对那些因自己过去行为的惩罚而忧伤悲痛压迫的人，要丰富地供以灵性知识的喜乐，如同\Quote{那使人心喜悦的酒}，并用\OriginalTerm{救恩之言}{word of salvation}的烈酒恢复他们，以免他们陷入持续的忧伤和致命的绝望，以致这类人\BibleQuote{被过度的忧伤所吞没。}\BibleRef{格后 2:7}至于那些仍处于冷淡和漠不关心、未受任何内心忧伤打击的人，我们读到以下的话：\BibleQuote{因为一个仁慈而无忧的人，必会缺乏。}\BibleRef{箴 14:23}因此要极其小心，避免因\OriginalTerm{爱慕虚荣}{love of vainglory}而自高自大，从而不能与先知所赞扬的那人共同分沾：\Quote{他没有将自己的钱财放债取利。}因为凡出于爱慕人的称赞而分施天主圣言的人（论到这圣言，经上说\Quote{主的言语是纯净的言语，如同银子在炉中炼过，除去地上杂质，精炼七次}），他就是将自己的钱财放债取利，为此他将不仅得不到赏报，反而应受惩罚。所以他选择用尽主人的钱财，为从暂时的利益中得益，而不是要主如经上所记载的\BibleQuote{当他来时，连本带利收回自己的钱。}\BibleRef{玛 25:27}\par

% structure: p0055 source=h2 target=subsection reason=relative-heading-rank; base=h2

\subsection{第十八章}

\Emph{论\OriginalTerm{神修学习}{spiritual learning}不结果的原因}\par

但可以肯定，神修之事教导无效有两个原因。要么是教师称颂他自己毫无经验的事，试图用空洞的言辞教导听者；要么是听者是个坏人，充满过错，无法在硬心中接受神修者神圣而救人的道理。关于这些人，先知说：\BibleQuote{因为这百姓的心是迟钝的，耳朵发沉，眼睛闭着；恐怕他们眼睛看见，耳朵听见，心里明白，回转过来，我就医治他们。}\BibleRef{依 6:10}\par

% structure: p0058 source=h2 target=subsection reason=relative-heading-rank; base=h2

\subsection{第十九章}

\Emph{甚至那些不配的人如何常能接受\OriginalTerm{救恩之言的恩宠}{grace of the saving word}}\par

但有时，在天主眷顾的慷慨中，\BibleQuote{祂愿意所有的人都得救，并得以认识真理，}\BibleRef{弟前 2:4}，恩赐给那未曾以无可指责的生活表明自己配得宣讲福音的人，使他获得神修教导的恩宠，以造福多人。至于主赐予医治的恩赐以驱逐魔鬼的方式，我们必须在类似的讨论中加以解释，但因为我们即将起身用晚餐，我们将留待晚间讨论，因为渐进而不过度耗费体力所领受的东西，总是更有效地被内心把握。\par

\SourceRecord{Translated by C.S. Gibson.；Nicene and Post-Nicene Fathers, Second Series；Vol. 11.；Edited by Philip Schaff and Henry Wace.；Buffalo, NY: Christian Literature Publishing Co.,；1894.；<http://www.newadvent.org/fathers/350814.htm>.}

% structure: p0001 source=h1 target=section reason=page-title

\section{第十五次会谈}

\Emph{圣若望·卡西安著}\par

阿爸内斯特罗斯的第二篇会谈。论天主的恩赐。\par

% structure: p0004 source=h2 target=subsection reason=relative-heading-rank; base=h2

\subsection{第一章}

\Emph{阿爸内斯特罗斯论恩赐的三重体系。}\par

晚祷后，我们照常一起坐在席上，准备听所许诺的叙述：我们因对长者的尊敬而沉默片刻后，他以这样的话打破了我们尊敬的沉默。先前我们谈话的顺序已使我们谈到属灵恩赐体系的阐释，我们从长者的传统中学到这是三重体系。第一重是为医治之故，当神迹的恩宠伴随着某些被拣选的义人，因他们圣德的功绩，正如显然宗徒们和许多圣人按照主的权柄行了神迹和奇事，主说：\BibleQuote{治好病人，复活死人，洁净癞病人，驱逐魔鬼：你们白白得来的，也要白白给人。}\BibleRef{玛 10:8}第二重是为了教会的建立，或由于带病人来的人的信德，或由于将被治愈之人的信德，健康的能力甚至从未自不配的人发出。救主在福音中说：\BibleQuote{在那一天，许多人要对我说：主啊，主啊，我们不是因你的名字说过预言，因你的名字驱逐魔鬼，因你的名字行了许多大能的事吗？那时我要向他们承认：我从来不认识你们，你们这些作恶的人，离开我吧！}\BibleRef{玛 7:22}另一方面，如果带病人来的人或病人的信德缺乏，就会阻止被赋予医治恩赐的人行使他们的医治能力。关于这一点，路加福音说：\BibleQuote{耶稣在那里不能行什么大能的事，因为他们的不信。}\BibleRef{谷 6:5}因此主自己说：\BibleQuote{在厄里叟先知的时候，以色列有许多癞病人，但他们中间没有一个得洁净的，只有叙利亚的纳阿曼。}\BibleRef{路 4:27}第三重医治的方法是魔鬼的欺骗和诡计模仿的，以便当一个人明显被罪奴役，却因他的奇迹被当作圣人和天主的仆人时，人们可能被说服效法他的罪，从而为诋毁打开入口，宗教的神圣可能被辱，或者，那相信自己拥有医治恩赐的人可能因心里的骄傲而自高自大，从而更严重地跌倒。因此，他们（魔鬼）呼求那些他们知道没有圣德功绩或任何属灵果实的人的名字，假装因他们的功绩而受扰并从所附的躯体中逃走。关于这一点，申命记说：\BibleQuote{如果你们中间兴起一位先知，或一个说自己见了梦的人，且预言一个征兆和奇迹，他所说的话应验了，并对你说：我们去跟随别的你们不认识的神，事奉他们吧：你们不要听那先知或那做梦者的话，因为上主你们的天主正在试探你们，要看看你们是否全心、全灵爱他。}\BibleRef{申 13:1}福音中说：\BibleQuote{将有许多假基督和假先知起来，显大神迹和大奇事，以致如果可能，连被选的人也要被迷惑。}\BibleRef{玛 24:24}\par

% structure: p0007 source=h2 target=subsection reason=relative-heading-rank; base=h2

\subsection{第二章}

\Emph{应在何处钦佩圣人。}\par

因此，我们绝不应该羡慕那些行这些事的人，即这些权能，反而要看他们在驱除所有罪行和改善品行上是否完美，因为这是赐给每个人，不是因别人的信德或各种原因，而是因他自己的热忱，借着天主恩宠的运作。因为这就是实践知识，宗徒又以另一名称称之为爱德，且借着宗徒的权威，它被置于众人和天使的各种语言，以及甚至能移山的充分信德确信，以及一切知识和预言，以及施舍所有财物，最后甚至殉道的光荣之上。因为他列举了各种恩赐后说：\BibleQuote{这一人因圣神得智慧之言，另一人因同一圣神得知识之言，另一人得信德，另一人得治病的恩赐，另一人得行奇迹，等等}\BibleRef{格前 12:8}；当他正要谈论爱德时，请注意他如何以几句话就将其置于所有恩赐之上：\BibleQuote{但我还要指示你们一条更高超的道路}\BibleRef{格前 12:31}。由此清楚表明，完美和真福的高度不在于行那些奇妙的工作，而在于爱德的纯洁。这并非没有理由。因为所有那些都要过去并消灭，而爱德却永存。因此，我们从未发现我们的父辈追求那些工作和神迹：相反，当他们借着圣神的恩宠拥有它们时，他们从不使用，除非极端和不可避免的需要迫使他们如此。\par

% structure: p0010 source=h2 target=subsection reason=relative-heading-rank; base=h2

\subsection{第三章}

\Emph{论阿爸玛卡里乌斯使一死人复活。}\par

我们也记得，一位死者曾被首批在\Place{赛特}旷野安家的\Person{玛加略}院长复活了。因为当一位追随\Person{欧诺弥}谬误的异端者，试图以诡辩精妙摧毁公教信德的纯朴，并已欺骗了许多人时，蒙福的\Person{玛加略}被一些公教徒恳求，他们因如此颠覆的恐怖而极为不安，要他解救全埃及的纯朴民众脱离不信的危险。为此他前来了。当异端者以诡辩术接近他，并想把他无知地拖入\Person{亚里士多德}的荆棘时，蒙福的玛加略以宗徒式的简洁止住了他的喋喋不休，说：\BibleQuote{天主的国不在于言词，而在于德能。} \BibleRef{格前 4:20}因此，让我们到坟墓那里去，让我们在我们找到的第一个死者身上呼号主的名，让我们，正如所记载的，\BibleQuote{用我们的行为将我们的信德显示出来，} \BibleRef{雅 2:14}好藉着祂的见证，显明正确信德的明显证据，并且我们不是以空洞的言词争论，而是以奇迹的德能和那不能受骗的判断，来证明清晰的真理。异端者听了这话，在在场的人们面前羞惭不堪，暂时假装同意提出的条件，并承诺次日再来。但次日，当所有人怀着更大的热忱聚集到约定地点，因渴望观看这场面时，他因自己缺乏信德的意识而恐惧，逃走了，并立即逃离了整个埃及。蒙福的玛加略与众人一起等候到第九时辰，见他因良心有愧而躲避自己，便带着那些已被异端者误导的民众，前往已选定的坟墓。在埃及，尼罗河的泛滥带来了这样的习俗：由于那国的整个宽度一年中有相当一部分时间被定期的洪水覆盖，如同大海一般，除了乘船通行外别无行动方式，死者的尸体经过防腐处理后，储存在较高的房间里。因为那地的土壤因持续潮湿而无法埋葬尸体。如果接受埋葬的尸体，就会因过多的洪水被迫将其抛到地面上。那时，蒙福的玛加略站定在一个最古老的尸体旁边，说：\Quote{人啊，如果那位异端者和丧亡之子曾与我同来这里，并且我站在他旁边，呼号并呼唤我天主基督的名，在你面前，这些几乎被他的诡计误导的人，你是否会复活？}于是他起来，以赞同的话语回答。然后玛加略院长问他，当他享受此生时，他从前是什么人，或者在人类的哪个时期生活过，或者他那时是否知道基督的名；他回答说，他生活在最古老的那些国王的时期，并声称那些日子里他从未听说过基督的名。玛加略院长再次对他说：\Quote{安睡罢，}他说，\Quote{与其他人一起安息在你们本位的秩序中，在终结时被基督唤醒。}那么，他内在的一切权能和恩宠，也许本来一直隐藏，若非全教省面临的危险，他对基督的完全献身和真诚的爱，迫使他行了这个奇迹。当然，并非炫耀荣耀，而是基督的爱和全体民众的利益，迫使他行了这奇迹。正如\WorkTitle{列王纪}中的段落向我们表明，蒙福的\Person{厄里亚}如此做了，他祈求火从天上降下在堆在柴堆上的祭品上，为的是解救全体民众的信德，这信德因假先知的诡计而面临危险。\par

% structure: p0013 source=h2 target=subsection reason=relative-heading-rank; base=h2

\subsection{第四章}

\Emph{亚巴郎院长在一位妇女乳房上行的奇迹}\par

何需提及亚巴郎院长的事迹呢？他被称为希腊语的\OriginalTerm{纯朴者}{\GreekText{ἁπλοῦς}}，由他生活的纯朴和纯真而得名。当他在五旬期离开旷野去埃及收割时，被一位妇女以眼泪和恳求纠缠，她带着一个因缺乏乳汁而憔悴半死的小孩；他给她一杯水，以十字圣号圣化了它；她喝下后，立刻奇迹般地，她那一直完全干涸的乳房流出了丰沛的乳汁。\par

% structure: p0016 source=h2 target=subsection reason=relative-heading-rank; base=h2

\subsection{第五章}

\Emph{同一位圣人对瘸子的治愈}\par

或者，当同一位圣人前往一个村庄时，被嘲笑的人群包围，他们讥讽他，并向他展示一个因膝盖抽搐多年失去行走能力的人，因长久虚弱而爬行；他们试探他说：\Quote{亚巴郎神父，请向我们显示，如果你是天主的仆人，就将这个人复原到以前的健康，好使我们相信你所敬拜的基督之名并非虚妄。}于是他立刻呼求基督的名，弯下腰抓住那人萎缩的脚，拉它。立刻，在他的触摸下，干枯弯曲的膝盖伸直了，他重新获得了在多年虚弱中遗忘如何使用双腿的能力，欢欢喜喜地走了。\par

% structure: p0019 source=h2 target=subsection reason=relative-heading-rank; base=h2

\subsection{第六章}

\Emph{不应以奇迹判断各人的功绩}\par

所以，这些人并不将施行如此奇迹的能力归于自己，因为他们承认，这些奇迹不是凭自己的功绩，而是靠主的怜悯完成的；他们借宗徒的话，拒绝了因他们的奇迹而发出的敬佩所带来的人世尊荣：\Quote{诸位同人，你们为什么惊奇这事？或者为什么注视我们，好像是我们凭自己的能力和虔诚使这人行走呢？}\BibleRef{宗 3:12}他们也不认为人应因天主的恩赐和奇事而闻名，而应因自己善行的果实而著称，这些果实是通过自己心灵的努力和行为的力量产生的。因为，正如上面所说，常常有心地败坏、在真理上可废弃的人，他们奉主的名驱逐魔鬼，并施行极大的奇迹。关于这些人，宗徒们曾抱怨说：\Quote{老师，我们看见一个人奉祢的名驱逐魔鬼，我们就禁止了他，因为他不与我们一同跟随祢。}然而基督当时回答他们说：\Quote{不要禁止他，因为不反对你们的，就是帮助你们的。}\BibleRef{路 9:49}尽管如此，当他们在末时说：\Quote{主啊，主啊，我们不是奉祢的名说过预言，奉祢的名驱逐魔鬼，奉祢的名行过许多异能吗？}祂将见证说：\Quote{我从来不认识你们，你们这些作恶的人，离开我去吧！}\BibleRef{玛 7:22}因此，祂甚至警告那些因自己的圣德而获得这奇迹和异能荣耀的人，免得他们因此自高自大，说：\Quote{不要因魔鬼服从了你们而欢喜，但要因你们的名字记录在天上而欢喜。}\BibleRef{路 10:20}\par

% structure: p0022 source=h2 target=subsection reason=relative-heading-rank; base=h2

\subsection{第七章}

\Emph{恩赐的卓越之处不在于奇迹，而在于谦卑。}\par

最后，一切奇迹和异能的创始者祂自己，当祂召叫门徒来学习祂的教导时，明确指出了那些真正的、特选的门徒主要应向祂学习什么，说：\Quote{你们来向我学习，}不是主要学习依靠天上的能力驱逐魔鬼，不是洁净癞病人，不是使瞎子复明，不是复活死人；因为即使我藉我的某些仆人行这些事，但人的地位不能插入天主的赞美，仆人和使者也绝不能在此为自己获得任何一份，因为那里只有天主性的荣耀。但你们要向我学习这一点，祂说：\Quote{因为我是良善心谦的。}\BibleRef{玛 11:28}因为这是所有人都可以学习并实践的，但施行奇迹和征兆并不总是必要的，也不是对所有人都有益，也不是赐予所有人的。因此，谦卑是一切德行之师，是天上建筑最稳固的基础，是救主特殊而辉煌的恩赐。因为凡能行基督所行的所有奇迹，却无自高自大危险的人，是那跟随温和的主，不是效法祂奇迹的伟大，而是效法祂忍耐和谦卑之德的人。但谁若企图命令邪魔，或赐予治病的恩赐，或在众人面前显示某种奇妙的奇迹，即使他在炫耀时呼求基督之名，他仍远离基督，因为他心中骄傲，没有跟随祂谦卑的老师。因为当祂将要回到父那里时，祂仿佛预备了遗嘱，留给门徒说：\Quote{我给你们一条新诫命：你们要彼此相爱；如同我爱了你们，你们也要照样彼此相爱。}接着祂立刻加上：\Quote{如果你们彼此相爱，众人因此就可认出你们是我的门徒。}\BibleRef{若 13:34}祂并没有说：\Quote{如果你们同样行征兆和奇迹，}而是说：\Quote{如果你们彼此相爱；}而这只有良善心谦的人才能持守，这是确定的。因此，我们的前辈从不认为那些在人前自称为驱魔者，并向惊叹的群众夸耀自己获得或声称拥有的恩宠的人是善修士，或免于虚荣之罪的人。但这是徒然的，因为\Quote{倚靠谎言的，是追风；那追逐飞鸟的，也是如此。}\BibleRef{箴 10:4}因为无疑，所罗门的箴言中所说的将要临到他们：\Quote{犹如风、云和雨显得清澈，同样，那些夸耀虚假恩赐的也是如此。}\BibleRef{箴 25:14}因此，若有人在我们在场时行这些事中的任何一件，他应得到我们的称赞，不是因对他的奇迹的钦佩，而是因他生活的美善，我们也不应问魔鬼是否服从他，而应问他是否拥有宗徒所描述的爱德特征。\par

% structure: p0025 source=h2 target=subsection reason=relative-heading-rank; base=h2

\subsection{第八章}

\Emph{从自身根除过错比从他人身上驱逐魔鬼更为奇妙。}\par

实际上，从自身的肉体根除淫欲的冲动，比从他人身体中驱逐邪魔是更大的奇迹；以忍耐之德抑制愤怒的强烈情感，比命令空中的权势是更非凡的征兆；从自己心中驱逐忧郁的吞噬之痛，比治愈他人的疾病和身体的热症是更伟大的事。总之，在许多方面，治愈自己灵魂的软弱比治愈他人身体的软弱，是更大的德行和更辉煌的成就。因为正如灵魂高于肉体，灵魂的得救也更为重要；灵魂的本性越珍贵卓越，其毁灭也越严重危险。\par

% structure: p0028 source=h2 target=subsection reason=relative-heading-rank; base=h2

\subsection{第九章}

\Emph{生活的正直比施行奇迹更为重要。}\par

关于那些治愈，曾对蒙福的宗徒们说：\Quote{不要因魔鬼服从了你们而欢喜。}\BibleRef{路 10:20}因为这不是凭他们自己的能力，而是凭所呼求之名的力量成就的。因此，他们受警告，不可因此声称自己有福或有荣耀，因为这纯粹是凭天主的权能和能力成就的，而只应因自己内心生命的纯洁而喜乐，因这纯洁使他们堪当名字记录在天上。\par

% structure: p0031 source=h2 target=subsection reason=relative-heading-rank; base=h2

\subsection{第十章}

\Emph{关于完美贞洁考验的启示。}\par

为了证明我们所说的，既借古人的见证，也借天主的圣言，最好用真福\Person{帕弗努提乌}本人之言及其经验，来说明他对奇迹的赞叹和贞洁恩宠的感受，或更确切地说，他从天使的启示中所学到的一切。因为此人多年因显著的克苦而闻名，以至于他自以为完全摆脱了肉欲的陷阱，因为他觉得自己已超越了他公开长期与之搏斗的魔鬼的一切攻击。当几位圣人来访时，他正为他们准备一种称为\OriginalTerm{阿特拉}{Athera}的扁豆粥，他的手不慎被炉中窜起的火焰烫伤。此事发生后，他深感羞愧，开始默默自忖：既然我与魔鬼更严重的争斗已经止息，为何火不与我平安相处？如果这外在的短暂火焰都不放过我，那么在那可畏的审判之日，那搜寻所有人功过的永不熄灭之火，又怎能放过我，不将我拘禁呢？正当他为此类思绪和烦恼所困扰时，突然睡意袭来，主的一位天使向他显现，说：\Quote{帕弗努提乌，你为什么因这地上的火尚未与你平安相处而烦恼呢？因为你的肢体中仍存留着一些未完全清除的肉欲骚动。只要这些根苗在你内茂盛，它们就不会让这物质之火与你平安相处。除非你通过这些证据发现你内的一切内在骚动都已消灭，否则你定无法感到它无害。去，拿一个裸体而最美丽的少女，如果你抱着她时，发现你心中的平安依然稳固，肉欲的热火在你内静止安宁，那么这可见火焰的接触也将温柔地、不加伤害地临到你身上，如同临到\Place{巴比伦}的三个青年一样。}于是这位长老被这启示所感动，却没有尝试这神圣指示给他的实验的危险，而是询问自己的良心，省察自己内心的纯洁；他猜想贞洁的分量尚不足以胜过这考验的力量，便说：难怪当我与不洁之灵交战时，仍感到那火焰的猛烈——我原以为这火焰不如魔鬼的残暴攻击重要。因为熄灭内在的肉欲，比藉主的标记和至高者的权能制服那从外冲来的恶鬼，或藉呼号神圣之名将它们从所附的身体驱逐出去，是更大的德行和更卓越的恩宠。至此，院长\Person{内斯特罗斯}结束了有关恩赐真实运作的叙述，陪同我们来到长老\Person{若瑟}的独修室——那里离他的独修室约有六英里远——因为我们渴望受教于他的训诲。\par

\SourceRecord{Translated by C.S. Gibson.；Nicene and Post-Nicene Fathers, Second Series；Vol. 11.；Edited by Philip Schaff and Henry Wace.；Buffalo, NY: Christian Literature Publishing Co.,；1894.；<http://www.newadvent.org/fathers/350815.htm>.}

% structure: p0001 source=h1 target=section reason=page-title

\section{第十六次讲论}

\Emph{圣若望·卡西安 著}\par

约瑟院长第一次讲论。论友谊。\par

% structure: p0004 source=h2 target=subsection reason=relative-heading-rank; base=h2

\subsection{第一章}

\Emph{约瑟院长首先问我们什么。}\par

有福的约瑟，其训诲和规诫现在即将陈述，他是我们在第一次讲论中提到的三位之一，出身于一个极为显赫的家族，是埃及一个名为\Place{特姆伊斯}的城市的主要官员，因此他不仅在埃及语方面，而且在希腊语的雄辩术方面都受过精心训练，以致他能以其本人的身份与我们或那些完全不懂埃及语的人进行令人钦佩的交谈，不像其他人那样通过翻译。当\Person{约瑟}发现我们渴望从他那里得到教导时，他首先问我们是否是亲兄弟，当他听说我们是在一种属灵的而非肉身的兄弟情谊中结合，并且从我们第一次开始舍弃世界起，我们一直以不断的纽带联结在一起，无论是在我们为属灵服务所进行的旅行中，还是在隐修院的追求中，他便开始如下讲论。\par

% structure: p0007 source=h2 target=subsection reason=relative-heading-rank; base=h2

\subsection{第二章}

\Emph{同一位长老关于不可靠的友谊的讲论。}\par

有许多种类的友谊和同伴关系，以极为不同的方式将人们结合在爱的纽带中。有些人的交往始于先前的推荐，先相识，甚至后来成为友谊。另一些人则因某种交易或接受和给予的协议而结合在爱的纽带中。还有一些人因事业、科学、艺术或学问的相似性和一致性而用友谊的锁链联结起来，即使是猛烈的心灵也因此彼此亲切相待，以至于那些在森林和山区以抢劫为乐、以人血为快的人，也会拥抱和珍惜他们罪行的同伙。但还有另一种爱，其结合来自本性的本能和血亲的法律，即同部落的人、妻子、父母、兄弟和儿女天生比其他人更受偏爱，我们发现这不仅是人类的情况，也是所有飞鸟和走兽的情况。因为在本能的驱使下，它们保护和捍卫自己的后代和幼雏，以至于常常为了它们而不惜使自己暴露于危险和死亡。确实，那些因难以忍受的凶猛或致命毒液而被与其他一切隔绝的走兽、蛇类和鸟类，如蛇怪、独角兽和兀鹫，虽然据说它们的面貌对每个人都是危险的，但彼此之间因同源和同类情感而保持和平无害。但我们看到，我们所说的所有这些爱，既然对善人和恶人、对走兽和蛇都是共同的，当然不能持久。因为地点的分隔常常打断并中止它们，同样，时间的流逝带来的遗忘、事务、商业和言语也会如此。因为它们通常是因为各种不同的联系，或是利益、欲望、亲属关系或商业，所以一旦出现任何分离的缘由，它们就会中断。\par

% structure: p0010 source=h2 target=subsection reason=relative-heading-rank; base=h2

\subsection{第三章}

\Emph{友谊如何不可分割。}\par

在所有这些中，有一种爱是不可分割的，这种结合不是由于推荐的恩惠、或某种大恩惠或礼物、或交易的理由、或本性的需要，而仅仅是由于德性的相似。我说，这就是不为任何偶然所破坏的，什么时间或空间的距离都不能分离或毁灭，甚至死亡本身也不能分开。这就是真实且不可打破的爱，它借着朋友的双重完美与良善而增长，一旦进入其纽带，没有趣味的不同、没有愿望的干扰性对立能够分离它。但我们知道许多立定此目标的人，虽然他们因对\Person{基督}炽热的爱而结伴在一起，却未能持续不断地保持它，因为虽然他们依靠一个良好的开端来建立友谊，却没有以同样一种热忱来保持他们所开始的意向，所以他们之间的爱只维持了一阵，因为它不是由双方的良善共同维持，而是由一方的耐心维持，因此尽管一方以不倦的英勇坚持着，但肯定会被另一方的狭隘所破坏。因为那些在寻求健全的完美境界上有些冷淡的人，他们的软弱，虽然强者可以耐心承受，但软弱者自己却不能容忍。因为他们内心已植入了不安的原因，不容他们平静，正如那些受身体软弱影响的人，通常把胃的娇弱和虚弱的健康归咎于厨子和仆人的疏忽，不论侍者如何细心地服侍他们，他们仍然把不适的原因归咎于那些身体健康的人，因为他们看不出这其实是由于他们自己健康不佳。因此，正如我们所说，确保无误且不可分割的友谊结合，其纽带只在于良善的相似。因为\BibleQuote{上主使人在家中同心合意。}\BibleRef{咏 68:7}因此，爱只会在那些有一致意向和心意去愿意和拒绝同样事物的人中持续不受打扰。如果你们也想保持这爱不破裂，就必须小心，先除掉自己的缺点，克制自己的欲望，以一致的热忱和意向，勤勉地实现那先知特别喜悦的事：\BibleQuote{看，兄弟们同居共处，多么快乐，多么幸福！}\BibleRef{咏 133:1}这应被理解为属灵的结合，而非地点的结合。因为性格和意向不同的人住在同一处是无用的，而基于同等良善的人被地点距离分开也无妨碍。因为在天主面前，是性格的结合，而非地点的结合，使弟兄们同住一处，凡是意志不同之处，平静的和平也不能维持。\par

% structure: p0013 source=h2 target=subsection reason=relative-heading-rank; base=h2

\subsection{第四章}

\Emph{一个疑问：即使是真正有益的事，若违兄弟之意，是否仍应施行？}\par

\Person{Germanus}：那么，若一方想要做某件他看作符合天主心意且有益的事，而另一方却不赞同，这事是应违兄弟之意而施行，还是应按他所愿而搁置？\par

% structure: p0016 source=h2 target=subsection reason=relative-heading-rank; base=h2

\subsection{第五章}

\Emph{回答：持久的友谊只能存在于成全者之间。}\par

\Person{Joseph}：因此我们曾说，完美且圆满的友谊恩宠只能存留在那些成全且同等良善的人中间，他们的同心同德和共同目的使他们从未，或至少极少，在关乎灵性进步的事上发生分歧或相左。但若他们开始因激烈争论而激动，那么显然他们从未按照我们上述的准则达到合一。然而，因为除了从根基开始的人外，无人能从成全起步，而你们的询问不是关乎成全的伟大，而是如何达到成全，我以为宜向你们略说明其准则以及你们脚步当循的道路，使你们能更容易获得忍耐与平安的福分。\par

% structure: p0019 source=h2 target=subsection reason=relative-heading-rank; base=h2

\subsection{第六章}

\Emph{借何种方式可保全合一不破。}\par

真友谊的首要基础在于轻视世俗财物并蔑视我们所拥有的一切。因为，若在摒弃了世俗及其一切虚幻之后，将剩余的可怜什物看得比最宝贵的弟兄之爱更重，那是完全错误且不可原谅的。第二，每个人要如此修剪自己的愿望，以致不把自己视为有智慧有经验的人，也不将自己的意见置于邻人的意见之上。第三，他当承认，一切事物，即使他认为有用且必要，也必须居于爱与平安的恩赐之后。第四，他当意识到，绝不应因任何理由（无论好坏）而动怒。第五，他当努力消除一个弟兄或许对他怀有的任何愤怒，无论多么不合理，如同消除自己的愤怒一样，深知他人的烦恼对他同样有害，如同他自己对他人动怒一样，除非他尽其所能从弟兄心中消除这烦恼。最后一点，无疑是普遍决定所有过犯的：即他当每日认识到自己将要离开这个世界；因为这样的认识不仅不容许任何烦恼滞留在心中，而且能抑制一切情欲和各样罪恶的冲动。凡掌握这一点的人，既不会遭受也不致引起苦毒的愤怒与不和。但若没有做到，一旦嫉妒爱者逐渐将烦恼的毒汁注入朋友心中，由于频繁的争吵，爱必将逐渐冷却，终有一天他会使长久激怒的恋爱者的心分离。因为若一个人行走在先前划定的路线上，他怎能与朋友不同呢？因为他若不为己主张什么，就完全断绝了争吵的首要原因（通常由琐碎和极不重要的事引发），他会尽力遵守我们在\WorkTitle{宗徒大事录}中读到的关于信徒合一的话：\BibleQuote{但众信徒都是一心一意，没有人说自己拥有的任何东西是自己的，而是他们凡物公用。} \BibleRef{宗 4:32} 那么，一个不是服侍自己的意愿而是服侍弟兄的意愿，并成为他的主和师傅的追随者的人，又怎能从之生出争论的种子呢？主和师傅曾以所取的人性说：\BibleQuote{我不是为执行我的旨意，而是为执行那派遣我者的旨意吗？} \BibleRef{若 6:38} 但一个决意不依赖自己的判断而依赖弟兄的裁决，不依靠自己的智慧和意图，而是按照弟兄的意愿赞同或反对自己的发现，并以虔诚的谦卑之心实践福音中这些话的人，又怎能激起争端的刺激呢？\BibleQuote{但不要照我的意愿，而照你的意愿。} \BibleRef{玛 26:39} 或者，一个认为没有什么比平安的恩赐更珍贵，并永不忘记主的这些话的人，又怎会接纳任何使弟兄忧愁的事呢？\BibleQuote{如果你们彼此相爱，众人因此就可认出你们是我的门徒。} \BibleRef{若 13:35} 因为藉此，基督愿意祂的羊群在这世上被认出，并可以说用这印记与其他一切分开。但他又凭什么能容忍自己心中接纳烦恼的积怨，或让这积怨留在别人心中呢？如果他坚信愤怒不可能有任何正当理由，因为愤怒是危险且错误的；并且当他的弟兄对他生气时，他不能祈祷，正如他自己对弟兄生气时一样；因为他常以谦卑的心记着我们的主和救主的话：\BibleQuote{所以，你若在祭台前献你的礼物时，在那里想起你的弟兄有什么怨你的事，就把你的礼物留在祭台前，先去与你的弟兄和好，然后再来献你的礼物。} \BibleRef{玛 5:23} 因为你宣称自己不生气，并相信自己遵守了命令——\BibleQuote{不可让太阳在你们的愤怒中落下}；以及 \BibleQuote{凡向弟兄发怒的，应受裁判}——是毫无用处的，如果你硬着心肠不顾别人的烦恼，而这烦恼你本可用你的善意予以平息。因为同样，你因违反主的命令而受罚。因为那说你不该向别人发怒的那一位，也说了你不该忽视别人的烦恼，因为在天主面前，无论你毁灭自己还是毁灭别人，都没有区别：\BibleQuote{祂愿意所有人得救}，\BibleRef{弟前 2:4}。因为任何人的死亡对天主都是同样的损失，同时对那以一切毁灭为乐者也是同样的收获，无论这毁灭是通过你的死亡还是你弟兄的死亡而造成的。最后，一个每日认识到自己即将离开这个世界的人，又怎能对弟兄保留丝毫的烦恼呢？\par

% structure: p0022 source=h2 target=subsection reason=relative-heading-rank; base=h2

\subsection{第七章}

\Emph{不应将任何事物置于爱之前，也不应置于愤怒之后。}\par

正如不应将任何事物置于爱之前，同样，也不应将任何事物置于暴怒和愤怒之下。因为一切事物，无论看似多么有用和必要，都当被轻视，以免引起扰乱之怒；一切我们以为不幸的事物，也当承当和忍受，以期爱与平安的宁静得以完整保存，因为我们应当认为没有什么比愤怒和烦恼更具破坏性，没有什么比爱更有益。\par

% structure: p0025 source=h2 target=subsection reason=relative-heading-rank; base=h2

\subsection{第八章}

\Emph{属神的人之间因何缘故会发生争论。}\par

因为我们的敌人借着突发怒气，因着一些琐碎和世俗之事，分离仍然软弱和属血肉的弟兄，同样，他也基于思想的差异，在属灵的人之间播下不和的种子，而宗徒所谴责的关于言语的争竞和争执，多半由此产生：因此我们恶毒和怀恨的敌人，在本来同心合意的弟兄之间播下不和的种子。因为智慧者撒罗满的话是真实的：\BibleQuote{争斗滋生仇恨，但友谊是众人的护卫。}\BibleRef{箴 10:12}\par

% structure: p0028 source=h2 target=subsection reason=relative-heading-rank; base=h2

\subsection{第九章}

\Emph{如何摆脱甚至属灵的不和起因。}\par

因此，为保存持久不懈的爱，仅仅移除了不和的第一起因——通常源于脆弱和世俗之事——或漠视一切属肉之事，并允许我们的弟兄无限制地分享我们所需的一切，是不够的，除非我们也同样切断第二个起因——通常假借属灵情感的外表出现；除非我们在一切事上获得谦卑的思想和和谐的心意。\par

% structure: p0031 source=h2 target=subsection reason=relative-heading-rank; base=h2

\subsection{第十章}

\Emph{论真理的最佳考验。}\par

因为我记得，当我年轻的时候，曾建议我找一个伙伴，这类思想常常混入我们的道德训练和圣经之中，以致我们幻想没有比这更真实或更合理的了：但当我们聚在一起，开始提出我们的想法时，在举行的一般讨论中，有些东西首先被其他人指出是错误的、危险的，然后立即被共同意见谴责和宣布为有害的；尽管此前它们似乎闪耀着魔鬼注入的光辉，以致它们很容易造成不和，若不是长上的训诫——像某种神圣的神谕一样被遵守——阻止了我们一切纷争；那训诫便是：他们几乎以法律的效力规定，如果我们任何人不想被魔鬼的诡计欺骗，就不应信赖自己的判断胜过兄弟的判断。\par

% structure: p0034 source=h2 target=subsection reason=relative-heading-rank; base=h2

\subsection{第十一章}

\Emph{一个信赖自己判断的人，如何不可能逃避魔鬼幻象的欺骗。}\par

因为常常证实，宗徒所说的话确实发生了。\BibleQuote{因为撒殚自身也伪装成光明的天使，}\BibleRef{格后 11:14}以致他诡诈地散布一种混乱和污秽的思想晦暗，取代真正知识的光明。除非这些思想被接受在一颗谦卑和温良的心中，并保留给一位更有经验的弟兄或被认可的长上考量，由他们的判断彻底审查后，被我们拒绝或接纳，我们必定会在思想中恭敬一位黑暗天使而非光明天使，并遭受严重的毁灭：这种伤害，任何信赖自己判断的人都不可能避免，除非他成为真正谦卑的爱好者和追随者，并以全心痛悔来实践宗徒特别祈求的：\BibleQuote{所以，如果你们在基督内有什么鼓励，有什么爱心的安慰，有什么圣神的交往，有什么慈悲和同情，你们就应使我喜乐圆满：你们要心思相同，爱德相同，同心合意，不做任何分争或虚荣的事；却要以谦卑彼此看作比自己高明；}\BibleQuote{论尊敬，彼此争先，}好使每个人看重同伴的知识和圣德，并认为真正明辨的最佳部分在于他人的判断，而非自己的。\par

% structure: p0037 source=h2 target=subsection reason=relative-heading-rank; base=h2

\subsection{第十二章}

\Emph{为何在会谈中不应轻视下级。}\par

因为常常发生，或因魔鬼的幻象，或因人为错误的出现（每个人在这生命中都有可能被欺骗），有时一个心智更敏锐、更有学问的人，头脑中会产生某种错误观念，而一个智力较迟钝、价值较低的人，反而更正确、更真实地把握了事情。因此，任何人，无论多么有学问，都不应在其空洞的虚荣中说服自己，他不需要与别人会谈。因为即使没有魔鬼的欺骗蒙蔽他的判断，他也无法避免骄傲和自负的有害陷阱。因为谁能毫无危险地将此归于自己呢？那蒙选的器皿，他声称基督亲自在他内说话，却宣称他上耶路撒冷去，纯粹只是为了这个原因：藉着主的启示和合作，就他所向外邦人宣讲的福音，与同是宗徒的弟兄们进行秘密讨论。由此事实我们被表明，我们不仅应当藉着这些诫命保持同心合意与和谐，而且我们不必害怕任何反对我们的魔鬼的诡计，或他幻象的陷阱。\par

% structure: p0040 source=h2 target=subsection reason=relative-heading-rank; base=h2

\subsection{第十三章}

\Emph{爱不仅属于天主，而且就是天主。}\par

最后，爱的德性被高举到如此地步，以致蒙福的宗徒若望宣称，它不仅属于天主，而且就是天主，说：\BibleQuote{天主是爱，那存留在爱内的，就存留在天主内，天主也存留在他内。}\BibleRef{若一 4:16}因为我们看到它是如此神圣，以致我们发现宗徒所说的话在我们内显然活生生地实现：\BibleQuote{因为天主的爱，藉着居住在我们内的圣神，已倾注在我们心中了。}\BibleRef{罗 5:5}这无异于说：天主的爱藉着居住在我们内的圣神，已倾注在我们心中；这位圣神，当我们不知道该如何祈求时，\BibleQuote{以不可言喻的叹息，为我们转求。但那洞察人心的，知道圣神的意愿是什么，因为祂按照天主，为圣徒祈求。}\BibleRef{罗 8:26}\par

% structure: p0043 source=h2 target=subsection reason=relative-heading-rank; base=h2

\subsection{第十四章}

\Emph{论爱的不同等级。}\par

所以，所有人都能够展示那称为\OriginalTerm{爱}{\GreekText{ἀγάπη}}的爱，蒙福的宗徒论及此爱说：\BibleQuote{因此，我们一有机会，就该向众人行善，尤其应向有同样信德的家人。}\BibleRef{迦 6:10} 这种爱应普遍地施予所有人，甚至主命令我们也应施予仇敌，因为祂说：\BibleQuote{爱你们的仇敌。}\BibleRef{玛 5:44} 但\OriginalTerm{倾向}{\GreekText{διάθεσις}}，即感情，却仅施予少数人，即那些与我们有相同性情或以善德相连的人；不过，感情似乎有许多不同程度的区别。因为我们爱父母的方式不同于爱妻子、爱兄弟、爱子女，这些情感的要求差别很大，父母对子女的爱也并非始终相等。正如圣祖雅各伯的例子所表明，他虽为十二子之父，并以父爱爱他们全体，却以更深的感情爱若瑟，圣经明确显示：\BibleQuote{但他的兄弟见他父亲爱他，就嫉妒他；}\BibleRef{创 37:4} 显然并不是那位善人父亲不爱其他子女，而是他在感情上更温柔、更纵容地依恋这一个，因为他是主的预象。我们也读到，这在若望宗徒身上表现得非常清楚，经上论及他说：\BibleQuote{那门徒是耶稣所爱的，}\BibleRef{若 13:23} 尽管主当然以祂特别的爱拥抱了祂同样拣选的其他十一位门徒，正如祂在福音中也以证言表明，祂说：\Quote{我怎样爱了你们，你们也该怎样彼此相爱；} 关于祂，别处也说：\Quote{祂既然爱了世上属于自己的人，就爱他们到底。} 但这偏爱一个人并不表示祂对其它门徒的爱有所冷淡，而只是对这一个的更加丰满和丰富的爱，这爱是由于他童贞的特权和肉身的纯洁所赋予他的。因此它被标记为特殊的对待，作为某种更崇高的事物，因为没有令人憎恶的与别人的比较，而是超额爱情的更丰盛的恩宠将其区分出来。我们在雅歌新娘的性格中也看到类似的事物，她说：\BibleQuote{请在我内把爱安排有序。}\BibleRef{歌 2:4} 因为这才是真正的、有序的爱：它虽不恨任何人，却因某些人的功德而更爱他们；它虽普遍爱所有人，却为自己选出一些能以特别感情拥抱的人；并且在这些特别且首要的爱慕对象中，又选出一些在感情上优于其他人的人。\par

% structure: p0046 source=h2 target=subsection reason=relative-heading-rank; base=h2

\subsection{第十五章}

\Emph{关于那些因隐藏而只会增加自己或弟兄的怨气的人。}\par

另一方面，我们知道（但愿我们不知道）有些弟兄如此固执顽强，以致当他们知道自己对弟兄动了怒气，或弟兄对他们动了怒气时，为了掩盖因对一方或另一方的怨愤而引起的内心烦恼，便离开那些本应通过谦卑接近和交谈来安抚的人；并开始吟唱圣咏中的几节经文。他们以为这样能软化心中升起的苦毒念头，却因傲慢的行为反而加剧了本可立即消除的怒气，如果它们愿意表现出更多的关怀和谦卑的话，因为适时的悔过本可治愈自己的情绪并软化弟兄的心。因为他们以那种计划滋养和助长卑劣或骄傲的罪，而非根除争吵的一切诱因，并忘记了主的命令，祂说：\BibleQuote{凡向弟兄发怒的，难免受审判；} 以及：\BibleQuote{若你记得你的弟兄有什么反对你的事，就把你的礼物留在祭台前，先去与你的弟兄和好，然后再来献你的礼物。}\BibleRef{玛 5:22}\par

% structure: p0049 source=h2 target=subsection reason=relative-heading-rank; base=h2

\subsection{第十六章}

\Emph{若我们的弟兄对我们有所不满，我们祈祷的礼物如何被主拒绝。}\par

所以，我们的主非常重视我们不可忽视他人的烦恼，以至于如果我们的弟兄有什么反对我们的事，祂就不接纳我们的奉献，即祂不允许我们向祂献上祈祷，除非我们通过迅速的补偿消除弟兄心中无论正当地或不正当地感到的烦恼。因为祂不是说：\Quote{如果你的弟兄有真正的理由抱怨你，就把你的礼物留在祭台前，先去与他和好；} 而是说：\Quote{若你记得你的弟兄有什么反对你的事，} 即，如果有一件无论多么微不足道或细小的事，因着它你弟兄的愤怒被激起，而这事突然记起回到你的脑海中，你就该知道，你不应献上你祈祷的精神礼物，直到你通过好意的补偿从弟兄心中消除了由任何原因引起的烦恼。既然福音的话命令我们向那些因过去且完全微不足道的争吵缘由而愤怒的人赔补，况且那些缘由是由于极小的事情引起的，那么，我们这些可怜的人将会怎样呢？我们以顽固的伪善忽视了最近的、至关重要的、且由我们自己的过错引起的冒犯；并被魔鬼自己的骄傲所膨胀，因耻于自卑，否认自己是弟兄烦恼的原因，并以反叛的精神不屑于服从主的命令，争辩说它们永远不应被遵守，也永远无法被实现。于是，随着我们断定祂命令了不可能和不适当的事，我们，用宗徒的话说，就成了\BibleQuote{不是法律的实行者，而是审判者。}\BibleRef{雅 4:11}\par

% structure: p0052 source=h2 target=subsection reason=relative-heading-rank; base=h2

\subsection{第十七章}

\Emph{关于那些认为忍耐应施予世俗人而非弟兄的人。}\par

这也应当令人痛心哀叹：即有些弟兄，当因一些责难之词而发怒时，若是有人以祈祷围攻他们，想要安抚他们，他们听到不应向弟兄容许或怀有恼怒，正如经上所载：\Quote{凡向弟兄发怒的，应受裁判；}以及：\BibleQuote{不可让太阳在你们含怒时落下，}\BibleRef{弗 4:26}时，便立即断言，若是一个外邦人或世俗之人说了或做了这事，便理当忍受。然而，谁能容忍一个弟兄犯下如此大的过犯，或用嘴唇说出如此无礼的责难呢！仿佛忍耐只应对外邦人和亵渎者显示，而不是对所有的人，或者仿佛愤怒若是对外邦人便是恶，但对弟兄却是善；而实际上，无论是针对谁，愤怒灵魂的顽固怒气对自身造成同样的伤害。但对他们来说，由于愚钝心智的麻木，不能辨别这些话的含义，这是多么可怕地顽固，甚至毫无理智啊！因为经上并没有说：\Quote{凡向陌生人发怒的，应受裁判，}——这或许可按照他们的解释，将那些与我们共享信仰和生活的人除外——但福音之言最明确地表达为：\Quote{凡向弟兄发怒的，应受裁判。}因此，尽管我们按真理的法则应视每一个人为弟兄，但在这段经文中，被称作弟兄的，是指一位信徒和共享我们生活方式的人，而非外邦人。\par

% structure: p0055 source=h2 target=subsection reason=relative-heading-rank; base=h2

\subsection{第十八章}

\Emph{论那些装作忍耐却以沉默激起弟兄愤怒的人}\par

然而，这是何等样的事：有时我们幻想自己耐心，因为受挑衅时不屑回答，却以阴郁的沉默或轻蔑的态度和手势嘲弄发怒的弟兄，以致我们沉默的神情更甚于愤怒的斥责会激怒他们，同时却以为自己在天主面前毫无罪过，因为我们没有脱口说出任何能在人的审判中标记或定我们罪的话。仿佛在天主眼中，只有言语，而非主要意志才被问责；仿佛只有罪恶的实际行为，而非意愿和意图才算作错；或者在审判中只问每人做了什么，而不问他意图做什么。因为不仅被激怒的愤怒性质是恶的，而且引发愤怒之人的意图也是恶的；因此，我们审判者的真实审查，将不是问争端如何挑起，而是问由谁的过错引起：因为必须考量的，是罪的意图，而非过犯被犯下的方式。因为一个人用刀亲手杀死弟兄，与用某种诡计逼他致死，有何区别？既然显然他是被其诡计和罪行所杀。仿佛亲手没有推瞎子下去就够了，尽管那本有能力救他、却在他跌倒即将落入沟壑时不屑救他的人同属有罪；或者仿佛只有用手抓住人的才有罪，而设下陷阱、或本可释放他却不肯的则无罪。因此，保持沉默并无益处，如果我们为此原因而强令自己沉默：借沉默达成原本呼喊所能达成的，装出某些姿态，使我们本该治愈的人越发愤怒，同时我们却因此被称赞，却致其受损和伤害——仿佛人不正因为此更加有罪，因为他试图从弟兄的堕落中为自己夺取荣耀？因为这样的沉默对两者同样有害：既增加他人心中的烦恼，又阻止它从自己心中消除。针对这样的人，先知诅咒有理地指向：\BibleQuote{祸哉，那给朋友喝酒，并掺入苦胆，使他醉倒，为要看见他的裸体的人。他充满了耻辱而非荣耀。}\BibleRef{哈 2:15} 还有另一位论及此类人的话：\BibleQuote{因为弟兄必欺骗弟兄，朋友必行诡诈。人必嘲笑自己的弟兄，不说真话；他们弯曲舌头如弓，为谎言而非真理。}\BibleRef{耶 9:4} 但常常虚假的忍耐比言语更强烈地激怒人，恶意的沉默超过最可怕的侮辱言语，敌人的创伤比嘲弄者的虚伪谄媚更易忍受。关于这一点，先知说得很好：\Quote{他们的言语比油更滑，却如利箭；} 别处又说：\Quote{狡猾人的话语温柔，却刺入腹中；} 以下这话也可恰当引用：\Quote{口对朋友说平安，却暗设网罗；} 然而欺骗者反被欺骗，因为\Quote{若人在朋友面前张设罗网，必缠住自己的脚；} 以及：\Quote{若人为邻舍掘坑，必自己掉进去。} 最后，当一大群人带着刀棒来捉拿主时，我们生命之主的凶手中，没有谁比那带头以假恭敬和问候、献上虚伪爱吻的人更残忍。主对他说：\BibleQuote{犹达斯，你用亲吻背叛人子吗？}\BibleRef{路 22:48} 即，你迫害和仇恨的苦毒，竟以这表达真爱甜蜜的举动为外衣。更明显、更有力地，主借着先知强调这种悲伤的力量，说：\Quote{倘若仇人辱骂我，我尚可忍受；倘若恨我的人向我夸口，我尚可躲避他。但却是你，与我同心之人，我的向导，我的密友；我们一起同食甘甜，在天主的殿中同心同行。}\par

% structure: p0058 source=h2 target=subsection reason=relative-heading-rank; base=h2

\subsection{第十九章}

\Emph{论因愤怒而禁食者。}\par

还有另一种恶性的烦恼，本不值得提及，唯因我们知道有些弟兄允许自己如此行：当他们被烦恼或激怒时，竟然坚持禁食不食，以致（我们羞于提及此事）那些平时声称自己不能将用餐延迟到第六时甚至第九时的人，一旦充满烦恼和愤怒时，即使禁食两天也不觉饥饿，并以这种禁食所导致的虚弱，借愤怒的过盛来支撑自己。在此他们显然犯有亵圣之罪，因为出于魔鬼自身的愤怒，他们忍受禁食，而这禁食本应专为谦卑心灵和净化罪恶而唯独献给天主；这如同向魔鬼而非向天主献上祈祷和祭献，从而配得聆听梅瑟的斥责：\BibleQuote{他们向魔鬼而非向天主献祭，向不认识的神明献祭。}\BibleRef{申 32:17}\par

% structure: p0061 source=h2 target=subsection reason=relative-heading-rank; base=h2

\subsection{第二十章}

\Emph{论有些人的虚假忍耐，他们提供另一边脸让人打。}\par

我们也不无知另一种癫狂，它在某些弟兄身上以虚假忍耐的伪装出现，就如在这种情形下，激起争端还不够，除非他们用激怒的言语挑衅以使自己挨打，而当他们被最轻微的打击碰触时，立即把身体的另一部分也献上挨打，仿佛他们能以这种方式完美地完成那条命令，它说：\BibleQuote{若有人打你的右颊，你把另一面也转给他；}\BibleRef{玛 5:39} 而他们完全忽略了该段落的含义和目的。因为他们幻想自己正在通过愤怒之罪实践福音的忍耐，而为了彻底根除愤怒，不仅禁止了报复的交换和争端的激怒，实际还命令我们通过忍受双倍的错待来减轻打人者的愤怒。\par

% structure: p0064 source=h2 target=subsection reason=relative-heading-rank; base=h2

\subsection{第二十一章}

一个疑问：如果我们服从基督的命令，怎能缺少福音的全德呢？\par

\Person{Germanus}：我们怎能责备一个满足福音命令的人，他不仅不报复，甚至准备承受双倍的错待？\par

% structure: p0067 source=h2 target=subsection reason=relative-heading-rank; base=h2

\subsection{第二十二章}

回答：基督不仅看行为，也看意愿。\par

\Person{Joseph}：正如稍前所说的，我们不仅要看所做的事，还要看心灵的性质和行者的意图。因此，如果你以内心的仔细审慎衡量各人所行的，并考虑它是怀着何种心意、从何种情感出发，你就会看到，忍耐和温良的德性不可能在相反的精神中实现，即不可能在急躁和暴怒中实现。因为我们的主救主在给我们彻底教导忍耐和温良的德行时（即教导我们不仅用嘴唇宣认，更要储存在灵魂的最深处），给了我们福音全德的概要，说：\BibleQuote{若有人打你的右颊，你把另一面也转给他}\BibleRef{玛 5:39}（无疑提到的\Emph{右}颊，因为另一个\Emph{右}颊除了在内在的人脸上，可以说是找不到的），祂借此希望从灵魂最深处彻底消除一切激怒的煽动，也就是说，如果你的外在右颊受到了打人者的打击，内在的人也谦卑地同意，献上自己的右颊挨打，与外在人的痛苦同情，并以某种方式将自己的身体向打人者的错待屈服并顺从，使内在的人甚至不在静默中外在人的打击所扰乱。那么你看，他们离福音的全德多么遥远，后者教导忍耐不仅要在言语中，更要在内心的内在安宁中持守，并命令我们无论遭遇什么邪恶都要保持忍耐，以便我们不仅始终使自己远离扰乱性的愤怒，而且通过屈服于他们的伤害，迫使那些因自身过错而扰乱的人，在受够打击后得以平静下来；从而以我们的温良战胜他们的暴怒。这样我们也将实现宗徒的这句话：\BibleQuote{不可为恶所胜，反要以善胜恶。}\BibleRef{罗 12:21} 显然，那些怀着如此精神和愤怒说出温良谦卑言语的人，不仅不能减弱已被点燃的怒火，反而在自己的情感和激怒弟兄的情感中使其燃烧得更猛烈，他们是无法实现这一点的。但这些人，即使他们能在某种程度上保持平静和安稳，也不会结出任何正义的果实，同时却以邻人的损失为自己赚取忍耐的光荣，从而完全远离那\BibleQuote{不求己益，}\BibleRef{格前 13:5} 而求他人益处的宗徒之爱。因为爱并不这样渴望财富，以致从邻人的损失中为自己获利，也不愿因他人的掠夺而获得任何东西。\par

% structure: p0070 source=h2 target=subsection reason=relative-heading-rank; base=h2

\subsection{第二十三章}

那让己意服从他人的人，如何才是强壮有力的人。\par

但你必定知道，通常那将自己的意志服从弟兄的人，比那被发现更固执地坚持己见的人，扮演着更强壮的角色。因为前者通过忍受和容忍邻人，获得了强壮有力的名声，而后者则获得了软弱有病的名声，需要被纵容和宠爱，以至于有时为了他的平安和安静，即使在必要的事上放松一些也是好事。确实，他不必幻想自己因此损失了任何全德，虽然由于让步他放弃了一些本意的严格，但相反，他确知自己因长久的容忍和忍耐之德获得了更多。因为这就是宗徒的命令：\BibleQuote{你们强壮者，应担代软弱者的软弱；}以及：\BibleQuote{你们要彼此担代重担，这样就满了基督的法律。}因为一个软弱的人决不能支持另一个软弱的人，也不能医治一个健康不佳的人，因为自己也同样受苦；但一个自己不患病的人，却给健康不佳者带来医治。因为对他所说的话是正确的：\BibleQuote{医生，医治你自己吧。}\BibleRef{路 4:23}\par

% structure: p0073 source=h2 target=subsection reason=relative-heading-rank; base=h2

\subsection{第二十四章}

软弱的人如何有害，且不能承受错待。\par

我们也必须注意这一事实：软弱的人本性总是如此——他们迅速而乐于抛出斥责，播下争端的种子，而他们自己却连最轻微错待的影子也不能承受；当他们粗暴地对待我们，随意抛掷罪名时，他们自己却连最细微琐碎的指责也承受不住。因此，按照长老们前述的意见，除非在具有相同志向和善德的人之间，爱不能稳固而不破裂。因为不论其中一方如何小心守护，终有一时它必然破裂。\par

% structure: p0076 source=h2 target=subsection reason=relative-heading-rank; base=h2

\subsection{第二十五章}

一个疑问：不常支持软弱者的人，如何能是强壮的？\par

\Person{热尔玛诺}：如果完美之人的忍耐不能经常容忍软弱者，那么它如何值得称赞？\par

% structure: p0079 source=h2 target=subsection reason=relative-heading-rank; base=h2

\subsection{第二十六章}

\Emph{答案：软弱者并不总是容许自己被容忍。}\par

\Person{若瑟}：我并不是说强壮有力之人的德行和忍耐会被克服，而是说软弱者的可怜境况，因完美之人的容忍而受到鼓励，且日趋恶化，必定会产生一些理由，使得他本身不再应该被容忍；或者，他狡黠地怀疑邻人的忍耐在某时显露并衬托出自己的不耐，因而选择逃离，而非总是被对方的高尚气度所容忍。那么，我们认为，那些想要保持同伴间感情不受损害的人，首先应当观察以下：即，当因任何冒犯而被激怒时，隐修士不仅应保持嘴唇不动，甚至应保持内心深处不动：但如果他发现它们有丝毫扰动，就应该以完全的缄默约束自己，并勤勉地注意圣咏作者所说的：\Quote{我忧闷不语；}以及：\Quote{我说：我要谨遵你的道路，免得我以口舌犯罪。当罪人站在我面前时，我给我的口设置了守卫。我默不作声，谦卑自抑，甚至对善事也缄默不语；}并且他不应留意自己当下的状态，也不应发泄暴怒所暗示的以及激怒的心灵此刻所表达的话，而应默想过去之爱的恩宠，或在心中期盼和平的更新与恢复，甚至在狂怒的时刻也默观之，好像它必然很快会回来一样。当他为即将到来的和谐之乐而保留自己时，他不会感到当前争吵的苦涩，并且会轻易地作出这样的回答：当爱恢复时，他将不能责备自己有罪或受到对方的指责；这样，他就实现了先知的话：\Quote{在忿怒中要记得仁慈。}\BibleRef{哈 3:2}\par

% structure: p0082 source=h2 target=subsection reason=relative-heading-rank; base=h2

\subsection{第二十七章}

\Emph{如何抑制忿怒。}\par

那么，我们应当在明智的指导下约束忿怒的每一个活动并加以节制，以免因盲目的暴怒而被卷入撒罗满所谴责的事：\Quote{恶人发泄其忿怒，但智者逐步平息它，}\BibleRef{箴 29:11}即，愚人被其忿怒的激情所燃烧以致报复自己；但智者，因其谋略的成熟和节制，逐渐减少并消除它。同样，宗徒也说了类似的话：\Quote{亲爱的诸位，你们不要自己报复，而是要给忿怒留地步，}\BibleRef{罗 12:19}即，不要在忿怒的强迫下进行报复，而是给忿怒留地步，即，不要让你们的良心被局限在不耐烦和怯懦的狭窄中，以至于当激情的猛烈风暴兴起时，你们无法忍受；但要在心中扩大，在那爱\Quote{凡事忍耐，凡事包容}\BibleRef{格前 13:7}的广阔深渊中接受忿怒的逆浪。这样，你的心灵将被宽广的坚忍和忍耐所扩大，并拥有其中安全的谋略的深处，忿怒的污浊烟雾将在其中被接收、扩散并立即消散；否则，这段话也可以这样理解：我们给忿怒留地步，每当我们以谦卑和平静的心灵屈服于他人的激情，并向激烈者的不耐烦低头，好像我们承认我们理应受到任何形式的伤害。但是，那些曲解宗徒所说的完美的含义，以致认为那些离开一个正在发怒的人就是给忿怒留地步的人，在我看来，并不是消除，反而是助长了争吵的煽动，因为除非邻居的怒火立即通过谦卑地赔补而被克服，否则人反倒因逃避而激起它而非避免它。撒罗满也说了类似的话：\Quote{不要让你的心灵轻易动怒，因为忿怒栖息在愚人的怀中；}以及：\Quote{不要急于加入争吵，免得最后因此后悔。}因为他并不是在责备急躁地表现出争吵和忿怒，从而称赞迟延的表现。同样也必须这样理解：\Quote{愚人在同一时刻就宣露忿怒，但明智人掩盖自己的羞辱。}\BibleRef{箴 12:16}因为他并没有规定明智人应当隐藏可耻的忿怒爆发，以至于他在责备急躁的忿怒爆发之时却没有禁止迟延的爆发，当然，如果因人性软弱而确实爆发出来，他意指应当为此原因隐藏它：即，虽然暂时被明智地掩盖，它可能被永远摧毁。因为忿怒的本性是这样的：当它被给予空间时，它就衰败并消失，但如果公开表现，它就燃烧得越来越烈。那么，心应当被扩大并敞开，以免它们被局限在怯懦的狭窄中，并被忿怒的汹涌浪潮所充满，这样我们就无法在我们狭窄的心中接受先知所谓的\Quote{极其宽广}的天主诫命，或与先知一同说：\Quote{我奔跑在你诫命的道路上，因为你宽畅了我的心。}因为圣经非常清楚的段落教导我们，坚忍就是智慧：因为\Quote{一个坚忍的人大有明智；但怯懦的人极其愚蠢。}\BibleRef{箴 14:29}因此圣经论及那位值得称赞地向上主祈求智慧恩赐的人说：\Quote{天主赐给撒罗满极大的智慧和明智，以及宽大如海边的沙般浩繁的心。}\BibleRef{列上 4:29}\par

% structure: p0085 source=h2 target=subsection reason=relative-heading-rank; base=h2

\subsection{第二十八章}

\Emph{通过阴谋建立的友谊不能持久。}\par

这一点也屡次被许多经验所证实：即那些从阴谋的起始就进入友谊束缚的人，不可能保持他们的和谐不受破坏；或者因为他们维持友谊并非出于对成全的渴望，亦非由于\OriginalTerm{宗徒的爱}{Apostolic love}的统御，而是出于世俗的爱，出于他们的需要和他们盟约的束缚；又或者因为我们那最狡猾的仇敌更加迅速地催促他们打破友谊的锁链，为的是使他们成为背誓者。因此，最谨慎之人的这个意见是最确实地建立起来的：即真正的和谐与不可分裂的合一，只存在于那些生活纯洁、具有相同良善和目的的人之间。\par

\Person{真福若瑟}就这样在他关于友谊的灵性谈话中论述了这些，并燃起了我们更热切的渴望，要保持我们团体的爱持久不变。\par

\SourceRecord{Translated by C.S. Gibson.；Nicene and Post-Nicene Fathers, Second Series；Vol. 11.；Edited by Philip Schaff and Henry Wace.；Buffalo, NY: Christian Literature Publishing Co.,；1894.；<http://www.newadvent.org/fathers/350816.htm>.}

% structure: p0001 source=h1 target=section reason=page-title

\section{第十七场会议}

\Emph{\Person{圣若望·迦西安}}\par

阿博特若瑟的第二场会议。论许愿。\par

% structure: p0004 source=h2 target=subsection reason=relative-heading-rank; base=h2

\subsection{第一章}

\Emph{论我们所忍受的守夜。}\par

先前那场会议结束时，连同中间夜晚的寂静，我们被圣洁的阿博特若瑟领到一间单独的小屋以求宁静，却彻夜未眠（因着他的话，我们心中燃起了火焰），我们走出小屋，退到约一百步远的地方，在一处隐蔽处坐下。因此，夜色提供了一个私下亲切交谈的机会，我们坐着时，阿博特革玛努斯沉重地叹息。\par

% structure: p0007 source=h2 target=subsection reason=relative-heading-rank; base=h2

\subsection{第二章}

\Emph{论阿博特革玛努斯因回忆我们的许诺而忧虑。}\par

他说道：\Quote{我们在做什么？因为我们看到自己陷入极大的困境，处境恶劣，因为理性本身和圣人的生活有效地教导我们，什么是灵性进步的最好方法，然而我们向长上们所作的许诺却不允许我们选择有益的事。因为，要不是我们许诺的条件迫使我们立即返回隐修院，我们可以借着如此伟大之人的榜样，被塑造为更完美的生活和目标。但如果我们返回那里，就再也不会得到再来这里的机会。但如果我们留在这里，选择实现我们的愿望，那么我们向长上们所发的誓言的信用何在？我们知道我们曾许诺尽快返回，以便我们能够匆匆巡视本省各隐修院和圣人。当处在此种动荡中，我们无法决定如何处置我们的救恩时，我们只能以叹息来证明我们处境的艰难，责备我们自己鲁莽的冒失，同时又憎恨我们天生的羞耻，受此重压，我们无法以其他方式抗拒那些阻拦我们、有损于我们利益和目的之人的恳求，除非用快速返回的许诺。我们确实哀哭，因为我们受那羞耻之过所苦，经上论及此说：\Quote{有一种羞耻带来罪恶。}\BibleRef{箴 26:11}}\par

% structure: p0010 source=h2 target=subsection reason=relative-heading-rank; base=h2

\subsection{第三章}

\Emph{我对此主题的看法。}\par

我回答说：长上的忠告，更确切地说，他的权威，可以迅速解决我们的困难；无论他的裁决如何，都如同神圣和天上的答复，可以终结我们所有的烦恼。我们不必怀疑，主为了他的功德也为了我们的信德，通过这位长上的口赐予我们什么。因为藉着他的恩赐，信徒常从不配的人那里得到有益的忠告，而不信者则从圣人那里得到，主或鉴于回答者的功德，或鉴于求问者的信德而赏赐此事。于是圣洁的阿博特革玛努斯热切地抓住这些话，仿佛我不是出于自己而是因主的催促而说的。我们稍待长上和夜间祈祷时辰的临近，在惯常的问候之后，诵读了适当数量的圣咏和祷词，便如常坐在我们睡觉时铺的同一草席上。\par

% structure: p0013 source=h2 target=subsection reason=relative-heading-rank; base=h2

\subsection{第四章}

\Emph{阿博特若瑟的问题及我们关于忧虑起因的回答。}\par

可敬的若瑟见我们情绪低落，猜想这必有其因，便用圣祖若瑟的这些话对我们说：\Quote{你们今日为何面带愁容？}\BibleRef{创 40:7}我们回答：我们不像法郎的那些仆人，见了梦却无人能解；但我承认，我们彻夜未眠，无人能减轻我们忧愁的重担，除非主藉你的智慧除去它们。他因功德和名字都堪当圣祖若瑟的杰出，便说：\Quote{人的困惑难道不是来自主吗？请说出来：因为神圣的慈悲能藉我们的建议，按照你们的信德，为它们提供救治。}\par

% structure: p0016 source=h2 target=subsection reason=relative-heading-rank; base=h2

\subsection{第五章}

\Emph{阿博特革玛努斯解释为何我们想留在埃及，却被拉回叙利亚。}\par

对此，革玛努斯说：\Quote{我们曾以为，我们会带着丰盈的灵性喜乐和裨益，亲眼见到你们的圣德后返回隐修院，回去后虽只能以微弱的竞赛，追随我们从你们教导中学到的东西。因为我们对长上们的爱促使我们向他们许诺，同时幻想我们在那隐修院中多少能效法你们崇高的生活和教义。因此，当我们以为由此会获得一切喜乐时，另一方面，我们被难以忍受的悲伤压垮，因为我们发现无法用这种方法得到对我们有益的东西。如今我们被两面夹击：如果我们愿意遵守在众弟兄前所作的许诺——那是在我主自己从童贞母胎的洞房中照耀出来的地方，祂亲自作了见证——，我们将蒙受灵性生活上最大的损失。但如果我们忽略许诺，留在此地，并选择将那誓言视为不如我们的成全重要，我们又担心说谎和背誓的可怕危险。即使我们采取另一方案，即通过极速返回履行誓言，再尽快回到此地，也无法减轻负担。因为对于追求良善和灵性进步的人来说，即使短暂的延迟也是危险有害的，但我们仍愿保持信德和许诺，尽管是不情愿地返回，除非我们确信会被长上们的权威和爱如此紧紧地束缚，以至于从此再没有机会回到此地。}\par

% structure: p0019 source=h2 target=subsection reason=relative-heading-rank; base=h2

\subsection{第六章}

\Emph{阿博特若瑟的问题：我们在埃及比在叙利亚得到更多益处吗？}\par

对此，有福的若瑟沉默片刻后说：\Quote{你确定你能在这个国家在灵性事务上获得更多益处吗？}\par

% structure: p0022 source=h2 target=subsection reason=relative-heading-rank; base=h2

\subsection{第七章}

\Emph{关于两地习俗差异的答复。}\par

格尔玛诺：虽然我们应当\Emph{极其}感谢那些从幼年就教导我们追求伟大事物、并以其卓越榜样在我们心中植入对成全之崇高渴望的人，然而，若我们的判断尚可倚靠，我们无法将这些习俗与我们在那里学到的习俗进行任何比较，以至于对您生活中无可比拟的纯洁缄口不言——我们相信，这纯洁不仅是因您心智与目标的专注，也因这地方本身的助佑和扶助而赐予您。因此，我们不怀疑，要追随您伟大的成全，仅凭给予我们的这番教导是不够的，除非我们也有生活的帮助，并且长期的教导训练逐渐化解我们内心的冷漠。\par

% structure: p0025 source=h2 target=subsection reason=relative-heading-rank; base=h2

\subsection{第八章}

\Emph{成全之人如何绝对不应随意许愿，以及决定能否在不犯罪的情况下撤销。}\par

若瑟：确实，我们有效地履行对任何许诺所决定之事，是好的、正当的，且完全符合我们的修道身份。因此，隐修士不应仓促许愿，以免被迫做不慎承诺之事，或因更明智见解的考虑而动摇，显得是背约者。但因目前我们之目的与其说是论述健康状态，不如说是医治疾病，故应以有益的劝言思考：首先考虑的不是你当初应当做什么，而是你如何能从这危险船难的礁石中逃脱。当毫无锁链阻碍、无条件限制时，在善事比较中，若提出选择，应优先选取最有益者；但若有某种损害与损失挡道，在有损之事比较中，应寻求使我们遭受最小损失者。此外，如你的陈述所示，当你不慎的许诺使你陷入这种境地，即无论哪种情况，你都必须承受某种严重损失与不便时，意志应倾向于选择造成较可忍受损失、或更易通过补赎补救的一方。因此，你若认为留在此处对你的灵性比在那所隐修院的生活方式带来的好处更多，且你的许诺条件无法在不损失大善的情况下履行，那么你宁可按受因撒谎和未履行诺言的损失（这是单次行为，无需重复或成为其他罪的根由），也不愿遭受那种损失（据你说，会使你的生活状态变得冷淡，且每日每时损害你）。因为粗心的许诺若能转向更好的道路，便可被宽恕甚至受赞扬，我们不应视之为不一致的失败，而是对轻率的纠正；凡有瑕疵的许诺被纠正时，皆应如此。这可由圣经最确定的见证证明：对许多人而言，履行许诺导致死亡；反之，对许多人而言，拒绝许诺是有益且有利的。\par

% structure: p0028 source=h2 target=subsection reason=relative-heading-rank; base=h2

\subsection{第九章}

\Emph{为何常常违背承诺比履行承诺更好。}\par

这两点由圣宗徒伯多禄和黑落德的例子清楚显明。前者因背离其明确决心（他曾以宣誓般确认说：\BibleQuote{你永远不可洗我的脚。}\BibleRef{若 13:8}）获得了与基督不朽的共融；反之，若他固执坚持己言，必会被断绝于这真福的恩宠。后者因固执于其不智誓言的承诺，成了主的先驱的血腥凶手，并因惧怕作伪证而陷入永罚与永死的惩罚。因此，凡事必须考虑结局，并依此行止与定志；若后来有更明智的劝言出现，发现它偏向较坏的一面，则宁可放弃不当的安排，转向更好的心意，也不可固执于承诺而陷入更重的罪。\par

% structure: p0031 source=h2 target=subsection reason=relative-heading-rank; base=h2

\subsection{第十章}

\Emph{我们关于在叙利亚隐修院所发誓言之恐惧的疑问。}\par

格尔玛诺：就我们为灵性益处而决心实行的心愿而言，我们本希望透过与您持续交往而受益。因为若我们返回隐修院，不仅肯定无法达成如此崇高的目标，而且还会因那里生活方式的平庸而遭受严重损失。但福音的命令令我们极其恐惧：\BibleQuote{你们的话，是就说是，非就说非；除此之外，都是出于邪恶者。}\BibleRef{玛 5:37} 因为我们坚持认为，无法以任何义德补偿违反如此重要命令的过失，且以不良开端之事最终不可能善终。\par

% structure: p0034 source=h2 target=subsection reason=relative-heading-rank; base=h2

\subsection{第十一章}

\Emph{答复：必须考虑行者的意向，而非事务的实行。}\par

若瑟：正如我们所说，在每件事上，我们不应看工作的进展，而要看工人的意向；不应先问一个人做了什么，而应问其目的何在。如此我们会发现，有些人因后来产生善果的行为而被定罪；反之，有些人藉着本身应受谴责的行为达到了义德的顶峰。在前者身上，行为的好结果于他们无益，因为他们以恶意接受了此事，并想促成——并非实际产生的善，而是相反之事；在后者身上，坏的开端并未害及他，因为他忍受了应受责备的开始之必要；这并非出于对天主的轻慢，或存心作恶，而是着眼于必要且神圣的结局。\par

% structure: p0037 source=h2 target=subsection reason=relative-heading-rank; base=h2

\subsection{第十二章}

\Emph{幸运的结局对作恶者毫无益处，而恶行也不会损害善人。}\par

为使我们借圣经中的实例阐明这些言论，还有什么比\Person{主}受难的救赎之药更有利于普世、更有益于世界呢？然而，它不仅毫无益处，反而对那借此成就此事的叛徒有害，以致论及他时确切地说：\BibleQuote{那人若没有生下来，为他更好。}\BibleRef{玛 26:24}因为他的劳苦的果实将不按照实际的结果回报给他，而是按照他想要做的、且相信自己会完成的事。再者，还有什么比诡计和欺骗——且不说对兄弟和父亲，即便对陌生人——更应受谴责呢？然而先祖\Person{雅各伯}不仅未因此类事情受到谴责或责备，反而获得了永久的祝福遗产。这并非无因，因为后者渴望那属于长子的祝福，并非出于贪图眼前利益，而是出于对永恒圣化的信德；而前者（犹达斯）将万民的救赎者交于死亡，并非为了人的救恩，而是出于贪财之罪。因此，在每种情况下，他们行为的果实皆按照心意的意向和意志的目的来估算，根据这一点，其一的目的并非行欺诈，另一的目的也非行救恩。因为各人将公正地按照他最初在心中所构思的、而不是按照由此产生的好或坏的结果——违背工作者意愿的——得到回报。所以，最公正的审判者认为那冒如此谎言者是可原谅的，甚至值得称赞，因为没有它，他无法获得长子的祝福；那源于祝福渴望的事不应算作罪。否则，上述先祖不仅对兄弟不公平，对父亲也是欺骗，对天主是亵渎，若还有其他方式能使他获得那祝福的恩赐，他却选择这条损害、伤害兄弟的道路。可见，天主所审查的并非行为的实施，而是心意的目的。有了这番预备返回所提出的问题（这一切皆为前言），我希望你首先告诉我，你们为何用那誓言的锁链束缚自己。\par

% structure: p0040 source=h2 target=subsection reason=relative-heading-rank; base=h2

\subsection{第十三章}

\Emph{我们关于为何要向我们起誓的回答。}\par

\Person{格尔玛诺斯}：第一个原因，如我们所说，是害怕激怒我们的长老、违抗他们的命令；第二个原因是，我们非常愚蠢地相信，若从你们那里学到了任何完美的或值得听视的事，当我们回到隐修院时，便能够实行它。\par

% structure: p0043 source=h2 target=subsection reason=relative-heading-rank; base=h2

\subsection{第十四章}

\Emph{长老的讲论：只要坚持实现善意，如何可以无过地改变行动方案。}\par

\Person{若瑟}：如我们先前所言，心意的意图带给人的是赏报或谴责，根据这段经文：\BibleQuote{他们的思念互相控告或辩护，在天主审判人隐秘事的日子。}以及这段：\BibleQuote{我要聚集万民及各种语言的民族，来观看我的荣耀。}因此，依我看来，你们是因追求成全而用这些誓言的锁链束缚自己，因为你们当时认为通过此方案能获得成全；而现在，更成熟的判断已出现，你们看到不能借此方式攀上其高峰。因此，那计划中的任何变动——若看似已发生——不会成为阻碍，只要那最初的意向没有随之改变。因为更换工具并不等于放弃工作，选择更短更直接的道路也不表示旅行者懒惰。故在此事中，对短浅安排的改善不可视为属灵承诺的背弃。凡出于对\Person{天主}的爱和对美善的渴望而做的事，既\BibleQuote{有今生和来生的应许}\BibleRef{弟前 4:8}，即使看似以艰难和不利的开端开始，不仅不应受责备，反而值得称赞。因此，轻率承诺的打破不会成为阻碍，只要每件事的结局，即所意向的美善目标得以保持。因为我们做这一切，正是为了能向\Person{天主}显出一颗清洁的心；若在此地更容易实现此目标，那么你们被强求的协议之变更对你们无害，只要那最初为使你们作出承诺的成全——即纯洁的成全——能更快地照主的旨意获得。\par

% structure: p0046 source=h2 target=subsection reason=relative-heading-rank; base=h2

\subsection{第十五章}

\Emph{一个问题：我们的见识为软弱的弟兄提供说谎的机会，是否可以无罪？}\par

\Person{格尔玛诺斯}：就那些经过合理慎重考虑的言辞的力量而言，我们关于承诺的疑虑本可轻易消除，若非我们极其惊恐：唯恐此例为某些较软弱的弟兄提供说谎的机会，若他们知道协议的信用可在某种程度上合法地打破，而这正是先知以如此严厉和警告性的言辞所禁止的：\Quote{上主，你必要灭绝一切说假话的人；}以及：\Quote{说谎的口，能杀害灵魂。}\par

% structure: p0049 source=h2 target=subsection reason=relative-heading-rank; base=h2

\subsection{第十六章}

\Emph{回答：圣经真理不可因软弱者的跌倒而更改。}\par

\Person{若瑟}：走向毁灭、更确切地说是急于自我毁灭的人，不可能缺少自我毁灭的机会与机缘；那些助长异端者乖僻、增加犹太人不信、或触犯外邦智慧骄傲的圣经段落，也不应被拒绝或完全从卷册中删去；相反，它们理应被虔信地接受、坚定地持守，并按真理的规则宣讲。因此，我们不应因他人的不信，而拒绝圣经所记载的先知和圣人的\OriginalTerm{经济}{\GreekText{οἰκονομίας}}，即\Quote{经纶}，免得我们一面想着该迁就他们的软弱，一面却不仅犯下说谎的罪，还犯下亵圣的罪。相反，正如我们所说，我们应按字面接受这些事，并解释它们为何是正当地做的。但对于那些心存恶意的人，如果我们试图完全否认或通过寓意解经来消解我们将要提出或已经提出的事件的真理，那么说谎的门路并不会因此被堵住。因为，如果他们的败坏意志本身就足以引他们犯罪，那么这些段落的权威又怎能伤害他们呢？\par

% structure: p0052 source=h2 target=subsection reason=relative-heading-rank; base=h2

\subsection{第十七章}

\Emph{圣人如何如同使用菟葵一样有益地使用谎言。}\par

因此，我们应当看待谎言并使用它，如同它的性质是菟葵一样；若在某种致命疾病威胁时服用它，它是有益的；但若没有重大危险的需要而服用，它便立即成为死亡的原因。因为我们也读到，圣人和那些最蒙天主悦纳的人曾使用谎言，不仅因此不招致罪的过犯，甚至获得了最大的善；如果欺骗能使他们得光荣，那么相反，真理岂不是只会给他们带来谴责吗？就如\Person{辣哈布}，圣经对她的记载不仅没有任何善行，甚至还有不贞；但仅仅因为她所撒的谎——她宁愿隐藏探子而不出卖他们——她得以与天主子民永远联合在祝福中。但如果她选择说实话而顾念市民的安全，无疑她和她全家都无法逃脱即将来临的毁灭，她也不会被特许列入我们主降生的先祖之中\BibleRef{玛 1:5}，被算在圣祖的名册里，并透过她后来的后裔，成为万民救主的母亲。相反，\Person{德里拉}为了保障同城居民的安全，出卖了她所发现的真理，却换来了永远的毁灭，只给众人留下了她的罪过的记忆。因此，当承认真理有重大危险时，我们必须以谎言作为避难所，但要以这样的方式：我们因谦卑良心的罪疚而感到困扰，却仍为得救而如此行。但若没有极大的必要，谎言就应当像致命之物一样被极其谨慎地避免：正如我们所说的一杯菟葵，只有在致命且不可避免的疾病威胁时，作为最后的手段服用才有益；若在身体健全健康时服用，其毒性会立刻攻向要害。这一点清楚地显现在耶里哥的\Person{辣哈布}和圣祖\Person{雅各伯}身上：前者唯有藉此药方才能逃脱死亡，后者如无此药方就不能获得长子的祝福。因为天主不仅是我们的言语和行为的审判者与监察者，祂也洞察它们的动机与目的。如果祂看到某人为永恒救恩的缘故做了或应许了某事，并显示出对神圣默观的洞见，即使这事在人看来是严苛或不公平的，祂仍会看内心的良善，看重意志的愿望胜于实际说出的话语，因为祂必须考虑工作的目的和行为者的心向；通过这一点，如前面所说，一个人可以因谎言而成义，而另一个人可能因说实话而犯了永死之罪。圣祖\Person{雅各伯}也关注这一点，他不怕用羊皮裹身来模仿他兄长身体的多毛外表，并值得称赞地顺从他母亲为达此目的而煽动的谎言。因为他看到，这样做将为他带来比保持单纯更高的祝福和正义的收益：他毫不怀疑，这谎言的污点会立即被父亲的祝福洪流洗净，并像小云彩一样被圣神的气息迅速吹散；而且通过他所采用的这种掩饰，他将获得比与他本性相符的真理更丰厚的功德奖赏。\par

% structure: p0055 source=h2 target=subsection reason=relative-heading-rank; base=h2

\subsection{第十八章}

\Emph{异议：唯有生活在法律之下的人使用谎言而免受惩罚。}\par

\Person{Germanus}：难怪这些计谋在旧约中被适当使用，有些圣人可称赞地或至少小罪地说了谎，正如我们看到许多更坏的事因时代的粗野而被允许给他们。因为为什么我们要惊奇，当有福的\Person{达味}逃离\Person{撒乌耳}时，回答司祭\Person{阿希默肋客}的问话说：\BibleQuote{为什么你独自一人，没有人与你在一起？}他回答说：\BibleQuote{王吩咐了我一件事，说：不要让人知道你从我这里奉派的事，因为我已吩咐我的仆人到某处地方。}又说：\BibleQuote{你这里有枪或刀吗？因为我没带自己的刀和武器，因为王的事紧急。}或者当他在\Place{加特}王\Person{阿基士}面前装疯时，\BibleQuote{在他们面前变了脸色，在他们手中滑倒，撞在城门上，唾沫流在胡须上。}当他们甚至被允许拥有众多妻妾，而且因此不归罪于他们，并且他们还常用自己的手流敌人的血，这不仅被认为无可指责，甚至可称赞吗？而这一切我们藉福音之光看到是完全禁止的，以致没有一件可以行而不犯大罪和过犯。同样，我们认为没有人可以使用谎言，我将不说正当地，甚至小罪也不可，即使它可能披着虔诚的外衣，如同主所说：\BibleQuote{你们的话当是：是，是；非，非；除此之外，都是出于那恶者。}宗徒也同意这话：\BibleQuote{不要彼此说谎。}\par

% structure: p0058 source=h2 target=subsection reason=relative-heading-rank; base=h2

\subsection{第十九章}

\Emph{答案：即使在旧盟约下也未曾给予说谎的许可，已被许多人正当地采用。}\par

\Person{若瑟}：随着时日的终结和人类繁衍的完成，在妻妾方面的一切自由，理应由福音的成全所割断，因为它已不再必要。因为在基督来临之前，原始谕令的祝福应该完全有效，其中说：\BibleQuote{你们要生育繁殖，充满大地。}\BibleRef{创 1:28}因此，按照那时期的安排，从在会堂中幸福生长的人类繁殖之根上，萌发出天使般的贞洁之芽，并在教会中产生节欲的芳香花朵，是相当合理的。但是，谎言在那时已经受到谴责，整部旧约经文清楚显示，正如它说：\Quote{你要消灭所有说谎言的人；}又说：\Quote{谎言之粮对人是甘甜的，但事后他的口却充满砂砾；}而立法者本身说：\Quote{你应避免谎言。}但我们说过，当时它被适当地作为最后手段使用，当某种需要或救恩计划与之相关联时，因此它不应受谴责。正如你所提及的\Person{达味}王的例子，当他逃避\Person{撒乌耳}的不义迫害时，对司铎\Person{阿希默肋客}用了谎言，并非为了任何利益，也非有意加害任何人，而只是为了从那次极不义的迫害中拯救自己；因为他不愿用自己的手沾染那敌对君王（尽管多次被天主交在他手中）的血，正如他所说：\BibleQuote{愿主怜悯我，使我决不对我主人，主的受傅者做这样的事，伸手加害他，因为他是主的受傅者。}\BibleRef{撒上 24:7}因此，我们听到旧约下的圣人们所采取的这些计谋，无论出于天主的旨意、为预表属灵的奥迹、或为拯救某些人，当必要迫使我们时，我们也不能完全拒绝，因为我们看到连宗徒们在需要有利考虑时也没有避免它们。这些我们暂时搁置，先讨论我们打算从旧约中提出的那些例子，之后我们会更恰当地引入它们，以便更易证明，在旧约和新约中，善良圣洁的人们在这些策略上完全一致。至于\Person{胡瑟}为拯救\Person{达味}王而对\Person{阿贝沙隆}所施的虔诚诡计，虽然由那欺骗者和骗子以看似善意的话语说出，且有悖于求问者自身的利益，却仍受到圣经权威的称赞，其中说：\BibleQuote{愿主的旨意使\Person{阿希托费耳}的有利计谋被挫败，好使主将灾祸降于\Person{阿贝沙隆}。}\BibleRef{撒下 17:14}那些为正义一方、出于正确目的和虔诚意向所做的，并通过神圣的掩饰来为一位虔诚中悦天主的人谋划拯救和胜利的事，也不能受责备。还有那妇人的事迹，她接待了上述\Person{胡瑟}派往\Person{达味}王的人，将他们藏在一口井里，并用一块布盖住井口，假装正在晾晒珍珠大麦，并说：\Quote{他们尝了一点水后就过去了。}借此发明将他们从追赶者手中救出，我们该说什么呢？所以，我请求你回答我，如果你现在生活在福音之下，遇到类似情况，你会怎么做？你会愿意用类似的谎话将他们藏起来，同样说：\Quote{他们尝了一点水后就过去了，}从而履行这条命令：\BibleQuote{解救被拉去处死的人，不要吝于赎回将被杀的人；}\BibleRef{箴 24:11}还是通过说实话，将那些躲藏的人交给要杀他们的人呢？那么，宗徒的话又如何呢？\Quote{各人不可寻求自己的利益，而要寻求别人的利益；}和\Quote{爱德不求自己的益，只求别人的利益；}他又论到自己说：\Quote{我不求自己的利益，只求众人的利益，为使他们得救。}因为如果我们寻求自己的利益，并固执地保留对自己有益的事，即使在这样紧急的情况下，我们也必须说实话，从而成为他人死亡的罪人；但如果我们把别人的益处置于自己的利益之上，并满足宗徒的要求，我们必然要忍受说谎的必要。因此，我们不能保持一颗完美的爱德之心，或像宗徒的成全所要求的那样寻求别人的事，除非我们在那些关乎我们自己生活严谨与成全的事情上稍作放松，并选择以随时准备的爱意屈尊俯就于对他人有益的事，从而与宗徒一起成为软弱的软弱之人，以便我们能赢得软弱之人。\par

% structure: p0061 source=h2 target=subsection reason=relative-heading-rank; base=h2

\subsection{第二十章}

\Emph{连宗徒们也认为谎言常常有益而真理有害。}\par

受这些例子的教导，蒙福的\Person{雅各伯}宗徒以及初期教会的所有首要首领，也都因软弱者的软弱，敦促\Person{保禄}宗徒屈从于一种虚构的安排，坚持要他按照法律的要求洁净自己、剃头并还愿，因为他们认为这种伪善所带来的当下损害不足挂齿，反而更看重由此带来的持续宣讲的益处。因为\Person{保禄}宗徒若固守严格，其益处也抵不过因他过早离世而给万民造成的损失。而当时整个教会若不借助这种良善而有益的伪善保护他继续宣讲福音，这损失就必然发生了。因为当我们说，讲真话会带来更大的危害，且讲真话所获得的善果无法抵消它所造成的损害时，我们才可以正当且可原谅地默认一个谎言的错误。\newline{}\newline{}此外，蒙福的宗徒在其他地方也用不同的言辞证明自己常持此心态；因为当他说：\Quote{对犹太人，我就成为犹太人，为要赢得犹太人；对在法律下的人，我就成为在法律下的，虽然我自己不在法律下，为要赢得在法律下的人；对没有法律的人，我就成为没有法律的，虽然我不是没有天主的法律，而是在基督的法律之下，为要赢得没有法律的人；对软弱的人，我就成为软弱的，为要赢得软弱的人：对一切人，我就成为一切，为要救一切人；}\BibleRef{格前9:20}——这岂不表明他常根据受教者的软弱和能力而降卑自己，放松一些完美的力度，不固守自己严谨生活所似乎要求的事，反而更偏爱软弱者之益所要求的事？\newline{}\newline{}为了更仔细地追溯这些事，并逐一列举宗徒善行的荣光，有人可能会问：蒙福的宗徒如何证明自己在一切事上适应了一切人？他何时对犹太人成了犹太人？当然是在这样的情形中：当他内心深处仍持有对迦拉达人所说的意见，即：\Quote{看，我保禄告诉你们：如果你们受割损，基督对你们就毫无益处了，}\BibleRef{迦5:2}然而，他通过给\Person{弟茂德}行割损，采取了犹太迷信的影儿。\newline{}\newline{}再者，他何时对在法律下的人成了在法律下的？正是在此情形中：\Person{雅各伯}和教会所有的长老，担心他会被犹太信徒，或者更确切地说，是被那些以仍受法律仪式礼仪束缚的方式接受了基督信仰的犹太化基督徒的群众攻击，于是用这个劝告和建议来帮助他摆脱困境，说：\Quote{弟兄，你看，犹太人中信主的有多少万，他们都是为法律热忱的人。他们听说你教训在外邦人中的犹太人去离弃\Person{梅瑟}，说他们不该给孩子行割损，也不必遵行古礼；……你就照我们对你说的去做吧：我们这里有四个人，他们都有愿在身；你带他们去，同他们一起行取洁礼，并替他们付费，叫他们剃头；这样，众人就知道先前所听到关于你的事都是虚的，而你自己也是一位遵守法律的人。}\BibleRef{宗21:20}于是，为了在法律下之人的益处，他暂时践踏了自己曾表达的严格意见：\Quote{因为我借着法律已死于法律，为能生活于天主；}\BibleRef{迦2:19}并被驱使剃头、按法律取洁、在圣殿中按\Person{梅瑟}的仪式还愿。\newline{}\newline{}你还要问：他何时为了完全不知天主法律之人的益处而自己成了没有法律的？请读他在\Place{雅典}讲道的开场白，那里异教的邪恶正盛行。他说：\Quote{我经过各处，看见你们所敬拜的神祇，也见到一座祭坛，上面写着：\InnerQuote{致未识之神}}然后他从他们的迷信出发，仿佛他自己也没有法律，借着那个亵渎的铭文引入了基督的信仰，说：\Quote{你们所不识而敬拜的，我就传告给你们。}稍后，他仿佛对神圣的法律一无所知，宁愿引用一位异教诗人的诗句，而不是\Person{梅瑟}或基督的言论，说：\Quote{正如你们自己的诗人也说过：我们也是祂的子孙。}他这样借他们自己不能拒绝的权威去接近他们，从而以虚假之事证实真理，又接着说：\Quote{既然我们是天主的子孙，就不该以为天主的神性仿佛与金银或石头雕刻的、由人技艺和思想制成的偶像相似。}\newline{}\newline{}\newline{}但他对软弱的人成了软弱的，当他以准许而非命令的方式允许那些不能自制的人再次结合时，\BibleRef{格前7:5}或者当他用奶而不是饭喂养格林多人，并说他在他们中间是软弱、恐惧和战栗的。但他成为一切人为了救一切人，当他说：\Quote{吃的，不要轻视不吃的；不吃的，也不要判断吃的；}以及：\Quote{嫁出自己女儿的人做得好，不嫁出的人做得更好；}别处又说：\Quote{谁软弱，我不软弱呢？谁跌倒，我不焦急呢？}这样，他履行了自己命令格林多人去做的事，他说：\Quote{你们不要成为犹太人、希腊人或天主的教会的绊脚石，如同我在一切事上取悦众人，不求自己的利益，而求多数人的利益，为叫他们得救。}\newline{}\newline{}因为不给\Person{弟茂德}行割损、不剃头、不接受犹太人的取洁、不赤脚行走、不履行法律上的愿，确实都是有益的；但他做了这一切，因为他不是求自己的利益，而是求多数人的利益。尽管这是在完全考虑天主的情况下做的，却并非没有掩饰。因为一个借着基督的法律死于法律而为天主生活的人，一个曾将自己在法律中的义——他在这方面曾无可指摘——视作废物以求获得基督的人，不可能以真诚的热心奉献属于法律的事；我们也不应该相信，那位说\Quote{如果我重建我所拆毁的，我就证明自己是违犯者}\BibleRef{迦2:18}的人，会自己陷入他所谴责的事。\newline{}\newline{}而且，是如此看重做事的动机而非实际行为，以至于另一方面，真相有时被发现对人有害，而谎言却有益。当\Person{撒乌耳}对他的仆从抱怨\Person{达味}的逃亡，说：\Quote{\Person{叶瑟}的儿子会赐给你们田地、葡萄园，立你们为千夫长、百夫长吗？你们竟都同谋对付我，没有人告诉我？}这时，厄东人\Person{多厄格}说了什么？他不是说出了真相吗？他告诉他：\Quote{我看见\Person{叶瑟}的儿子来到\Place{诺布}，到了\Person{阿希突布}的儿子\Person{阿彼默肋客}司祭那里；司祭为\Person{达味}求问了主，给了他食物，又给了他培肋舍特人\Person{哥肋雅}的刀。}因为这个真实的故事，他该当从活人之地被拔除，先知论及他说：\Quote{因此天主必永远毁灭你，将你拔除，将你从你的帐幕中撕出，将你的根从活人之地拔除。}他因显明真相而被永远拔除，从那个地中拔出，而妓女\Person{辣哈布}却因她的谎言而家族被移栽到那地中。同样，我们也记得\Person{三松}最有害地将长久以谎言隐藏的真相泄露给了他那邪恶的妻子，因此这轻率泄露的真相成了他自己被骗的原因，因为他忽略了先知的命令：\Quote{要谨守你口的门，不要向睡在你怀中的妇人泄露。}\BibleRef{米2:7}\par

% structure: p0064 source=h2 target=subsection reason=relative-heading-rank; base=h2

\subsection{第二十一章}

\Emph{是否应当在询问者面前公开秘密的斋戒而不必撒谎，以及一旦拒绝的事是否可以重新着手。}\par

再举一些我们无法避免且几乎是日常所需的例子，这些需要即使我们百般小心，也永远无法防范到不被驱使其发生，无论我们愿意与否：请问，当我们打算推迟晚餐时，一位弟兄前来问我们是否已经用过餐：我们是应当隐瞒斋戒，隐藏这善工的斋戒，还是应当说实话将其宣扬呢？如果我们隐瞒，为遵守主的诫命说：\BibleQuote{你不可对人显现为斋戒，只对你的在隐秘中的父；} 又说：\BibleQuote{不要让你的左手知道右手所做的，} 我们就必须立刻说谎。如果我们彰显这善工的斋戒，福音的话语合理地劝阻我们：\BibleQuote{我实在告诉你们，他们已经得到了他们的赏报。} 但如果有人坚决拒绝了某位弟兄递来的杯，完全否认他会接受那弟兄因他的到来而欢喜地恳求他收下的东西，又该如何呢？他应当强迫自己向那位跪地叩首、以为只有这样才能表达爱心的弟兄让步吗？还是应当固执地坚持自己的言语和意向？\par

% structure: p0067 source=h2 target=subsection reason=relative-heading-rank; base=h2

\subsection{第二十二章}

\Emph{异议：斋戒应当隐藏，但已拒绝的东西不应再接受。}\par

\Person{日耳曼努斯}：在第一种情形中，我们以为毫无疑问，隐藏斋戒比向询问者显露更好，在这类情况下我们也承认谎言不可避免。但在第二种情形中，我们无需说谎，首先因为我们可以拒绝弟兄服侍所给的东西，而不将自己束缚于任何决定的捆绑中，其次因为一旦拒绝，我们可以保持我们的意见不变。\par

% structure: p0070 source=h2 target=subsection reason=relative-heading-rank; base=h2

\subsection{第二十三章}

\Emph{答复：此种决定的固执是不合理的。}\par

\Person{若瑟}：毫无疑问，这些决定来自那些隐修院——根据你的叙述，你弃绝的初学阶段是在那里受训的——它们的院长惯于将自己的意愿置于弟兄的晚餐之上，且最固执地坚持自己曾定意的事。但我们的长老们——他们的信德有宗徒能力的标记作证，他们以精神的判断和审慎而非心灵固执来处理一切——规定：那些迁就他人软弱的人，所得的果实远比那些坚持自己决定的人更为丰硕，并宣称：以这种必需的谦卑谎言来隐藏斋戒，远胜于以骄傲的真理炫耀它。\par

% structure: p0073 source=h2 target=subsection reason=relative-heading-rank; base=h2

\subsection{第二十四章}

\Emph{院长\Person{皮亚蒙}如何选择隐藏他的斋戒。}\par

最后，院长\Person{皮亚蒙}在二十五年后毫不犹豫地接受了某位弟兄献上的葡萄和酒，并立刻违反他的规矩品尝了所献之物，而不愿显露他对所有人都保密的斋戒。因为如果我们还记得长老们常做的事——他们为了在会议中教导初学修士而不得不提出自己善功的奇事和行为时，常常假托他人之名来隐藏——那么除了公开的谎言，我们还能认为这是什么吗？但愿我们也有任何配得上的事，可以用来激发初学修士的信德！那么，我们无疑会效仿他们那种虚构。因为用那种比喻的颜色说谎，胜过为了保持那不合时宜的真理，要么用不明智的沉默隐藏可能有益于听众的事，要么以真实但以自己的身份说出来而陷入招人厌恶的虚荣。外邦人的导师也用他的教导清楚地教给了我们同样的道理，因为他选择以他人的口吻说出他所得的大启示，说：\BibleQuote{我认识一个在基督内的人，或在身内，或在身外，我不知道，天主知道；他被提到第三层天上。我也认识这样一个人，他被提到乐园里，听到了不可言说的话，是人不可说出的。} \BibleRef{格后 12:2}\par

% structure: p0076 source=h2 target=subsection reason=relative-heading-rank; base=h2

\subsection{第二十五章}

\Emph{圣经关于改变决定的证据。}\par

我们无法简短地一一列举一切。因为谁能数算几乎所有的圣祖和无数的圣人呢？他们中有些是为了保存生命，有些是出于对祝福的渴望，有些是出于怜悯，有些是为了隐藏某个秘密，有些是出于对天主的热情，有些是为了寻求真理，可以说成了谎言的庇护者。既然不能全部列举，也不应完全略过。因为虔诚迫使荣福若瑟甚至以国王的性命起誓，虚假地指控他的兄弟们，说：\Quote{你们是探子，来窥看这地的薄弱之处；}随后又说：\Quote{你们中派一人去，将你们的兄弟带来；你们要留在这里，直到你们的话得以证实，看你们是否说实话；不然，以法老的性命担保，你们是探子。} \BibleRef{创 42:9} 因为若他不以怜悯之心藉此谎言惊吓他们，他就无法再见他的父亲和兄弟，也无法在饥荒的巨大危险中保全他们，更无法使他的兄弟们摆脱出卖他的罪恶感。所以，藉谎言使兄弟们恐惧的行为，并不那么应受谴责，反而是神圣可赞的，因为他藉伪装危险敦促他的敌人和寻求者作有益的补赎。最后，当他们被这严重指控的重压所困时，他们良心受责，不是因为虚假的指控，而是想到先前的罪行，彼此说：\Quote{我们受这苦是应该的，因为我们得罪了我们的兄弟，看见他灵魂的痛苦，他恳求我们，我们却不听从他；因此这一切灾祸临到我们身上。}我们认为，这次认罪以最有益的谦卑赎了他们的重罪，不仅是得罪兄弟（他们曾以恶毒残忍得罪了他），也是得罪了天主。至于撒罗满，他在第一次审判中，仅仅藉使用谎言，就彰显了从天主领受的智慧恩赐？因为为了揭露被那妇人谎言隐藏的真相，他甚至使用了一个极其巧妙编造的谎言的帮助，说：\Quote{给我拿刀来，将活孩子劈成两半，一半给这一个，一半给那一个。}当这假装的残酷刺透了真正母亲的心，却得到非真正母亲的赞同时，他终于藉这极其明智的真相发现，作出了每个人都认为是天主启发的判决，说：\Quote{将活孩子给她，不可杀他；他是她的母亲。} \BibleRef{列上 3:24} 此外，圣经的其他段落也更充分地教导我们，我们不能也不应执行我们因心平气和或心烦意乱而定的一切决定，因为我们常听到圣人、天使甚至全能天主自己改变了他们先前决定的事。荣福达味决定并发誓说：\Quote{愿天主这样对待达味的仇敌，并且加倍，如果我到早晨还给纳霸耳留下任何一个男子。}随后，当他的妻子阿彼盖耳为他求情时，他放弃威胁，减轻判决，宁愿被视为食言，也不愿因残酷执行誓约而遵守誓言，说：\Quote{上主永在！若不是你赶快来迎接我，到早晨，纳霸耳必不留一个男子。}正如我们不认为他那出于愤怒和心烦意乱的鲁莽宣誓应被我们效法，我们也认为他后来的宽恕和改变决定是值得效法的。\Quote{特选的器皿}在致格林多人书中，无条件地应许要回来，说：\Quote{但我从马其顿经过时，必到你们那里去；因为我要经过马其顿。但我要在你们那里住下，甚至过冬，好叫你们送我往我所要去的地方。因为我不愿仅仅路过见到你们；我希望住些时候。} \BibleRef{格前 16:5} 他在第二封书信中记起这事，这样说：\Quote{我既然有这样的信心，就打算先到你们那里，使你们再蒙恩惠，然后经过你们往马其顿去，再从马其顿回到你们那里，由你们送我往犹太去。}但一个更好的计划出现，他坦承不实现他所承诺的。\Quote{既然如此，我计划这事时，难道我是轻率吗？或者我计划的事，是按血气计划，以致在我有\Quote{是，是}，又有\Quote{非，非}吗？}最后，他甚至以誓言声明，为什么他宁愿搁置诺言，也不愿因他的露面给门徒带来负担和忧伤：\Quote{我呼求天主作证，我以我的灵魂起誓，是为了宽容你们，我才没有到格林多去。因为我自己决定，不再在忧愁中到你们那里去。}当天使们拒绝进入索多玛的罗特家，向他说：\Quote{我们不在家里过夜，要留在街上，}他们立即被他的恳求所迫，改变了决定，正如圣经接着说：\Quote{罗特极力邀请他们，他们便转身进到他那里。} \BibleRef{创 19:2} 如果他们知道他们会进去，他们是以虚假的借口拒绝他的请求；但如果他们的借口是真的，那么他们显然改变了主意。我们确信圣神将此记在圣经中，无非是藉他们的榜样教导我们，我们不应固执于自己的决定，而应使它们服从我们的意志，使我们的判断摆脱一切法律的束缚，好能随时随从好忠告的召唤，毫不迟延地转向健全辨别所认为更好的选择。为了提升到更高的例子，当希则克雅王卧病在床，患重病时，先知依撒意亚以天主的名义对他说：\Quote{上主这样说：料理你的家事，因为你必死，不能活。}希则克雅就转脸朝墙，恳求上主说：\Quote{恳求你，上主，求你记念我怎样以真诚和全心在你面前行走，怎样作了你眼中视为正义的事。}希则克雅就放声大哭。此后，上主的话又传给他：\Quote{你回去，告诉我民的首领希则克雅说：你祖先达味的天主上主这样说：我听见了你的祈祷，看见了你的眼泪；看，我要增加你十五年的寿数，且要把你从亚述王手中救出，并为了我自己和我仆人达味的缘故，保护这座城。} \BibleRef{列下 20:1} 有什么比这证明更清楚呢？主出于慈悲和良善，宁愿打破自己的话，不按预定死亡期限，而将祈祷者的寿命延长十五年，也不愿因祂不可改变的旨意而显得决绝？同样，天主对尼尼微人说的话：\Quote{还有三天，尼尼微就要毁灭了；}随即这严厉而突兀的判决因他们的补赎和斋戒而缓和，转向慈悲和易于祈求的良善。但如果有人认为，主威胁毁城（尽管祂预知他们将悔改）是为了激发他们作有益的补赎，那么，管理弟兄的人，若有需要，也可以毫无说谎之咎，对那些需要改进的人，威胁以比实际将施加的更严厉的处置。但如果有人说，天主因他们的补赎而撤销了那严厉的判决，按照祂藉厄则克耳所说的：\Quote{如果我对恶人说：你必要死；而他悔改自己的罪，行公平正义，他必要活，不会死；} \BibleRef{则 33:14} 我们同样受教，不应固执于自己的决定，而应以温和的怜惜缓和因需要而发出的威胁。为了不让我们以为主只对尼尼微人特准此事，祂不断藉耶肋米亚声明，祂将对所有人普遍如此行，并应许立即根据我们的功过改变判决，说：\Quote{我忽然论到一邦一国，说要拔除、拆毁、毁灭；但那邦若悔改我所指摘的恶，我也必后悔，不降我原想降给他们的灾祸。我忽然论到一邦一国，说要建立、栽植；但那邦若行我眼中为恶的事，不听我的声音，我也必后悔，不将我曾想赐给他们的福赐给他们。}对厄则克耳也说：\Quote{一句话也不可省略，也许他们会听从，各人回头离开邪道，使我能后悔我因他们邪恶行为原想加给他们的灾祸。}这些段落宣告，我们不应固执于自己的决定，而应以理性和判断修改它们，并应始终采纳和选择更好的途径，毫不迟延地转向被认为更有益的途径。尤其是那无价的句子教导我们，因为虽然每个人的结局在出生前就已为祂预知，但祂却以计划和方式为所有人安排一切，并顾及人的性情，祂决定一切不是单凭自己的能力，也不是凭着祂预知拥有的不可言喻的知识，而是根据人现时的行为，接纳或拒绝每一个人，每天赐予或收回祂的恩宠。撒乌耳的拣选也向我们表明这一点，天主的预知当然不能不知道他悲惨的结局，但祂却从以色列千万人中拣选他，傅他为王，赏报他当时生活的功绩，而不考虑他将来的堕落之罪，以致当他成为弃者后，天主几乎以人的口吻，带着人的情感，好像后悔自己的拣选，说：\Quote{我后悔立了撒乌耳为王，因为他离弃了我，没有遵行我的话；}又说：\Quote{撒慕尔因撒乌耳悲伤，因为上主后悔立了撒乌耳为以色列王。}最后，主也藉先知厄则克耳宣布，祂将在日常审判中对所有人如此行，说：\Quote{如果我对义人说：你必要活；而他依仗自己的义行，却犯下不义，那么他的一切义行都不再被记念，他所犯的不义，他必因此而死。如果我对恶人说：你必要死；而他悔改自己的罪，行公平正义，归还抵押品，偿还所抢的，遵行生命的诫命，不再行不义，他必要活，不会死。他所犯的一切罪过都不归咎于他。} \BibleRef{则 33:13} 最后，当主因百姓迅速堕落而要转离祂慈颜离开祂从万民中所拣选的百姓时，颁赐法律者替他们求情，呼喊说：\Quote{恳求你，上主，这百姓犯了重罪，为自己铸了金神像。现在，你若赦免他们的罪，就赦免；不然，就求你从你所写的册子上删除我的名字。}上主回答：\Quote{谁得罪了我，我就从我的册子上删除谁。} \BibleRef{出 32:31} 达味也在先知精神中抱怨犹达斯和主的迫害者，说：\Quote{愿他们从生命册上被涂去；}因为他们因重罪之咎不配得到救恩的补赎，他接着说：\Quote{愿他们不被列入义人之中。}最后，在犹达斯本人身上，先知的诅咒清楚应验了，因为他完成了致死之罪后，自缢而死，以免名字被涂去后，还能悔改，配得再次被写在义人之中。因此，我们不应怀疑，当犹达斯被基督拣选，获得宗徒地位时，他的名字写在生命册上，他也和其余人一样听到：\Quote{不要因为魔鬼屈服于你们而欢喜，而要因为你们的名字已登记在天上而欢喜。} \BibleRef{路 10:20} 但因他被贪婪的瘟疫腐蚀，名字从那天上名单中被除去，先知论及他和类似的人恰当地说：\Quote{上主，愿所有离弃你的人都蒙羞。愿那些离开你的人都写在地上，因为他们离弃了上主，这活水的泉源。}别处又说：\Quote{他们不得归入我百姓的集会，不得写在以色列家的册上，也不得进入以色列地。}\par

% structure: p0079 source=h2 target=subsection reason=relative-heading-rank; base=h2

\subsection{第二十六章}

\Emph{圣德之人如何不会固执刚硬。}\par

我们也不可忽略那命令的价值，因为即使我们在愤怒或其他情绪的驱使下（隐修士绝不应如此），以某种誓言约束了自己，仍应用健全的理智审慎权衡双方情况，将已定之策与所劝之法相比较，并毫不犹豫地采纳在更明智的考量下被判定为最佳者。因为舍弃承诺，胜过失去某种美好且更可取的利益。最后，我们从不记得那些可敬且公认的教父在此类决定上刚硬不化；相反，他们如同热力下的蜡，被理性所改变，当更明智的劝告占上风时，他们会毫不犹豫地屈从于更好的方面。但我们所见固执于己见者，则一概视作非理性且缺乏判断。\par

% structure: p0082 source=h2 target=subsection reason=relative-heading-rank; base=h2

\subsection{第二十七章}

\Emph{问题：\Quote{我已发誓并决意}这句话是否与上述观点相悖？}\par

\Person{Germanus}：就这一已清楚充分论述的考虑而言，隐修士绝不应决定任何事，以免他要么成为食言者，要么固执。那么，圣咏作者的这句话：\Quote{我已发誓并决意遵守你正义的典章}，我们该如何理解？\Quote{发誓并决意}除了一成不变地持守自己的决定，还能是什么？\par

% structure: p0085 source=h2 target=subsection reason=relative-heading-rank; base=h2

\subsection{第二十八章}

\Emph{回答：在何种情况下决定必须坚定不移，在何种情况下如有必要可以背弃。}\par

\Person{Joseph}：我们并非针对那些根本性的诫命，因为没有它们我们的救恩根本无法存在；而是针对那些我们可以在不危及自身状态的情况下或放松或持守的事，例如不间断的严格斋戒、完全戒酒或戒油、完全禁止离开自己的斗室、持续专注于阅读与默想——所有这些都可以随意实践，无损于我们的职业和目标，必要时也可不受责备地放弃。但对于那些根本性的诫命，我们必须最坚决地决心遵守，甚至必要时为此殉难也在所不惜；对此我们必须坚定不移地宣告：\Quote{我已发誓并决意。}这样做是为了维护爱德，为此应忽略其它一切，免得其宁静之美与成全受到玷污。同样，我们应为贞洁的纯洁而发誓，也应为信德、节制和正义而如此行；这些我们都必须以不变的恒心持守，稍加背离便应受责备。但在那些身体锻炼的事上（经上说\Quote{操练身体益处甚少}，\BibleRef{弟前 4:8}），我们应按所说那样决定：倘若出现某个更确定的善行机会，使我们放松它们，我们就不必受任何规则束缚，可以放弃它们而自由采纳更有益的事。因为在身体锻炼方面，暂时放弃并无危险；但放弃另一些即使片刻也足以致命。\par

% structure: p0088 source=h2 target=subsection reason=relative-heading-rank; base=h2

\subsection{第二十九章}

\Emph{应当如何处理那些必须保密的事。}\par

你们同样须谨慎留意：倘若无意中有某句欲保密的话从你口中溜出，莫要让那听者因保密叮嘱而困扰。因为若是随意简单地说出，反而更可能不被留意；兄弟（无论他是谁）将不会受泄露的诱惑而苦恼，因为他会视之为闲聊中不小心脱口的小事，正因此事更不重要，因为它并非以严格保密叮嘱托付给听者。即使你以誓言束缚了他的忠诚，也不须怀疑它很快就会泄露；因为魔鬼会更加猛烈地攻击他，既要烦扰你、出卖你，又使他尽快背誓。\par

% structure: p0091 source=h2 target=subsection reason=relative-heading-rank; base=h2

\subsection{第三十章}

\Emph{不可就关乎共同生活需要的事做出任何决定。}\par

因此，隐修士不应急于对仅涉及身体锻炼的事做出任何承诺，以免刺激仇敌更猛烈地攻击他所遵守的（如同受律法约束的），从而更容易被迫打破它。因为凡生活在恩宠的自由之下，却为自己立了法的人，便将自己束缚于危险的奴役中；倘若偶然的迫使他做那本可合法地、甚至可赞地、怀着感恩去做的事，他被迫成为违法者，陷入罪恶：\Quote{因为没有律法，就没有过犯。}\BibleRef{罗 4:15}\par

藉着这训导和真福\Person{Joseph}的教导，我们如同得到神谕一般得到坚固，决心留在\Place{Egypt}。尽管此后我们对承诺不再那么忧虑，但七年届满时，我们很乐意履行它。因为我们急忙回到我们的隐修院，那时我们自信能获准返回旷野；首先我们适当地向长老们致敬；接着，我们重新燃起他们心中从前的爱（因为他们因热忱的爱并未因我们无数封安抚信而软化）；最后，我们完全除去了背誓之刺，返回\Place{Scete}旷野的幽深处，他们自己也欢欣地送我们前行。\par

这段学习和教义来自杰出的教父们，圣弟兄啊，我们的无知尽其所能地向你说明了。如果也许我们笨拙的文体反而使之混乱而非条理分明，我相信，我们笨拙应得的指责不会干扰对这些伟大人物应得的赞美。因为在我们看来，在我们的审判者面前，更稳妥的做法是即使用朴素的文体陈述这光辉的教义，而不是对此缄默不语，因为如果他考虑思想的高超，那么我们文体的笨拙使他烦恼这一事实，就不应损害读者的益处，而我们自己更关心它的实用性而非被赞美。至少我要嘱咐所有可能读到这本小册子的人：他们必须知道，其中令他们喜欢的任何部分都属于教父们，而他们不喜欢的都是我们自己的。\par

\SourceRecord{Translated by C.S. Gibson.；Nicene and Post-Nicene Fathers, Second Series；Vol. 11.；Edited by Philip Schaff and Henry Wace.；Buffalo, NY: Christian Literature Publishing Co.,；1894.；<http://www.newadvent.org/fathers/350817.htm>.}

% structure: p0001 source=h1 target=section reason=page-title

\section{第十八讲}

\Emph{圣若望·卡西安作}\par

院长皮亚蒙的讲论。论三种隐修士。\par

我靠基督的恩宠之助发表了十篇教父讲论，这些讲论是应最受祝福的\Person{赫拉狄乌斯}和\Person{莱昂提乌斯}的急切请求而编撰的；之后我又将另外七篇献给\Person{奥诺拉图斯}（一位名实相符的主教）以及基督的神圣仆人\Person{欧凯里乌斯}。现在我认为同样数量的讲论也应奉献给你们，神圣的弟兄们：\Person{约维尼安努斯}、\Person{米涅尔维乌斯}、\Person{莱昂提乌斯}和\Person{特奥多雷}。因为你们中的最后一位在高卢省以古代德行的严格创立了那座神圣而辉煌的隐修会规，而你们其余几位则通过你们的教诲激励隐修士们不仅首先追求集体隐修院的共同生活，甚至热切渴望那崇高的隐居者生活。这些最优秀教父的讲论编排如此精心，各方面考虑周全，以至于它们适用于两种生活模式，你们藉此不仅使西方各国，甚至使各岛屿因众多弟兄的群体而繁荣；也就是说，不仅那些仍以值得称赞的服从生活在集体中的人，而且那些退隐到离你们隐修院不远、试图实行隐居者会规的人，都能根据地方的性质和其状况的特点得到更充分的教育。你们之前的努力和劳作对此尤其做出了贡献：他们既然已在这些操练中有所准备和练习，就能更容易地接受长者的规诫和规范，并将讲论的作者连同讲论本身的卷帙一同带进他们的静室，仿佛通过每日问答的方式与他们交谈，从而不至于独自面对那条艰难且在此地几乎不为人所知的道路——即使在那些不乏旧路和无数先行者的地方，这条道路也充满危险——而是学习遵循由那些古代传统、勤奋和长期经验已完全教导过的榜样所传授的隐居者生活规范。\par

% structure: p0005 source=h2 target=subsection reason=relative-heading-rank; base=h2

\subsection{第一章}

\Emph{我们如何来到狄奥尔科斯并受到院长皮亚蒙的接待。}\par

在拜访了那三位长者并与他们交谈之后（我们应兄弟\Person{欧凯里乌斯}的请求尝试记述了他们的讲论），我们仍更加热切地渴望探访埃及的更远地区，那里居住着更大、更成全的圣人团体。我们来到一个名为狄奥尔科斯的村庄，它位于尼罗河的七条支流之一。当我们听说那里有许多由古代教父建立的著名隐修院时，如同最热切的商人，我们立刻踏上了不确定的旅程，被更大收获的希望所催促。我们在那里徘徊了一段时间，用好奇的目光审视那些因崇高高度而显眼的德行之山，周围人的目光首先锁定了院长\Person{皮亚蒙}，他是当时居住在那里的所有隐居者中的长者兼他们的司铎，仿佛一座高大的灯塔。他被安置在一座高山上，就如福音中那座城\BibleRef{玛 5:14}，立刻将他的光芒照在我们脸上；他的德行和奇迹（在我们眼皮底下由他施行，神圣恩宠以此证明他的卓越），如果我们不打算超出本书的计划和范围，就必须沉默地略过。因为我们承诺记下的是所能回忆的，不是天主的奇迹，而是圣人的规范和追求，以便仅为读者提供成全生活的必要教导，而非提供无聊无益的惊叹而无任何过失的改正。因此，当院长\Person{皮亚蒙}愉快地接待了我们，并以恰当的和善款待我们之后，因为他意识到我们不是同一乡里的人，他首先关切地问我们从哪里来、为何访问埃及；当他发现我们出于对成全的渴望从叙利亚的一座隐修院来到这里时，他开始如下讲话：——\par

% structure: p0008 source=h2 target=subsection reason=relative-heading-rank; base=h2

\subsection{第二章}

\Emph{院长皮亚蒙的讲话：初学隐修士应当如何被长者的榜样教导。}\par

我的孩子们，无论谁渴望在任何技艺上达到熟练，除非他以最大的努力和细心投入到学习他想掌握的那套体系中，并遵守那项工作或科学的最佳大师的规则和秩序，否则他是在以虚妄的愿望徒劳地希望达到任何与他所避免模仿其辛劳和努力之人的相似。因为我们知道有一些人从你们的乡来到这些地方，只是为了走访隐修院以认识弟兄们，并无意采用他们远道而来所要学习的规则和条例，也无意识退回到静室并努力将所见所闻付诸实践。这些人保留了他们已习惯的性格和追求，正如有些人所指责的，他们被认为改变乡里不是为了自己的益处，而是由于逃避匮乏的需要。因为他们倔强固执的心态，不仅什么也学不到，甚至不愿再在这些地方停留。因为如果他们既没有改变禁食的方法，也没有改变圣咏的安排，甚至连衣服的样式都没有改变，那么我们除了认为他们在这个国家是为了获取粮食供应之外，还能认为他们是为了什么呢？\par

% structure: p0011 source=h2 target=subsection reason=relative-heading-rank; base=h2

\subsection{第三章}

\Emph{年幼者不应议论长者的命令。}\par

因此，如果我们相信，是天主的原因吸引你们努力效法我们的知识，那么你们必须完全忽视你们早期所受训练的一切规则，并要以极大的谦逊，追随你们所见我们的\OriginalTerm{长老}{Elders}所做或所教导的一切。即使暂时不明白任何行为或行动的原因或理由，也不要困惑、被牵引或偏离效法它，因为如果人们对一切事物持有良好而纯朴的观念，并急于忠实地效法他们所见他们的\OriginalTerm{长老}{Elders}教导或所做的一切，而不加以讨论，那么一切事物的知识将通过工作的经验而获得。但一个以讨论开始学习的人，永远不会进入真理的理性；因为仇敌见他信赖自己的判断超过父辈的判断，便轻易地鼓动他，直至那些特别有用和有益的事物在他看来变成不必要的或有害的；那狡猾的仇敌如此利用他的傲慢，使他固执地坚持自己的意见，说服自己只有他在执迷不悟的谬误中认为良善和正确的事物才是神圣的。\par

% structure: p0014 source=h2 target=subsection reason=relative-heading-rank; base=h2

\subsection{第四章}

\Emph{论埃及存在的三种隐修士。}\par

因此，你们应当首先听一听，我们的修会制度与起源是如何或从何而来的。因为一个人只有当他了解了他所希望学习的任何艺术的作者和创始人的荣耀时，他才能有效地被培养，并勤勉地实践它。埃及有三种隐修士，其中两种值得称赞，第三种则很差劲，必须完全避免。第一种是\OriginalTerm{隐修者}{Cœnobites}，他们共同生活在团体中，受一位\OriginalTerm{长老}{Elder}的指导；在整个埃及，大多数隐修士属于此类。第二种是\OriginalTerm{独修者}{anchorites}，他们先在隐修院中受训，然后在实践生活中达到成全后，选择了旷野的隐居处；我们也希望在这一行列中占有一席之地。第三种是可责斥的\OriginalTerm{撒拉巴依特}{Sarabaites}。我们将逐一更详细地论述它们。因此，正如我们所说，你们应当首先了解这三种修会的创始人。因为由此立刻会产生对应当避免的修会的憎恶，或对应当追随的修会的渴望，因为每一种方式必定将其追随者引向其创始人和发现者所达到的终点。\par

% structure: p0017 source=h2 target=subsection reason=relative-heading-rank; base=h2

\subsection{第五章}

\Emph{论创始隐修会秩序的创始人。}\par

因此，隐修会的制度起源于宗徒宣讲的时代。因为那时在\Place{耶路撒冷}的全部信徒群众，正如\WorkTitle{宗徒大事录}所描述的：\Quote{但信徒群众都一心一意，凡各人所有的，没有人说是自己的，却归公用。他们变卖财产和所有物，分给众人，按每人的需要。}又说：\Quote{因为他们中间没有一个缺乏的，因为凡有田地或房屋的，都变卖了，把变卖得来的钱，放在宗徒脚前，照每人的需要分配。}我说，整个教会当时就像现在隐修院中难得找到的少数人一样。但当宗徒们去世后，信徒群众开始冷淡，尤其是从各种外邦民族归信基督的群众；宗徒们顾念他们信仰的幼稚和根深蒂固的异教习惯，只要求他们\Quote{戒绝祭偶像之物、奸淫、勒死的牲畜和血}，\BibleRef{宗 15:29}因此，由于新诞生信仰的软弱而准许给外邦人的那种自由，逐渐损害了耶路撒冷教会的成全，早期信仰的热忱因本地人和外邦人的日益增多而冷却；不仅那些接受了基督信仰的人，就连教会的领袖也放松了某些严格。因为有些人认为，他们所见因软弱而准许给外邦人的，也允许他们自己；他们想，如果他们在追随基督的信仰和宣认的同时，保留自己的财产和所有物，也不会遭受损失。但那些仍然保持宗徒热忱的人，怀念那最初的成全，离开城市和与那些认为疏忽和松弛的生活对自己和天主的教会是许可之人的交往，开始生活在乡村和更僻静的地方；在那里私下按自己的方式，实践他们所学到宗徒们在整个教会中普遍命令的事。因此，我们所说的整个制度，是从那些与蔓延的邪恶分离的门徒中成长起来的。这些人随着时间推移，逐渐从广大信徒中分离出来；由于他们远离婚姻，断绝与亲属和世俗生活的交往，被称为隐修士或独修者，因他们孤独独居生活的严格。由此，由于他们的共同生活，被称为隐修者，他们的居所被称为隐修院。因此，那最早的隐修士类型，不仅在时间上首先，在恩宠上也首先，它持续了很长时期，直到\Person{保禄}院长和\Person{安当}的时代；至今我们在严格的隐修院中仍能看见其痕迹。\par

% structure: p0020 source=h2 target=subsection reason=relative-heading-rank; base=h2

\subsection{第六章}

\Emph{论独修者的体制及其起源。}\par

从这些成全者的数目中——如果可以这样表达——从这最丰硕的圣人之根上，后来也产生了\OriginalTerm{独修士}{anchorites}的花朵和果实。我们听说，这一修会的创始者正是我们刚才提到的那几位：圣保禄和圣安当。他们常出入旷野深处，并非像某些人那样出于胆怯和不耐烦的恶习，而是出于对更高\OriginalTerm{成全}{perfection}和神圣\OriginalTerm{默观}{contemplation}的渴望，尽管前者据说是因形势所迫才进入旷野，因为在迫害时期，他躲避邻人的阴谋。因此，从我们所说的那种制度中，产生了另一种成全的方式，其追随者被正确地称为独修士，即\Quote{退隐者}，因为他们绝不满足于在人间生活时所取得的胜过魔鬼隐秘陷阱的胜利，而是渴望与魔鬼公开交战、正面搏斗，因此他们不畏惧深入旷野的广阔幽深之处，效法洗者若翰——他一生都在旷野中度过——以及厄里亚、厄里叟和那些宗徒所描述的人：\BibleQuote{他们披着绵羊皮和山羊皮，处于贫乏、困苦、磨难中，这世界原不配他们，他们流落在旷野、山岭、山洞和地穴中。}\BibleRef{希 11:37-38} 关于这些人，主也以比喻对约伯说：\BibleQuote{谁放了野驴自由，解开了它的缰绳？我以旷野为它的家，以盐碱地为它的居处。它嗤笑城市的喧嚣，不听赶车人的吆喝；它环顾山间的牧草，寻找一切青绿的东西。}\BibleRef{约 39:5-8} 在圣咏中也说：\BibleQuote{愿上主所救赎的人这样述说，就是那些他从仇敌手中所救赎的人；}\BibleRef{咏 107:2} 稍后又说：\BibleQuote{他们在旷野中，在无水之地流荡，找不到可居住的城邑。他们又饥又渴，他们的灵魂在他们内衰颓。他们在患难中呼号上主，他即从困厄中拯救他们。}\BibleRef{咏 107:4-6} 耶肋米亚也这样描述：\BibleQuote{人自幼负轭，是有福的。他应孤独静坐，因为轭已加在他身上。}\BibleRef{哀 3:27-28} 他们在心中和行为上咏唱圣咏作者的话：\BibleQuote{我成了旷野中的鹈鹕；我警醒，成了屋顶上孤寂的麻雀。}\BibleRef{咏 102:6-7}\par

% structure: p0023 source=h2 target=subsection reason=relative-heading-rank; base=h2

\subsection{第七章}

\Emph{论\OriginalTerm{撒拉巴依特}{Sarabaites}的起源及其生活方式。}\par

当基督信仰因这两种隐修会士而欢欣时——尽管这一体制已逐渐衰落——后来又出现了那种令人厌恶且不忠信的隐修会士；更确切地说，那邪恶的植株重新萌芽并生长，这种植株在初期教会中首次出现在\Person{阿纳尼雅}和\Person{撒斐辣}身上时，便被宗徒\Person{伯多禄}的严厉所剪除。长期以来，这种隐修会士在隐修士中被视为可憎可恶的，只要信友们心中仍铭记着那极其严厉判决的恐惧——在判决中，蒙福的宗徒不仅拒绝让前述新罪行的始作俑者通过补赎或任何赔偿得到治愈，反而以他们的速死消灭了最危险的祸根。当这一在\Person{阿纳尼雅}和\Person{撒斐辣}身上以宗徒之严厉受罚的先例，由于长期疏忽和岁月流逝而逐渐从一些人心中淡忘时，便兴起了沙辣巴依特派。在埃及语中，他们被称为沙辣巴依特，这称号恰如其分，因为他们脱离了公修团体的集合，各顾各的事务。这些沙辣巴依特出自我们之前提到的那类人，他们并非真正追求福音的成全，而是出于竞争或人们对完全弃绝基督者胜过一切财富的赞美，才试图效仿而非真正立志。这些人内心软弱，假装极为良善，因必需被迫加入这一派别，他们渴望被视为隐修士之名，却不效仿隐修士的追求；他们完全不奉行纪律，不顺从长上之命，也不受传统教导而学习克制己意，或接受并恰当学习任何明智的分辨规则。他们只是公开宣誓弃绝，即在人面前以此为名，要么仍留在自己家中从事与先前相同的职业，只是冠以此名号；要么为自己建造小屋并称之为隐修院，在其中完全自由、自作主宰，从不服从福音的诫命——那些诫命禁止他们为日用的食物操心或为家务烦恼。只有那些完全舍弃今世一切财物、服从公修院上司的人，才能以毫无疑虑的信心履行这些诫命，因为他们不再承认自己是自己的主人。而那些如我们所说逃避隐修院严厉管束的人，两三人一起住在自己的小屋里，不愿接受院长的照管和管辖，而是主要安排自己随心所欲、四处游荡、任意而行，这些人日夜忙于日常事务，甚至比公修院的隐修士更甚，但并非怀着同样的信德和目的。因为这些沙辣巴依特这样做，不是要将劳动果实交给管家的意愿，而是为了积攒钱财。请看他们之间有何区别：公修院的隐修士不为明日忧虑，将最中悦天主的劳动果实奉献给天主；而沙辣巴依特不仅为明日，甚至为多年之后的事，怀着不信的忧虑，妄想天主要么虚假，要么无能，无法或不愿赐予他们应许的饮食和衣着。前者在所有祈祷中寻求的，是获得\OriginalTerm{一无所有}{\GreekText{ἀκτημοσύνην}}，即失去一切、恒久贫困；后者则要确保各种物资的丰裕。前者热切超越每日工作的规定，以便将隐修院圣事所需之外的一切，由院长分配于监狱、客房、病院或穷人；后者则将每日所余用于奢侈需求，或出于贪婪之罪而囤积。最后，即使我们承认他们不善意收集之物可以比我们所说的更好的方式处置，他们也未能达到良善和成全的功德。因为前者将这样的收入带入隐修院，每日报告并持续谦卑服从，以至于他们对自己劳动所得之权利也被剥夺，正如对自己本身的权利被剥夺一样，他们不断更新最初弃绝的热诚，每日剥夺自己劳动的果实；而后者却因向穷人施舍某事而自高自大，每日跌入罪恶。前者以忍耐和严格，虔诚持守他们一旦接受的规则，永不实现自己的意愿，每日被钉死于世俗，成为活的殉道者；后者则因意愿的冷淡而被投入地狱。这两种隐修士在这个省份几乎数量相当；但在其他省份——由于信德的需要，我被迫走访——我们发现第三类沙辣巴依特盛行，几乎只此一类。在\Person{瓦伦斯}皇帝时代，阿里乌谬信的主教\Person{路济乌斯}在位期间，我们带着救济品前往\Place{埃及}和\Place{提瓦依}的弟兄那里，他们因坚守公教信德而被流放到\Place{彭托斯}和\Place{亚美尼亚}的矿坑；我们虽在一些城市中发现了零星的公修院体制，却从未听说独修者之名。\par

% structure: p0026 source=h2 target=subsection reason=relative-heading-rank; base=h2

\subsection{第八章}

\Emph{论第四种隐修士。}\par

然而还有另一种、即第四类隐修士，近来在我们中间出现，他们以独修隐士的外表和形式自诩，早期似乎以短暂的热情寻求\OriginalTerm{隐修院}{Cœnobium}的成全，但很快就冷淡下来。由于他们不愿改掉旧日的习惯和过错，也不满足于长久背负谦卑与忍耐的轭，并蔑视服从长老的规矩，便寻求单独的住所，想要独自留下，好使因无人激怒他们，在人眼中被视为忍耐、温和、谦卑的。而这种安排——或者说这种冷淡——一旦占据了某人，就决不容他接近成全。因为这样，他们的过错非但未被根除，反而愈演愈烈，犹如某种致死的内在毒素，无人激惹它，却因隐藏愈深而愈深入内里，给病人带来不治之症。因为出于对各自住所的尊重，无人敢斥责独居者的过错，他们宁愿忽略这些过错也不愿医治。再者，德行不是通过隐藏过错，而是通过驱逐过错才得以产生的。\par

% structure: p0029 source=h2 target=subsection reason=relative-heading-rank; base=h2

\subsection{第九章}

\Emph{问题：\OriginalTerm{隐修院}{Cœnobium}与\OriginalTerm{隐修院}{monastery}之间有何区别？}\par

格尔马努斯：\OriginalTerm{隐修院}{Cœnobium}与\OriginalTerm{隐修院}{monastery}之间是否有区别？还是说两者同指一物？\par

% structure: p0032 source=h2 target=subsection reason=relative-heading-rank; base=h2

\subsection{第十章}

\Emph{答复。}\par

皮亚蒙：虽然许多人随便地把\OriginalTerm{隐修院}{monastery}称为\OriginalTerm{隐修院}{Cœnobium}，但两者实有区别：\OriginalTerm{隐修院}{monastery}是住所的名称，仅指地方，即隐修士的居所；而\OriginalTerm{隐修院}{Cœnobium}则描述生活的性质及其制度。\OriginalTerm{隐修院}{monastery}可以指单个隐修士的住所，而\OriginalTerm{隐修院}{Cœnobium}则只能指一大群人共同生活的联合团体。然而，那些萨拉巴依特隐士群居之处也被称为\OriginalTerm{隐修院}{monastery}。\par

% structure: p0035 source=h2 target=subsection reason=relative-heading-rank; base=h2

\subsection{第十一章}

\Emph{论真谦卑，以及院长塞拉皮翁如何揭露某人的假谦卑。}\par

因此，既然我看到你们从最优秀的那类隐修士——即从\OriginalTerm{隐修院}{Cœnobium}的卓越学校出发，正朝着独修隐士规矩的高峰迈进——你们就应以真心的情感追求谦卑与忍耐的德行，我相信你们在那里学到了这德行；而不要像某些人那样，通过言语上的假谦卑，或通过某些身体劳作上造作而不必要的准备来假装谦卑。这种假谦卑，院长塞拉皮翁曾极其出色地嘲笑过。因为有一个人来到他那里，通过衣着和言语大肆炫耀自己的卑微，老人按惯例请他\Quote{领祷}，他却不肯依从要求，反而贬低自己，声称自己犯有种种罪恶，甚至不配呼吸与众人共有的空气，连草席也不肯用，宁可直接坐在地上。但当他在洗脚一事上表现得更加不情愿时，塞拉皮翁院长在晚餐结束后，利用惯例的会谈机会，温和而善意地劝他不要以轻浮的心意到处游荡，无所事事，尤其在他年轻力壮之时，而应依照长老的规矩留在自己的住所，选择靠自己的劳作而非他人的施舍维生；这一点连宗徒保禄也不允许，尽管他在为福音劳苦时，本可轻易获得供给，但他却宁愿日夜劳作，为自己和那些侍奉他却不能亲手工作的人提供日粮。对此，那人竟充满厌烦和憎恶，以致无法从脸色掩饰内心的恼怒。老人对他说：我的儿子，到现在为止，你一直用各种罪过的重负压在自己身上，不惧怕因承认如此可怕的罪而给名誉带来耻辱；请问，如今我简单劝诫你，其中并无责备，只是出于为你建树和爱德的感受，为何你竟如此恼怒，以致无法从脸色掩饰，也无法以平静的外表隐藏？或许你贬低自己时，正希望从我们口中听到这句话：\BibleQuote{义人在发言之初就控告自己。}\BibleRef{箴 18:17}此外，必须持守内心的真谦卑，它并非来自身体和言语上的造作谦卑，而是来自灵魂内在的谦卑；只有当人不再夸耀无人会相信的罪过，反而在他人无礼指责时无动于衷，并以温和的内心平静忍受所受的委屈时，这种谦卑才会以忍耐的明证彰显出来。\par

% structure: p0038 source=h2 target=subsection reason=relative-heading-rank; base=h2

\subsection{第十二章}

\Emph{问题：如何获得真正的忍耐？}\par

格尔马努斯：我们希望知道，那种平静如何能获得并保持，以至于当我们被命令静默时，便闭上口门，禁止言语，同样地，我们也能保持内心的温和——有时即使舌头被约束，内心却失去了平静状态。因此，我们认为温和之福只能由住在偏远住所和独居处的人才能保持。\par

% structure: p0041 source=h2 target=subsection reason=relative-heading-rank; base=h2

\subsection{第十三章}

\Emph{答复。}\par

\Person{皮亚蒙}：真正的忍耐与宁静，若没有深切的内心谦逊，既不能获得，也不能保持；若它源自这一根源，便既不需要独修室的好处，也不需要旷野的隐居。因为若它拥有其母与守护者——谦逊之德的内在支持，它就不会寻求任何外在的帮助。但若我们被任何人攻击时感到不安，就表明谦逊的根基尚未在我们内稳固奠定，因此即使爆发一场小风暴，我们的整座建筑也会摇撼，遭到破坏性的扰动。因为忍耐若只在没有敌人箭矢攻击时才保持其目的性的宁静，就不值得赞美和钦佩；但它是崇高而光荣的，因为当诱惑的风暴击打它时，它仍岿然不动。因为人在何处被认为因逆境而烦恼和受伤，就在何处更加坚强；他在何处被认为烦恼，就在何处更多地受锻炼。因为人人都知道，忍耐之名来自激情和忍受，所以显然，只有能无烦恼地忍受加给他的一切侮辱的人，才能被称为有忍耐的人。因此，\Person{所罗门}赞美他不是没有理由的：\Quote{忍耐者优于壮士，控制怒气者胜于攻取城池；}又说：\Quote{宽宏的人大有明智，懦弱的人极其愚蠢。}因此，当任何人被冤屈所压倒，怒火中烧时，我们不应认为加给他的侮辱的苦涩是他罪的\Emph{原因}，而更应是隐藏软弱的\Emph{表现}，这符合我们的主和救主关于两座房屋的比喻，祂说一座建在磐石上，另一座建在沙土上，风雨洪水和风暴同样冲击两者；但建在坚固磐石上的那座，冲击的猛烈丝毫未伤及它，而建在松散流动沙土上的那座则立刻倒塌。显然，它倒塌不是因为被风暴和洪水的冲击所击打，而是因为不明智地建在沙土上。因为圣人与罪人的区别不在于他同样不受诱惑，而在于即使面对强大的攻击他也不被击败，而后者连轻微的诱惑也会使他跌倒。因为正如我们所说，任何好人的刚毅若没有受到诱惑而获胜，就不值得赞美，因为最确定的是，没有争斗和冲突就没有胜利的余地：因为\Quote{忍受试探的人是有福的，因为他既被证明，就必得生命的冠冕，这是天主应许给爱祂的人的。}\BibleRef{雅 1:12}按照宗徒\Person{保禄}所说的，\Quote{能力得以成全}，不是在安逸与享乐中，而是\Quote{在软弱中。}\Quote{因为看，}祂说，\Quote{我今天使你成为坚城、铁柱、铜墙，对抗全地，对抗犹大君王、首领、司祭和当地的人民。他们要攻击你，却不能战胜你，因为我与你同在，上主说，要拯救你。}\par

% structure: p0044 source=h2 target=subsection reason=relative-heading-rank; base=h2

\subsection{第十四章}

\Emph{论一位虔诚妇女所立的忍耐榜样}\par

关于这忍耐，我现在至少给你两个例子：第一个关于一位虔诚的妇女，她如此急切地追求忍耐的德行，以至于她不仅不回避诱惑的袭击，反而为自己制造烦恼的时机，为的是更频繁地受考验。因为这位妇女住在\Place{亚历山大}，出身并非低微，且虔诚地在她父母留给她的家中服侍主，她来到可敬的\Person{亚纳削}主教那里，恳求他给她另一个需要供养的寡妇——这寡妇是教会出资供养的。并且，用她自己的话来表达她的请求：\Quote{给我，}她说，\Quote{一个姊妹来照顾。}当时主教称赞了这妇女的用意，因为他看出她非常乐意从事慈悲的工作，于是他下令从全体寡妇中挑选一位，其品德之良善、生活之庄重有序超越所有人，唯恐她施恩的愿望反被接受者的过错所破坏，而这位想从穷人身上获益的人会因那寡妇的恶劣品性而厌恶，从而在信德上受损。当那寡妇被带回家后，她以各种服事照顾她，发现她极其谦逊温良，并看到她每一次因善心服务都得到感谢，于是在几天后，她回到上述主教那里，说：我请求你吩咐给我一个女人，让我供养并以顺从的殷勤服事她。当他尚未理解这妇女的目的和愿望，以为她的请求被上司的欺骗所忽视，于是带着内心的些许愤怒询问拖延的原因时，他立刻发现分配给她的是全团体中最好的寡妇，于是他秘密下令给她一个最差的，就是说，一个在愤怒、争吵、酗酒和多言方面超过所有被这些毛病控制的人。当这位最差的很容易被找到并交给她后，她便开始留她在家里，以同样甚至更大的细心照顾她，而她从她那里得到的服务报酬仅仅是：不断遭受不应有的伤害，不断被辱骂和指责所困扰，因为她抱怨她，用恶毒贬损的话责备她，说她把她从主教那里要来，不是为了她的休息，而是为了她的折磨和烦恼，把她从休息带到劳苦，而非从劳苦带到休息。当这持续的责备如此爆发，以至于那个放荡的妇女不惜对她动手时，另一位则更加倍地服务，以更卑微的服事，并学会不是通过抵抗，而是通过更谦卑地自抑来战胜那个悍妇，以致当她被各种侮辱激怒时，她能用温和与仁慈安抚那泼妇的疯狂。当她通过这些操练被彻底坚固，达到了她所渴望的忍耐的完美德行时，她来到上述主教面前，感谢他的决定和选择，以及她操练的祝福，因为他终于如她所愿地为她的忍耐提供了最宝贵的女主人，通过这位女主人持续的烦扰，如同被摔跤用的油每日加强，她达到了心灵的完全忍耐。最后，她说，你终于给了我一个要供养的人，因为前一位更像是以她的服务尊敬我、使我得到休息。这些关于女性性别的叙述已经足够，通过这个故事，我们不仅得到建造，甚至感到羞愧，因为我们无法保持忍耐，除非我们像野兽一样被隔离在山洞和静室中。\par

% structure: p0047 source=h2 target=subsection reason=relative-heading-rank; base=h2

\subsection{第十五章}

\Emph{关于\Person{帕弗努息}院父所给的忍耐的榜样。}\par

现在让我们举出帕夫努提乌斯院长另一个实例。他始终如此热忱地留在那著名且远扬的斯刻提斯旷野的深处，如今他在那里担任司铎，以至于其他隐修士称他为布巴利斯，因为他仿佛带着一种天生的喜好，始终喜好在旷野中居住。他自少年时代就如此良善且充满恩宠，以至于当时那些著名和伟大的人物都钦佩他的庄重和坚定的恒心；尽管他年龄较小，却因德行受到尊重，被置于与长老同等的地位，且被认为适宜被接纳到他们的行列中。那从前曾激起弟兄们对圣祖若瑟嫉妒的同一嫉妒，在他一位弟兄的心中燃起了炽热而吞噬的嫉妒之火。此人想用某种瑕疵或污点玷污他的美德，便想出了这样一种恶毒的计谋：趁帕夫努提乌斯在星期天离开他的小室去教堂的机会，偷偷进入他的小室，狡猾地将他自己的书藏在帕夫努提乌斯通常用棕榈枝编织的树枝中，然后自信计谋已定，自己仿佛怀着一颗纯洁无瑕的良心前往教堂。当整个礼仪照常结束后，他在众弟兄面前向这旷野的司铎圣依西多尔（他先于帕夫努提乌斯担任此职）提出控诉，声称他的书从小室中被盗。他的控诉扰乱了所有弟兄，尤其是司铎的心绪，使他们不知首先怀疑或想什么，因为众人都因这如此新奇且前所未闻的罪行而极度震惊——在此以前无人记得这旷野中曾发生过这样的事，此后也再未发生。提出此事作为控告者的人敦促将所有人留在教堂，并选派若干选定的弟兄逐个搜查各人的小室。此事由司铎托付给三位长老后，他们翻遍了所有人的卧室，最终在帕夫努提乌斯的小室中，在棕榈枝（他们称之为\OriginalTerm{编织物}{\GreekText{σειρά}}）中找到了那本书，正如设计者所隐藏的那样。当搜查者立即将书带回教堂，并在众人面前出示时，帕夫努提乌斯，尽管他良心的纯真完全清白，却像一个承认盗窃之罪的人，完全投身于补赎，并谦卑地请求一个忏悔的计划，因为他如此珍视自己的羞耻和谦逊，深怕若试图用言语洗刷盗窃的污点，还会被指为说谎者，因为除了已发现的事实外，无人会相信其他。他立即离开教堂，不是心情沮丧，而是信靠天主的审判，他在祈祷中不断流泪，禁食比从前加倍频繁，并在人前以全心的谦卑俯伏在地。当他这样怀着身灵的全然痛悔，几乎两周之久如此顺服，以至于在星期六和星期日的清晨，他前来不是领受圣体，而是俯伏在教堂门槛上，谦卑地请求宽恕时，那洞察并知晓一切隐秘的主，不再容许他继续受试炼或被他人诽谤。因为那犯罪者——自己财物的邪恶窃贼、他人信誉的狡猾诽谤者——所行的事，在无人见证之下，主却藉着那罪恶本身的煽动者——魔鬼——使之显明。因为那人被最凶猛的恶魔附身，吐露了他隐秘计谋的一切诡诈，那原本捏造控告和欺诈的人自己揭露了它。但他被那不洁之神如此长久而痛苦地折磨，以至于甚至无法通过居住在那里的圣人们的祈祷（他们藉着神圣恩赐能命令魔鬼）而康复，连司铎圣依西多尔本人的特殊恩宠也无法从他身上驱除那残酷的折磨者，尽管藉着主的恩惠，他获得了如此能力，以至于任何被附身的人被带到他的门前，无不即刻痊愈；因为基督将这一荣耀保留给年轻的帕夫努提乌斯，使那人只有通过他所谋害之人的祈祷才能被洁净，且那嫉妒的仇敌只有宣告那人的名字（他原以为能中伤其名誉）才能获得对其过犯的宽恕和当前惩罚的终结。他于是在早年青春时代已显出这些未来品格的迹象，甚至在少年岁月中便勾勒了那将在成熟年龄成长起来的成全的轮廓。因此，若我们想达到他那样的德行高度，就必须打下同样的基础。\par

% structure: p0050 source=h2 target=subsection reason=relative-heading-rank; base=h2

\subsection{第十六章}

\Emph{论忍耐的成全}\par

然而，促使我讲述此事的理由有双重：其一，让我们掂量此人的坚定与恒心，当我们受到敌人较轻的诡计攻击时，能更好地保持更大的平静与忍耐；其二，让我们以坚决的决心持守：若我们将自己平安的一切保障和所有信心，不是建立在内在之人的力量上，而是建立在斗室之门、沙漠的隐蔽处、圣人的陪伴或任何外在事物的庇护上，我们就无法免受诱惑的风暴和魔鬼的攻击而安全。因为，除非我们因祂保护之力的德能而坚强——祂在福音中说：\BibleQuote{天主的国就在你们心里}\BibleRef{路 17:21}——否则，我们若幻想能借与我们同住之人的帮助打败空中敌人的阴谋，或借地点的距离避开它们，或借墙壁的保护排除它们，都是徒然。因为这些对圣帕弗努提乌而言一样都不缺，但诱惑者还是找到了攻击他的途径；环绕的墙壁、沙漠的孤寂以及会众中所有那些圣人的功德，都无法击退那最不洁的恶魔。然而，因为天主的圣仆将他内心的希望不是寄托于外在事物，而是寄托于鉴察一切隐秘的那一位，他即使面对如此攻击的阴谋也不会动摇。另一方面，那被嫉妒驱入如此严重罪恶的人，岂不也享受了孤寂的好处、隐修居所的保护，以及与蒙福的院长兼司铎依西铎及其他圣人的交往？然而，因为魔鬼在他身上掀起的风暴发现他是在沙土上，不仅冲垮了他的房屋，而且彻底颠覆了它。因此，我们不必在外物中寻求平安，也不必幻想别人的忍耐能弥补我们急躁的过失。因为正如\BibleQuote{天主的国就在你们心里}，同样，\BibleQuote{人的仇敌就是自己的家人}\BibleRef{玛 10:36}。因为没有人比我的内心更是我的仇敌，它确实是我最亲近的家人。因此，若我们谨慎，就不可能被内在的敌人伤害。因为当我们的家人不与我们为敌时，那里天主的国也在心灵的平安中得以稳固。因为，若你仔细探究此事，我不能被任何怨恨我的人伤害，除非我以好战之心与自己争斗。但我若受伤，过错不在于他人的攻击，而在于我自己的急躁。因为正如强壮而坚固的食物对健康的人有益，对生病的人则有害。但它不能伤害取用之人，除非接受者的软弱赋予了它伤害之力。因此，若任何类似的诱惑在弟兄们中出现，我们绝不可动摇我们平常的步调，给世俗之人亵渎的咆哮以可乘之机，也不可惊奇有些邪恶可憎的人秘密混入圣人的行列，因为只要我们在这世界的打谷场上被践踏、被踩踏，注定被投入永火的糠秕必然与精选的麦子混杂。最后，若我们记起撒殚是在天使中被拣选的，犹达斯是在宗徒中被拣选的，而尼古拉（可憎异端的创始人）是在执事中被拣选的，那么在圣人的行列中发现最卑劣的人就不足为奇了。虽然有人认为此尼古拉并非被宗徒们选为执行职务者，但他们无法否认他曾是门徒中的一员——而这些门徒在品性和成全上，明显与我们现在在隐修院中难以发现的少数人无异。让我们因此提出的，不是上述那位弟兄的跌倒（他竟在沙漠中如此严重地倒塌），也不是他后来用大量忏悔的眼泪洗去的可怕污点，而是蒙福的帕弗努提乌的榜样；我们不要因前者的堕落而毁灭（他根深蒂固的嫉妒之罪因他伪装的虔诚而加剧并恶化），但让我们尽全力效法后者（帕弗努提乌）的谦卑——他的谦卑并非由沙漠的寂静突然产生，而是在人群中已获得，并由独处生活得以成全和完美。然而，你应当知道，嫉妒之恶比其他毛病更难治愈，因为我几乎可以说，一旦被它毒素的祸害玷污的人无药可救。因为这是先知借比喻所说的灾祸：\BibleQuote{看，我必打发毒蛇和蛇蜥到你们中间，是没法念咒的，它们要咬伤你们}\BibleRef{耶 8:17}。因此，先知将嫉妒的刺比作蛇蜥的致命毒液是恰当的，因为一切毒素的首创者和首领就是因它而丧亡死亡。因为他先于他所嫉妒的人杀死了自己，在向人倾倒死亡的毒液之前先毁灭了自己：因为\BibleQuote{因魔鬼的嫉妒，死亡进入了世界；凡属于他的人就效法他}\BibleRef{智 2:24}。正如那首先被这邪恶瘟疫腐蚀的人，不接受任何忏悔的补救或医治的药膏，同样那些让自己受到同一毒刺刺伤的人，也排除了神圣施咒者的一切帮助；因为他们所受的折磨，不是因他们所嫉妒之人的过错，而是因这些人的兴旺，他们羞于显露真实情况，反而寻找一些外在的、不必要的、微不足道的冒犯理由；这些理由完全是假的，故治愈的希望是徒然的，而他们不愿排出的致命毒液却潜伏在血管中。智者对此恰当地说：\BibleQuote{蛇不出声而咬人，念咒的便毫无效力}\BibleRef{训 10:2}。因为那些无声的咬伤，只有智者之药无法救治。因为这邪恶如此无可救药，以至于关切使它变得更糟，服事使它增加，礼物激怒它，因为如同同一撒罗满所说：\BibleQuote{嫉妒什么也容忍不了}\BibleRef{箴 27:4}。因为随着别人在谦卑顺服、忍耐之德或慷慨之功德上的进步，那人就会因更强烈的嫉妒之刺而激动，因为他所求的莫过于他所嫉妒之人的毁灭或死亡。最后，那十一位族长对无辜兄弟的任何顺服都未能软化他们的嫉妒，以至于经上论到他们说：\BibleQuote{他的兄弟们见父亲爱他，就嫉妒他，不能与他和平相处}\BibleRef{创 37:4}，直到他们的嫉妒（对那顺服和谦逊的兄弟的一切恳求充耳不闻）渴望他死亡，并且几乎连出卖兄弟的罪都不能满足。可见，嫉妒比一切毛病都更糟糕，更难以根除，因为它被那些使其他毛病毁灭的补救反而点燃。例如，一个因自己所受损失而忧伤的人，可通过丰厚的补偿而治愈；一个因所受伤害而痛苦的人，可通过谦卑的赔礼而平息。你能对一个因见你更谦卑、更仁慈而越发被冒犯的人做什么呢？他并非被任何可用贿赂平息的贪欲所激怒，也非被任何伤害性的攻击或复仇之心（这可通过服侍而克服）所激动，而只是被别人的成功和幸福所激怒。然而，有谁为了满足一个嫉妒他的人，愿意放弃自己的好运、失去自己的兴隆或陷入某种灾祸呢？因此，我们必须不断恳求天主助佑——对祂而言没有不可能的事——以避免蛇凭借这邪恶的一咬，摧毁我们内里一切蓬勃生长的、仿佛被圣神的生命和活力所滋养的事物。因为蛇的其他毒液，即肉性的罪和过失（人性软弱易陷其中，也易从中净化），在肉身上会留下一些伤口的痕迹，尽管这尘世的身体会被最危险地灼烧，但若任何精通神圣咒语的施咒者施以治疗和解毒剂，或施以得救话语的补救，这毒害的邪恶不会触及灵魂的永死。但嫉妒之毒，仿佛由蛇蜥喷出，甚至在身体感到伤口之前，就摧毁了宗教信仰的生命。因为那些忌恨兄弟时只挑剔他的幸福，不挑剔任何过错，从而不指责人的毛病，而只指责天主的判断——这样的人不是反对人，而是亵渎地反对天主。因此，这就是那\BibleQuote{生出苦根来}\BibleRef{希 12:15}的，它高升至天上，倾向于责备那将恩赐赐予人的施主本身。任何人不当因天主威胁要打发\BibleQuote{毒蛇和蛇蜥}\BibleRef{耶 8:17}来咬那些以罪行冒犯祂的人而感到不安。因为虽然确知天主不能是嫉妒的源头，然而这符合天主的公义且与之相称：当恩赐赐予谦卑者，而拒绝给骄傲和可弃绝之人时，那些按照宗徒的话\BibleQuote{配受邪僻之心}\BibleRef{罗 1:28}的人，应被祂如同派遣的嫉妒所咬伤和消灭，正如这段经文所说：\BibleQuote{他们以不是神的激起了我的妒火，我也以不是子民的激起他们的妒火}\BibleRef{申 32:21}。\par

通过这篇讲话，真福\Person{丕亚孟}更加激起了我们的渴望，我们已开始从隐修院的初级学校提升到隐修生活的第二标准。因为正是在他的指导下，我们初次开始了独居生活，后来我们在\Place{斯塞提斯}更深入地学习了这一知识。\par

\SourceRecord{Translated by C.S. Gibson.；Nicene and Post-Nicene Fathers, Second Series；Vol. 11.；Edited by Philip Schaff and Henry Wace.；Buffalo, NY: Christian Literature Publishing Co.,；1894.；<http://www.newadvent.org/fathers/350818.htm>.}

% structure: p0001 source=h1 target=section reason=page-title

\section{\WorkTitle{第十九次讲论}}

\Emph{圣若望·卡西安 著}\par

约翰院长讲论集。论\OriginalTerm{会院修士}{Coenobite}与\OriginalTerm{独修士}{Hermit}的目标。\par

% structure: p0004 source=h2 target=subsection reason=relative-heading-rank; base=h2

\subsection{第一章}

\Emph{论\Person{保罗}\OriginalTerm{院长}{Abbot}的\OriginalTerm{隐修院}{Cœnobium}与某位兄弟的忍耐}\par

我们提到这次集会，是为了简要谈谈某位兄弟的忍耐，他在全体会众面前表现出了不可动摇的温柔。因为虽然这部作品的目标涉及另一个人，即我们可以呈现\Person{约翰}\OriginalTerm{院长}{Abbot}的言论，他离开了旷野，以极大的良善和谦卑使自己服从于那\OriginalTerm{隐修院}{Cœnobium}，但我们认为，毫不荒谬地以最简洁的语言叙述我们认为对渴求良善的人最有教益的事。于是，当全体修士分成每十二人一组坐在宽阔的庭院中时，有一位兄弟在上菜时动作较慢，上述保罗院长正在忙碌地穿梭于侍奉的弟兄们中间，看见了，就当众在他摊开的掌心上打了一下，那手掌被击打的声音甚至传到了那些背对着、坐在远处的人耳中。但这位极具忍耐的青年以如此平静的心态承受了它，不仅没有吐出一个字，也没有因嘴唇无声的动作而有丝毫抱怨的迹象，甚至脸色丝毫未变，或（失去）唇边谦和宁静的表情。这一事实不仅令我们——我们刚从\Place{叙利亚}的一座隐修院来，尚未通过如此明显的榜样学习到这种忍耐的祝福——感到惊讶，也令所有那些并非没有这种热诚经验的人感到惊讶，以至于它给那些甚至已经很有修为的人上了一堂重要的课，因为即使这种慈父般的纠正没有扰乱他的忍耐，那么多人的在场也没有给他的脸颊带来丝毫颜色的变化。\par

% structure: p0007 source=h2 target=subsection reason=relative-heading-rank; base=h2

\subsection{第二章}

\Emph{论\Person{约翰}\OriginalTerm{院长}{Abbot}的谦逊与我们的问题}\par

那时，我们在\Person{保罗}院长的\OriginalTerm{隐修院}{Cœnobium}中看到这位老人，首先被他的年龄和那人所拥有的恩宠所打动，于是目光低垂，开始恳求他惠允向我们解释，为什么他放弃了旷野的自由和那崇高的事业——他在其中声誉卓著，超越其他选择同样生活的人——，又为什么他选择进入隐修院的轭下。他说，由于他不能胜任\OriginalTerm{独修士}{anchorite}的体系，也不配达到如此完美的境界，所以他回到了初级学校，为的是学会按照生活的要求来实践那里教授的课程。当我们的恳求未能满足，我们拒绝接受这种谦卑的回答时，他终于开始如下所述。\par

% structure: p0010 source=h2 target=subsection reason=relative-heading-rank; base=h2

\subsection{第三章}

\Emph{\Person{约翰}\OriginalTerm{院长}{Abbot}的回答：他为何离开旷野}\par

\OriginalTerm{独修士}{anchorite}的体系，你们惊奇我离开，我不但不拒绝、不抛弃，反而拥抱并极其崇敬它。我曾在\OriginalTerm{隐修院}{Cœnobium}中生活了三十年，之后又在该体系中度过了二十年，我为此感到喜乐，以至于我永远不会被那些半心半意尝试它的人指责为懈怠。但是，由于它的纯洁——我曾对此略有体验——有时会被对属世之事的焦虑所玷污，所以似乎最好是回到隐修院，以便更易于达到所承担的更容易的目标，并且比冒险过卑微独修者的更高生活风险更小。因为宁可在小诺言上显得恳切，也不要在大事上轻率。因此，如果我说得有些傲慢，甚至有些过于直率，我恳求你们不要以为是出于夸耀的罪过，而是出于我希望你们受益；并且，既然我想，当你们如此恳切地问时，不应向你们隐瞒任何真理，请你们将它归因于爱德，而非夸耀。因为我想，如果我放下谦卑，单纯地揭示我目标的全部真理，就能给你们一些教导。因为我相信，我既不会因直言而从你们那里招致虚荣的责备，也不会因隐瞒真理而良心上承担虚假的指控。\par

% structure: p0013 source=h2 target=subsection reason=relative-heading-rank; base=h2

\subsection{第四章}

\Emph{论上述老者在\OriginalTerm{独修士}{anchorite}体系中所展现的卓越}\par

若有人喜爱沙漠的隐僻，愿忘记一切人际交往，并与耶肋米亚一同说：\Quote{我不曾渴慕人的日子，你都知道，}\BibleRef{耶 17:16}我承认，因着天主恩宠的祝福，我也获得了，或至少努力获得了这一点。因此，藉着主的恩赐，我记得我常常被提升到这样的\OriginalTerm{出神}{ecstasy}之中，以至于忘记自己披戴着软弱肉身的重担，我的灵魂突然忘记了一切外在的概念，完全断绝与一切物质对象的联系，以至于我的眼睛和耳朵都不再执行其本有的功能。我的灵魂如此充满神圣的默想和属神的沉思，以至于常常在傍晚时我不知是否用过餐，次日也甚怀疑昨天是否守了斋。为此，一周的食粮，即七组饼干，被存放在一种手提篮中，在周六预备好，这样便不致怀疑何时错过了晚餐；这个计划也避免了因遗忘而产生的另一种错误，因为当饼数用完时，便表明一周的周期已过，同一日子的礼仪已来临，而集会的主日、庆节和礼仪便不会逃过隐修者的注意。但即使我们所说的那种心神出神偶尔会干扰这一安排，但每日的工作方式仍会显示日数并纠正错误。至于沙漠的其他益处（因为我们并非要讨论其数目和分量，而是要探讨独居和\OriginalTerm{修院}{Cœnobium}的目标），我将简要说明我为何选择离开它，这也是你们想要知道的，并将以简明的论述略微提及我所说的独居的所有果实，并表明它们在另一方面应被视为低于哪些更大的益处。\par

% structure: p0016 source=h2 target=subsection reason=relative-heading-rank; base=h2

\subsection{第五章}

\Emph{论沙漠的益处。}\par

那时，由于当时居住在沙漠中的人数稀少，我们在更广阔的荒野中享有更大的自由，被提升到那些属天的出神中，并未因大量来访的弟兄而受压迫，因此因着接待客旅的需要而受各种重大挂虑分心，我曾以难以满足的渴望和全心投入那些沙漠的宁静隐退之处，以及那种只堪比天使福乐的生活。但是，当如我所说，更大数量的弟兄开始寻求在那一沙漠中居住，并且因着拥挤而限制了广大荒野的自由，不仅使那神圣沉思之火变得冷淡，而且以许多方式将心灵缠缚于属肉之事物的锁链中，我决心在这体制中实行我的目的，而不愿在那崇高的生活方式中因提供属肉的需要而变得冷淡；因此，即便那自由和那些属灵出神被拒绝于我，但既然避免了对于明日的一切挂虑，我可以藉着履行福音的诫命而自慰，并且我在沉思崇高中所失掉的，可以通过服从和听命而得补偿。因为一个人宣称学习某种技艺或事业却从未达到其完美，是可悲的。\par

% structure: p0019 source=h2 target=subsection reason=relative-heading-rank; base=h2

\subsection{第六章}

\Emph{论\OriginalTerm{修院}{Cœnobium}的便利。}\par

因此，我将简要说明我在这种生活方式中如今所享有的益处。\newline{}\Emph{您}必须思量我的话，并判断旷野的那些益处是否胜过这些慰藉；由此你们也将能证明，我选择被局限在隐修院的狭窄界限内，究竟是出于对独修生活的厌恶还是渴望那种纯洁。\newline{}在此生活中，没有为日间的劳作操心，没有买卖的纷扰，没有对一年粮食的无可避免的牵挂，没有对肉身之事的焦虑——因此人不仅要为自己所需，还要为许多访客的需要预备必要的物品；最后，也没有来自人赞美的虚荣，这比所有这一切更糟，且在天主眼中，有时甚至抵消了旷野中巨大努力的好处。\newline{}但是，且不提独修士生活中属灵骄傲的波涛和虚浮荣耀的致命危险，让我们回到这个影响所有人的普遍负担，即：供应食物的日常焦虑。这焦虑已远远超出——我不说那完全不用油的古代严格尺度——甚至开始对当代的宽缓标准也不再满足；按照那标准，一年所有食物的需求本可由预备的一品脱油和一摩底豆子（为访客之用）而得到满足；但现在，所需食物的供应几乎难以由两三倍之数满足。\newline{}这一危险松懈的力量已在某些人中增长到如此程度：当他们混合醋和酱汁时，他们不添加那一滴油——我们的前辈（他们以更强的苦修能力遵循旷野的规则）为了躲避虚浮荣耀而惯常只倒入一滴——反而为了奢侈而掰碎埃及奶酪，倒上超出所需的油，以此在一道美味调品之下享用两种不同风味食物；每种食物单独在隐修士的各个时期本应是一道愉快的提神小食。\newline{}然而，\OriginalTerm{物质的获取}{\GreekText{ὑλικὴ} \GreekText{κτῆσις}}竟已发展到如此地步：实际上，以款待和欢迎客人为借口，独修士开始在他们的斗室里存放一条毯子——我提及此事都感到羞愧——更不用说那些更妨碍敬畏并专注于属灵默想之心的事情：即弟兄们的聚集、接待来客和送别离客的职责、相互拜访以及各种谈话和事务带来的无尽烦恼；由于这些习惯性干扰的持续性质，即使在那些麻烦似乎停止的时候，对其的预期也使心处于紧绷状态。\newline{}结果便是：独修生活的自由被这些束缚所阻碍，以至于它永远无法上升到那种无法言喻的心灵锐敏，从而失去了其隐修生活的果实。\newline{}即使现在我在团体中与众人同住，此事被剥夺，至少内心平安与免于一切事务的心灵安宁并不缺乏。\newline{}除非那些生活在旷野中也随时拥有这安宁，否则他们确实要承受独修生活的劳苦，却将失去那只能在心灵的平静安稳中获得的果实。\newline{}最后，即使我在隐修院中生活时，内心的纯洁有所减少，我也将满足于以交换持守那一条福音的诫命——它当然不比旷野的所有那些果实更为次要——即：我不应为明日忧虑，并且完全服从于院长，从而在某种程度上效仿\par

那一位，经上论祂说：\Quote{祂贬抑自己，听命至死。}如此我就能谦逊地使用祂的话说：\Quote{因为我从天降下，不是为执行我的旨意，而是为执行那派遣我的圣父的旨意。}\par

% structure: p0023 source=h2 target=subsection reason=relative-heading-rank; base=h2

\subsection{第七章}

\Emph{关于隐修院与旷野果实的问题。}\par

\Person{杰尔马努斯}说：既然您显然不像许多人那样只是触及每种生活方式的边缘，而是登上了其顶峰，我们愿意知道隐修院生活的终向是什么，以及隐修士生活的终向是什么。因为无人能怀疑：一个人若被长期的经验与实践所教导，两者都追随过，就能以更大的充实和忠实来论述这些主题，从而通过真实的教导向我们展示它们的价值与目的。\par

% structure: p0026 source=h2 target=subsection reason=relative-heading-rank; base=h2

\subsection{第八章}

\Emph{对所提问题的答复。}\par

\Person{若望}：我绝对应该坚持，同一个人不可能在两种生活中达到成全，除非我被少数几个人的榜样所阻碍。既然在两者之一中找到一位成全的人都不是小事，显而易见，要在两方面都彻底卓越要困难得多——我几乎要说完是不可能的。如果这样的事曾经发生过，也不能归入任何一般规则。因为一般规则必须基于例外实例，即基于极少数人的经验，而是基于多数人——或许所有人——的能力所能及。但那种只在一两个人身上偶尔达到的，且超越普通善性能力的，必须被排除在一般规则之外，作为在人类软弱的状况和本性之外所允许的事，应按作奇迹而非榜样来提出。因此，我将尽自己浅薄的能力，简要地透露您想知道的事。团体隐修士的目标，确实是死化和钉死自己的一切欲望，并根据福音成全的那条有益的命令，不为明日挂虑。显然，这种成全除了团体隐修士外，无人能够达到，正是先知依撒意亚所描述、祝福并赞美的那类人，他说：\Quote{若你在安息日收回你的脚，在我的圣日不做你自己的意愿，并光荣祂，当你不行你自己的道路，你自己的意愿也不发声言语，那时你将在主内喜乐，我必将你提升于地的高处之上，并以你祖先雅各伯的产业养育你，因为上主的口说了这话。} \BibleRef{依 58:13} 而独修隐修士的成全，是使他的心灵摆脱一切尘世事物，并按人性软弱所允许的程度，将心灵与基督结合；先知耶肋米亚描述这样的人说：\Quote{福哉，从小就背负轭的人。他应独坐无语，因为他已担起这轭；} 圣咏作者也说：\Quote{我成了旷野中的鹈鹕。我警醒，成了屋顶上孤独的麻雀。} 因此，对于我们所描述的每一种生活的目标，除非每个人达成了它，那么前者（团体隐修士）徒然接受了团体隐修制度，后者（独修隐修士）徒然接受了独修制度；因为他们二人都不会从其生活方式中获得益处。\par

% structure: p0029 source=h2 target=subsection reason=relative-heading-rank; base=h2

\subsection{第九章}

\Emph{论真实且圆满的成全。}\par

但这是\OriginalTerm{部分的}{\GreekText{μερική}}，即是说，不是彻底又全面的成全，而只是部分的成全。成全确实非常罕见，并靠着天主的恩赐仅赐予极少数人。因为真正而非部分成全的人，能同样不动心地忍受旷野中的荒凉，以及团体中弟兄们的软弱。因此，很难找到一位在两种生活中都成全的人，因为独修隐修士不能彻底获得\OriginalTerm{舍弃物质}{\GreekText{ἀκτημοσύνη}}，即轻视并剥除物质之物，而团体隐修士也不能在静观中获得纯洁，尽管我们知道院长\Person{梅瑟}、\Person{巴弗纽提乌斯}和两位\Person{玛加略}都精通两方面。所以他们在每一种生活中都是成全的；他们比所有旷野居民退得更远，无间地在旷野退隐中喜乐，并且尽己所能，从不寻求与别人交往，然而他们却能忍耐来到他们身边的人及其软弱，以致当大批弟兄为看见他们并获得益处而来时，他们以不动心的忍耐忍受了这几乎是连续不断的接待烦恼。人们以为他们一生从未学习或实践过任何事情，除了向来者表示通常的礼遇。因此，所有人都困惑，不知他们的热忱主要显示在何种生活中，即他们的伟大更能适应独修院的纯洁还是共同生活。\par

% structure: p0032 source=h2 target=subsection reason=relative-heading-rank; base=h2

\subsection{第十章}

\Emph{论那些尚未成全便退入旷野者。}\par

但有些人有时如此被终日持续的旷野寂静所吸引，以致完全畏惧与人交往；当他们因某些弟兄来访的偶然事件而稍稍打破退隐习惯时，就内心显著地焦躁沸腾，并显露出恼怒的明显迹象。这尤其发生在那些没有成熟心志、没有在团体中受过彻底训练便投身独修生活的人身上；这些人总是未成全、易动摇，随着烦恼的风暴所驱，偏向一侧或另一侧。因为他们在与弟兄交往或谈话时不耐烦地沸腾，同样，当他们生活在独居中时，又无法忍受自己所寻求的那片广大的寂静，因为他们自己甚至不知道为何要渴望和寻求独处，却想象这生活的价值和主要部分在于此：即避免与弟兄交往，单纯地避开并厌恶看见人。\par

% structure: p0035 source=h2 target=subsection reason=relative-heading-rank; base=h2

\subsection{第十一章}

\Emph{一个问题：如何医治那些草率离开团体隐修团体的人。}\par

\Person{格玛诺}：经由何种治疗，才能给予我们或其他如此软弱且仅止于此的人帮助？我们尚未摆脱自己的过失，就开始渴望独居，但在团体隐修制度中只接受了很少的指导。或藉由什么方式，我们才能获得不动心的恒常心智和不可摇动的忍耐坚定？我们过早地放弃了团体的共同生活，并离开了这些操练的学校和训练场——我们的原则本应先在这些场所受到彻底训练和成全的。那么，我们现在独居时，如何能获得长期受苦和忍耐的成全？或者，良心——那探查内在动机的——如何能发现这些美德在我们内是存在还是缺失？因我们与人断绝交往，不受他们的任何挑衅而激怒，我们可能被虚假观念欺骗，幻想自己已经获得了那种不动心的平安。\par

% structure: p0038 source=h2 target=subsection reason=relative-heading-rank; base=h2

\subsection{第十二章}

\Emph{解答如何能发现自身过犯的独修士。}\par

若望：对于真正寻求救治的人，灵魂的真正医生绝不会缺少医治的良方；尤其是那些不忽视自身病况（或因绝望而不理会，或因漠不关心而不在意），也不向伤口隐藏危险，更不因放荡的心拒绝补赎之药，反而以谦卑谨慎的心，为因无知、错误或必要而染上的疾病投奔天上的医生的人，更会获得。因此，我们应当知道，如果我们退隐到独修或隐秘之处，但过犯尚未先被治愈，那么它们的活动只是被抑制，而感受它们的力量并未被熄灭。因为万罪之根未被铲除，仍潜伏在我们内，甚至渐渐蔓延，我们仍能通过这些迹象知道它存活。例如，若我们在独修中接待某位弟兄的来访，或他们稍有停留，我们心中便感到焦虑或烦躁，就应认识到极其急躁的不耐烦的诱因仍存在于我们内。但若我们期待一位弟兄的来临，而他却因某种原因稍有迟延，我们心中便暗责他的缓慢，且对这不便的等待感到恼火，扰乱我们的心神，那么良心的省察将表明愤怒与烦恼之罪仍明显存留于我们内。再者，若一位弟兄请求借阅我们的书或其他物品使用，他的请求令我们烦恼，或我们的拒绝使他反感，那无疑我们仍被悭吝或贪心的网罗缠住。但若一个突然的念头或一段圣经经文勾起对一位妇女的记忆，我们感到对她有所眷恋，就应知道淫乱之火尚未在我们内熄灭。但若以我们自己的严格与他人的松懈相比，哪怕一丝的自负试探我们的心神，就显明我们患了可怕骄傲的瘟疫。因此，当我们在心中察觉这些过犯的迹象时，就应清楚认识到，我们只是被剥夺了犯罪的机会，而非犯罪的激情。而这些激情，若我们任何时候介入普通人的生活，就会立刻从思想中的潜伏处跳出来，证明它们并非在爆发时才初次产生，而是因为长久潜伏才最终暴露。所以，即使是一位独修士，也能通过确切的迹象觉察到各过犯的根仍植于他内，只要他不试图在人前显示自己的纯洁，而是在祂——一切心内秘密都无法隐藏的那位——眼中保持纯洁而无玷污。\par

% structure: p0041 source=h2 target=subsection reason=relative-heading-rank; base=h2

\subsection{第十三章}

\Emph{问题：一个进入独修而未根除过犯的人如何可被治愈？}\par

格尔玛努斯：我们非常清楚明白地看见了推断病症迹象的证明和辨别疾病的方法，即如何发现隐藏在我们内的过犯：因为我们每日的经验和思想的日常动态向我们展示了所有这些，正如所述。那么，既然我们的病症之证明与原因已以最清晰的方式向我们揭示，它们的疗法和医治也当被显示。因为无人能怀疑，一个首先发现了疾病之根源与开端，并得到患者良心见证的人，最能讲述其疗法。因此，虽然您的圣德的教导已经揭露了我们伤口的秘密，使我们敢于对疗法抱有些许希望，因为如此清晰的疾病诊断预示了治愈的希望，然而，正如您所说，救恩的初步要素是在会院中获得的，除非首先通过会院的药物得到治愈，否则人无法在独修中保持健康，我们便又陷入绝望的危险状态，唯恐由于我们离开会院时尚不完善，如今在旷野也无法完成成全。\par

% structure: p0044 source=h2 target=subsection reason=relative-heading-rank; base=h2

\subsection{第十四章}

\Emph{关于其疗法的解答。}\par

\Person{若望}：对于那些渴望治愈自己疾病的人，有效的救治肯定是不匮乏的，因此，应该通过发现每种过错的迹象的同一方法来寻求疗法。因为正如我们所说，普通生活中的过错并不缺少独修者，所以我们也不否认，一切对德行的热忱和一切治疗的方法，都供那些脱离普通人生活的人使用。因此，当一个人通过我们上面描述的那些迹象，发现自己受到不耐烦或愤怒的爆发攻击时，他应当经常练习相反和对立的事物，并通过把各种伤害和冤屈（好像是由别人加给他的）放在自己面前，使自己的心灵习惯于以完全的谦卑承受邪恶所能带来的一切；并通过常常向自己呈现各种粗鲁和难以忍受的事物，以心中的全部悲伤不断思考他应以何种温柔来面对它们。这样，通过注视所有圣人的苦难，甚或主自身的苦难，他会承认各种责难和惩罚都是他应得的，并准备好承受各种痛苦。偶尔当他因某个邀请被召回弟兄们的集会——即使是旷野中最严格的隐修士也不可避免地会偶尔发生——如果他发现自己的心灵即使为小事也被悄悄搅乱，他就应当像一位对其隐秘情感严厉的审查者，要求自己完全忍受那些他通过每日默想训练自己完美承受的各种艰难的冤屈，并自责如下，说：我的好人，你是那个在\OriginalTerm{独修}{solitude}的操练场上训练时，最坚决地认为自己能战胜所有坏品性的人吗？你刚才，当你在心中呈现不仅各种辛辣的责难，还有无法忍受的惩罚时，幻想自己相当坚强，能抵挡一切风暴吗？你那不可战胜的忍耐如何被第一声轻言的试探就推翻了？你那座你幻想建在坚固磐石上的房屋，如何被一阵微风就摇动了？你那些在和平时期凭愚昧的信心渴望战争时宣告的话在哪里？\Quote{我已准备好，并不惊惶；}以及你常与先知一起说的：\Quote{主啊，求你考验我，试探我：求你把我的脏腑与我的心都搜索出来；}以及：\Quote{主啊，求你考验我，认识我的心：问我，认识我的路径；并看看我内里有没有邪恶的道路。}一个微小的敌人幽灵把你的巨大战争准备吓成什么样了？人应当用这样的责备和忏悔来谴责自己，不允许那使他动摇的突然试探不受惩罚，而是通过以更严厉的禁食和守夜惩罚肉体，并通过以不断克苦的痛苦惩罚自己轻浮的过错，同时在独修生活中，在这操练的火中消耗那些他本应在\OriginalTerm{修院}{Cœnobium}生活中彻底驱除的东西。无论如何，为了获得持久且不断的忍耐，我们必须坚定地持守这一点：即对于我们——在神圣法律中不仅禁止报复，连想到伤害也被禁止——因某种损失或烦恼而发怒是不允许的。因为对灵魂还有比这更大的伤害吗？即由于愤怒导致的突然盲目，失去真实永恒之光的明亮，并且无法看见那\BibleQuote{良善心谦的}？\BibleRef{玛 11:29}请问，还有比这更危险或更尴尬的事吗？一个人会失去判断善良的能力，以及真正明辨的标准和规则，并且一个清醒的人会做出连醉酒的人和愚人都不会被原谅的事？因此，一个人若仔细考虑这些以及其他类似的伤害，就会很容易忍受并忽视不仅各种损失，还有最残忍的人所能施加的任何冤屈和惩罚，因为他会认为没有比愤怒更具破坏性的，也没有比心灵的平安和持续的纯洁更宝贵的，为了这些，我们应当不仅认为肉体之事的利益无关紧要，而且那些看似精神的事物，如果不能在保持这种宁静的情况下获得或实行，也应置之度外。\par

% structure: p0047 source=h2 target=subsection reason=relative-heading-rank; base=h2

\subsection{第十五章}

\Emph{是否贞洁应像其他情感一样被确定？}\par

\Person{葛尔曼}：正如其他疾病的治疗，即愤怒、烦恼和急躁，已被证明在于以相反之物对抗它们，同样，我们也希望学习针对不洁之灵该使用何种治疗方法：我的意思是，是否情欲之火可以通过像在其他情况下那样，呈现更大的诱惑和刺激物来熄灭？因为我们相信，不仅增加我们内里情欲的诱惑，甚至以心灵的一瞥触碰它们，对贞洁而言是完全致命的。\par

% structure: p0050 source=h2 target=subsection reason=relative-heading-rank; base=h2

\subsection{第十六章}

\Emph{给出识别它的凭据的答案。}\par

\Person{若望}：您敏锐的问题已经预见了这个主题，即使您什么也没说，它也必然从我们的谈话中出现，因此我不怀疑它会有效地被你们的心灵把握，因为你们敏锐的才智已经超出了我们的指导。任何问题的困惑都很容易消除，当询问先于答案，并且首先踏上它将要跟随的道路。因此，对于上面我们谈到的那些过犯，与他人的交往不仅不是障碍，反而是一种相当有益的帮助，因为他们的不耐烦暴露得越频繁，就越给失败者带来更深的悔恨和痛悔，也越给那些与之奋斗的人带来更快的康复。因此，即使当我们生活在独居中，虽然刺激和诱因不能来自人，但我们仍应当有意地默想这些刺激，以便我们在思想中持续战斗，找到更快的治疗。然而，对付邪淫之神的方法是不同的，方式也改变了。因为正如我们必须剥夺身体的欲望机会和肉体接触，同样我们也必须剥夺思想对它的记忆。因为对于仍然软弱和脆弱的心灵，容忍对这种情欲最轻微的回忆也是足够危险的，以至于有时在回忆圣善妇女，或读圣经中的故事时，会激起危险的兴奋。正因如此，我们的长老们过去常有意地省略这类段落，当有年轻人时。然而，对于那些完美且建立了贞洁感受的人，不缺乏自我检验的证据，他们可以通过自己良心无腐败的判断来确立内心的正直。因此，对于彻底坚定的人，甚至对于这情欲也有类似的检验，以至于一个确信自己已完全铲除此恶根的人，为了确证自己的贞洁，可以有意地唤起某种图片，如同带着放荡的思想。但对于仍软弱的人，这样的检验绝不应当尝试（因为对他们而言，这将是危险而非有益的），\OriginalTerm{他们在心中处理女性的结合以及某种温柔且极为柔软的触摸}{ut conjunctionem femineam et palpationem quodammodo teneram atque mollissimam corde pertractent.} 因此，一个以完美德行为基础的人，如果在思想中虚构的最温柔的兴奋刺激中，觉察不到自己有任何心灵的同意、任何肉体的骚动，那么他将有一个非常确凿的纯洁证明，以至于他锻炼自己达到这种坚定的纯洁，不仅在心里拥有贞洁和无染的福祉，而且即使被迫触摸女子的身体，他也会对此感到惊骇。\par

\Person{若望}院长看到正是第九时辰休息的时候，就把他的讲论结束了。\par

\SourceRecord{Translated by C.S. Gibson.；Nicene and Post-Nicene Fathers, Second Series；Vol. 11.；Edited by Philip Schaff and Henry Wace.；Buffalo, NY: Christian Literature Publishing Co.,；1894.；<http://www.newadvent.org/fathers/350819.htm>.}

% structure: p0001 source=h1 target=section reason=page-title

\section{第二十次会谈}

\Emph{\Person{圣若望·卡西安} 著}\par

\Person{院长皮努斐}的会谈：论补赎之目的与赔补之标记。\par

% structure: p0004 source=h2 target=subsection reason=relative-heading-rank; base=h2

\subsection{第一章}

\Emph{论\Person{院长皮努斐}的谦卑及其隐居之所。}\par

如今我将要叙述那位杰出而非凡的人物、\Person{院长皮努斐}关于补赎之目的的训诲，我以为可以节省大量篇幅，若顾虑使读者厌倦，而在此缄默不谈他的谦卑之颂扬——我在\WorkTitle{隐修院规则}第四卷中曾以简短言论略述此事，该书标题为\Quote{论弃绝者当守之规则}——尤其因为许多未读过那部著作的人，或许会读到本书，那时若对发言者的德行毫无记叙，则其言论的权威势必削弱。\newline{} 此人曾任院长与司铎，主持一座距埃及城市\Place{帕内菲西斯}不远的大隐修院，当全境因其德行与奇迹而赞颂他至天际，以致他似乎已因人的赞誉之酬报而领取了他劳苦的果实，但他唯恐他所特别厌恶的民众宠爱的虚浮会妨碍永赏的果实，便秘密逃离隐修院，前往\Place{塔本纳}隐修士们最偏远的角落。在那里他选择的不是旷野的孤寂，也不是独居生活所提供的那种无忧无虑——甚至那些不完美的、不能忍受隐修院中服从所要求的劳苦的人，有时也以骄傲的僭妄追求这种生活——而是选择投身于一座极负盛名的隐修院。然而，为免因着装标记而暴露，他换上世俗服装，在门口流泪跪伏多日——按当地惯例——并抱着所有人的膝头，每日被那些为考验其决心而声称他已至暮年、寻求这圣洁生活并非真诚而是因缺乏食物所驱的人斥退，终于得以入院。在那里他被派去协助一位负责园圃的年轻弟兄。他不仅以如此惊人而圣洁的谦卑完成了长上命令他的一切，以及所托付的工作所需之事，还在夜间以隐秘劳作完成某些其余人因厌恶而回避的必要事务，以致天亮时，全体会众为这些有益的工程而欣喜，却不知其作者。他在那里将近三年，以自己所渴望却备受不公对待的劳苦为乐。恰巧，有位认识他的弟兄从同属埃及的地区来到那里。此人一度犹豫，因他衣着简陋、职务卑微而不能立即认出他，但仔细端详后，俯伏在他脚前，首先使众弟兄惊讶，随后当他泄露其名时——这名字因他特殊圣德的名声也为他们所知——他们便因曾指派一位如此德行之人、一位司铎从事如此卑贱的职务而悲痛与痛悔。但他泪流不止，将这泄露的意外归咎于魔鬼的严重嫉妒，由众弟兄簇拥着，在庄严的护送中被带回隐修院。他在那里停留不久，再次因尊位与品级所受的敬重而困扰，便秘密登船前往\Place{叙利亚}的\Place{巴勒斯坦}省，在那里被当作初学新手接入我们当时所住的隐修院，并由院长安排住在我们的斗室中。然而，即使在彼处，他的德行与功绩也不能长久隐藏。因为他以同样方式被发现并泄露，最终以最大的荣誉与尊敬被带回自己的隐修院。\par

% structure: p0007 source=h2 target=subsection reason=relative-heading-rank; base=h2

\subsection{第二章}

\Emph{论我们来到他那里。}\par

不久之后，渴求神圣训诲的愿望也促使我们前往埃及拜访他。我们以极大的热忱与虔诚寻访他，他如此和蔼谦恭地接待我们，以至于他竟尊崇我们为昔日同室者，让我们住在他建在园子最角落的斗室里。在那里，当全体弟兄在场时，他曾向一位遵守隐修院规矩的弟兄传授了足够艰难而崇高的训导——正如我所说，我在\WorkTitle{隐修院规则}第四卷中已尽量简要地概括了这些训导——那时，真正的弃绝之高峰在我们看来如此难以企及、如此奇妙，以至于我们这些卑微之人竟以为永远无法攀登。因此，我们陷入绝望，并未在面容上掩饰内心的痛苦思绪，怀着颇为焦虑的心回到那位有福的长老面前。他立刻询问我们为何如此忧愁，\Person{院长日尔曼诺}深深叹息，回答如下。\par

% structure: p0010 source=h2 target=subsection reason=relative-heading-rank; base=h2

\subsection{第三章}

\Emph{论补赎之目的与赔补之标记的问题。}\par

您伟大而辉煌的新教义阐述，为我们开辟了一条通往最光荣弃绝的更艰难道路，并移开了我们眼前的鳞片，向我们展示了其高耸于天际的顶峰；因此我们相应地被更沉重的绝望所压倒。因为当我们将其广阔与我们微薄的力量相较量，并将我们无知之极端卑微与所展示的无限德行之高度相比较时，我们感到如此渺小，不仅无法达至，而且肯定在已有的方面也必不足。因为我们在过分绝望的重压下，竟从最深处跌入更深处。因此，唯一能治愈我们创伤的支柱就是：让我们学习关于补赎之目的，尤其是关于赔补之标记的知识，以便我们确知已往罪过的赦免，从而受到激励，攀登上述完美之高峰。\par

% structure: p0013 source=h2 target=subsection reason=relative-heading-rank; base=h2

\subsection{第四章}

\Emph{我们请求所表现的谦卑的回答。}\par

\Person{皮努斐乌斯}说：我确实非常欣喜于你们谦卑的丰硕果实，这果实我曾以并非漠不关心的态度看到，那时我先前被接纳居住在你们那小室的住所；我很高兴你们如此恭敬地接受我们——所有基督徒中最小的——所给的训诫，以及我冒昧说出的话，以至于如果我没有弄错的话，你们一听到我们说出就立刻实行；并且，虽然我记得，这些话的重要性几乎不值得你们所付出的努力，但你们如此隐藏你们德行的功绩，仿佛从未有任何关于你们日常实践之事的消息传到你们耳中。但因为这一事实值得最高的赞扬，即：你们声称那些圣人会规对你们仍是未知，仿佛你们仍是初学者，我们将尽可能简短地总结你们如此热切向我们请求的事。因为我们必须甚至超越我们的能力和才干，服从像你们这样的老朋友的命令。因此，关于补赎的价值和平息之能，许多人已发表了许多，不仅口头上而且文字上，表明它是多么有益、多么坚强、充满恩宠，以致当天主因我们过去的罪过而受冒犯，并即将对这些冒犯施加最公正的惩罚时，它（如果这样说不算错的话）在某种程度上阻止了他，并且，如果可以这样说的话，甚至违背他的意愿而制止了报复者的右手。但我毫不怀疑，这一切你们都很熟悉，或是出于你们天生的智慧，或是出于你们对圣经的不倦研读，以至于从这些——可以说——初生的嫩芽，你们的皈依便萌发了。最后，你们所关切的不是补赎的性质，而是它的终结和赔补的标记，因此通过一个非常巧妙的问题，你们问到了别人所遗漏的事。\par

% structure: p0016 source=h2 target=subsection reason=relative-heading-rank; base=h2

\subsection{第五章}

\Emph{论补赎的方法与赦免的证明。}\par

因此，为了尽可能简短地满足你们的渴望和问题，补赎的圆满而完整的描述是：决不再向那些我们为之做补赎、或良心为之刺痛的同一些罪过屈服。但赔补与赦免的证明在于：我们已经从心中驱逐了对它们的爱。因为只要任何所犯之罪或类似之罪的影像还在眼前跳动，且我不说喜悦，而是对它们的回忆在灵魂深处萦绕，同时他正为这些罪过致力于赔补和流泪，那么任何人都不能确信自己已从先前的罪过中得获自由。因此，一个正警惕地做赔补的人，当他从未感到自己的心被同一些罪过的诱惑和幻想所激动时，他那时才可以确信自己已从罪过中得获自由，并已获得对过去过失的赦免。因此，补赎最真实的验证和赦免的见证在于我们自己的良心，这良心甚至在审判和认知的日子之前，当我们仍在肉身内时，就显明了我们脱离罪责的裁定，并揭示了赔补的终结和赦免的恩宠。为使所说的话更意味深长，我们只应在以下情况下相信过去罪过的污点已被赦免：当对当前享乐的欲望以及情欲都已被从心中驱逐出去之时。\par

% structure: p0019 source=h2 target=subsection reason=relative-heading-rank; base=h2

\subsection{第六章}

\Emph{一个问题：我们是否应当出于内心的痛悔而记住自己的罪过？}\par

\Person{日耳曼努斯}说：那么，如何才能在心中激起这种神圣而有益的、出于谦卑的痛悔呢？这种痛悔在忏悔者的口中是如此描述的：\BibleQuote{我承认了我的罪，没有隐藏我的不义。我说：我要向主承认我的不义}，这样我们才能有效地说出接下来的话：\BibleQuote{你就赦免了我心中的罪孽}；或者，当我们跪下祈祷时，如何才能激起自己产生告解的眼泪，借以获得对我们过犯的赦免——正如这些话所说的：\BibleQuote{每夜我要洗涤我的床铺，用眼泪浸透我的卧榻}——如果我们从心中驱逐了对所有过失的回忆？虽然相反，我们却被命要小心地保存对它们的记忆，正如主所说：\BibleQuote{你的罪过我不再记念，但你，你要记念它们吗？} \BibleRef{依 43:25} 因此，不仅在我工作时，而且在我祈祷时，我特意努力召回对罪的记忆，以便更有效地导向真正的谦卑和心灵的痛悔，并敢于与先知一同说：\BibleQuote{求你看顾我的谦卑和劳苦，赦免我所有的罪。}\par

% structure: p0022 source=h2 target=subsection reason=relative-heading-rank; base=h2

\subsection{第七章}

\Emph{回答：我们应当在多大程度上保存对先前行为的记忆。}\par

\Person{Pinufius}：你的问题，正如上面已经说过的，不是关乎补赎的本质，而是关乎它的目的和赔补的标记。对此，我认为已经给出了合理且相关的答复。但你所说的关于罪的记忆，对于那些仍在做补赎的人而言，是足够有用和必要的，好使他们不断捶胸说：\Quote{因为我承认我的邪恶，我的罪常在我面前；}以及：\Quote{我要思念我的罪过。}因此，当我们做补赎时，仍因回忆过失行为而悲伤，那因承认我们的过错而引发的泪水之雨，必定会熄灭我们良心的火焰。然而，当一个人仍处于这种心灵谦卑和心神痛悔的状态，并继续劳苦哭泣时，这些事的记忆渐渐消退，良心的荆棘因天主的恩宠从他内心深处被拔出，那么显然他已达到了赔补的终点和赦免的赏报，并从所犯之罪的污秽中被洁净。要达到这种遗忘的状态，只有通过抹去我们过去的罪过和偏好，以及达到完全彻底的心灵纯洁才能实现。而这绝不会为那些因懈怠或疏忽而未能清除自己过犯的人所获得，只唯独为那因不断悲叹和忧伤之叹息而除去旧日污点的人，并因他良善的心和劳苦向主宣告：\Quote{我承认了我的罪，没有隐藏我的不义；}以及：\Quote{我的眼泪昼夜成为我的饮食；}以致最终他堪当听到这些话：\Quote{让你的声音停止哭泣，你的眼睛止住眼泪：因你的劳苦有报酬，主说；}并且主的声音也论他说：\Quote{我如云彩抹去了你的过犯，如薄雾消除了你的罪恶；}又说：\Quote{我，惟有我为自己的缘故涂抹了你的过犯，你的罪孽我不再记念；}这样，当他由\Quote{罪恶的绳索}——\Quote{人人都被捆绑}——中获得释放时\BibleRef{箴 5:22}，他将以全备的感恩向主歌唱：\Quote{你折断了我的锁链，我要向你献上赞美的祭献。}\par

% structure: p0025 source=h2 target=subsection reason=relative-heading-rank; base=h2

\subsection{第八章}

\Emph{论补赎的各种果实。}\par

因为在普遍的圣洗恩宠之后，以及那通过血洗而获得的极其珍贵的殉道恩赐之后，还有许多补赎的果实，我们可以藉此赎罪。因为永生的应许不仅赐予补赎本身，正如真福宗徒伯多禄所说：\Quote{你们悔改并归正，好使你们的罪被赦免。}；洗者若翰和主自己也说：\Quote{悔改吧，因为天国临近了。}而且，爱德之情也能压倒我们罪过的重量：因为\BibleQuote{爱德遮盖许多罪过。}\BibleRef{伯前 4:8}同样，施舍的果实也为我们的创伤提供了良药，因为\BibleQuote{水能熄灭烈火，施舍能赎罪过。}\BibleRef{德 3:33}流泪也能洗去罪愆，因为\Quote{每夜我要洗涤我的床榻，以眼泪浸湿我的铺盖。}最后，为表明它们并非徒然，他接着说道：\Quote{你们所有作恶的人，离开我吧！因为上主听了我哭泣的声音。}此外，藉着告解，罪过得赦免：因为\Quote{我说：我要向上主承认我的过犯；你就赦免了我心中的罪孽。}又说道：\Quote{首先你应声明你的不义，好使你得到成义。}同样，克苦内心和身体也能获赦所犯之罪，因为他说：\Quote{求你垂顾我的卑微和劳苦，赦免我的一切罪恶。}尤其藉着生活的改善：他说：\Quote{从我的眼前除去你们思想的邪恶。停止作恶，学习行善。寻求正义，解救受压迫者：为孤儿伸冤，保护寡妇。然后来，让我们彼此辩论——上主说：你们的罪虽似朱红，也将变得洁白如雪；虽红得像丹砂，也将白得像羊毛。}有时，圣人的代祷也能获得罪过的赦免，因为\Quote{谁若看见自己的弟兄犯了不至于死的罪，就应当祈求，并赐给他生命，给那些犯罪不至于死的人。}又说道：\Quote{你们中间有患病的吗？他该请教会的长老来，他们该为他祈祷，因主的名给他傅油。信德的祈祷必能拯救病者，主必使他起来；如果他犯了罪，也必得赦免。}有时，慈悲和信德之德也能除去罪的污点，正如这段经文：\BibleQuote{以慈悲和信德，罪过得以洗净。}\BibleRef{箴 15:27}常常，藉着因我们的劝诫和宣讲而得救之人的悔改和得救，也是如此：\BibleQuote{那引导罪人从迷途回头的人，必救自己的灵魂免于死亡，并遮盖许多罪过。}\BibleRef{雅 5:20}此外，我们若宽免和原谅别人，也能获得罪过的赦免：\BibleQuote{因为你们若宽免别人的过犯，你们的天父也必宽免你们的过犯。}\BibleRef{玛 6:14}你看，救主的慈爱为我们开启了何等伟大的获救途径，使任何渴望得救的人都不致因绝望而丧气，因为他看见自己被如此众多的良方呼唤归向生命。因为如果你申辩由于肉身的软弱不能藉斋戒除去罪过，你不能说：\Quote{我的双膝因斋戒而无力，我的肉体因油而改变；因为我以灰烬当饭，以眼泪和饮料同饮。}那么，你就用大量的施舍来补偿吧。如果你没有什么可以给穷人的（虽然贫穷和匮乏并不排除任何人行此善举，因为寡妇的两个小钱比富人的丰厚赠予更被看重，并且主应许连一杯凉水也将得到赏报），你至少可以藉着生活的改善来洗除它们。如果你不能藉根除所有缺点而达至完善的善，你可以为别人的善和得救而展现虔诚的关怀。如果你抱怨不能胜任此职，你可以藉着爱德之情来遮盖你的罪恶。如果在这方面，心灵的怠惰使你软弱，你至少应该谦卑地怀着谦逊之情，藉着圣人的祈祷和代祷，祈求医治你创伤的良方。最后，有谁不能谦卑地说：\Quote{我承认了我的罪，没有隐瞒我的不义。}，以致他能藉此认罪再加上这句：\Quote{你就赦免了我心中的罪孽。}但如果羞耻阻止了你，你耻于在人类面前揭露它们，你应当不断地向那位无法隐瞒的主恳切认罪，并对他说：\Quote{我承认我的过犯，我的罪恶常在我眼前。我唯独得罪了你，在你面前行了恶。}因为祂常愿医治它们而不带任何使人羞耻的公布，赦免罪过而不加任何谴责。此外，除了那现成而可靠的帮助，天主的俯就还给了我们另一个更容易的帮助，并将这良药托付给我们自己的意志，使我们能从自身的情感中推断出罪过的赦免，当我们向祂说：\BibleQuote{求你宽免我们的罪债，如同我们也宽免得罪我们的人。}\BibleRef{玛 6:12}因此，谁若想获得罪赦，就应当努力通过这些方式来预备自己。不要让顽石般的心硬使任何人背离这拯救的良方和如此良善的源头，因为即使我们做了这一切，它们也无法赎我们的罪，除非被上主的善良和仁慈所涂抹。当祂看到我们以谦卑之心所奉献的虔诚努力时，就以极大的美善支持我们微小无力的努力，并说：\BibleQuote{是我，是我为了自己涂抹了你的过犯，不再记念你的罪恶。}\BibleRef{依 43:25}因此，凡追求我们所提及的这种境界的人，必将通过每日的斋戒和心灵与身体的克苦来寻求补赎的恩宠，因为正如经上所记：\BibleQuote{若不流血，就不得赦免。}\BibleRef{希 9:22}而这并非没有理由。因为\BibleQuote{血肉之人不能承受天主的国。}\BibleRef{格前 15:50}因此，谁若将\BibleQuote{圣神的剑，即天主的话}\BibleRef{弗 6:17}从这流血中收回，必然落在耶肋米亚的诅咒之下；因为他说：\BibleQuote{那制止刀剑不沾血的人，是可诅咒的。}\BibleRef{耶 48:10}因为这把剑为了我们的益处而让那恶血流出——我们罪恶的材料藉这恶血而存活；并切除和削去在灵魂肢体中所发现的一切肉性和属地之物；使人在罪上死，在天主内活，并以属神的德行而兴旺。于是，他将不再为回忆从前的罪而哭泣，而是为来世的希望而哭泣，对已往的罪恶想得较少，而对将来的好事想得较多，流下眼泪，不是因罪过悲伤，而是因永生的喜悦，并且\Quote{忘记背后}\BibleRef{斐 3:13}，即肉体的罪，努力向着\BibleQuote{前面的}\BibleRef{斐 3:13}，即属神的恩赐和德行。\par

% structure: p0028 source=h2 target=subsection reason=relative-heading-rank; base=h2

\subsection{第九章}

\Emph{对成全者而言，忘记罪过何其宝贵。}\par

但关于你刚才所说的一点，即你刻意回想过去的罪过，这必定不应该做；不，如果它强行侵袭你，必须立即驱除。因为它极大地阻碍灵魂默观纯洁，尤其对于独居者，因为它将他缠缚于世俗的污秽，淹没在可憎的罪过中。因为当你回忆那些因无知或放纵、顺从今世之王所行的事时，我承认当你沉浸在这些思绪中时，没有快感潜入，但至少旧日污秽的沾染必定以其臭气败坏你的灵魂，并阻挡善德的神圣芬芳，即甘饴的馨香。因此，当过去罪过的回忆浮现于你的脑海时，你必须像正人君子那样避开它，正如一位正直的人在公共场合被一个无耻放荡的女子以言语或拥抱纠缠时，他必逃脱。而且，除非他立即脱离与她的接触，如果他允许自己稍留在不洁的言谈中，即使他拒绝同意可耻的快乐，也无法避免在所有路人评判中背负恶名和鄙视的烙印。同样，我们若被有害的回忆引向此类思念，理应立刻停止沉浸其中，并履行撒罗满所吩咐我们的：\BibleQuote{但你要出去，}他说，\BibleQuote{不要在她的地方停留，也不要定睛看她；}\BibleRef{箴 9:18}免得天使看见我们沉溺于不洁和污秽的思念，在路过时不能对我们说：\Quote{愿主的祝福临于你。}因为当内心主要部分被污秽和属地的思虑占据时，灵魂不可能持续于善念。因为撒罗满的这句话是真实的：\BibleQuote{你的眼睛注视淫妇时，你的口就会说出邪恶；你必像躺在大海之中，又像躺在桅杆上，在风暴中。你要说：他们打了我，我不觉痛；他们嘲弄我，我不知觉。}\BibleRef{箴 23:33}因此我们不但要离弃所有污秽，甚至一切属地的思念，并常将我们灵魂的渴望提升至天上的事物，依从我们救主的话：\BibleQuote{我在哪里，}祂说，\BibleQuote{我的仆人也要在那里。}\BibleRef{若 12:26}因为常常发生这样的事：当有人出于怜悯在思想中回顾自己的跌倒或其他犯错之人的跌倒时，他受到快感的影响而同意这最巧妙的攻击，那以貌似良善开始和发动的事，终以污秽和有害的结局告终，因为\BibleQuote{有一条路，在人看来是正直的，但其终局却通向地狱的深处。}\BibleRef{箴 16:25}\par

% structure: p0031 source=h2 target=subsection reason=relative-heading-rank; base=h2

\subsection{第十章}

\Emph{如何避免回忆我们自己的罪过。}\par

因此，我们必须努力使自己振作，以达到这值得称赞的痛悔，藉着追求德行和对天国神圣的渴望，而非藉着危险的罪过回忆；因为一个人只要选择站在粪坑上或搅动其污秽，就必然被其瘟疫般的气味窒息而死。\par

% structure: p0034 source=h2 target=subsection reason=relative-heading-rank; base=h2

\subsection{第十一章}

\Emph{论补赎的标记与过去罪过的消除。}\par

然而我们知道，正如我们常常所说，只有当引致我们负罪、应当为之痛悔的那些动机和情感从心中根除时，我们才算为过去的罪过作了补赎。但无人应妄想他能获得这一点，除非他首先以全部心灵的热忱切断导致他堕入这些罪过的机会和场合；例如，如果由于与女子危险地交往而陷入奸淫或通奸，他必须竭尽全力避免甚至看她一眼；或者如果他因酗酒暴食而被制伏，他应当以最严厉的方式克制对无节制饮食的贪欲。再者，如果他被贪图财物的欲望和喜爱引入歧途，陷入伪证、偷窃、杀人或亵渎，他应当切断诱骗他的贪婪之机会。如果他被骄傲的激情驱使而陷入愤怒之罪，他应当以谦卑的全部德行除去傲慢的诱因。如此，为使每一个罪过被毁灭，必须先除去其所以或为之而犯的机会和场合。因为通过这种治疗之法，我们定能达到对所犯罪过的遗忘。\par

% structure: p0037 source=h2 target=subsection reason=relative-heading-rank; base=h2

\subsection{第十二章}

\Emph{哪些罪过我们只需作暂时补赎；哪些可以永无止境。}\par

但是，所说关于遗忘的这种描述，只涉及大罪，这些大罪也受梅瑟法律的谴责；对它们的倾向被良好的生活摧毁并终结，因此对它们的补赎也有终结。但为那些小罪，如经上所载：\BibleQuote{义人七次跌倒，仍必再起}\BibleRef{箴 24:16}，补赎永不止息。因为或因无知，或因遗忘，或因思想，或因言语，或因意外，或因需要，或因肉身的软弱，或因梦中的玷污，我们每日常常故意或不自愿地跌倒；这些过犯，\Person{达味}也曾向主祈求，恳求洁净和宽恕，说：\Quote{谁能认清自己的过犯？求你赦免我未觉察的罪愆；也使你的仆人免受故意犯的罪；}宗徒也说：\Quote{我所愿意的善，我不去行；我所不愿意的恶，我却去作。}为此，同一位宗徒叹息喊道：\Quote{我真不幸！谁能救我脱离这该死的肉身呢？}因为我们是如此轻易地陷入这些过犯，仿佛出于本性的法则，无论我们如何小心谨慎地提防，总不能完全避免。既然这是耶稣所爱的那位门徒所宣告并绝对定论的，说：\BibleQuote{如果我们说我们没有罪，便是欺骗自己，祂的话就不在我们内。}\BibleRef{若一 1:8}再者，对于一个渴望达到成全高峰的人，达到补赎的终结，也就是约束自己不犯不法之事，并没有多大帮助，除非他始终以不懈的奔跑催促自己迈向那些德行，借此我们才能达到赔补的标志。因为一个人仅仅使自己摆脱主所憎恨的那些罪恶的污秽是不够的，除非他还以心灵的纯洁和成全的宗徒的爱德，获得了主所喜悦的那种德行的甘饴芬芳。以上就是\Person{皮努斐熙}院长关于赔补的标志和补赎的终结的论述。虽然他怀着焦急的爱催促我们决定留在他的隐修院团体里，但当我们因受到\Place{赛格特}旷野名声的激励而无法留下时，他便送我们上路。\par

\SourceRecord{Translated by C.S. Gibson.；Nicene and Post-Nicene Fathers, Second Series；Vol. 11.；Edited by Philip Schaff and Henry Wace.；Buffalo, NY: Christian Literature Publishing Co.,；1894.；<http://www.newadvent.org/fathers/350820.htm>.}

% structure: p0001 source=h1 target=section reason=page-title

\section{第二十一篇会谈}

由圣若望·卡西安著\par

院长特奥纳的第一篇会谈。论五十天的放松。\par

% structure: p0004 source=h2 target=subsection reason=relative-heading-rank; base=h2

\subsection{第一章}

特奥纳如何来到院长若望处。\par

在我们开始陈述与那位卓越人物院长特奥纳的这次会谈的言论之前，我认为最好先简要叙述他归化的起始；因为从此读者将能更清楚地看到此人的卓越与恩宠。他当时还很年轻，在父母的愿望和命令下，结成了婚姻的纽带；因为他们出于虔诚的忧虑，关心他的贞洁，担心在危险的年龄出现关键的跌倒，便以为可通过合法婚姻的补救来预防青春的激情。于是他与妻子生活了五年后，来到了院长若望那里；若望当时因其非凡的圣德被选主持赈济事务。因为并非任何喜欢的人出于自己的意愿或野心被晋升到这一职务，而是只有全体长老聚会认为，从他的年龄优势以及对其信德与德行的见证来看，他比所有其他人都更卓越、更超越的人。这位有福的若望，上述年轻人怀着虔诚的热忱来到他那里，在那些急于向我所提及的那位老人奉献他们财产的十一与初果的其它主人中，他也带来了虔诚的礼物；当老人看到他们带着许多礼物涌向自己，并想为他们的奉献作出某种回报时，他便如宗徒所说，开始向他们播撒属灵的事物，因为他正收获他们属肉体的礼物。\BibleRef{格前 9:11}最后，他这样开始了劝言的话语。\par

% structure: p0007 source=h2 target=subsection reason=relative-heading-rank; base=h2

\subsection{第二章}

院长若望对特奥纳及与他同来的其他人的劝言。\par

我的孩子们，我确实为你们礼物的虔诚慷慨而喜悦；我也欣然接受你们虔诚的奉献，其处置托付给了我，因为你们将初果和十一奉献出来，为穷人的利益和使用，作为馨香的祭献献给主，相信通过这件奉献，你们的果实和所有财富——你们从中取出了这些献给主——将丰盛地蒙福，并且你们自己将按照祂命令的信德，甚至在此世获得各样美物的丰盛：\BibleQuote{你要以正义的劳役尊崇主，将你正义的果实献给他；这样你的仓库必充满丰盈的麦子，你的酒池必溢出新酒。}\BibleRef{箴 3:9}既然你们忠实地履行这项服务，你们可以知道，你们已经实现了旧约的义德；在旧约之下，那些当时生活的人若违犯它，必无可避免地陷入罪债，而若他们履行它，却不能达到完美的顶点。\par

% structure: p0010 source=h2 target=subsection reason=relative-heading-rank; base=h2

\subsection{第三章}

论十一奉献和初果的奉献。\par

因为确实，按主的命令，十一奉献被分别出来，为肋未人的服务所用，而祭品和初果则为司祭所用。但初果的法律是这样：即果物或牲畜的五十分之一应献给圣殿和司祭的服务；有些人不信实、粗心大意地减少了这一比例，而那些非常虔诚的人则增加了它，以至于有人只给六十分之一，有人则给四十分之一的果实。因为义人，法律不是为他们制定的，这样便表明他们不在法律之下，因为他们不仅试图实现，甚至超越法律的义德；他们的虔敬大于法律的要求，因为它超越了诫命的遵守，并凭自己的自由意志增加了所应尽的义务。\par

% structure: p0013 source=h2 target=subsection reason=relative-heading-rank; base=h2

\subsection{第四章}

亚巴郎、达味及其他圣人如何超越了法律的要求。\par

因为这样我们读到，\Person{亚巴郎}超越了那后来将要赐下的法律的要求，当他战胜四位国王之后，他不肯触碰\Place{索多玛}的任何战利品，那些战利品作为征服者是理应归他的，而且确实是他所救出战利品的那位国王本人向他提供的；并用天主之名发誓说：\Quote{我向上天下地的主、至高者天主举手起誓：凡是你的东西，就是一根线、一根鞋带，我都不拿。}\BibleRef{创 14:22}同样，我们知道\Person{达味}超越了法律的要求，因为虽然\Person{梅瑟}命令要报复仇敌\BibleRef{出 21:24}，他不仅没有这样做，反而以爱拥抱迫害他的人，虔诚地为他们祈求上主，当他们被杀时还痛哭流涕并为他们报仇。同样，我们确信\Person{厄里亚}和\Person{耶肋米亚}并不在法律之下，因为他们虽然可以无可指摘地利用合法的婚姻，却宁愿保持童贞。同样，我们读到\Person{厄里叟}和其他过着同样生活方式的人超越了\Person{梅瑟}的诫命，关于他们，宗徒这样说道：\Quote{他们披着绵羊皮、山羊皮，受压迫，受磨难，受贫困，世界配不上他们，他们在荒野、山岭、山洞和地穴中漂流。}\BibleRef{希 11:37}关于\Person{勒加布}的儿子\Person{约纳达布}的子孙，我该说什么呢？关于他们，我们读到，当先知\Person{耶肋米亚}奉主的命令给他们酒喝时，他们回答说：\Quote{我们不喝酒：因为我们的父亲、\Person{勒加布}的儿子\Person{约纳达布}命令我们说：\InnerQuote{你们和你们的子孙永远不可喝酒：你们不可盖房，不可撒种，不可栽植或拥有葡萄园：但你们要终生住在帐幕里。}}因此，他们也得以从同一位先知那里听到这些话：\Quote{万军的上主天主，以色列的天主这样说：\Person{勒加布}的儿子\Person{约纳达布}的后裔，总不会缺少在我面前侍立的人。}因为他们所有人都不满足于仅献出财产的十分之一，而是实际放弃了财产，反而将自己和自己的灵魂献给天主——人无法为自己的灵魂付出任何赎价，正如主在福音中所证实的：\Quote{人还能拿什么来交换自己的灵魂？}\BibleRef{玛 16:25}\par

% structure: p0016 source=h2 target=subsection reason=relative-heading-rank; base=h2

\subsection{第五章}

\Emph{生活在福音恩宠下的人应如何超越法律的要求。}\par

因此，我们应当知道，我们这些人虽然不再受法律的苛求，但福音的话每日在我们耳中回响：\BibleQuote{你若愿意成为成全的，去，变卖你所有的，施舍给穷人，你将有宝藏在天上，然后来跟随我。}\BibleRef{玛 19:21}当我们向上主献上财产的十分之一时，我们仍多少背负着法律的负担，无法上升到福音的高峰；而那些遵行福音的人，不仅在此生蒙福，也将获得未来的赏报。因为法律向遵守它的人所应许的，不是天国的赏报，而只是此生的安慰，它说：\BibleQuote{遵行这些事的人，必因此生活。}\BibleRef{肋 18:5}但主向祂的门徒说：\BibleQuote{神贫的人是有福的，因为天国是他们的。}又说：\BibleQuote{凡为我的名，舍弃了房屋、或兄弟、或姊妹、或父亲、或母亲、或妻子、或儿女、或田地的，必要领取百倍的赏报，并承受永生。}这合情合理。因为我们不应当因禁止之事而受称赞，而应因合法之事而受称赞，并且出于对那因我们的软弱而允许我们使用这些事者的敬畏，而不使用它们。因此，即使那些忠心献上果实十分之一、遵守主更古老诫命的人，尚且不能攀登福音的高峰，你就能清楚地看见，那些连这一点也不做的人，相差何其远。因为那些忽视履行法律中较轻松的命令（而法律赋予者庄重的话语证明了这些命令的容易）的人，怎能分享福音的恩宠呢？法律甚至对那些不履行的人发出了诅咒，因为它说：\BibleQuote{凡不持守这法律书上一切的话去遵行的，是可咒骂的。}\BibleRef{申 27:26}但在这里，由于诫命的卓越和崇高，它说：\BibleQuote{谁能领会，就领会吧！}\BibleRef{玛 19:12}在那里，立法者的强制显示了诫命的容易；因为他（\Person{梅瑟}）说：\BibleQuote{我今天请天地向你们作证，如果你们不遵守上主你们天主的诫命，你们必从这地上消灭。}\BibleRef{申 4:26}在这里，崇高命令的伟大性正体现在祂不是\Emph{命令}，而是\Emph{劝勉}，说：\BibleQuote{你若愿意成全，去}做这个或那个。在那里，\Person{梅瑟}将无法拒绝的重担加在不情愿的人身上；在这里，\Person{保禄}以劝言迎接那些愿意且渴望成全的人。因为那因神奇而崇高、并非所有人都能履行的，不应作为普遍要求，或（若可以这样说）作为常规规则强制于人；但藉着劝言，所有人都被激励走向恩宠，使伟大的人因其德行的成全而该当得冠冕，而那些渺小、无法达到\BibleQuote{达到基督圆满年龄的程度}\BibleRef{弗 4:13}的人，虽然似乎被更大的星辰的光芒所遮蔽和隐藏，却仍可免于法律诅咒的黑暗，不被判定遭受现世的灾祸，也不受永罚的责打。因此，基督并不以命令的强制来约束任何人走向那崇高的良善高处，而是藉着自由意志的力量激发他们，并以智慧的劝言和成全的渴望推动他们。因为哪里有命令，哪里就有义务，因而也有惩罚。但那些因法律所建立的严厉而被驱策遵守这些事的人，仅能避免他们所受到的威胁，却不能获得赏报和酬劳。\par

% structure: p0019 source=h2 target=subsection reason=relative-heading-rank; base=h2

\subsection{第六章}

\Emph{福音的恩宠如何扶持软弱的人，使他们能获得宽恕，正如它保证成全的人获得天国一样。}\par

就如福音的话将强壮的人提升到崇高卓越的高峰，同样，它也不让软弱的人被拖入深渊，因为它为成全的人保证了福分的圆满，并为那些因软弱而跌倒的人带来了宽恕。因为法律将履行其命令的人置于一种介于两者之间的状态，使他们既脱离违命者的定罪，又远离成全者的荣耀。但这状态是何等可怜和悲惨，你可以从今生的状况看出：在现世，一个人若只为不被好人视为有罪而流汗劳作，而不求被视为富有、尊贵和有名望，那实在是一件非常可悲的事。\par

% structure: p0022 source=h2 target=subsection reason=relative-heading-rank; base=h2

\subsection{第七章}

\Emph{如何在于我们自己的选择，是要留在福音的恩宠之下，还是留在法律的恐吓之下。}\par

因此，今天在于我们自己的权力，选择是要生活在福音的恩宠之下，还是生活在法律的恐吓之下：因为每个人必须按照他行为的特点倾向于其中一边，要么基督的恩宠欢迎那些超越法律的人，要么法律仍然对那些软弱者（作为它的负债人并在其掌握中）保持控制。因为一个对法律诫命有罪的人，永远无法达到福音的成全，即使他空虚地夸耀自己是基督徒，并藉主的恩宠得了释放；因为我们必须不仅把拒绝履行法律所命令的人看作仍在法律之下，而且把那些满足于仅仅遵守法律命令、从不结出相称于其召叫和基督恩宠的果实的人，也看作仍在法律之下；这里所说的不是：\BibleQuote{你应向上主你的天主献上十分之一和初果}，而是：\BibleQuote{去，变卖你所有的，施舍给穷人，然后来跟随我}；在这里，由于成全的崇高，连埋葬父亲的最短时间也不允许门徒的请求，因为人性的爱德被天主的爱德胜过。\par

% structure: p0025 source=h2 target=subsection reason=relative-heading-rank; base=h2

\subsection{第八章}

\Emph{\Person{德奥纳斯}如何劝勉他的妻子也作出她的弃绝。}\par

\Person{真福德奥纳斯}听了这话，便被一种无法抑制的对福音成全的热望所点燃，仿佛将所领受的圣言之种，存于结实的心田中，深埋在胸中破碎的犁沟里；他大为谦卑，良心自责，因为那位老人不仅说他未能达到福音的成全，甚至还说他几乎连法律的命令都未能履行；虽然他惯常每年将果实的十分之一作为施舍缴纳，却悲哀地承认自己从未听说过初果的法律；即便他也同样履行了此法律，他仍谦卑地承认，在老人看来，他仍离福音的成全甚远。于是他悲伤地回家，充满那生出悔改以致救恩的忧愁，\BibleRef{格后 7:10}并出于自己的意志和决心，将他妻子全部的心思和忧虑转向救恩；他开始激励她产生同样的热切渴望，正如他自己所被激发的，用同样的劝勉，日夜流泪催促她，使他们能一同在圣洁和贞洁中事奉天主，告诉她，他们皈依更好生活不应延迟，因为对青春的虚妄希望不能作为对抗突然死亡不可避免性的理由，这死亡同样带走孩童、青年和老年人。\par

% structure: p0028 source=h2 target=subsection reason=relative-heading-rank; base=h2

\subsection{第九章}

\Emph{当他的妻子不同意时，他如何逃到一座隐修院。}\par

当他的妻子心硬，不肯同意他，而他一再坚持这种恳求时，她却说自己正值青春年华，无法完全缺少丈夫的安慰，而且如果她被遗弃而陷入罪恶，那么破裂婚姻纽带的罪责更应归于他。对此，他长久地申述人性的状况（人性如此软弱和不确定，再与肉欲和其行为纠缠是危险的），并断言：任何人都不应割舍那他已经知道无论如何都当持守的德行，发现善行后忽视它比未发现前不爱它更为危险；而且，若他在发现了如此伟大的天上恩赐后仍偏爱卑微的世间之物，他已然陷入了堕落的罪责。此外，成全的伟大向每个年龄和性别开放，教会所有成员都被催促攀登天上善行的高峰，因为宗徒说：\Quote{你们要这样跑，好能得到奖赏；}\BibleRef{格前 9:24}那些准备好并热切追求的人不应因为缓慢和拖延之人的迟延而退缩，因为促使懒惰者被前面奔跑的人催促，要比让努力者被懒惰者妨碍更好。他此外还决定并下定决心弃绝世界，死于世界，以活于天主；若他不能获得这种幸福，即与妻子一同进入与基督的结合，他宁愿损失一个肢体而得救，像残缺的人进入天国，也不愿身体完整而被定罪。但他也补充说了以下的话：如果梅瑟因人心硬准许休妻，为什么基督不能因渴求贞洁而准许呢？尤其当同一主在其他情感中——即对父母、儿女的情感（不仅法律，连祂自己也命令要予以应有的尊重，但为了祂的名和成全的渴望，祂命令不仅要不予理会，而且要实际憎恨）——对这些，我说，祂也加上了妻子的提及，说：\Quote{凡为我的名舍弃了房屋、或兄弟、或姊妹、或父亲、或母亲、或妻子、或儿女的，必要领取百倍的赏报，并承受永生。}\BibleRef{玛 19:29}因此，祂不但不准许任何事物反对祂所宣布的成全，反而命令为爱祂而打破和忽略父母的情感纽带，尽管根据宗徒，这是第一条带有恩许的诫命，即：\Quote{应孝敬你的父亲和你的母亲——这是第一条带有恩许的诫命——为使你获得幸福，并在地上延年益寿。}\BibleRef{弗 6:2}既然福音的话谴责那些在未犯奸淫罪时破坏婚姻约束的人，它也清楚地应许那些为爱基督和渴求贞洁而摆脱肉欲之轭的人要获得百倍的赏报。因此，如果可以让你听从理性，并与我一同转向这个最可取的抉择——即我们一同事奉主，逃避地狱的痛苦——我不会拒绝婚姻的情感，反而要以更大的爱来拥抱它。因为我承认并尊敬主的话所赐予我的助手，我不拒绝在基督内以不破裂的爱之绳索与她结合，也不将主藉起初状况的法律与我联结的分离。\BibleRef{创 2:18}只要你愿意成为造物主所命定的你。但如果你不愿做助手，反而宁愿使自己成为欺骗者，成为仇敌的助手而非我的助手，并以为婚姻圣事赐给你是为了让你剥夺自己获得救恩的机会，并阻止我作为门徒跟随救主，那么我将坚定地抓住由若望院长，或更确切地说由基督亲口所说的话，使任何血肉之情不能把我从属神的恩赐中夺走，因为祂说：\Quote{谁不恼恨自己的父亲、母亲、妻子、儿女、兄弟、姊妹、田产，甚至自己的性命，不能做我的门徒。}\BibleRef{路 14:26}当这些话和类似的话不能打动这妇人的心意，她仍坚持同样固执的心硬时，真福德奥纳说：\Quote{若我不能将你从死亡中拉出，你也不能将我与基督分离；离开一个人比离开天主更为安全。}于是，靠着天主恩宠的助佑，他立即着手执行他的决定，不让他的热切愿望因任何迟延而冷却。他立即舍弃所有世俗财物，逃往一座隐修院，在那里，他在极短的时间内因圣德和谦卑的光辉而如此出名，以致当若望真福离世归主，圣厄里亚——一位不亚于前任的伟大人物——也同样去世后，德奥纳在众人评断中被选为第三任继承他们管理施舍工作的人。\par

% structure: p0031 source=h2 target=subsection reason=relative-heading-rank; base=h2

\subsection{第十章}

\Emph{解释：以免我们显得是建议与妻子分离。}\par

但不要让人以为我们编造了这一点是为了鼓励离婚，因为我们不仅绝不谴责婚姻，反而遵循宗徒的话说：\BibleQuote{婚姻应受众人尊重，床笫也应是纯洁的。}\BibleRef{希 13:4}而是为了忠实地向读者展示这位伟人奉献给天主的转变的根源。我恳请读者善意地允许，不论他是否喜欢这一点，两种情况我都无过，并将此行为的赞或咎归于真正的作者。但至于我，既然我没有提出自己的意见，而是单纯叙述事实的历史，那么，正如我不要求那些赞同所行之事的人的称赞，同样我也不应受到那些不赞同之人的仇恨攻击。因此，正如我们所说，让每个人对此事有自己的看法。但我建议他在考虑时克制批评，以免他认为自己比天主的审判更公正、更圣洁，天主审判甚至赋予他（即\Person{特敖纳}）宗徒德性的标记，更不用说如此伟大的教父们的意见，显然他的行为不仅未被指责，反而受到赞扬，以至于在选举赈济员时，他们宁愿选他而不选那些卓越杰出的人。而且我认为，那么多属神的人以天主为作者而作出的判断是不会错的，因为这判断，如上所述，被如此奇妙的标记所证实。\par

% structure: p0034 source=h2 target=subsection reason=relative-heading-rank; base=h2

\subsection{第十一章}

\Emph{关于为何在埃及他们在整个（复活期）五十天内不禁食也不在祈祷时屈膝的探询。}\par

但现在是时候履行所应许的讲话计划了。所以，当\Person{特敖纳}院长在复活期晚祷结束后，来我们的小室拜访我们时，我们在地上坐了一小会儿，开始热切地思考为什么他们如此谨慎，以至于在整个五十天内没有人在祈祷时屈膝或敢守斋到第九时辰，我们更加迫切地询问，因为我们从未在叙利亚的隐修院中看到这种习俗被如此谨慎地遵守。\par

% structure: p0037 source=h2 target=subsection reason=relative-heading-rank; base=h2

\subsection{第十二章}

\Emph{关于善、恶与中性事物本性的回答。}\par

对此，\Person{特敖纳}院长开始如此回答。确实，即使我们看不出理由，我们也应当顺从教父们的权威和我们的前辈延续了这么多年直到我们时代的习俗，并像从古代传下来的一样以不断的关心和敬意来遵守它。但是既然你们想知道此事的理由和根据，请听我们简略地讲述我们从长老们那里听来的关于此事的传承。但在引证圣经的权威之前，如果你们愿意，我们先略谈一下守斋的本性和特征，以便随后圣经的权威可以支持我们的话。天主的智慧在\BibleBook{训道}{篇}中指出，每件事——即所有快乐的事或那些被认为不幸和不快乐的事——都有其适当的时候，说：\BibleQuote{凡事都有定时，天下万务都有定期。生有时，死有时；栽种有时，拔出所栽种的也有时；杀戮有时，医治有时；拆毁有时，建造有时；哭泣有时，欢笑有时；哀恸有时，跳舞有时；抛掷石头有时，堆聚石头有时；拥抱有时，不拥抱有时；获得有时，失落有时；保存有时，舍弃有时；撕裂有时，缝补有时；静默有时，言语有时；爱慕有时，憎恨有时；战争有时，和平有时；}下面又说：\BibleQuote{因为有一个时间，}它说，\BibleQuote{对每件事和每一种行为。}所以，这些事物中的任何一个，它都没有设定为总是善的，而只有当它们恰当地并在适当的时候被施行时才是善的；因此，那些在某个时候，在适当的时刻施行时结果良好之事，若在错误或不适宜的时间妄行，就会被发现是无用或有害的；只有那些在其自身本性中是善或恶的事物例外，它们永远不能变成相反的，例如，正义、明智、刚毅、节制以及其他德行，或者另一方面，那些其描述不可能改变或归入另一类的恶习。但那些有时可以产生两种结果的事物，按照使用它们的人的性格，它们被发现是善或恶，我们认为这些事物在其本性上不绝对有用或有害，而只是根据行者的心意和时间的适宜性。\par

% structure: p0040 source=h2 target=subsection reason=relative-heading-rank; base=h2

\subsection{第十三章}

\Emph{守斋之为善，究竟是怎样的善？}\par

因此，我们现在必须探究，我们应如何理解斋戒的状态：我们是否意指它如同义德、智德、勇德和节德那样，本质上是善，绝不可能变成别的什么；或者它是一件中性的事，有时做是有益的，有时可能不遭谴责地省略；并且有时做是错的，有时省略是值得称许的。因为如果我们认为斋戒包括在那些德性之列，以致禁食被视为本身为善之事，那么享用食物必定是恶的、错的。因为任何与本身为善之物相反的东西，必定被认为本身为恶。但圣经的权威不允许我们如此规定。因为如果我们怀着这样的思想和意图禁食，以致认为进食是犯罪，那么我们不仅不会从禁食中获得任何益处，反而会招致严重的罪责，陷入不敬之罪，如宗徒所说：\Quote{禁止食用天主为信者和认识真理的人所造、要怀着感恩之心领受的食物。因为天主所造的一切都是好的，如果怀着感恩之心领受，就没有一样是可弃绝的。}因为\Quote{若有人以为某物是不洁的，那物对他来说便是不洁的。}因此，我们从未读到有人纯粹因进食而受谴责，只有当某物与之相连或随后发生，他因此应受谴责时，才会如此。\par

% structure: p0043 source=h2 target=subsection reason=relative-heading-rank; base=h2

\subsection{第十四章}

\Emph{斋戒如何在本质上并非善}\par

因此，这也清楚地表明它是一件中立之事：因为正如守斋带来成义，破斋也不带来定罪；除非也许违抗命令比进食带来惩罚。但就一件本身为善的事而言，不应有任何时刻缺少它，以致人可以离了它而活；因为若它停止，疏忽它的人必定陷入邪恶。同样，也没有任何时间给予本身为恶的事，因为有害之事若纵容它，它必造成伤害，也绝不可能成为值得称赞的。此外，明显可见，那些我们看到有条件和时间规定的事物，它们被遵守时圣化人，被忽略时也不败坏我们，这些都是中立之事，例如：婚姻、耕作、财富、退隐旷野、守夜、阅读和默想圣经，以及我们讨论所由起的守斋本身。所有这一切，天主诫命和圣经的权威规定，不应如此不停地追求或如此持续遵守，以至于暂时中断它们便成为错误。因为任何绝对命令的事，若未完成，便带来死亡；但那些我们被劝勉而非命令的事，做了有益，不做不带来惩罚。因此，对于所有这些或其中某些事，我们的先辈命令我们要么审慎地去做，要么按照原因、地点、方式、时间仔细遵守，因为若相宜地去做，便是合宜合用的；若不相宜，便成为愚拙有害的。若在兄弟来到时——他本应在兄弟身上以礼遇款待基督并最亲善地欢迎他——人却选择严守斋戒，难道他岂不是更该被定为失礼，而非获得虔诚的称赞或赏报吗？或者，若肉体的衰竭或软弱需要进食以恢复力量，人却不肯放松其刻苦的严规，难道他不是该被视为自身身体的残忍凶手，而非关心自己救恩的人吗？同样，当节期允许适当放宽饮食而必须丰盛用餐时，若人执意坚守严格守斋，他必须被视为不虔敬，反是粗野无理。但那些藉守斋寻求人称赞，并因愚昧显露苍白而得圣德名声的人，这些事也会成为恶事。福音之言告诉我们，他们已得了他们的赏报，主也藉先知憎恶他们的守斋。先知首先代他们向主提出异议说：\Quote{我们为什么守斋而祢不理睬？我们为什么刻苦灵魂而祢不知道？} 然后他立即回答并解释他们不应被俯听的原因：\Quote{看哪，} 他说，\Quote{在你们守斋的日子，你们寻欢作乐，且压榨你们的债户。看哪，你们守斋是为争辩和争斗，以凶恶的拳头打人。不要像你们今日这样守斋，使你们的声音高高可闻。难道我所拣选的斋戒是这样的吗？使人一天刻苦灵魂？难道这是：像圆圈般环绕头颅，铺上麻衣和灰烬？你能称这为斋戒，为主所悦纳的日子吗？} 接着他继续教导守斋者的刻苦如何成为可接纳的，并清楚指出守斋本身不能单独为善，只有当附带以下原因时才行：\Quote{岂不是这样吗？} 他说，\Quote{我所拣选的斋戒吗？解开邪恶的绳锁，解除轭索的重压，让被压迫者自由，折断一切轭。把你的饼分给饥饿的人，将无家可归的穷人领进你的家；见到赤身露体者给他衣服，不要轻视你的骨肉。那时，你的光将如晨光升起，你的健康将迅速恢复，你的正义必在你面前行走，主的荣耀必作你的后盾。那时你呼求，主必应允；你哀号，祂必说：我在这里。} \BibleRef{依 58:3} 那么你看，守斋确实不被主视为本身为善的事，因为它成为善和悦乐天主的事，不是靠自身，而是靠其他行为；并且，根据周围环境，它可能被视为不仅虚妄，甚至可憎，如主说：\Quote{他们守斋时，我不听他们的祈祷。} \BibleRef{耶 14:12}\par

% structure: p0046 source=h2 target=subsection reason=relative-heading-rank; base=h2

\subsection{第十五章}

\Emph{本身为善的事不应为了某种较小的善而去做。}\par

我们不应为了斋戒而实行怜悯、忍耐和爱德，以及上述诸德的规诫（其中存在在本性上为善的事物），而应为了它们而实行斋戒。因为我们的努力必须是通过斋戒获得那些真正为善的德行，而不是行那些德行以达到斋戒为其目的。因此，肉身的刻苦是有用的，节制的疗法也必须被采用；即，借着它我们能够成功达至爱德，其中有不变化的善，且持续不断，没有时间的例外。因为医药、金匠的艺术以及世上其他艺术的体系，并不是为了特定工作的工具而存在；而是工具为实践艺术而准备。正如它们对懂得它们的人有用，对不懂那艺术体系的人则无用；正如它们对依靠其帮助完成工作的人是大帮助，对不知道它们为何目的而造、只满足于拥有它们的人则毫无用处，因为他们将一切价值归于单纯拥有，而非完成工作。那么，那在本性上是最好的事物，中立之事为其而做，但那至善却不为了任何其他事物，而是因其本身的内在善性。\par

% structure: p0049 source=h2 target=subsection reason=relative-heading-rank; base=h2

\subsection{\Emph{第十六章}}

\Emph{如何区分本性为善的事物与其他为善的事物？}\par

这可以从我们称为中立的那些事物中以这些方式区分：如果一事物本身为善，而非因其他事物；如果它因其自身有用，而非因其他事物；如果它不变地且在任何时候为善，且总是保持其特性，而绝不会变成任何不同的事物；如果它的移除或停止不可能不产生最大的危害；如果它的对立面同样在本性上是恶的，且绝不会变成善。这些描述，凭之可以区分本性为善的事物之性质，不能应用于斋戒，因为斋戒本身不是善的，也不是因其自身有用，因为它被明智地用于获得心灵与身体的纯洁，以便肉身的刺痛被缓解，灵魂得以平静并与造物主和好，也不是不变地且在任何时候为善，因为常常我们并不因停止斋戒而受损，而且有时如果不合理地实行，它反而变得有害。其对立面，即进食，似乎是恶的，但并非本性为恶，因为进食本为自然适意，不能被视为恶，除非无节制和奢侈或其他过失是它的结果；因为\BibleQuote{不是入于口的使人污秽，而是出于口的才使人污秽。}\BibleRef{玛 15:11}所以人若轻视本性为善的事物，且不恰当或无罪地对待它，若他不为其自身而为之，而为其他事物，则一切其他事物都应为其而做，但它本身应单独寻求。\par

% structure: p0052 source=h2 target=subsection reason=relative-heading-rank; base=h2

\subsection{\Emph{第十七章}}

\Emph{论斋戒的理由及其价值。}\par

因此，让我们时常记住这关于斋戒性质的描述，并以灵魂的一切力量始终朝向它，这样认识到只有我们在其中持守时间、性质和程度，才适合我们，而且不将我们希望的目的设定于其上，而是借它成功达至心灵的纯洁和宗徒的爱德。因此，由此清楚可知，斋戒——不仅有其特定季节应实行或放松，还有条件与规则制定——在本性上不是善的，而是中立的。但那些被诫命的权威命令为善或被禁止为恶的事物，永远不受任何时间例外的影响，以致有时我们应做被禁止的事或省略被命令的事。因为对正义、忍耐、节制、谦逊、爱德没有限制，另一方面也从未允许不义、急躁、愤怒、不谦逊、嫉妒和骄傲。\par

% structure: p0055 source=h2 target=subsection reason=relative-heading-rank; base=h2

\subsection{\Emph{第十八章}}

\Emph{斋戒如何不总是适合。}\par

因此，既然我们事先提出了关于斋戒条件的这点，似乎适合附加圣经的权威，借此将更清楚地证明斋戒既不能也不应常守。在福音中，当法利塞人与若翰的门徒一起斋戒时，宗徒们作为天上新郎的朋友和伴侣，尚未遵守斋戒的规矩，若翰的门徒（他们以为凭自己的斋戒获得了完美的义德，因为他们跟随那位伟大的悔改宣讲者，他以自己的榜样给所有人提供了模范，不仅拒绝供人使用的各类食物，而且实际上完全不吃众人共同的面包）向主抱怨说：\Quote{为什么我们和法利塞人屡次斋戒，而你的门徒却不斋戒？}主在回答中清楚表明，斋戒并非在所有时候都适合或必要，当任何节期或爱德的机会介入并容许饮食时，说：\BibleQuote{新郎的伴郎岂能当新郎与他们在一起时悲哀？但日子将到，新郎要从他们中被带走，那时他们就要斋戒；}\BibleRef{玛 9:14}这些话虽然是在祂身体复活之前说的，却特别指向复活节期，在此期祂复活后四十天与门徒一同进食，他们因祂每日同在的喜乐而不斋戒。\par

% structure: p0058 source=h2 target=subsection reason=relative-heading-rank; base=h2

\subsection{\Emph{第十九章}}

\Emph{为何我们在整个复活节期间破斋的问题？}\par

\Person{格玛诺}：那么，为什么我们在整个五十天中放松饮食的刻苦，而基督复活后只与门徒同在了四十天？\par

% structure: p0061 source=h2 target=subsection reason=relative-heading-rank; base=h2

\subsection{第二十章}

\Emph{答复。}\par

你那切题的问题值得被告知完美的真实理由。在我们的救主复活后第四十天发生的升天之后，宗徒们从橄榄山上下来，祂在回归圣父时曾允许他们在该山看见祂，正如\WorkTitle{宗徒大事录}所证实的，他们进入\Place{耶路撒冷}，并被告知等待圣神降临十天，当这些日子在第五十天满期时，他们欣喜地领受了圣神。因此，如此这个节日的数目就清楚地构成了，我们读到它在旧约中也有预表的预象，其中当七个星期满期时，司祭奉命向主奉献初熟之果的饼：\BibleRef{申 16:9}；这确实被显示为奉献给主，通过宗徒们在那天向百姓所作的宣讲；真正的初熟之果的饼，当从新教义的训导中产生时，将犹太人的初果祝圣为基督徒子民奉献给主，五千人因那食物的恩赐而饱足。因此，这十天应与前四十天一样以同等的隆重和喜乐来遵守。而这关于这节日的传统，由宗徒的传人传给我们，应以同样的统一性来遵守。因此，在那些日子里，他们在祈祷时不屈膝，因为屈膝是补赎和哀悼的标记。因此，在这些日子里，我们在一切事上遵守与主日相同的隆重礼仪，我们的前辈教导说，在主日人不应当斋戒或屈膝，出于对主复活的敬意。\par

% structure: p0064 source=h2 target=subsection reason=relative-heading-rank; base=h2

\subsection{第二十一章}

\Emph{一个疑问：放宽斋戒是否不利于身体的贞洁。}\par

\Person{杰曼诺斯}：难道肉体被如此长节期的异乎寻常的享乐所吸引，不会从已除去的罪的诱因中产生出一些荆棘吗？或者灵魂被不习惯的宴饮的消耗所压垮，难道不会削弱其对仆役身体管治的严苛吗？尤其是，在我们这个成熟的年龄，如果我们斗胆比平时更多量地取用惯常食物，或者比平时更自由地取用不习惯的食物，我们的顺从肢体可能会迅速反叛？\par

% structure: p0067 source=h2 target=subsection reason=relative-heading-rank; base=h2

\subsection{第二十二章}

\Emph{关于保持节制之道的答复。}\par

\Person{德奥纳斯}：若我们以理智的合理判断衡量我们所做的一切，并常以心中的纯洁咨询我们自己的良心而非他人的意见，那么，这休息的间隔就肯定不会干扰我们应有的严谨，只要，如所说，我们纯洁的头脑公正地考虑放纵与克己的适当限度，公平地抑制任何一方的过度，并以真正的辨别力分辨美酒佳肴的重量是否成为我们精神的负担，或者过度的克己是否压低了另一边，即身体，并使它看到的那边因被抬起或压低而相应地降低或升高。因为我们的主不愿意任何事为他的尊荣和光荣而行却不经过判断的调和，因为\Quote{君王的光荣爱慕判断}，因此最聪明的\Person{所罗门}劝我们不要让我们的判断偏向任何一边，说：\BibleQuote{你要以公义的劳苦尊崇天主，并将你义德的果实献给他。}\BibleRef{箴 3:9}因为在我们良心中居住着一位不受腐蚀的真正审判者，他有时在众人皆错时，是唯一不被我们纯洁状态所欺骗的人。因此，我们应竭尽全力在我们谨慎的心中保持一个恒定的目的，以免如果我们明智的判断失误，我们会被不当克制的欲望点燃，或被过度放纵的愿望诱惑，从而在不公平天平的舌头上称量我们力量之实质；但我们应将灵魂的纯洁放在一个秤盘中，将身体的力气放在另一个秤盘中，并以良心的真正判断来衡量它们，使我们不至于因对其中一方的过分欲望而不义地将公平的天平倾向一边，或倾向不当的严格，或倾向过度的放纵，从而因过度的严格或放纵而有人对我们说：\Quote{你若奉献正当，却未正当分割，岂不犯了罪？}因为那些我们因粗暴地撕裂自己的肠腑而思虑不周地强求来的斋戒奉献，并幻想我们正确地奉献给了主，这些却被那\Quote{爱慕仁慈与判断}的主所憎恶，他说：\Quote{我主爱慕判断，却厌恶燔祭中的抢夺。}\BibleRef{耶 48:10}那些取走他们奉献的主要部分，即他们的职责和行动，用于有益自己的肉身，而将剩余部分和一小部分留给主的人，天主圣言如此谴责他们是欺诈的工人：\BibleQuote{凡欺诈地行主的工的，是可咒骂的。}\BibleRef{耶 48:10}因此，主责备那些以不公平的考量自欺欺人的人，并非没有理由，他说：\Quote{世人尽是虚荣：世人在天平上撒谎，好欺骗人。}因此，蒙福的宗徒警告我们要握住明智的缰绳，不被过度所吸引而摇摆到任何一边，他说：\BibleQuote{你们合理的敬礼。}\BibleRef{罗 12:1}而律法的颁布者也同样禁止同一件事，他说：\BibleQuote{天平要公正，砝码要相等，公平的布色尔和公平的塞克斯塔留斯，}\BibleRef{肋 19:36}\Person{所罗门}也就此事给出了类似的意见：\BibleQuote{大砝码和小砝码，以及双重的量器，在主前都是不洁的，使用它们的人必被他的计谋所阻挠。}\BibleRef{箴 20:10}此外，不仅在我们所说的方法上，而且在这件事上我们也必须努力，不要在我们心中有偏私的砝码，也不要在我们良心的仓房中有双重的量器，即不要用过于严格和比我们自己所能承受的更重的诫命来压倒那些我们应该向他们宣讲主道的人，同时我们却想当然地认为，对我们自己，那些与严格规则有关的事可以通过更自由地允许放松而变得缓和。因为当我们这样做时，除了用双重的重量和尺度来称量和测量主命令的财富和果实之外，还能是什么呢？因为我们若用一种方式为自己施行，又用另一种方式为我们的弟兄施行，我们就理所当然地被主责备，因为我们有不公平的天平和双重的量器，符合\Person{所罗门}的话告诉我们：\Quote{双重的砝码是主所憎恶的，欺骗性的天平在他眼中不为美善。}这样，我们也明显地犯下了使用欺骗性砝码和双重量器的罪，如果出于贪图人的称赞，我们在弟兄们面前表现出比我们在自己斗室私下实行时更大的严格，试图在人前显得比在天主前更克己和更圣洁，这是一种我们不仅应该避免而且应该深恶痛绝的恶。但同时，既然我们已经从眼前的问题偏离了一段距离，让我们回到我们开始的地方。\par

% structure: p0070 source=h2 target=subsection reason=relative-heading-rank; base=h2

\subsection{第二十三章}

\Emph{关于休息的时间和尺度}\par

为此，我们应以这样的方式遵守所提及的节期，使所允许的松懈对身心的益处大于害处，因为任何节日的喜乐都不能削弱肉身的刺，那个凶恶的仇敌也不会因对日子的尊重而平息。因此，为了既保持为节期所规定的习俗，又不超越最有益的禁食规则，我们只需将允许的松懈限制在这一点上：出于对节期的尊重，我们可以在第六时辰稍早进食本应在第九时辰吃的食物，但条件是食物的通常数量和性质不得改变，以免复活期的松懈导致四旬期禁食所获的身心纯洁与正直丧失，并且我们因斋戒所获得的无益于我们，而粗心的饱足又使我们立刻失去它，尤其因为我们的仇敌狡猾，常在庆祝某个节日时，见我们对纯洁的守护稍有松懈，便主要攻击我们纯洁的堡垒。因此，我们必须极其警惕，以免灵魂的力量因诱人的奉承而软弱，并且我们因四旬期的持续努力而获得的贞洁纯洁，因复活期的休息和粗心而丧失。因此，食物的质量或数量绝不应有任何增加，即使在最大的节日上，我们也应同样禁食那些在普通日子因禁食而保持正直的食物，以免节日的喜乐在我们内心激起最致命的肉身情欲冲突，从而转为悲伤，并终结那因纯洁之喜乐而欢欣的最卓越的心灵节日；在短暂的肉身喜乐展示之后，我们便开始以持久的忏悔之痛哀悼失去的心灵纯洁。此外，我们应当努力，免得先知劝诫的警告徒然针对我们说出：\Quote{犹大啊，庆祝你的节日，并偿还你的誓愿。}因为若节期的出现不妨碍我们禁食的连续性，我们将不断享受属灵的节日，这样，当我们停止奴役性的工作时，\BibleQuote{将有一月又一月，安息日又安息日。}\BibleRef{依 66:23}\par

% structure: p0073 source=h2 target=subsection reason=relative-heading-rank; base=h2

\subsection{第二十四章}

\Emph{关于守四旬期的不同方式的一个问题。}\par

\Person{杰曼努斯}说：为什么四旬期守六个星期，而在某些国家，对宗教的更热忱关怀似乎又增加了第七个星期，尽管减去星期日和星期六后，两个数字都未达到四十天？因为这些星期只包含三十六天。\par

% structure: p0076 source=h2 target=subsection reason=relative-heading-rank; base=h2

\subsection{第二十五章}

\Emph{关于四旬期的斋戒与年度什一之物相关的回答。}\par

\Person{泰奥纳斯}说：虽然某些人虔诚的单纯会搁置对这一问题的探讨，但因为你们在考察他人认为不值得询问的事情上更为谨慎，并想知道我们这一习俗的全部真理及其秘密，你们也将得到对此的非常清楚的解释，以便更明显地确信我们前辈所教导的并非不合理。按照梅瑟的法律，向全体民众颁布的命令是：\BibleQuote{你要将你的什一之物和初果献给上主你的天主。}\BibleRef{出 22:29} 因此，既然我们奉命献上我们财产和所有果实的什一之物，那么，献上我们生命、日常工作和行动的什一之物就更为必要，这一点在四旬期的计算中确实清楚地安排了。因为一年循环中所有日期的什一之物是三十六天半。但在七个星期中，减去星期日和星期六，剩下三十五天的斋戒。但若加上复活节前夕（星期六的斋戒延长到复活节破晓鸡鸣时），不仅三十六天的数目得以补足，而且关于似乎多出的五天的什一之物，若将所加的那段夜间考虑在内，总数便毫无欠缺。\par

% structure: p0079 source=h2 target=subsection reason=relative-heading-rank; base=h2

\subsection{第二十六章}

\Emph{我们应当如何也将我们的初果献给主。}\par

但是，对于那每日由所有忠心事奉基督的人所献上的初果，我该说什么呢？因为当人们从睡眠中醒来，经过休息后重新振作起来，在接纳任何冲动或思想到心中，或者想起任何事务或考虑之前，将他们最早的思想作为神圣的奉献而祝圣，他们所行的难道不正是通过大司祭耶稣基督将他们的出产初果献上，为今生的享受和每日复活的预像吗？同样，当他们以同样的方式从睡眠中醒来时，向天主献上喜乐的祭品，并以舌头最初的动作呼求祂，赞美祂的名，首先张开嘴唇向祂歌唱赞颂，将他们口中的职务献给天主；同样地，当他们从床上起身，站立祈祷，在为自己的目的使用肢体的服务之前，他们也以同样的方式将手足的最初奉献献给祂；他们不为自己取用任何服务，而是为祂的光荣前进脚步，为赞美祂而迈步，如此藉着伸开手、弯曲膝盖、俯伏全身而奉献他们所有运动之初果。因为，除非我们在睡眠后的休息中，如我们以上所说的，从黑暗和死亡似的被恢复到这光明中时，我们有勇气不开始将灵魂和肉体的一切服务取为己用，否则我们无法成就我们在圣咏中所歌唱的：\Quote{我提前破晓而哀号；}以及：\Quote{我的眼睛向着祢提前清晨，为能默思祢的言语；}以及：\Quote{清晨我的祈祷要提前临到祢面前；}因为先知所\InnerQuote{提前}的，或者我们应当以同样方式提前的，没有别的清晨，除非是我们自己，即我们的职务、情感和地上的挂虑，没有它们我们不能存在——或者是仇敌最微妙的暗示，当我们仍在休息并被睡眠战胜时，他试图通过虚诞之梦的幻像向我们暗示，以这些幻像，当我们随即醒来时，他将充满我们的心思并占据我们，以便他首先夺取并掠走我们初果的掠物。因此，我们必须极其谨慎（如果我们想要在行为上实现上述引文的含义），以便一种焦急的警觉关注我们最早清晨的思想，免得它们被我们嫉妒的仇敌仓促夺取而预先玷污，从而使我们的初果被主当作无价值和平凡的而拒绝。如果我们不以警觉的谨慎心态提前防备他，他不会放弃他那可怜的抢先于我们的习惯，也不会日复一日地停止用他的诡计抢先于我们。因此，如果我们想要从我们心灵的果实中献上悦纳而令天主喜悦的初果，我们应当给予不寻常的注意，以保持我们身体的一切感官，尤其是在清晨的时刻，作为向主奉献的神圣全燔祭，纯全而无玷污。这种虔诚，甚至许多生活在世俗中的人也极其谨慎地遵守，他们在天亮或清晨之前起身，并不立即混杂于世界日常必要的事务，而是先赶往教堂，努力在天主面前奉献他们所有行为与事迹的初果。\par

% structure: p0082 source=h2 target=subsection reason=relative-heading-rank; base=h2

\subsection{第二十七章}

\Emph{为什么许多人以不同的天数守四旬期。}\par

进一步，至于你所说的，即在某些地方四旬期以不同的方式遵守，即六周或七周，其实是同一制度，同样的斋戒方式，只是通过不同的周数而保持。因为那些认为星期六也应当守斋的人，已决定遵守六周。因此他们每周七天中守斋六天，这样重复六次便构成三十六天。因此，正如我们所说，虽是同一制度，同样的斋戒方式，尽管周数似乎有所不同。\par

% structure: p0085 source=h2 target=subsection reason=relative-heading-rank; base=h2

\subsection{第二十八章}

\Emph{为何它被称为\OriginalTerm{四旬期}{Quadragesima}，尽管守斋只守三十六天。}\par

然而，由于人的疏忽，这理由被忽略了，这个季节——如上所述，一年中的什一奉献通过三十六天半的斋戒来奉献——便被称为\OriginalTerm{四旬期}{Quadragesima}，他们或许认为应当为此如此称呼；即，据说梅瑟、厄里亚和我们的主耶稣基督本人曾斋戒四十天。与此数字的奥妙相称的，还有以色列在旷野中居住的四十年，以及他们以奥秘意义所经过的四十站。或者，什一奉献从海关的惯例中得到了四旬期的名称。因为那国家税通常如此称呼，其中同样比例的增收被指定给君王使用，正如法律上的四旬期税，万世之王为我们的生命使用向我们索取的。无论如何，尽管这与所提出的问题无关，但我认为我不应省略这一事实：我们的长辈常常见证，特别是在这些日子里，整个隐修士团体按照敌对民众的古老习俗受到攻击，并被更强烈地催促离弃他们的家园，原因如下：按照这预象，埃及人从前以严重的苦难压迫以色列子民，如今的精神埃及人也试图以艰难卑贱的劳作压倒真正的以色列，即隐修士群体，以免借着天主所喜爱的平安，我们离开埃及地，为我们的益处穿越到德行之旷野，以至于法郎向我们发怒说：\Quote{他们是懒惰的，因此他们喊叫说：让我们去向我们的主天主献祭。让他们被劳役压迫，在他们的工作中受骚扰，他们就不会被虚妄的话所骚扰。} \BibleRef{出 5:8} 因为他们的愚妄确实认为，在纯洁心灵的旷野中所献上的主的圣祭是愚妄的顶峰，因为\Quote{宗教为罪人是可憎的。} \BibleRef{德 50:24}\par

% structure: p0088 source=h2 target=subsection reason=relative-heading-rank; base=h2

\subsection{第二十九章}

\Emph{那些成全者如何超越四旬期的固定规则。}\par

因此，正直成全的人不受这四旬期法律的约束，也不满足于仅仅服从教会首领为那些整年沉浸在享乐或事务中的人所设立的微小规则，以便他们能被这法律要求所约束，并在这几天内至少为主腾出时间，将他们生命日子的什一奉献给主，而他们本会将其全部作为利润消耗。但是，义人——法律不是为他们设立的——他们不是将一小部分，即十分之一，而是将整个生命的时间用于属灵职责，因为他们在法律上免除了什一奉献的负担，因此，如果有任何值得和虔诚的机缘临到他们，他们便准备好毫不犹豫地放松他们的站斋。因为在他们的情形中，所减少的不是微不足道的什一奉献，因为他们将所有的一切都与自己一起献给主。而一个人若没有自愿奉献任何东西给天主，却被那不容推辞的严酷法律强迫缴纳什一奉献，却不能这样做而不犯下严重的不义。因此，显然可见，法律的仆人不能成为成全的，他只避开那些禁止的事，并做那些命令的事；但那些真正成全的人，甚至法律所允许的事也不利用。这样，虽然关于梅瑟法律说\Quote{法律未曾使什么成为成全的} \BibleRef{希 7:19}，但我们读到旧约中有些圣人曾是成全的，因为他们超越了法律的命令，生活在福音的成全之下：\Quote{因为知道法律不是为义人设立的，乃是为不法和不顺服的人，为不虔敬和犯罪的人，为不圣洁和亵渎的人等。} \BibleRef{弟前 1:9}\par

% structure: p0091 source=h2 target=subsection reason=relative-heading-rank; base=h2

\subsection{第三十章}

\Emph{论四旬期的起源和开端。}\par

然而，你应当知道，当初期教会保持其成全不变时，这四旬期的遵守并不存在。因为他们不受这秩序的要求或任何法律规定的约束，也不被局限在非常狭窄的斋戒界限内，因为斋戒同样涵盖全年。但当信徒的群体日渐从那种宗徒的热忱中衰退，开始照顾自己的财富，并按照宗徒的安排不为所有信徒的利益而分配，反而着眼于自己的私用，试图不仅保留财富，甚至增加它，不满足于仿效阿纳尼亚和撒斐辣的榜样时，所有司铎便认为，那些被世俗忧虑所困、几乎（我可以说）不知道节制和痛悔的人，应当通过教会法令规定的斋戒被召回虔敬的本分，并被缴纳法律什一奉献的必要性所约束，因为这无疑对软弱的弟兄有益，也不能损害那些生活在福音恩宠下、以自愿的奉献超越法律而得以达到宗徒所说的真福的成全者：\Quote{因为罪必不能主宰你们，因为你们不在法律之下，而在恩宠之下。} \BibleRef{罗 6:14} 的确，罪不能主宰那在恩宠的自由中忠信生活的人。\par

% structure: p0094 source=h2 target=subsection reason=relative-heading-rank; base=h2

\subsection{第三十一章}

\Emph{一个问题：我们应如何理解宗徒的话：\Quote{罪必不能主宰你们}。}\par

\Person{杰尔曼诺}：因为宗徒的这句话，不仅向隐修士，也向所有基督徒普遍应许了无忧无虑的自由，它不会误导我们，但对我们来说似乎有些晦涩。因为他坚持所有信福音的人是自由的，脱离了罪的轭和主宰，那么罪的统治如何几乎在所有领洗者身上有力地掌权，正如主的话所说的：\Quote{凡犯罪的，就是罪的奴隶}？ \BibleRef{若 8:34}\par

% structure: p0097 source=h2 target=subsection reason=relative-heading-rank; base=h2

\subsection{第三十二章}

\Emph{论恩宠与法律诫命之间的区别。}\par

\Person{德奥纳}说：你的询问再次在我们面前提出一个范围不小的问题。尽管我知道，对无经验者而言，这个问题无法通过教导来传授或理解，但我会尽力用言语加以阐述并简要解释，只要你们的心智能跟随并实践我们所说的话。因为凡是凭经验而非教导所知的事，正如无经验者无法教导，同样，没有通过类似学习和训练打下基础的人也无法领会或接受。因此，我认为我们首先需要稍微仔细地探究，法律的目的或意义是什么，以及恩宠的体系和完美是什么，以便得以理解罪的权势以及如何驱逐它。因此，法律主要命令人寻求婚姻的纽带，说：\Quote{凡在熙雍有后裔，在耶路撒冷有家室的人，是有福的；}以及：\Quote{不生育、不怀孕的妇人是有祸的。}\BibleRef{约 24:21}另一方面，恩宠邀请我们追求永恒贞洁的纯洁与有福童贞的无玷状态，说：\Quote{不生育的和没有乳哺过的有福了；}以及：\Quote{凡不恨自己的父亲、母亲和妻子的，不能做我的门徒；}还有宗徒的话：\Quote{今后有妻子的，要像没有妻子一样。}法律说：\Quote{不可迟延奉献你们的什一和初熟之物；}恩宠说：\Quote{你若愿意成全，去变卖你所有的一切，施舍给穷人；}法律不禁让人报复冤仇、以牙还牙，说：\Quote{以眼还眼，以牙还牙。}恩宠却愿意我们的忍耐通过加倍的伤害和打击得到证明，并命令我们准备承受双倍损失，说：\Quote{有人打你的右脸，连左脸也转过来由他打；有人想要与你争讼，拿你的里衣，连外衣也由他拿去。}一条命令我们恨仇敌，另一条却命令我们爱仇敌，以至于认为我们也应当常为他们向天主祈祷。\par

% structure: p0100 source=h2 target=subsection reason=relative-heading-rank; base=h2

\subsection{第三十三章}

\Emph{论福音的诫命较法律诫命更为温和。}\par

因此，凡登上福音成全之顶峰者，立刻因如此德行之功劳而被提升于一切法律之上，视梅瑟所命令的一切为微不足道，并承认自己只服从救主的恩宠；他知道赖其助佑，自己已达至那最崇高的境界。因此罪不能管辖他，因为\BibleQuote{天主的爱，借着所赐与我们的圣神，已倾注在我们心中了}\BibleRef{罗 5:5}排除了对其他一切事物的挂虑，既不能贪求禁止的事，也不忽视所命令的事，因为其全部目标与渴望都常集中于天主的爱；它如此不被无价值之事的快乐所捕获，甚至实际上不利用那些允许的事。但在法律之下，虽遵守合法的婚姻，放荡的游荡被约束，局限于一个女子，但肉欲的刺激仍强烈；为使有明确供应的燃料的火被限制在预先规定的界限内而不蔓延、不烧及所触之物，实属不易。即使这火遇到这样的障碍：不可在这些界限之外燃烧，但即便被抑制，它仍在燃烧，因为意志本身有过失，而肉体交合的习惯使之仓促陷入过于频繁的奸淫过犯。然而，那些被救主的恩宠以贞洁的神圣之爱点燃的人，在主爱的火中烧尽一切肉欲的荆棘，以致没有罪恶的余烬干扰他们贞洁的冷静。法律的奴仆因使用合法之事而堕入非法；恩宠的分享者既忽视合法之事，对非法之事一无所知。但正如罪在爱婚姻者身上是活跃的，在那些仅仅满足于缴纳什一和初果的人身上也是如此。因为，当他懒散或疏忽时，他必然在什一和初果的质、量或每日分配上犯罪。因为他被命令不疲倦地将自己的财物分施给有需要的人，虽然他可能以最大的信德和虔诚来分施，但他还是难以不常常陷入罪的网罗。但是，对于那些没有藐视主的劝告，反而将全部财产分给穷人，背起十字架跟随恩宠赐予者的人，罪不能管辖他们。因为，那虔诚地分施并散尽已献给基督、不再视为己有的财富的人，不会因获取食物而有无信德的忧虑；也不会有吝啬的犹豫减损他施舍的喜乐，因为他无需考虑自己的需要或担心食物短缺，而是在分施那已一次性完全献给天主、不再视为己有的财物；他确信，一旦如愿以偿地舍弃一切，天主必会喂养他，远超过喂养天空的飞鸟。另一方面，那保留世上财物的人，或被旧法律的规定束缚，分施其出产的什一、初果或部分收入的人，虽然他可能借此施舍的甘露在一定程度上熄灭罪恶之火，但无论他多么慷慨地施舍，他也不可能完全摆脱罪的管辖，除非藉救主的恩宠，连同他的财物，也除去一切占有之爱。同样，凡选择按法律所命以眼还眼、以牙还牙或恨仇敌的人，不能不服于罪恶的血腥统治；因为当他渴望以报复来补偿所受的伤害，并对仇敌怀有苦毒的仇恨时，他必定常被忿怒和暴怒的激情所燃烧。但是，凡生活在福音恩宠之光下，以不抵抗而忍受来战胜邪恶，并自愿毫不迟疑地给打他右颊的人也伸出左颊，给想控告他拿他外衣的人连斗篷也给他，且爱仇敌、为毁谤他的人祈祷的人，他已折断罪的轭，挣断罪的锁链。因为他已不生活在法律之下，法律不消灭罪的种子（因此宗徒并非无理由论及法律说：\Quote{先前的诫命因它的软弱无能而被废除，因为法律不能成就什么完全的事}；主也藉先知说：\Quote{我给了他们不良的诫命和不能使人生活的法令}），而是生活在恩宠之下，恩宠不仅砍掉邪恶的枝条，而且确实拔除邪恶意志的根。\par

% structure: p0103 source=h2 target=subsection reason=relative-heading-rank; base=h2

\subsection{第三十四章}

\Emph{一个人如何能显示自己处于恩宠之下}\par

凡努力达到福音教导之成全的人，这人在恩宠之下，不被罪恶的权势所压迫，因为处于恩宠之下，就是去行恩宠所命令的事。但谁若不愿服从福音成全的全面要求，就不可不知道，虽然他看似受了洗，也成了隐修士，他却不在恩宠之下，反而仍被法律的锁链所束缚，被罪恶的重担所压。因为那藉收养之恩接纳所有接受祂的人的那一位，其目的不是要摧毁，而是要建造，不是要废除，而是要成全梅瑟的法律。但有些人对此一无所知，并漠视基督那卓越的劝告与鼓励，竟因一种过早假设的自由的疏忽而如此放纵，以致他们不仅不遵行基督的命令，仿佛这些命令太难，反而鄙视梅瑟法律给初学和孩童的命令，视其为过时，并在这种危险的自由中说出了宗徒所谴责的话：\BibleQuote{我们犯了罪，因为我们不在法律之下，而在恩宠之下。}\BibleRef{罗 6:15}那么，那既不在恩宠之下（因为他从未攀登主教导的高峰），又不在法律之下（因为他连法律的那些微小命令也没有接受）的人，在罪恶的双重统治下受折磨，竟幻想自己已经领受了基督的恩宠，而唯一的目的就是凭这危险的自由使自己不属于祂，从而陷入伯多禄宗徒警告我们要避免的状态，他说：\BibleQuote{你们要作自由的人，却不可用这自由作恶的掩饰。}蒙福的保禄宗徒也说：\BibleQuote{弟兄们，你们蒙召是为得自由，}即脱离罪恶的统治，\BibleQuote{只是不可用这自由作为放纵肉欲的机会，}即以为废除法律的命令就是犯罪的许可。但这自由，保禄宗徒教导我们，唯在主居住之处才有，因为他说：\BibleQuote{主就是那圣神；主的圣神在哪里，哪里就有自由。}\BibleRef{格后 3:17}因此，我不知能否如同那些有经验的人一样表达和解释蒙福宗徒的意思；我只知道，对于所有完全获得了\OriginalTerm{实践训练}{\GreekText{πρακτική}}，即实际操练的人，即使没有人解释，这也非常清楚地显明了。因为他们在讨论中无需费力就能理解他们已由实践学到的东西。\par

% structure: p0106 source=h2 target=subsection reason=relative-heading-rank; base=h2

\subsection{第三十五章}

\Emph{一个问题：为何有时我们比平时更严格守斋，反而比平时更强烈地被肉欲所困扰？}\par

\Person{格尔玛诺}说：您非常清楚地解释了一个极其困难的问题，一个我们认为许多人所不知的问题。因此我们恳求您也为我们加上这点好处，并仔细解释为何有时我们比平时更严格守斋，筋疲力尽时，反而激起更严重的肉身挣扎。因为常常在睡醒时，若发现自己被玷污，我们便如此心灰意冷，甚至不敢忠信地起来祈祷。\par

% structure: p0109 source=h2 target=subsection reason=relative-heading-rank; base=h2

\subsection{第三十六章}

\Emph{答复：说明这个问题应保留到将来的会谈中。}\par

\Person{泰奥纳斯}说：你们的热情，即渴望达到成全之路，不仅是一时的，而是完全而圆满地，促使我们不知疲倦地继续这场讨论。因为你们所焦虑询问的，不是外在的贞洁或外表的割损，而是那隐秘的，因为你们知道，圆满的成全并不在于这可见的肉身节制（这节制甚至外教人也可以因强制或伪善而达到），而在于那自愿且不可见的心灵纯洁，蒙福的宗徒如此描述：\BibleQuote{因为外表做犹太人的，并不是真犹太人；外表肉身的割损，也不是真割损；唯有内心的才是真犹太人，心的割损是因圣神而非因文字，其称赞不是从人而来，而是从天主而来，}\BibleRef{罗 2:28}唯独祂洞察人心的隐秘。但由于你们的心愿不能完全满足（因为剩下的短暂今夜不足以探究这个极其困难的问题），我认为最好暂时搁置。因为这些问题，既应由我们以安静且完全摆脱一切纷扰思虑的心来提出，也应被你们接受到心中；正如为共同纯洁的缘故而应从事探讨，同样，一个缺乏正直恩赐的人无法学习或获得它们。因为我们所寻求的，不是虚浮言辞的论证，而是良心的内在信德和真理的更大力量所能说服的。因此，关于这净化之道，除了有经验的人，无人能提出什么；也唯有那最热切、最真诚喜爱真理本身的人，才可交付给他；他并不期望仅凭空言提问而获得，而是全力以赴，不好无益的空谈，而渴望内在的洁净。\par

\SourceRecord{Translated by C.S. Gibson.；Nicene and Post-Nicene Fathers, Second Series；Vol. 11.；Edited by Philip Schaff and Henry Wace.；Buffalo, NY: Christian Literature Publishing Co.,；1894.；<http://www.newadvent.org/fathers/350821.htm>.}

% structure: p0001 source=h1 target=section reason=page-title

\section{第二十三次会谈}

\Emph{圣若望·卡西安}\par

阿爸\Person{特奥纳斯}的第三次会谈。论\OriginalTerm{无罪}{Sinlessness}。\par

% structure: p0004 source=h2 target=subsection reason=relative-heading-rank; base=h2

\subsection{第一章}

\Emph{阿爸\Person{特奥纳斯}关于宗徒话语的论述：\Quote{因为我不做我所愿意的善事。}}\par

因此，当曙光初现时，老人被我们急切的追问所迫，深入探讨宗徒的主题，他这样说道：至于你们试图用以证明宗徒\Person{保禄}不是以自己的身份，而是以罪人的身份所说的那些经文：\BibleQuote{因为我不做我所愿意的善，反倒做我所憎恨的恶}；或者这段：\BibleQuote{但若我做我不愿意的，就不再是我做的，而是住在我内的罪做的}；以及随后的话：\BibleQuote{因为我按内心的人，喜悦天主的法律；但我看到在我肢体中另有一条法律，反对我理智的法律，并将我掳去，隶属于那在我肢体中罪的法律}——这些经文恰恰清楚表明，它们根本不可能适用于罪人的身份，而所说的内容只能适用于那些成全的人，并且只适合那些追随宗徒善表之人的贞洁。否则，这些话怎能适用于罪人的身份：\BibleQuote{因为我不做我所愿意的善，反倒做我所憎恨的恶}？甚至这句：\BibleQuote{但若我做我不愿意的，就不再是我做的，而是住在我内的罪做的}？因为哪个罪人会不情愿地因奸淫和淫乱玷污自己？谁会违背自己的意愿对邻人设谋？谁会被不可避免的必要所驱使，以假见证压迫人，或以偷窃欺骗人，或贪图他人财物，或流人血？反而，如圣经所说：\BibleQuote{人自幼就倾向邪恶}。\BibleRef{创 8:21} 因为所有人都被对罪的爱和行其所好的欲望所燃烧，以至他们警惕地寻找作恶的机会，害怕自己太慢享受私欲，并以自己的羞耻和众多罪恶为荣，如宗徒在谴责中所说，\BibleRef{斐 3:19} 并从自己的混乱中寻求荣誉。关于他们，耶肋米亚先知也断言，他们不仅不是不情愿地，也不是轻松地，而是以辛苦的努力去犯下那丑恶的罪行，以至于他们劳力去行恶：\BibleQuote{他们劳力作恶}。\BibleRef{耶 9:5} 还有谁会认为这话适用于罪人：\BibleQuote{这样，我自己以理智服从天主的法律，但以肉身服从罪的法律}？因为显然他们既不理智也不肉身地事奉天主。或者，那些以肉身犯罪的人怎能以理智事奉天主？当肉身从心接受罪的煽动，而两种本性的创造者亲自宣告罪的源头和泉源从心流出，说：\BibleQuote{从心里发出来的是恶念、奸淫、淫乱、偷窃、假见证等}。\BibleRef{玛 15:19} 因此，清楚显示，这话绝不能用于罪人的身份，他们不仅不憎恨恶，反而喜爱恶，并且既不是以理智也不是以肉身事奉天主，以至于他们在以肉身犯罪之前，已在理智和思想中犯了罪，在行肉身的快乐之前，已在心灵和思想中被罪征服。\par

% structure: p0007 source=h2 target=subsection reason=relative-heading-rank; base=h2

\subsection{第二章}

\Emph{宗徒如何完成了许多善行。}\par

所以我们必须以说话者最深切的感受来揣度其含义和主旨，讨论真福宗徒所称为善的和他所比较为恶的，不是凭字面的意思，而是凭他同样的洞察力，并按照说话者的价值和良善来探讨其含义。因为只有当我们衡量了说话者的地位和品格，并且我们自己也被赋予同样的感受（不是凭言语，而是凭经验），才能按照那品格（所有思想和意见无疑都是按此产生和发表的）来理解天主所默感的话语。因此，让我们仔细考虑，那宗徒想做却做不到的善，主要是什么。因为我们知道，有许多善事，我们不能否认真福宗徒和所有像他一样的人或者生来就有，或者因恩宠获得。因为贞洁是善的，节制是可赞美的，明智是令人钦佩的，仁慈是慷慨的，清醒是谨慎的，克己是谦逊的，怜悯是良善的，正义是圣洁的：所有这些，我们毫不怀疑在宗徒\Person{保禄}和他的同伴身上完美而完全地存在，以至于他们更以德行的榜样而非言语来教导信仰。如果他们始终为所有教会的不断关怀和警醒的焦虑所消耗，那又如何？这怜悯是何等大的善，\BibleQuote{为那些跌倒的人燃烧，与软弱的人一同软弱}！\BibleRef{格后 11:29} 如果宗徒如此充满这些善，那么除非我们达到他说话时的心境，否则我们无法认识那宗徒所缺乏的完美之善究竟是什么。因此，所有我们说他拥有的那些德行，尽管它们像最灿烂珍贵的珍珠，但与福音中商人所寻求并想通过变卖所有来获得的那颗最美丽独特的珍珠相比，它们的价值就显得贫乏而微不足道，以至于如果毫不迟疑地舍弃它们，单单拥有一件善事就能使卖掉无数善事的人富有。\par

% structure: p0010 source=h2 target=subsection reason=relative-heading-rank; base=h2

\subsection{第三章}

\Emph{宗徒所见证的他不能行的善究竟是什么。}\par

那么，那件独一无二、远胜于那些伟大且无数的善事、以至于所有其他善事都被藐视和抛弃、唯独应获取的，是什么呢？无疑，那是真正更好的一份，主如此描述其伟大而持久的性质，那时玛利亚不顾待客之礼与殷勤，选择了它：\Quote{玛尔大，玛尔大，你为许多事操心忙碌，其实需要的只有一件，甚至唯一的一件。玛利亚选择了那更好的一份，是不能从她夺去的。}因此，默观，即默想天主，是那一件事，其价值超越我们一切义行的功德、一切对德行的追求。而我们在宗徒保禄身上看到的一切，不仅是善的、有益的，甚至是伟大而辉煌的。但例如，被认为有些用处的合金，当考虑到银子时就变得无价值；银子的价值与金子相比也消失；金子与宝石相比也被忽视；而一堆宝石无论多么灿烂，也比不上一颗珍珠的光辉。同样，一切圣德的功德，虽然不仅在今世生活中是善的、有用的，而且保证了永恒的恩赐，但若与默观天主的功德相比，则被视为微不足道，可以说，如同可卖之物。为以圣经的权威支持这一比喻，圣经不是普遍论及天主所创造的一切，说：\Quote{看，天主所造的一切都极好}；又说：\Quote{天主所造的一切，在其季节内都是好的}？这些事物，在现世被称为不单单是善的，而是强调\Quote{极好}（因为它们对我们现世生活确实便利，或为生活用途，或为身体疗治，或因某种未知的有用性，或者它们确实是\Quote{极好}，因为它们使我们\Quote{从受造的世界看见天主不可见的事，藉着所造之物得以理解，即祂永远的大能和天主性} \BibleRef{罗 1:20}，从这伟大而有序的世界构造中；并藉着其中万物之存在来默观它们），但没有一样会保持善的名称，如果它们被对照那未来的世界来看，在那里无需惧怕善事的任何变化或真福的失落。那世界的福乐如此描述：\Quote{月亮的光将如同太阳的光，太阳的光将七倍于七日的光} \BibleRef{依 30:26}。这些事物，虽值得观看、神奇而奇妙，若与信德所期盼的未来承诺相比，立刻显得虚幻，如达味所说：\Quote{它们都要像衣服一样旧，你要像更换外衣一样更换它们，它们就被更换。但祢却是一样的，祢的岁月不会穷尽。}因此，既然没有什么是本身恒久、不变、善的，唯独天主性除外，而每个受造物要获得永恒和不变的祝福，并非靠其本性，而是靠参与造物主及其恩宠，所以当与造物主相比时，它们无法保持其善的性质。\par

% structure: p0013 source=h2 target=subsection reason=relative-heading-rank; base=h2

\subsection{第四章}

\Emph{人的善与义，若与天主的善与义相比，就不是善。}\par

但如果我们还愿意藉着更清楚的证据来确立这意见的力量，难道不是这样吗？我们在福音中读到许多事物被称为善的，诸如好树、宝库、善人、忠信的仆人，因为祂说：\Quote{好树不能结坏果子；}以及：\Quote{善人从他心里所存的善宝库中发出善来；}以及：\Quote{好！忠信的仆人；}而且毫无疑问，这些中没有一个是本身为善的；但如果我们考虑到天主的良善，它们中没有一个会被称为善的，正如主说：\BibleQuote{除了天主一个外，没有善的。}\BibleRef{路 18:19}在祂眼中，就连宗徒们自己——他们因蒙召的崇高在很多方面超越了人类的良善——也被称为恶的，因为主这样对他们说：\BibleQuote{你们纵然不善，尚且知道把好的东西给你们的儿女，何况你们在天上的父，岂不更要把好的东西给求祂的人吗？}\BibleRef{玛 7:11}最后，正如我们的良善在至高者眼中变为邪恶，同样，我们的义德若与天主的义德相比，也如同月经布一般，正如依撒意亚先知所说：\BibleQuote{你们的全部义德都像月经布。}\BibleRef{依 64:6}并且为提出更明显的事，就连法律本身的必要规条——即据说是\Quote{藉天使经中保之手所颁布的}，而同一位宗徒也说：\Quote{这样看来，法律本是圣的，诫命也是圣的、公义的和善的}——当它们与福音的完美相比时，却被天主的圣言称为不是善的：因为祂说：\BibleQuote{我也给了他们不美的诫命，以及他们不能凭以生活的条例。}\BibleRef{则 20:25}宗徒也证实，法律的荣耀被新约的光辉所湮没，以至于他宣称与福音的光彩相比，它不算荣耀，说：\BibleQuote{那先前有过光荣的，因这更超越的光荣，才不算光荣了。}\BibleRef{格后 3:10}圣经从另一方也保持着这种比较，即在衡量罪人的功过时，与恶人相比，它使犯罪较轻的人成义，说：\Quote{索多玛比你更正义；}又说：\Quote{你的妹妹索多玛犯了什么罪？}以及\Quote{悖逆的以色列在背信的犹大面前显为正义。}因此，上面我所列举的一切德性的功绩，虽然本身是善的和宝贵的，但与默观的璀璨相比，就变得暗淡无光。因为它们极大地阻碍和迟延了那些专注于属地目标（即使是在善工上）的圣人们，使他们无法默观那崇高的善。\par

% structure: p0016 source=h2 target=subsection reason=relative-heading-rank; base=h2

\subsection{第五章}

\Emph{无人能持续专注于那至善。}\par

有谁在\Quote{从强暴者手中解救穷人，从剥削弱者手中解救穷困之人}时，在\Quote{折断恶人的牙关，从其齿间夺回猎物}时，能在实际干预的工作中，以平静的心默观天主威严的光荣呢？有谁在向穷人提供支援，或慈爱地接待前来投奔他的群众时，在当下因兄弟的急需而焦虑不安时，能默观至高福乐的广大无边，且在被现世生活的烦恼和牵挂所摇撼时，能以超脱尘世污秽的心，期待来世的状态呢？因此，圣达味在断言惟有这才是对人有益时，渴望恒常依附于天主，并说：\Quote{亲近天主于我有益，我将我的希望寄托于上主。}\BibleBook{}{训道篇}也宣称，没有哪位圣人能做到这一点而无过，说：\BibleQuote{世上没有行善而不犯罪的义人。}\BibleRef{训 7:21}因为，即使他是所有义人和圣人的首脑，我们岂能认为，他在被现世锁链束缚时，能如此获得这至高之善，以至于从不中断对天主的默观，或被认为即使短暂地被世俗思想从他——那独一为善者——那里拉开呢？有谁从不关心食物，不关心衣物或其他属肉体的事物，或在为接待兄弟、更换地点、建造小屋而焦虑时，从未渴望过人的帮助，也从未因缺乏和困苦的折磨而受到上主这样的斥责：\BibleQuote{不要为你们的生命忧虑吃什么，也不要为你们的身体忧虑穿什么。}\BibleRef{玛 6:23}此外，我们确信，即使宗徒保禄本人——他在所受的苦难次数上超过所有圣人的辛劳——也不可能实现这一点，正如他本人在\BibleBook{宗徒}{大事录}中向门徒们作证说：\Quote{你们自己知道，这双手供应了我的需要，也供应了与我同住之人的需要。}或者，他在\BibleBook{得撒洛尼}{前书}中作证说，他\Quote{昼夜辛苦劳碌地工作。}虽然为此，他的功德有极大的赏报预备着，然而他的心思，无论多么神圣崇高，有时也不免因专注于地上的劳作而从那天上的默观中被拉开。此外，当他看到自己被如此丰硕的实践果实所充满，同时在心中思考默想之善，仿佛在一端称量所有这些劳苦的利益，另一端称量神圣默观的喜悦，当他长期调整心中的平衡时，一边劳苦的巨大奖赏使他喜悦，另一边与基督合一及密不可分相伴的渴望促使他离开此世，最后，他在困惑中呼喊说：\BibleQuote{我不知道该选择什么。我夹在两者之间，渴望离世与基督同在，因为那样好得多；但留在肉身中，为你们更有必要。}\BibleRef{斐 1:22}虽然如此，他在许多方面宁愿将这卓越之善置于他所有宣讲的果实之上，但出于爱——没有爱，没有人能获得主——他顺服自己；为了那些迄今他以福音的乳液哺育的人，他不拒绝与基督分离，这对他自己不利，却对他人有益。因为他更是被那种超乎寻常的善意所驱使，为了兄弟的救恩，他愿意，如果可能，甚至承受\OriginalTerm{绝罚}{Anathema}这最后的不幸。他说：\BibleQuote{我切愿自己为我的兄弟——按肉身是我的亲属，即以色列人——的缘故，从基督受绝罚。}\BibleRef{罗 9:3}也就是说，我宁愿不仅忍受暂时的，甚至永久的惩罚，只要所有人，如果可能，都能享受与基督的共融：因为我确信，众人的救恩对基督和我而言，比我们自己的救恩更好。因此，为使宗徒能完美地获得这至高之善，即享受对天主的瞻仰并持续与基督结合，他愿意与这身体分离——此身体因其软弱，并被其脆弱性的诸多需要所阻碍，不得不与基督结合分离：因为心志被如此频繁的忧虑所困扰，被如此多样且恼人的麻烦所阻碍，不可能常常享受神圣的瞻仰。因为圣人的目标能有如此持久吗？目的能有如此崇高，以至于那狡猾的阴谋者有时不能破坏它吗？有谁如此深入旷野的隐僻处，并避开与所有人的交往，以至于从不因不必要的思想而跌倒，也不因观看事物或忙于地上的行动而偏离那真正独一的善——默观天主？有谁能保持如此的精神热忱，以至于有时不从对祈祷的专注中游走于飘忽的思绪，突然从天上坠入地下？我们之中有谁（且不提其他飘荡的时刻）甚至在他将灵魂高举向天主祈祷的那一刻，不会陷入某种迟钝，并甚至违心地在希望从那里得到罪的赦免的那件事上犯罪呢？我问，有谁如此警觉和警惕，以至于在向天主咏唱圣咏时，从不允许自己的心思游离于圣经的意义之外？有谁如此亲密且紧密地与天主结合，以至于能为自己在一天内实现了宗徒的命令——他命我们不断祈祷——而庆幸？\BibleRef{得前 5:17}虽然这一切，在那些陷于更严重罪的人看来，可能微不足道，且与罪全然无关，但对于那些知道成全价值的人而言，即使极小事情的累积也变得极其严重。\par

% structure: p0019 source=h2 target=subsection reason=relative-heading-rank; base=h2

\subsection{第六章}

\Emph{那些自以为无罪的人，如同半盲之人。}\par

譬如我们设想有两个人，一人视力敏锐、目力完美，另一人则因视力模糊而目光昏暗，二人一同进入一所堆满许多包裹、器具和器皿的大房子。视力昏暗者因看不见一切，便会断言那里只有箱子、床、长凳、桌子，以及一切触及他摸索的手指而非看见的眼睛之物；而另一人却以清晰明亮的眼睛探查隐藏之物，宣称那里有许多极其细微、几乎难以计数的物件；若将这些物件收聚成一堆，其数量将等于或超过那视力昏暗者所触摸的少数物件。同样，圣人——容我这样说，那些有\Emph{看见}能力的人——他们的目标是最极致的成全，他们敏锐地在自身中察觉那些我们因心智之眼仿佛昏暗而看不见的事物，并极其严厉地谴责它们，以至于那些在我们看来因疏忽而未以丝毫罪污玷污其身体洁白如雪的人，在他们自己眼中却布满许多污点——我且不说任何邪恶或虚妄的思想潜入他们心灵的门口，就连必须咏唱的圣咏的回忆，在祈祷时刻也会分散跪祷者的注意力。因为，他们说，当我们向某位大人物恳求——我不说是为我们的生命和救恩，而是为某些利益和好处——我们将所有身心的注意力都集中在他身上，以战栗的期待悬于他的点头示意，怀着不小的恐惧，唯恐某个愚昧或不恰当的话语会打消听者的怜悯；再者，当我们站在尘世审判者的法庭或公堂上，对手站在对面时，若在控诉和审理过程中咳嗽、吐痰、发笑、打哈欠或瞌睡临到我们，那时刻警惕的对手会何等恶意地激起审判者的严厉来损害我们：那么，当我们恳求那知晓一切隐秘的天主时，因着永死迫在眉睫的危险，我们更当以恳切而忧虑的祈祷祈求审判者的仁慈，尤其是另一边站着那既是狡猾的诱惑者又是控告者的人！谁若在向天主倾泻祈祷时，立刻离开他的视线，仿佛离开一位既看不见也听不到者的目光，并追随邪恶思想的虚妄，那他必非因小罪而被定罪，而是因邪恶的重罪。但那些以自己罪恶的厚幔遮掩心灵之眼的人，正如救主所说：\BibleQuote{他们看却看不见，听却听不见，也不明白}\BibleRef{玛 13:13}，他们即使在内心的最深隐秘处，也几乎察觉不到那些重大致命的过错，无法以清澈的目光注视任何欺骗性的思想，甚至那些以轻微而微妙的暗示击中心灵的模糊隐秘欲望，以及灵魂的俘虏；他们总在污秽的思想中游荡，不知在从那特殊的默观中分心时如何悲伤，也无法因失去什么而哀痛，因为他们随心所欲地向任何思想的进入敞开自己的心灵，没有明确的东西呈在他们面前作为主旨或理应渴慕的目标。\par

% structure: p0022 source=h2 target=subsection reason=relative-heading-rank; base=h2

\subsection{第七章}

\Emph{持守人可无罪之说者身负双重错误。}\par

然而，驱使我们陷入这一错误的原因是：由于我们完全不知无罪之德，便幻想自己不会因思想中那些空虚随意的游荡而沾染任何罪责；反因迟钝而愚笨，仿佛被盲症击中，我们只能看到自身中那些大罪，以为只需避开那些也受世俗法律严厉谴责之事；若发现自己甚至短时间内免于这些，便立刻想象自己毫无罪过。因此，我们与那些有\Emph{看见}能力的人有别，因为我们看不到自身中聚集的许多小污点，也不受有益痛悔的打击——若烦恼的病症侵袭我们的思想，我们不为所动；我们不因受虚荣建议的打击而悲伤；不为献上如此迟缓和冷漠的祈祷而哭泣；不视之为过错——当我们咏唱或祈祷时，除了祈祷或圣咏本身之外，有别的思想充满我们的心思；我们也不因在心中构思许多我们羞于在人前说出或做出的事而惊骇——明知那颗心暴露在天主的目光之下；我们也不以大量泪水的洗涤来净化污秽之梦的玷污；也不因在施舍的虔诚行为中——当我们援助弟兄的需要或供应穷人的支持时，因吝啬的延迟而使喜乐的容光暗淡而忧伤；也不认为当我们忘记天主、思念暂时而腐朽之事时，自己蒙受了任何损失。因此，所罗门的这些话恰好适用于我们：\BibleQuote{他们打我，我却不忧伤；他们戏弄我，我却不察觉。}\BibleRef{箴 23:35}\par

% structure: p0025 source=h2 target=subsection reason=relative-heading-rank; base=h2

\subsection{第八章}

\Emph{只有少数人蒙受理解何为罪的恩赐。}\par

另一方面，那些将一切喜乐、欢愉和福乐的总和完全归于对神圣和属神之事的默观的人，如果因强行闯入的思想而被迫暂时离开它们，便会惩罚这如同一种亵渎，并立即以惩罚来报复；他们因将某种无价值的受造物（他们的心思转向了它）置于造物主之上而感到悲伤，便归咎自己（我几乎要说）不虔诚的罪责；尽管他们以极快的速度转动心灵的眼睛注视天主光荣的光辉，却仍无法容忍片刻血肉思想的黑暗，并诅咒任何阻挡他们灵魂视线离开真光的事物。最后，当蒙福的宗徒若望要将这情感灌输给众人时，他说：\BibleQuote{孩子们，不要爱世界，也不要爱世界上的事。谁若爱世界，天父的爱就不在他内；因为世界上的一切，肉身的贪欲、眼目的贪欲和生活的骄矜，都不是出于父，而是出于世界。世界和它的贪欲都要逝去，但那履行天主旨意的人，却永远存留。}\BibleRef{若一 2:15}因此，圣人们藐视世界所依凭的一切，但他们不可能不被思想的短暂偏差牵引到它们那里；就连如今，除了我们的主和救主以外，没有人能使自己本会游荡的心思永远专注于默观天主，以至于从不因爱慕这个世界上的任何事物而被牵引离开；正如圣经所说：\BibleQuote{连星辰在他眼前也不洁净}，又说：\BibleQuote{他连自己的圣者也不信任，在天使身上也找到错谬}，或按照更准确的译文：\BibleQuote{看，在他的圣者中，没有一位是不变的，诸天在他眼前也不纯洁}。\par

% structure: p0028 source=h2 target=subsection reason=relative-heading-rank; base=h2

\subsection{第九章}

\Emph{论隐修士应当如何谨慎地保持对天主的记念。}\par

那么，我要说，那些紧守对天主的记念，仿佛脚步悬于高处伸展的绳索上被带往前行的圣人们，正可被比作走绳者（通常称为\OriginalTerm{走绳索者}{funambuli}），他们在这条极窄的绳索上冒险，将全部安全和生命置于险境，毫不怀疑只要脚稍有偏离或绊倒，或越过那唯一安全的路线，就会立即遭遇最可怕的死亡。他们以惊人的技巧在空中施展轻盈的步伐，如果他们没有以细心焦虑的调节将脚步保持在那唯一安全的极窄小径上，那么对于所有人而言本是自然基座、最坚固最安全的基础的大地，便会立即成为他们明显的危险，不是因为其本性改变，而是因为他们因自身重量而头朝下坠落其上。同样，天主不懈的良善和祂不变的本性确实不伤害任何人，但我们自己从高处坠落，趋向深渊，成了自己死亡的制造者，或者更确切地说，坠落本身成了坠落者的死亡。因为经上说：\BibleQuote{祸哉，因为他们离开了我；他们必被荒废，因为他们背叛了我。}又说：\BibleQuote{祸哉，当我离开他们的时候。}因为\BibleQuote{你自己的邪恶要指证你，你的背信要责备你。要知道并看见，你离开上主你的天主是多么邪恶和痛苦。}因为\BibleQuote{每个人都被自己罪恶的绳索束缚。}主恰当地将这责备指向他们：\BibleQuote{看，}祂说，\BibleQuote{所有点火、被火焰包围的人，行走在你们的火把和所点燃的火焰的光中。}又说：\BibleQuote{点火行恶的人，必被其毁灭。}\par

% structure: p0031 source=h2 target=subsection reason=relative-heading-rank; base=h2

\subsection{第十章}

\Emph{那些走向成全的人如何真正谦卑，并感到自己一直需要天主的恩宠。}\par

因此，当圣人们感到被世俗思想的重量压迫，从心灵的高远坠落，并违背己意（不如说不知不觉）被牵入罪恶和死亡的法律，并且（姑且不提其他）被我上面描述的那些虽属世俗却良善正当的行为拦阻，无法默观天主时，他们便有不断向主叹息的事，有真正谦卑的事，并在忏悔中不仅以言语，更以内心承认自己是罪人；为此，当他们不断因肉身的软弱而被击败，每日所犯的过犯祈求主的恩宠宽恕时，他们该不断流出真诚的忏悔之泪；因为他们看见，即便直到生命终结，仍被相同的困扰所缠绕，并因这些困扰被磨炼而不断悲伤，甚至无法在不受烦扰思想的情况下献上祈祷。于是，他们凭经验知道，因肉身的负担阻碍，无法凭人力达到所愿的终点，也不能按心愿与那至高至善者结合，反而被俘虏到属世的事物，远离对它的默观，他们便投奔天主的恩宠，\BibleQuote{祂使不虔敬的人成义}\BibleRef{罗 4:5}，并与宗徒一同呼喊：\BibleQuote{我这个人真不幸啊！谁能救我脱离这死亡的身体？感谢天主，借着我们的主耶稣基督。}\BibleRef{罗 7:24}因为他们感到无法行所愿的善，却总是在行所不愿并憎恨的恶，即游荡的思想和属血肉的挂虑。\par

% structure: p0034 source=h2 target=subsection reason=relative-heading-rank; base=h2

\subsection{第十一章}

\Emph{对\Quote{就内心而言，我喜悦天主的法律}等语句的解释。}\par

而且，他们的确\Quote{喜爱}\Quote{按内在的人，喜爱天主的法律}，它超乎一切有形之物，常力求唯独与天主结合；但他们\Quote{在肢体中见到另一条法律}，即植入其自然人性状况中的法律，它\Quote{抗拒理智的法律}，将他们的思想俘虏到罪的强暴法律之下，迫使他们离弃那至善，而屈从于世俗的观念；这些观念虽然当为某种宗教需要而采纳时，可能显得必要且有用，但一旦与那令众圣人目眩神迷的至善相比，便被他们视为恶而应当避免，因为藉着它们，他们以某种方式并在短时间内从那完美福乐的喜悦中被拉走。因为罪的法律实在就是其始祖的堕落藉着他的过犯带给人类的，对此最公正的审判者宣布了这样的判决：\Quote{因你的作为，土地受诅咒；它要给你生出荆棘和蒺藜，你必汗流满面，才有饭吃。}我说，这就是法律，植入所有凡人的肢体中，它抗拒我们理智的法律，阻止我们看见天主；并且，正如在认识善恶之后，土地因我们的作为而受诅咒，它就开始产生思想的荆棘和蒺藜，其尖锐的刺痛扼杀了德性的自然种子，以致没有额上的汗水，我们就不能吃那\Quote{从天降下}并\Quote{坚固人心}的食粮。因此，整个人类普遍地、毫无例外地服从这一法律。因为无论多么圣洁的人，没有谁不是在额上的汗水和内心的焦虑努力中吃上述食粮的。但许多富人，如我们所见的，却不用额上流汗就能吃那普通食粮。\par

% structure: p0037 source=h2 target=subsection reason=relative-heading-rank; base=h2

\subsection{第十二章}

\Emph{论此亦：\Quote{但我们知道法律是属神的，}等。}\par

宗徒也称呼这一法律为属神的，说：\BibleQuote{但我们知道法律是属神的，但我却是属血肉的，被卖于罪恶之下。}\BibleRef{罗 7:14} 因为这一法律是属神的，它命令我们汗流满面地吃那\BibleQuote{从天上降下的真食粮}\BibleRef{若 6:33}。我问，这是什么罪？或谁的罪？无疑是亚当的罪，因他的堕落——我似乎可以说，因那毁灭性的交易和欺诈的买卖——我们被卖了。因为当他被蛇的劝说所迷惑时，他使所有后代都陷入了永久奴役的轭下，因吃了禁果而疏远。因为买卖双方通常遵守这样的惯例：一个人若想把自己交给别人管辖，就从买主那里得到一笔丧失自由的代价，并被判为永久奴隶。我们很清楚地看到这在亚当和蛇之间发生了。因吃了禁树上的果子，他从蛇那里得到了自由的代价，放弃了自然自由，选择将自己交给那从那里得到禁果致命代价的人永久为奴；从此他被这个条件所束缚，并有理由地使所有后裔永久服事他已成了其奴仆的那一位。因为奴隶中的婚姻能产生什么除了奴隶？那么，那个狡猾奸诈的买主是否从真正合法的领主那里夺走了所有权？不是。因为既不是他通过一次欺骗的诡计战胜了天主的所有财产，以致真正的主人失去了所有权——尽管买主本人是个叛逆和背教者，却用奴役的轭压迫他——而是因为造物主赐予一切理性受造物自由意志，祂不愿违背那些因贪婪情欲之罪不义地出卖自己之人的意愿，恢复他们自然的自由。因为凡违反良善和公平的事，都为公义和虔诚的创造主所憎恶。因为祂撤销已赐予的自由祝福是不妥当的，用祂的能力压制自由的人，抓住他而不让他行使所接受自由的特权，是不公平的；因为祂为后世保留祂的救恩，好使指定时期的圆满在适时得以完成。因为他的后裔应在古老条件下如此长久地存留，直到藉着祂自己血的代价，主的恩宠将他们从原始枷锁中赎回，并在原始自由状态中释放他们；尽管当时祂也能拯救他们，但祂不愿意，因为公平禁止祂违犯自己法令的条款。你想知道你被卖的原因吗？请听你的救赎主亲自通过依撒意亚先知公开宣告：\Quote{你们母亲休书的字据在哪里，我凭什么休了她？还是我的债主是谁，我将你们卖给了他？看，你们是因你们的不义被卖，因你们的邪恶行为我休弃了你们的母亲。}你想明白为什么当你被交给奴役的轭时，祂不愿用祂自己的力量救赎你吗？请听祂在前面的话之后加上的内容，以及祂如何同样指控这些罪的奴仆自愿出卖的理由。\BibleQuote{我的手岂是缩短而不能救赎？或我无力拯救？}\BibleRef{依 50:1} 但总是阻碍祂全能怜悯的东西，同一先知表明：\BibleQuote{看，主的手并非缩短不能拯救，耳朵并非沉重不能听见；但你们的罪过使你们与你们的天主隔绝，你们的罪恶使祂掩面不垂听你们。}\BibleRef{依 59:1}\par

% structure: p0040 source=h2 target=subsection reason=relative-heading-rank; base=h2

\subsection{第十三章}

\Emph{论此亦：\Quote{但我知道在我内，即在我肉身上，没有良善。}}\par

因为那时天主的原始诅咒已使我们成了属血肉的，且判定我们遭受荆棘和蒺藜；我们的父亲通过那不幸的交易将我们出卖，以致我们不能行我们所愿的善，因为我们被从对至高天主的记念中夺去，被迫思想属于人性软弱的事；尽管我们燃烧着对贞洁的热爱，却常常甚至违背自己的意愿，被自然的欲望所困扰，而这些欲望我们宁愿一无所知；我们知道，在我们肉身之内，毫无良善\BibleRef{罗 7:18}，即我们所说的这种默想的持久和平；但在我们身上却发生了那种可怜可悲的离异：当我们立意用心灵服事天主的法律时，由于我们总不愿将目光从神圣的光辉上移开，却被肉性的黑暗所包围，被迫因某种罪的法律放弃我们所知的善，从心灵的高处坠落到世俗的忧虑和思念中。罪的法律，即天主对第一个犯法者所宣告的判决，并非无根据地判定了我们。因此，蒙福的宗徒虽然公开承认他和诸圣都被这罪的约束所束缚，却大胆断言他们中无人会因此被定罪，说：\BibleQuote{所以现在，凡在基督耶稣内的，已不被定罪了：因为生命之圣神的法律在基督耶稣内已释放了我，脱离了罪与死的法律，}\BibleRef{罗 8:1} 即基督的恩宠日日释放祂的诸圣脱离这罪与死的法律；他们在这种法律下，每当为主祈求赦免过犯时，都不得不勉强屈服。可见，蒙福的宗徒是以真正圣人、完人（而非罪人）的身份说了这话：\BibleQuote{因为我不行我所愿的善，反行我所憎恨的恶；}以及：\BibleQuote{我发觉在我肢体内另有一条法律，反对我心灵的法律，掳掠我，使我隶属于那在我肢体内作祟的罪的法律。}\BibleRef{罗 7:19}\par

% structure: p0043 source=h2 target=subsection reason=relative-heading-rank; base=h2

\subsection{第十四章}

\Emph{异议：\BibleQuote{因为我不行我所愿的善}等语，所指的既非不信者，亦非圣人。}\par

\Person{格尔玛诺}：我们说，这并不适用于那些身犯大罪的人，也不适用于宗徒及已达到他尺度的人；我们以为它应恰当地适用于那些在领受天主的恩宠及认识真理之后，本应戒避肉性之罪，但却因古代习俗如同自然法律在肢体中强有力地统治着他们，而被牵引到根深蒂固的情欲中的人。因为犯罪的习惯与频繁，已变得如同自然法律一般，植入人软弱的肢体中，引导那在一切德行追求上尚未受教的灵魂——仍可说是一种未受训的、娇嫩的贞洁——被罪掳掠，并以古老法律使其臣服于死亡，将它们置于统治它们的罪的轭下，不容许它们获得所爱的清洁之善，反而强制它们行所憎恨的恶。\par

% structure: p0046 source=h2 target=subsection reason=relative-heading-rank; base=h2

\subsection{第十五章}

\Emph{对上述异议的答复。}\par

\Person{德奥纳}：尔之意念微不足道；因尔等现今已开始主张，此不可能适用于彻底为罪人者，而正当地适用于那些努力保持自己远离肉身之罪者。既然尔等已将这些人与罪人分离，则尔等不久必得承认他们进入信众及圣者之列。盖尔等言，他们能犯何种罪，以至在领受圣洗恩宠后若陷于此等罪，能借基督每日之恩宠获得释放？或者，宗徒所言\BibleQuote{谁能救我脱离这取死之身？感谢天主，藉着我们的主耶稣基督}\BibleRef{罗 7:24}，我们该认为这是指何种死亡之身？岂非显然，如真理迫使尔等也承认，此言并非论及那等大罪之肢体——即谋杀、奸淫、私通、酗酒、偷窃、抢劫，其工价乃是永死——而是论及前述之身，即基督每日之恩宠所助者？因为凡在领洗并认识天主之后陷入那死亡者，必须知道，他不是藉基督每日之恩宠——即一种轻易的宽恕，我主在随时受祈求时惯常赐予我们过错的——得洁净，而是藉终生的补赎苦行与忏悔之罚，否则将来必为这些罪受永火之罚，如同一宗徒所宣告：\BibleQuote{你们不要自欺：无论是放荡的、拜偶像的、奸淫的、变童的、好男色的、偷窃的、贪婪的、酗酒的、辱骂的、勒索的，都不能承继天主的国。}\BibleRef{格前 6:9}或者，那在我们肢体中交战、反抗我们理智之律，并将我们这些抗拒者掳去归服罪与死之律，使我们以肉身服事它，却仍容许我们以理智服事天主之律的，是什么律？因为我不认为这罪之律指大罪，或可理解为上述之过犯；人若犯此等过犯，便不以理智服事天主之律，因他在以肉身犯任何一项之前，必先在心中离弃了那律。因为服事罪之律，岂非就是行罪所命之事？那么，是何等样的罪，竟使如此伟大的圣德与成全感到自己被掳，却又确信必因基督之恩宠从中得释，说：\Quote{我真不幸啊！谁能救我脱离这取死之身？感谢天主，藉着我们的主耶稣基督}？请问，你会主张在我们肢体中植有何律，它使我们偏离天主之律、被掳归于罪之律，使我们成为可怜而非有罪，以致不归于永罚，却仍为未中断的福乐而叹息，并寻求一位助者恢复我们，与宗徒同呼：\Quote{我真不幸啊！谁能救我脱离这取死之身？}因为被掳归于罪之律，岂非就是继续行罪、犯罪？或者，除了那至善——如上所述，与之相比，其余皆非善——之外，还有什么其他至善是圣者们不能成就的？诚然，我们知道世上许多事物是善的，尤其是谦逊、节制、清醒、谦卑、正义、仁慈、温和、虔诚；但所有这些都达不到那至善，它们不仅能被宗徒，甚至能被普通人做到；而那些不行这些事的人，或受永罚惩戒，或如前述，藉补赎的巨大努力——而非基督每日之恩宠——得释放。那么，我们只得承认，宗徒此言正当地仅适用于圣者们：他们每日落入我们所述之罪律（非大罪之律）之下，对自己的得救有把握，不陷于恶行，但如常言，被从对天主的默观中拉向肉身思想的痛苦，且常被剥夺那真福之福。因为若他们觉得自己因肢体之律每日被绑于大罪，他们抱怨的就不会是福乐之丧失，而是纯洁之丧失；保禄宗徒就不会说\Quote{我真不幸啊}，而会说\Quote{我真不洁}或\Quote{我真邪恶}，并渴望解脱的就不是这取死之身（即此有死之境），而是这肉身的罪行与恶事。但因他出于人性软弱之境，感到自己被掳——即被引向肉身的挂虑与焦虑，即罪与死之律所导致的——他为此罪之律悲叹，他虽不愿却已陷入其中，并立即投靠基督，因祂恩宠的现世救赎而得救。因此，这罪之律自然在宗徒心田中生出的烦恼——即世俗思想与挂虑的荆棘与蒺藜——恩宠之律立时拔除。因为他说：\BibleQuote{在基督耶稣内生命之灵的法律，已释放我脱离了罪与死的法律。}\BibleRef{罗 8:2}\par

% structure: p0049 source=h2 target=subsection reason=relative-heading-rank; base=h2

\subsection{第十六章}

\Emph{何为罪之身体。}\par

这，就是那死亡之身，我们无法逃脱；被囚其中，那些成全的人，曾尝过\Quote{主是何等甘饴}，每天与先知一同感受\BibleQuote{人离弃主——他的天主，是何等自招祸患与痛苦}\BibleRef{耶 2:19}。这就是那死亡之身，它阻止我们瞻仰天上异象，将我们拖回属地之事；使人在咏唱圣咏、跪地祈祷时，心思充满人形、交谈、事务或不必要的举动。这就是那死亡之身，因着它，那些本想效法天使的圣德、长久渴望依附于天主的人，仍不能达到这善的成全，因为死亡之身挡在路上，他们做那不愿意的恶，即他们的心神被拖到与他们在德行上的进步与成全毫无关系的事物上。最后，蒙福的宗徒为清楚表明他是就圣善成全的人、以及像他自己这样的人说这话，就用某种方式用手指指着自己，立刻接着说：\Quote{这样我自己}，即说这话的我，揭露的是我自己的良心秘密，而非他人的。无论如何，宗徒惯于使用这种表达方式，每当他特别指自己时，就如这里：\Quote{我\Person{保禄}自己以基督的温良和谦虚恳求你们}；又说：\Quote{除了我自己没有加重你们负担}；又说：\Quote{就这样吧：我自己没有加重你们负担}；别处：\Quote{我\Person{保禄}自己告诉你们：如果你受割损，基督对你们就毫无益处}；致\WorkTitle{罗马书}：\Quote{为我弟兄，我宁愿自己因基督而被诅咒}。但这样理解亦非无理：\Quote{这样我自己}是明确带着强调说的，即你们知道我是基督的宗徒，你们对之极表尊敬，相信我是最高尚完美的人，基督在我内说话，但我用理智侍奉天主的法律，却用肉身承认我侍奉罪的法律，即因我人类境况的挂虑，有时从天上被拖到属地之事，我心灵的崇高被降到对卑微事物的忧虑水平。并且因这罪的法律，我发现每时每刻我被如此俘虏，虽然我以坚定不渝的渴望持守天主的法律，却绝无法逃脱这囚禁的权势，除非我总是投奔救主的恩宠。\par

% structure: p0052 source=h2 target=subsection reason=relative-heading-rank; base=h2

\subsection{第十七章}

\Emph{所有圣人如何真实地承认自己是不洁且有罪的。}\par

因此，所有圣人日日叹息，为他们本性的软弱而悲伤，并省察自己流变的思想、秘密和良心的最深隐秘处时，恳求道：\BibleQuote{求祢不要与祢的仆人进入审判，因为在祢面前，凡活着的，没有一人得以成义；}以及：\BibleQuote{谁能自夸有纯洁的心？谁又能自信自己无罪？}并且：\BibleQuote{世上没有一个行善而不犯罪的义人；}还有：\BibleQuote{谁能认清自己的过犯？}因此，他们认识到人的义德是软弱、不完善的，常需要天主的仁慈，以致于那位天主用从祭坛上取来的祂话语的活炭除去了其不义和罪恶的人，在经历了那奇妙的天主异象，看到了高天的色辣芬，并获启示天上奥迹之后，说道：\BibleQuote{祸哉我！因为我是一个唇舌不洁的人，住在唇舌不洁的人民中间。}\BibleRef{依 6:5}我想，也许此前他还不曾感到自己嘴唇的不洁，除非在天主显现中得以认识那真正而完美的纯洁；一见天主，他立刻觉察到自己先前所不知的不洁。因为当他说：\BibleQuote{祸哉我！因为我是一个唇舌不洁的人，}是表明他随后告明的涉及自己的嘴唇，而非人民的不洁：\BibleQuote{我住在唇舌不洁的人民中间。}但即使他在祈祷中承认所有罪人的不洁，他的普遍恳求不仅包含恶人，也包含善人，说道：\BibleQuote{看，祢发怒，我们犯罪：在它们中我们一直存在，我们仍将得救。我们都变成不洁，我们的义德都像污秽的衣服。}\BibleRef{依 64:5}请问，有什么比这话更清楚的呢？先知在其中不仅包括某一种，而是我们所有的义德，环顾所有被视为不洁和可憎的事物，因他在人类生活中找不到比污秽衣服更龌龊或不洁的东西，遂选择以之相比。那么，针对这极其清楚的真理而提出的纠缠异议的锋锐是徒然的，正如你不久前所说：\Quote{如果没有人无罪，就没有人是圣的；如果没有人是圣的，就没有人会得救。}因为这一问题的难题可通过先知的证言解决。他说：\BibleQuote{看，}他说，\BibleQuote{祢发怒，我们犯罪，}即当祢拒绝我们心中的骄傲或疏忽，并剥夺祢的助佑时，我们罪恶的深渊立刻吞噬了我们，就如一个人对灿烂的太阳说：看，你落下，于是幽暗的黑暗立刻笼罩了我们。然而，尽管他在这里说圣人们犯了罪，而且不仅犯了罪，还长久停留在他们的罪中，他并未完全绝望于救恩，而是加上：\BibleQuote{在它们中我们一直存在，我们仍将得救。}我要将\BibleQuote{看，祢发怒，我们犯罪}这句话与宗徒的那句话相比：\BibleQuote{我真不幸！谁能救我脱离这死亡之身？}同样，先知接着说的：\BibleQuote{在它们中我们一直存在，我们仍将得救，}与宗徒接下来的话相应：\BibleQuote{感谢天主，借着我们的主耶稣基督。}同样，同一位先知的这段经文：\BibleQuote{祸哉我！因为我是一个唇舌不洁的人，住在唇舌不洁的人民中间，}似乎与上面所引的话一致：\BibleQuote{我真不幸！谁能救我脱离这死亡之身？}先知随后说到：\BibleQuote{看，有一个色辣芬飞到我面前，手中拿着用钳子从祭坛上取下的红炭。他触摸我的口说：看，我以这炭沾了你的口唇，你的不义已除，你的罪已赎。}\BibleRef{依 6:6}这正像是出自保禄口中的话，他说：\BibleQuote{感谢天主，借着我们的主耶稣基督。}由此你看到，所有圣人如何真实地——与其说是代表人民，不如说是代表自己——承认自己是罪人，却绝不绝望于自己的救恩，而是期望从主恩宠和仁慈中获得完全的成义（他们并不认为因人性软弱的状态就不能获得成义）。\par

% structure: p0055 source=h2 target=subsection reason=relative-heading-rank; base=h2

\subsection{第十八章}

\Emph{连良善圣洁的人也不是无罪的。}\par

但是，无论何人，即便圣洁，在此生中绝不会免于过犯和罪恶，这也是救主的教导晓谕我们的。祂将完美祈祷的范式赐给祂的门徒，并在那些崇高神圣的命令——因其仅赐予圣者和成全者，故不适用于恶人和不信者——中，吩咐加入这一句：\BibleQuote{求你宽免我们的罪债，如同我们也宽免得罪我们的人。} \BibleRef{玛 6:12} 倘若这确实是圣者所献的真诚祈祷——我们理当确信无疑——那么，谁会如此顽梗放肆、被魔鬼般的狂傲冲昏头脑，竟声称自己无罪，不仅自以为大于宗徒，甚至还指责救主无知或愚昧，好像祂不知道有些人可以免于罪债，或是徒然教导那些祂明知无需这祈祷之补救的人？然而，凡是完全遵守君王诫命的圣者，每日都说：\BibleQuote{求你宽免我们的罪债}；若他们说真话，则确实无人无罪；若他们说谎，则同样真实的是，他们也未脱离谎言之罪。因此，那至为智慧的\BibleBook{}{训道篇}作者在脑海中省察人类一切行为与意向后，毫无例外地宣告：\BibleQuote{世上没有一个义人，只行善而不犯罪。} \BibleRef{训 7:21} 也就是说，在这世上，从未有人、也永不会有人如此圣洁、如此勤勉、如此热忱，以致能持续持守那真实而唯一的善，而不日复一日地感到自己偏离了它、跌倒了。但是，尽管他坚持自己不能免于过失，我们仍不可否认他是义人。\par

% structure: p0058 source=h2 target=subsection reason=relative-heading-rank; base=h2

\subsection{第十九章}

\Emph{即使祈祷时刻，也几乎无法避免罪恶。}\par

那么，凡是将无罪性归于人性者，必须与之对抗的不是空言，而是我们自身良心的见证与证明；只有当一个人发现自己未曾从这至高的善中被撕离时，才应坚持自己无罪。不，更确切地说，凡审视自己良心——姑且不提别的——发现自己即使举行过一次祈祷，也未因一言、一行或一念而分心者，可以说自己无罪。再者，因为我们承认，人心那游移不定的轻浮无法摆脱这些空虚无用之物，所以我们就随之真诚地宣认我们并非无罪。因为，无论一个人如何谨慎地守护自己的心，由于肉性本性的抗拒，他永远无法按自己精神的意愿来守护它。因为，无论人心在默观的更精细洁净方面有了多大进步和进展，它就越会看到自己的不洁，如同在它的洁净之镜中；因为当灵魂为更崇高的景象而提升，向外瞻望时，它渴慕比自己所能做到的更伟大的事物，就必然会贬低自己所涉入的那些事物为低劣和无价值。因为更敏锐的眼力会注意到更多；而无过犯的生活在被指摘时反而产生更大的悲伤；生活的改善、对善的热切追求会倍增哀叹和叹息。因为无人能安于自己所达到的阶段；无论一个人在思想上多么被净化，他就越会看到自己的污秽，找到谦卑而非骄傲的理由；无论他多么迅速地攀登到更高处，他就越会看到在自己之上他正趋向的所在。最后，那位被选的宗徒，\BibleQuote{耶稣所爱的}，\BibleRef{若 13:23} 正靠在祂的胸膛上，说出了这句仿佛出自主内心的话：\BibleQuote{如果我们说自己没有罪，就是欺骗自己，真理也不在我们内。} \BibleRef{若一 1:8} 因此，若我们说没有罪，就没有真理——即基督——在我们内，那么除了藉此宣认证明自己是罪人中的罪魁和恶人，我们还能做什么呢？\par

% structure: p0061 source=h2 target=subsection reason=relative-heading-rank; base=h2

\subsection{第二十章}

\Emph{从谁那里我们可以学到罪的毁灭和善的成全。}\par

最后，如果你愿意更深入地考察人性是否可能达到无罪性，除了从那些\Quote{已将肉体连同它的毛病与情欲钉在十字架上}的人，以及那些\Quote{世界已真正被钉在十字架上}的人以外，我们还能从谁那里更清楚地学到呢？这些人不仅完全根除了心中一切毛病，而且还试图连对罪的念想和回忆都摒弃，却仍日复一日忠实地坚持，他们连一小时都不能免于罪的玷污。\par

% structure: p0064 source=h2 target=subsection reason=relative-heading-rank; base=h2

\subsection{第二十一章}

\Emph{虽然我们承认不能无罪，但仍不应疏离主的共融。}\par

然而，我们不应因承认自己是罪人而远离主的圣体，反而应更加热切地奔赴它，以医治自己的灵魂、净化自己的心神，并以谦卑的心和信德寻求伤口的药方，因为我们自认不配领受如此伟大的恩宠。否则，我们甚至无法一年一次相称地领受圣体，如同某些在隐修院中生活的人所行的那样，他们如此看重天上圣事的尊贵、神圣和价值，以至于认为只有圣人和无玷者才敢领受它们，而不认为领受它们会使我们成为圣洁纯洁的人。由此，他们陷入了比他们自以为避免的更大的狂妄和傲慢，因为在他们领受之时，他们自以为配得领受。但远更好的做法是，每周主日都带着内心的谦卑领受它们以医治我们的软弱，由此我们相信并承认自己永远无法相称地触摸那些神圣的奥迹，而不是被内心的愚蠢信念所膨胀，相信在年终时我们配得领受。因此，为使我们能掌握并富有成果地持守这一点，让我们更热切地恳求主的仁慈帮助我们实行此事；此事不像其他人类技艺那样通过事先的口头解说习得，而是通过经验和行动引导；而且，除非经常在属灵之人的座谈会中思考并推敲，并通过日常的经验和考验焦虑地筛析，否则它就会因疏忽而变得过时，或因懈怠的遗忘而消亡。\par

\SourceRecord{Translated by C.S. Gibson.；Nicene and Post-Nicene Fathers, Second Series；Vol. 11.；Edited by Philip Schaff and Henry Wace.；Buffalo, NY: Christian Literature Publishing Co.,；1894.；<http://www.newadvent.org/fathers/350823.htm>.}

% structure: p0001 source=h1 target=section reason=page-title

\section{第二十四次会议}

\Signature{由\Person{圣若望·卡西安}著}\par

\Person{亚巴郎}院长的会议。论\OriginalTerm{克己}{mortification}。\par

% structure: p0004 source=h2 target=subsection reason=relative-heading-rank; base=h2

\subsection{第一章}

我们如何向亚巴郎院长袒露心中的秘密。\par

承蒙基督的恩宠，现在呈现亚巴郎院长的第二十四次会议，它总结了所有长者的传统与决定；并凭借你们祈祷的援助得以完成，因其数目神秘地对应了神圣默示录中所述向羔羊献上冠冕的二十四位长老\BibleRef{默 4:4}，我们想我们将偿还所有许诺的债务。从此以后，如果我们这二十四位长者因教导而获得任何荣耀，他们将俯首将荣耀献给那为世界的救恩而被宰杀的羔羊：因为正是祂为了祂圣名的尊荣，赐予他们如此崇高的思想，并赐予我们必要的言辞以阐述如此深奥的道理。祂恩赐的功绩必须归于万善之源，所欠越多，所偿也越多。因此，我们怀着焦虑的坦白，向这位亚巴郎袒露了我们思想的冲动，这些冲动因每日内心的困惑而驱使我们返回故乡、探望亲友。我们渴望的最大缘由在于，我们记得亲友们具备如此虔诚与良善，深信他们绝不会妨碍我们的志向，并不断思量：我们应从他们的热忱中获得更多益处，不会受物质顾虑和食物供应的困扰，因为他们会欣然充足供应我们所需的一切；此外，我们还以空幻喜乐的希望滋养灵魂，以为通过我们的榜样和训导，能使许多人悔改并走向救恩之路，从而获得最大的益处。除此之外，祖传产业的所在地和邻近地区的宜人景色也浮现在我们眼前：它多么宜人而适当地延伸至旷野，林中幽静不仅令隐修士心旷神怡，还提供丰富的食物。当我们坦率地按良心的信德将这一切向上述长者阐明，并以涌流的泪水表明，若天主的恩宠不藉他的医治来援救，我们再也无法抗拒这股冲动时，他长时沉默，最终深深叹息，说道：\par

% structure: p0007 source=h2 target=subsection reason=relative-heading-rank; base=h2

\subsection{第二章}

长者如何揭露我们的错误。\par

你们思想的软弱表明你们尚未弃绝世间的欲望，也未克治昔日的贪情。因为你们欲望的游移不定证明了你们心灵的懈怠，而你们本应以内心忍受的这次朝圣和离乡背井，却\Emph{确实}只是以肉身忍受。因为如果你们掌握了正确的弃绝方法，以及我们居住这独处的根本原因，所有这些本应被埋葬并彻底逐出你们的心。所以我看出你们正受那懒惰之病的折磨，正如《箴言》中所描述的：\Quote{懒惰人时常有所欲望；} 又说：\Quote{欲望杀害懒惰人。} 因为在我们身上，你们所描述的那些世俗便利的供应也不会缺乏，如果我们相信它们适合于我们的圣召，或认为我们能从那享乐和愉悦中获得与这荒野的贫瘠和身体的苦楚同样多的益处。我们也没有被剥夺亲人的安慰，以至于那些乐于以财物支持我们的人会缺位，若不是救主的那句话临到我们，排斥了一切有助于滋养肉身的事物，祂说：\BibleQuote{谁若不离开（或憎恨）父母、儿女和兄弟，就不能作我的门徒。} \BibleRef{路 14:26} 但若我们完全失去了父母的庇护，这世上的王侯们的服务也不会缺乏，他们会极其感恩地以宽厚的慷慨满足我们的急需。靠着他们的恩赐，我们本可免于预备食物的忧虑，若不是先知的那诅咒可怕地惊吓了我们。因为 \Quote{诅咒，} 他说，\Quote{那信赖世人的人；} 以及：\Quote{不要信赖王侯。} 我们也本可将我们的斗室安置在尼罗河畔，使水就在门旁，免得我们需将水扛在肩上走四英里，若不是蒙福的宗徒使我们孜孜不倦地忍受这劳苦，并用他的话鼓励我们，说：\BibleQuote{各人要照自己的劳苦，获得自己的酬报。} \BibleRef{格前 3:8} 我们也不是不知道，在我们的地方甚至还有宜人的幽静之处，那里有丰富的果实、悦目的花园和肥沃的土地，只需稍费体力就能提供我们所需的食物，若不是我们害怕那针对福音中富人的责备会临到我们：\BibleQuote{因为你已在今生得了你的安慰。} \BibleRef{路 16:25} 但我们轻视这一切，连同世上的一切享乐一同鄙视，我们只喜爱这荒芜之地，更胜于一切奢华，我们不将任何肥沃土壤的财富与这贫瘠的沙地相比，因为我们追求的不是肉身暂时的收益，而是永恒的灵的赏报。因为一个隐修士只做一次弃绝是不够的，即在他初皈依时无视今世，除非他每日持续弃绝。因为直到生命终结，我们都必须同先知这样说：\BibleQuote{至于我，我没有切慕世人的日子，你知道。} \BibleRef{耶 17:16} 因此主也在福音中说：\BibleQuote{若有人愿意跟从我，该当否认自己，天天背着自己的十字架来跟从我。} \BibleRef{路 9:23}\par

% structure: p0010 source=h2 target=subsection reason=relative-heading-rank; base=h2

\subsection{第三章}

\Emph{独修者应寻求何种地域。}\par

因此，那关心自己内心洁净的人，应寻求这样的地域：它们不会因自身的多产和肥沃而引诱他的心思去操劳耕种，也不会驱使他离开他在斗室中固定不移的位置，迫使他到户外去从事某些工作，于是，他的思绪仿佛被公开倾泻，随风飘散了他所有心神的专注和他对目标的锐利视线。这一点，任何再谨慎、再警醒的人都无法防范或察觉，除非他持续地将自己的灵魂和肉体关闭、禁锢在墙壁之内，以便如同一位卓越的渔夫，借着宗徒的技艺为自己寻找食物，他急切而不移动地捕捉在心灵平静深处游动的思想群体，从高处像从岩石上好奇地俯视深水，明智而狡猾地决定：哪些他应用救命的钩子诱捕，哪些他当作坏鱼和臭鱼忽略和抛弃。\par

% structure: p0013 source=h2 target=subsection reason=relative-heading-rank; base=h2

\subsection{第四章}

\Emph{独修者应选择何种工作。}\par

因此，凡恒常持守这种警醒的人，必能有效实现哈巴谷先知清楚表达的话：\Quote{我要站在我的守望所，立在磐石上，窥看祂要对我说什么，以及我该如何回答那指责我的人。} \BibleRef{哈 2:1}然而，这件事有多艰难和烦人，从那些生活在卡拉穆斯或波菲里翁旷野的人的经历中，可清楚看出。因为虽然他们穿越比斯塞特旷野更长的荒漠，远离一切城市与人的住所（穿透广大荒野的深处，行走七八天路程，才勉强到达他们的藏身之所和斗室），但由于他们在那里从事农耕，丝毫不受隐修院的约束，每当我们来到我们所居住的这些简陋地区或斯塞特时，他们便被如此折磨人的思想和如此焦虑的心境所困扰，以致仿佛他们是初学者、从未对独居操练给予丝毫注意的人，不能忍受斗室的生活及其平安与宁静，立即被驱赶出去，不得不离开，如同缺乏经验和初学的人。因为他们没有学会止息内在之人的动荡，并以焦虑的关怀和不懈的努力平息思想的狂风暴雨，反而日复一日在户外劳作，终日不仅在肉体上，也在心神上空荡荡地游荡；随着身体四处移动，他们的思想也随处倾泻。因此，他们察觉不到自己心灵在渴求许多事物时的愚昧，也无法抑制其漫无目的的散漫；由于不能忍受精神的忧伤，他们以为持续静默是无法忍受的，那些在田野辛劳从不疲倦的人，竟被静默打败，被长久的休息耗尽。\par

% structure: p0016 source=h2 target=subsection reason=relative-heading-rank; base=h2

\subsection{第五章}

\Emph{心思的焦虑不仅不会因身体的躁动而减轻，反而会加剧。}\par

如果一个住在斗室中的人，他的思想仿佛被聚集在狭窄的隐修院内，却受到众多焦虑的压迫，这些焦虑随着人本身从居所的束缚中爆发出来，立刻像野马一样四散奔逃，这并不奇怪。然而，当它们暂时从马厩中自由驰骋时，能获得片刻短暂而可悲的慰藉；但在身体回到自己的斗室之后，整个思想队伍又退回到其适当的住所，长期放纵的习惯引发了更剧烈的痛苦。因此，那些无力也不懂得如何与自己幻想的冲动争战的人，当他们被更猛烈的\OriginalTerm{倦怠}{accidie}攻击胸中时，若放松严格的纪律，允许自己更频繁地外出，他们就会因自己误以为的补救方法而激起更严重的灾祸：就像人们以为用最冷的水可以平息内热的猛烈，但实际上，这反而更点燃了火焰，而非熄灭，因为短暂的缓解之后，更糟糕的发作随之而来。\par

% structure: p0019 source=h2 target=subsection reason=relative-heading-rank; base=h2

\subsection{第六章}

\Emph{一个比喻：隐修士应当如何看守自己的思想。}\par

因此，隐修士的全部注意力应固定在一个点上，所有思想的兴起与循环都要严格地限制于此：即对天主的忆念，犹如一个人急于建造一座圆拱的高顶，必须不断地从其精确中心画线，并依据其可靠的标准，按建筑学的规律发现所需的均匀与圆整。但如果有人试图在不确定其中心的情况下完成它——尽管他对自己的技艺和能力极有信心，他不可能使圆周保持均匀而无差错，也无法单凭观察发现他因失误从真正圆形的美中减去了多少，除非他始终求助于真理的试金石，并依其判断修正其工作的内外边缘，从而将高大建筑完成到精确点上。同样，我们的心灵，除非以对主之爱作为不可动摇的固定中心，围绕它运行，在我们所有的工作和谋划的整个环境中，借由爱的出色圆规（如果可以这样说的话）来适配或拒绝我们所有思想的特质，否则绝不可能以卓越的技能建造起那属于保禄这位建筑师的属灵建筑的架构，也无法拥有那有福的达味切愿向主显示并说\Quote{我爱慕祢殿宇的美丽，祢荣耀居住之处}的美丽房屋；反而会轻率地在心中建造一座不美、且不配圣神的房屋，这房屋将立即倒塌，因此不会因接受有福的住主而获得荣耀，反而会因房屋的倒塌而悲惨地毁灭。\par

% structure: p0022 source=h2 target=subsection reason=relative-heading-rank; base=h2

\subsection{第七章}

\Emph{一个问题：为何与亲属为邻被认为妨碍我们，而对生活在埃及的人却无妨碍？}\par

杰尔曼努斯：就斗室内可行的工作种类而言，您所给予的规则十分有用且必要。因为我们不仅从您圣德的榜样（基于对宗徒德行的效法），而且从我们自身的经验，多次证明了其价值。但为什么我们应当如此彻底地避免与亲属为邻，而您并未完全拒绝这一点，却不够清楚。因为我们看到您们无可指责地行走在完美的道路上，不仅有人居住在本乡，甚至有人未远离自己的村庄，那么，对您们无害的事，为何对我们却被视为有害？\par

% structure: p0025 source=h2 target=subsection reason=relative-heading-rank; base=h2

\subsection{第八章}

\Emph{回答：并非所有事都适合所有人。}\par

\Person{亚巴郎}说：\newline{}\Quote{有时我们从善事中看到坏榜样。若有人效法他人做同样的事，却没有同样的心意和目的，或没有同样的善德，便会立刻因那些使他人获得永生果实之事，陷入欺骗和死亡的罗网。正如那位臂力强壮的少年，若穿上只适合成人的\Person{撒乌耳}的重甲，与好战的巨人交锋，必定遭遇危险；而年长者凭此甲可击倒无数敌军，对少年却只会带来确定的危险，除非他明智地选择适合自己年龄的武器，不用别人所用的甲胄和盾牌，而用自己能作战的武器来武装自己，以对抗那可憎的敌人。因此，我们每人首先当仔细衡量自己的力量，并按其限度选择自己所喜欢的修道方式，因为虽然一切都是善的，但并非一切适合众人。我们并非因隐修生活是善的，就断言它适合每一个人，因为对许多人而言，它不仅无益，甚至有害。同样，因我们正确选择了公修团体的生活以及对弟兄们虔诚可赞的照顾，便认为这当被众人遵循。同样，照顾旅客的果实虽极丰盛，但并非人人都能毫无耐心地承担。此外，你们本国和此地的修道方式必须先互相衡量；然后，从众人美德或恶习的持续表现中得出的力量，须分别放在相反的天平上权衡。因为很可能，对一个民族的人而言困难或不可能的事，在另一个民族中却因根深蒂固的习惯而化为本性，正如一些地区相隔遥远的民族，能赤身裸露承受极大的严寒或酷暑，而这在其他未经历那种严酷气候的人，无论多么强壮，也绝不可能忍受。同样，你们在此地以身心极大努力试图在诸多方面克服本国本性的人，也当认真考虑：在那些据报冰封、因极度不信的寒冷而束缚的地区，你们能否忍受这所谓赤身裸体的状态。因为对我们而言，长久持续的圣洁生活几乎自然而然地赋予了我们这种坚毅；若我们看到你们在德性与恒心上与我们相等，你们同样不必回避亲族和弟兄的邻近。}\par

% structure: p0028 source=h2 target=subsection reason=relative-heading-rank; base=h2

\subsection{第九章}

\Emph{凡能效法阿颇罗院长刻苦生活者，不必惧怕亲族邻人。}\par

为使你们能由某种严格的考验公平地衡量自己的力量，我将指出一位长者所作的事，即\Person{阿颇罗}院长。这样，若你们内心的细察断定自己在这人的志向和善德上并不逊色，便可放心留在本国，住在亲族附近而不损害自己的志向或生活方式，并确信亲情的感受或对本地的喜爱都不能干扰这谦卑身分的严格，这严格不仅是你们自己的意愿，也是你们在此地朝圣的需要所强加于你们的。当时，这位长者的亲兄弟在深夜来到他那里，恳求他从隐修院出来片刻，帮忙救一头牛，那牛悲伤地抱怨说陷在附近沼泽的泥中，他一人无法救出。\Person{阿颇罗}院长冷淡地回答说：\Quote{你为何不请我们那位更年轻的弟兄？他离你更近，你经过时本该请他的。}那兄弟以为他忘了早已埋葬的兄弟之死，并且因过度的斋戒与持续的独居而脑筋不清，便说：\Quote{我怎能召唤一位十五年前已死去的人？}\Person{阿颇罗}院长说：\Quote{难道你不知道，我也已向这个世界死了二十年，无法从我这斗室的坟墓中，为你今生的事务提供任何帮助？基督绝不让我为救你的牛而稍微放松我已开始的刻苦志向，祂甚至不允许我为父亲的葬礼有一刻的间断，那本可更妥善、更虔诚地办理的。}所以，你们现在要探查自己内心的隐密，仔细省察，是否也能对亲族保持如此持续的内心严格。当你们发现自己在灵魂的刻苦上与他相似，那时你们便可知道，同样地，亲族和弟兄的邻近不会伤害你们，即当你们虽与他们极亲近，却对他们如同死人，以致既不容他们因你们的援助得益，也不容因对他们的职责而懈怠。\par

% structure: p0031 source=h2 target=subsection reason=relative-heading-rank; base=h2

\subsection{第十章}

\Emph{问题：隐修士由亲族供应所需是否有害？}\par

\Person{杰尔曼努斯}说：\newline{}\Quote{关于此事，你们无疑没有留下任何不确定的余地。因为我们确信，在亲族附近，我们无法保持目前这可怜的装束，也无法每日赤足行走；在那里，我们甚至不能以同样的劳力获取生活必需品，如同在此地，我们实际上必须背负三里的重担取水。因为我们的羞耻和他们的羞耻，绝不会允许我们在他们面前这样做。然而，若由他们供应一切，我们完全免去准备食物的焦虑，专心于读经和祈祷，从而因除去如今使我们分心的劳作，我们能更热切地仅致力于属神的事，这对我们生活计划有何损害呢？}\par

% structure: p0034 source=h2 target=subsection reason=relative-heading-rank; base=h2

\subsection{第十一章}

\Emph{论圣安当对此问题所定的答案。}\par

亚巴郎说：我不会就此事给你我自己的意见，而是给出\Person{真福安当}的意见，他以此驳斥了某位弟兄的懒惰（这位弟兄被你所说的这种冷淡所征服），同时也解开了你问题的结。因为当一个人来到所述的长老面前，说\OriginalTerm{安道}{Anchorite}生活根本不值得钦佩，声称生活在人群中间比在旷野中更能使人实践成全所需的德行时，\Person{真福安当}问他住在哪里；当他说自己住在亲戚附近，并夸耀因他们的供应，他摆脱了每日工作的一切挂虑和焦虑，专务读经和祈祷，毫无心神涣散时，\Person{真福安当}又说：\Quote{请告诉我，好朋友，你是否因他们的忧苦和不幸而忧苦，也同样因他们的好运而喜乐？}他承认自己两者都分担。长老对他说：\Quote{你应当知道，在来世，你也要与那些你在今生与他们共担得失、同忧同乐的人一同受审。}真福安当对此回答仍不满意，便进一步展开讨论说：\Quote{这种生活方式和这种极冷淡的状况，不仅带给你我所说的那种损害（虽然你现在没有感觉到，正如你不知不觉地按照《箴言》所说：\InnerQuote{他们打我，我却未感伤痛；他们讥笑我，我却不自知。}\BibleRef{箴 23:35}或先知所说：\InnerQuote{外邦人耗尽了他的精力，他却不知道。}\BibleRef{欧 7:9}），因为日复一日地它将你的心不断拖向世俗之事，并随命运的变迁而改变；而且它也剥夺了你双手的果实和你自己努力应有的报酬，因为它不允许你靠这些供应生活，也不允许你按真福宗徒的规则亲手为自己挣得每日食粮；当他对厄弗所教会首领们作最后嘱托时，他声称虽然自己忙于宣讲福音的神圣职分，但不仅供应了自己，也供应了那些因职务必要事务而无法工作的人，说：\InnerQuote{你们自己知道，这双手供应了我和同我一起的人的需要。}\BibleRef{宗 20:34}但为了表明他这样做是为我们立榜样，他在别处又说：\InnerQuote{我们在你们中间从未游手好闲，也没有白吃任何人的饭，而是日夜辛苦劳作，免得成为你们任何人的负担。这并不是因为我们没有权利，而是为给你们自己榜样，叫你们效法我们。}\BibleRef{得后 3:8-9}}\par

% structure: p0037 source=h2 target=subsection reason=relative-heading-rank; base=h2

\subsection{\Emph{第十二章}}

\Emph{论工作之价值与懒惰之害。}\par

因此，虽然我们也可以得到亲戚的保护，但我们却宁愿选择他的克己胜过一切财富，选择用自己的劳动来获取每日肉身的食粮，而不是依靠亲戚可靠的供应，因为我们不那么倾向于你所提到的对圣经的闲散默想和无益的阅读，而是倾向于这种艰苦的贫穷。当然，我们非常乐意追求前者，如果宗徒们的权威通过他们的榜样教导我们这对我们更好，或者长老们的规则为我们的益处如此规定的话。但你必须知道，你所受的这种影响并不亚于我上面所说的那种损害，因为虽然你的身体可能健康强壮，但你却靠别人的施舍维持生活，这本应是属于软弱之人的。因为整个人类，除了那类按照宗徒的命令靠双手每日劳动生活的隐修士之外，都仰赖他人的怜悯之爱。因此，显然不仅那些夸耀自己靠亲戚的财富、仆人的劳动或农场的出产生活的人，而且这个世界上的国王，都是靠爱德生活的。至少，这包含在我们的前辈的定义中，他们规定：凡为每日食粮所需而取用之物，若不是靠自己的双手劳动取得和准备的，都应归于爱德，正如宗徒所教导的，他完全禁止懒惰者接受他人的慷慨，说：\BibleQuote{如果一个人不工作，就不该吃饭。}\BibleRef{得后 3:10} 真福安当用这些话驳斥了某人，也通过他教导的榜样教训我们，要躲避我们的亲戚以及所有为我们提供必要爱德食粮之人的有害诱惑，以及舒适家庭的乐趣，而宁愿选择这个世界上一切财富之外的、荒凉的自然贫瘠的沙地和被活生生的硬壳所覆盖的地区，因此不受人的控制和管辖，这样我们不仅为了荒芜的旷野而避免人群，而且肥沃土壤的特性也永远不会引诱我们从事耕种的纷扰，从而将心神从内心的主要侍奉中召回，使之对属灵的目的毫无用处。\par

% structure: p0040 source=h2 target=subsection reason=relative-heading-rank; base=h2

\subsection{\Emph{第十三章}}

\Emph{一个理发师的报酬故事，为认清魔鬼的幻觉而引入。}\par

因为正如你希望也能拯救他人，并渴望怀着更大收益的希望返回本国，也请听关于\Person{玛加略}院长的这则故事，它编得十分巧妙优美，是他在一个人处于类似欲望骚动时，为医治他而讲的一个非常恰当的故事。\Quote{有位，}他说，\Quote{在某座城里有一位非常精明的理发师，他经常为每个人刮脸只收三便士，靠这可怜而微薄的报酬来工作，并从中获取每日食物所需，在照顾身体一切需要之后，每天还能在口袋里存入一百便士。但当他勤勉地积累这笔收益时，他听说在很远的一座城里，每个人付给理发师一先令作为报酬。他发现这事后，\InnerQuote{我还要忍受这种乞讨多久，}他说，\InnerQuote{以便凭我的劳动赚取三便士的报酬，而到那里去我就可以靠大量先令的收益积累财富？}于是他立即带上自己的手艺工具，并耗尽长期积攒和节省下来的所有积蓄作为路费，艰难地前往那座最赚钱的城市。到达的那天，他按照所听说的，从每个人那里得到了劳动报酬，傍晚看到自己赚得了大量先令，便高兴地到肉铺去买晚餐所需的食物。当他开始用一大笔先令购买时，他花掉了所有赚来的先令，只买了一点儿肉，甚至没有带回家一个多余的便士。当他看到自己的收益就这样每天用光，不仅无法储蓄任何东西，甚至连每日食物所需都难以获得时，他心中思考说：\InnerQuote{我要回到我的城市，重新寻求那些非常适度的利润，从那里，在满足身体一切需要之后，每日的盈余会为我的老年积累增长的数目；这虽然看似微小而琐碎，但通过不断积累，竟达到了不小的数目。事实上，铜币的收益比这名义上的先令收益对我更有益，因为从后者不仅没有多余的可以储蓄，而且连日常食物的必需都难以满足。}}因此，对我们来说，最好以不间断的坚持，追求隐修中这份非常微薄的收益——世俗的忧虑、世间的搅扰、虚妄的骄傲与虚荣都不能减少它，日常需要的压力也不能削减它（因为\Quote{义人拥有的微小之物，胜过恶人的大量财富} ）而不是去追求那些更大的收益，即使它们是通过许多人最有价值的悔改而获得的，却仍被世俗生活的要求和日常分心的渗漏所耗尽。因为，正如撒罗满所说：\BibleQuote{一把安逸，胜过两把劳碌和心烦。}\BibleRef{训 4:6}在所有这些暗示和不便中，所有软弱的人必然会被缠住，因为他们甚至对自己的救恩都存疑，自己还需要他人的教导和指引，却被魔鬼的诡计煽动去使他人悔改和引导他人；而且即使他们从某些人的悔改中获得了任何益处，也会因自己的不耐烦和粗鲁行为而浪费所获得的一切。因为先知哈盖所描述的事必会发生在他们身上：\BibleQuote{那积聚财富的，将之放入有洞的袋子。}\BibleRef{盖 1:6}因为一个人若因缺乏自制和日常的心神散漫，失去他看似从他人悔改中所获得的一切，那他就是将自己的收益放入有洞的袋子。结果，当他们幻想通过教导他人可以获得更大收益时，实际上却被剥夺了自己的进步。因为\Quote{有人自充富人，其实一无所有；有人自甘卑微，却拥有大量财富；}以及：\Quote{一个在卑微地位中服侍自己的人，胜过那为自己获取尊荣却缺少面包的人。}\par

% structure: p0043 source=h2 target=subsection reason=relative-heading-rank; base=h2

\subsection{第十四章}

\Emph{问题是，这种错误观念是如何潜入我们心中的。}\par

\Person{杰尔玛诺}：你的讨论通过这个比喻非常恰当地指出了这些幻觉的错误；但我们同样希望被告知其根源和治愈方法，并且同样渴望知道这种欺骗是如何占据我们的。因为每个人都必须明白，除非已经诊断出疾病的真正根源，否则没有人能够对健康不良施以治疗。\par

% structure: p0046 source=h2 target=subsection reason=relative-heading-rank; base=h2

\subsection{第十五章}

\Emph{关于灵魂三重运动的答复。}\par

亚巴郎：一切恶习只有一个根源和来源，但情欲和败坏的名称却各不相同，根据灵魂中受到伤害的那部分（若可如此称之）或肢体的特征而定：正如有时也由身体疾病的情况所示，虽然病因同一种，却因受影响的肢体性质不同而分为不同的病种。因为当有害体液的力量占领身体的堡垒，即头部时，它引起头痛的感觉；但当它影响耳朵或眼睛时，就变成耳痛或眼炎的疾病；当它扩散到关节和手部末端时，称为关节痛风或手痛风；但当它下到脚部末端时，名称改变，称为足痛风：而原初同一种有害体液，随着它所影响的不同肢体而有同样多的名称。同样，从可见之物转到不可见之物，我们当认为，每一种恶习的倾向存在于我们灵魂的各个部分（若可如此称之）和肢体中。并且，正如一些极有智慧的人所设定的，灵魂的力量有三种：要么是\OriginalTerm{理性的}{\GreekText{λογικόν}}，即合理的；要么是\OriginalTerm{易怒的}{\GreekText{θυμικόν}}；要么是\OriginalTerm{可欲的}{\GreekText{ἐπιθυμητικόν}}，即受欲望支配的，它必定会受到某种攻击的困扰。于是，当有害情欲的力量因这些情感而占据一个人时，恶习的名称便根据受影响的部分而赋予。因为如果罪恶的瘟疫侵入了它的理性部分，就会产生虚荣、自负、嫉妒、骄傲、妄断、纷争、异端等罪。如果它伤害了易怒的情感，就会生出愤怒、急躁、愠怒、懈怠、怯懦和残忍。如果它影响了受欲望支配的部分，就会成为贪食、邪淫、贪婪、悭吝，以及有害和世俗的欲望之母。\par

% structure: p0049 source=h2 target=subsection reason=relative-heading-rank; base=h2

\subsection{第十六章}

\Emph{我们灵魂的理性部分已败坏。}\par

因此，若你想发现这恶习的根源和来源，你必须承认你理智和灵魂的理性部分已败坏，即大部分滋生妄断和虚荣恶习的那部分。此外，你灵魂的这第一个肢体（可谓如此）必须通过正确辨别的判断和谦逊的德性来医治；因为当它受伤时，你自以为不仅能攀登完美的顶峰，甚至能教导他人，并认为自己有能力且足以指导别人，因着虚荣的骄傲，你被这些虚妄的游荡所夺走，你的告明揭示了这些。然后，若你如我所说，建立在真正明辨的谦卑中，并心中悲伤地学到，我们每个人拯救自己的灵魂是多么艰难困难的事，且从内心深处承认，你不仅远离那种教导的骄傲，而且实际上仍然需要一位老师的帮助，你便能毫不费力地摆脱这些。\par

% structure: p0052 source=h2 target=subsection reason=relative-heading-rank; base=h2

\subsection{第十七章}

\Emph{灵魂较弱的那部分如何首先向魔鬼的诱惑屈服。}\par

然后，你应将真正谦卑的良药应用于我们描述为特别受伤的那一成员或灵魂部分；因为，就所见而言，它在你灵魂的其他能力中较为软弱，它必定首先屈服于魔鬼的攻击。正如当一些伤害临到我们，无论是因劳苦还是因恶劣的天气所致，通常在人身体中，较弱的部位首先屈服于这些际遇，当疾病特别侵袭它们时，也会以同样的祸患影响身体健康的部位；同样，当罪的毒风向我们吹来，我们每个人的灵魂必定首先被那情欲所诱惑，因为其中较软弱、较薄弱的部位对仇敌强烈的攻击没有做出如此顽强的抵抗，并冒着被那些因疏忽大意而更容易背叛之物俘虏的风险。因为巴兰如此聚集，认为天主的子民肯定能被一种确定的方法欺骗，他建议，在他知道以色列子民软弱之处，为他们设置危险的陷阱，因为他毫不怀疑，当提供妇女时，他们会立刻跌倒，因奸淫而毁灭，因为他知道他们灵魂中易受欲望支配的部分已经败坏。因此，属灵的邪恶以狡猾的恶意试探我们每个人，特别针对他们所见我们灵魂中软弱的情感暗中设下陷阱，例如，若他们看见我们灵魂的理性部分受到影响，他们试图以经上记载阿哈布王被那些叙利亚人欺骗的方式欺骗我们，他们说：\Quote{我们知道以色列的君王是仁慈的：因此让我们把麻布束在腰间，把绳子套在头上，出去见以色列君王，对他说：你的仆人本哈达得说：我求你，让我的灵魂存活。} 由此他并非受真正良善的影响，而是受对他仁慈的空虚称赞的影响，就说：\Quote{他若还活着，他就是我的兄弟；} 这样，他们也能够借着那理性部分的错误欺骗我们，使我们因那本希望从中获得赏报并接受良善回报的事物而招致天主的愤怒，而对我们也可说同样的责备：\Quote{因为你从你手中放走了一个该死的人，你的性命必代替他的性命，你的百姓必代替他的百姓。} 或者当污灵说：\BibleQuote{我要出去，在他众先知口中作谎言之灵，} \BibleRef{列上 22:22}，他当然借着那他知道易受其致命诡计影响的理性情感散布了欺骗的网罗。同样的灵在我们的主身上也期望如此，当祂在灵魂的这三种情感中试探祂时，他知道全人类都被这些俘虏了，但他的诡诈伎俩却一无所获。因为他接近祂心智中易受欲望支配的部分，当他说：\Quote{吩咐这些石头变成饼；} 易怒的部分，当他试图煽动祂寻求今世的生命和今世万国的权柄；理性部分，当他说：\Quote{你若是天主子，就从这里跳下去。} 在这些事上，他的欺骗无效，原因如下：因为他发现祂身上没有任何受损之处，这与他从错误观念所形成的假设相符。因此，当仇敌的诡计试探时，祂的灵魂没有一部分屈服，祂说：\BibleQuote{因为看，这世界的元首来了，在我内一无所获。} \BibleRef{若 14:30}\par

% structure: p0055 source=h2 target=subsection reason=relative-heading-rank; base=h2

\subsection{第十八章}

\Emph{一个疑问：我们是否应按正当的愿望，为更大的静默而被带回祖国。}\par

\Person{格尔玛努斯}说：在我们其他各类幻觉和错误中（这些错误借着对神益的空虚应许，燃起了我们对祖国的渴望，正如您的圣德以敏锐的心智所察觉的），这作为主要原因突显出来：有时候我们被弟兄们缠扰，无法像我们希望的那样继续无间断的静独和持续的静默。因此，我们日常克苦的进程和尺度——我们总想为克制身体而保持不受干扰——在弟兄们到来时肯定会受影响。我们肯定觉得，这在我们祖国永远不会发生，在那里不可能找到任何人，或几乎没有人采纳这种生活方式。\par

% structure: p0058 source=h2 target=subsection reason=relative-heading-rank; base=h2

\subsection{第十九章}

\Emph{回答：关于魔鬼的幻觉，因为他应许我们更广阔静独的平安。}\par

\Person{亚巴郎}：根本不被人们拜访，是缺少理性、考虑不周的严厉，或者说是最冷淡的标志。因为如果一个人以过于缓慢的步伐行走在他所进入的路上，并按照旧人生活，那么不仅没有圣人，而且没有一个人应该拜访他，这是对的。但是，如果你被对我们主的真实而完美的爱所点燃，并以全部的热诚跟随确实是爱的天主，那么你肯定会被人们拜访，无论你逃到多么难以到达的地方，而且，根据神圣的爱之火使你更接近天主的程度，更多的圣洁弟兄们也将聚集到你身边。因为主说：\BibleQuote{建在山上的城是不能隐藏的。}\BibleRef{玛 5:14}因为主说：\BibleQuote{尊重我的，我必尊重他；藐视我的，他必被轻视。}\BibleRef{撒上 2:30}但是你应该知道，这是魔鬼最巧妙的诡计，是他精心隐藏的陷阱，他把一些可怜的、粗心的人推入其中，这样，在他向他们许诺更伟大的东西的同时，却夺走了他们日常进益的必要优势，说服他们去寻找更遥远、更荒凉的旷野，并在他们心中把这些旷野描绘成仿佛布满奇妙的快乐。此外，他还捏造一些未知的、不存在的地方，假装这些地方是众所周知的、合适的，已经交给我们权力，可以毫无困难地获得。他还假称那个国家的人们是温顺的、跟随救恩之路的，以便在他许诺那里为灵魂结出更丰盛的果实时，狡猾地毁坏我们目前的收益。因为当每个人因这虚妄的希望而与长老们的共同生活分离，并被剥夺了他在心中空想的一切时，他仿佛从最深的沉睡中醒来，醒来时将发现他梦想的一切都没有。因此，由于他被今生更大的需要和无法挣脱的罗网所束缚，魔鬼甚至不让他渴望他曾向自己许诺的那些东西，而且，由于他不再受到他先前所避免的弟兄们那种稀少的、属灵的拜访，反而受到世俗之人每日的打扰，他将永远不能回到独修者生活的适中安宁和秩序中。\par

% structure: p0061 source=h2 target=subsection reason=relative-heading-rank; base=h2

\subsection{第二十章}

\Emph{弟兄到来时的放松是何等有益。}\par

那种因弟兄到来而有时会介入的、最令人振奋的放松和好客的小插曲，虽然在我们看来可能令人厌烦，是我们应该避免的，但它对我们的身体和灵魂是多么有用和有益，你必须耐心地听我简略说明。我所说的，不仅对初学者和软弱的人，甚至对经验最丰富、最成全的人，也常常发生：除非他们通过一些变化来放松精神的紧张和压力，否则他们要么陷入精神冷淡，要么至少陷入对身体健康最危险的状态。因此，当甚至频繁的弟兄访问发生时，明智而成全的人不仅应该耐心忍受，而且应该感恩欢迎；首先，因为它们激励我们更热切地渴望旷野的隐退（因为某种方式上，它们被认为阻碍我们的进步，其实却使我们的进步保持不懈、不断裂，如果它从未被任何障碍阻挠，就无法以坚定不移的毅力坚持到底）；其次，因为它们给了我们休息身体的机会，连同友善的好处，同时，身体最愉快的放松，反而给我们带来了比守斋带来的疲劳更大的好处。关于这一点，我将简要地给出一个最恰当的比喻，它出自一个古老的故事。\par

% structure: p0064 source=h2 target=subsection reason=relative-heading-rank; base=h2

\subsection{第二十一章}

\Emph{圣史若望如何展示放松的价值。}\par

据说\Person{真福若望}正用手轻轻地抚弄一只山鹑时，忽然看见一位穿着猎人服装的哲学家走近，他对这样一位大名鼎鼎、享誉盛名的人竟屈尊从事如此微不足道、毫无意义的消遣感到惊讶，说道：\Quote{你能是那位若望吗？你伟大而著名的声誉也使我极其渴望认识你。为何你竟从事如此低劣的消遣？} \Person{真福若望}对他说：\Quote{那是什么？}他说，\Quote{你手里拿着什么？} 那人答道：\Quote{一把弓。} \Quote{为什么？}他说，\Quote{你不总是把它弯曲着随身带吗？} 那人答道：\Quote{不行，因为如果一直弯曲，它的张力就会松弛，会减弱甚至毁掉，等到需要用更强劲的箭射向野兽时，它的劲道会因过度持续的绷紧而丧失，就不可能射出更有力的箭了。} \Quote{我的孩子，}\Person{真福若望}说，\Quote{不要让我内心这一短暂轻微的放松打扰你，因为除非它有时通过一些消遣来缓解和放松其意志的严格，否则精神会因持续的紧绷而失去弹性，当需要时，就无法坚定地追随正确的事。}\par

% structure: p0067 source=h2 target=subsection reason=relative-heading-rank; base=h2

\subsection{第二十二章}

\Emph{一个问题：我们应如何理解福音所说的\BibleQuote{我的轭是容易的，我的担子是轻省的。}}\par

\Person{日耳曼努斯}说：既然您给了我们医治一切妄想的药方，并且借天主的恩宠，您教导所揭露的、困扰我们的魔鬼的一切诡计，我们恳求您也为我们解释福音中所说的这句话：\BibleQuote{我的轭是容易的，我的担子是轻省的。}\BibleRef{玛 11:30}因为这似乎相当违背先知所说的那句话：\Quote{因了祢嘴唇的言语，我遵守了艰难的道路。}而且宗徒也说：\Quote{凡愿意在基督内虔敬生活的，都受迫害。}但是，凡是艰难且充满迫害的，就不能是容易和轻省的。\par

% structure: p0070 source=h2 target=subsection reason=relative-heading-rank; base=h2

\subsection{第二十三章}

\Emph{回答与该说法的解释。}\par

亚巴郎：我们可用自己经验的简易教训证明，主和救主的这句话是千真万确的，只要我们按基督的旨意，正当地走上成全的道路，克制我们的一切私欲，断绝有害的偏爱，不仅不让世上任何财物留在我们身边（否则我们的仇敌就能随意找到毁灭和伤害我们的机会），并且真正认识到我们并非自己的主人，诚心实意地把宗徒的话当作自己的话：\Quote{我生活，不是我生活，而是基督在我内生活。}\BibleRef{迦 2:20} 因为，对于全心拥抱基督的轭、建立在真谦逊中、常注视主的苦难、为基督所受的一切冤屈而喜乐的人，有什么是沉重的或艰难的呢？他说：\Quote{为此，我为基督的缘故，喜欢在软弱中、在凌辱中、在急需中、在迫害中、在困苦中，因为我几时软弱，正是我有能力的时候。}\BibleRef{格后 12:10} 请问，那夸耀完全弃绝、为基督自愿弃绝世上一切荣华、视世间所有欲望如同粪土、为赢得基督而不断默想这句福音训言、蔑视并摆脱因任何损失而起的激动的人，还有什么常物之失能伤害他呢？\Quote{人纵然赚得了全世界，却赔上了自己的灵魂，为他有什么益处？或者人还能拿什么作为自己灵魂的代价？}\BibleRef{玛 16:26} 至于那认出一切能从别人手中夺取的都不是自己的、并以不屈的勇气宣告\Quote{我们没有带什么到世界上来，也不能带什么去}\BibleRef{弟前 6:7}的人，他会因什么损失而烦恼呢？那知道如何不带\Quote{路上行囊和钱袋}\BibleRef{玛 10:9}、并像宗徒那样夸耀\Quote{多次禁食、忍饥受渴、忍受寒冷和赤身露体}\BibleRef{格后 11:27}的人，什么缺乏的需要能挫败他的勇气？那没有自己意志、不但忍耐而且感恩地接受命令、并效法救主寻求不成就自己的意愿而成就父的旨意的人，如主自己向父所说的：\Quote{但不要照我意，而要照祢意。}\BibleRef{玛 26:39}，什么辛劳或长老的严命能扰乱他内心的平安？再者，何等冤屈、何等迫害能吓倒他？相反，有什么刑罚不能使他感到愉快？因为他常与宗徒一同在鞭打中喜乐，并渴求因基督的名受辱。\par

% structure: p0073 source=h2 target=subsection reason=relative-heading-rank; base=h2

\subsection{第二十四章}

\Emph{为什么主的轭感觉沉重，祂的担子觉得沉重。}\par

但事实上，对我们而言，基督的轭似乎既不轻也不易，这必须正确地归咎于我们的乖僻，因为我们被不信和缺乏信德所压倒，并愚蠢地顽固抵抗祂的命令，或更准确地说，劝诫，祂说：\BibleQuote{你若愿意成为成全的，去变卖（或舍弃）你所有的一切，然后来跟随我，}\BibleRef{玛 19:21} 因为我们保留了世俗财物的实质。并且，当魔鬼用这些财物的罗网紧紧捆绑我们的灵魂，那么，当他想要切断我们与神慰的联系时，他岂不会通过减少和剥夺这些来折磨我们，并以其诡诈的伎俩谋算，使祂的轭之甘饴与祂担子之轻松，因我们腐朽欲望的邪恶而变得令我们痛苦，当我们被那同一财产与实质的锁链所缠住——我们曾为安慰和慰藉保留这些——他就可以用世俗忧虑的鞭笞不断折磨我们，从我们自身榨出那折磨我们的东西？因为\Quote{各人被自己的罪恶绳索所捆绑，}并从先知那里听到：\Quote{看哪，你们所有点火、被火焰包围的人，要在你们之火的光中行走，并在你们所点燃的火焰中行走。}既然，如撒罗满所见证的，\Quote{各人将因他所犯的罪而受罚。}因为我们享受的快乐本身变成了折磨，而这肉身的欢愉和享乐，如同刽子手般转向其始作俑者，因为一个依靠先前财富和财产的人肯定不能接受心灵的完全谦卑，也不能完全克制危险的快感。但是，在所有这些美善工具发挥作用的地方，今生的所有考验，以及仇敌所能设计的一切损失，不仅以极大的忍耐承受，而且以真正的愉快承受；然而，当这些美善工具缺乏时，如此危险的骄傲便滋生，以至于我们因最轻微的责备就受到不耐烦的致命打击，并且耶肋米亚先知可能会对我们说：\BibleQuote{现在你在埃及的路上做什么，要喝混浊之水呢？你在亚述的路上做什么，要喝大河之水呢？你自己的邪恶必责备你，你的背教必谴责你。你要知道并看见，你离弃上主你的天主，是一件邪恶和痛苦的事，并且我的敬畏不在你内，上主说。}\BibleRef{耶 2:18}既然如此，主的轭之奇妙甘饴为何被尝为苦涩，不正是因为我们嫌恶的苦涩损害了它吗？天主的担子何以变得极其沉重，不正是因为我们固执的傲慢藐视了那背负它的一位，尤其是圣经本身明确地为这件事作证说：\BibleQuote{因为如果他们行走在正途，他们必定会发现正义的道路是平坦的}\BibleRef{箴 2:20}？我说，显然正是我们，用我们欲望的龌龊硬石使主正直平坦的道路变得崎岖；我们最愚蠢地放弃了那条铺满宗徒与先知之燧石、被诸位圣人及主自己脚步所踏平的王道，而寻找无路且有荆棘之处，并被眼前欢愉的诱惑弄瞎了眼睛，带着撕裂的腿和破损的婚袍，穿越长满罪恶荆棘的黑暗路径，以致不仅被荆棘的锐刺刺穿，而且实际上被潜伏在那里的致命蛇蝎之咬伤所击倒。因为\BibleQuote{邪道上有荆棘和蒺藜，但敬畏上主的人必远离它们。}\BibleRef{箴 22:5}关于这样的人，主也藉先知在别处说：\BibleQuote{我的百姓忘记了，空献祭，并在他们的路上绊倒，在古道上，行在一条未踏过的路上。}\BibleRef{耶 18:15}因为按照撒罗满的格言：\BibleQuote{不劳动之人的道路布满荆棘，但强健之人的道路被踏平。}\BibleRef{箴 15:19}因此，当他们偏离王道，他们就永远无法到达那座大城——我们的航向应该永不偏离地指向那里。\BibleBook{}{训道篇}也相当有意义地表达了这一点，说：\Quote{愚人的劳苦使那些不知道如何去城的人疲惫；}即，\Quote{天上的耶路撒冷，她是众人的母亲。}但是，凡真正舍弃此世、背起基督的轭并向祂学习，且在每日受屈的练习中被训练的人——因为祂是\BibleQuote{良善心谦的，}\BibleRef{玛 11:29}——将永远不受一切诱惑的搅扰，并且\BibleQuote{万事互相效力，使他得益处。}\BibleRef{罗 8:28}因为正如亚北底亚先知所说，天主的话是\Quote{对行为正直的人是美善的；}并且又说：\Quote{因为上主的道路是正直的，义人必行在其中，但悖逆的人必在其中跌倒。}\par

% structure: p0076 source=h2 target=subsection reason=relative-heading-rank; base=h2

\subsection{第二十五章}

\Emph{论诱惑的袭击所带来的益处}\par

正因如此，藉着与诱惑的搏斗，救主的仁慈恩宠赐予我们更大的赞美赏报，远胜过它完全除去我们争战的必要。因为被逼迫和考验包围时仍能坚定不移，忠信勇敢地站在天主的壁垒上，并在人的攻击中，仿佛佩带着不可征服之德性的武器，光荣地战胜不耐烦，并从软弱中获取力量，这标志着更高尚更伟大的德行，因为\Quote{软弱中见成全}。\BibleQuote{因为看，我使你成了铁柱和铜墙，在全地之上，面对犹大君王、首领、司祭和境内的人民，——主说。他们要攻击你，却不能得胜，因为我与你同在，拯救你——主说。}\BibleRef{耶 1:18}因此，按照主的明训，君王的大道平易顺遂，尽管可能感觉艰难坎坷：因为那些虔诚忠信地事奉他的人，当他们负起主的轭，并从他学会了他是良善心谦的，立刻以某种方式卸下世俗情欲的重担，并发现不是劳苦，而是灵魂的安息，藉着主的恩赐，正如他亲口通过耶肋米亚先知作证说：\Quote{你们应站在路上，查看，询问古旧的道路，哪一条是好的道路，就走在上面；这样，你们必会找到你们灵魂的安息。}因为对他们而言，立刻\Quote{弯曲的将成为正直，崎岖的将成为平坦}；他们必\Quote{尝尝并观看主是何等甘饴}，并且当他们听到基督在福音中宣告：\BibleQuote{凡劳苦和负重担的，你们都到我这里来，我要使你们安息}，他们便将卸下罪恶的重担，并意识到随后的话：\BibleQuote{因为我的轭是良善的，我的担子是轻省的。}\BibleRef{玛 11:28}因此，如果遵守他的法律，主的道路就有安息。而是我们因烦扰的分心，把忧伤和困苦带给自己，当我们甚至竭尽全力艰难地追随这世界弯曲悖谬的道路。但当我们这样使主的轭对我们变得沉重艰难时，我们立刻以亵渎的精神抱怨这轭本身的坚硬和崎岖，或者抱怨那将轭加给我们的基督，正与这段经文相符：\Quote{人的愚昧败坏他的道路，他的心却抱怨天主。}并且正如哈盖先知所说，当我们说\Quote{主的道路不公平}时，主恰当地回答我们：\Quote{我的道路岂不公平？你们的行为岂不是邪恶？}的确，如果你将童贞的馨香之花与贞洁的温柔纯净，与情欲的污秽恶臭的泥沼作比较；将隐修士的宁静安稳，与世俗之人所陷入的危险和损失作比较；将我们神贫的平安，与财富的啃噬烦恼和焦虑挂虑（他们日夜在其中消耗，且不无生命危险）作比较；那么你就会证明基督的轭是何等轻省，祂的担子是是何等轻快。\par

% structure: p0079 source=h2 target=subsection reason=relative-heading-rank; base=h2

\subsection{第二十六章}

\Emph{论完美弃绝者如何获赐今世百倍赏报}\par

此外，那赏报的酬劳，主许诺给那些彻底舍弃一切的人在今世获得百倍，并说：\BibleQuote{凡为我的名，舍弃了房屋、或兄弟、或姊妹、或父亲、或母亲、或妻子、或儿女、或田地的，在现今要得到百倍，并要承受永生。}\BibleRef{玛 19:29} 这一许诺的确可以正确且真实地按同一意义理解，而不致动摇信德。因为许多人藉此言论，固执地以粗俗的智慧认为这些将在千禧年中按肉体实现，尽管他们必须承认，他们所说的复活后的时代不可能被理解为\Quote{现今}。因此，更可信且更清楚的是：一个人因基督的劝勉而轻视了任何世俗情感或财物，就会从与他因属神纽带相连的弟兄和生命伴侣那里，甚至在今生就得到百倍更好的爱；因为可以肯定，在父母、儿女、兄弟、妻子和亲属之间，若纽带仅由交往形成，或联合的纽带由亲属关系的要求所系，那么这种爱是相当短暂且易破碎的。最后，即使是善良孝顺的儿女，长大成人后有时也会被父母逐出家园和产业，有时婚姻的纽带会因一个真正合理的理由而断裂，兄弟之间的争吵分裂也会摧毁财产。唯独隐修士在亲密关系中保持持久的合一，并拥有一切公有，因为他们认为凡属于弟兄的，都是自己的；凡属于自己的，都是弟兄的。因此，如果把我们的爱之恩宠与那些以肉体之爱为联合纽带的感情相比，无疑它要甘美和美好百倍。确实，从节制的婚姻生活中获得的愉悦，也要比从两性结合中产生的愉悦大百倍。而且，一个人若拥有继承天主儿子的名分，并像那位真子那样在心里和灵魂里宣称：\BibleQuote{凡父所有的，都是我的；}\BibleRef{若 16:15} 且不再被分散注意和忧虑的罪恶烦恼所缠扰，而是无忧无虑、心中喜乐地处处继承自己的产业，每日听见宗徒向他宣告：\Quote{一切都是你们的，无论是世界、是现世、是来世；}以及撒罗满所说：\Quote{忠信的人拥有整个世界的财富。}那么，他原本因拥有一块田地或一所房屋而体验到的快乐，将被百倍更大的财富之乐所取代。因此，你已获得那百倍的酬报，由价值的伟大和性质的无比差异所产生。因为如果用固定重量的黄铜、铁或某种更普通的金属，交换同样重量的黄金，那么他给予的远不止百倍。同样，当以属神的喜乐和最宝贵之爱的欢乐来酬报对世俗享乐和情感的蔑视时，即使实际数量相同，也仍是百倍更好、更伟大。为了通过反复强调使这一点更清楚：我从前在欲望的邪情中有一个妻子，如今我在光荣的圣化和基督的真实爱中有一个妻子。女人是同一位，但爱的价值增长了百倍。但如果你以持续的温柔和忍耐取代不信任的愤怒和暴躁，以平安和无忧取代焦虑和烦恼的紧张，以有益的忧伤之苦取代世间无益且有罪的烦恼，以属神之乐的丰盈取代短暂之乐的虚荣，你就会在这些情感的转变中看到百倍的酬报。如果我们将各罪的短暂易逝之乐与相反德行的益处相比，那增长的喜乐证明这些是百倍更好的。因为当你屈指计算时，你将百的数字从左手移到右手，虽然手指的排列似乎相同，但数量的价值却大大增加。结果将是，我们这些似乎具有左边山羊形象的人，将被移开并获得右边绵羊的赏报。现在让我们继续思考基督因我们轻视世俗利益而在今世归还给我们的那些事物的性质，特别是根据马尔谷福音，他说：\BibleQuote{没有人为了我和福音的缘故，舍弃了房屋、或兄弟、或姊妹、或母亲、或儿女、或田地，而不在现今得到百倍的：房屋、兄弟、姊妹、母亲、儿女、田地，并受迫害，且在来世获得永生。}\BibleRef{谷 10:29} 因为谁为了基督的名而舍弃对一位父亲或母亲或儿女的爱，全心投入对一切事奉基督者最纯净的爱，他必将得到百倍的兄弟和亲属；因为代替一个，他将会拥有那么多因更热忱、更可钦佩的情感而与他相连的父亲和兄弟。谁因爱基督而放弃一栋房屋，也必因拥有无数隐修院中的家而富裕，无论他退隐到世界何处，都如同自己的家。他如何能不获得百倍？如果对我们的主的话稍加补充也不算错的话，他舍弃了十个或二十个奴隶的不信和强迫服役，却赢得如此众多高贵且生来自由的人的自愿侍奉，这岂止是百倍？而事实如此，你们可以用自己的经验证明：既然你们每个人都只离开了一位父亲、一位母亲和一个家，却毫不费力、无忧无虑地在世界任何你们所到之处，获得了无数的父亲、母亲和兄弟，以及房屋、田地和最忠实的仆人，他们待你们如同主人，热忱地欢迎、尊敬和照顾你们。但是，我说，圣人们应当配得且自信地享受这种服侍，如果他们首先将自己和自己所有的一切自愿奉献给弟兄们的服侍。因为，如主所说，他们将白白收回他们自己给予别人的。但若一个人没有首先以真正的谦卑将此奉献给同伴，他怎能平静地忍受别人向他奉献呢？因为他知道，这些服侍与其说帮助他，不如说成为他的负担，因为他宁愿接受弟兄们的服侍，而不愿服侍他们。\par

但是，这一切他领受时，并非带着漫不经心的怠惰和懒散的喜悦，而是按照主的话，\Quote{也带着迫害，}即此世的压力和来自自身情欲的可怕痛苦，因为正如智者所证：\Quote{行事散漫、无忧无虑的人，必将陷于匮乏。}因为不是懒惰、粗心、柔弱或娇嫩的人强夺天国，而是强悍的人。那么强悍的人是谁呢？他们正是那些不是向他人，而是向自己的灵魂展现卓越暴力的人，他们以可赞许的力量，剥夺灵魂对现世一切享乐的贪恋，并被主的口宣布为卓越的掠夺者，凭借这般的强夺，暴力地夺取天国。因为，正如主所说：\BibleQuote{天国遭受暴力，强悍的人以暴力夺取它。}\BibleRef{玛 11:12}那些向自己的毁灭施暴的人，无疑可赞为强悍，因为\Quote{一个在忧苦中的人，}正如经上记载，\Quote{为自己劳作，并向自己的毁灭施暴。}因为我们的毁灭在于对今世生命的享乐，更确切地说，在于实现自己的喜好和欲望。如果一个人将这些从灵魂中除去并予以克制，他便立刻对自己的毁灭施加了光荣而宝贵的暴力，前提是他拒绝满足灵魂最愉悦的意愿——而天主的话常常借先知斥责这些意愿，说：\Quote{因为在你禁食的日子，仍寻自己的意愿；}又说：\Quote{若你转离安息日的脚步，在我的圣日不行你的意愿，并荣耀祂，不循你自己的道路，不寻你自己的意愿，不说一句话。}先知随即加上了应许给他的极大福分。他说：\Quote{那时，你将在主内喜乐，我将使你凌驾于地上高处之上，并以你祖先雅各伯的产业喂养你。因为主的口已说了。}因此，我们的主和救主，为给我们树立放弃自己意愿的榜样，说：\Quote{我来不是为行我的意愿，而是为行那派遣我者的意愿；}又说：\Quote{不要照我的意愿，而要照你的意愿。}而这一美德，那些生活在隐修院内、受长老规章管理的人尤为彰显，他们不行自己的选择，而是将自己的意愿服于院长的意愿。最后，为结束这场讨论，我问你们：那些忠信地事奉基督的人，难道不正是最清楚地在此世领受百倍的恩宠吗？他们为了祂的圣名而受到最伟大君王的尊敬，尽管他们并不寻求人的称赞，却在迫害的考验中受人敬仰——而他们卑微的身份（或因出身低微，或因身为奴隶）若仍在世间生活，或许连平民都会轻视。但由于事奉基督，无人敢对他们的高贵身份加以诽谤，或当面指责他们出身的卑微。相反，正是通过那令他人羞愧和窘迫的卑微身份的羞辱，基督的仆人们反而更光荣地被尊崇。这一点，我们可以用住在靠近 \Place{吕科斯}城的旷野中的 \Person{若望}院长的实例清楚说明。他出身卑微，但因基督的圣名，几乎全人类都认识了他，以致那些创造之主——那些掌握着此世和帝国缰绳、令所有权力和君王畏惧的人——都视他为主人，从遥远的地方寻求他的建议，并将他们帝国的冠冕、安全的状况和战争的命运托付于他的祈祷和功德。\par

可敬的 \Person{亚巴郎}用这样的言辞论述了我们幻觉的根源与疗方，并将魔鬼所发起和暗示的狡猾思想展现在我们眼前，激起了我们对真正克修生活的渴望——我们期望这渴望也能点燃许多人，即使这一切是用较为朴素的文体写成的。因为尽管我们言词中的余烬掩盖了诸位伟大教父的炽热思想，但我们希望，对于那些试图拨开我们言词的余烬、将隐藏的思绪煽成火焰的许多人，他们的冷漠将会化为热忱。然而，圣善的弟兄们，我并非因僭越之灵而自高自大，以致把这火（就是主来投到世上的，并急切盼望它燃烧起来的火 \BibleRef{路 12:49}）交给你们，想借这炽热来点燃你们已经白热化的心志；而是为了使你们在子女面前更具权威，如果加之最伟大、最古老的教父们的训诫——不是通过言辞的死板声音，而是通过你们活生生的榜样——支持你们所教导的。剩下的只是：我这迄今被最危险的风暴所抛掷的人，愿藉你们祈祷的属灵之风，被吹入静默的安全港湾。\par

\SourceRecord{Translated by C.S. Gibson.；Nicene and Post-Nicene Fathers, Second Series；Vol. 11.；Edited by Philip Schaff and Henry Wace.；Buffalo, NY: Christian Literature Publishing Co.,；1894.；<http://www.newadvent.org/fathers/350824.htm>.}
